%% Generated by Sphinx.
\def\sphinxdocclass{report}
\documentclass[letterpaper,10pt,english]{sphinxmanual}
\ifdefined\pdfpxdimen
   \let\sphinxpxdimen\pdfpxdimen\else\newdimen\sphinxpxdimen
\fi \sphinxpxdimen=.75bp\relax
\ifdefined\pdfimageresolution
    \pdfimageresolution= \numexpr \dimexpr1in\relax/\sphinxpxdimen\relax
\fi
%% let collapsible pdf bookmarks panel have high depth per default
\PassOptionsToPackage{bookmarksdepth=5}{hyperref}

\PassOptionsToPackage{booktabs}{sphinx}
\PassOptionsToPackage{colorrows}{sphinx}

\PassOptionsToPackage{warn}{textcomp}
\usepackage[utf8]{inputenc}
\ifdefined\DeclareUnicodeCharacter
% support both utf8 and utf8x syntaxes
  \ifdefined\DeclareUnicodeCharacterAsOptional
    \def\sphinxDUC#1{\DeclareUnicodeCharacter{"#1}}
  \else
    \let\sphinxDUC\DeclareUnicodeCharacter
  \fi
  \sphinxDUC{00A0}{\nobreakspace}
  \sphinxDUC{2500}{\sphinxunichar{2500}}
  \sphinxDUC{2502}{\sphinxunichar{2502}}
  \sphinxDUC{2514}{\sphinxunichar{2514}}
  \sphinxDUC{251C}{\sphinxunichar{251C}}
  \sphinxDUC{2572}{\textbackslash}
\fi
\usepackage{cmap}
\usepackage[T1]{fontenc}
\usepackage{amsmath,amssymb,amstext}
\usepackage{babel}



\usepackage{tgtermes}
\usepackage{tgheros}
\renewcommand{\ttdefault}{txtt}



\usepackage[Bjarne]{fncychap}
\usepackage{sphinx}

\fvset{fontsize=auto}
\usepackage{geometry}


% Include hyperref last.
\usepackage{hyperref}
% Fix anchor placement for figures with captions.
\usepackage{hypcap}% it must be loaded after hyperref.
% Set up styles of URL: it should be placed after hyperref.
\urlstyle{same}

\addto\captionsenglish{\renewcommand{\contentsname}{Contents:}}

\usepackage{sphinxmessages}
\setcounter{tocdepth}{1}



\title{Quantum Navigation}
\date{Mar 25, 2025}
\release{1.1.0}
\author{University of Liverpool}
\newcommand{\sphinxlogo}{\vbox{}}
\renewcommand{\releasename}{Release}
\makeindex
\begin{document}

\ifdefined\shorthandoff
  \ifnum\catcode`\=\string=\active\shorthandoff{=}\fi
  \ifnum\catcode`\"=\active\shorthandoff{"}\fi
\fi

\pagestyle{empty}
\sphinxmaketitle
\pagestyle{plain}
\sphinxtableofcontents
\pagestyle{normal}
\phantomsection\label{\detokenize{index::doc}}


\noindent{\hspace*{\fill}\sphinxincludegraphics[scale=0.12]{{uol_logo}.png}\hspace*{\fill}}


\chapter{About}
\label{\detokenize{index:about}}
\sphinxAtStartPar
A Python based navigation toolbox designed to be used for experimenting
with multi\sphinxhyphen{}source navigation. This library provides the necessary services
to conduct in\sphinxhyphen{}depth realistic simulations, fully configurable by the user.
These can extend from employing simple numerical methods to track movement
along a basic pre\sphinxhyphen{}generated racetrack, to using kalman filters with
tightly\sphinxhyphen{}coupled GPS fusion to track a custom vehicle along a bespoke
trajectory. The implementation and object association is easy to
understand an supports third\sphinxhyphen{}party additions. In summary, the toolbox
presents the following key features:
\begin{itemize}
\item {} 
\sphinxAtStartPar
Trajectory generation with vehicle constraints

\item {} 
\sphinxAtStartPar
Configurable accelerometers, gyroscopes and altimeters

\item {} 
\sphinxAtStartPar
External clock allowing for controlled timing errors

\item {} 
\sphinxAtStartPar
Selection of Inertial Navigation (INS) Fusion Methods

\item {} 
\sphinxAtStartPar
GPS satellite simulation with pseudo range generation

\item {} 
\sphinxAtStartPar
GPS fusion (fixed gain, loosely\sphinxhyphen{}coupled, tightly\sphinxhyphen{}coupled)

\item {} 
\sphinxAtStartPar
Inertial based quantum sensors and fusion methods

\item {} 
\sphinxAtStartPar
Gravity gradient based quantum sensors and fusion methods

\item {} 
\sphinxAtStartPar
Generating and plotting interactive results in web browser

\item {} 
\sphinxAtStartPar
Handling of various input and output file types for data

\item {} 
\sphinxAtStartPar
Support for platforms with limited hardware resources

\end{itemize}

\sphinxstepscope


\section{Project Structure}
\label{\detokenize{usage/structure:project-structure}}\label{\detokenize{usage/structure::doc}}
\sphinxAtStartPar
This project contains a handful of files, directories and sub\sphinxhyphen{}directories; each with their own purpose.
This section outlines the role of all project folders and files. Each are grouped and outlined below:


\subsection{Project Folders}
\label{\detokenize{usage/structure:project-folders}}\begin{description}
\sphinxlineitem{/src/}
\sphinxAtStartPar
All source code files that the toolbox are composed of. This directory is identified as a “Python Directory”
and houses all the developed code. This folder contains the project \sphinxcode{\sphinxupquote{qnav}} and its series of sub\sphinxhyphen{}packages.

\sphinxlineitem{/test/}
\sphinxAtStartPar
All pytest test functions and classes. This can be executed to re\sphinxhyphen{}perform all automated tests used to evaluate
the correctness and robustness of the project.

\sphinxlineitem{/docs/}
\sphinxAtStartPar
A collection of scripts and configuration files used for auto\sphinxhyphen{}generating the user documentation
(in HTML offline webpages, LaTeX and PDF formats). This folder mostly consists of ReStructuredText (.rst)
files which describe the information to be included within the documentation.

\sphinxlineitem{/config/}
\sphinxAtStartPar
A series of sub\sphinxhyphen{}configuration profiles that can be imported by main configuration profiles (such as
\sphinxcode{\sphinxupquote{default\_settings.ini}}). The purpose for these is to reduce the complexity of the main configuration script
by partitioning sets of frequently used variables into sub\sphinxhyphen{}profiles. This folder contains sub\sphinxhyphen{}profiles for various
aspects such as vehicle types, hardware models and waypoint trajectories.

\sphinxlineitem{/examples/}
\sphinxAtStartPar
A collection of helpful example scripts providing useful sample code that can be executed to demonstrate the usage of
various features contained within the toolbox. These examples including: generating trajectories,
loading and altering configuration files and using various databases.

\sphinxlineitem{/output/}
\sphinxAtStartPar
The default output directory produced by the toolbox. Here output results (data and report figures) are saved
to auto\sphinxhyphen{}generated folders named by their time and date of creation. Check this folder after running a simulation
to see the results produced.

\end{description}

\begin{DUlineblock}{0em}
\item[] 
\end{DUlineblock}


\subsection{Project Files}
\label{\detokenize{usage/structure:project-files}}\begin{description}
\sphinxlineitem{default\_settings.ini}
\sphinxAtStartPar
The default configuration settings automatically used by the project if a configuration file of choice is not
presented. When making custom configuration files it is recommended to create a copy of \sphinxcode{\sphinxupquote{default\_settings.ini}}
to use as a fully commented template.

\sphinxlineitem{README.md}
\sphinxAtStartPar
A short description about the purpose of the project and how to use it. Only minor details are covered to provide
a quick\sphinxhyphen{}start guide. This file is automatically detected and previewed by some software.

\sphinxlineitem{pyproject.toml}
\sphinxAtStartPar
A standardised configuration file for the project, namely specifying build system requirements. Included to solve
build\sphinxhyphen{}tool dependency problems and install necessary build tools in a virtual environment. Also required for
configuring various development tools.

\sphinxlineitem{requirements.txt}
\sphinxAtStartPar
A list of all the required Python packages along with the minimum version number. This can be used for
automatically installing all requirements via applications such as \sphinxhref{https://docs.python.org/3/installing/}{PIP}.

\sphinxlineitem{LICENSE}
\sphinxAtStartPar
Standard MIT license included with the project

\end{description}

\sphinxstepscope


\section{Software Usage}
\label{\detokenize{usage/execution:software-usage}}\label{\detokenize{usage/execution::doc}}
\sphinxAtStartPar
The navigation software can be used in many ways. Firstly, and most commonly, it can simply be executed via command
line with a selected configuration file. The toolbox can also be executed through Integrated Development Environments
(IDEs) such as PyCharm, Spyder or IDLE. These are particularly useful for viewing the output values in detail and/or
performing additional processing steps such as calculating or plotting statistics. Lastly, the toolbox can be
integrated and employed by other Python projects. The most common ways to use the software are discussed
in this section.

\begin{sphinxadmonition}{important}{Important:}
\sphinxAtStartPar
During development it was noted that the interface and layouts for the various IDEs used evolved with
time. There is also small variations between platforms such as Microsoft Windows and MacOS. The names and
locations of the buttons, menus and settings discussed in this section could have changed slightly.
\end{sphinxadmonition}

\begin{DUlineblock}{0em}
\item[] 
\end{DUlineblock}


\subsection{Command Line Usage}
\label{\detokenize{usage/execution:command-line-usage}}
\noindent{\hspace*{\fill}\sphinxincludegraphics[width=105\sphinxpxdimen]{{terminal_icon}.png}\hspace*{\fill}}

\sphinxAtStartPar
For standalone usage, the software simulation can be simply executed from the command line when within the project’s
root directory. This usage has the advantage of not requiring the software to be installed or invoked via
Python programming. To run a simulation directly, \sphinxstyleemphasis{employing the default configuration settings}, the following single
command can be used. This launch command is identical for Windows, Linux and MacOS systems.
\sphinxSetupCaptionForVerbatim{Execute PNT Toolbox with default settings}
\def\sphinxLiteralBlockLabel{\label{\detokenize{usage/execution:id3}}}
\begin{sphinxVerbatim}[commandchars=\\\{\}]
\PYG{n}{python3} \PYG{n}{run}\PYG{o}{.}\PYG{n}{py}
\end{sphinxVerbatim}

\sphinxAtStartPar
If multiple versions of Python are installed on the host machine it is always best to explicitly state which version
is to be used to avoid any compatibility issues. To use Python 3.12 replace ‘\sphinxstylestrong{python3}’ with ‘\sphinxstylestrong{python3.12}’.

\sphinxAtStartPar
On execution, the software will extract its settings and user preferences from a provided configuration file.
Configuration files are used to describe all input parameters including the simulated environment, hardware/sensor
activity and vehicle limitations/restrictions. A requested configuration can be state by supplying its name (and
location) as the first argument to the run script. If a configuration file is not provided, the default
configuration \sphinxcode{\sphinxupquote{default\_settings.ini}} will be used automatically.

\sphinxAtStartPar
An example of using a custom created configuration file named \sphinxcode{\sphinxupquote{custom\_config.ini}} (that is located in the project’s
root directory) has been generated below.
\sphinxSetupCaptionForVerbatim{Execute PNT  Toolbox with custom settings}
\def\sphinxLiteralBlockLabel{\label{\detokenize{usage/execution:id4}}}
\begin{sphinxVerbatim}[commandchars=\\\{\}]
\PYG{n}{python3} \PYG{n}{run}\PYG{o}{.}\PYG{n}{py} \PYG{n}{custom\PYGZus{}settings}\PYG{o}{.}\PYG{n}{ini}
\end{sphinxVerbatim}

\sphinxAtStartPar
The expected structure and format of configuration files is discussed in detail under
{\hyperref[\detokenize{usage/configuration:user-configuration}]{\sphinxcrossref{\DUrole{std,std-ref}{User Configuration}}}} (discussed next).

\begin{DUlineblock}{0em}
\item[] 
\end{DUlineblock}


\subsection{Anaconda/Spyder Usage}
\label{\detokenize{usage/execution:anaconda-spyder-usage}}
\noindent{\hspace*{\fill}\sphinxincludegraphics[height=100\sphinxpxdimen]{{spyder_icon}.png}\hspace*{\fill}}

\sphinxAtStartPar
To run the software through \sphinxhref{https://www.anaconda.com/products/individual}{Anaconda/Spyder} simply open Spyder and
select “\sphinxstylestrong{New Project…}” under the “\sphinxstylestrong{Project}” dropdown menu at the top of the page. From there, create a new
project with an appropriate name such as ‘Nav Toolbox’. After this step, copy or move the entire contents of the
toolbox’s root directory to the location of the created Spyder project. Once complete, the files should appear on
the left hand\sphinxhyphen{}side in Spyder. To launch the software \sphinxstylestrong{right\sphinxhyphen{}click} the script \sphinxcode{\sphinxupquote{run.py}} and select “\sphinxstylestrong{run}”.
The software should now be running with the default settings.

\begin{sphinxadmonition}{note}{Note:}
\sphinxAtStartPar
It is best to create a new project rather than attempting to open the toolbox as an existing one, since
creating a new project results in Spyder producing new hidden auxiliary and settings files (based on your
environment rather than trusting the included ones).
\end{sphinxadmonition}

\sphinxAtStartPar
If you wish to select a different configuration file to use rather than the default settings, this can be done so by
providing the location of a configuration file as the first command line argument. To do this in Spyder simply
select “\sphinxstylestrong{Configuration per file…}” from the “\sphinxstylestrong{Run}” dropdown menu at the top. Ensure the execution file
\sphinxcode{\sphinxupquote{run.py}} is selected under the dropdown box titled “\sphinxstylestrong{Select a run configuration:}”. Then under
“\sphinxstylestrong{General settings}” enable “\sphinxstylestrong{Command line options:}” and enter the name/location of the configuration file to use.
To launch the toolbox with these settings click the “\sphinxstylestrong{run}” button located at the bottom of the menu.

\sphinxAtStartPar
For an example, lets assume we wish to use a configuration file called \sphinxcode{\sphinxupquote{custom\_settings.ini}} which is located in the
same directory as \sphinxcode{\sphinxupquote{default\_settings.ini}}. To do this, we would enter \sphinxstylestrong{custom\_settings.ini} into the
“\sphinxstylestrong{Command line options:}” box. The toolbox will then automatically find the configuration file to use.

\begin{DUlineblock}{0em}
\item[] 
\end{DUlineblock}


\subsection{PyCharm Usage}
\label{\detokenize{usage/execution:pycharm-usage}}
\noindent{\hspace*{\fill}\sphinxincludegraphics[height=80\sphinxpxdimen]{{pycharm_icon}.png}\hspace*{\fill}}

\sphinxAtStartPar
To use the software with \sphinxhref{https://www.jetbrains.com/pycharm/download/}{JetBrain’s PyCharm} simply create a new
project. If not prompted at the home screen, this can be achieved by selected “\sphinxstylestrong{New Project…}” from the
“\sphinxstylestrong{File}” dropdown menu at the top of the screen. From there ensure the project is being created in a sensible
location and given an appropriate name. During the new project’s creation you will be asked to select the
“\sphinxstylestrong{Python Interpreter}” to use. It is recommended to select “\sphinxstylestrong{New Environment Using:}” one of the followings:
\begin{enumerate}
\sphinxsetlistlabels{\arabic}{enumi}{enumii}{}{.}%
\item {} 
\sphinxAtStartPar
\sphinxstylestrong{Virtualenv} \sphinxhyphen{} Isolated virtual environment (highly recommended).

\item {} 
\sphinxAtStartPar
\sphinxstylestrong{Pipenv} \sphinxhyphen{} Creates and manages a virtualenv with Pipfiles.

\item {} 
\sphinxAtStartPar
\sphinxstylestrong{Conda} \sphinxhyphen{} Anaconda Python distribution (requires Conda to be installed).

\end{enumerate}

\begin{sphinxadmonition}{note}{Note:}
\sphinxAtStartPar
More information of configuring a Python interpreter in PyCharm is available
\sphinxhref{https://www.jetbrains.com/help/pycharm/configuring-python-interpreter.html}{here}!
\end{sphinxadmonition}

\sphinxAtStartPar
With the project created copy the Quantum Navigation files to the its location. Once complete, the files should
appear within the left hand\sphinxhyphen{}side panel. To launch the software \sphinxstylestrong{right\sphinxhyphen{}click} the script \sphinxcode{\sphinxupquote{run.py}}
inside the root folder and select “\sphinxstylestrong{Run ‘run\_simulation’}”. The software should now be running with
the default settings.

\sphinxAtStartPar
If you wish to select a different configuration file to use rather than the default settings, this can be done so by
providing the location of a configuration file as the first command line argument. To do this in PyCharm simply
\sphinxstylestrong{right\sphinxhyphen{}click} the script and select the \sphinxstylestrong{More Run/Debug} menu, then select “\sphinxstylestrong{Modify Configurations…}”.
From this menu enter the name/location of the configuration file to use in the box named “\sphinxstylestrong{Script Parameters}” and
then click “\sphinxstylestrong{Apply}”. Each time the \sphinxcode{\sphinxupquote{run.py}} is executed these settings will now be used.

\sphinxAtStartPar
Note these settings may reset or change between sessions or when executing other scripts.

\begin{DUlineblock}{0em}
\item[] 
\end{DUlineblock}


\subsection{Eclipse/PyDev Usage}
\label{\detokenize{usage/execution:eclipse-pydev-usage}}
\noindent{\hspace*{\fill}\sphinxincludegraphics[height=60\sphinxpxdimen]{{pydev_icon}.png}\hspace*{\fill}}

\sphinxAtStartPar
The project can also be used in \sphinxhref{https://www.pydev.org/}{PyDev} (an \sphinxhref{https://www.eclipse.org/ide/}{Eclipse}
plug\sphinxhyphen{}in for Python development). To use the Quantum Navigation Toolbox in PyDev, simply create a new project. This can
be done by going to the “\sphinxstylestrong{File}” Dropdown menu and selecting “\sphinxstylestrong{New}”, “\sphinxstylestrong{Project}”, “\sphinxstylestrong{PyDev}”, then
“\sphinxstylestrong{PyDev project}”.
From the PyDev Project creation window ensure the new project is given a suitable name and location. Ensure
“\sphinxstylestrong{Python}” is selected as the “\sphinxstylestrong{Project type}” and that a compatible version of Python is chosen
(i.e. Python 3.8 or above).

\begin{sphinxadmonition}{note}{Note:}
\sphinxAtStartPar
More information of creating and importing Python projects in PyDev is available
\sphinxhref{https://www.pydev.org/manual\_101\_project\_conf.html}{here}!
\end{sphinxadmonition}

\sphinxAtStartPar
With the project created copy the toolbox files to the its location. Once complete, the files should
appear within the left hand\sphinxhyphen{}side panel. To launch the software within PyDev \sphinxstylestrong{open the script} \sphinxcode{\sphinxupquote{run.py}}
inside the root folder. With the script \sphinxcode{\sphinxupquote{run.py}} being displayed in the editor, go to the “\sphinxstylestrong{Run}” dropdown menu
and select “\sphinxstylestrong{Run As}” followed by “\sphinxstylestrong{Python Run}”. Alternatively the shortcut key \sphinxstylestrong{F9} can be used to run the
script currently open. The software should now be running with the default settings.

\sphinxAtStartPar
If you wish to select a different configuration file to use rather than the default settings, this can be done so by
providing the location of a configuration file as the first command line argument. To do this in PyDev firstly open the
main run file \sphinxcode{\sphinxupquote{run.py}} so it is being displayed in the editor. Then go to the “\sphinxstylestrong{Run}” dropdown menu and select
“\sphinxstylestrong{Run Configurations…}” located under the “Run As”. From there ensure the current script is selected
on the \sphinxstylestrong{left\sphinxhyphen{}hand side} under “\sphinxstylestrong{Python Run}”. Finally, under the “\sphinxstylestrong{Arguments}” tab enter the name/location of the
configuration file to use in the box titled “\sphinxstylestrong{Program arguments:}” then click “\sphinxstylestrong{apply}” followed by “\sphinxstylestrong{Run}”.
If needed, alternative steps for setting runtime arguments are available
\sphinxhref{https://www.dev2qa.com/how-to-pass-command-line-parameters-to-python-script-in-eclipse-pydev/}{here}!

\sphinxstepscope


\section{Technical Overview}
\label{\detokenize{usage/technical:technical-overview}}\label{\detokenize{usage/technical::doc}}
\sphinxAtStartPar
This section discusses the technical aspects of the project; detailing the internal dataflow. This is namely used to
describe the key stages during typical execution of the software. Moreover, this section also outlines how the project
correctly enforces separations between the simulated “real” world, extracted measurements and the estimated states.

\noindent{\hspace*{\fill}\sphinxincludegraphics[width=500\sphinxpxdimen]{{sim_meas_est}.png}\hspace*{\fill}}

\begin{DUlineblock}{0em}
\item[] 
\end{DUlineblock}


\subsection{Trajectory Generation}
\label{\detokenize{usage/technical:trajectory-generation}}
\sphinxAtStartPar
Firstly, during a typical run\sphinxhyphen{}cycle, the software will firstly calculate the ground truth trajectory. This is performed
using a two\sphinxhyphen{}step approach. During the first step a trajectory path is constructed to reach the given base points while
complying with the selected vehicle model. This is generated at a frequency subject to the vehicles speed. During the
next step the trajectory properties are calculated at each step (being the position, velocity, acceleration,
attitude and angle rate records). These produced values can then serve as a look\sphinxhyphen{}up table with interpolation,
using the simulation time in seconds as the key.


\subsection{Estimation Initialisation}
\label{\detokenize{usage/technical:estimation-initialisation}}
\sphinxAtStartPar
Secondly, with the ground truth trajectory generated, the initial estimation states can be declared. In the simplest
usage these can be read directly from the first record of the ground truth trajectory (at 0 seconds). But
realistically, initial estimation errors will be applied on top of these. Other differences, such as expected
state errors and a non\sphinxhyphen{}identical gravity model (purposely unlike the one used to generate the ground truth)
can be set here.


\subsection{Sensor Configuration}
\label{\detokenize{usage/technical:sensor-configuration}}
\sphinxAtStartPar
Thirdly, using the initial estimated states, the sensors and their corresponding fusion methods can now be
configured and paired. Individual sensors, such as the accelerometer and gyroscopes, can be configured with custom
error profiles, axis orientations and random number generation (used for generating pseudo\sphinxhyphen{}random noise and errors).
These will then assigned to one or more fusion methods. For instance, the accelerometer and gyroscopes could be
be allocated to a configured Kalman Filter for INS fusion. In addition, the same sensors could be shared and also
assigned to other fusion methods, such as quantum sensors.


\subsection{Simulation Loop}
\label{\detokenize{usage/technical:simulation-loop}}
\sphinxAtStartPar
With everything configured, the main simulation loop can now be executed with all over the above. At a minimum, a
simulation loop requires a ground truth trajectory, initial estimated state and a series of sensor\sphinxhyphen{}fusion pairs.
Optionally, a custom clock model can be included to invoke purposeful timing errors. Custom display and result
collectors can also be passed to obtain information with configured levels of filtering.

\sphinxAtStartPar
During the first step of the simulation, a scheduler is initialised for handling the sensors and fusion, calling the
next due senor or fusion method on time. When a sensor is ready for a measurement, the simulation time is move up to
that time, with clock errors applied. The sensor is then given the ground truth record for that time, which is then
uses to produce a noisy measurement (which is held/cached internally). When a fusion method is ready, depending on
its trigger mode, it collects the noisy measurements from its sensors and processes then before updating the
estimated state.

\sphinxAtStartPar
This approach allows sensors and fusion methods to be dynamically included in simulations, while enforcing complete
separation between ground truth (solely seen by sensors) and the estimated state (solely updated by fusion methods
with noisy measurements).


\subsection{Capturing Results}
\label{\detokenize{usage/technical:capturing-results}}
\sphinxAtStartPar
After the simulation loop is complete, the captured results (either from the default or given results collector)
will be returned. Results are firstly written to disk using the projects own binary file format (optimised for
sequentially reading and writing of records). After this, results will then be exported in any of the requested file
formats (such as csv, binary, numpy, MATLAB save) with various levels of compression (both lossly and lossless).
Separate files are made for both the ground truth (optional) and estimation state records.


\subsection{Displaying Results}
\label{\detokenize{usage/technical:displaying-results}}
\sphinxAtStartPar
If enabled, after results have been written to disk, their data will be displayed in a series of interactive
graphs and plots comparing the ground truth and estimation history against time. The supported figure types are
as listed:
\begin{itemize}
\item {} 
\sphinxAtStartPar
Summary Table Plot

\item {} 
\sphinxAtStartPar
2D Open Street Map Plot

\item {} 
\sphinxAtStartPar
2D Lat\sphinxhyphen{}Lon Plot

\item {} 
\sphinxAtStartPar
3D Lat\sphinxhyphen{}Lon\sphinxhyphen{}Alt Plot

\item {} 
\sphinxAtStartPar
Position Axes Plot

\item {} 
\sphinxAtStartPar
Velocity Axes Plot

\item {} 
\sphinxAtStartPar
Acceleration Axes Plot

\item {} 
\sphinxAtStartPar
Attitude Axes Plot

\item {} 
\sphinxAtStartPar
Angle Rate Axes Plot

\end{itemize}


\subsection{Recreation}
\label{\detokenize{usage/technical:recreation}}
\sphinxAtStartPar
After the simulation, the same results can be recreated by setting the same parameters and random seeds. To aid with
this, a complete configuration file containing all parameters used is dumped into the generated results folder.

\sphinxstepscope


\section{Trajectory Generation}
\label{\detokenize{usage/trajectory:trajectory-generation}}\label{\detokenize{usage/trajectory::doc}}
\sphinxAtStartPar
Ground truth records are represented as a trajectory, a series of waypoint records generated at fixed frequency that
support interpolation look\sphinxhyphen{}up. The project allows a trajectory is constructed of a vehicle model, a described path,
and a gravity model.


\subsection{Vehicle Model}
\label{\detokenize{usage/trajectory:vehicle-model}}
\sphinxAtStartPar
A vehicle model describes the limitations of the vehicle platform, stating its expected speed, moving capabilities,
acceleration rates, maneuvering restrictions and expected behaviour. These essentially define how trajectories are
generated with realistic behaviours. For example, the path generated by a land vehicle and large ship vessel will
be very different.

\begin{sphinxadmonition}{important}{Important:}
\sphinxAtStartPar
A vehicle model is required to restrict the simulated maneuvering of the vehicle platform to within realistic
bounds. Ensure one is configured correctly for the requested trajectory.
\end{sphinxadmonition}


\subsection{Path Generation}
\label{\detokenize{usage/trajectory:path-generation}}
\sphinxAtStartPar
Trajectory paths are configured either as a arithmetical shape or series of base points that will form a shape.
From these, a series waypoint positions are generated (using constraints set by the vehicle model). This operates by
simulating controlled movement to reach each base point, or get as close to as possible. The result from this is a
table of positions at fixed frequency. These can be generated using one of the three methods:


\subsubsection{INI Configuration}
\label{\detokenize{usage/trajectory:ini-configuration}}
\sphinxAtStartPar
Properties in the INI configuration file can describe the the latitude, longitude and altitude base points for a
trajectory to cover. Moreover, a sub\sphinxhyphen{}configuration file can be referenced, allowing for a complex trajectory to be
shared. Indepth details for INI configuration is covered by it own section.


\subsubsection{CSV Configuration}
\label{\detokenize{usage/trajectory:csv-configuration}}
\sphinxAtStartPar
Similar to the above, a comma\sphinxhyphen{}separate\sphinxhyphen{}values (CSV) file can be used to describe the base points. For this a CSV file
containing either three or four columns can be used. If three columns are used, the records for each line will
assume to be the latitude, longitude and altitude (given in decimal degrees, decimal degrees and metres respectively).
Alternatively, if four columns are used, then the first will be the arrival time at each base point in seconds. If this
is not provided, the vehicle model’s average speed will be used to determine this.


\subsubsection{Racetrack Configuration}
\label{\detokenize{usage/trajectory:racetrack-configuration}}
\sphinxAtStartPar
A Racetrack is the name referred to for a looping arithmetical trajectory. The advantage to using one of these instead
of the above base point methods is that a trajectory path of specified duration and lap time can be quickly generated
at any location, which is useful for rapid prototyping.


\subsection{Gravity Model}
\label{\detokenize{usage/trajectory:gravity-model}}
\sphinxAtStartPar
A gravity model is required for correctly calculating the gravitational and coriolis effects on the trajectory values.
The selected model can range from something as simple as the Somigliana formula, to something as complex as a large
scale database. Typically, a far more accurate model is used to generate the ground truth than what is used during
simulation estimation.

\begin{sphinxadmonition}{note}{Note:}
\sphinxAtStartPar
More information on {\hyperref[\detokenize{usage/gravity_modelling:gravity-modelling}]{\sphinxcrossref{\DUrole{std,std-ref}{gravity models}}}} is given under its own section.
\end{sphinxadmonition}

\sphinxstepscope


\section{Gravity Modelling}
\label{\detokenize{usage/gravity_modelling:gravity-modelling}}\label{\detokenize{usage/gravity_modelling::doc}}
\sphinxAtStartPar
The project include a small collection of gravity models. Each is equipped with their own methods for calculating
various values, such as the vector acceleration, vertical gravity gradient and expected anomaly or disturbance values
for given positions.

\sphinxAtStartPar
Gravity models are grouped into two sets, being functions and correction maps. Gravity functions serve as a simple
and quick approximation for the gravity at a position. These can be applied for any location and are used as base
values. Extending from these are the gravity correction maps, which as suggested apply a correction on top of a
selected base gravity function. These provide correction records, but can only be applied in areas they cover. If
data for the requested position is unavailable, no correction to the value produced by the base function will be
applied.

\begin{sphinxadmonition}{note}{Note:}
\sphinxAtStartPar
For simple usage, only a basic gravity function will be needed. However, for applying realistic gravity behaviour,
which is required for some forms of sensing, a high resolution correction map will also be needed.
\end{sphinxadmonition}


\subsection{Gravity Functions}
\label{\detokenize{usage/gravity_modelling:gravity-functions}}

\subsubsection{Fixed Constant Value}
\label{\detokenize{usage/gravity_modelling:fixed-constant-value}}
\sphinxAtStartPar
Essentially not a function, serving a fixed value of 9.81 used for vertical acceleration at any position.
This is primarily used for debugging.


\subsubsection{Simple Uniform Sphere}
\label{\detokenize{usage/gravity_modelling:simple-uniform-sphere}}
\sphinxAtStartPar
Approximates the Earth as uniform sphere with height extrapolation. A slight improvement over the Fixed Constant model,
but is again primarily used for debugging.


\subsubsection{Somigliana Formula}
\label{\detokenize{usage/gravity_modelling:somigliana-formula}}
\sphinxAtStartPar
The standardised “Normal” gravity function with height extrapolation (named after Carlo Somigliana (1860\textendash{}1955). Also
known as the International Gravity Formula 1980. This is far more accurate than the previous models, and also prooduces
east component values in calculated gravity vectors.


\subsubsection{NIMA’s WGS\sphinxhyphen{}84 Gravity}
\label{\detokenize{usage/gravity_modelling:nimas-wgs-84-gravity}}
\sphinxAtStartPar
Calculations provided by National Imagery and Mapping Agency (NIMA) in “Department of Defense World Geodetic System
1984, Its Definition and Relationship with Local Geodetic Systems.”. This shares many similarities with the Somigliana
model. This was namely included due to it being used in software such as MATLAB as standard.


\subsection{Gravity Correction Maps}
\label{\detokenize{usage/gravity_modelling:gravity-correction-maps}}

\subsubsection{EGM Models}
\label{\detokenize{usage/gravity_modelling:egm-models}}
\sphinxAtStartPar
Application of 7 Earth Gravitation Models defining geoid undulations in raster form. These essentially provide the
geoid height correction on top of the WGS\sphinxhyphen{}84 reference ellipsoid. Supported EGM models include:
\begin{itemize}
\item {} 
\sphinxAtStartPar
EGM\sphinxhyphen{}1984 (in 30 and 15 arc\sphinxhyphen{}minute resolutions)

\item {} 
\sphinxAtStartPar
EGM\sphinxhyphen{}1996 (in 15 and 5 arc\sphinxhyphen{}minute resolutions)

\item {} 
\sphinxAtStartPar
EGM\sphinxhyphen{}2008 (in 5, 2.5 and 1 arc\sphinxhyphen{}minute resolutions)

\end{itemize}

\begin{sphinxadmonition}{note}{Note:}
\sphinxAtStartPar
While EGM\sphinxhyphen{}2020 has been technically released, it is currently unavailable for public access.
\end{sphinxadmonition}


\subsubsection{Geosat\sphinxhyphen{}44}
\label{\detokenize{usage/gravity_modelling:geosat-44}}
\sphinxAtStartPar
the Geosat\sphinxhyphen{}44 gravity model provided data collected in 1999 from satellite altimeter measurements, which contain
precise geoid and gravity anomaly profiles constructed from the average of 44 repeat cycles. In total, this dataset
contains 987,755 ascending records and 991,313 descending records between the latitude bounds +/\sphinxhyphen{} 72 degrees
north/south. This model is represented as scatter points, which are triangulated and used in interpolation.


\subsubsection{Marine Gravity Map}
\label{\detokenize{usage/gravity_modelling:marine-gravity-map}}
\sphinxAtStartPar
The marine gravity map provides gravity anomaly values for acquired from satellite altimeter data in 2013. It is based
on previous data from CryoSat\sphinxhyphen{}2 and Jason\sphinxhyphen{}1 satellite projects. The resulting model is a 9,600 x 21,600 raster grid,
between the latitude bounds +/\sphinxhyphen{} 80 degrees north, with interpolation performed for land space.


\subsubsection{GGM Plus}
\label{\detokenize{usage/gravity_modelling:ggm-plus}}
\sphinxAtStartPar
The Global Gravity Map Plus (GGM Plus) model is the result of a joint research initiative by Curtin University and
TU Munich University in 2013. It provides records of gravitational acceleration and distributable at a 7.2 arc\sphinxhyphen{}second
(approximately 200 metres) spatial resolution, for all land and near\sphinxhyphen{}coastal areas. The resulting model consists of
3,062,677,383 points, partitioned into 5 x 5 degree tiles, covering latitude bounds +/\sphinxhyphen{} 60 degrees north/south.


\subsubsection{SRTM2Gravity}
\label{\detokenize{usage/gravity_modelling:srtm2gravity}}
\sphinxAtStartPar
The SRTM2Gravity model extends from GGM Plus by performing the conversion of Shuttle Radar Topography Mission\sphinxhyphen{}based
digital elevation data to implied gravity effects. It provides records for full\sphinxhyphen{}scale and residual gravity at 3
arc\sphinxhyphen{}second (approximately 90 metres) spatial resolution, for all land and near\sphinxhyphen{}coastal areas. The resulting model is
partitioned into 1 x 1 degree tiles, covering latitude bounds +/\sphinxhyphen{} 60 degrees north/south.


\subsubsection{Irish Sea Model}
\label{\detokenize{usage/gravity_modelling:irish-sea-model}}
\sphinxAtStartPar
A bespoke marine gravity anomaly map covering an area a north western region Irish sea. Support for this model
was included by request.

\sphinxstepscope


\section{Sensor Processing}
\label{\detokenize{usage/measurement:sensor-processing}}\label{\detokenize{usage/measurement::doc}}
\sphinxAtStartPar
To acquire information from the ground truth and produce measurement of control noise and errors, the project presents
a handful of sensors. Each of these implement the base sensor class and feature methods for generating measurements,
updating internal states and indicating when their next measurement is due. These sensors are categorised below:


\subsection{Classic Inertia}
\label{\detokenize{usage/measurement:classic-inertia}}
\sphinxAtStartPar
Classes representing conventional inertial based sensors are the accelerometer and gyroscope. These are implemented
as generic sensors, as they are reused internally by other sensors within the toolbox. These are responsible for
obtaining  acceleration and angle rate measurements respectively.
\begin{enumerate}
\sphinxsetlistlabels{\arabic}{enumi}{enumii}{}{.}%
\item {} 
\sphinxAtStartPar
{\hyperref[\detokenize{modules/measurement/accelerometer:accelerometer-py}]{\sphinxcrossref{\DUrole{std,std-ref}{Generic Accelerometer}}}} (\sphinxcode{\sphinxupquote{Accelerometer}})

\item {} 
\sphinxAtStartPar
{\hyperref[\detokenize{modules/measurement/gyroscope:gyroscope-py}]{\sphinxcrossref{\DUrole{std,std-ref}{Generic Gyroscope}}}} (\sphinxcode{\sphinxupquote{Gyroscope}})

\end{enumerate}


\subsection{Altitude}
\label{\detokenize{usage/measurement:altitude}}
\sphinxAtStartPar
For obtaining altitude measurements, an altimeter sensor has been implemented. This simple obtains the measured
vertical position, following configured noise profiles.
\begin{enumerate}
\sphinxsetlistlabels{\arabic}{enumi}{enumii}{}{.}%
\item {} 
\sphinxAtStartPar
{\hyperref[\detokenize{modules/measurement/altimeter:altimeter-py}]{\sphinxcrossref{\DUrole{std,std-ref}{Generic Altimeter}}}} (\sphinxcode{\sphinxupquote{Altimeter}})

\end{enumerate}


\subsection{GPS/GNSS}
\label{\detokenize{usage/measurement:gps-gnss}}
\sphinxAtStartPar
For simulating GPS measurements two classes have been implemented. The first class is satellite, which represents a
satellite in orbit. While not specifically a sensor class, multiple satellite instances are used by the GPS sensor
class, which is responsible for producing measurements. Measurement information include clock times, pseudo range
information and estimated position and velocity data.
\begin{enumerate}
\sphinxsetlistlabels{\arabic}{enumi}{enumii}{}{.}%
\item {} 
\sphinxAtStartPar
{\hyperref[\detokenize{modules/gps/satellite:satellite-py}]{\sphinxcrossref{\DUrole{std,std-ref}{GPS Satellite}}}} (\sphinxcode{\sphinxupquote{Satellite}})

\item {} 
\sphinxAtStartPar
{\hyperref[\detokenize{modules/gps/sensor:gps-sensor-py}]{\sphinxcrossref{\DUrole{std,std-ref}{GPS Sensor}}}} (\sphinxcode{\sphinxupquote{GPSSensor}})

\end{enumerate}


\subsection{Quantum Inertia}
\label{\detokenize{usage/measurement:quantum-inertia}}
\sphinxAtStartPar
Two cold\sphinxhyphen{}atom quantum based inertia sensors have been implemented. The first one is a basic concept, simply extending
the use of the generic altimeter by simulating fairly accurate average inertia measurements over a duty cycle. The
second one improves on this by simulating and processing the movement of atoms to gain measurements.
\begin{enumerate}
\sphinxsetlistlabels{\arabic}{enumi}{enumii}{}{.}%
\item {} 
\sphinxAtStartPar
{\hyperref[\detokenize{modules/quantum/base:base-py}]{\sphinxcrossref{\DUrole{std,std-ref}{Concept Quantum IMU}}}} (\sphinxcode{\sphinxupquote{ConceptQuantumImu}})

\item {} 
\sphinxAtStartPar
{\hyperref[\detokenize{modules/quantum/realistic:realistic-py}]{\sphinxcrossref{\DUrole{std,std-ref}{Realistic Quantum IMU}}}} (\sphinxcode{\sphinxupquote{QuantumIMU}})

\end{enumerate}


\subsection{Quantum Gravity}
\label{\detokenize{usage/measurement:quantum-gravity}}
\sphinxAtStartPar
One cold\sphinxhyphen{}atom gravity gradient sensor has been implemented. This simulates a dual\sphinxhyphen{}interferometer, acquiring
measurement signals from the releasing of two vertically stacked clouds of atoms that share a phase. A series of
signals collected across multiple duty\sphinxhyphen{}cycles can be acquired to form an ellipse, used to estiamted the
vertical gravity gradient.
\begin{enumerate}
\sphinxsetlistlabels{\arabic}{enumi}{enumii}{}{.}%
\item {} 
\sphinxAtStartPar
{\hyperref[\detokenize{modules/quantum/base:base-py}]{\sphinxcrossref{\DUrole{std,std-ref}{Quantum Gravity Gradiometer}}}} (\sphinxcode{\sphinxupquote{GravityGradiometer}})

\end{enumerate}

\sphinxstepscope


\section{Fusion Methods}
\label{\detokenize{usage/fusion:fusion-methods}}\label{\detokenize{usage/fusion::doc}}
\sphinxAtStartPar
To utilise measurements acquired from the sensors described in the previous section, a collection of fusion methods
have been implemented. These are grouped and categorised as before.


\subsection{Classic Inertia}
\label{\detokenize{usage/fusion:classic-inertia}}
\sphinxAtStartPar
Numerical integration methods for the INS include: Simple Numerical, 4th Order Runge\sphinxhyphen{}Kutta, Adams Bashforth and
Kalman Filter Integration. For certain cases, Kalman Filter Integration will be required to leverage its covariance
matrices. For example, some GPS and Quantum Sensor fusion methods will require access to these.
\begin{enumerate}
\sphinxsetlistlabels{\arabic}{enumi}{enumii}{}{.}%
\item {} 
\sphinxAtStartPar
{\hyperref[\detokenize{modules/fusion/ins:ins-py}]{\sphinxcrossref{\DUrole{std,std-ref}{Simple Numerical}}}} (\sphinxcode{\sphinxupquote{NumericalINS}})

\item {} 
\sphinxAtStartPar
{\hyperref[\detokenize{modules/fusion/ins:ins-py}]{\sphinxcrossref{\DUrole{std,std-ref}{Runge Kutta}}}} (\sphinxcode{\sphinxupquote{RungeKuttaINS}})

\item {} 
\sphinxAtStartPar
{\hyperref[\detokenize{modules/fusion/ins:ins-py}]{\sphinxcrossref{\DUrole{std,std-ref}{Adams Bashforth}}}} (\sphinxcode{\sphinxupquote{AdamsBashforthINS}})

\item {} 
\sphinxAtStartPar
{\hyperref[\detokenize{modules/fusion/ins:ins-py}]{\sphinxcrossref{\DUrole{std,std-ref}{Kalman Filter}}}} (\sphinxcode{\sphinxupquote{KalmanINS}})

\end{enumerate}


\subsection{Altitude}
\label{\detokenize{usage/fusion:altitude}}
\sphinxAtStartPar
For processing altimeter measurements two fusion methods have been implemented. The first is simple fixed\sphinxhyphen{}gain, which
updates the current estimated altitude by using a weighted average between the measurement and current state. The
second is alpha\sphinxhyphen{}beta fusion, which extends on the previous by also correcting the vertical velocity in a similar way.
\begin{enumerate}
\sphinxsetlistlabels{\arabic}{enumi}{enumii}{}{.}%
\item {} 
\sphinxAtStartPar
{\hyperref[\detokenize{modules/measurement/altimeter:altimeter-py}]{\sphinxcrossref{\DUrole{std,std-ref}{Fixed\sphinxhyphen{}Gained Fusion}}}} (\sphinxcode{\sphinxupquote{FixedGainAltimeter}})

\item {} 
\sphinxAtStartPar
{\hyperref[\detokenize{modules/measurement/altimeter:altimeter-py}]{\sphinxcrossref{\DUrole{std,std-ref}{Alpha\sphinxhyphen{}Beta Fusion}}}} (\sphinxcode{\sphinxupquote{AlphaBetaAltimeter}})

\end{enumerate}


\subsection{GPS/GNSS}
\label{\detokenize{usage/fusion:gps-gnss}}
\sphinxAtStartPar
For processing GPS measurements three fusion methods are available. Firstly, fixed gain simple updates the estimated
position and velocity by using a fixed weighted gain for the measurements produced by the satellites. Secondly, Loosely
coupled fusion improves on this by using a Kalman filter to integrate INS parameters with the received GPS position
and velocity measurements. Lastly, Tightly coupled fusion further extends this by operating on the raw GPS measurements
from the receiver, such as the satellite pseudo\sphinxhyphen{}ranges and Doppler observables.
\begin{enumerate}
\sphinxsetlistlabels{\arabic}{enumi}{enumii}{}{.}%
\item {} 
\sphinxAtStartPar
{\hyperref[\detokenize{modules/gps/fusion:fusion-py}]{\sphinxcrossref{\DUrole{std,std-ref}{GPS Fixed Gain}}}} (\sphinxcode{\sphinxupquote{GpsFixedGainFusion}})

\item {} 
\sphinxAtStartPar
{\hyperref[\detokenize{modules/gps/fusion:fusion-py}]{\sphinxcrossref{\DUrole{std,std-ref}{GPS Loosely Coupled}}}} (\sphinxcode{\sphinxupquote{GpsLooseFusion}})

\item {} 
\sphinxAtStartPar
{\hyperref[\detokenize{modules/gps/fusion:fusion-py}]{\sphinxcrossref{\DUrole{std,std-ref}{GPS Tightly Coupled}}}} (\sphinxcode{\sphinxupquote{GpsTightFusion}})

\end{enumerate}


\subsection{Quantum Inertia}
\label{\detokenize{usage/fusion:quantum-inertia}}
\sphinxAtStartPar
For processing cold\sphinxhyphen{}atom quantum inertia measurements three fusion methods are implemented. The first method is
called Concept Fusion, which mostly serves as a base class. From the collected measurements this simple calculates
the believed sensor biases and corrects the estimated state. The second method is a particle filter method, which
applies a particle filter approach for estimating the acceleration and angle rate biases independently. The third
method extends on this further by also accounting for any sensor misalignment.
\begin{enumerate}
\sphinxsetlistlabels{\arabic}{enumi}{enumii}{}{.}%
\item {} 
\sphinxAtStartPar
{\hyperref[\detokenize{modules/quantum/base:base-py}]{\sphinxcrossref{\DUrole{std,std-ref}{Concept Quantum Fusion}}}} (\sphinxcode{\sphinxupquote{ConceptQuantumFusion}})

\item {} 
\sphinxAtStartPar
{\hyperref[\detokenize{modules/quantum/realistic:realistic-py}]{\sphinxcrossref{\DUrole{std,std-ref}{Particle Filter Quantum Fusion}}}} (\sphinxcode{\sphinxupquote{QuantumPFFusion}})

\item {} 
\sphinxAtStartPar
{\hyperref[\detokenize{modules/quantum/realistic:realistic-py}]{\sphinxcrossref{\DUrole{std,std-ref}{Particle Filter and Misaligned Quantum Fusion}}}} (\sphinxcode{\sphinxupquote{QuantumPFMFusion}})

\end{enumerate}


\subsection{Quantum Gravity}
\label{\detokenize{usage/fusion:quantum-gravity}}
\sphinxAtStartPar
Currently, there is only one fusion method for processing the gravity gradient sensor, by performing gravity
gradient map matching. This employs a particle filter approach, exploring a search space for best fitting
position, velocity and attitude candidates.
\begin{enumerate}
\sphinxsetlistlabels{\arabic}{enumi}{enumii}{}{.}%
\item {} 
\sphinxAtStartPar
{\hyperref[\detokenize{modules/quantum/gravity_gradient:gravity-gradient-py}]{\sphinxcrossref{\DUrole{std,std-ref}{Quantum Gravity Gradient Fusion}}}} (\sphinxcode{\sphinxupquote{GravityGradientPF}})

\end{enumerate}

\sphinxstepscope


\section{User Configuration}
\label{\detokenize{usage/configuration:user-configuration}}\label{\detokenize{usage/configuration::doc}}
\sphinxAtStartPar
Configuration is primarily performed by providing the location of a configuration file to use as the
first command line argument. Examples of how to do this is cover under {\hyperref[\detokenize{usage/execution:software-usage}]{\sphinxcrossref{\DUrole{std,std-ref}{Software Usage}}}}.
This section discusses the various sections contained with a standard configuration files, along with their
purpose and key variables.

\sphinxAtStartPar
In total, there is over 400 customisable variables that can be changed. Except for a few cases, each variable has its
own default value which is used if it is not included in the current configuration file. In the event a configuration
file is not provided, the default configuration \sphinxcode{\sphinxupquote{default\_settings.ini}} will be used automatically. It is recommended
that a copy of \sphinxcode{\sphinxupquote{default\_settings.ini}} is made and renamed. This copy can serve as a fully commented template
configuration file to use and modify.

\begin{sphinxadmonition}{note}{Note:}
\sphinxAtStartPar
Configuration files used by the project are formatted as “initialisation” files (.ini format). Because .ini
files are a popular and commonly used method for software configuration, many editor programs support syntax
highlighting for them. To avoid making mistakes or typos, it is recommended to use an editor that supports
initialisation file format.
\end{sphinxadmonition}

\begin{sphinxadmonition}{important}{Important:}\begin{itemize}
\item {} 
\sphinxAtStartPar
Both \sphinxcode{\sphinxupquote{=}} and \sphinxcode{\sphinxupquote{:}} can be used for assigning values.

\item {} 
\sphinxAtStartPar
The symbols \sphinxcode{\sphinxupquote{\#}} and \sphinxcode{\sphinxupquote{;}} are used to state comments.

\item {} 
\sphinxAtStartPar
String values do not require quotes such as \sphinxcode{\sphinxupquote{"}} or \sphinxcode{\sphinxupquote{\textquotesingle{}}}.

\item {} 
\sphinxAtStartPar
In\sphinxhyphen{}line comments (comments after variables) are \sphinxstylestrong{not} supported.

\item {} 
\sphinxAtStartPar
For numerical values anything after a string character is ignored.

\item {} 
\sphinxAtStartPar
File paths can be given in either Windows or POSIX (Linux) formats.

\item {} 
\sphinxAtStartPar
Valid boolean values are \sphinxcode{\sphinxupquote{true}}, \sphinxcode{\sphinxupquote{false}}, \sphinxcode{\sphinxupquote{yes}}, \sphinxcode{\sphinxupquote{no}}, \sphinxcode{\sphinxupquote{on}}, \sphinxcode{\sphinxupquote{off}}, \sphinxcode{\sphinxupquote{1}} or \sphinxcode{\sphinxupquote{0}}.

\item {} 
\sphinxAtStartPar
Mathematics expressions such as \sphinxcode{\sphinxupquote{+}}, \sphinxcode{\sphinxupquote{\sphinxhyphen{}}}, \sphinxcode{\sphinxupquote{*}}, \sphinxcode{\sphinxupquote{/}}, \sphinxcode{\sphinxupquote{//}}, \sphinxcode{\sphinxupquote{\textasciicircum{}}} and \sphinxcode{\sphinxupquote{e}} are supported.

\item {} 
\sphinxAtStartPar
To reduce the likely\sphinxhyphen{}hood of errors and typo, variable names are \sphinxstylestrong{case insensitive}.

\end{itemize}
\end{sphinxadmonition}


\subsection{Configuration Sections}
\label{\detokenize{usage/configuration:configuration-sections}}
\sphinxAtStartPar
A configuration file is composed of 15 sections in total: \sphinxstylestrong{Output}, \sphinxstylestrong{Database}, \sphinxstylestrong{Waypoints}, \sphinxstylestrong{Racetrack},
\sphinxstylestrong{Vehicle}, \sphinxstylestrong{Gravity}, \sphinxstylestrong{Geoid}, \sphinxstylestrong{Measurement}, \sphinxstylestrong{Estimation}, \sphinxstylestrong{Altimeter}, \sphinxstylestrong{GPS}, \sphinxstylestrong{Quantum},
\sphinxstylestrong{Holonomics}, \sphinxstylestrong{Random} and \sphinxstylestrong{Misc}. The use of each of these section is discussed below along with a
description for the most common variables used:


\subsubsection{Output Settings}
\label{\detokenize{usage/configuration:output-settings}}
\sphinxAtStartPar
The \sphinxstylestrong{Output} section holds all variables used for controlling the output of the simulation. This includes
settings for stating what is to be displayed during runtime, what files at to be produced, where files are to be saved,
any down\sphinxhyphen{}sampling or compression that should be applied.


\begin{savenotes}\sphinxattablestart
\sphinxthistablewithglobalstyle
\centering
\sphinxcapstartof{table}
\sphinxthecaptionisattop
\sphinxcaption{Key Output Section Variables}\label{\detokenize{usage/configuration:id1}}
\sphinxaftertopcaption
\begin{tabular}[t]{\X{1}{4}\X{2}{4}\X{1}{4}}
\sphinxtoprule
\sphinxstyletheadfamily 
\sphinxAtStartPar
Variable Name
&\sphinxstyletheadfamily 
\sphinxAtStartPar
Description
&\sphinxstyletheadfamily 
\sphinxAtStartPar
Type
\\
\sphinxmidrule
\sphinxtableatstartofbodyhook
\sphinxAtStartPar
\sphinxcode{\sphinxupquote{saveResults}}
&
\sphinxAtStartPar
Should summary results be saved to disk?
&
\sphinxAtStartPar
Yes \sphinxstyleemphasis{or} No
\\
\sphinxhline
\sphinxAtStartPar
\sphinxcode{\sphinxupquote{saveFigures}}
&
\sphinxAtStartPar
Should figures be generated and saved to disk?
&
\sphinxAtStartPar
Yes \sphinxstyleemphasis{or} No
\\
\sphinxhline
\sphinxAtStartPar
\sphinxcode{\sphinxupquote{showFigures}}
&
\sphinxAtStartPar
Should figures be generated in the web\sphinxhyphen{}browser?
&
\sphinxAtStartPar
Yes \sphinxstyleemphasis{or} No
\\
\sphinxhline
\sphinxAtStartPar
\sphinxcode{\sphinxupquote{copyConfig}}
&
\sphinxAtStartPar
Should a copy of the configuration file be created?
&
\sphinxAtStartPar
Yes \sphinxstyleemphasis{or} No
\\
\sphinxhline
\sphinxAtStartPar
\sphinxcode{\sphinxupquote{resultsDownSampleRate}}
&
\sphinxAtStartPar
A factor to down\sphinxhyphen{}sample the collection of results by.
&
\sphinxAtStartPar
Integer
\\
\sphinxhline
\sphinxAtStartPar
\sphinxcode{\sphinxupquote{outputDownSampleRate}}
&
\sphinxAtStartPar
A factor to down\sphinxhyphen{}sample the exporting of results data by.
&
\sphinxAtStartPar
Integer
\\
\sphinxhline
\sphinxAtStartPar
\sphinxcode{\sphinxupquote{resultsPath}}
&
\sphinxAtStartPar
The location where all results are to be saved.
&
\sphinxAtStartPar
Filepath
\\
\sphinxhline
\sphinxAtStartPar
\sphinxcode{\sphinxupquote{appendDateFolder}}
&
\sphinxAtStartPar
Should the date be appended to the output?
&
\sphinxAtStartPar
Yes \sphinxstyleemphasis{or} No
\\
\sphinxbottomrule
\end{tabular}
\sphinxtableafterendhook\par
\sphinxattableend\end{savenotes}

\begin{sphinxadmonition}{note}{Note:}
\sphinxAtStartPar
Toolbox supports the following save file formats:
\begin{itemize}
\item {} 
\sphinxAtStartPar
\sphinxcode{\sphinxupquote{csv}} \sphinxhyphen{} Comma Seperated Values

\item {} 
\sphinxAtStartPar
\sphinxcode{\sphinxupquote{csv\_gz}} \sphinxhyphen{} Comma Seperated Values with G\sphinxhyphen{}Zip compression

\item {} 
\sphinxAtStartPar
\sphinxcode{\sphinxupquote{numpy}} \sphinxhyphen{} Numpy binary array file (generated with np.save())

\item {} 
\sphinxAtStartPar
\sphinxcode{\sphinxupquote{numpy\_gz}} \sphinxhyphen{} Numpy array file with G\sphinxhyphen{}Zip compression

\item {} 
\sphinxAtStartPar
\sphinxcode{\sphinxupquote{binary}} \sphinxhyphen{} Standard float64 big\sphinxhyphen{}endian binary file

\item {} 
\sphinxAtStartPar
\sphinxcode{\sphinxupquote{mat}} \sphinxhyphen{} MATLAB v7.2 save file without compression

\item {} 
\sphinxAtStartPar
\sphinxcode{\sphinxupquote{mat\_comp}} \sphinxhyphen{} MATLAB v7.2 save file with compression

\end{itemize}
\end{sphinxadmonition}


\subsubsection{Database Settings}
\label{\detokenize{usage/configuration:database-settings}}
\sphinxAtStartPar
The \sphinxstylestrong{Database} section controls the location where datasets are read from and where they will be written to.


\begin{savenotes}\sphinxattablestart
\sphinxthistablewithglobalstyle
\centering
\sphinxcapstartof{table}
\sphinxthecaptionisattop
\sphinxcaption{Key Database Section Variables}\label{\detokenize{usage/configuration:id2}}
\sphinxaftertopcaption
\begin{tabular}[t]{\X{1}{4}\X{2}{4}\X{1}{4}}
\sphinxtoprule
\sphinxstyletheadfamily 
\sphinxAtStartPar
Variable Name
&\sphinxstyletheadfamily 
\sphinxAtStartPar
Description
&\sphinxstyletheadfamily 
\sphinxAtStartPar
Type
\\
\sphinxmidrule
\sphinxtableatstartofbodyhook
\sphinxAtStartPar
\sphinxcode{\sphinxupquote{geoidDatabase}}
&
\sphinxAtStartPar
Directory for the geoid dataset files.
&
\sphinxAtStartPar
Filepath
\\
\sphinxhline
\sphinxAtStartPar
\sphinxcode{\sphinxupquote{gpsDatabase}}
&
\sphinxAtStartPar
Directory for the GPS Rinex dataset files.
&
\sphinxAtStartPar
Filepath
\\
\sphinxhline
\sphinxAtStartPar
\sphinxcode{\sphinxupquote{geosatGravityDatabase}}
&
\sphinxAtStartPar
Directory for the Geosat\sphinxhyphen{}44 gravity model files.
&
\sphinxAtStartPar
Filepath
\\
\sphinxhline
\sphinxAtStartPar
\sphinxcode{\sphinxupquote{marineGravityDatabase}}
&
\sphinxAtStartPar
Directory for the Marine gravity model files.
&
\sphinxAtStartPar
Filepath
\\
\sphinxhline
\sphinxAtStartPar
\sphinxcode{\sphinxupquote{ggmPlusGravityDatabase}}
&
\sphinxAtStartPar
Directory for the GGM Plus gravity model files.
&
\sphinxAtStartPar
Filepath
\\
\sphinxhline
\sphinxAtStartPar
\sphinxcode{\sphinxupquote{srtm2GravityDatabase}}
&
\sphinxAtStartPar
Directory for the SRTM2Gravity gravity model files.
&
\sphinxAtStartPar
Filepath
\\
\sphinxhline
\sphinxAtStartPar
\sphinxcode{\sphinxupquote{irishSeaGravityDatabase}}
&
\sphinxAtStartPar
Directory for the Irish Sea gravity model files.
&
\sphinxAtStartPar
Filepath
\\
\sphinxbottomrule
\end{tabular}
\sphinxtableafterendhook\par
\sphinxattableend\end{savenotes}


\subsubsection{Waypoints Settings}
\label{\detokenize{usage/configuration:waypoints-settings}}
\sphinxAtStartPar
The \sphinxstylestrong{Waypoints} section controls how the true trajectory is generated and used. This is where information of the
effective trajectory path is stated. This also controls the timings and frequency of the data.


\begin{savenotes}\sphinxattablestart
\sphinxthistablewithglobalstyle
\centering
\sphinxcapstartof{table}
\sphinxthecaptionisattop
\sphinxcaption{Key Waypoint Section Variables}\label{\detokenize{usage/configuration:id3}}
\sphinxaftertopcaption
\begin{tabular}[t]{\X{1}{4}\X{2}{4}\X{1}{4}}
\sphinxtoprule
\sphinxstyletheadfamily 
\sphinxAtStartPar
Variable Name
&\sphinxstyletheadfamily 
\sphinxAtStartPar
Description
&\sphinxstyletheadfamily 
\sphinxAtStartPar
Type
\\
\sphinxmidrule
\sphinxtableatstartofbodyhook
\sphinxAtStartPar
\sphinxcode{\sphinxupquote{waypointFrequency}}
&
\sphinxAtStartPar
The waypoint/simulated world’s update frequency to use.
&
\sphinxAtStartPar
Float
\\
\sphinxhline
\sphinxAtStartPar
\sphinxcode{\sphinxupquote{waypointInterpolation}}
&
\sphinxAtStartPar
The interpolation method to use for ground truth look\sphinxhyphen{}up.
&
\sphinxAtStartPar
cubic, pchip \sphinxstyleemphasis{or} akima
\\
\sphinxhline
\sphinxAtStartPar
\sphinxcode{\sphinxupquote{waypointMode}}
&
\sphinxAtStartPar
The mode for reading base waypoints.
&
\sphinxAtStartPar
config, csv \sphinxstyleemphasis{or} racetrack
\\
\sphinxhline
\sphinxAtStartPar
\sphinxcode{\sphinxupquote{waypointConfig}}
&
\sphinxAtStartPar
A waypoint sub\sphinxhyphen{}configuration file.
&
\sphinxAtStartPar
Yes \sphinxstyleemphasis{or} No
\\
\sphinxhline
\sphinxAtStartPar
\sphinxcode{\sphinxupquote{waypointCSV}}
&
\sphinxAtStartPar
A waypoint CVS file.
&
\sphinxAtStartPar
Yes \sphinxstyleemphasis{or} No
\\
\sphinxhline
\sphinxAtStartPar
\sphinxcode{\sphinxupquote{useVirtualMemory}}
&
\sphinxAtStartPar
Should disk space be used to hold waypoint data?
&
\sphinxAtStartPar
Yes \sphinxstyleemphasis{or} No
\\
\sphinxbottomrule
\end{tabular}
\sphinxtableafterendhook\par
\sphinxattableend\end{savenotes}


\subsubsection{Racetrack Settings}
\label{\detokenize{usage/configuration:racetrack-settings}}
\sphinxAtStartPar
The \sphinxstylestrong{Racetrack} section controls how arithmetical looping trajectories will be generated, if the waypoint method
is set to racetrack. This can be useful for quickly experimenting with different settings or profiles, without
having to present a planned route to use. The produced track is a looping trajectory which will be followed until
the simulations maximum time is reached.


\begin{savenotes}\sphinxattablestart
\sphinxthistablewithglobalstyle
\centering
\sphinxcapstartof{table}
\sphinxthecaptionisattop
\sphinxcaption{Key Waypoint Section Variables}\label{\detokenize{usage/configuration:id4}}
\sphinxaftertopcaption
\begin{tabular}[t]{\X{1}{4}\X{2}{4}\X{1}{4}}
\sphinxtoprule
\sphinxstyletheadfamily 
\sphinxAtStartPar
Variable Name
&\sphinxstyletheadfamily 
\sphinxAtStartPar
Description
&\sphinxstyletheadfamily 
\sphinxAtStartPar
Type
\\
\sphinxmidrule
\sphinxtableatstartofbodyhook
\sphinxAtStartPar
\sphinxcode{\sphinxupquote{circuitTime}}
&
\sphinxAtStartPar
The time (in seconds) to complete a single lap of the circuit.
&
\sphinxAtStartPar
Float
\\
\sphinxhline
\sphinxAtStartPar
\sphinxcode{\sphinxupquote{distanceX}}
&
\sphinxAtStartPar
The distance span in the X (North) direction (in metres).
&
\sphinxAtStartPar
Float
\\
\sphinxhline
\sphinxAtStartPar
\sphinxcode{\sphinxupquote{distanceY}}
&
\sphinxAtStartPar
The distance span in the Y (East) direction (in metres).
&
\sphinxAtStartPar
Float
\\
\sphinxhline
\sphinxAtStartPar
\sphinxcode{\sphinxupquote{distanceZ}}
&
\sphinxAtStartPar
The distance span in the Z (Down) direction (in metres).
&
\sphinxAtStartPar
Float
\\
\sphinxhline
\sphinxAtStartPar
\sphinxcode{\sphinxupquote{numVariationsY}}
&
\sphinxAtStartPar
The number of loops in the Y axis with respect to the X axis.
&
\sphinxAtStartPar
Float
\\
\sphinxhline
\sphinxAtStartPar
\sphinxcode{\sphinxupquote{numVariationsZ}}
&
\sphinxAtStartPar
The number of loops in the Z axis with respect to the X axis.
&
\sphinxAtStartPar
Float
\\
\sphinxbottomrule
\end{tabular}
\sphinxtableafterendhook\par
\sphinxattableend\end{savenotes}


\subsubsection{Vehicle Settings}
\label{\detokenize{usage/configuration:vehicle-settings}}
\sphinxAtStartPar
The \sphinxstylestrong{Vehicle} section lists all the values used for stating the limitations of the vehicle (model) that is to be use.
This values effect the shape of the requested trajectory, the manoeuvring limitations and the movement noise of the
vehicle.


\begin{savenotes}\sphinxattablestart
\sphinxthistablewithglobalstyle
\centering
\sphinxcapstartof{table}
\sphinxthecaptionisattop
\sphinxcaption{Key Vehicle Section Variables}\label{\detokenize{usage/configuration:id5}}
\sphinxaftertopcaption
\begin{tabular}[t]{\X{1}{4}\X{2}{4}\X{1}{4}}
\sphinxtoprule
\sphinxstyletheadfamily 
\sphinxAtStartPar
Variable Name
&\sphinxstyletheadfamily 
\sphinxAtStartPar
Description
&\sphinxstyletheadfamily 
\sphinxAtStartPar
Type
\\
\sphinxmidrule
\sphinxtableatstartofbodyhook
\sphinxAtStartPar
\sphinxcode{\sphinxupquote{vehicleConfig}}
&
\sphinxAtStartPar
The location of the vehicle sub\sphinxhyphen{}profile to use.
&
\sphinxAtStartPar
Filepath
\\
\sphinxbottomrule
\end{tabular}
\sphinxtableafterendhook\par
\sphinxattableend\end{savenotes}


\subsubsection{Gravity Settings}
\label{\detokenize{usage/configuration:gravity-settings}}
\sphinxAtStartPar
The \sphinxstylestrong{Gravity} section controls all gravity related calculations, specifically the gravity functions and gravity
correction maps used in both ground truth and estimation.


\begin{savenotes}\sphinxattablestart
\sphinxthistablewithglobalstyle
\centering
\sphinxcapstartof{table}
\sphinxthecaptionisattop
\sphinxcaption{Key Gravity Section Variables}\label{\detokenize{usage/configuration:id6}}
\sphinxaftertopcaption
\begin{tabular}[t]{\X{1}{4}\X{2}{4}\X{1}{4}}
\sphinxtoprule
\sphinxstyletheadfamily 
\sphinxAtStartPar
Variable Name
&\sphinxstyletheadfamily 
\sphinxAtStartPar
Description
&\sphinxstyletheadfamily 
\sphinxAtStartPar
Type
\\
\sphinxmidrule
\sphinxtableatstartofbodyhook
\sphinxAtStartPar
\sphinxcode{\sphinxupquote{autoDownload}}
&
\sphinxAtStartPar
Should missing gravity data be downloaded automatically?
&
\sphinxAtStartPar
Yes \sphinxstyleemphasis{or} No
\\
\sphinxhline
\sphinxAtStartPar
\sphinxcode{\sphinxupquote{trueGravityFunction}}
&
\sphinxAtStartPar
The gravity function to use for the ground truth.
&
\sphinxAtStartPar
See Below
\\
\sphinxhline
\sphinxAtStartPar
\sphinxcode{\sphinxupquote{trueGravityCorrectionMap}}
&
\sphinxAtStartPar
The optional gravity correction map to use for the ground truth.
&
\sphinxAtStartPar
See Below
\\
\sphinxhline
\sphinxAtStartPar
\sphinxcode{\sphinxupquote{estimatedGravityFunction}}
&
\sphinxAtStartPar
The gravity function to use for the simulation estimation.
&
\sphinxAtStartPar
See Below
\\
\sphinxhline
\sphinxAtStartPar
\sphinxcode{\sphinxupquote{estimatedGravityCorrectionMap}}
&
\sphinxAtStartPar
The optional gravity correction map to use for the simulation estimation.
&
\sphinxAtStartPar
See Below
\\
\sphinxbottomrule
\end{tabular}
\sphinxtableafterendhook\par
\sphinxattableend\end{savenotes}

\begin{sphinxadmonition}{note}{Note:}
\sphinxAtStartPar
Toolbox supports the following gravity function options:
\begin{itemize}
\item {} 
\sphinxAtStartPar
\sphinxcode{\sphinxupquote{fixed}} \sphinxhyphen{} Simple fixed constant value

\item {} 
\sphinxAtStartPar
\sphinxcode{\sphinxupquote{uniform}} \sphinxhyphen{} Uniform sphere with height interpolation

\item {} 
\sphinxAtStartPar
\sphinxcode{\sphinxupquote{somigliana}} \sphinxhyphen{} Somigliana formula for theoretical gravity

\item {} 
\sphinxAtStartPar
\sphinxcode{\sphinxupquote{wgs84}} \sphinxhyphen{} NIMA’s WGS\sphinxhyphen{}84 ellipsoid gravity model (used by MATLAB).

\end{itemize}
\end{sphinxadmonition}

\begin{sphinxadmonition}{note}{Note:}
\sphinxAtStartPar
Toolbox supports the following gravity correction map options:
\begin{itemize}
\item {} 
\sphinxAtStartPar
\sphinxcode{\sphinxupquote{none}} \sphinxhyphen{} Do not use a map, simple use the function alone.

\item {} 
\sphinxAtStartPar
\sphinxcode{\sphinxupquote{geoid}} \sphinxhyphen{} Use the EGM model (defined under Geoid section).

\item {} 
\sphinxAtStartPar
\sphinxcode{\sphinxupquote{geosat}} \sphinxhyphen{} Use the geosat\sphinxhyphen{}44 gravity correction map.

\item {} 
\sphinxAtStartPar
\sphinxcode{\sphinxupquote{marine}} \sphinxhyphen{} Use the marine gravity correction map.

\item {} 
\sphinxAtStartPar
\sphinxcode{\sphinxupquote{ggmPlusAcc}} \sphinxhyphen{} Use GGMPlus gravity acceleration map.

\item {} 
\sphinxAtStartPar
\sphinxcode{\sphinxupquote{ggmPlusDist}} \sphinxhyphen{} Use GGMPlus gravity disturbance map.

\item {} 
\sphinxAtStartPar
\sphinxcode{\sphinxupquote{srtm2gravityFS}} \sphinxhyphen{} Use SRTM2Gravity full\sphinxhyphen{}scale gravity map.

\item {} 
\sphinxAtStartPar
\sphinxcode{\sphinxupquote{srtm2gravityRes}} \sphinxhyphen{} Use SRTM2Gravity residual gravity map.

\item {} 
\sphinxAtStartPar
\sphinxcode{\sphinxupquote{irishSea}} \sphinxhyphen{} Use bespoke Irish sea gravity anomaly map.

\end{itemize}
\end{sphinxadmonition}


\subsubsection{Geoid Settings}
\label{\detokenize{usage/configuration:geoid-settings}}
\sphinxAtStartPar
The \sphinxstylestrong{Geoid} section controls all the selection of geoid model for the ground truth and simulation estimation,


\begin{savenotes}\sphinxattablestart
\sphinxthistablewithglobalstyle
\centering
\sphinxcapstartof{table}
\sphinxthecaptionisattop
\sphinxcaption{Key Geoid Section Variables}\label{\detokenize{usage/configuration:id7}}
\sphinxaftertopcaption
\begin{tabular}[t]{\X{1}{4}\X{2}{4}\X{1}{4}}
\sphinxtoprule
\sphinxstyletheadfamily 
\sphinxAtStartPar
Variable Name
&\sphinxstyletheadfamily 
\sphinxAtStartPar
Description
&\sphinxstyletheadfamily 
\sphinxAtStartPar
Type
\\
\sphinxmidrule
\sphinxtableatstartofbodyhook
\sphinxAtStartPar
\sphinxcode{\sphinxupquote{trueGeoidModel}}
&
\sphinxAtStartPar
The geoid model to apply when calculating the ground truth.
&
\sphinxAtStartPar
See Below
\\
\sphinxhline
\sphinxAtStartPar
\sphinxcode{\sphinxupquote{estimatedGeoidModel}}
&
\sphinxAtStartPar
The geoid model to apply during estimation simulations.
&
\sphinxAtStartPar
See Below
\\
\sphinxbottomrule
\end{tabular}
\sphinxtableafterendhook\par
\sphinxattableend\end{savenotes}

\begin{sphinxadmonition}{note}{Note:}
\sphinxAtStartPar
Toolbox supports the following geoid model options:
\begin{itemize}
\item {} 
\sphinxAtStartPar
\sphinxcode{\sphinxupquote{none}} \sphinxhyphen{} Do not use a geoid model

\item {} 
\sphinxAtStartPar
\sphinxcode{\sphinxupquote{egm84\sphinxhyphen{}30}} or \sphinxcode{\sphinxupquote{egm84\sphinxhyphen{}15}} \sphinxhyphen{} For the EGM 1984 model (at 30 or 15 minute resolutions).

\item {} 
\sphinxAtStartPar
\sphinxcode{\sphinxupquote{egm96\sphinxhyphen{}15}} or \sphinxcode{\sphinxupquote{egm96\sphinxhyphen{}5}} \sphinxhyphen{} For the EGM 1996 model (at 15 or 5 minute resolutions).

\item {} 
\sphinxAtStartPar
\sphinxcode{\sphinxupquote{egm2008\sphinxhyphen{}5}}, \sphinxcode{\sphinxupquote{egm2008\sphinxhyphen{}2\_5}} or \sphinxcode{\sphinxupquote{egm2008\sphinxhyphen{}1}} \sphinxhyphen{} For the EGM 2008 model (at 5, 2.5 or 1 minute resolutions).

\end{itemize}
\end{sphinxadmonition}


\subsubsection{Measurement Settings}
\label{\detokenize{usage/configuration:measurement-settings}}
\sphinxAtStartPar
The \sphinxstylestrong{Measurement} section contains mostly hardware configurations for the sensors to use. This includes stating
measurement frequencies and related errors.


\begin{savenotes}\sphinxattablestart
\sphinxthistablewithglobalstyle
\centering
\sphinxcapstartof{table}
\sphinxthecaptionisattop
\sphinxcaption{Key Measurement Section Variables}\label{\detokenize{usage/configuration:id8}}
\sphinxaftertopcaption
\begin{tabular}[t]{\X{1}{4}\X{2}{4}\X{1}{4}}
\sphinxtoprule
\sphinxstyletheadfamily 
\sphinxAtStartPar
Variable Name
&\sphinxstyletheadfamily 
\sphinxAtStartPar
Description
&\sphinxstyletheadfamily 
\sphinxAtStartPar
Type
\\
\sphinxmidrule
\sphinxtableatstartofbodyhook
\sphinxAtStartPar
\sphinxcode{\sphinxupquote{imuMeasurementFreq}}
&
\sphinxAtStartPar
The virtual hardware measurement frequency in Hertz.
&
\sphinxAtStartPar
Float
\\
\sphinxhline
\sphinxAtStartPar
\sphinxcode{\sphinxupquote{altimeterMeasurementFreq}}
&
\sphinxAtStartPar
The altimeter measurement frequency in Hertz (if enabled).
&
\sphinxAtStartPar
Float
\\
\sphinxhline
\sphinxAtStartPar
\sphinxcode{\sphinxupquote{gpsMeasurementFreq}}
&
\sphinxAtStartPar
The GPS measurement frequency in Hertz (if enabled).
&
\sphinxAtStartPar
Float
\\
\sphinxhline
\sphinxAtStartPar
\sphinxcode{\sphinxupquote{accelerometerConfig}}
&
\sphinxAtStartPar
The accelerometer profile to use for measurements.
&
\sphinxAtStartPar
Filepath
\\
\sphinxhline
\sphinxAtStartPar
\sphinxcode{\sphinxupquote{gyroscopeConfig}}
&
\sphinxAtStartPar
The accelerometer profile to use for measurements.
&
\sphinxAtStartPar
Filepath
\\
\sphinxhline
\sphinxAtStartPar
\sphinxcode{\sphinxupquote{altimeterConfig}}
&
\sphinxAtStartPar
The altimeter profile to use for measurements.
&
\sphinxAtStartPar
Filepath
\\
\sphinxhline
\sphinxAtStartPar
\sphinxcode{\sphinxupquote{clockConfig}}
&
\sphinxAtStartPar
The clock profile to use for measurements.
&
\sphinxAtStartPar
Filepath
\\
\sphinxbottomrule
\end{tabular}
\sphinxtableafterendhook\par
\sphinxattableend\end{savenotes}


\subsubsection{Estimation Settings}
\label{\detokenize{usage/configuration:estimation-settings}}
\sphinxAtStartPar
The \sphinxstylestrong{Estimation} section contains values relating to estimation methods that will be used for predicting the
vehicle’s current position, orientation, velocity and other factors.


\begin{savenotes}\sphinxattablestart
\sphinxthistablewithglobalstyle
\centering
\sphinxcapstartof{table}
\sphinxthecaptionisattop
\sphinxcaption{Key Estimation Section Variables}\label{\detokenize{usage/configuration:id9}}
\sphinxaftertopcaption
\begin{tabular}[t]{\X{1}{4}\X{2}{4}\X{1}{4}}
\sphinxtoprule
\sphinxstyletheadfamily 
\sphinxAtStartPar
Variable Name
&\sphinxstyletheadfamily 
\sphinxAtStartPar
Description
&\sphinxstyletheadfamily 
\sphinxAtStartPar
Type
\\
\sphinxmidrule
\sphinxtableatstartofbodyhook
\sphinxAtStartPar
\sphinxcode{\sphinxupquote{integrationMethod}}
&
\sphinxAtStartPar
The navigation integration method to use.
&
\sphinxAtStartPar
Integration ID
\\
\sphinxhline
\sphinxAtStartPar
\sphinxcode{\sphinxupquote{accelerometerConfig}}
&
\sphinxAtStartPar
The accelerometer profile to use for estimating errors.
&
\sphinxAtStartPar
Filepath
\\
\sphinxhline
\sphinxAtStartPar
\sphinxcode{\sphinxupquote{gyroscopeConfig}}
&
\sphinxAtStartPar
The gyroscope profile to use for estimating errors.
&
\sphinxAtStartPar
Filepath
\\
\sphinxhline
\sphinxAtStartPar
\sphinxcode{\sphinxupquote{altimeterConfig}}
&
\sphinxAtStartPar
The altimeter profile to use for estimating errors.
&
\sphinxAtStartPar
Filepath
\\
\sphinxhline
\sphinxAtStartPar
\sphinxcode{\sphinxupquote{clockConfig}}
&
\sphinxAtStartPar
The clock profile to use for estimating errors.
&
\sphinxAtStartPar
Filepath
\\
\sphinxhline
\sphinxAtStartPar
\sphinxcode{\sphinxupquote{gpsFusionMethod}}
&
\sphinxAtStartPar
The GPS fusion method to use for processing GPS data.
&
\sphinxAtStartPar
GPS Method ID
\\
\sphinxhline
\sphinxAtStartPar
\sphinxcode{\sphinxupquote{useVirtualMemory}}
&
\sphinxAtStartPar
Should disk space be used to hold estimation results?
&
\sphinxAtStartPar
Yes \sphinxstyleemphasis{or} No
\\
\sphinxbottomrule
\end{tabular}
\sphinxtableafterendhook\par
\sphinxattableend\end{savenotes}

\begin{sphinxadmonition}{note}{Note:}
\sphinxAtStartPar
The following INS integration methods are supported:
\begin{description}
\sphinxlineitem{Integration Methods:}
\sphinxAtStartPar
\sphinxcode{\sphinxupquote{Numerical}}, \sphinxcode{\sphinxupquote{Kalman}}, \sphinxcode{\sphinxupquote{Runge\_Kutta}}, \sphinxcode{\sphinxupquote{Adams\_Bashforth}}

\end{description}
\end{sphinxadmonition}


\subsubsection{Altimeter Settings}
\label{\detokenize{usage/configuration:altimeter-settings}}
\sphinxAtStartPar
The \sphinxstylestrong{Altimeter} section controls all altimeter related features. This include the altimeter frequency,
error profile and fusion method.


\begin{savenotes}\sphinxattablestart
\sphinxthistablewithglobalstyle
\centering
\sphinxcapstartof{table}
\sphinxthecaptionisattop
\sphinxcaption{Key Altimeter Section Variables}\label{\detokenize{usage/configuration:id10}}
\sphinxaftertopcaption
\begin{tabular}[t]{\X{1}{4}\X{2}{4}\X{1}{4}}
\sphinxtoprule
\sphinxstyletheadfamily 
\sphinxAtStartPar
Variable Name
&\sphinxstyletheadfamily 
\sphinxAtStartPar
Description
&\sphinxstyletheadfamily 
\sphinxAtStartPar
Type
\\
\sphinxmidrule
\sphinxtableatstartofbodyhook
\sphinxAtStartPar
\sphinxcode{\sphinxupquote{altimeterEnabled}}
&
\sphinxAtStartPar
Should the altimeter be used in the simulation?
&
\sphinxAtStartPar
Yes \sphinxstyleemphasis{or} No
\\
\sphinxhline
\sphinxAtStartPar
\sphinxcode{\sphinxupquote{altimeterFusion}}
&
\sphinxAtStartPar
Should the altimeter be used in the simulation?
&
\sphinxAtStartPar
See Below
\\
\sphinxhline
\sphinxAtStartPar
\sphinxcode{\sphinxupquote{altimeterGainAmount}}
&
\sphinxAtStartPar
The fixed gain amount to use for fixed\sphinxhyphen{}gain filtering.
&
\sphinxAtStartPar
Float
\\
\sphinxhline
\sphinxAtStartPar
\sphinxcode{\sphinxupquote{altimeterAlphaAmount}}
&
\sphinxAtStartPar
The alpha value to use for alpha\sphinxhyphen{}beta filtering.
&
\sphinxAtStartPar
Float
\\
\sphinxhline
\sphinxAtStartPar
\sphinxcode{\sphinxupquote{altimeterBetaAmount}}
&
\sphinxAtStartPar
The beta value to use for alpha\sphinxhyphen{}beta filtering.
&
\sphinxAtStartPar
Float
\\
\sphinxbottomrule
\end{tabular}
\sphinxtableafterendhook\par
\sphinxattableend\end{savenotes}

\begin{sphinxadmonition}{note}{Note:}
\sphinxAtStartPar
The following altimeter fusion methods are supported:
\begin{description}
\sphinxlineitem{Integration Methods:}
\sphinxAtStartPar
\sphinxcode{\sphinxupquote{fixed\_gain}} or \sphinxcode{\sphinxupquote{alpha\_beta}}

\end{description}
\end{sphinxadmonition}


\subsubsection{GPS Settings}
\label{\detokenize{usage/configuration:gps-settings}}
\sphinxAtStartPar
The \sphinxstylestrong{GPS} section controls all GPS related features. This includes if GPS is to be used in a simulation, the GPS
satellite data to use and areas where GPS signals are available.


\begin{savenotes}\sphinxattablestart
\sphinxthistablewithglobalstyle
\centering
\sphinxcapstartof{table}
\sphinxthecaptionisattop
\sphinxcaption{Key GPS Section Variables}\label{\detokenize{usage/configuration:id11}}
\sphinxaftertopcaption
\begin{tabular}[t]{\X{1}{4}\X{2}{4}\X{1}{4}}
\sphinxtoprule
\sphinxstyletheadfamily 
\sphinxAtStartPar
Variable Name
&\sphinxstyletheadfamily 
\sphinxAtStartPar
Description
&\sphinxstyletheadfamily 
\sphinxAtStartPar
Type
\\
\sphinxmidrule
\sphinxtableatstartofbodyhook
\sphinxAtStartPar
\sphinxcode{\sphinxupquote{gpsEnabled}}
&
\sphinxAtStartPar
Should the GPS be used in the simulation?
&
\sphinxAtStartPar
Yes \sphinxstyleemphasis{or} No
\\
\sphinxhline
\sphinxAtStartPar
\sphinxcode{\sphinxupquote{gpsFusionMethod}}
&
\sphinxAtStartPar
The GPS fusion method to use
&
\sphinxAtStartPar
See Below
\\
\sphinxhline
\sphinxAtStartPar
\sphinxcode{\sphinxupquote{gpsTimeYear}}
&
\sphinxAtStartPar
The 4\sphinxhyphen{}digit year to use for the satellite information.
&
\sphinxAtStartPar
Integer
\\
\sphinxhline
\sphinxAtStartPar
\sphinxcode{\sphinxupquote{gpsTimeMonth}}
&
\sphinxAtStartPar
The 2\sphinxhyphen{}digit month to use for the satellite information.
&
\sphinxAtStartPar
Integer
\\
\sphinxhline
\sphinxAtStartPar
\sphinxcode{\sphinxupquote{gpsTimeYear}}
&
\sphinxAtStartPar
The 2\sphinxhyphen{}digit day to use for the satellite information.
&
\sphinxAtStartPar
Integer
\\
\sphinxhline
\sphinxAtStartPar
\sphinxcode{\sphinxupquote{gpsAutoDownload}}
&
\sphinxAtStartPar
Can GPS information be automatically downloaded?
&
\sphinxAtStartPar
Yes \sphinxstyleemphasis{or} No
\\
\sphinxbottomrule
\end{tabular}
\sphinxtableafterendhook\par
\sphinxattableend\end{savenotes}

\begin{sphinxadmonition}{note}{Note:}
\sphinxAtStartPar
The following GPS fusion methods are supported:
\begin{description}
\sphinxlineitem{GPS Fusion Methods:}
\sphinxAtStartPar
\sphinxcode{\sphinxupquote{Fixed}}, \sphinxcode{\sphinxupquote{Loose}}, \sphinxcode{\sphinxupquote{Tight}}

\end{description}
\end{sphinxadmonition}


\subsubsection{Quantum Settings}
\label{\detokenize{usage/configuration:quantum-settings}}
\sphinxAtStartPar
The \sphinxstylestrong{Quantum} section controls all Quantum Sensor related features. This covers both the cold\sphinxhyphen{}atom quantum
inertia and gravity gradient sensors and fusion methods.


\begin{savenotes}\sphinxattablestart
\sphinxthistablewithglobalstyle
\centering
\sphinxcapstartof{table}
\sphinxthecaptionisattop
\sphinxcaption{Key Quantum Section Variables}\label{\detokenize{usage/configuration:id12}}
\sphinxaftertopcaption
\begin{tabular}[t]{\X{1}{4}\X{2}{4}\X{1}{4}}
\sphinxtoprule
\sphinxstyletheadfamily 
\sphinxAtStartPar
Variable Name
&\sphinxstyletheadfamily 
\sphinxAtStartPar
Description
&\sphinxstyletheadfamily 
\sphinxAtStartPar
Type
\\
\sphinxmidrule
\sphinxtableatstartofbodyhook
\sphinxAtStartPar
\sphinxcode{\sphinxupquote{quantumImuEnabled}}
&
\sphinxAtStartPar
Should the quantum IMU be used in the simulation?
&
\sphinxAtStartPar
Yes \sphinxstyleemphasis{or} No
\\
\sphinxhline
\sphinxAtStartPar
\sphinxcode{\sphinxupquote{quantumImuSensor}}
&
\sphinxAtStartPar
The type of quantum inertia sensor to use.
&
\sphinxAtStartPar
concept \sphinxstyleemphasis{or} realistic
\\
\sphinxhline
\sphinxAtStartPar
\sphinxcode{\sphinxupquote{quantumImuFusion}}
&
\sphinxAtStartPar
The fusion method to use for the quantum inertia sensor.
&
\sphinxAtStartPar
basic, particle\_filter \sphinxstyleemphasis{or} misaligned
\\
\sphinxhline
\sphinxAtStartPar
\sphinxcode{\sphinxupquote{quantumImuConfig}}
&
\sphinxAtStartPar
The quantum IMU sub\sphinxhyphen{}configuration file.
&
\sphinxAtStartPar
Filepath
\\
\sphinxhline
\sphinxAtStartPar
\sphinxcode{\sphinxupquote{quantumGravityEnabled}}
&
\sphinxAtStartPar
Should the quantum gravity map\sphinxhyphen{}matcher be used?
&
\sphinxAtStartPar
Yes \sphinxstyleemphasis{or} No
\\
\sphinxhline
\sphinxAtStartPar
\sphinxcode{\sphinxupquote{quantumGravitySensor}}
&
\sphinxAtStartPar
The type of quantum gravity sensor to use.
&
\sphinxAtStartPar
dual\_interferometer
\\
\sphinxhline
\sphinxAtStartPar
\sphinxcode{\sphinxupquote{quantumGravityFusion}}
&
\sphinxAtStartPar
The fusion method to use for the quantum gravity map\sphinxhyphen{}matching.
&
\sphinxAtStartPar
particle\_filter
\\
\sphinxhline
\sphinxAtStartPar
\sphinxcode{\sphinxupquote{quantumGravityConfig}}
&
\sphinxAtStartPar
The quantum gravity map\sphinxhyphen{}matching sub\sphinxhyphen{}configuration file.
&
\sphinxAtStartPar
Filepath
\\
\sphinxbottomrule
\end{tabular}
\sphinxtableafterendhook\par
\sphinxattableend\end{savenotes}


\subsubsection{Holonomic Settings}
\label{\detokenize{usage/configuration:holonomic-settings}}
\sphinxAtStartPar
The \sphinxstylestrong{Holonomics} section controls the optional use of holonomic constraints for preventing numerical
sensitivities impacting simulation accuracy. Essentially these are included to combat unwanted numerical rounding
errors in simulations.


\begin{savenotes}\sphinxattablestart
\sphinxthistablewithglobalstyle
\centering
\sphinxcapstartof{table}
\sphinxthecaptionisattop
\sphinxcaption{Key Holonomics Section Variables}\label{\detokenize{usage/configuration:id13}}
\sphinxaftertopcaption
\begin{tabulary}{\linewidth}[t]{TTT}
\sphinxtoprule
\sphinxstyletheadfamily 
\sphinxAtStartPar
Variable Name
&\sphinxstyletheadfamily 
\sphinxAtStartPar
Description
&\sphinxstyletheadfamily 
\sphinxAtStartPar
Type
\\
\sphinxmidrule
\sphinxtableatstartofbodyhook
\sphinxAtStartPar
\sphinxcode{\sphinxupquote{useHolonomics}}
&
\sphinxAtStartPar
Should holonomic constraints be used?
&
\sphinxAtStartPar
Yes \sphinxstyleemphasis{or} No
\\
\sphinxhline
\sphinxAtStartPar
\sphinxcode{\sphinxupquote{correctionFreq}}
&
\sphinxAtStartPar
The update rate at which corrections are applied (in Hz).
&
\sphinxAtStartPar
Float
\\
\sphinxhline
\sphinxAtStartPar
\sphinxcode{\sphinxupquote{positionGain}}
&
\sphinxAtStartPar
Holonomic gain for position.
&
\sphinxAtStartPar
Float
\\
\sphinxhline
\sphinxAtStartPar
\sphinxcode{\sphinxupquote{positionAxis}}
&
\sphinxAtStartPar
Holonomic axis percentages for position.
&
\sphinxAtStartPar
CSV (3\sphinxhyphen{}values)
\\
\sphinxhline
\sphinxAtStartPar
\sphinxcode{\sphinxupquote{velocityGain}}
&
\sphinxAtStartPar
Holonomic gain for velocity.
&
\sphinxAtStartPar
Float
\\
\sphinxhline
\sphinxAtStartPar
\sphinxcode{\sphinxupquote{velocityAxis}}
&
\sphinxAtStartPar
Holonomic axis percentages for position.
&
\sphinxAtStartPar
CSV (3\sphinxhyphen{}values)
\\
\sphinxhline
\sphinxAtStartPar
\sphinxcode{\sphinxupquote{accelerationGain}}
&
\sphinxAtStartPar
Holonomic gain for acceleration.
&
\sphinxAtStartPar
Float
\\
\sphinxhline
\sphinxAtStartPar
\sphinxcode{\sphinxupquote{accelerationAxis}}
&
\sphinxAtStartPar
Holonomic axis percentages for acceleration.
&
\sphinxAtStartPar
CSV (3\sphinxhyphen{}values)
\\
\sphinxhline
\sphinxAtStartPar
\sphinxcode{\sphinxupquote{attitudeGain}}
&
\sphinxAtStartPar
Holonomic gain for attitude.
&
\sphinxAtStartPar
Float
\\
\sphinxhline
\sphinxAtStartPar
\sphinxcode{\sphinxupquote{attitudeAxis}}
&
\sphinxAtStartPar
Holonomic axis percentages for attitude.
&
\sphinxAtStartPar
CSV (3\sphinxhyphen{}values)
\\
\sphinxhline
\sphinxAtStartPar
\sphinxcode{\sphinxupquote{angleRateGain}}
&
\sphinxAtStartPar
Holonomic gain for angle rates.
&
\sphinxAtStartPar
Float
\\
\sphinxhline
\sphinxAtStartPar
\sphinxcode{\sphinxupquote{angleRateAxis}}
&
\sphinxAtStartPar
Holonomic axis percentages for angle rates.
&
\sphinxAtStartPar
CSV (3\sphinxhyphen{}values)
\\
\sphinxbottomrule
\end{tabulary}
\sphinxtableafterendhook\par
\sphinxattableend\end{savenotes}


\subsubsection{Random Settings}
\label{\detokenize{usage/configuration:random-settings}}
\sphinxAtStartPar
The \sphinxstylestrong{Random} section controls the generation of all random numbers within waypoint generation, object initialisation
and during simulations. Essentially this section allows random aspects such as noise, drift, sensor errors etc. to be
managed via random seeds. If the same random seed is reused, the randomness will always be the same.


\begin{savenotes}\sphinxattablestart
\sphinxthistablewithglobalstyle
\centering
\sphinxcapstartof{table}
\sphinxthecaptionisattop
\sphinxcaption{Key Random Section Variables}\label{\detokenize{usage/configuration:id14}}
\sphinxaftertopcaption
\begin{tabulary}{\linewidth}[t]{TTT}
\sphinxtoprule
\sphinxstyletheadfamily 
\sphinxAtStartPar
Variable Name
&\sphinxstyletheadfamily 
\sphinxAtStartPar
Description
&\sphinxstyletheadfamily 
\sphinxAtStartPar
Type
\\
\sphinxmidrule
\sphinxtableatstartofbodyhook
\sphinxAtStartPar
\sphinxcode{\sphinxupquote{initialRandomSeed}}
&
\sphinxAtStartPar
The random seed to use during initialisation.
&
\sphinxAtStartPar
Integer
\\
\sphinxhline
\sphinxAtStartPar
\sphinxcode{\sphinxupquote{autoGenerateInitialSeed}}
&
\sphinxAtStartPar
Should the value be autogenerated?
&
\sphinxAtStartPar
Yes \sphinxstyleemphasis{or} No
\\
\sphinxhline
\sphinxAtStartPar
\sphinxcode{\sphinxupquote{waypointRandomSeed}}
&
\sphinxAtStartPar
The random seed to use during waypoint creation.
&
\sphinxAtStartPar
Integer
\\
\sphinxhline
\sphinxAtStartPar
\sphinxcode{\sphinxupquote{autoGenerateWaypointSeed}}
&
\sphinxAtStartPar
Should the value be autogenerated?
&
\sphinxAtStartPar
Yes \sphinxstyleemphasis{or} No
\\
\sphinxhline
\sphinxAtStartPar
\sphinxcode{\sphinxupquote{simulationRandomSeed}}
&
\sphinxAtStartPar
The random seed to use during simulations.
&
\sphinxAtStartPar
Integer
\\
\sphinxhline
\sphinxAtStartPar
\sphinxcode{\sphinxupquote{autoGenerateSimulationSeed}}
&
\sphinxAtStartPar
Should the value be autogenerated?
&
\sphinxAtStartPar
Yes \sphinxstyleemphasis{or} No
\\
\sphinxbottomrule
\end{tabulary}
\sphinxtableafterendhook\par
\sphinxattableend\end{savenotes}


\subsubsection{Misc Settings}
\label{\detokenize{usage/configuration:misc-settings}}
\sphinxAtStartPar
The \sphinxstylestrong{Misc} section contains the remaining miscellaneous settings that do not fit into any of the above sections.
These namely control the extraction process of values from configuration files and whether any automatic
behaviour is to be displayed.


\begin{savenotes}\sphinxattablestart
\sphinxthistablewithglobalstyle
\centering
\sphinxcapstartof{table}
\sphinxthecaptionisattop
\sphinxcaption{Key Misc Section Variables}\label{\detokenize{usage/configuration:id15}}
\sphinxaftertopcaption
\begin{tabulary}{\linewidth}[t]{TTT}
\sphinxtoprule
\sphinxstyletheadfamily 
\sphinxAtStartPar
Variable Name
&\sphinxstyletheadfamily 
\sphinxAtStartPar
Description
&\sphinxstyletheadfamily 
\sphinxAtStartPar
Type
\\
\sphinxmidrule
\sphinxtableatstartofbodyhook
\sphinxAtStartPar
\sphinxcode{\sphinxupquote{allowConfigOverloading}}
&
\sphinxAtStartPar
Should sub\sphinxhyphen{}config values overload those in the main config?
&
\sphinxAtStartPar
Yes \sphinxstyleemphasis{or} No
\\
\sphinxhline
\sphinxAtStartPar
\sphinxcode{\sphinxupquote{displayFallbackWarnings}}
&
\sphinxAtStartPar
Should warnings about missing variables be displayed?
&
\sphinxAtStartPar
Yes \sphinxstyleemphasis{or} No
\\
\sphinxbottomrule
\end{tabulary}
\sphinxtableafterendhook\par
\sphinxattableend\end{savenotes}

\begin{DUlineblock}{0em}
\item[] 
\end{DUlineblock}


\subsection{Sub\sphinxhyphen{}configuration Files}
\label{\detokenize{usage/configuration:sub-configuration-files}}
\sphinxAtStartPar
To reduce the complexity of the configuration file and to allow for a more modular design, some sections will
reference sub\sphinxhyphen{}configuration files which contain variables to use. For example, the variable \sphinxcode{\sphinxupquote{vehicleConfig}} points
to the vehicle sub\sphinxhyphen{}configuration file to use, allowing pre\sphinxhyphen{}created profiles to be re\sphinxhyphen{}used. For instance,
\sphinxcode{\sphinxupquote{vehicleConfig}} can be used to load values belonging to different vehicle models. This has many advantages
including allowing the changing of multiple parameters by only changing a single value.

\sphinxAtStartPar
When a sub\sphinxhyphen{}configuration file is loaded, the values extracted from it will be overwritten by pre\sphinxhyphen{}existing values in
the main configuration file. For instance, if \sphinxcode{\sphinxupquote{vehicleName}} is set to “\sphinxstylestrong{Cargo Plane}” in the main configuration
file, but is set to “\sphinxstylestrong{Commercial Plane}” in the vehicle sub\sphinxhyphen{}configuration file, “\sphinxstylestrong{Cargo Plane}” will be used.

\begin{sphinxadmonition}{note}{Note:}
\sphinxAtStartPar
This overwriting priority behaviour can be altered by changing the value of \sphinxcode{\sphinxupquote{allowConfigOverloading}}
(set to \sphinxcode{\sphinxupquote{false}} by default) located in the \sphinxcode{\sphinxupquote{Misc}} section.
\end{sphinxadmonition}

\sphinxstepscope


\section{Software Components}
\label{\detokenize{usage/packages:software-components}}\label{\detokenize{usage/packages::doc}}
\sphinxAtStartPar
The functionality of the produced software has been logically categorised into Python sub\sphinxhyphen{}packages. This allows means
for grouping related elements together and organising the project into manageable components (rather than large amount
of files and resources). The modules belonging to each sub\sphinxhyphen{}package of the project have been outlined below.


\subsection{Earth}
\label{\detokenize{usage/packages:earth}}
\sphinxAtStartPar
Modules relating to core calculations involving the earth:

\sphinxstepscope


\subsubsection{qnav.util.transformations}
\label{\detokenize{modules/util/transformations:module-qnav.util.transformations}}\label{\detokenize{modules/util/transformations:qnav-util-transformations}}\label{\detokenize{modules/util/transformations::doc}}\index{module@\spxentry{module}!qnav.util.transformations@\spxentry{qnav.util.transformations}}\index{qnav.util.transformations@\spxentry{qnav.util.transformations}!module@\spxentry{module}}

\paragraph{transformations.py}
\label{\detokenize{modules/util/transformations:transformations-py}}\begin{quote}\begin{description}
\sphinxlineitem{summary}
\sphinxAtStartPar
This module includes all general transformation calculations that
are required for the Navigation Simulation Study. The output from the
presented functions has been compared with other implementations such as
Python’s NavPy and MATLAB’s mapping toolbox.

\sphinxlineitem{authors}
\begin{DUlineblock}{0em}
\item[] Prof. Jason Ralph \sphinxhyphen{} \sphinxhref{mailto:jfralph@liverpool.ac.uk}{jfralph@liverpool.ac.uk}
\item[] Michael Wright \sphinxhyphen{} \sphinxhref{mailto:mjwright@liverpool.ac.uk}{mjwright@liverpool.ac.uk}
\end{DUlineblock}

\sphinxlineitem{version}
\sphinxAtStartPar
1.0.0 \sphinxhyphen{} First full implementation of this module.

\end{description}\end{quote}
\index{deg2rad() (in module qnav.util.transformations)@\spxentry{deg2rad()}\spxextra{in module qnav.util.transformations}}

\begin{fulllineitems}
\phantomsection\label{\detokenize{modules/util/transformations:qnav.util.transformations.deg2rad}}
\pysigstartsignatures
\pysiglinewithargsret{\sphinxcode{\sphinxupquote{qnav.util.transformations.}}\sphinxbfcode{\sphinxupquote{deg2rad}}}{\sphinxparam{\DUrole{n}{theta\_deg}\DUrole{p}{:}\DUrole{w}{ }\DUrole{n}{float}}}{{ $\rightarrow$ float}}
\pysigstopsignatures
\sphinxAtStartPar
A wrapper function for converting from degrees to radians.
\begin{quote}\begin{description}
\sphinxlineitem{Parameters}
\sphinxAtStartPar
\sphinxstyleliteralstrong{\sphinxupquote{theta\_deg}} (\sphinxstyleliteralemphasis{\sphinxupquote{float}}) \textendash{} An input value in degrees.

\sphinxlineitem{Returns}
\sphinxAtStartPar
The converted input value in radians.

\sphinxlineitem{Return type}
\sphinxAtStartPar
float

\end{description}\end{quote}

\end{fulllineitems}

\index{rad2deg() (in module qnav.util.transformations)@\spxentry{rad2deg()}\spxextra{in module qnav.util.transformations}}

\begin{fulllineitems}
\phantomsection\label{\detokenize{modules/util/transformations:qnav.util.transformations.rad2deg}}
\pysigstartsignatures
\pysiglinewithargsret{\sphinxcode{\sphinxupquote{qnav.util.transformations.}}\sphinxbfcode{\sphinxupquote{rad2deg}}}{\sphinxparam{\DUrole{n}{theta\_rad}\DUrole{p}{:}\DUrole{w}{ }\DUrole{n}{float}}}{{ $\rightarrow$ float}}
\pysigstopsignatures
\sphinxAtStartPar
A wrapper function for converting from radians to degrees.
\begin{quote}\begin{description}
\sphinxlineitem{Parameters}
\sphinxAtStartPar
\sphinxstyleliteralstrong{\sphinxupquote{theta\_rad}} (\sphinxstyleliteralemphasis{\sphinxupquote{float}}) \textendash{} An input value in radians.

\sphinxlineitem{Returns}
\sphinxAtStartPar
The converted input value in degrees.

\sphinxlineitem{Return type}
\sphinxAtStartPar
float

\end{description}\end{quote}

\end{fulllineitems}

\index{deg2dec() (in module qnav.util.transformations)@\spxentry{deg2dec()}\spxextra{in module qnav.util.transformations}}

\begin{fulllineitems}
\phantomsection\label{\detokenize{modules/util/transformations:qnav.util.transformations.deg2dec}}
\pysigstartsignatures
\pysiglinewithargsret{\sphinxcode{\sphinxupquote{qnav.util.transformations.}}\sphinxbfcode{\sphinxupquote{deg2dec}}}{\sphinxparam{\DUrole{n}{degrees}\DUrole{p}{:}\DUrole{w}{ }\DUrole{n}{float}}\sphinxparamcomma \sphinxparam{\DUrole{n}{minutes}\DUrole{p}{:}\DUrole{w}{ }\DUrole{n}{float}}\sphinxparamcomma \sphinxparam{\DUrole{n}{seconds}\DUrole{p}{:}\DUrole{w}{ }\DUrole{n}{float}}}{{ $\rightarrow$ float}}
\pysigstopsignatures
\sphinxAtStartPar
Converts a degrees, minutes, seconds coordinate to decimal format.
\begin{quote}\begin{description}
\sphinxlineitem{Parameters}\begin{itemize}
\item {} 
\sphinxAtStartPar
\sphinxstyleliteralstrong{\sphinxupquote{degrees}} (\sphinxstyleliteralemphasis{\sphinxupquote{float}}) \textendash{} The degrees of the given coordinate.

\item {} 
\sphinxAtStartPar
\sphinxstyleliteralstrong{\sphinxupquote{minutes}} (\sphinxstyleliteralemphasis{\sphinxupquote{float}}) \textendash{} The minutes of the given coordinate.

\item {} 
\sphinxAtStartPar
\sphinxstyleliteralstrong{\sphinxupquote{seconds}} (\sphinxstyleliteralemphasis{\sphinxupquote{float}}) \textendash{} The seconds of the given coordinate.

\end{itemize}

\sphinxlineitem{Returns}
\sphinxAtStartPar
The decimal equivalent of the given coordinate.

\sphinxlineitem{Return type}
\sphinxAtStartPar
float

\end{description}\end{quote}

\end{fulllineitems}

\index{dec2deg() (in module qnav.util.transformations)@\spxentry{dec2deg()}\spxextra{in module qnav.util.transformations}}

\begin{fulllineitems}
\phantomsection\label{\detokenize{modules/util/transformations:qnav.util.transformations.dec2deg}}
\pysigstartsignatures
\pysiglinewithargsret{\sphinxcode{\sphinxupquote{qnav.util.transformations.}}\sphinxbfcode{\sphinxupquote{dec2deg}}}{\sphinxparam{\DUrole{n}{theta\_deg}\DUrole{p}{:}\DUrole{w}{ }\DUrole{n}{float}}}{{ $\rightarrow$ tuple\DUrole{p}{{[}}float\DUrole{p}{,}\DUrole{w}{ }float\DUrole{p}{,}\DUrole{w}{ }float\DUrole{p}{{]}}}}
\pysigstopsignatures
\sphinxAtStartPar
Converts a decimal coordinate to degrees, minutes, seconds format.
\begin{quote}\begin{description}
\sphinxlineitem{Parameters}
\sphinxAtStartPar
\sphinxstyleliteralstrong{\sphinxupquote{theta\_deg}} (\sphinxstyleliteralemphasis{\sphinxupquote{float}}\sphinxstyleliteralemphasis{\sphinxupquote{ or }}\sphinxstyleliteralemphasis{\sphinxupquote{NumPy Array}}) \textendash{} The decimal format of the given coordinate.

\sphinxlineitem{Returns}
\sphinxAtStartPar
The degrees, minutes and seconds of the given coordinate.

\sphinxlineitem{Return type}
\sphinxAtStartPar
float or NumPy Array

\end{description}\end{quote}

\end{fulllineitems}

\index{radius() (in module qnav.util.transformations)@\spxentry{radius()}\spxextra{in module qnav.util.transformations}}

\begin{fulllineitems}
\phantomsection\label{\detokenize{modules/util/transformations:qnav.util.transformations.radius}}
\pysigstartsignatures
\pysiglinewithargsret{\sphinxcode{\sphinxupquote{qnav.util.transformations.}}\sphinxbfcode{\sphinxupquote{radius}}}{\sphinxparam{\DUrole{n}{lat}\DUrole{p}{:}\DUrole{w}{ }\DUrole{n}{float}}}{{ $\rightarrow$ float}}
\pysigstopsignatures
\sphinxAtStartPar
This function calculates the radius (in metres) from the centre of the
Earth for a given latitude position, using WGS84 parameters.
\begin{quote}\begin{description}
\sphinxlineitem{Parameters}
\sphinxAtStartPar
\sphinxstyleliteralstrong{\sphinxupquote{lat}} (\sphinxstyleliteralemphasis{\sphinxupquote{float}}) \textendash{} The latitude of the points of interest.

\sphinxlineitem{Returns}
\sphinxAtStartPar
The calculated radius in metres for the given position.

\sphinxlineitem{Return type}
\sphinxAtStartPar
float

\end{description}\end{quote}

\end{fulllineitems}

\index{radius\_vec() (in module qnav.util.transformations)@\spxentry{radius\_vec()}\spxextra{in module qnav.util.transformations}}

\begin{fulllineitems}
\phantomsection\label{\detokenize{modules/util/transformations:qnav.util.transformations.radius_vec}}
\pysigstartsignatures
\pysiglinewithargsret{\sphinxcode{\sphinxupquote{qnav.util.transformations.}}\sphinxbfcode{\sphinxupquote{radius\_vec}}}{\sphinxparam{\DUrole{n}{lat}\DUrole{p}{:}\DUrole{w}{ }\DUrole{n}{ndarray}}}{{ $\rightarrow$ ndarray}}
\pysigstopsignatures
\sphinxAtStartPar
The vectorised version of the
{\hyperref[\detokenize{modules/util/transformations:qnav.util.transformations.radius}]{\sphinxcrossref{\sphinxcode{\sphinxupquote{radius()}}}}} function. This
calculates the radius (in metres) from the centre of the Earth for a large
series of given latitude positions, using WGS84 parameters.
\begin{quote}\begin{description}
\sphinxlineitem{Parameters}
\sphinxAtStartPar
\sphinxstyleliteralstrong{\sphinxupquote{lat}} (\sphinxstyleliteralemphasis{\sphinxupquote{numpy.ndarray}}) \textendash{} The latitude of the points of interest.

\sphinxlineitem{Returns}
\sphinxAtStartPar
The calculated radius in metres for the given position.

\sphinxlineitem{Return type}
\sphinxAtStartPar
numpy.ndarray

\end{description}\end{quote}

\end{fulllineitems}

\index{rotation\_rate() (in module qnav.util.transformations)@\spxentry{rotation\_rate()}\spxextra{in module qnav.util.transformations}}

\begin{fulllineitems}
\phantomsection\label{\detokenize{modules/util/transformations:qnav.util.transformations.rotation_rate}}
\pysigstartsignatures
\pysiglinewithargsret{\sphinxcode{\sphinxupquote{qnav.util.transformations.}}\sphinxbfcode{\sphinxupquote{rotation\_rate}}}{\sphinxparam{\DUrole{n}{lat}\DUrole{p}{:}\DUrole{w}{ }\DUrole{n}{float}}}{{ $\rightarrow$ ndarray}}
\pysigstopsignatures
\sphinxAtStartPar
Returns the rotation rate of earth at given latitude position.
Calculates the rotation rate vector of the earth the given latitude.
\begin{quote}\begin{description}
\sphinxlineitem{Parameters}
\sphinxAtStartPar
\sphinxstyleliteralstrong{\sphinxupquote{lat}} (\sphinxstyleliteralemphasis{\sphinxupquote{float}}) \textendash{} The latitude position (in decimal degrees).

\sphinxlineitem{Returns}
\sphinxAtStartPar
The rotation vector in degrees/second.

\sphinxlineitem{Return type}
\sphinxAtStartPar
np.ndarray

\end{description}\end{quote}

\end{fulllineitems}

\index{rotation\_rate\_vec() (in module qnav.util.transformations)@\spxentry{rotation\_rate\_vec()}\spxextra{in module qnav.util.transformations}}

\begin{fulllineitems}
\phantomsection\label{\detokenize{modules/util/transformations:qnav.util.transformations.rotation_rate_vec}}
\pysigstartsignatures
\pysiglinewithargsret{\sphinxcode{\sphinxupquote{qnav.util.transformations.}}\sphinxbfcode{\sphinxupquote{rotation\_rate\_vec}}}{\sphinxparam{\DUrole{n}{lat}\DUrole{p}{:}\DUrole{w}{ }\DUrole{n}{ndarray}}}{{ $\rightarrow$ ndarray}}
\pysigstopsignatures
\sphinxAtStartPar
Returns the rotation rate of earth at given latitude position.
Calculates the rotation rate vector of the earth the given latitude.
\begin{quote}\begin{description}
\sphinxlineitem{Parameters}
\sphinxAtStartPar
\sphinxstyleliteralstrong{\sphinxupquote{lat}} (\sphinxstyleliteralemphasis{\sphinxupquote{np.ndarray}}) \textendash{} The latitude position (in decimal degrees).

\sphinxlineitem{Returns}
\sphinxAtStartPar
The rotation vector in degrees/second.

\sphinxlineitem{Return type}
\sphinxAtStartPar
np.ndarray

\end{description}\end{quote}

\end{fulllineitems}

\index{radius84() (in module qnav.util.transformations)@\spxentry{radius84()}\spxextra{in module qnav.util.transformations}}

\begin{fulllineitems}
\phantomsection\label{\detokenize{modules/util/transformations:qnav.util.transformations.radius84}}
\pysigstartsignatures
\pysiglinewithargsret{\sphinxcode{\sphinxupquote{qnav.util.transformations.}}\sphinxbfcode{\sphinxupquote{radius84}}}{\sphinxparam{\DUrole{n}{lat}\DUrole{p}{:}\DUrole{w}{ }\DUrole{n}{float}}}{{ $\rightarrow$ float}}
\pysigstopsignatures
\sphinxAtStartPar
This function calculates the radius (in metres) from the centre of the
Earth for a given latitude position, using WGS84 parameters.
\begin{quote}\begin{description}
\sphinxlineitem{Parameters}
\sphinxAtStartPar
\sphinxstyleliteralstrong{\sphinxupquote{lat}} (\sphinxstyleliteralemphasis{\sphinxupquote{float}}) \textendash{} The latitude of the points of interest.

\sphinxlineitem{Returns}
\sphinxAtStartPar
The calculated radius in metres for the given position.

\sphinxlineitem{Return type}
\sphinxAtStartPar
float

\end{description}\end{quote}

\end{fulllineitems}

\index{radius84\_vec() (in module qnav.util.transformations)@\spxentry{radius84\_vec()}\spxextra{in module qnav.util.transformations}}

\begin{fulllineitems}
\phantomsection\label{\detokenize{modules/util/transformations:qnav.util.transformations.radius84_vec}}
\pysigstartsignatures
\pysiglinewithargsret{\sphinxcode{\sphinxupquote{qnav.util.transformations.}}\sphinxbfcode{\sphinxupquote{radius84\_vec}}}{\sphinxparam{\DUrole{n}{lat}\DUrole{p}{:}\DUrole{w}{ }\DUrole{n}{ndarray}}}{{ $\rightarrow$ ndarray}}
\pysigstopsignatures
\sphinxAtStartPar
The vectorised version of the
{\hyperref[\detokenize{modules/util/transformations:qnav.util.transformations.radius84}]{\sphinxcrossref{\sphinxcode{\sphinxupquote{radius84()}}}}} function. This
calculates the radius (in metres) from the centre of the Earth for a large
series of given latitude positions, using WGS84 parameters.
\begin{quote}\begin{description}
\sphinxlineitem{Parameters}
\sphinxAtStartPar
\sphinxstyleliteralstrong{\sphinxupquote{lat}} (\sphinxstyleliteralemphasis{\sphinxupquote{numpy.ndarray}}) \textendash{} The latitude of the points of interest.

\sphinxlineitem{Returns}
\sphinxAtStartPar
The calculated radius in metres for the given position.

\sphinxlineitem{Return type}
\sphinxAtStartPar
numpy.ndarray

\end{description}\end{quote}

\end{fulllineitems}

\index{rotate\_3d() (in module qnav.util.transformations)@\spxentry{rotate\_3d()}\spxextra{in module qnav.util.transformations}}

\begin{fulllineitems}
\phantomsection\label{\detokenize{modules/util/transformations:qnav.util.transformations.rotate_3d}}
\pysigstartsignatures
\pysiglinewithargsret{\sphinxcode{\sphinxupquote{qnav.util.transformations.}}\sphinxbfcode{\sphinxupquote{rotate\_3d}}}{\sphinxparam{\DUrole{n}{psi}\DUrole{p}{:}\DUrole{w}{ }\DUrole{n}{float}}\sphinxparamcomma \sphinxparam{\DUrole{n}{theta}\DUrole{p}{:}\DUrole{w}{ }\DUrole{n}{float}}\sphinxparamcomma \sphinxparam{\DUrole{n}{phi}\DUrole{p}{:}\DUrole{w}{ }\DUrole{n}{float}}}{{ $\rightarrow$ ndarray}}
\pysigstopsignatures
\sphinxAtStartPar
This function produces a three\sphinxhyphen{}dimensional rotation matrix corresponding
to the Euler angles psi (Heading), theta (Pitch), phi (Roll/Bank).
\begin{quote}\begin{description}
\sphinxlineitem{Parameters}\begin{itemize}
\item {} 
\sphinxAtStartPar
\sphinxstyleliteralstrong{\sphinxupquote{psi}} (\sphinxstyleliteralemphasis{\sphinxupquote{float}}) \textendash{} The heading angle in radians.

\item {} 
\sphinxAtStartPar
\sphinxstyleliteralstrong{\sphinxupquote{theta}} (\sphinxstyleliteralemphasis{\sphinxupquote{float}}) \textendash{} The pitch angle in radians.

\item {} 
\sphinxAtStartPar
\sphinxstyleliteralstrong{\sphinxupquote{phi}} (\sphinxstyleliteralemphasis{\sphinxupquote{float}}) \textendash{} The roll/bank angle in radians.

\end{itemize}

\sphinxlineitem{Returns}
\sphinxAtStartPar
A matrix corresponding to three rotations in the above order.

\sphinxlineitem{Return type}
\sphinxAtStartPar
numpy.ndarray (3\sphinxhyphen{}by\sphinxhyphen{}3 elements)

\end{description}\end{quote}

\end{fulllineitems}

\index{rotate\_3d\_vec() (in module qnav.util.transformations)@\spxentry{rotate\_3d\_vec()}\spxextra{in module qnav.util.transformations}}

\begin{fulllineitems}
\phantomsection\label{\detokenize{modules/util/transformations:qnav.util.transformations.rotate_3d_vec}}
\pysigstartsignatures
\pysiglinewithargsret{\sphinxcode{\sphinxupquote{qnav.util.transformations.}}\sphinxbfcode{\sphinxupquote{rotate\_3d\_vec}}}{\sphinxparam{\DUrole{n}{psi}\DUrole{p}{:}\DUrole{w}{ }\DUrole{n}{ndarray}}\sphinxparamcomma \sphinxparam{\DUrole{n}{theta}\DUrole{p}{:}\DUrole{w}{ }\DUrole{n}{ndarray}}\sphinxparamcomma \sphinxparam{\DUrole{n}{phi}\DUrole{p}{:}\DUrole{w}{ }\DUrole{n}{ndarray}}}{{ $\rightarrow$ ndarray}}
\pysigstopsignatures
\sphinxAtStartPar
This function is the vectorised version of the
{\hyperref[\detokenize{modules/util/transformations:qnav.util.transformations.rotate_3d}]{\sphinxcrossref{\sphinxcode{\sphinxupquote{rotate\_3d()}}}}} function. This
produces multiple three\sphinxhyphen{}dimensional rotation matrices corresponding
to the Euler angles psi (Heading), theta (Pitch), phi (Roll/Bank).
\begin{quote}\begin{description}
\sphinxlineitem{Parameters}\begin{itemize}
\item {} 
\sphinxAtStartPar
\sphinxstyleliteralstrong{\sphinxupquote{psi}} (\sphinxstyleliteralemphasis{\sphinxupquote{numpy.ndarray}}) \textendash{} The heading angles in radians.

\item {} 
\sphinxAtStartPar
\sphinxstyleliteralstrong{\sphinxupquote{theta}} (\sphinxstyleliteralemphasis{\sphinxupquote{numpy.ndarray}}) \textendash{} The pitch angles in radians.

\item {} 
\sphinxAtStartPar
\sphinxstyleliteralstrong{\sphinxupquote{phi}} (\sphinxstyleliteralemphasis{\sphinxupquote{numpy.ndarray}}) \textendash{} The roll/bank angles in radians.

\end{itemize}

\sphinxlineitem{Returns}
\sphinxAtStartPar
The matrices corresponding to three rotations in the above order.

\sphinxlineitem{Return type}
\sphinxAtStartPar
numpy.ndarray (n\sphinxhyphen{}3\sphinxhyphen{}by\sphinxhyphen{}3 elements)

\end{description}\end{quote}

\end{fulllineitems}

\index{lla2ned() (in module qnav.util.transformations)@\spxentry{lla2ned()}\spxextra{in module qnav.util.transformations}}

\begin{fulllineitems}
\phantomsection\label{\detokenize{modules/util/transformations:qnav.util.transformations.lla2ned}}
\pysigstartsignatures
\pysiglinewithargsret{\sphinxcode{\sphinxupquote{qnav.util.transformations.}}\sphinxbfcode{\sphinxupquote{lla2ned}}}{\sphinxparam{\DUrole{n}{lla}\DUrole{p}{:}\DUrole{w}{ }\DUrole{n}{Buffer\DUrole{w}{ }\DUrole{p}{|}\DUrole{w}{ }\_SupportsArray\DUrole{p}{{[}}dtype\DUrole{p}{{[}}Any\DUrole{p}{{]}}\DUrole{p}{{]}}\DUrole{w}{ }\DUrole{p}{|}\DUrole{w}{ }\_NestedSequence\DUrole{p}{{[}}\_SupportsArray\DUrole{p}{{[}}dtype\DUrole{p}{{[}}Any\DUrole{p}{{]}}\DUrole{p}{{]}}\DUrole{p}{{]}}\DUrole{w}{ }\DUrole{p}{|}\DUrole{w}{ }bool\DUrole{w}{ }\DUrole{p}{|}\DUrole{w}{ }int\DUrole{w}{ }\DUrole{p}{|}\DUrole{w}{ }float\DUrole{w}{ }\DUrole{p}{|}\DUrole{w}{ }complex\DUrole{w}{ }\DUrole{p}{|}\DUrole{w}{ }str\DUrole{w}{ }\DUrole{p}{|}\DUrole{w}{ }bytes\DUrole{w}{ }\DUrole{p}{|}\DUrole{w}{ }\_NestedSequence\DUrole{p}{{[}}bool\DUrole{w}{ }\DUrole{p}{|}\DUrole{w}{ }int\DUrole{w}{ }\DUrole{p}{|}\DUrole{w}{ }float\DUrole{w}{ }\DUrole{p}{|}\DUrole{w}{ }complex\DUrole{w}{ }\DUrole{p}{|}\DUrole{w}{ }str\DUrole{w}{ }\DUrole{p}{|}\DUrole{w}{ }bytes\DUrole{p}{{]}}}}\sphinxparamcomma \sphinxparam{\DUrole{n}{ref\_lla}\DUrole{p}{:}\DUrole{w}{ }\DUrole{n}{Buffer\DUrole{w}{ }\DUrole{p}{|}\DUrole{w}{ }\_SupportsArray\DUrole{p}{{[}}dtype\DUrole{p}{{[}}Any\DUrole{p}{{]}}\DUrole{p}{{]}}\DUrole{w}{ }\DUrole{p}{|}\DUrole{w}{ }\_NestedSequence\DUrole{p}{{[}}\_SupportsArray\DUrole{p}{{[}}dtype\DUrole{p}{{[}}Any\DUrole{p}{{]}}\DUrole{p}{{]}}\DUrole{p}{{]}}\DUrole{w}{ }\DUrole{p}{|}\DUrole{w}{ }bool\DUrole{w}{ }\DUrole{p}{|}\DUrole{w}{ }int\DUrole{w}{ }\DUrole{p}{|}\DUrole{w}{ }float\DUrole{w}{ }\DUrole{p}{|}\DUrole{w}{ }complex\DUrole{w}{ }\DUrole{p}{|}\DUrole{w}{ }str\DUrole{w}{ }\DUrole{p}{|}\DUrole{w}{ }bytes\DUrole{w}{ }\DUrole{p}{|}\DUrole{w}{ }\_NestedSequence\DUrole{p}{{[}}bool\DUrole{w}{ }\DUrole{p}{|}\DUrole{w}{ }int\DUrole{w}{ }\DUrole{p}{|}\DUrole{w}{ }float\DUrole{w}{ }\DUrole{p}{|}\DUrole{w}{ }complex\DUrole{w}{ }\DUrole{p}{|}\DUrole{w}{ }str\DUrole{w}{ }\DUrole{p}{|}\DUrole{w}{ }bytes\DUrole{p}{{]}}}}}{{ $\rightarrow$ ndarray}}
\pysigstopsignatures
\sphinxAtStartPar
This function converts a reference Latitude\sphinxhyphen{}Longitude\sphinxhyphen{}Altitude (LLA)
position from a given centre LLA point and translates its location to
North\sphinxhyphen{}East\sphinxhyphen{}Down (NED) coordinates.
\begin{quote}\begin{description}
\sphinxlineitem{Parameters}\begin{itemize}
\item {} 
\sphinxAtStartPar
\sphinxstyleliteralstrong{\sphinxupquote{lla}} (\sphinxstyleliteralemphasis{\sphinxupquote{numpy.ndarray}}\sphinxstyleliteralemphasis{\sphinxupquote{ (}}\sphinxstyleliteralemphasis{\sphinxupquote{3\sphinxhyphen{}elements}}\sphinxstyleliteralemphasis{\sphinxupquote{)}}) \textendash{} The centre location (array containing {[}lat, lon, alt{]}) with
lat\sphinxhyphen{}lon given in degrees and alt given in metres.

\item {} 
\sphinxAtStartPar
\sphinxstyleliteralstrong{\sphinxupquote{ref\_lla}} (\sphinxstyleliteralemphasis{\sphinxupquote{numpy.ndarray}}\sphinxstyleliteralemphasis{\sphinxupquote{ (}}\sphinxstyleliteralemphasis{\sphinxupquote{3\sphinxhyphen{}elements}}\sphinxstyleliteralemphasis{\sphinxupquote{)}}) \textendash{} The reference location (array containing {[}lat, lon, alt{]})
with lat\sphinxhyphen{}lon given in degrees and alt given in metres.

\end{itemize}

\sphinxlineitem{Returns}
\sphinxAtStartPar
An array containing the calculated North, East and Down
coordinates from the centre position, were all values are given in
metres.

\sphinxlineitem{Return type}
\sphinxAtStartPar
numpy.ndarray (3\sphinxhyphen{}elements)

\end{description}\end{quote}

\end{fulllineitems}

\index{lla2ned\_vec() (in module qnav.util.transformations)@\spxentry{lla2ned\_vec()}\spxextra{in module qnav.util.transformations}}

\begin{fulllineitems}
\phantomsection\label{\detokenize{modules/util/transformations:qnav.util.transformations.lla2ned_vec}}
\pysigstartsignatures
\pysiglinewithargsret{\sphinxcode{\sphinxupquote{qnav.util.transformations.}}\sphinxbfcode{\sphinxupquote{lla2ned\_vec}}}{\sphinxparam{\DUrole{n}{lla}\DUrole{p}{:}\DUrole{w}{ }\DUrole{n}{ndarray}}\sphinxparamcomma \sphinxparam{\DUrole{n}{ref\_lla}\DUrole{p}{:}\DUrole{w}{ }\DUrole{n}{ndarray}}}{{ $\rightarrow$ ndarray}}
\pysigstopsignatures
\sphinxAtStartPar
This function is the vectorised version of the
{\hyperref[\detokenize{modules/util/transformations:qnav.util.transformations.lla2ned}]{\sphinxcrossref{\sphinxcode{\sphinxupquote{lla2ned()}}}}} function. This
allows means for converting reference Latitude\sphinxhyphen{}Longitude\sphinxhyphen{}Altitude (LLA)
positions from given centre LLA points and translates their location to
North\sphinxhyphen{}East\sphinxhyphen{}Down (NED) coordinates.
\begin{quote}\begin{description}
\sphinxlineitem{Parameters}\begin{itemize}
\item {} 
\sphinxAtStartPar
\sphinxstyleliteralstrong{\sphinxupquote{lla}} (\sphinxstyleliteralemphasis{\sphinxupquote{numpy.ndarray}}\sphinxstyleliteralemphasis{\sphinxupquote{ (}}\sphinxstyleliteralemphasis{\sphinxupquote{n\sphinxhyphen{}by\sphinxhyphen{}3 elements}}\sphinxstyleliteralemphasis{\sphinxupquote{)}}) \textendash{} The centre locations (array rows containing {[}lat, lon, alt{]})
with lat\sphinxhyphen{}lon given in degrees and alt given in metres.

\item {} 
\sphinxAtStartPar
\sphinxstyleliteralstrong{\sphinxupquote{ref\_lla}} (\sphinxstyleliteralemphasis{\sphinxupquote{numpy.ndarray}}\sphinxstyleliteralemphasis{\sphinxupquote{ (}}\sphinxstyleliteralemphasis{\sphinxupquote{n\sphinxhyphen{}by\sphinxhyphen{}3 elements}}\sphinxstyleliteralemphasis{\sphinxupquote{)}}) \textendash{} The reference location (array rows containing
{[}lat, lon, alt{]}) with lat\sphinxhyphen{}lon given in degrees and alt given in
metres.

\end{itemize}

\sphinxlineitem{Returns}
\sphinxAtStartPar
An array containing the calculated North, East and Down
coordinates from the centre position, were all values are given in
metres.

\sphinxlineitem{Return type}
\sphinxAtStartPar
numpy.ndarray (n\sphinxhyphen{}by\sphinxhyphen{}3 elements)

\end{description}\end{quote}

\end{fulllineitems}

\index{ned2lla() (in module qnav.util.transformations)@\spxentry{ned2lla()}\spxextra{in module qnav.util.transformations}}

\begin{fulllineitems}
\phantomsection\label{\detokenize{modules/util/transformations:qnav.util.transformations.ned2lla}}
\pysigstartsignatures
\pysiglinewithargsret{\sphinxcode{\sphinxupquote{qnav.util.transformations.}}\sphinxbfcode{\sphinxupquote{ned2lla}}}{\sphinxparam{\DUrole{n}{ned}\DUrole{p}{:}\DUrole{w}{ }\DUrole{n}{ndarray}}\sphinxparamcomma \sphinxparam{\DUrole{n}{ref\_lla}\DUrole{p}{:}\DUrole{w}{ }\DUrole{n}{ndarray}}}{{ $\rightarrow$ ndarray}}
\pysigstopsignatures
\sphinxAtStartPar
This function converts a reference North\sphinxhyphen{}East\sphinxhyphen{}Down (NED) position from a
given centre position and translates its location to Latitude\sphinxhyphen{}Longitude\sphinxhyphen{}
Altitude (LLA) coordinates.
\begin{quote}\begin{description}
\sphinxlineitem{Parameters}\begin{itemize}
\item {} 
\sphinxAtStartPar
\sphinxstyleliteralstrong{\sphinxupquote{ned}} (\sphinxstyleliteralemphasis{\sphinxupquote{numpy.ndarray}}\sphinxstyleliteralemphasis{\sphinxupquote{ (}}\sphinxstyleliteralemphasis{\sphinxupquote{3\sphinxhyphen{}elements}}\sphinxstyleliteralemphasis{\sphinxupquote{)}}) \textendash{} The centre location (array containing {[}north, east, down{]})
with all values being given in metres.

\item {} 
\sphinxAtStartPar
\sphinxstyleliteralstrong{\sphinxupquote{ref\_lla}} (\sphinxstyleliteralemphasis{\sphinxupquote{numpy.ndarray}}\sphinxstyleliteralemphasis{\sphinxupquote{ (}}\sphinxstyleliteralemphasis{\sphinxupquote{3\sphinxhyphen{}elements}}\sphinxstyleliteralemphasis{\sphinxupquote{)}}) \textendash{} The reference location (array containing {[}lat, lon, alt{]})
with lat\sphinxhyphen{}lon given in degrees and alt given in metres.

\end{itemize}

\sphinxlineitem{Returns}
\sphinxAtStartPar
An array containing the calculated Latitude, Longitude and
Altitude coordinates from the centre position, with lat\sphinxhyphen{}lon is given
in degrees and alt given in metres.

\sphinxlineitem{Return type}
\sphinxAtStartPar
numpy.ndarray (3\sphinxhyphen{}elements)

\end{description}\end{quote}

\end{fulllineitems}

\index{ned2lla\_vec() (in module qnav.util.transformations)@\spxentry{ned2lla\_vec()}\spxextra{in module qnav.util.transformations}}

\begin{fulllineitems}
\phantomsection\label{\detokenize{modules/util/transformations:qnav.util.transformations.ned2lla_vec}}
\pysigstartsignatures
\pysiglinewithargsret{\sphinxcode{\sphinxupquote{qnav.util.transformations.}}\sphinxbfcode{\sphinxupquote{ned2lla\_vec}}}{\sphinxparam{\DUrole{n}{ned}\DUrole{p}{:}\DUrole{w}{ }\DUrole{n}{ndarray}}\sphinxparamcomma \sphinxparam{\DUrole{n}{ref\_lla}\DUrole{p}{:}\DUrole{w}{ }\DUrole{n}{ndarray}}}{{ $\rightarrow$ ndarray}}
\pysigstopsignatures
\sphinxAtStartPar
This function is the vectorised version of the
{\hyperref[\detokenize{modules/util/transformations:qnav.util.transformations.ned2lla}]{\sphinxcrossref{\sphinxcode{\sphinxupquote{ned2lla()}}}}} function. This
converts reference North\sphinxhyphen{}East\sphinxhyphen{}Down (NED) positions from a given centre
positions and translates their location to Latitude\sphinxhyphen{}Longitude\sphinxhyphen{}Altitude
(LLA) coordinates.
\begin{quote}\begin{description}
\sphinxlineitem{Parameters}\begin{itemize}
\item {} 
\sphinxAtStartPar
\sphinxstyleliteralstrong{\sphinxupquote{ned}} (\sphinxstyleliteralemphasis{\sphinxupquote{numpy.ndarray}}\sphinxstyleliteralemphasis{\sphinxupquote{ (}}\sphinxstyleliteralemphasis{\sphinxupquote{n\sphinxhyphen{}by\sphinxhyphen{}3 elements}}\sphinxstyleliteralemphasis{\sphinxupquote{)}}) \textendash{} The centre locations (array rows containing
{[}north, east, down{]}) with all values being given in metres.

\item {} 
\sphinxAtStartPar
\sphinxstyleliteralstrong{\sphinxupquote{ref\_lla}} (\sphinxstyleliteralemphasis{\sphinxupquote{numpy.ndarray}}\sphinxstyleliteralemphasis{\sphinxupquote{ (}}\sphinxstyleliteralemphasis{\sphinxupquote{n\sphinxhyphen{}by\sphinxhyphen{}3 elements}}\sphinxstyleliteralemphasis{\sphinxupquote{)}}) \textendash{} The reference location (array rows containing
{[}lat, lon, alt{]}) with lat\sphinxhyphen{}lon given in degrees and alt given in
metres.

\end{itemize}

\sphinxlineitem{Returns}
\sphinxAtStartPar
Array rows containing the calculated Latitude, Longitude and
Altitude coordinates from the centre positions, with lat\sphinxhyphen{}lon is given
in degrees and alt given in metres.

\sphinxlineitem{Return type}
\sphinxAtStartPar
numpy.ndarray (n\sphinxhyphen{}by\sphinxhyphen{}3 elements)

\end{description}\end{quote}

\end{fulllineitems}

\index{lla2ecef() (in module qnav.util.transformations)@\spxentry{lla2ecef()}\spxextra{in module qnav.util.transformations}}

\begin{fulllineitems}
\phantomsection\label{\detokenize{modules/util/transformations:qnav.util.transformations.lla2ecef}}
\pysigstartsignatures
\pysiglinewithargsret{\sphinxcode{\sphinxupquote{qnav.util.transformations.}}\sphinxbfcode{\sphinxupquote{lla2ecef}}}{\sphinxparam{\DUrole{n}{lla}\DUrole{p}{:}\DUrole{w}{ }\DUrole{n}{ndarray}}}{{ $\rightarrow$ ndarray}}
\pysigstopsignatures
\sphinxAtStartPar
This function converts a reference Latitude\sphinxhyphen{}Longitude\sphinxhyphen{}Altitude (LLA)
position to Earth\sphinxhyphen{}Centred Earth\sphinxhyphen{}Fixed coordinates.
\begin{quote}\begin{description}
\sphinxlineitem{Parameters}
\sphinxAtStartPar
\sphinxstyleliteralstrong{\sphinxupquote{lla}} (\sphinxstyleliteralemphasis{\sphinxupquote{numpy.ndarray}}\sphinxstyleliteralemphasis{\sphinxupquote{ (}}\sphinxstyleliteralemphasis{\sphinxupquote{3\sphinxhyphen{}elements}}\sphinxstyleliteralemphasis{\sphinxupquote{)}}) \textendash{} The reference location (array containing {[}lat, lon, alt{]}) with
lat\sphinxhyphen{}lon given in degrees and alt given in metres.

\sphinxlineitem{Returns}
\sphinxAtStartPar
An array containing the calculated X, Y and Z coordinates from
the centre of the Earth, with all values being given in metres.

\sphinxlineitem{Return type}
\sphinxAtStartPar
numpy.ndarray (3\sphinxhyphen{}elements)

\end{description}\end{quote}

\end{fulllineitems}

\index{lla2ecef\_vec() (in module qnav.util.transformations)@\spxentry{lla2ecef\_vec()}\spxextra{in module qnav.util.transformations}}

\begin{fulllineitems}
\phantomsection\label{\detokenize{modules/util/transformations:qnav.util.transformations.lla2ecef_vec}}
\pysigstartsignatures
\pysiglinewithargsret{\sphinxcode{\sphinxupquote{qnav.util.transformations.}}\sphinxbfcode{\sphinxupquote{lla2ecef\_vec}}}{\sphinxparam{\DUrole{n}{lla}\DUrole{p}{:}\DUrole{w}{ }\DUrole{n}{ndarray}}}{{ $\rightarrow$ ndarray}}
\pysigstopsignatures
\sphinxAtStartPar
This function is the vectorised version of the
{\hyperref[\detokenize{modules/util/transformations:qnav.util.transformations.lla2ecef}]{\sphinxcrossref{\sphinxcode{\sphinxupquote{lla2ecef()}}}}} function. This
converts a series of reference Latitude\sphinxhyphen{}Longitude\sphinxhyphen{}Altitude (LLA) positions
to Earth\sphinxhyphen{}Centred Earth\sphinxhyphen{}Fixed coordinates.
\begin{quote}\begin{description}
\sphinxlineitem{Parameters}
\sphinxAtStartPar
\sphinxstyleliteralstrong{\sphinxupquote{lla}} (\sphinxstyleliteralemphasis{\sphinxupquote{numpy.ndarray}}\sphinxstyleliteralemphasis{\sphinxupquote{ (}}\sphinxstyleliteralemphasis{\sphinxupquote{n\sphinxhyphen{}by\sphinxhyphen{}3 elements}}\sphinxstyleliteralemphasis{\sphinxupquote{)}}) \textendash{} The reference locations (array rows containing
{[}lat, lon, alt{]}) with lat\sphinxhyphen{}lon given in degrees and alt given in
metres.

\sphinxlineitem{Returns}
\sphinxAtStartPar
An array containing the calculated X, Y and Z coordinates from
the centre of the Earth, with all values being given in metres.

\sphinxlineitem{Return type}
\sphinxAtStartPar
numpy.ndarray (n\sphinxhyphen{}by\sphinxhyphen{}3 elements)

\end{description}\end{quote}

\end{fulllineitems}

\index{euler\_to\_quaternion() (in module qnav.util.transformations)@\spxentry{euler\_to\_quaternion()}\spxextra{in module qnav.util.transformations}}

\begin{fulllineitems}
\phantomsection\label{\detokenize{modules/util/transformations:qnav.util.transformations.euler_to_quaternion}}
\pysigstartsignatures
\pysiglinewithargsret{\sphinxcode{\sphinxupquote{qnav.util.transformations.}}\sphinxbfcode{\sphinxupquote{euler\_to\_quaternion}}}{\sphinxparam{\DUrole{n}{angles}\DUrole{p}{:}\DUrole{w}{ }\DUrole{n}{ndarray}}}{{ $\rightarrow$ ndarray}}
\pysigstopsignatures
\sphinxAtStartPar
Converts given attitude angles to quaternion format.
Given an array representing euler angles, this function converts them to
quaternion format. On succession, an array containing the a, b, c and d
quaternion coefficients (i.e. where a + bi + cj + dk) will be returned.
\begin{quote}\begin{description}
\sphinxlineitem{Parameters}
\sphinxAtStartPar
\sphinxstyleliteralstrong{\sphinxupquote{angles}} (\sphinxstyleliteralemphasis{\sphinxupquote{numpy.ndarray}}\sphinxstyleliteralemphasis{\sphinxupquote{ (}}\sphinxstyleliteralemphasis{\sphinxupquote{3\sphinxhyphen{}elements}}\sphinxstyleliteralemphasis{\sphinxupquote{)}}) \textendash{} Euler angles given in degrees

\sphinxlineitem{Returns}
\sphinxAtStartPar
The a, b, c, and d quaternion coefficients in columns.

\sphinxlineitem{Return type}
\sphinxAtStartPar
numpy.ndarray (4\sphinxhyphen{}elements)

\end{description}\end{quote}

\end{fulllineitems}

\index{quaternion\_to\_euler() (in module qnav.util.transformations)@\spxentry{quaternion\_to\_euler()}\spxextra{in module qnav.util.transformations}}

\begin{fulllineitems}
\phantomsection\label{\detokenize{modules/util/transformations:qnav.util.transformations.quaternion_to_euler}}
\pysigstartsignatures
\pysiglinewithargsret{\sphinxcode{\sphinxupquote{qnav.util.transformations.}}\sphinxbfcode{\sphinxupquote{quaternion\_to\_euler}}}{\sphinxparam{\DUrole{n}{quat}\DUrole{p}{:}\DUrole{w}{ }\DUrole{n}{ndarray}}}{{ $\rightarrow$ ndarray}}
\pysigstopsignatures
\sphinxAtStartPar
Converts given quaternions to Euler angle format.
Given an array representing coefficients for quaternions (i.e.
values a, b, c and d where a + bi + cj + dk), this function converts them
to euler angle format. On succession, an array containing the Euler angles
will be returned.
\begin{quote}\begin{description}
\sphinxlineitem{Parameters}
\sphinxAtStartPar
\sphinxstyleliteralstrong{\sphinxupquote{quat}} (\sphinxstyleliteralemphasis{\sphinxupquote{numpy.ndarray}}\sphinxstyleliteralemphasis{\sphinxupquote{ (}}\sphinxstyleliteralemphasis{\sphinxupquote{4\sphinxhyphen{}elements}}\sphinxstyleliteralemphasis{\sphinxupquote{)}}) \textendash{} The a, b, c, and d quaternion coefficients in columns.

\sphinxlineitem{Returns}
\sphinxAtStartPar
The converted Euler angles given in degrees.

\sphinxlineitem{Return type}
\sphinxAtStartPar
numpy.ndarray (3\sphinxhyphen{}elements)

\end{description}\end{quote}

\end{fulllineitems}

\index{euler\_to\_quaternion\_vec() (in module qnav.util.transformations)@\spxentry{euler\_to\_quaternion\_vec()}\spxextra{in module qnav.util.transformations}}

\begin{fulllineitems}
\phantomsection\label{\detokenize{modules/util/transformations:qnav.util.transformations.euler_to_quaternion_vec}}
\pysigstartsignatures
\pysiglinewithargsret{\sphinxcode{\sphinxupquote{qnav.util.transformations.}}\sphinxbfcode{\sphinxupquote{euler\_to\_quaternion\_vec}}}{\sphinxparam{\DUrole{n}{angles}\DUrole{p}{:}\DUrole{w}{ }\DUrole{n}{ndarray}}}{{ $\rightarrow$ ndarray}}
\pysigstopsignatures
\sphinxAtStartPar
Converts given attitude angles to quaternion format.
Given an array representing euler angles, this function converts them to
quaternion format. On succession, an array containing the a, b, c and d
quaternion coefficients (i.e. where a + bi + cj + dk) will be returned.
\begin{quote}\begin{description}
\sphinxlineitem{Parameters}
\sphinxAtStartPar
\sphinxstyleliteralstrong{\sphinxupquote{angles}} (\sphinxstyleliteralemphasis{\sphinxupquote{numpy.ndarray}}\sphinxstyleliteralemphasis{\sphinxupquote{ (}}\sphinxstyleliteralemphasis{\sphinxupquote{n\sphinxhyphen{}by\sphinxhyphen{}3 elements}}\sphinxstyleliteralemphasis{\sphinxupquote{)}}) \textendash{} Euler angles given in degrees

\sphinxlineitem{Returns}
\sphinxAtStartPar
The a, b, c, and d quaternion coefficients in columns.

\sphinxlineitem{Return type}
\sphinxAtStartPar
numpy.ndarray (n\sphinxhyphen{}by\sphinxhyphen{}4 elements)

\end{description}\end{quote}

\end{fulllineitems}

\index{quaternion\_to\_euler\_vec() (in module qnav.util.transformations)@\spxentry{quaternion\_to\_euler\_vec()}\spxextra{in module qnav.util.transformations}}

\begin{fulllineitems}
\phantomsection\label{\detokenize{modules/util/transformations:qnav.util.transformations.quaternion_to_euler_vec}}
\pysigstartsignatures
\pysiglinewithargsret{\sphinxcode{\sphinxupquote{qnav.util.transformations.}}\sphinxbfcode{\sphinxupquote{quaternion\_to\_euler\_vec}}}{\sphinxparam{\DUrole{n}{quat}\DUrole{p}{:}\DUrole{w}{ }\DUrole{n}{ndarray}}}{{ $\rightarrow$ ndarray}}
\pysigstopsignatures
\sphinxAtStartPar
Converts given quaternions to Euler angle format.
Given an array representing coefficients for quaternions (i.e.
values a, b, c and d where a + bi + cj + dk), this function converts them
to euler angle format. On succession, an array containing the Euler angles
will be returned.
\begin{quote}\begin{description}
\sphinxlineitem{Parameters}
\sphinxAtStartPar
\sphinxstyleliteralstrong{\sphinxupquote{quat}} (\sphinxstyleliteralemphasis{\sphinxupquote{numpy.ndarray}}\sphinxstyleliteralemphasis{\sphinxupquote{ (}}\sphinxstyleliteralemphasis{\sphinxupquote{n\sphinxhyphen{}by\sphinxhyphen{}4 elements}}\sphinxstyleliteralemphasis{\sphinxupquote{)}}) \textendash{} The a, b, c, and d quaternion coefficients in columns.

\sphinxlineitem{Returns}
\sphinxAtStartPar
The converted Euler angles given in degrees.

\sphinxlineitem{Return type}
\sphinxAtStartPar
numpy.ndarray (n\sphinxhyphen{}by\sphinxhyphen{}3 elements)

\end{description}\end{quote}

\end{fulllineitems}

\index{haversine() (in module qnav.util.transformations)@\spxentry{haversine()}\spxextra{in module qnav.util.transformations}}

\begin{fulllineitems}
\phantomsection\label{\detokenize{modules/util/transformations:qnav.util.transformations.haversine}}
\pysigstartsignatures
\pysiglinewithargsret{\sphinxcode{\sphinxupquote{qnav.util.transformations.}}\sphinxbfcode{\sphinxupquote{haversine}}}{\sphinxparam{\DUrole{n}{lat1}\DUrole{p}{:}\DUrole{w}{ }\DUrole{n}{float}}\sphinxparamcomma \sphinxparam{\DUrole{n}{lon1}\DUrole{p}{:}\DUrole{w}{ }\DUrole{n}{float}}\sphinxparamcomma \sphinxparam{\DUrole{n}{alt1}\DUrole{p}{:}\DUrole{w}{ }\DUrole{n}{float}}\sphinxparamcomma \sphinxparam{\DUrole{n}{lat2}\DUrole{p}{:}\DUrole{w}{ }\DUrole{n}{float}}\sphinxparamcomma \sphinxparam{\DUrole{n}{lon2}\DUrole{p}{:}\DUrole{w}{ }\DUrole{n}{float}}\sphinxparamcomma \sphinxparam{\DUrole{n}{alt2}\DUrole{p}{:}\DUrole{w}{ }\DUrole{n}{float}}}{{ $\rightarrow$ float}}
\pysigstopsignatures
\sphinxAtStartPar
Returns the estimated distance in metres between two sets of
latitude\sphinxhyphen{}longitude\sphinxhyphen{}altitude positions. This function uses the Haversine
formula to calculate the distance between latitude\sphinxhyphen{}longitude points, then
employs pythagoras to estimate the actual distance by considering change
in altitude between positions.
\begin{quote}\begin{description}
\sphinxlineitem{Parameters}\begin{itemize}
\item {} 
\sphinxAtStartPar
\sphinxstyleliteralstrong{\sphinxupquote{lat1}} (\sphinxstyleliteralemphasis{\sphinxupquote{numpy.ndarray}}\sphinxstyleliteralemphasis{\sphinxupquote{ (}}\sphinxstyleliteralemphasis{\sphinxupquote{n\sphinxhyphen{}elements}}\sphinxstyleliteralemphasis{\sphinxupquote{)}}) \textendash{} The latitude position(s) belonging to the first set.

\item {} 
\sphinxAtStartPar
\sphinxstyleliteralstrong{\sphinxupquote{lon1}} (\sphinxstyleliteralemphasis{\sphinxupquote{numpy.ndarray}}\sphinxstyleliteralemphasis{\sphinxupquote{ (}}\sphinxstyleliteralemphasis{\sphinxupquote{n\sphinxhyphen{}elements}}\sphinxstyleliteralemphasis{\sphinxupquote{)}}) \textendash{} The longitude position(s) belonging to the first set.

\item {} 
\sphinxAtStartPar
\sphinxstyleliteralstrong{\sphinxupquote{alt1}} (\sphinxstyleliteralemphasis{\sphinxupquote{numpy.ndarray}}\sphinxstyleliteralemphasis{\sphinxupquote{ (}}\sphinxstyleliteralemphasis{\sphinxupquote{n\sphinxhyphen{}elements}}\sphinxstyleliteralemphasis{\sphinxupquote{)}}) \textendash{} The altitude position(s) belonging to the first set.

\item {} 
\sphinxAtStartPar
\sphinxstyleliteralstrong{\sphinxupquote{lat2}} (\sphinxstyleliteralemphasis{\sphinxupquote{numpy.ndarray}}\sphinxstyleliteralemphasis{\sphinxupquote{ (}}\sphinxstyleliteralemphasis{\sphinxupquote{n\sphinxhyphen{}elements}}\sphinxstyleliteralemphasis{\sphinxupquote{)}}) \textendash{} The latitude position(s) belonging to the second set.

\item {} 
\sphinxAtStartPar
\sphinxstyleliteralstrong{\sphinxupquote{lon2}} (\sphinxstyleliteralemphasis{\sphinxupquote{numpy.ndarray}}\sphinxstyleliteralemphasis{\sphinxupquote{ (}}\sphinxstyleliteralemphasis{\sphinxupquote{n\sphinxhyphen{}elements}}\sphinxstyleliteralemphasis{\sphinxupquote{)}}) \textendash{} The longitude position(s) belonging to the second set.

\item {} 
\sphinxAtStartPar
\sphinxstyleliteralstrong{\sphinxupquote{alt2}} (\sphinxstyleliteralemphasis{\sphinxupquote{numpy.ndarray}}\sphinxstyleliteralemphasis{\sphinxupquote{ (}}\sphinxstyleliteralemphasis{\sphinxupquote{n\sphinxhyphen{}elements}}\sphinxstyleliteralemphasis{\sphinxupquote{)}}) \textendash{} The latitude position(s) belonging to the second set.

\end{itemize}

\sphinxlineitem{Returns}
\sphinxAtStartPar
The distance between the given positions in metres.

\sphinxlineitem{Return type}
\sphinxAtStartPar
numpy.array (n\sphinxhyphen{}elements)

\end{description}\end{quote}

\end{fulllineitems}

\index{haversine\_vec() (in module qnav.util.transformations)@\spxentry{haversine\_vec()}\spxextra{in module qnav.util.transformations}}

\begin{fulllineitems}
\phantomsection\label{\detokenize{modules/util/transformations:qnav.util.transformations.haversine_vec}}
\pysigstartsignatures
\pysiglinewithargsret{\sphinxcode{\sphinxupquote{qnav.util.transformations.}}\sphinxbfcode{\sphinxupquote{haversine\_vec}}}{\sphinxparam{\DUrole{n}{lat1}\DUrole{p}{:}\DUrole{w}{ }\DUrole{n}{Buffer\DUrole{w}{ }\DUrole{p}{|}\DUrole{w}{ }\_SupportsArray\DUrole{p}{{[}}dtype\DUrole{p}{{[}}Any\DUrole{p}{{]}}\DUrole{p}{{]}}\DUrole{w}{ }\DUrole{p}{|}\DUrole{w}{ }\_NestedSequence\DUrole{p}{{[}}\_SupportsArray\DUrole{p}{{[}}dtype\DUrole{p}{{[}}Any\DUrole{p}{{]}}\DUrole{p}{{]}}\DUrole{p}{{]}}\DUrole{w}{ }\DUrole{p}{|}\DUrole{w}{ }bool\DUrole{w}{ }\DUrole{p}{|}\DUrole{w}{ }int\DUrole{w}{ }\DUrole{p}{|}\DUrole{w}{ }float\DUrole{w}{ }\DUrole{p}{|}\DUrole{w}{ }complex\DUrole{w}{ }\DUrole{p}{|}\DUrole{w}{ }str\DUrole{w}{ }\DUrole{p}{|}\DUrole{w}{ }bytes\DUrole{w}{ }\DUrole{p}{|}\DUrole{w}{ }\_NestedSequence\DUrole{p}{{[}}bool\DUrole{w}{ }\DUrole{p}{|}\DUrole{w}{ }int\DUrole{w}{ }\DUrole{p}{|}\DUrole{w}{ }float\DUrole{w}{ }\DUrole{p}{|}\DUrole{w}{ }complex\DUrole{w}{ }\DUrole{p}{|}\DUrole{w}{ }str\DUrole{w}{ }\DUrole{p}{|}\DUrole{w}{ }bytes\DUrole{p}{{]}}}}\sphinxparamcomma \sphinxparam{\DUrole{n}{lon1}\DUrole{p}{:}\DUrole{w}{ }\DUrole{n}{Buffer\DUrole{w}{ }\DUrole{p}{|}\DUrole{w}{ }\_SupportsArray\DUrole{p}{{[}}dtype\DUrole{p}{{[}}Any\DUrole{p}{{]}}\DUrole{p}{{]}}\DUrole{w}{ }\DUrole{p}{|}\DUrole{w}{ }\_NestedSequence\DUrole{p}{{[}}\_SupportsArray\DUrole{p}{{[}}dtype\DUrole{p}{{[}}Any\DUrole{p}{{]}}\DUrole{p}{{]}}\DUrole{p}{{]}}\DUrole{w}{ }\DUrole{p}{|}\DUrole{w}{ }bool\DUrole{w}{ }\DUrole{p}{|}\DUrole{w}{ }int\DUrole{w}{ }\DUrole{p}{|}\DUrole{w}{ }float\DUrole{w}{ }\DUrole{p}{|}\DUrole{w}{ }complex\DUrole{w}{ }\DUrole{p}{|}\DUrole{w}{ }str\DUrole{w}{ }\DUrole{p}{|}\DUrole{w}{ }bytes\DUrole{w}{ }\DUrole{p}{|}\DUrole{w}{ }\_NestedSequence\DUrole{p}{{[}}bool\DUrole{w}{ }\DUrole{p}{|}\DUrole{w}{ }int\DUrole{w}{ }\DUrole{p}{|}\DUrole{w}{ }float\DUrole{w}{ }\DUrole{p}{|}\DUrole{w}{ }complex\DUrole{w}{ }\DUrole{p}{|}\DUrole{w}{ }str\DUrole{w}{ }\DUrole{p}{|}\DUrole{w}{ }bytes\DUrole{p}{{]}}}}\sphinxparamcomma \sphinxparam{\DUrole{n}{alt1}\DUrole{p}{:}\DUrole{w}{ }\DUrole{n}{Buffer\DUrole{w}{ }\DUrole{p}{|}\DUrole{w}{ }\_SupportsArray\DUrole{p}{{[}}dtype\DUrole{p}{{[}}Any\DUrole{p}{{]}}\DUrole{p}{{]}}\DUrole{w}{ }\DUrole{p}{|}\DUrole{w}{ }\_NestedSequence\DUrole{p}{{[}}\_SupportsArray\DUrole{p}{{[}}dtype\DUrole{p}{{[}}Any\DUrole{p}{{]}}\DUrole{p}{{]}}\DUrole{p}{{]}}\DUrole{w}{ }\DUrole{p}{|}\DUrole{w}{ }bool\DUrole{w}{ }\DUrole{p}{|}\DUrole{w}{ }int\DUrole{w}{ }\DUrole{p}{|}\DUrole{w}{ }float\DUrole{w}{ }\DUrole{p}{|}\DUrole{w}{ }complex\DUrole{w}{ }\DUrole{p}{|}\DUrole{w}{ }str\DUrole{w}{ }\DUrole{p}{|}\DUrole{w}{ }bytes\DUrole{w}{ }\DUrole{p}{|}\DUrole{w}{ }\_NestedSequence\DUrole{p}{{[}}bool\DUrole{w}{ }\DUrole{p}{|}\DUrole{w}{ }int\DUrole{w}{ }\DUrole{p}{|}\DUrole{w}{ }float\DUrole{w}{ }\DUrole{p}{|}\DUrole{w}{ }complex\DUrole{w}{ }\DUrole{p}{|}\DUrole{w}{ }str\DUrole{w}{ }\DUrole{p}{|}\DUrole{w}{ }bytes\DUrole{p}{{]}}}}\sphinxparamcomma \sphinxparam{\DUrole{n}{lat2}\DUrole{p}{:}\DUrole{w}{ }\DUrole{n}{Buffer\DUrole{w}{ }\DUrole{p}{|}\DUrole{w}{ }\_SupportsArray\DUrole{p}{{[}}dtype\DUrole{p}{{[}}Any\DUrole{p}{{]}}\DUrole{p}{{]}}\DUrole{w}{ }\DUrole{p}{|}\DUrole{w}{ }\_NestedSequence\DUrole{p}{{[}}\_SupportsArray\DUrole{p}{{[}}dtype\DUrole{p}{{[}}Any\DUrole{p}{{]}}\DUrole{p}{{]}}\DUrole{p}{{]}}\DUrole{w}{ }\DUrole{p}{|}\DUrole{w}{ }bool\DUrole{w}{ }\DUrole{p}{|}\DUrole{w}{ }int\DUrole{w}{ }\DUrole{p}{|}\DUrole{w}{ }float\DUrole{w}{ }\DUrole{p}{|}\DUrole{w}{ }complex\DUrole{w}{ }\DUrole{p}{|}\DUrole{w}{ }str\DUrole{w}{ }\DUrole{p}{|}\DUrole{w}{ }bytes\DUrole{w}{ }\DUrole{p}{|}\DUrole{w}{ }\_NestedSequence\DUrole{p}{{[}}bool\DUrole{w}{ }\DUrole{p}{|}\DUrole{w}{ }int\DUrole{w}{ }\DUrole{p}{|}\DUrole{w}{ }float\DUrole{w}{ }\DUrole{p}{|}\DUrole{w}{ }complex\DUrole{w}{ }\DUrole{p}{|}\DUrole{w}{ }str\DUrole{w}{ }\DUrole{p}{|}\DUrole{w}{ }bytes\DUrole{p}{{]}}}}\sphinxparamcomma \sphinxparam{\DUrole{n}{lon2}\DUrole{p}{:}\DUrole{w}{ }\DUrole{n}{Buffer\DUrole{w}{ }\DUrole{p}{|}\DUrole{w}{ }\_SupportsArray\DUrole{p}{{[}}dtype\DUrole{p}{{[}}Any\DUrole{p}{{]}}\DUrole{p}{{]}}\DUrole{w}{ }\DUrole{p}{|}\DUrole{w}{ }\_NestedSequence\DUrole{p}{{[}}\_SupportsArray\DUrole{p}{{[}}dtype\DUrole{p}{{[}}Any\DUrole{p}{{]}}\DUrole{p}{{]}}\DUrole{p}{{]}}\DUrole{w}{ }\DUrole{p}{|}\DUrole{w}{ }bool\DUrole{w}{ }\DUrole{p}{|}\DUrole{w}{ }int\DUrole{w}{ }\DUrole{p}{|}\DUrole{w}{ }float\DUrole{w}{ }\DUrole{p}{|}\DUrole{w}{ }complex\DUrole{w}{ }\DUrole{p}{|}\DUrole{w}{ }str\DUrole{w}{ }\DUrole{p}{|}\DUrole{w}{ }bytes\DUrole{w}{ }\DUrole{p}{|}\DUrole{w}{ }\_NestedSequence\DUrole{p}{{[}}bool\DUrole{w}{ }\DUrole{p}{|}\DUrole{w}{ }int\DUrole{w}{ }\DUrole{p}{|}\DUrole{w}{ }float\DUrole{w}{ }\DUrole{p}{|}\DUrole{w}{ }complex\DUrole{w}{ }\DUrole{p}{|}\DUrole{w}{ }str\DUrole{w}{ }\DUrole{p}{|}\DUrole{w}{ }bytes\DUrole{p}{{]}}}}\sphinxparamcomma \sphinxparam{\DUrole{n}{alt2}\DUrole{p}{:}\DUrole{w}{ }\DUrole{n}{Buffer\DUrole{w}{ }\DUrole{p}{|}\DUrole{w}{ }\_SupportsArray\DUrole{p}{{[}}dtype\DUrole{p}{{[}}Any\DUrole{p}{{]}}\DUrole{p}{{]}}\DUrole{w}{ }\DUrole{p}{|}\DUrole{w}{ }\_NestedSequence\DUrole{p}{{[}}\_SupportsArray\DUrole{p}{{[}}dtype\DUrole{p}{{[}}Any\DUrole{p}{{]}}\DUrole{p}{{]}}\DUrole{p}{{]}}\DUrole{w}{ }\DUrole{p}{|}\DUrole{w}{ }bool\DUrole{w}{ }\DUrole{p}{|}\DUrole{w}{ }int\DUrole{w}{ }\DUrole{p}{|}\DUrole{w}{ }float\DUrole{w}{ }\DUrole{p}{|}\DUrole{w}{ }complex\DUrole{w}{ }\DUrole{p}{|}\DUrole{w}{ }str\DUrole{w}{ }\DUrole{p}{|}\DUrole{w}{ }bytes\DUrole{w}{ }\DUrole{p}{|}\DUrole{w}{ }\_NestedSequence\DUrole{p}{{[}}bool\DUrole{w}{ }\DUrole{p}{|}\DUrole{w}{ }int\DUrole{w}{ }\DUrole{p}{|}\DUrole{w}{ }float\DUrole{w}{ }\DUrole{p}{|}\DUrole{w}{ }complex\DUrole{w}{ }\DUrole{p}{|}\DUrole{w}{ }str\DUrole{w}{ }\DUrole{p}{|}\DUrole{w}{ }bytes\DUrole{p}{{]}}}}}{{ $\rightarrow$ ndarray}}
\pysigstopsignatures
\sphinxAtStartPar
Returns the estimated distance in metres between two sets of
latitude\sphinxhyphen{}longitude\sphinxhyphen{}altitude positions. This function uses the Haversine
formula to calculate the distance between latitude\sphinxhyphen{}longitude points, then
employs pythagoras to estimate the actual distance by considering change
in altitude between positions.
\begin{quote}\begin{description}
\sphinxlineitem{Parameters}\begin{itemize}
\item {} 
\sphinxAtStartPar
\sphinxstyleliteralstrong{\sphinxupquote{lat1}} (\sphinxstyleliteralemphasis{\sphinxupquote{numpy.ndarray}}\sphinxstyleliteralemphasis{\sphinxupquote{ (}}\sphinxstyleliteralemphasis{\sphinxupquote{n\sphinxhyphen{}elements}}\sphinxstyleliteralemphasis{\sphinxupquote{)}}) \textendash{} The latitude position(s) belonging to the first set.

\item {} 
\sphinxAtStartPar
\sphinxstyleliteralstrong{\sphinxupquote{lon1}} (\sphinxstyleliteralemphasis{\sphinxupquote{numpy.ndarray}}\sphinxstyleliteralemphasis{\sphinxupquote{ (}}\sphinxstyleliteralemphasis{\sphinxupquote{n\sphinxhyphen{}elements}}\sphinxstyleliteralemphasis{\sphinxupquote{)}}) \textendash{} The longitude position(s) belonging to the first set.

\item {} 
\sphinxAtStartPar
\sphinxstyleliteralstrong{\sphinxupquote{alt1}} (\sphinxstyleliteralemphasis{\sphinxupquote{numpy.ndarray}}\sphinxstyleliteralemphasis{\sphinxupquote{ (}}\sphinxstyleliteralemphasis{\sphinxupquote{n\sphinxhyphen{}elements}}\sphinxstyleliteralemphasis{\sphinxupquote{)}}) \textendash{} The altitude position(s) belonging to the first set.

\item {} 
\sphinxAtStartPar
\sphinxstyleliteralstrong{\sphinxupquote{lat2}} (\sphinxstyleliteralemphasis{\sphinxupquote{numpy.ndarray}}\sphinxstyleliteralemphasis{\sphinxupquote{ (}}\sphinxstyleliteralemphasis{\sphinxupquote{n\sphinxhyphen{}elements}}\sphinxstyleliteralemphasis{\sphinxupquote{)}}) \textendash{} The latitude position(s) belonging to the second set.

\item {} 
\sphinxAtStartPar
\sphinxstyleliteralstrong{\sphinxupquote{lon2}} (\sphinxstyleliteralemphasis{\sphinxupquote{numpy.ndarray}}\sphinxstyleliteralemphasis{\sphinxupquote{ (}}\sphinxstyleliteralemphasis{\sphinxupquote{n\sphinxhyphen{}elements}}\sphinxstyleliteralemphasis{\sphinxupquote{)}}) \textendash{} The longitude position(s) belonging to the second set.

\item {} 
\sphinxAtStartPar
\sphinxstyleliteralstrong{\sphinxupquote{alt2}} (\sphinxstyleliteralemphasis{\sphinxupquote{numpy.ndarray}}\sphinxstyleliteralemphasis{\sphinxupquote{ (}}\sphinxstyleliteralemphasis{\sphinxupquote{n\sphinxhyphen{}elements}}\sphinxstyleliteralemphasis{\sphinxupquote{)}}) \textendash{} The latitude position(s) belonging to the second set.

\end{itemize}

\sphinxlineitem{Returns}
\sphinxAtStartPar
The distance between the given positions in metres.

\sphinxlineitem{Return type}
\sphinxAtStartPar
numpy.array (n\sphinxhyphen{}elements)

\end{description}\end{quote}

\end{fulllineitems}

\index{get\_time\_pos\_vel() (in module qnav.util.transformations)@\spxentry{get\_time\_pos\_vel()}\spxextra{in module qnav.util.transformations}}

\begin{fulllineitems}
\phantomsection\label{\detokenize{modules/util/transformations:qnav.util.transformations.get_time_pos_vel}}
\pysigstartsignatures
\pysiglinewithargsret{\sphinxcode{\sphinxupquote{qnav.util.transformations.}}\sphinxbfcode{\sphinxupquote{get\_time\_pos\_vel}}}{\sphinxparam{\DUrole{n}{time}\DUrole{p}{:}\DUrole{w}{ }\DUrole{n}{datetime}}}{{ $\rightarrow$ tuple\DUrole{p}{{[}}ndarray\DUrole{p}{,}\DUrole{w}{ }ndarray\DUrole{p}{{]}}}}
\pysigstopsignatures
\sphinxAtStartPar
This function calculates the position and velocity matrices for a given
timestamp.
\begin{quote}\begin{description}
\sphinxlineitem{Parameters}
\sphinxAtStartPar
\sphinxstyleliteralstrong{\sphinxupquote{time}} (\sphinxstyleliteralemphasis{\sphinxupquote{Python datetime instance}}) \textendash{} A datetime for the requested time.

\sphinxlineitem{Returns}
\sphinxAtStartPar
The position and velocity matrices.

\sphinxlineitem{Return type}
\sphinxAtStartPar
tuple (mat\_pos, mat\_vel)

\end{description}\end{quote}

\end{fulllineitems}

\index{ecef2eci() (in module qnav.util.transformations)@\spxentry{ecef2eci()}\spxextra{in module qnav.util.transformations}}

\begin{fulllineitems}
\phantomsection\label{\detokenize{modules/util/transformations:qnav.util.transformations.ecef2eci}}
\pysigstartsignatures
\pysiglinewithargsret{\sphinxcode{\sphinxupquote{qnav.util.transformations.}}\sphinxbfcode{\sphinxupquote{ecef2eci}}}{\sphinxparam{\DUrole{n}{time}\DUrole{p}{:}\DUrole{w}{ }\DUrole{n}{datetime}}\sphinxparamcomma \sphinxparam{\DUrole{n}{r\_ecef}\DUrole{p}{:}\DUrole{w}{ }\DUrole{n}{ndarray}}\sphinxparamcomma \sphinxparam{\DUrole{n}{v\_ecef}\DUrole{p}{:}\DUrole{w}{ }\DUrole{n}{ndarray}}}{{ $\rightarrow$ tuple\DUrole{p}{{[}}ndarray\DUrole{p}{,}\DUrole{w}{ }ndarray\DUrole{p}{{]}}}}
\pysigstopsignatures
\sphinxAtStartPar
This function converts a position and velocity from the ECI frame to the
ECEF frame.
\begin{quote}\begin{description}
\sphinxlineitem{Parameters}\begin{itemize}
\item {} 
\sphinxAtStartPar
\sphinxstyleliteralstrong{\sphinxupquote{time}} (\sphinxstyleliteralemphasis{\sphinxupquote{Python datetime instance}}) \textendash{} A datetime for the requested time.

\item {} 
\sphinxAtStartPar
\sphinxstyleliteralstrong{\sphinxupquote{r\_ecef}} (\sphinxstyleliteralemphasis{\sphinxupquote{numpy.ndarray}}\sphinxstyleliteralemphasis{\sphinxupquote{ (}}\sphinxstyleliteralemphasis{\sphinxupquote{3\sphinxhyphen{}elements}}\sphinxstyleliteralemphasis{\sphinxupquote{)}}) \textendash{} Position in ECI (given in metres).

\item {} 
\sphinxAtStartPar
\sphinxstyleliteralstrong{\sphinxupquote{v\_ecef}} \textendash{} Velocity in ECI (given in metres/second).

\end{itemize}

\sphinxlineitem{Returns}
\sphinxAtStartPar
The position and velocity vectors in ECEF. Position is given in
metres and velocity is given in metres/second.

\sphinxlineitem{Return type}
\sphinxAtStartPar
tuple (r\_eci, v\_eci)

\end{description}\end{quote}

\end{fulllineitems}

\index{eci2ecef() (in module qnav.util.transformations)@\spxentry{eci2ecef()}\spxextra{in module qnav.util.transformations}}

\begin{fulllineitems}
\phantomsection\label{\detokenize{modules/util/transformations:qnav.util.transformations.eci2ecef}}
\pysigstartsignatures
\pysiglinewithargsret{\sphinxcode{\sphinxupquote{qnav.util.transformations.}}\sphinxbfcode{\sphinxupquote{eci2ecef}}}{\sphinxparam{\DUrole{n}{time}\DUrole{p}{:}\DUrole{w}{ }\DUrole{n}{datetime}}\sphinxparamcomma \sphinxparam{\DUrole{n}{r\_eci}\DUrole{p}{:}\DUrole{w}{ }\DUrole{n}{ndarray}}\sphinxparamcomma \sphinxparam{\DUrole{n}{v\_eci}\DUrole{p}{:}\DUrole{w}{ }\DUrole{n}{ndarray}}}{{ $\rightarrow$ tuple\DUrole{p}{{[}}ndarray\DUrole{p}{,}\DUrole{w}{ }ndarray\DUrole{p}{{]}}}}
\pysigstopsignatures
\sphinxAtStartPar
This function converts a position and velocity from the ECEF frame to the
ECI frame.
\begin{quote}\begin{description}
\sphinxlineitem{Parameters}\begin{itemize}
\item {} 
\sphinxAtStartPar
\sphinxstyleliteralstrong{\sphinxupquote{time}} (\sphinxstyleliteralemphasis{\sphinxupquote{Python datetime instance}}) \textendash{} A datetime for the requested time.

\item {} 
\sphinxAtStartPar
\sphinxstyleliteralstrong{\sphinxupquote{r\_eci}} (\sphinxstyleliteralemphasis{\sphinxupquote{numpy.ndarray}}\sphinxstyleliteralemphasis{\sphinxupquote{ (}}\sphinxstyleliteralemphasis{\sphinxupquote{3\sphinxhyphen{}elements}}\sphinxstyleliteralemphasis{\sphinxupquote{)}}) \textendash{} Position in ECEF (given in metres).

\item {} 
\sphinxAtStartPar
\sphinxstyleliteralstrong{\sphinxupquote{v\_eci}} (\sphinxstyleliteralemphasis{\sphinxupquote{numpy.ndarray}}\sphinxstyleliteralemphasis{\sphinxupquote{ (}}\sphinxstyleliteralemphasis{\sphinxupquote{3\sphinxhyphen{}elements}}\sphinxstyleliteralemphasis{\sphinxupquote{)}}) \textendash{} Velocity in ECEF (given in metres/second).

\end{itemize}

\sphinxlineitem{Returns}
\sphinxAtStartPar
The position and velocity vectors in ECI. Position is given in
metres and velocity is given in metres/second.

\sphinxlineitem{Return type}
\sphinxAtStartPar
tuple (r\_ecef, v\_ecef)

\end{description}\end{quote}

\end{fulllineitems}

\index{ecef2lla() (in module qnav.util.transformations)@\spxentry{ecef2lla()}\spxextra{in module qnav.util.transformations}}

\begin{fulllineitems}
\phantomsection\label{\detokenize{modules/util/transformations:qnav.util.transformations.ecef2lla}}
\pysigstartsignatures
\pysiglinewithargsret{\sphinxcode{\sphinxupquote{qnav.util.transformations.}}\sphinxbfcode{\sphinxupquote{ecef2lla}}}{\sphinxparam{\DUrole{n}{xyz}\DUrole{p}{:}\DUrole{w}{ }\DUrole{n}{ndarray}}}{{ $\rightarrow$ ndarray}}
\pysigstopsignatures
\sphinxAtStartPar
Converts the given XYZ Earth Centred Earth Fixed coordinates to
longitude latitude altitude coordinates.
\begin{quote}\begin{description}
\sphinxlineitem{Parameters}
\sphinxAtStartPar
\sphinxstyleliteralstrong{\sphinxupquote{xyz}} (\sphinxstyleliteralemphasis{\sphinxupquote{numpy.ndarray}}\sphinxstyleliteralemphasis{\sphinxupquote{ (}}\sphinxstyleliteralemphasis{\sphinxupquote{3\sphinxhyphen{}elements}}\sphinxstyleliteralemphasis{\sphinxupquote{)}}) \textendash{} A position in ECEF (m\sphinxhyphen{}m\sphinxhyphen{}m)

\sphinxlineitem{Returns}
\sphinxAtStartPar
The given position converted to LLA (deg\sphinxhyphen{}deg\sphinxhyphen{}m).

\sphinxlineitem{Return type}
\sphinxAtStartPar
numpy.ndarray (3\sphinxhyphen{}elements)

\end{description}\end{quote}

\end{fulllineitems}

\index{ecef2lla\_vec() (in module qnav.util.transformations)@\spxentry{ecef2lla\_vec()}\spxextra{in module qnav.util.transformations}}

\begin{fulllineitems}
\phantomsection\label{\detokenize{modules/util/transformations:qnav.util.transformations.ecef2lla_vec}}
\pysigstartsignatures
\pysiglinewithargsret{\sphinxcode{\sphinxupquote{qnav.util.transformations.}}\sphinxbfcode{\sphinxupquote{ecef2lla\_vec}}}{\sphinxparam{\DUrole{n}{xyz}\DUrole{p}{:}\DUrole{w}{ }\DUrole{n}{ndarray}}}{{ $\rightarrow$ ndarray}}
\pysigstopsignatures
\sphinxAtStartPar
This function is the vectorised version of the
{\hyperref[\detokenize{modules/util/transformations:qnav.util.transformations.ecef2lla}]{\sphinxcrossref{\sphinxcode{\sphinxupquote{ecef2lla()}}}}} function.
This converts the given XYZ Earth Centred Earth Fixed coordinates to
longitude latitude altitude coordinates.
\begin{quote}\begin{description}
\sphinxlineitem{Parameters}
\sphinxAtStartPar
\sphinxstyleliteralstrong{\sphinxupquote{xyz}} (\sphinxstyleliteralemphasis{\sphinxupquote{numpy.ndarray}}\sphinxstyleliteralemphasis{\sphinxupquote{ (}}\sphinxstyleliteralemphasis{\sphinxupquote{n\sphinxhyphen{}by\sphinxhyphen{}3 elements}}\sphinxstyleliteralemphasis{\sphinxupquote{)}}) \textendash{} A series of positions in ECEF (m\sphinxhyphen{}m\sphinxhyphen{}m)

\sphinxlineitem{Returns}
\sphinxAtStartPar
The given positions converted to LLA (deg\sphinxhyphen{}deg\sphinxhyphen{}m).

\sphinxlineitem{Return type}
\sphinxAtStartPar
numpy.ndarray (n\sphinxhyphen{}by\sphinxhyphen{}3 elements)

\end{description}\end{quote}

\end{fulllineitems}

\index{ecef2ned\_transform() (in module qnav.util.transformations)@\spxentry{ecef2ned\_transform()}\spxextra{in module qnav.util.transformations}}

\begin{fulllineitems}
\phantomsection\label{\detokenize{modules/util/transformations:qnav.util.transformations.ecef2ned_transform}}
\pysigstartsignatures
\pysiglinewithargsret{\sphinxcode{\sphinxupquote{qnav.util.transformations.}}\sphinxbfcode{\sphinxupquote{ecef2ned\_transform}}}{\sphinxparam{\DUrole{n}{ref\_lla}\DUrole{p}{:}\DUrole{w}{ }\DUrole{n}{ndarray}}}{{ $\rightarrow$ ndarray}}
\pysigstopsignatures
\sphinxAtStartPar
This function generates a transformation matrix for ECEF to
North\sphinxhyphen{}East\sphinxhyphen{}Down (NED) for an input Latitude\sphinxhyphen{}Longitude\sphinxhyphen{}Altitude (LLA)
reference.
\begin{quote}\begin{description}
\sphinxlineitem{Parameters}
\sphinxAtStartPar
\sphinxstyleliteralstrong{\sphinxupquote{ref\_lla}} (\sphinxstyleliteralemphasis{\sphinxupquote{numpy.ndarray}}\sphinxstyleliteralemphasis{\sphinxupquote{ (}}\sphinxstyleliteralemphasis{\sphinxupquote{3\sphinxhyphen{}elements}}\sphinxstyleliteralemphasis{\sphinxupquote{)}}) \textendash{} The reference position in LLA, with lat\sphinxhyphen{}lon given in
degrees and alt given in metres.

\sphinxlineitem{Returns}
\sphinxAtStartPar
The calculated 3\sphinxhyphen{}by\sphinxhyphen{}3 transformation matrix.

\sphinxlineitem{Return type}
\sphinxAtStartPar
NumPy Matrix (3\sphinxhyphen{}by\sphinxhyphen{}3 elements)

\end{description}\end{quote}

\end{fulllineitems}

\index{ecef2ned\_transform\_vec() (in module qnav.util.transformations)@\spxentry{ecef2ned\_transform\_vec()}\spxextra{in module qnav.util.transformations}}

\begin{fulllineitems}
\phantomsection\label{\detokenize{modules/util/transformations:qnav.util.transformations.ecef2ned_transform_vec}}
\pysigstartsignatures
\pysiglinewithargsret{\sphinxcode{\sphinxupquote{qnav.util.transformations.}}\sphinxbfcode{\sphinxupquote{ecef2ned\_transform\_vec}}}{\sphinxparam{\DUrole{n}{ref\_lla}\DUrole{p}{:}\DUrole{w}{ }\DUrole{n}{ndarray}}}{{ $\rightarrow$ ndarray}}
\pysigstopsignatures
\sphinxAtStartPar
This function is the vectorised version of the
{\hyperref[\detokenize{modules/util/transformations:qnav.util.transformations.ecef2ned_transform}]{\sphinxcrossref{\sphinxcode{\sphinxupquote{ecef2ned\_transform()}}}}}
function. This generates transformation matrices for ECEF to
North\sphinxhyphen{}East\sphinxhyphen{}Down (NED) for input Latitude\sphinxhyphen{}Longitude\sphinxhyphen{}Altitude (LLA)
references.
\begin{quote}\begin{description}
\sphinxlineitem{Parameters}
\sphinxAtStartPar
\sphinxstyleliteralstrong{\sphinxupquote{ref\_lla}} (\sphinxstyleliteralemphasis{\sphinxupquote{numpy.ndarray}}\sphinxstyleliteralemphasis{\sphinxupquote{ (}}\sphinxstyleliteralemphasis{\sphinxupquote{n\sphinxhyphen{}by\sphinxhyphen{}3 elements}}\sphinxstyleliteralemphasis{\sphinxupquote{)}}) \textendash{} The reference position in LLAs, with lat\sphinxhyphen{}lon given in
degrees and alt given in metres.

\sphinxlineitem{Returns}
\sphinxAtStartPar
The calculated n\sphinxhyphen{}by\sphinxhyphen{}3\sphinxhyphen{}by\sphinxhyphen{}3 transformation matrix.

\sphinxlineitem{Return type}
\sphinxAtStartPar
numpy.ndarray (n\sphinxhyphen{}3\sphinxhyphen{}3 elements)

\end{description}\end{quote}

\end{fulllineitems}

\index{ecef2ned() (in module qnav.util.transformations)@\spxentry{ecef2ned()}\spxextra{in module qnav.util.transformations}}

\begin{fulllineitems}
\phantomsection\label{\detokenize{modules/util/transformations:qnav.util.transformations.ecef2ned}}
\pysigstartsignatures
\pysiglinewithargsret{\sphinxcode{\sphinxupquote{qnav.util.transformations.}}\sphinxbfcode{\sphinxupquote{ecef2ned}}}{\sphinxparam{\DUrole{n}{xyz}\DUrole{p}{:}\DUrole{w}{ }\DUrole{n}{Buffer\DUrole{w}{ }\DUrole{p}{|}\DUrole{w}{ }\_SupportsArray\DUrole{p}{{[}}dtype\DUrole{p}{{[}}Any\DUrole{p}{{]}}\DUrole{p}{{]}}\DUrole{w}{ }\DUrole{p}{|}\DUrole{w}{ }\_NestedSequence\DUrole{p}{{[}}\_SupportsArray\DUrole{p}{{[}}dtype\DUrole{p}{{[}}Any\DUrole{p}{{]}}\DUrole{p}{{]}}\DUrole{p}{{]}}\DUrole{w}{ }\DUrole{p}{|}\DUrole{w}{ }bool\DUrole{w}{ }\DUrole{p}{|}\DUrole{w}{ }int\DUrole{w}{ }\DUrole{p}{|}\DUrole{w}{ }float\DUrole{w}{ }\DUrole{p}{|}\DUrole{w}{ }complex\DUrole{w}{ }\DUrole{p}{|}\DUrole{w}{ }str\DUrole{w}{ }\DUrole{p}{|}\DUrole{w}{ }bytes\DUrole{w}{ }\DUrole{p}{|}\DUrole{w}{ }\_NestedSequence\DUrole{p}{{[}}bool\DUrole{w}{ }\DUrole{p}{|}\DUrole{w}{ }int\DUrole{w}{ }\DUrole{p}{|}\DUrole{w}{ }float\DUrole{w}{ }\DUrole{p}{|}\DUrole{w}{ }complex\DUrole{w}{ }\DUrole{p}{|}\DUrole{w}{ }str\DUrole{w}{ }\DUrole{p}{|}\DUrole{w}{ }bytes\DUrole{p}{{]}}}}\sphinxparamcomma \sphinxparam{\DUrole{n}{ref\_lla}\DUrole{p}{:}\DUrole{w}{ }\DUrole{n}{Buffer\DUrole{w}{ }\DUrole{p}{|}\DUrole{w}{ }\_SupportsArray\DUrole{p}{{[}}dtype\DUrole{p}{{[}}Any\DUrole{p}{{]}}\DUrole{p}{{]}}\DUrole{w}{ }\DUrole{p}{|}\DUrole{w}{ }\_NestedSequence\DUrole{p}{{[}}\_SupportsArray\DUrole{p}{{[}}dtype\DUrole{p}{{[}}Any\DUrole{p}{{]}}\DUrole{p}{{]}}\DUrole{p}{{]}}\DUrole{w}{ }\DUrole{p}{|}\DUrole{w}{ }bool\DUrole{w}{ }\DUrole{p}{|}\DUrole{w}{ }int\DUrole{w}{ }\DUrole{p}{|}\DUrole{w}{ }float\DUrole{w}{ }\DUrole{p}{|}\DUrole{w}{ }complex\DUrole{w}{ }\DUrole{p}{|}\DUrole{w}{ }str\DUrole{w}{ }\DUrole{p}{|}\DUrole{w}{ }bytes\DUrole{w}{ }\DUrole{p}{|}\DUrole{w}{ }\_NestedSequence\DUrole{p}{{[}}bool\DUrole{w}{ }\DUrole{p}{|}\DUrole{w}{ }int\DUrole{w}{ }\DUrole{p}{|}\DUrole{w}{ }float\DUrole{w}{ }\DUrole{p}{|}\DUrole{w}{ }complex\DUrole{w}{ }\DUrole{p}{|}\DUrole{w}{ }str\DUrole{w}{ }\DUrole{p}{|}\DUrole{w}{ }bytes\DUrole{p}{{]}}}}}{{ $\rightarrow$ ndarray}}
\pysigstopsignatures
\sphinxAtStartPar
This function takes an ECEF position and a reference
Latitude\sphinxhyphen{}Longitude\sphinxhyphen{}Altitude (LLA) position and returns the relative
position of the former from the latter in North\sphinxhyphen{}East\sphinxhyphen{}Down (NED) reference
frame.
\begin{quote}\begin{description}
\sphinxlineitem{Parameters}\begin{itemize}
\item {} 
\sphinxAtStartPar
\sphinxstyleliteralstrong{\sphinxupquote{xyz}} (\sphinxstyleliteralemphasis{\sphinxupquote{numpy.ndarray}}\sphinxstyleliteralemphasis{\sphinxupquote{ (}}\sphinxstyleliteralemphasis{\sphinxupquote{3\sphinxhyphen{}elements}}\sphinxstyleliteralemphasis{\sphinxupquote{)}}) \textendash{} Position in ECEF, given in metres.

\item {} 
\sphinxAtStartPar
\sphinxstyleliteralstrong{\sphinxupquote{ref\_lla}} (\sphinxstyleliteralemphasis{\sphinxupquote{numpy.ndarray}}\sphinxstyleliteralemphasis{\sphinxupquote{ (}}\sphinxstyleliteralemphasis{\sphinxupquote{3\sphinxhyphen{}elements}}\sphinxstyleliteralemphasis{\sphinxupquote{)}}) \textendash{} The reference LLA position, with lat\sphinxhyphen{}lon given in degrees
and alt given in metres.

\end{itemize}

\sphinxlineitem{Returns}
\sphinxAtStartPar
The given position converted to NED, given in metres.

\sphinxlineitem{Return type}
\sphinxAtStartPar
numpy.ndarray (3\sphinxhyphen{}elements)

\end{description}\end{quote}

\end{fulllineitems}

\index{ned2ecef() (in module qnav.util.transformations)@\spxentry{ned2ecef()}\spxextra{in module qnav.util.transformations}}

\begin{fulllineitems}
\phantomsection\label{\detokenize{modules/util/transformations:qnav.util.transformations.ned2ecef}}
\pysigstartsignatures
\pysiglinewithargsret{\sphinxcode{\sphinxupquote{qnav.util.transformations.}}\sphinxbfcode{\sphinxupquote{ned2ecef}}}{\sphinxparam{\DUrole{n}{ned}\DUrole{p}{:}\DUrole{w}{ }\DUrole{n}{ndarray}}\sphinxparamcomma \sphinxparam{\DUrole{n}{ref\_lla}\DUrole{p}{:}\DUrole{w}{ }\DUrole{n}{ndarray}}}{{ $\rightarrow$ ndarray}}
\pysigstopsignatures
\sphinxAtStartPar
This function takes an origin Latitude\sphinxhyphen{}Longitude\sphinxhyphen{}Altitude (LLA) position
and a relative position in North\sphinxhyphen{}East\sphinxhyphen{}Down (NED) from this origin and
returns the latter position in ECEF.
\begin{quote}\begin{description}
\sphinxlineitem{Parameters}\begin{itemize}
\item {} 
\sphinxAtStartPar
\sphinxstyleliteralstrong{\sphinxupquote{ned}} (\sphinxstyleliteralemphasis{\sphinxupquote{numpy.ndarray}}\sphinxstyleliteralemphasis{\sphinxupquote{ (}}\sphinxstyleliteralemphasis{\sphinxupquote{3\sphinxhyphen{}elements}}\sphinxstyleliteralemphasis{\sphinxupquote{)}}) \textendash{} Position in NED, given in metres.

\item {} 
\sphinxAtStartPar
\sphinxstyleliteralstrong{\sphinxupquote{ref\_lla}} (\sphinxstyleliteralemphasis{\sphinxupquote{numpy.ndarray}}\sphinxstyleliteralemphasis{\sphinxupquote{ (}}\sphinxstyleliteralemphasis{\sphinxupquote{3\sphinxhyphen{}elements}}\sphinxstyleliteralemphasis{\sphinxupquote{)}}) \textendash{} The reference LLA position, with lat\sphinxhyphen{}lon given in degrees
and alt given in metres.

\end{itemize}

\sphinxlineitem{Returns}
\sphinxAtStartPar
The given position converted to ECEF, given in metres.

\sphinxlineitem{Return type}
\sphinxAtStartPar
numpy.ndarray (3\sphinxhyphen{}elements)

\end{description}\end{quote}

\end{fulllineitems}

\index{ned2ecef\_vec() (in module qnav.util.transformations)@\spxentry{ned2ecef\_vec()}\spxextra{in module qnav.util.transformations}}

\begin{fulllineitems}
\phantomsection\label{\detokenize{modules/util/transformations:qnav.util.transformations.ned2ecef_vec}}
\pysigstartsignatures
\pysiglinewithargsret{\sphinxcode{\sphinxupquote{qnav.util.transformations.}}\sphinxbfcode{\sphinxupquote{ned2ecef\_vec}}}{\sphinxparam{\DUrole{n}{ned}\DUrole{p}{:}\DUrole{w}{ }\DUrole{n}{ndarray}}\sphinxparamcomma \sphinxparam{\DUrole{n}{ref\_lla}\DUrole{p}{:}\DUrole{w}{ }\DUrole{n}{ndarray}}}{{ $\rightarrow$ ndarray}}
\pysigstopsignatures
\sphinxAtStartPar
This function takes a North\sphinxhyphen{}East\sphinxhyphen{}Down (NED) position along with its
relative Latitude\sphinxhyphen{}Longitude\sphinxhyphen{}Altitude (LLA) reference position and
returns the latter position in ECEF.
\begin{quote}\begin{description}
\sphinxlineitem{Parameters}\begin{itemize}
\item {} 
\sphinxAtStartPar
\sphinxstyleliteralstrong{\sphinxupquote{ned}} (\sphinxstyleliteralemphasis{\sphinxupquote{numpy.ndarray}}\sphinxstyleliteralemphasis{\sphinxupquote{  (}}\sphinxstyleliteralemphasis{\sphinxupquote{n\sphinxhyphen{}by\sphinxhyphen{}3 elements}}\sphinxstyleliteralemphasis{\sphinxupquote{)}}) \textendash{} Positions in NED, given in metres.

\item {} 
\sphinxAtStartPar
\sphinxstyleliteralstrong{\sphinxupquote{ref\_lla}} (\sphinxstyleliteralemphasis{\sphinxupquote{numpy.ndarray}}\sphinxstyleliteralemphasis{\sphinxupquote{ (}}\sphinxstyleliteralemphasis{\sphinxupquote{n\sphinxhyphen{}by\sphinxhyphen{}3 elements}}\sphinxstyleliteralemphasis{\sphinxupquote{)}}) \textendash{} The reference LLA positions, with lat\sphinxhyphen{}lon given in
degrees and alt given in metres.

\end{itemize}

\sphinxlineitem{Returns}
\sphinxAtStartPar
The given position converted to ECEF, given in metres.

\sphinxlineitem{Return type}
\sphinxAtStartPar
numpy.ndarray (n\sphinxhyphen{}by\sphinxhyphen{}3 elements)

\end{description}\end{quote}

\end{fulllineitems}

\index{vned2vecef() (in module qnav.util.transformations)@\spxentry{vned2vecef()}\spxextra{in module qnav.util.transformations}}

\begin{fulllineitems}
\phantomsection\label{\detokenize{modules/util/transformations:qnav.util.transformations.vned2vecef}}
\pysigstartsignatures
\pysiglinewithargsret{\sphinxcode{\sphinxupquote{qnav.util.transformations.}}\sphinxbfcode{\sphinxupquote{vned2vecef}}}{\sphinxparam{\DUrole{n}{vel\_ned}\DUrole{p}{:}\DUrole{w}{ }\DUrole{n}{ndarray}}\sphinxparamcomma \sphinxparam{\DUrole{n}{pos\_lla}\DUrole{p}{:}\DUrole{w}{ }\DUrole{n}{ndarray}}}{{ $\rightarrow$ ndarray}}
\pysigstopsignatures
\sphinxAtStartPar
This function takes an origin Latitude\sphinxhyphen{}Longitude\sphinxhyphen{}Altitude (LLA) position
and a relative velocity in North\sphinxhyphen{}East\sphinxhyphen{}Down (NED) components from this
origin and returns the velocity in ECEF format.
\begin{quote}\begin{description}
\sphinxlineitem{Parameters}\begin{itemize}
\item {} 
\sphinxAtStartPar
\sphinxstyleliteralstrong{\sphinxupquote{vel\_ned}} (\sphinxstyleliteralemphasis{\sphinxupquote{numpy.ndarray}}\sphinxstyleliteralemphasis{\sphinxupquote{ (}}\sphinxstyleliteralemphasis{\sphinxupquote{3\sphinxhyphen{}elements}}\sphinxstyleliteralemphasis{\sphinxupquote{)}}) \textendash{} Velocity in NED, given in metres/second.

\item {} 
\sphinxAtStartPar
\sphinxstyleliteralstrong{\sphinxupquote{pos\_lla}} (\sphinxstyleliteralemphasis{\sphinxupquote{numpy.ndarray}}\sphinxstyleliteralemphasis{\sphinxupquote{ (}}\sphinxstyleliteralemphasis{\sphinxupquote{3\sphinxhyphen{}elements}}\sphinxstyleliteralemphasis{\sphinxupquote{)}}) \textendash{} The reference LLA position, with lat\sphinxhyphen{}lon given in degrees
and alt given in metres.

\end{itemize}

\sphinxlineitem{Returns}
\sphinxAtStartPar
The given velocity vector converted to ECEF, given in metres/sec.

\sphinxlineitem{Return type}
\sphinxAtStartPar
NumPy Matrix (3\sphinxhyphen{}by\sphinxhyphen{}3 elements)

\end{description}\end{quote}

\end{fulllineitems}

\index{utc2gps() (in module qnav.util.transformations)@\spxentry{utc2gps()}\spxextra{in module qnav.util.transformations}}

\begin{fulllineitems}
\phantomsection\label{\detokenize{modules/util/transformations:qnav.util.transformations.utc2gps}}
\pysigstartsignatures
\pysiglinewithargsret{\sphinxcode{\sphinxupquote{qnav.util.transformations.}}\sphinxbfcode{\sphinxupquote{utc2gps}}}{\sphinxparam{\DUrole{n}{time}\DUrole{p}{:}\DUrole{w}{ }\DUrole{n}{datetime}}}{{ $\rightarrow$ dict}}
\pysigstopsignatures
\sphinxAtStartPar
This function converts a provided time given in Coordinated Universal Time
(UTC) to a Python dictionary containing the corresponding week number, day
number and second.
\begin{quote}\begin{description}
\sphinxlineitem{Parameters}
\sphinxAtStartPar
\sphinxstyleliteralstrong{\sphinxupquote{time}} (\sphinxstyleliteralemphasis{\sphinxupquote{Python Datetime instance}}) \textendash{} A requested time in Coordinated Universal Time (UTC) format.

\sphinxlineitem{Returns}
\sphinxAtStartPar
The provided time as a Python dictionary that can be used for
GPS functions.

\sphinxlineitem{Return type}
\sphinxAtStartPar
dict

\end{description}\end{quote}

\end{fulllineitems}

\index{ned2azel() (in module qnav.util.transformations)@\spxentry{ned2azel()}\spxextra{in module qnav.util.transformations}}

\begin{fulllineitems}
\phantomsection\label{\detokenize{modules/util/transformations:qnav.util.transformations.ned2azel}}
\pysigstartsignatures
\pysiglinewithargsret{\sphinxcode{\sphinxupquote{qnav.util.transformations.}}\sphinxbfcode{\sphinxupquote{ned2azel}}}{\sphinxparam{\DUrole{n}{ned}\DUrole{p}{:}\DUrole{w}{ }\DUrole{n}{ndarray}}}{{ $\rightarrow$ tuple\DUrole{p}{{[}}float\DUrole{p}{,}\DUrole{w}{ }float\DUrole{p}{{]}}}}
\pysigstopsignatures
\sphinxAtStartPar
This function takes an origin North\sphinxhyphen{}East\sphinxhyphen{}Down (NED) position and converts
it into an azimuth and elevation value.
\begin{quote}\begin{description}
\sphinxlineitem{Parameters}
\sphinxAtStartPar
\sphinxstyleliteralstrong{\sphinxupquote{ned}} (\sphinxstyleliteralemphasis{\sphinxupquote{numpy.ndarray}}\sphinxstyleliteralemphasis{\sphinxupquote{ (}}\sphinxstyleliteralemphasis{\sphinxupquote{3\sphinxhyphen{}elements}}\sphinxstyleliteralemphasis{\sphinxupquote{)}}) \textendash{} The centre location (array containing {[}north, east, down{]})
with all values being given in metres.

\sphinxlineitem{Returns}
\sphinxAtStartPar
A tuple containing two elements, the corresponding azimuth and
elevation in radians.

\sphinxlineitem{Return type}
\sphinxAtStartPar
tuple{[}float, float{]}

\end{description}\end{quote}

\end{fulllineitems}

\index{get\_transport\_rate() (in module qnav.util.transformations)@\spxentry{get\_transport\_rate()}\spxextra{in module qnav.util.transformations}}

\begin{fulllineitems}
\phantomsection\label{\detokenize{modules/util/transformations:qnav.util.transformations.get_transport_rate}}
\pysigstartsignatures
\pysiglinewithargsret{\sphinxcode{\sphinxupquote{qnav.util.transformations.}}\sphinxbfcode{\sphinxupquote{get\_transport\_rate}}}{\sphinxparam{\DUrole{n}{position\_lla}\DUrole{p}{:}\DUrole{w}{ }\DUrole{n}{ndarray}}\sphinxparamcomma \sphinxparam{\DUrole{n}{velocity\_ned}\DUrole{p}{:}\DUrole{w}{ }\DUrole{n}{ndarray}}}{{ $\rightarrow$ ndarray}}
\pysigstopsignatures
\sphinxAtStartPar
This function computes the transport rate (in Local NED axes) for given
position and velocity vectors.
\begin{quote}\begin{description}
\sphinxlineitem{Parameters}\begin{itemize}
\item {} 
\sphinxAtStartPar
\sphinxstyleliteralstrong{\sphinxupquote{position\_lla}} (\sphinxstyleliteralemphasis{\sphinxupquote{numpy.ndarray}}\sphinxstyleliteralemphasis{\sphinxupquote{ (}}\sphinxstyleliteralemphasis{\sphinxupquote{3\sphinxhyphen{}elements}}\sphinxstyleliteralemphasis{\sphinxupquote{)}}) \textendash{} The current position vector stating the latitude,
longitude and altitude (degrees\sphinxhyphen{}degrees\sphinxhyphen{}metres).

\item {} 
\sphinxAtStartPar
\sphinxstyleliteralstrong{\sphinxupquote{velocity\_ned}} (\sphinxstyleliteralemphasis{\sphinxupquote{numpy.ndarray}}\sphinxstyleliteralemphasis{\sphinxupquote{ (}}\sphinxstyleliteralemphasis{\sphinxupquote{3\sphinxhyphen{}elements}}\sphinxstyleliteralemphasis{\sphinxupquote{)}}) \textendash{} The current velocity vector stating the velocity
North, East and Down (given in metres/second).

\end{itemize}

\sphinxlineitem{Returns}
\sphinxAtStartPar
The transport rate for the given position and velocity.

\sphinxlineitem{Return type}
\sphinxAtStartPar
NumPy Matrix (3\sphinxhyphen{}by\sphinxhyphen{}3 elements)

\end{description}\end{quote}

\end{fulllineitems}

\index{get\_transport\_rate\_vec() (in module qnav.util.transformations)@\spxentry{get\_transport\_rate\_vec()}\spxextra{in module qnav.util.transformations}}

\begin{fulllineitems}
\phantomsection\label{\detokenize{modules/util/transformations:qnav.util.transformations.get_transport_rate_vec}}
\pysigstartsignatures
\pysiglinewithargsret{\sphinxcode{\sphinxupquote{qnav.util.transformations.}}\sphinxbfcode{\sphinxupquote{get\_transport\_rate\_vec}}}{\sphinxparam{\DUrole{n}{position\_lla}\DUrole{p}{:}\DUrole{w}{ }\DUrole{n}{ndarray}}\sphinxparamcomma \sphinxparam{\DUrole{n}{velocity\_ned}\DUrole{p}{:}\DUrole{w}{ }\DUrole{n}{ndarray}}}{{ $\rightarrow$ ndarray}}
\pysigstopsignatures
\sphinxAtStartPar
This function is the vectorised version of the
{\hyperref[\detokenize{modules/util/transformations:qnav.util.transformations.get_transport_rate}]{\sphinxcrossref{\sphinxcode{\sphinxupquote{get\_transport\_rate()}}}}}
function. This computes the transport rate (in Local NED axes)
for a series of given positions and velocity vectors.
\begin{quote}\begin{description}
\sphinxlineitem{Parameters}\begin{itemize}
\item {} 
\sphinxAtStartPar
\sphinxstyleliteralstrong{\sphinxupquote{position\_lla}} (\sphinxstyleliteralemphasis{\sphinxupquote{numpy.ndarray}}\sphinxstyleliteralemphasis{\sphinxupquote{ (}}\sphinxstyleliteralemphasis{\sphinxupquote{n\sphinxhyphen{}by\sphinxhyphen{}3 elements}}\sphinxstyleliteralemphasis{\sphinxupquote{)}}) \textendash{} The current position vectors stating the latitude,
longitude and altitude (degrees\sphinxhyphen{}degrees\sphinxhyphen{}metres).

\item {} 
\sphinxAtStartPar
\sphinxstyleliteralstrong{\sphinxupquote{velocity\_ned}} (\sphinxstyleliteralemphasis{\sphinxupquote{numpy.ndarray}}\sphinxstyleliteralemphasis{\sphinxupquote{ (}}\sphinxstyleliteralemphasis{\sphinxupquote{n\sphinxhyphen{}by\sphinxhyphen{}3 elements}}\sphinxstyleliteralemphasis{\sphinxupquote{)}}) \textendash{} The current velocity vectors stating the velocity
North, East and Down (given in metres/second).

\end{itemize}

\sphinxlineitem{Returns}
\sphinxAtStartPar
The transport rate for the given positions and velocities.

\sphinxlineitem{Return type}
\sphinxAtStartPar
NumPy Matrix (n\sphinxhyphen{}3\sphinxhyphen{}3 elements)

\end{description}\end{quote}

\end{fulllineitems}

\index{vel\_ecef2vel\_ned() (in module qnav.util.transformations)@\spxentry{vel\_ecef2vel\_ned()}\spxextra{in module qnav.util.transformations}}

\begin{fulllineitems}
\phantomsection\label{\detokenize{modules/util/transformations:qnav.util.transformations.vel_ecef2vel_ned}}
\pysigstartsignatures
\pysiglinewithargsret{\sphinxcode{\sphinxupquote{qnav.util.transformations.}}\sphinxbfcode{\sphinxupquote{vel\_ecef2vel\_ned}}}{\sphinxparam{\DUrole{n}{vel\_ecef}\DUrole{p}{:}\DUrole{w}{ }\DUrole{n}{ndarray}}\sphinxparamcomma \sphinxparam{\DUrole{n}{lla}\DUrole{p}{:}\DUrole{w}{ }\DUrole{n}{ndarray}}}{{ $\rightarrow$ ndarray}}
\pysigstopsignatures
\sphinxAtStartPar
A function to take ECEF velocity and an LLA location and convert it into
the NED velocity.
\begin{quote}\begin{description}
\sphinxlineitem{Parameters}\begin{itemize}
\item {} 
\sphinxAtStartPar
\sphinxstyleliteralstrong{\sphinxupquote{vel\_ecef}} (\sphinxstyleliteralemphasis{\sphinxupquote{numpy.ndarray}}\sphinxstyleliteralemphasis{\sphinxupquote{ (}}\sphinxstyleliteralemphasis{\sphinxupquote{3\sphinxhyphen{}elements}}\sphinxstyleliteralemphasis{\sphinxupquote{)}}) \textendash{} The velocity given in Earth\sphinxhyphen{}Centred Earth\sphinxhyphen{}Fixed format.

\item {} 
\sphinxAtStartPar
\sphinxstyleliteralstrong{\sphinxupquote{lla}} \textendash{} The current latitude, longitude and altitude (in deg\sphinxhyphen{}deg\sphinxhyphen{}m).

\end{itemize}

\sphinxlineitem{Returns}
\sphinxAtStartPar
The given Earth\sphinxhyphen{}Centred Earth\sphinxhyphen{}Fixed (ECEF) velocity converted
into local North\sphinxhyphen{}East\sphinxhyphen{}Down (NED), subject to the given reference
position.

\sphinxlineitem{Return type}
\sphinxAtStartPar
numpy.ndarray (3\sphinxhyphen{}elements)

\end{description}\end{quote}

\end{fulllineitems}

\index{cross\_prod\_xy() (in module qnav.util.transformations)@\spxentry{cross\_prod\_xy()}\spxextra{in module qnav.util.transformations}}

\begin{fulllineitems}
\phantomsection\label{\detokenize{modules/util/transformations:qnav.util.transformations.cross_prod_xy}}
\pysigstartsignatures
\pysiglinewithargsret{\sphinxcode{\sphinxupquote{qnav.util.transformations.}}\sphinxbfcode{\sphinxupquote{cross\_prod\_xy}}}{\sphinxparam{\DUrole{n}{vec1}\DUrole{p}{:}\DUrole{w}{ }\DUrole{n}{ndarray}}\sphinxparamcomma \sphinxparam{\DUrole{n}{vec2}\DUrole{p}{:}\DUrole{w}{ }\DUrole{n}{ndarray}}}{{ $\rightarrow$ ndarray}}
\pysigstopsignatures
\sphinxAtStartPar
An optimised method for computing the cross\sphinxhyphen{}product of two 1\sphinxhyphen{}by\sphinxhyphen{}3 arrays.
In most cases this method is far quicker that the Numpy cross function.
\begin{quote}\begin{description}
\sphinxlineitem{Parameters}\begin{itemize}
\item {} 
\sphinxAtStartPar
\sphinxstyleliteralstrong{\sphinxupquote{vec1}} (\sphinxstyleliteralemphasis{\sphinxupquote{numpy.ndarray}}\sphinxstyleliteralemphasis{\sphinxupquote{ (}}\sphinxstyleliteralemphasis{\sphinxupquote{3 elements}}\sphinxstyleliteralemphasis{\sphinxupquote{)}}) \textendash{} The first 1\sphinxhyphen{}by\sphinxhyphen{}3 tuple/array.

\item {} 
\sphinxAtStartPar
\sphinxstyleliteralstrong{\sphinxupquote{vec2}} (\sphinxstyleliteralemphasis{\sphinxupquote{numpy.ndarray}}\sphinxstyleliteralemphasis{\sphinxupquote{ (}}\sphinxstyleliteralemphasis{\sphinxupquote{3 elements}}\sphinxstyleliteralemphasis{\sphinxupquote{)}}) \textendash{} The first 1\sphinxhyphen{}by\sphinxhyphen{}3 tuple/array.

\end{itemize}

\sphinxlineitem{Returns}
\sphinxAtStartPar
The cross product of the two tuples.

\sphinxlineitem{Return type}
\sphinxAtStartPar
numpy.ndarray (3 elements)

\end{description}\end{quote}

\end{fulllineitems}


\sphinxstepscope


\subsubsection{qnav.earth.dted}
\label{\detokenize{modules/earth/dted:module-qnav.earth.dted}}\label{\detokenize{modules/earth/dted:qnav-earth-dted}}\label{\detokenize{modules/earth/dted::doc}}\index{module@\spxentry{module}!qnav.earth.dted@\spxentry{qnav.earth.dted}}\index{qnav.earth.dted@\spxentry{qnav.earth.dted}!module@\spxentry{module}}

\paragraph{dted.py}
\label{\detokenize{modules/earth/dted:dted-py}}\begin{quote}\begin{description}
\sphinxlineitem{summary}
\sphinxAtStartPar
A module for managing Digital Terrain Elevation Data (DTED).
This is responsible for the reading and dynamic loading of DTED files for
extracting and interpolating elevation from a locally stored Digital
Elevation Map (DEM). This allows the extraction of the terrain height
above the WSG84 ellipsoid, excluding man\sphinxhyphen{}made features. This is namely
used for waypoint/trajectory generation of man\sphinxhyphen{}made features.

\sphinxlineitem{authors}
\begin{DUlineblock}{0em}
\item[] Michael Wright \sphinxhyphen{} \sphinxhref{mailto:mjwright@liverpool.ac.uk}{mjwright@liverpool.ac.uk}
\end{DUlineblock}

\sphinxlineitem{version}
\sphinxAtStartPar
1.0.0 \sphinxhyphen{} First full implementation of this module.

\end{description}\end{quote}
\index{DtedTile (class in qnav.earth.dted)@\spxentry{DtedTile}\spxextra{class in qnav.earth.dted}}

\begin{fulllineitems}
\phantomsection\label{\detokenize{modules/earth/dted:qnav.earth.dted.DtedTile}}
\pysigstartsignatures
\pysiglinewithargsret{\sphinxbfcode{\sphinxupquote{class\DUrole{w}{ }}}\sphinxcode{\sphinxupquote{qnav.earth.dted.}}\sphinxbfcode{\sphinxupquote{DtedTile}}}{\sphinxparam{\DUrole{n}{origin\_lat}\DUrole{p}{:}\DUrole{w}{ }\DUrole{n}{int}}\sphinxparamcomma \sphinxparam{\DUrole{n}{origin\_lon}\DUrole{p}{:}\DUrole{w}{ }\DUrole{n}{int}}\sphinxparamcomma \sphinxparam{\DUrole{n}{dted\_dir}\DUrole{p}{:}\DUrole{w}{ }\DUrole{n}{Path}}\sphinxparamcomma \sphinxparam{\DUrole{n}{dted\_version}\DUrole{p}{:}\DUrole{w}{ }\DUrole{n}{int}\DUrole{w}{ }\DUrole{o}{=}\DUrole{w}{ }\DUrole{default_value}{2}}\sphinxparamcomma \sphinxparam{\DUrole{n}{interp\_method}\DUrole{p}{:}\DUrole{w}{ }\DUrole{n}{str}\DUrole{w}{ }\DUrole{o}{=}\DUrole{w}{ }\DUrole{default_value}{\textquotesingle{}linear\textquotesingle{}}}\sphinxparamcomma \sphinxparam{\DUrole{n}{is\_empty\_on\_error}\DUrole{p}{:}\DUrole{w}{ }\DUrole{n}{bool}\DUrole{w}{ }\DUrole{o}{=}\DUrole{w}{ }\DUrole{default_value}{True}}}{}
\pysigstopsignatures
\sphinxAtStartPar
Holds and handles the read content for a single DTED file.
Manages the data contained in a DTED tile (1\sphinxhyphen{}degree by 1\sphinxhyphen{}degree region).
This allows local interpolation and other relevant calculations. Instances
of this are used by \sphinxcode{\sphinxupquote{DtedHandler}} for dynamic loading and
management of terrain information.
\index{get\_elevation() (qnav.earth.dted.DtedTile method)@\spxentry{get\_elevation()}\spxextra{qnav.earth.dted.DtedTile method}}

\begin{fulllineitems}
\phantomsection\label{\detokenize{modules/earth/dted:qnav.earth.dted.DtedTile.get_elevation}}
\pysigstartsignatures
\pysiglinewithargsret{\sphinxbfcode{\sphinxupquote{get\_elevation}}}{\sphinxparam{\DUrole{n}{lat}\DUrole{p}{:}\DUrole{w}{ }\DUrole{n}{float}}\sphinxparamcomma \sphinxparam{\DUrole{n}{lon}\DUrole{p}{:}\DUrole{w}{ }\DUrole{n}{float}}\sphinxparamcomma \sphinxparam{\DUrole{n}{method}\DUrole{p}{:}\DUrole{w}{ }\DUrole{n}{str}\DUrole{w}{ }\DUrole{o}{=}\DUrole{w}{ }\DUrole{default_value}{None}}}{{ $\rightarrow$ float}}
\pysigstopsignatures
\sphinxAtStartPar
Interpolates the terrain elevation for the requested position.
Given a latitude and longitude, the interpolated terrain elevation
will be return (in metres). Optionally, an overriding interpolation
method to use can be provided.
\begin{quote}\begin{description}
\sphinxlineitem{Parameters}\begin{itemize}
\item {} 
\sphinxAtStartPar
\sphinxstyleliteralstrong{\sphinxupquote{lat}} (\sphinxstyleliteralemphasis{\sphinxupquote{float}}) \textendash{} The latitude coordinate of the reference position.

\item {} 
\sphinxAtStartPar
\sphinxstyleliteralstrong{\sphinxupquote{lon}} (\sphinxstyleliteralemphasis{\sphinxupquote{float}}) \textendash{} The longitude coordinate of the reference position.

\item {} 
\sphinxAtStartPar
\sphinxstyleliteralstrong{\sphinxupquote{method}} (\sphinxstyleliteralemphasis{\sphinxupquote{str}}) \textendash{} 
\sphinxAtStartPar
An alternative interpolation method to use. Supported
values are: ‘linear’, ‘nearest’, ‘slinear’, ‘cubic’,
\begin{quote}

\sphinxAtStartPar
’quintic’ and ‘pchip’.
\end{quote}


\end{itemize}

\sphinxlineitem{Returns}
\sphinxAtStartPar
The interpolated terrain elevation (in metres)

\sphinxlineitem{Return type}
\sphinxAtStartPar
float

\end{description}\end{quote}

\end{fulllineitems}

\index{get\_elevation\_vec() (qnav.earth.dted.DtedTile method)@\spxentry{get\_elevation\_vec()}\spxextra{qnav.earth.dted.DtedTile method}}

\begin{fulllineitems}
\phantomsection\label{\detokenize{modules/earth/dted:qnav.earth.dted.DtedTile.get_elevation_vec}}
\pysigstartsignatures
\pysiglinewithargsret{\sphinxbfcode{\sphinxupquote{get\_elevation\_vec}}}{\sphinxparam{\DUrole{n}{lat\_points}\DUrole{p}{:}\DUrole{w}{ }\DUrole{n}{Buffer\DUrole{w}{ }\DUrole{p}{|}\DUrole{w}{ }\_SupportsArray\DUrole{p}{{[}}dtype\DUrole{p}{{[}}Any\DUrole{p}{{]}}\DUrole{p}{{]}}\DUrole{w}{ }\DUrole{p}{|}\DUrole{w}{ }\_NestedSequence\DUrole{p}{{[}}\_SupportsArray\DUrole{p}{{[}}dtype\DUrole{p}{{[}}Any\DUrole{p}{{]}}\DUrole{p}{{]}}\DUrole{p}{{]}}\DUrole{w}{ }\DUrole{p}{|}\DUrole{w}{ }bool\DUrole{w}{ }\DUrole{p}{|}\DUrole{w}{ }int\DUrole{w}{ }\DUrole{p}{|}\DUrole{w}{ }float\DUrole{w}{ }\DUrole{p}{|}\DUrole{w}{ }complex\DUrole{w}{ }\DUrole{p}{|}\DUrole{w}{ }str\DUrole{w}{ }\DUrole{p}{|}\DUrole{w}{ }bytes\DUrole{w}{ }\DUrole{p}{|}\DUrole{w}{ }\_NestedSequence\DUrole{p}{{[}}bool\DUrole{w}{ }\DUrole{p}{|}\DUrole{w}{ }int\DUrole{w}{ }\DUrole{p}{|}\DUrole{w}{ }float\DUrole{w}{ }\DUrole{p}{|}\DUrole{w}{ }complex\DUrole{w}{ }\DUrole{p}{|}\DUrole{w}{ }str\DUrole{w}{ }\DUrole{p}{|}\DUrole{w}{ }bytes\DUrole{p}{{]}}}}\sphinxparamcomma \sphinxparam{\DUrole{n}{lon\_points}\DUrole{p}{:}\DUrole{w}{ }\DUrole{n}{Buffer\DUrole{w}{ }\DUrole{p}{|}\DUrole{w}{ }\_SupportsArray\DUrole{p}{{[}}dtype\DUrole{p}{{[}}Any\DUrole{p}{{]}}\DUrole{p}{{]}}\DUrole{w}{ }\DUrole{p}{|}\DUrole{w}{ }\_NestedSequence\DUrole{p}{{[}}\_SupportsArray\DUrole{p}{{[}}dtype\DUrole{p}{{[}}Any\DUrole{p}{{]}}\DUrole{p}{{]}}\DUrole{p}{{]}}\DUrole{w}{ }\DUrole{p}{|}\DUrole{w}{ }bool\DUrole{w}{ }\DUrole{p}{|}\DUrole{w}{ }int\DUrole{w}{ }\DUrole{p}{|}\DUrole{w}{ }float\DUrole{w}{ }\DUrole{p}{|}\DUrole{w}{ }complex\DUrole{w}{ }\DUrole{p}{|}\DUrole{w}{ }str\DUrole{w}{ }\DUrole{p}{|}\DUrole{w}{ }bytes\DUrole{w}{ }\DUrole{p}{|}\DUrole{w}{ }\_NestedSequence\DUrole{p}{{[}}bool\DUrole{w}{ }\DUrole{p}{|}\DUrole{w}{ }int\DUrole{w}{ }\DUrole{p}{|}\DUrole{w}{ }float\DUrole{w}{ }\DUrole{p}{|}\DUrole{w}{ }complex\DUrole{w}{ }\DUrole{p}{|}\DUrole{w}{ }str\DUrole{w}{ }\DUrole{p}{|}\DUrole{w}{ }bytes\DUrole{p}{{]}}}}\sphinxparamcomma \sphinxparam{\DUrole{n}{method}\DUrole{p}{:}\DUrole{w}{ }\DUrole{n}{str}\DUrole{w}{ }\DUrole{o}{=}\DUrole{w}{ }\DUrole{default_value}{None}}}{{ $\rightarrow$ ndarray}}
\pysigstopsignatures
\sphinxAtStartPar
Interpolates the terrain elevation for the requested positions.
Given latitude and longitude coordinates, the interpolated terrain
elevations will be return (in metres). Optionally, an overriding
interpolation method to use can be provided.
\begin{quote}\begin{description}
\sphinxlineitem{Parameters}\begin{itemize}
\item {} 
\sphinxAtStartPar
\sphinxstyleliteralstrong{\sphinxupquote{lat\_points}} (\sphinxstyleliteralemphasis{\sphinxupquote{npt.ArrayLike}}) \textendash{} The latitude coordinates of the reference positions.

\item {} 
\sphinxAtStartPar
\sphinxstyleliteralstrong{\sphinxupquote{lon\_points}} (\sphinxstyleliteralemphasis{\sphinxupquote{npt.ArrayLike}}) \textendash{} The longitude coordinates of the reference positions.

\item {} 
\sphinxAtStartPar
\sphinxstyleliteralstrong{\sphinxupquote{method}} (\sphinxstyleliteralemphasis{\sphinxupquote{str}}) \textendash{} 
\sphinxAtStartPar
An alternative interpolation method to use. Supported
values are: ‘linear’, ‘nearest’, ‘slinear’, ‘cubic’,
\begin{quote}

\sphinxAtStartPar
’quintic’ and ‘pchip’.
\end{quote}


\end{itemize}

\sphinxlineitem{Returns}
\sphinxAtStartPar
The interpolated terrain elevations (in metres)

\sphinxlineitem{Return type}
\sphinxAtStartPar
npt.ArrayLike

\end{description}\end{quote}

\end{fulllineitems}

\index{is\_position\_covered() (qnav.earth.dted.DtedTile method)@\spxentry{is\_position\_covered()}\spxextra{qnav.earth.dted.DtedTile method}}

\begin{fulllineitems}
\phantomsection\label{\detokenize{modules/earth/dted:qnav.earth.dted.DtedTile.is_position_covered}}
\pysigstartsignatures
\pysiglinewithargsret{\sphinxbfcode{\sphinxupquote{is\_position\_covered}}}{\sphinxparam{\DUrole{n}{lat}\DUrole{p}{:}\DUrole{w}{ }\DUrole{n}{float}}\sphinxparamcomma \sphinxparam{\DUrole{n}{lon}\DUrole{p}{:}\DUrole{w}{ }\DUrole{n}{float}}}{{ $\rightarrow$ bool}}
\pysigstopsignatures
\sphinxAtStartPar
A method for checking if a requested point is covered by the tile.
Performs a check to see if the position is encompassed by the tile,
allowing its terrain heights to be interpolated.
\begin{quote}\begin{description}
\sphinxlineitem{Parameters}\begin{itemize}
\item {} 
\sphinxAtStartPar
\sphinxstyleliteralstrong{\sphinxupquote{lat}} (\sphinxstyleliteralemphasis{\sphinxupquote{float}}) \textendash{} The latitude coordinate of the reference position.

\item {} 
\sphinxAtStartPar
\sphinxstyleliteralstrong{\sphinxupquote{lon}} (\sphinxstyleliteralemphasis{\sphinxupquote{float}}) \textendash{} The longitude coordinate of the reference position.

\end{itemize}

\sphinxlineitem{Returns}
\sphinxAtStartPar
True if the point is covered by the tile.

\sphinxlineitem{Return type}
\sphinxAtStartPar
bool

\end{description}\end{quote}

\end{fulllineitems}

\index{is\_position\_covered\_vec() (qnav.earth.dted.DtedTile method)@\spxentry{is\_position\_covered\_vec()}\spxextra{qnav.earth.dted.DtedTile method}}

\begin{fulllineitems}
\phantomsection\label{\detokenize{modules/earth/dted:qnav.earth.dted.DtedTile.is_position_covered_vec}}
\pysigstartsignatures
\pysiglinewithargsret{\sphinxbfcode{\sphinxupquote{is\_position\_covered\_vec}}}{\sphinxparam{\DUrole{n}{lat}\DUrole{p}{:}\DUrole{w}{ }\DUrole{n}{Buffer\DUrole{w}{ }\DUrole{p}{|}\DUrole{w}{ }\_SupportsArray\DUrole{p}{{[}}dtype\DUrole{p}{{[}}Any\DUrole{p}{{]}}\DUrole{p}{{]}}\DUrole{w}{ }\DUrole{p}{|}\DUrole{w}{ }\_NestedSequence\DUrole{p}{{[}}\_SupportsArray\DUrole{p}{{[}}dtype\DUrole{p}{{[}}Any\DUrole{p}{{]}}\DUrole{p}{{]}}\DUrole{p}{{]}}\DUrole{w}{ }\DUrole{p}{|}\DUrole{w}{ }bool\DUrole{w}{ }\DUrole{p}{|}\DUrole{w}{ }int\DUrole{w}{ }\DUrole{p}{|}\DUrole{w}{ }float\DUrole{w}{ }\DUrole{p}{|}\DUrole{w}{ }complex\DUrole{w}{ }\DUrole{p}{|}\DUrole{w}{ }str\DUrole{w}{ }\DUrole{p}{|}\DUrole{w}{ }bytes\DUrole{w}{ }\DUrole{p}{|}\DUrole{w}{ }\_NestedSequence\DUrole{p}{{[}}bool\DUrole{w}{ }\DUrole{p}{|}\DUrole{w}{ }int\DUrole{w}{ }\DUrole{p}{|}\DUrole{w}{ }float\DUrole{w}{ }\DUrole{p}{|}\DUrole{w}{ }complex\DUrole{w}{ }\DUrole{p}{|}\DUrole{w}{ }str\DUrole{w}{ }\DUrole{p}{|}\DUrole{w}{ }bytes\DUrole{p}{{]}}}}\sphinxparamcomma \sphinxparam{\DUrole{n}{lon}\DUrole{p}{:}\DUrole{w}{ }\DUrole{n}{Buffer\DUrole{w}{ }\DUrole{p}{|}\DUrole{w}{ }\_SupportsArray\DUrole{p}{{[}}dtype\DUrole{p}{{[}}Any\DUrole{p}{{]}}\DUrole{p}{{]}}\DUrole{w}{ }\DUrole{p}{|}\DUrole{w}{ }\_NestedSequence\DUrole{p}{{[}}\_SupportsArray\DUrole{p}{{[}}dtype\DUrole{p}{{[}}Any\DUrole{p}{{]}}\DUrole{p}{{]}}\DUrole{p}{{]}}\DUrole{w}{ }\DUrole{p}{|}\DUrole{w}{ }bool\DUrole{w}{ }\DUrole{p}{|}\DUrole{w}{ }int\DUrole{w}{ }\DUrole{p}{|}\DUrole{w}{ }float\DUrole{w}{ }\DUrole{p}{|}\DUrole{w}{ }complex\DUrole{w}{ }\DUrole{p}{|}\DUrole{w}{ }str\DUrole{w}{ }\DUrole{p}{|}\DUrole{w}{ }bytes\DUrole{w}{ }\DUrole{p}{|}\DUrole{w}{ }\_NestedSequence\DUrole{p}{{[}}bool\DUrole{w}{ }\DUrole{p}{|}\DUrole{w}{ }int\DUrole{w}{ }\DUrole{p}{|}\DUrole{w}{ }float\DUrole{w}{ }\DUrole{p}{|}\DUrole{w}{ }complex\DUrole{w}{ }\DUrole{p}{|}\DUrole{w}{ }str\DUrole{w}{ }\DUrole{p}{|}\DUrole{w}{ }bytes\DUrole{p}{{]}}}}}{{ $\rightarrow$ ndarray}}
\pysigstopsignatures
\sphinxAtStartPar
A vectorised method for checking if points are covered by the tile.
Performs a check to see which positions are encompassed by the tile,
supporting their terrain heights to be interpolated.
\begin{quote}\begin{description}
\sphinxlineitem{Parameters}\begin{itemize}
\item {} 
\sphinxAtStartPar
\sphinxstyleliteralstrong{\sphinxupquote{lat}} (\sphinxstyleliteralemphasis{\sphinxupquote{npt.ArrayLike}}) \textendash{} The latitude coordinates of the reference positions.

\item {} 
\sphinxAtStartPar
\sphinxstyleliteralstrong{\sphinxupquote{lon}} (\sphinxstyleliteralemphasis{\sphinxupquote{npt.ArrayLike}}) \textendash{} The longitude coordinates of the reference positions.

\end{itemize}

\sphinxlineitem{Returns}
\sphinxAtStartPar
A boolean array indicating which points are covered by the
tile. True is used for points inside the tile.

\sphinxlineitem{Return type}
\sphinxAtStartPar
np.ndarray.

\end{description}\end{quote}

\end{fulllineitems}

\index{distance\_to\_centre() (qnav.earth.dted.DtedTile method)@\spxentry{distance\_to\_centre()}\spxextra{qnav.earth.dted.DtedTile method}}

\begin{fulllineitems}
\phantomsection\label{\detokenize{modules/earth/dted:qnav.earth.dted.DtedTile.distance_to_centre}}
\pysigstartsignatures
\pysiglinewithargsret{\sphinxbfcode{\sphinxupquote{distance\_to\_centre}}}{\sphinxparam{\DUrole{n}{lat}\DUrole{p}{:}\DUrole{w}{ }\DUrole{n}{float}}\sphinxparamcomma \sphinxparam{\DUrole{n}{lon}\DUrole{p}{:}\DUrole{w}{ }\DUrole{n}{float}}}{{ $\rightarrow$ float}}
\pysigstopsignatures
\sphinxAtStartPar
Returns approximation to distance from centre of tile (in metres).
Calculates the approximate to distance from centre of tile for the
given position. Namely used to identify the most distant tile.
\begin{quote}\begin{description}
\sphinxlineitem{Parameters}\begin{itemize}
\item {} 
\sphinxAtStartPar
\sphinxstyleliteralstrong{\sphinxupquote{lat}} (\sphinxstyleliteralemphasis{\sphinxupquote{float}}) \textendash{} The latitude coordinate of the reference position.

\item {} 
\sphinxAtStartPar
\sphinxstyleliteralstrong{\sphinxupquote{lon}} (\sphinxstyleliteralemphasis{\sphinxupquote{float}}) \textendash{} The longitude coordinate of the reference position.

\end{itemize}

\sphinxlineitem{Returns}
\sphinxAtStartPar
The distance from the centre of the tile (in metres).

\sphinxlineitem{Return type}
\sphinxAtStartPar
float

\end{description}\end{quote}

\end{fulllineitems}


\end{fulllineitems}

\index{DtedHandler (class in qnav.earth.dted)@\spxentry{DtedHandler}\spxextra{class in qnav.earth.dted}}

\begin{fulllineitems}
\phantomsection\label{\detokenize{modules/earth/dted:qnav.earth.dted.DtedHandler}}
\pysigstartsignatures
\pysiglinewithargsret{\sphinxbfcode{\sphinxupquote{class\DUrole{w}{ }}}\sphinxcode{\sphinxupquote{qnav.earth.dted.}}\sphinxbfcode{\sphinxupquote{DtedHandler}}}{\sphinxparam{\DUrole{n}{dted\_dir}\DUrole{p}{:}\DUrole{w}{ }\DUrole{n}{Path}}\sphinxparamcomma \sphinxparam{\DUrole{n}{dted\_version}\DUrole{p}{:}\DUrole{w}{ }\DUrole{n}{int}\DUrole{w}{ }\DUrole{o}{=}\DUrole{w}{ }\DUrole{default_value}{2}}\sphinxparamcomma \sphinxparam{\DUrole{n}{max\_tiles}\DUrole{p}{:}\DUrole{w}{ }\DUrole{n}{int}\DUrole{w}{ }\DUrole{o}{=}\DUrole{w}{ }\DUrole{default_value}{None}}\sphinxparamcomma \sphinxparam{\DUrole{n}{method}\DUrole{p}{:}\DUrole{w}{ }\DUrole{n}{str}\DUrole{w}{ }\DUrole{o}{=}\DUrole{w}{ }\DUrole{default_value}{None}}}{}
\pysigstopsignatures
\sphinxAtStartPar
Manages the dynamic loading and handling of regions of DTED data.
This class is use dynamically load DTED data for extracting elevation.
This automatically determines the required DTED files and loads into
memory. Constraints on the maximum amount of data to load can be set to
conserve resources.
\index{get\_elevation() (qnav.earth.dted.DtedHandler method)@\spxentry{get\_elevation()}\spxextra{qnav.earth.dted.DtedHandler method}}

\begin{fulllineitems}
\phantomsection\label{\detokenize{modules/earth/dted:qnav.earth.dted.DtedHandler.get_elevation}}
\pysigstartsignatures
\pysiglinewithargsret{\sphinxbfcode{\sphinxupquote{get\_elevation}}}{\sphinxparam{\DUrole{n}{lat}\DUrole{p}{:}\DUrole{w}{ }\DUrole{n}{float}}\sphinxparamcomma \sphinxparam{\DUrole{n}{lon}\DUrole{p}{:}\DUrole{w}{ }\DUrole{n}{float}}\sphinxparamcomma \sphinxparam{\DUrole{n}{method}\DUrole{p}{:}\DUrole{w}{ }\DUrole{n}{str}\DUrole{w}{ }\DUrole{o}{=}\DUrole{w}{ }\DUrole{default_value}{None}}}{{ $\rightarrow$ float}}
\pysigstopsignatures
\sphinxAtStartPar
Interpolates the terrain elevation for the requested position.
Performed interpolation of the terrain elevation for the given
latitude and longitude location. Optionally an overriding
interpolation method can be used.
\begin{quote}\begin{description}
\sphinxlineitem{Parameters}\begin{itemize}
\item {} 
\sphinxAtStartPar
\sphinxstyleliteralstrong{\sphinxupquote{lat}} (\sphinxstyleliteralemphasis{\sphinxupquote{float}}) \textendash{} The latitude coordinates of the reference positions.

\item {} 
\sphinxAtStartPar
\sphinxstyleliteralstrong{\sphinxupquote{lon}} (\sphinxstyleliteralemphasis{\sphinxupquote{float}}) \textendash{} The longitude coordinates of the reference positions.

\item {} 
\sphinxAtStartPar
\sphinxstyleliteralstrong{\sphinxupquote{method}} (\sphinxstyleliteralemphasis{\sphinxupquote{str}}) \textendash{} The overriding interpolation method to use. Supported
methods are: ‘linear’, ‘nearest’, ‘slinear’, ‘cubic’,
‘quintic’ and ‘pchip’.

\end{itemize}

\sphinxlineitem{Returns}
\sphinxAtStartPar
The interpolated terrain elevation (in metres).

\sphinxlineitem{Return type}
\sphinxAtStartPar
float

\end{description}\end{quote}

\end{fulllineitems}

\index{get\_elevation\_vec() (qnav.earth.dted.DtedHandler method)@\spxentry{get\_elevation\_vec()}\spxextra{qnav.earth.dted.DtedHandler method}}

\begin{fulllineitems}
\phantomsection\label{\detokenize{modules/earth/dted:qnav.earth.dted.DtedHandler.get_elevation_vec}}
\pysigstartsignatures
\pysiglinewithargsret{\sphinxbfcode{\sphinxupquote{get\_elevation\_vec}}}{\sphinxparam{\DUrole{n}{lat\_points}\DUrole{p}{:}\DUrole{w}{ }\DUrole{n}{ndarray}}\sphinxparamcomma \sphinxparam{\DUrole{n}{lon\_points}\DUrole{p}{:}\DUrole{w}{ }\DUrole{n}{ndarray}}\sphinxparamcomma \sphinxparam{\DUrole{n}{method}\DUrole{p}{:}\DUrole{w}{ }\DUrole{n}{str}\DUrole{w}{ }\DUrole{o}{=}\DUrole{w}{ }\DUrole{default_value}{None}}}{{ $\rightarrow$ ndarray}}
\pysigstopsignatures
\sphinxAtStartPar
Interpolates the terrain elevation at given positions.
Performed vectorised interpolation of the terrain elevation for the
given latitude and longitude locations. Optionally an overriding
\begin{quote}

\sphinxAtStartPar
interpolation method can be used.
\end{quote}
\begin{quote}\begin{description}
\sphinxlineitem{Parameters}\begin{itemize}
\item {} 
\sphinxAtStartPar
\sphinxstyleliteralstrong{\sphinxupquote{lat\_points}} (\sphinxstyleliteralemphasis{\sphinxupquote{np.ndarray}}) \textendash{} The latitude coordinates of the reference positions.

\item {} 
\sphinxAtStartPar
\sphinxstyleliteralstrong{\sphinxupquote{lon\_points}} (\sphinxstyleliteralemphasis{\sphinxupquote{np.ndarray}}) \textendash{} The longitude coordinates of the reference positions.

\item {} 
\sphinxAtStartPar
\sphinxstyleliteralstrong{\sphinxupquote{method}} (\sphinxstyleliteralemphasis{\sphinxupquote{str}}) \textendash{} The overriding interpolation method to use. Supported
methods are: ‘linear’, ‘nearest’, ‘slinear’, ‘cubic’,
‘quintic’ and ‘pchip’.

\end{itemize}

\sphinxlineitem{Returns}
\sphinxAtStartPar
The interpolated terrain elevation (in metres).

\sphinxlineitem{Return type}
\sphinxAtStartPar
np.ndarray

\end{description}\end{quote}

\end{fulllineitems}

\index{set\_max\_tiles() (qnav.earth.dted.DtedHandler method)@\spxentry{set\_max\_tiles()}\spxextra{qnav.earth.dted.DtedHandler method}}

\begin{fulllineitems}
\phantomsection\label{\detokenize{modules/earth/dted:qnav.earth.dted.DtedHandler.set_max_tiles}}
\pysigstartsignatures
\pysiglinewithargsret{\sphinxbfcode{\sphinxupquote{set\_max\_tiles}}}{\sphinxparam{\DUrole{n}{value}\DUrole{p}{:}\DUrole{w}{ }\DUrole{n}{int}}}{}
\pysigstopsignatures
\sphinxAtStartPar
A setter method for the maximum number of tiles.
Can use the value None to remove limit on number of tiles to hold.
\begin{quote}\begin{description}
\sphinxlineitem{Parameters}
\sphinxAtStartPar
\sphinxstyleliteralstrong{\sphinxupquote{value}} (\sphinxstyleliteralemphasis{\sphinxupquote{int}}) \textendash{} The maximum number of DTED tiles to hold in the cache.

\end{description}\end{quote}

\end{fulllineitems}

\index{clear\_cache() (qnav.earth.dted.DtedHandler method)@\spxentry{clear\_cache()}\spxextra{qnav.earth.dted.DtedHandler method}}

\begin{fulllineitems}
\phantomsection\label{\detokenize{modules/earth/dted:qnav.earth.dted.DtedHandler.clear_cache}}
\pysigstartsignatures
\pysiglinewithargsret{\sphinxbfcode{\sphinxupquote{clear\_cache}}}{}{}
\pysigstopsignatures
\sphinxAtStartPar
Clears all cached DTED from memory.
Simply removes all DTED tiles from the cache.

\end{fulllineitems}


\end{fulllineitems}

\index{get\_dted\_file\_path() (in module qnav.earth.dted)@\spxentry{get\_dted\_file\_path()}\spxextra{in module qnav.earth.dted}}

\begin{fulllineitems}
\phantomsection\label{\detokenize{modules/earth/dted:qnav.earth.dted.get_dted_file_path}}
\pysigstartsignatures
\pysiglinewithargsret{\sphinxcode{\sphinxupquote{qnav.earth.dted.}}\sphinxbfcode{\sphinxupquote{get\_dted\_file\_path}}}{\sphinxparam{\DUrole{n}{dted\_dir}\DUrole{p}{:}\DUrole{w}{ }\DUrole{n}{Path}}\sphinxparamcomma \sphinxparam{\DUrole{n}{origin\_lat}\DUrole{p}{:}\DUrole{w}{ }\DUrole{n}{int}}\sphinxparamcomma \sphinxparam{\DUrole{n}{origin\_lon}\DUrole{p}{:}\DUrole{w}{ }\DUrole{n}{int}}\sphinxparamcomma \sphinxparam{\DUrole{n}{dted\_version}\DUrole{p}{:}\DUrole{w}{ }\DUrole{n}{int}}}{{ $\rightarrow$ Path}}
\pysigstopsignatures
\sphinxAtStartPar
Give a latitude, longitude and optional DTED version number,
this function generates the file path of the DTED file that would
contain information for the given position.
\begin{quote}\begin{description}
\sphinxlineitem{Parameters}\begin{itemize}
\item {} 
\sphinxAtStartPar
\sphinxstyleliteralstrong{\sphinxupquote{dted\_dir}} (\sphinxstyleliteralemphasis{\sphinxupquote{Path}}) \textendash{} The parent database\_dir of the DTED files.

\item {} 
\sphinxAtStartPar
\sphinxstyleliteralstrong{\sphinxupquote{origin\_lat}} (\sphinxstyleliteralemphasis{\sphinxupquote{int}}) \textendash{} The integer latitude value (south\sphinxhyphen{}most coordinate).

\item {} 
\sphinxAtStartPar
\sphinxstyleliteralstrong{\sphinxupquote{origin\_lon}} (\sphinxstyleliteralemphasis{\sphinxupquote{int}}) \textendash{} The integer longitude value (west\sphinxhyphen{}most coordinate).

\item {} 
\sphinxAtStartPar
\sphinxstyleliteralstrong{\sphinxupquote{dted\_version}} (\sphinxstyleliteralemphasis{\sphinxupquote{int}}) \textendash{} The version of the DTED data.

\end{itemize}

\end{description}\end{quote}

\sphinxAtStartPar
:return the location of a corresponding DTED file.
:rtype: Path

\end{fulllineitems}

\index{read\_dted\_file() (in module qnav.earth.dted)@\spxentry{read\_dted\_file()}\spxextra{in module qnav.earth.dted}}

\begin{fulllineitems}
\phantomsection\label{\detokenize{modules/earth/dted:qnav.earth.dted.read_dted_file}}
\pysigstartsignatures
\pysiglinewithargsret{\sphinxcode{\sphinxupquote{qnav.earth.dted.}}\sphinxbfcode{\sphinxupquote{read\_dted\_file}}}{\sphinxparam{\DUrole{n}{file\_dir}\DUrole{p}{:}\DUrole{w}{ }\DUrole{n}{Path}}}{{ $\rightarrow$ dict}}
\pysigstopsignatures
\sphinxAtStartPar
Extracts elevation information from a provided DTED file (in metres).
When provided with the location of a DTED file (of any level), this
function will return an elevation matrix and latitude\sphinxhyphen{}longitude spacing
of each point. This function reads DTED files of format version 1.1.
\begin{quote}\begin{description}
\sphinxlineitem{Parameters}
\sphinxAtStartPar
\sphinxstyleliteralstrong{\sphinxupquote{file\_dir}} (\sphinxstyleliteralemphasis{\sphinxupquote{Path}}) \textendash{} The location of the DTED file to be read.

\sphinxlineitem{Returns}
\sphinxAtStartPar
A dictionary containing the contents of the DTED file.
This contains the extracted origin coordinates, step sizes,
coordinate ticks and elevation matrix.

\sphinxlineitem{Return type}
\sphinxAtStartPar
dict

\end{description}\end{quote}

\end{fulllineitems}

\index{get\_template\_dted\_data() (in module qnav.earth.dted)@\spxentry{get\_template\_dted\_data()}\spxextra{in module qnav.earth.dted}}

\begin{fulllineitems}
\phantomsection\label{\detokenize{modules/earth/dted:qnav.earth.dted.get_template_dted_data}}
\pysigstartsignatures
\pysiglinewithargsret{\sphinxcode{\sphinxupquote{qnav.earth.dted.}}\sphinxbfcode{\sphinxupquote{get\_template\_dted\_data}}}{\sphinxparam{\DUrole{n}{origin\_lat}\DUrole{p}{:}\DUrole{w}{ }\DUrole{n}{int}}\sphinxparamcomma \sphinxparam{\DUrole{n}{origin\_lon}\DUrole{p}{:}\DUrole{w}{ }\DUrole{n}{int}}\sphinxparamcomma \sphinxparam{\DUrole{n}{num\_lat\_points}\DUrole{p}{:}\DUrole{w}{ }\DUrole{n}{int}\DUrole{w}{ }\DUrole{o}{=}\DUrole{w}{ }\DUrole{default_value}{2}}\sphinxparamcomma \sphinxparam{\DUrole{n}{num\_lon\_points}\DUrole{p}{:}\DUrole{w}{ }\DUrole{n}{int}\DUrole{w}{ }\DUrole{o}{=}\DUrole{w}{ }\DUrole{default_value}{2}}}{{ $\rightarrow$ dict}}
\pysigstopsignatures
\sphinxAtStartPar
Generates a template DTED data dictionary as if read from a file.
Simulates the output of \sphinxcode{\sphinxupquote{read\_dted\_file()}} by generating
identical output structure with flat elevation. This is to be used as
substitution for missing/unavailable DTED files.
\begin{quote}\begin{description}
\sphinxlineitem{Parameters}\begin{itemize}
\item {} 
\sphinxAtStartPar
\sphinxstyleliteralstrong{\sphinxupquote{origin\_lat}} (\sphinxstyleliteralemphasis{\sphinxupquote{int}}) \textendash{} The integer latitude value (south\sphinxhyphen{}most coordinate).

\item {} 
\sphinxAtStartPar
\sphinxstyleliteralstrong{\sphinxupquote{origin\_lon}} (\sphinxstyleliteralemphasis{\sphinxupquote{int}}) \textendash{} The integer longitude value (west\sphinxhyphen{}most coordinate).

\item {} 
\sphinxAtStartPar
\sphinxstyleliteralstrong{\sphinxupquote{num\_lat\_points}} (\sphinxstyleliteralemphasis{\sphinxupquote{int}}) \textendash{} The number of latitude points (rows) to generate.

\item {} 
\sphinxAtStartPar
\sphinxstyleliteralstrong{\sphinxupquote{num\_lon\_points}} (\sphinxstyleliteralemphasis{\sphinxupquote{int}}) \textendash{} The number of longitude points (columns) to generate.

\end{itemize}

\sphinxlineitem{Returns}
\sphinxAtStartPar
A dictionary containing contents of a DTED file.
This contains the imitated origin coordinates, step sizes,
coordinate ticks and elevation matrix that would be read.

\sphinxlineitem{Return type}
\sphinxAtStartPar
dict

\end{description}\end{quote}

\end{fulllineitems}


\sphinxstepscope


\subsubsection{qnav.earth.geoid}
\label{\detokenize{modules/earth/geoid:qnav-earth-geoid}}\label{\detokenize{modules/earth/geoid::doc}}\index{module@\spxentry{module}!qnav.earth.geoid@\spxentry{qnav.earth.geoid}}\index{qnav.earth.geoid@\spxentry{qnav.earth.geoid}!module@\spxentry{module}}

\paragraph{geoid.py}
\label{\detokenize{modules/earth/geoid:geoid-py}}\label{\detokenize{modules/earth/geoid:module-qnav.earth.geoid}}\begin{quote}\begin{description}
\sphinxlineitem{summary}
\sphinxAtStartPar
A module for interpolating geoid heights from presented geoid PGM files.
This module provides classes and functions used to extract geoid
information from given PGM encoded geoid files. The most common use for
this module is for interpolating the geoid height at requested positions.

\sphinxlineitem{authors}
\begin{DUlineblock}{0em}
\item[] Michael Wright \sphinxhyphen{} \sphinxhref{mailto:mjwright@liverpool.ac.uk}{mjwright@liverpool.ac.uk}
\end{DUlineblock}

\sphinxlineitem{version}
\sphinxAtStartPar
1.0.0 \sphinxhyphen{} First full implementation of this module.

\end{description}\end{quote}
\index{GeoidModel (class in qnav.earth.geoid)@\spxentry{GeoidModel}\spxextra{class in qnav.earth.geoid}}

\begin{fulllineitems}
\phantomsection\label{\detokenize{modules/earth/geoid:qnav.earth.geoid.GeoidModel}}
\pysigstartsignatures
\pysigline{\sphinxbfcode{\sphinxupquote{class\DUrole{w}{ }}}\sphinxcode{\sphinxupquote{qnav.earth.geoid.}}\sphinxbfcode{\sphinxupquote{GeoidModel}}}
\pysigstopsignatures
\sphinxAtStartPar
Abstract base class for geoid representations.
Defines the method that implemented geoid representations must supply.
\index{get\_height() (qnav.earth.geoid.GeoidModel method)@\spxentry{get\_height()}\spxextra{qnav.earth.geoid.GeoidModel method}}

\begin{fulllineitems}
\phantomsection\label{\detokenize{modules/earth/geoid:qnav.earth.geoid.GeoidModel.get_height}}
\pysigstartsignatures
\pysiglinewithargsret{\sphinxbfcode{\sphinxupquote{abstract\DUrole{w}{ }}}\sphinxbfcode{\sphinxupquote{get\_height}}}{\sphinxparam{\DUrole{n}{lat}\DUrole{p}{:}\DUrole{w}{ }\DUrole{n}{float}}\sphinxparamcomma \sphinxparam{\DUrole{n}{lon}\DUrole{p}{:}\DUrole{w}{ }\DUrole{n}{float}}}{{ $\rightarrow$ float}}
\pysigstopsignatures
\sphinxAtStartPar
Returns geoid height for single latitude and longitude position.
\begin{quote}\begin{description}
\sphinxlineitem{Parameters}\begin{itemize}
\item {} 
\sphinxAtStartPar
\sphinxstyleliteralstrong{\sphinxupquote{lat}} (\sphinxstyleliteralemphasis{\sphinxupquote{float}}) \textendash{} A latitude coordinate in decimal degrees.

\item {} 
\sphinxAtStartPar
\sphinxstyleliteralstrong{\sphinxupquote{lon}} (\sphinxstyleliteralemphasis{\sphinxupquote{float}}) \textendash{} A longitude coordinate in decimal degrees.

\end{itemize}

\sphinxlineitem{Returns}
\sphinxAtStartPar
The geoid height for the corresponding location (in metres).

\sphinxlineitem{Return type}
\sphinxAtStartPar
float

\end{description}\end{quote}

\end{fulllineitems}

\index{get\_height\_vec() (qnav.earth.geoid.GeoidModel method)@\spxentry{get\_height\_vec()}\spxextra{qnav.earth.geoid.GeoidModel method}}

\begin{fulllineitems}
\phantomsection\label{\detokenize{modules/earth/geoid:qnav.earth.geoid.GeoidModel.get_height_vec}}
\pysigstartsignatures
\pysiglinewithargsret{\sphinxbfcode{\sphinxupquote{abstract\DUrole{w}{ }}}\sphinxbfcode{\sphinxupquote{get\_height\_vec}}}{\sphinxparam{\DUrole{n}{lat}\DUrole{p}{:}\DUrole{w}{ }\DUrole{n}{ndarray}}\sphinxparamcomma \sphinxparam{\DUrole{n}{lon}\DUrole{p}{:}\DUrole{w}{ }\DUrole{n}{ndarray}}}{{ $\rightarrow$ ndarray}}
\pysigstopsignatures
\sphinxAtStartPar
Returns geoid height for given latitude and longitude positions.
\begin{quote}\begin{description}
\sphinxlineitem{Parameters}\begin{itemize}
\item {} 
\sphinxAtStartPar
\sphinxstyleliteralstrong{\sphinxupquote{lat}} (\sphinxstyleliteralemphasis{\sphinxupquote{float}}) \textendash{} The latitude coordinates in decimal degrees.

\item {} 
\sphinxAtStartPar
\sphinxstyleliteralstrong{\sphinxupquote{lon}} (\sphinxstyleliteralemphasis{\sphinxupquote{float}}) \textendash{} The longitude coordinates in decimal degrees.

\end{itemize}

\sphinxlineitem{Returns}
\sphinxAtStartPar
The geoid height for the corresponding locations (in metres).

\sphinxlineitem{Return type}
\sphinxAtStartPar
float

\end{description}\end{quote}

\end{fulllineitems}

\index{get\_dov() (qnav.earth.geoid.GeoidModel method)@\spxentry{get\_dov()}\spxextra{qnav.earth.geoid.GeoidModel method}}

\begin{fulllineitems}
\phantomsection\label{\detokenize{modules/earth/geoid:qnav.earth.geoid.GeoidModel.get_dov}}
\pysigstartsignatures
\pysiglinewithargsret{\sphinxbfcode{\sphinxupquote{abstract\DUrole{w}{ }}}\sphinxbfcode{\sphinxupquote{get\_dov}}}{\sphinxparam{\DUrole{n}{lat}\DUrole{p}{:}\DUrole{w}{ }\DUrole{n}{float}}\sphinxparamcomma \sphinxparam{\DUrole{n}{lon}\DUrole{p}{:}\DUrole{w}{ }\DUrole{n}{float}}}{{ $\rightarrow$ tuple\DUrole{p}{{[}}float\DUrole{p}{,}\DUrole{w}{ }float\DUrole{p}{{]}}}}
\pysigstopsignatures
\sphinxAtStartPar
Returns the Deflection of the Vertical (DoV) for requested position.
Calculates and returns the DoV gradients for the north and east
components (in metres) for requested location.
\begin{quote}\begin{description}
\sphinxlineitem{Parameters}\begin{itemize}
\item {} 
\sphinxAtStartPar
\sphinxstyleliteralstrong{\sphinxupquote{lat}} (\sphinxstyleliteralemphasis{\sphinxupquote{float}}) \textendash{} A latitude coordinate in decimal degrees.

\item {} 
\sphinxAtStartPar
\sphinxstyleliteralstrong{\sphinxupquote{lon}} (\sphinxstyleliteralemphasis{\sphinxupquote{float}}) \textendash{} A longitude coordinate in decimal degrees.

\end{itemize}

\sphinxlineitem{Returns}
\sphinxAtStartPar
Tuple containing the north and east components respectively.

\sphinxlineitem{Return type}
\sphinxAtStartPar
tuple{[}float, float{]}

\end{description}\end{quote}

\end{fulllineitems}

\index{get\_dov\_vec() (qnav.earth.geoid.GeoidModel method)@\spxentry{get\_dov\_vec()}\spxextra{qnav.earth.geoid.GeoidModel method}}

\begin{fulllineitems}
\phantomsection\label{\detokenize{modules/earth/geoid:qnav.earth.geoid.GeoidModel.get_dov_vec}}
\pysigstartsignatures
\pysiglinewithargsret{\sphinxbfcode{\sphinxupquote{abstract\DUrole{w}{ }}}\sphinxbfcode{\sphinxupquote{get\_dov\_vec}}}{\sphinxparam{\DUrole{n}{lat}\DUrole{p}{:}\DUrole{w}{ }\DUrole{n}{ndarray}}\sphinxparamcomma \sphinxparam{\DUrole{n}{lon}\DUrole{p}{:}\DUrole{w}{ }\DUrole{n}{ndarray}}}{{ $\rightarrow$ tuple\DUrole{p}{{[}}ndarray\DUrole{p}{,}\DUrole{w}{ }ndarray\DUrole{p}{{]}}}}
\pysigstopsignatures
\sphinxAtStartPar
Returns the Deflection of the Vertical (DoV) for requested positions.
Calculates and returns the DoV gradients for the north and east
components (in metres) for requested location.
\begin{quote}\begin{description}
\sphinxlineitem{Parameters}\begin{itemize}
\item {} 
\sphinxAtStartPar
\sphinxstyleliteralstrong{\sphinxupquote{lat}} (\sphinxstyleliteralemphasis{\sphinxupquote{np.ndarray}}) \textendash{} The latitude coordinates in decimal degrees.

\item {} 
\sphinxAtStartPar
\sphinxstyleliteralstrong{\sphinxupquote{lon}} (\sphinxstyleliteralemphasis{\sphinxupquote{np.ndarray}}) \textendash{} The longitude coordinates in decimal degrees.

\end{itemize}

\sphinxlineitem{Returns}
\sphinxAtStartPar
Tuple containing the north and east components respectively.

\sphinxlineitem{Return type}
\sphinxAtStartPar
tuple{[}float, float{]}

\end{description}\end{quote}

\end{fulllineitems}


\end{fulllineitems}

\index{GeoidPGM (class in qnav.earth.geoid)@\spxentry{GeoidPGM}\spxextra{class in qnav.earth.geoid}}

\begin{fulllineitems}
\phantomsection\label{\detokenize{modules/earth/geoid:qnav.earth.geoid.GeoidPGM}}
\pysigstartsignatures
\pysiglinewithargsret{\sphinxbfcode{\sphinxupquote{class\DUrole{w}{ }}}\sphinxcode{\sphinxupquote{qnav.earth.geoid.}}\sphinxbfcode{\sphinxupquote{GeoidPGM}}}{\sphinxparam{\DUrole{n}{pgm\_file}\DUrole{p}{:}\DUrole{w}{ }\DUrole{n}{Path}}}{}
\pysigstopsignatures
\sphinxAtStartPar
A class for handling and managing reading the Geoid height information
from a selected PGM (Pixel Gray Map) file and providing functionality for
extracting required information from it. Methods for reading or
interpolating the Geoid height at given coordinates, as well as
visualising the loaded data are all contained within this class.
\index{get\_height() (qnav.earth.geoid.GeoidPGM method)@\spxentry{get\_height()}\spxextra{qnav.earth.geoid.GeoidPGM method}}

\begin{fulllineitems}
\phantomsection\label{\detokenize{modules/earth/geoid:qnav.earth.geoid.GeoidPGM.get_height}}
\pysigstartsignatures
\pysiglinewithargsret{\sphinxbfcode{\sphinxupquote{get\_height}}}{\sphinxparam{\DUrole{n}{lat}\DUrole{p}{:}\DUrole{w}{ }\DUrole{n}{float}}\sphinxparamcomma \sphinxparam{\DUrole{n}{lon}\DUrole{p}{:}\DUrole{w}{ }\DUrole{n}{float}}}{{ $\rightarrow$ float}}
\pysigstopsignatures
\sphinxAtStartPar
Interpolates geoid height for given latitude and longitude.
For a single coordinate point, interpolates and returns the geoid
elevation height in metres.
\begin{quote}\begin{description}
\sphinxlineitem{Parameters}\begin{itemize}
\item {} 
\sphinxAtStartPar
\sphinxstyleliteralstrong{\sphinxupquote{lat}} (\sphinxstyleliteralemphasis{\sphinxupquote{float}}) \textendash{} A latitude coordinate in decimal degrees.

\item {} 
\sphinxAtStartPar
\sphinxstyleliteralstrong{\sphinxupquote{lon}} (\sphinxstyleliteralemphasis{\sphinxupquote{float}}) \textendash{} A longitude coordinate in decimal degrees.

\end{itemize}

\sphinxlineitem{Returns}
\sphinxAtStartPar
The geoid height for the corresponding location (in metres).

\sphinxlineitem{Return type}
\sphinxAtStartPar
float

\end{description}\end{quote}

\end{fulllineitems}

\index{get\_height\_vec() (qnav.earth.geoid.GeoidPGM method)@\spxentry{get\_height\_vec()}\spxextra{qnav.earth.geoid.GeoidPGM method}}

\begin{fulllineitems}
\phantomsection\label{\detokenize{modules/earth/geoid:qnav.earth.geoid.GeoidPGM.get_height_vec}}
\pysigstartsignatures
\pysiglinewithargsret{\sphinxbfcode{\sphinxupquote{get\_height\_vec}}}{\sphinxparam{\DUrole{n}{lat}\DUrole{p}{:}\DUrole{w}{ }\DUrole{n}{ndarray}}\sphinxparamcomma \sphinxparam{\DUrole{n}{lon}\DUrole{p}{:}\DUrole{w}{ }\DUrole{n}{ndarray}}}{{ $\rightarrow$ ndarray}}
\pysigstopsignatures
\sphinxAtStartPar
Interpolates geoid height for given latitude and longitude points.
For a single coordinate point, interpolates and returns the geoid
elevation height in metres.
\begin{quote}\begin{description}
\sphinxlineitem{Parameters}\begin{itemize}
\item {} 
\sphinxAtStartPar
\sphinxstyleliteralstrong{\sphinxupquote{lat}} (\sphinxstyleliteralemphasis{\sphinxupquote{float}}) \textendash{} The latitude coordinates in decimal degrees.

\item {} 
\sphinxAtStartPar
\sphinxstyleliteralstrong{\sphinxupquote{lon}} (\sphinxstyleliteralemphasis{\sphinxupquote{float}}) \textendash{} The longitude coordinates in decimal degrees.

\end{itemize}

\sphinxlineitem{Returns}
\sphinxAtStartPar
The geoid height for the corresponding locations (in metres).

\sphinxlineitem{Return type}
\sphinxAtStartPar
float

\end{description}\end{quote}

\end{fulllineitems}

\index{get\_dov() (qnav.earth.geoid.GeoidPGM method)@\spxentry{get\_dov()}\spxextra{qnav.earth.geoid.GeoidPGM method}}

\begin{fulllineitems}
\phantomsection\label{\detokenize{modules/earth/geoid:qnav.earth.geoid.GeoidPGM.get_dov}}
\pysigstartsignatures
\pysiglinewithargsret{\sphinxbfcode{\sphinxupquote{get\_dov}}}{\sphinxparam{\DUrole{n}{lat}\DUrole{p}{:}\DUrole{w}{ }\DUrole{n}{float}}\sphinxparamcomma \sphinxparam{\DUrole{n}{lon}\DUrole{p}{:}\DUrole{w}{ }\DUrole{n}{float}}\sphinxparamcomma \sphinxparam{\DUrole{n}{step\_size}\DUrole{p}{:}\DUrole{w}{ }\DUrole{n}{float}\DUrole{w}{ }\DUrole{o}{=}\DUrole{w}{ }\DUrole{default_value}{100}}}{{ $\rightarrow$ tuple\DUrole{p}{{[}}float\DUrole{p}{,}\DUrole{w}{ }float\DUrole{p}{{]}}}}
\pysigstopsignatures
\sphinxAtStartPar
Returns the Deflection of the Vertical (DoV) for requested position.
Calculates and returns the DoV gradients for the north and east
components (in metres) for requested location.
\begin{quote}\begin{description}
\sphinxlineitem{Parameters}\begin{itemize}
\item {} 
\sphinxAtStartPar
\sphinxstyleliteralstrong{\sphinxupquote{lat}} (\sphinxstyleliteralemphasis{\sphinxupquote{float}}) \textendash{} A latitude coordinate in decimal degrees.

\item {} 
\sphinxAtStartPar
\sphinxstyleliteralstrong{\sphinxupquote{lon}} (\sphinxstyleliteralemphasis{\sphinxupquote{float}}) \textendash{} A longitude coordinate in decimal degrees.

\item {} 
\sphinxAtStartPar
\sphinxstyleliteralstrong{\sphinxupquote{step\_size}} (\sphinxstyleliteralemphasis{\sphinxupquote{float}}) \textendash{} The step size for calculating the gradient.

\end{itemize}

\sphinxlineitem{Returns}
\sphinxAtStartPar
Tuple containing the north and east components respectively.

\sphinxlineitem{Return type}
\sphinxAtStartPar
tuple{[}float, float{]}

\end{description}\end{quote}

\end{fulllineitems}

\index{get\_dov\_vec() (qnav.earth.geoid.GeoidPGM method)@\spxentry{get\_dov\_vec()}\spxextra{qnav.earth.geoid.GeoidPGM method}}

\begin{fulllineitems}
\phantomsection\label{\detokenize{modules/earth/geoid:qnav.earth.geoid.GeoidPGM.get_dov_vec}}
\pysigstartsignatures
\pysiglinewithargsret{\sphinxbfcode{\sphinxupquote{get\_dov\_vec}}}{\sphinxparam{\DUrole{n}{lat}\DUrole{p}{:}\DUrole{w}{ }\DUrole{n}{ndarray}}\sphinxparamcomma \sphinxparam{\DUrole{n}{lon}\DUrole{p}{:}\DUrole{w}{ }\DUrole{n}{ndarray}}\sphinxparamcomma \sphinxparam{\DUrole{n}{step\_size}\DUrole{p}{:}\DUrole{w}{ }\DUrole{n}{float}\DUrole{w}{ }\DUrole{o}{=}\DUrole{w}{ }\DUrole{default_value}{100}}}{{ $\rightarrow$ tuple\DUrole{p}{{[}}ndarray\DUrole{p}{,}\DUrole{w}{ }ndarray\DUrole{p}{{]}}}}
\pysigstopsignatures
\sphinxAtStartPar
Returns the Deflection of the Vertical (DoV) for requested position.
Calculates and returns the DoV gradients for the north and east
components (in metres) for requested location.
\begin{quote}\begin{description}
\sphinxlineitem{Parameters}\begin{itemize}
\item {} 
\sphinxAtStartPar
\sphinxstyleliteralstrong{\sphinxupquote{lat}} (\sphinxstyleliteralemphasis{\sphinxupquote{np.ndarray}}) \textendash{} The latitude coordinates in decimal degrees.

\item {} 
\sphinxAtStartPar
\sphinxstyleliteralstrong{\sphinxupquote{lon}} (\sphinxstyleliteralemphasis{\sphinxupquote{np.ndarray}}) \textendash{} The longitude coordinates in decimal degrees.

\item {} 
\sphinxAtStartPar
\sphinxstyleliteralstrong{\sphinxupquote{step\_size}} (\sphinxstyleliteralemphasis{\sphinxupquote{float}}) \textendash{} The step size for calculating the gradient.

\end{itemize}

\sphinxlineitem{Returns}
\sphinxAtStartPar
Tuple containing the north and east components respectively.

\sphinxlineitem{Return type}
\sphinxAtStartPar
tuple{[}float, float{]}

\end{description}\end{quote}

\end{fulllineitems}

\index{resolution (qnav.earth.geoid.GeoidPGM property)@\spxentry{resolution}\spxextra{qnav.earth.geoid.GeoidPGM property}}

\begin{fulllineitems}
\phantomsection\label{\detokenize{modules/earth/geoid:qnav.earth.geoid.GeoidPGM.resolution}}
\pysigstartsignatures
\pysigline{\sphinxbfcode{\sphinxupquote{property\DUrole{w}{ }}}\sphinxbfcode{\sphinxupquote{resolution}}\sphinxbfcode{\sphinxupquote{\DUrole{p}{:}\DUrole{w}{ }tuple\DUrole{p}{{[}}int\DUrole{p}{,}\DUrole{w}{ }int\DUrole{p}{{]}}}}}
\pysigstopsignatures
\sphinxAtStartPar
The north and east resolutions of the grid data.
Can be used for determining the accuracy of the model.
\begin{quote}\begin{description}
\sphinxlineitem{Returns}
\sphinxAtStartPar
The north and east steps of the grid data.

\sphinxlineitem{Return type}
\sphinxAtStartPar
tuple{[}int, int{]}

\end{description}\end{quote}

\end{fulllineitems}


\end{fulllineitems}

\index{GeoidDummy (class in qnav.earth.geoid)@\spxentry{GeoidDummy}\spxextra{class in qnav.earth.geoid}}

\begin{fulllineitems}
\phantomsection\label{\detokenize{modules/earth/geoid:qnav.earth.geoid.GeoidDummy}}
\pysigstartsignatures
\pysigline{\sphinxbfcode{\sphinxupquote{class\DUrole{w}{ }}}\sphinxcode{\sphinxupquote{qnav.earth.geoid.}}\sphinxbfcode{\sphinxupquote{GeoidDummy}}}
\pysigstopsignatures\index{get\_height() (qnav.earth.geoid.GeoidDummy method)@\spxentry{get\_height()}\spxextra{qnav.earth.geoid.GeoidDummy method}}

\begin{fulllineitems}
\phantomsection\label{\detokenize{modules/earth/geoid:qnav.earth.geoid.GeoidDummy.get_height}}
\pysigstartsignatures
\pysiglinewithargsret{\sphinxbfcode{\sphinxupquote{get\_height}}}{\sphinxparam{\DUrole{n}{lat}\DUrole{p}{:}\DUrole{w}{ }\DUrole{n}{float}}\sphinxparamcomma \sphinxparam{\DUrole{n}{lon}\DUrole{p}{:}\DUrole{w}{ }\DUrole{n}{float}}}{{ $\rightarrow$ float}}
\pysigstopsignatures
\sphinxAtStartPar
Returns geoid height for single latitude and longitude position.
\begin{quote}\begin{description}
\sphinxlineitem{Parameters}\begin{itemize}
\item {} 
\sphinxAtStartPar
\sphinxstyleliteralstrong{\sphinxupquote{lat}} (\sphinxstyleliteralemphasis{\sphinxupquote{float}}) \textendash{} A latitude coordinate in decimal degrees.

\item {} 
\sphinxAtStartPar
\sphinxstyleliteralstrong{\sphinxupquote{lon}} (\sphinxstyleliteralemphasis{\sphinxupquote{float}}) \textendash{} A longitude coordinate in decimal degrees.

\end{itemize}

\sphinxlineitem{Returns}
\sphinxAtStartPar
The geoid height for the corresponding location (in metres).

\sphinxlineitem{Return type}
\sphinxAtStartPar
float

\end{description}\end{quote}

\end{fulllineitems}

\index{get\_height\_vec() (qnav.earth.geoid.GeoidDummy method)@\spxentry{get\_height\_vec()}\spxextra{qnav.earth.geoid.GeoidDummy method}}

\begin{fulllineitems}
\phantomsection\label{\detokenize{modules/earth/geoid:qnav.earth.geoid.GeoidDummy.get_height_vec}}
\pysigstartsignatures
\pysiglinewithargsret{\sphinxbfcode{\sphinxupquote{get\_height\_vec}}}{\sphinxparam{\DUrole{n}{lat}\DUrole{p}{:}\DUrole{w}{ }\DUrole{n}{ndarray}}\sphinxparamcomma \sphinxparam{\DUrole{n}{lon}\DUrole{p}{:}\DUrole{w}{ }\DUrole{n}{ndarray}}}{{ $\rightarrow$ ndarray}}
\pysigstopsignatures
\sphinxAtStartPar
Returns geoid height for given latitude and longitude positions.
\begin{quote}\begin{description}
\sphinxlineitem{Parameters}\begin{itemize}
\item {} 
\sphinxAtStartPar
\sphinxstyleliteralstrong{\sphinxupquote{lat}} (\sphinxstyleliteralemphasis{\sphinxupquote{float}}) \textendash{} The latitude coordinates in decimal degrees.

\item {} 
\sphinxAtStartPar
\sphinxstyleliteralstrong{\sphinxupquote{lon}} (\sphinxstyleliteralemphasis{\sphinxupquote{float}}) \textendash{} The longitude coordinates in decimal degrees.

\end{itemize}

\sphinxlineitem{Returns}
\sphinxAtStartPar
The geoid height for the corresponding locations (in metres).

\sphinxlineitem{Return type}
\sphinxAtStartPar
float

\end{description}\end{quote}

\end{fulllineitems}

\index{get\_dov() (qnav.earth.geoid.GeoidDummy method)@\spxentry{get\_dov()}\spxextra{qnav.earth.geoid.GeoidDummy method}}

\begin{fulllineitems}
\phantomsection\label{\detokenize{modules/earth/geoid:qnav.earth.geoid.GeoidDummy.get_dov}}
\pysigstartsignatures
\pysiglinewithargsret{\sphinxbfcode{\sphinxupquote{get\_dov}}}{\sphinxparam{\DUrole{n}{lat}\DUrole{p}{:}\DUrole{w}{ }\DUrole{n}{float}}\sphinxparamcomma \sphinxparam{\DUrole{n}{lon}\DUrole{p}{:}\DUrole{w}{ }\DUrole{n}{float}}}{{ $\rightarrow$ tuple\DUrole{p}{{[}}float\DUrole{p}{,}\DUrole{w}{ }float\DUrole{p}{{]}}}}
\pysigstopsignatures
\sphinxAtStartPar
Returns the Deflection of the Vertical (DoV) for requested position.
Calculates and returns the DoV gradients for the north and east
components (in metres) for requested location.
\begin{quote}\begin{description}
\sphinxlineitem{Parameters}\begin{itemize}
\item {} 
\sphinxAtStartPar
\sphinxstyleliteralstrong{\sphinxupquote{lat}} (\sphinxstyleliteralemphasis{\sphinxupquote{float}}) \textendash{} A latitude coordinate in decimal degrees.

\item {} 
\sphinxAtStartPar
\sphinxstyleliteralstrong{\sphinxupquote{lon}} (\sphinxstyleliteralemphasis{\sphinxupquote{float}}) \textendash{} A longitude coordinate in decimal degrees.

\end{itemize}

\sphinxlineitem{Returns}
\sphinxAtStartPar
Tuple containing the north and east components respectively.

\sphinxlineitem{Return type}
\sphinxAtStartPar
tuple{[}float, float{]}

\end{description}\end{quote}

\end{fulllineitems}

\index{get\_dov\_vec() (qnav.earth.geoid.GeoidDummy method)@\spxentry{get\_dov\_vec()}\spxextra{qnav.earth.geoid.GeoidDummy method}}

\begin{fulllineitems}
\phantomsection\label{\detokenize{modules/earth/geoid:qnav.earth.geoid.GeoidDummy.get_dov_vec}}
\pysigstartsignatures
\pysiglinewithargsret{\sphinxbfcode{\sphinxupquote{get\_dov\_vec}}}{\sphinxparam{\DUrole{n}{lat}\DUrole{p}{:}\DUrole{w}{ }\DUrole{n}{ndarray}}\sphinxparamcomma \sphinxparam{\DUrole{n}{lon}\DUrole{p}{:}\DUrole{w}{ }\DUrole{n}{ndarray}}}{{ $\rightarrow$ tuple\DUrole{p}{{[}}ndarray\DUrole{p}{,}\DUrole{w}{ }ndarray\DUrole{p}{{]}}}}
\pysigstopsignatures
\sphinxAtStartPar
Returns the Deflection of the Vertical (DoV) for requested positions.
Calculates and returns the DoV gradients for the north and east
components (in metres) for requested location.
\begin{quote}\begin{description}
\sphinxlineitem{Parameters}\begin{itemize}
\item {} 
\sphinxAtStartPar
\sphinxstyleliteralstrong{\sphinxupquote{lat}} (\sphinxstyleliteralemphasis{\sphinxupquote{np.ndarray}}) \textendash{} The latitude coordinates in decimal degrees.

\item {} 
\sphinxAtStartPar
\sphinxstyleliteralstrong{\sphinxupquote{lon}} (\sphinxstyleliteralemphasis{\sphinxupquote{np.ndarray}}) \textendash{} The longitude coordinates in decimal degrees.

\end{itemize}

\sphinxlineitem{Returns}
\sphinxAtStartPar
Tuple containing the north and east components respectively.

\sphinxlineitem{Return type}
\sphinxAtStartPar
tuple{[}float, float{]}

\end{description}\end{quote}

\end{fulllineitems}


\end{fulllineitems}

\index{read\_pgm\_data() (in module qnav.earth.geoid)@\spxentry{read\_pgm\_data()}\spxextra{in module qnav.earth.geoid}}

\begin{fulllineitems}
\phantomsection\label{\detokenize{modules/earth/geoid:qnav.earth.geoid.read_pgm_data}}
\pysigstartsignatures
\pysiglinewithargsret{\sphinxcode{\sphinxupquote{qnav.earth.geoid.}}\sphinxbfcode{\sphinxupquote{read\_pgm\_data}}}{\sphinxparam{\DUrole{n}{pgm\_file}\DUrole{p}{:}\DUrole{w}{ }\DUrole{n}{Path}}\sphinxparamcomma \sphinxparam{\DUrole{n}{pgm\_offset}\DUrole{p}{:}\DUrole{w}{ }\DUrole{n}{float}\DUrole{w}{ }\DUrole{o}{=}\DUrole{w}{ }\DUrole{default_value}{\sphinxhyphen{}108}}\sphinxparamcomma \sphinxparam{\DUrole{n}{pgm\_scale}\DUrole{p}{:}\DUrole{w}{ }\DUrole{n}{float}\DUrole{w}{ }\DUrole{o}{=}\DUrole{w}{ }\DUrole{default_value}{0.003}}}{{ $\rightarrow$ ndarray}}
\pysigstopsignatures
\sphinxAtStartPar
Reads and returns the numerical contents of given PGM file.
For a given PGM file, this function will read it and extract the Geoid
height information (relative to the Earth’s Ellipsoid) from it. The
returned matrix’s origin point’s coordinates will be {[}90, \sphinxhyphen{}180{]} with the
last point’s coordinates being {[}\sphinxhyphen{}90, 180{]}.
\begin{quote}\begin{description}
\sphinxlineitem{Parameters}\begin{itemize}
\item {} 
\sphinxAtStartPar
\sphinxstyleliteralstrong{\sphinxupquote{pgm\_file}} (\sphinxstyleliteralemphasis{\sphinxupquote{str}}\sphinxstyleliteralemphasis{\sphinxupquote{ or }}\sphinxstyleliteralemphasis{\sphinxupquote{Path}}) \textendash{} The location of the PGM file containing Geoid data.

\item {} 
\sphinxAtStartPar
\sphinxstyleliteralstrong{\sphinxupquote{pgm\_offset}} (\sphinxstyleliteralemphasis{\sphinxupquote{float}}) \textendash{} (Optional) The offset values for the values within the
PGM file as PGM values normally start from 0 (default = \sphinxhyphen{}108).

\item {} 
\sphinxAtStartPar
\sphinxstyleliteralstrong{\sphinxupquote{pgm\_scale}} (\sphinxstyleliteralemphasis{\sphinxupquote{float}}) \textendash{} (Optional) The scale for the values within the PGM file
as PGM values can only be integers (default = 0.003).

\end{itemize}

\sphinxlineitem{Returns}
\sphinxAtStartPar
A matrix containing the Geoid height (relative to the Earth’s
Ellipsoid) for various locations along the earth.

\sphinxlineitem{Return type}
\sphinxAtStartPar
np.ndarray

\end{description}\end{quote}

\end{fulllineitems}

\index{read\_pgm\_header() (in module qnav.earth.geoid)@\spxentry{read\_pgm\_header()}\spxextra{in module qnav.earth.geoid}}

\begin{fulllineitems}
\phantomsection\label{\detokenize{modules/earth/geoid:qnav.earth.geoid.read_pgm_header}}
\pysigstartsignatures
\pysiglinewithargsret{\sphinxcode{\sphinxupquote{qnav.earth.geoid.}}\sphinxbfcode{\sphinxupquote{read\_pgm\_header}}}{\sphinxparam{\DUrole{n}{pgm\_file}\DUrole{p}{:}\DUrole{w}{ }\DUrole{n}{Path}}}{{ $\rightarrow$ dict}}
\pysigstopsignatures
\sphinxAtStartPar
Returns values present in the header of a given PGM file.
PGM files can contain comments which describe information such as the
offset and scale values for the Geoid data they contain. For scenarios
where the offset and scale are contained within these comments, this
function will find and extract their values.
\begin{quote}\begin{description}
\sphinxlineitem{Parameters}
\sphinxAtStartPar
\sphinxstyleliteralstrong{\sphinxupquote{pgm\_file}} (\sphinxstyleliteralemphasis{\sphinxupquote{Path}}) \textendash{} The location of the PGM file containing Geoid data.

\sphinxlineitem{Returns}
\sphinxAtStartPar
A dictionary containing the found offset and scale values
(or None/Null if the values could not be found within the comments).

\sphinxlineitem{Return type}
\sphinxAtStartPar
dict

\end{description}\end{quote}

\end{fulllineitems}


\sphinxstepscope


\subsubsection{qnav.earth.utm}
\label{\detokenize{modules/earth/utm:module-qnav.earth.utm}}\label{\detokenize{modules/earth/utm:qnav-earth-utm}}\label{\detokenize{modules/earth/utm::doc}}\index{module@\spxentry{module}!qnav.earth.utm@\spxentry{qnav.earth.utm}}\index{qnav.earth.utm@\spxentry{qnav.earth.utm}!module@\spxentry{module}}\index{utm2lla() (in module qnav.earth.utm)@\spxentry{utm2lla()}\spxextra{in module qnav.earth.utm}}

\begin{fulllineitems}
\phantomsection\label{\detokenize{modules/earth/utm:qnav.earth.utm.utm2lla}}
\pysigstartsignatures
\pysiglinewithargsret{\sphinxcode{\sphinxupquote{qnav.earth.utm.}}\sphinxbfcode{\sphinxupquote{utm2lla}}}{\sphinxparam{\DUrole{n}{x}\DUrole{p}{:}\DUrole{w}{ }\DUrole{n}{float}}\sphinxparamcomma \sphinxparam{\DUrole{n}{y}\DUrole{p}{:}\DUrole{w}{ }\DUrole{n}{float}}\sphinxparamcomma \sphinxparam{\DUrole{n}{zone}\DUrole{p}{:}\DUrole{w}{ }\DUrole{n}{int}}}{{ $\rightarrow$ tuple\DUrole{p}{{[}}float\DUrole{p}{,}\DUrole{w}{ }float\DUrole{p}{{]}}}}
\pysigstopsignatures
\sphinxAtStartPar
This function coverts from Universal Transverse Mercator (UTM) to
Latitude\sphinxhyphen{}Longitude (LLA) coordinates. This function takes the X (eastern)
and Y (northern) offset from a given UTM integer zone (1 to 60) and
return a tuple containing the corresponding latitude\sphinxhyphen{}longitude
coordinates, respectively.
\begin{quote}\begin{description}
\sphinxlineitem{Parameters}\begin{itemize}
\item {} 
\sphinxAtStartPar
\sphinxstyleliteralstrong{\sphinxupquote{x}} (\sphinxstyleliteralemphasis{\sphinxupquote{float}}) \textendash{} The X/Eastern offset from the given UTM zone (in metres).

\item {} 
\sphinxAtStartPar
\sphinxstyleliteralstrong{\sphinxupquote{y}} (\sphinxstyleliteralemphasis{\sphinxupquote{float}}) \textendash{} The Y/Northern offset from the given UTM zone (in metres).

\item {} 
\sphinxAtStartPar
\sphinxstyleliteralstrong{\sphinxupquote{zone}} (\sphinxstyleliteralemphasis{\sphinxupquote{int}}) \textendash{} The integer UTM zone (in range 1 t 60).

\end{itemize}

\sphinxlineitem{Returns}
\sphinxAtStartPar
The given UTM coordinate converted to Latitude and Longitude
coordinates (contained within a tuple).

\sphinxlineitem{Return type}
\sphinxAtStartPar
tuple

\end{description}\end{quote}

\end{fulllineitems}



\subsection{Estimation}
\label{\detokenize{usage/packages:estimation}}
\sphinxAtStartPar
Modules related to state estimation:

\sphinxstepscope


\subsubsection{qnav.estimation.state}
\label{\detokenize{modules/estimation/state:module-qnav.estimation.state}}\label{\detokenize{modules/estimation/state:qnav-estimation-state}}\label{\detokenize{modules/estimation/state::doc}}\index{module@\spxentry{module}!qnav.estimation.state@\spxentry{qnav.estimation.state}}\index{qnav.estimation.state@\spxentry{qnav.estimation.state}!module@\spxentry{module}}

\paragraph{state.py}
\label{\detokenize{modules/estimation/state:state-py}}\begin{quote}\begin{description}
\sphinxlineitem{summary}
\sphinxAtStartPar
Supported base estimation state classes used to hold estimated data.
This module contains classes that handle the estimated state of the
platform during simulations. These are shared between the employed
fusion methods. Extensions of the base class are required for the
handling of addition data (such as error states and kalman matrices).

\sphinxlineitem{authors}
\begin{DUlineblock}{0em}
\item[] Prof. Jason Ralph \sphinxhyphen{} \sphinxhref{mailto:jfralph@liverpool.ac.uk}{jfralph@liverpool.ac.uk}
\item[] Michael Wright \sphinxhyphen{} \sphinxhref{mailto:mjwright@liverpool.ac.uk}{mjwright@liverpool.ac.uk}
\end{DUlineblock}

\sphinxlineitem{version}
\sphinxAtStartPar
1.0.0 \sphinxhyphen{} First full implementation of module.

\end{description}\end{quote}
\index{EstimatedState (class in qnav.estimation.state)@\spxentry{EstimatedState}\spxextra{class in qnav.estimation.state}}

\begin{fulllineitems}
\phantomsection\label{\detokenize{modules/estimation/state:qnav.estimation.state.EstimatedState}}
\pysigstartsignatures
\pysiglinewithargsret{\sphinxbfcode{\sphinxupquote{class\DUrole{w}{ }}}\sphinxcode{\sphinxupquote{qnav.estimation.state.}}\sphinxbfcode{\sphinxupquote{EstimatedState}}}{\sphinxparam{\DUrole{n}{ground\_truth: \textasciitilde{}qnav.waypoints.trajectory.GroundTruth}}\sphinxparamcomma \sphinxparam{\DUrole{n}{gravity\_model: \textasciitilde{}qnav.gravity.base.GravityModel = \textless{}qnav.gravity.simple.SimpleUniform object\textgreater{}}}}{}
\pysigstopsignatures
\sphinxAtStartPar
The base estimated state holding only general state information.
Simply maintains the estimated time, position, velocity, acceleration,
attitude and angle rates of the platform. Additionally, this also holds
a shared gravity model for calculating the expected gravity dynamics.
\index{timestamp (qnav.estimation.state.EstimatedState property)@\spxentry{timestamp}\spxextra{qnav.estimation.state.EstimatedState property}}

\begin{fulllineitems}
\phantomsection\label{\detokenize{modules/estimation/state:qnav.estimation.state.EstimatedState.timestamp}}
\pysigstartsignatures
\pysigline{\sphinxbfcode{\sphinxupquote{property\DUrole{w}{ }}}\sphinxbfcode{\sphinxupquote{timestamp}}\sphinxbfcode{\sphinxupquote{\DUrole{p}{:}\DUrole{w}{ }float}}}
\pysigstopsignatures
\sphinxAtStartPar
Returns the current estimated duration of the simulation (in seconds).
\begin{quote}\begin{description}
\sphinxlineitem{Returns}
\sphinxAtStartPar
The estimated time (in seconds).

\sphinxlineitem{Return type}
\sphinxAtStartPar
float

\end{description}\end{quote}

\end{fulllineitems}

\index{position (qnav.estimation.state.EstimatedState property)@\spxentry{position}\spxextra{qnav.estimation.state.EstimatedState property}}

\begin{fulllineitems}
\phantomsection\label{\detokenize{modules/estimation/state:qnav.estimation.state.EstimatedState.position}}
\pysigstartsignatures
\pysigline{\sphinxbfcode{\sphinxupquote{property\DUrole{w}{ }}}\sphinxbfcode{\sphinxupquote{position}}\sphinxbfcode{\sphinxupquote{\DUrole{p}{:}\DUrole{w}{ }ndarray}}}
\pysigstopsignatures
\sphinxAtStartPar
Returns the estimated position (in degrees\sphinxhyphen{}degrees\sphinxhyphen{}metres).
This is the estimated position represented by an array of latitude,
longitude and altitude values respectively.
\begin{quote}\begin{description}
\sphinxlineitem{Returns}
\sphinxAtStartPar
The estimated position (in degrees\sphinxhyphen{}degrees\sphinxhyphen{}metres).

\sphinxlineitem{Return type}
\sphinxAtStartPar
np.ndarray

\end{description}\end{quote}

\end{fulllineitems}

\index{velocity (qnav.estimation.state.EstimatedState property)@\spxentry{velocity}\spxextra{qnav.estimation.state.EstimatedState property}}

\begin{fulllineitems}
\phantomsection\label{\detokenize{modules/estimation/state:qnav.estimation.state.EstimatedState.velocity}}
\pysigstartsignatures
\pysigline{\sphinxbfcode{\sphinxupquote{property\DUrole{w}{ }}}\sphinxbfcode{\sphinxupquote{velocity}}\sphinxbfcode{\sphinxupquote{\DUrole{p}{:}\DUrole{w}{ }ndarray}}}
\pysigstopsignatures
\sphinxAtStartPar
Returns the estimated velocity (in meters/second).
This is the estimated velocity in body axis represented as an array
of along, across and down velocities respectively.
\begin{quote}\begin{description}
\sphinxlineitem{Returns}
\sphinxAtStartPar
The estimated velocity (in meters/second).

\sphinxlineitem{Return type}
\sphinxAtStartPar
np.ndarray

\end{description}\end{quote}

\end{fulllineitems}

\index{acceleration (qnav.estimation.state.EstimatedState property)@\spxentry{acceleration}\spxextra{qnav.estimation.state.EstimatedState property}}

\begin{fulllineitems}
\phantomsection\label{\detokenize{modules/estimation/state:qnav.estimation.state.EstimatedState.acceleration}}
\pysigstartsignatures
\pysigline{\sphinxbfcode{\sphinxupquote{property\DUrole{w}{ }}}\sphinxbfcode{\sphinxupquote{acceleration}}\sphinxbfcode{\sphinxupquote{\DUrole{p}{:}\DUrole{w}{ }ndarray}}}
\pysigstopsignatures
\sphinxAtStartPar
Returns the estimated acceleration (in meters/second\textasciicircum{}2).
This is the estimated acceleration in body axis represented as an
array of along, across and down accelerations respectively.
\begin{quote}\begin{description}
\sphinxlineitem{Returns}
\sphinxAtStartPar
The estimated acceleration (in meters/second\textasciicircum{}2).

\sphinxlineitem{Return type}
\sphinxAtStartPar
np.ndarray

\end{description}\end{quote}

\end{fulllineitems}

\index{attitude (qnav.estimation.state.EstimatedState property)@\spxentry{attitude}\spxextra{qnav.estimation.state.EstimatedState property}}

\begin{fulllineitems}
\phantomsection\label{\detokenize{modules/estimation/state:qnav.estimation.state.EstimatedState.attitude}}
\pysigstartsignatures
\pysigline{\sphinxbfcode{\sphinxupquote{property\DUrole{w}{ }}}\sphinxbfcode{\sphinxupquote{attitude}}\sphinxbfcode{\sphinxupquote{\DUrole{p}{:}\DUrole{w}{ }ndarray}}}
\pysigstopsignatures
\sphinxAtStartPar
Returns the estimated attitude (in degrees).
This is the estimated attitude in NED axis represented as an array
of heading, pitch and yaw angles respectively.
\begin{quote}\begin{description}
\sphinxlineitem{Returns}
\sphinxAtStartPar
The estimated attitude (in degrees).

\sphinxlineitem{Return type}
\sphinxAtStartPar
np.ndarray

\end{description}\end{quote}

\end{fulllineitems}

\index{angle\_rates (qnav.estimation.state.EstimatedState property)@\spxentry{angle\_rates}\spxextra{qnav.estimation.state.EstimatedState property}}

\begin{fulllineitems}
\phantomsection\label{\detokenize{modules/estimation/state:qnav.estimation.state.EstimatedState.angle_rates}}
\pysigstartsignatures
\pysigline{\sphinxbfcode{\sphinxupquote{property\DUrole{w}{ }}}\sphinxbfcode{\sphinxupquote{angle\_rates}}\sphinxbfcode{\sphinxupquote{\DUrole{p}{:}\DUrole{w}{ }ndarray}}}
\pysigstopsignatures
\sphinxAtStartPar
Returns the estimated angle rates (in degrees/second).
This is the estimated angle rates in body axis represented as an
array of P, Q and R euler angles respectively.
\begin{quote}\begin{description}
\sphinxlineitem{Returns}
\sphinxAtStartPar
The estimated angle rates (in degrees/second).

\sphinxlineitem{Return type}
\sphinxAtStartPar
np.ndarray

\end{description}\end{quote}

\end{fulllineitems}

\index{gravity\_model (qnav.estimation.state.EstimatedState property)@\spxentry{gravity\_model}\spxextra{qnav.estimation.state.EstimatedState property}}

\begin{fulllineitems}
\phantomsection\label{\detokenize{modules/estimation/state:qnav.estimation.state.EstimatedState.gravity_model}}
\pysigstartsignatures
\pysigline{\sphinxbfcode{\sphinxupquote{property\DUrole{w}{ }}}\sphinxbfcode{\sphinxupquote{gravity\_model}}\sphinxbfcode{\sphinxupquote{\DUrole{p}{:}\DUrole{w}{ }{\hyperref[\detokenize{modules/gravity/base:qnav.gravity.base.GravityModel}]{\sphinxcrossref{GravityModel}}}}}}
\pysigstopsignatures
\sphinxAtStartPar
Returns the shared gravity model.
This is the gravity model instance used for estimating the
gravitational effect during estimation and fusion. For consistency,
this is held in the estimated state so it is shared.
\begin{quote}\begin{description}
\sphinxlineitem{Returns}
\sphinxAtStartPar
The estimation gravity model.

\sphinxlineitem{Return type}
\sphinxAtStartPar
{\hyperref[\detokenize{modules/gravity/base:qnav.gravity.base.GravityModel}]{\sphinxcrossref{GravityModel}}}

\end{description}\end{quote}

\end{fulllineitems}

\index{update\_estimates() (qnav.estimation.state.EstimatedState method)@\spxentry{update\_estimates()}\spxextra{qnav.estimation.state.EstimatedState method}}

\begin{fulllineitems}
\phantomsection\label{\detokenize{modules/estimation/state:qnav.estimation.state.EstimatedState.update_estimates}}
\pysigstartsignatures
\pysiglinewithargsret{\sphinxbfcode{\sphinxupquote{update\_estimates}}}{\sphinxparam{\DUrole{n}{timestamp}\DUrole{p}{:}\DUrole{w}{ }\DUrole{n}{float}\DUrole{w}{ }\DUrole{o}{=}\DUrole{w}{ }\DUrole{default_value}{None}}\sphinxparamcomma \sphinxparam{\DUrole{n}{position}\DUrole{p}{:}\DUrole{w}{ }\DUrole{n}{ndarray}\DUrole{w}{ }\DUrole{o}{=}\DUrole{w}{ }\DUrole{default_value}{None}}\sphinxparamcomma \sphinxparam{\DUrole{n}{velocity}\DUrole{p}{:}\DUrole{w}{ }\DUrole{n}{ndarray}\DUrole{w}{ }\DUrole{o}{=}\DUrole{w}{ }\DUrole{default_value}{None}}\sphinxparamcomma \sphinxparam{\DUrole{n}{acceleration}\DUrole{p}{:}\DUrole{w}{ }\DUrole{n}{ndarray}\DUrole{w}{ }\DUrole{o}{=}\DUrole{w}{ }\DUrole{default_value}{None}}\sphinxparamcomma \sphinxparam{\DUrole{n}{attitude}\DUrole{p}{:}\DUrole{w}{ }\DUrole{n}{ndarray}\DUrole{w}{ }\DUrole{o}{=}\DUrole{w}{ }\DUrole{default_value}{None}}\sphinxparamcomma \sphinxparam{\DUrole{n}{angle\_rates}\DUrole{p}{:}\DUrole{w}{ }\DUrole{n}{ndarray}\DUrole{w}{ }\DUrole{o}{=}\DUrole{w}{ }\DUrole{default_value}{None}}}{}
\pysigstopsignatures
\sphinxAtStartPar
Overwrites selected current estimated states.
This is namely used when estimation fusion is applied and the
internal states must be updated. To prevent accidental overwriting of
estimated states or the setting of invalid data, all changes must be
submitted through this method.
\begin{quote}\begin{description}
\sphinxlineitem{Parameters}\begin{itemize}
\item {} 
\sphinxAtStartPar
\sphinxstyleliteralstrong{\sphinxupquote{timestamp}} (\sphinxstyleliteralemphasis{\sphinxupquote{float}}) \textendash{} The updated timestamp marking the estimated current
duration of the simulation (in seconds).

\item {} 
\sphinxAtStartPar
\sphinxstyleliteralstrong{\sphinxupquote{position}} (\sphinxstyleliteralemphasis{\sphinxupquote{numpy.ndarray}}\sphinxstyleliteralemphasis{\sphinxupquote{ (}}\sphinxstyleliteralemphasis{\sphinxupquote{3\sphinxhyphen{}elements}}\sphinxstyleliteralemphasis{\sphinxupquote{)}}) \textendash{} The updated position vector containing the
latitude (in degrees), longitude (in degrees) and altitude
(in metres) respectively.

\item {} 
\sphinxAtStartPar
\sphinxstyleliteralstrong{\sphinxupquote{velocity}} (\sphinxstyleliteralemphasis{\sphinxupquote{numpy.ndarray}}\sphinxstyleliteralemphasis{\sphinxupquote{ (}}\sphinxstyleliteralemphasis{\sphinxupquote{3\sphinxhyphen{}elements}}\sphinxstyleliteralemphasis{\sphinxupquote{)}}) \textendash{} The updated velocity vector for the
North, East and Down components respectively (in m/s).

\item {} 
\sphinxAtStartPar
\sphinxstyleliteralstrong{\sphinxupquote{acceleration}} (\sphinxstyleliteralemphasis{\sphinxupquote{numpy.ndarray}}\sphinxstyleliteralemphasis{\sphinxupquote{ (}}\sphinxstyleliteralemphasis{\sphinxupquote{3\sphinxhyphen{}elements}}\sphinxstyleliteralemphasis{\sphinxupquote{)}}) \textendash{} The updated acceleration vector for
the North, East and Down components respectively (in m/s\textasciicircum{}2).

\item {} 
\sphinxAtStartPar
\sphinxstyleliteralstrong{\sphinxupquote{attitude}} (\sphinxstyleliteralemphasis{\sphinxupquote{numpy.ndarray}}\sphinxstyleliteralemphasis{\sphinxupquote{ (}}\sphinxstyleliteralemphasis{\sphinxupquote{3\sphinxhyphen{}elements}}\sphinxstyleliteralemphasis{\sphinxupquote{)}}) \textendash{} The updated attitude vector for the Heading,
Pitch and Roll components respectively (in degrees).

\item {} 
\sphinxAtStartPar
\sphinxstyleliteralstrong{\sphinxupquote{angle\_rates}} (\sphinxstyleliteralemphasis{\sphinxupquote{numpy.ndarray}}\sphinxstyleliteralemphasis{\sphinxupquote{ (}}\sphinxstyleliteralemphasis{\sphinxupquote{3\sphinxhyphen{}elements}}\sphinxstyleliteralemphasis{\sphinxupquote{)}}) \textendash{} The updated angle rate vector for the
P, Q and R body rates respectively (in deg/s).

\end{itemize}

\end{description}\end{quote}

\end{fulllineitems}

\index{as\_dict() (qnav.estimation.state.EstimatedState method)@\spxentry{as\_dict()}\spxextra{qnav.estimation.state.EstimatedState method}}

\begin{fulllineitems}
\phantomsection\label{\detokenize{modules/estimation/state:qnav.estimation.state.EstimatedState.as_dict}}
\pysigstartsignatures
\pysiglinewithargsret{\sphinxbfcode{\sphinxupquote{as\_dict}}}{}{{ $\rightarrow$ dict\DUrole{p}{{[}}str\DUrole{p}{,}\DUrole{w}{ }float\DUrole{w}{ }\DUrole{p}{|}\DUrole{w}{ }ndarray\DUrole{p}{{]}}}}
\pysigstopsignatures
\sphinxAtStartPar
Represents the state and its content as a dictionary.
Namely used when a simple representation is needed.
\begin{quote}\begin{description}
\sphinxlineitem{Returns}
\sphinxAtStartPar
States in a python dictionary.

\sphinxlineitem{Return type}
\sphinxAtStartPar
dict{[}str, float | np.ndarray{]}

\end{description}\end{quote}

\end{fulllineitems}

\index{as\_numpy() (qnav.estimation.state.EstimatedState method)@\spxentry{as\_numpy()}\spxextra{qnav.estimation.state.EstimatedState method}}

\begin{fulllineitems}
\phantomsection\label{\detokenize{modules/estimation/state:qnav.estimation.state.EstimatedState.as_numpy}}
\pysigstartsignatures
\pysiglinewithargsret{\sphinxbfcode{\sphinxupquote{as\_numpy}}}{}{{ $\rightarrow$ ndarray}}
\pysigstopsignatures
\sphinxAtStartPar
Represents the state and its content as a numpy array.
Namely used when a simple representation is needed.
\begin{quote}\begin{description}
\sphinxlineitem{Returns}
\sphinxAtStartPar
States in concatenated numpy array.

\sphinxlineitem{Return type}
\sphinxAtStartPar
numpy.ndarray

\end{description}\end{quote}

\end{fulllineitems}

\index{clone() (qnav.estimation.state.EstimatedState method)@\spxentry{clone()}\spxextra{qnav.estimation.state.EstimatedState method}}

\begin{fulllineitems}
\phantomsection\label{\detokenize{modules/estimation/state:qnav.estimation.state.EstimatedState.clone}}
\pysigstartsignatures
\pysiglinewithargsret{\sphinxbfcode{\sphinxupquote{clone}}}{}{}
\pysigstopsignatures
\sphinxAtStartPar
Generates a cloned instance of itself.
\begin{quote}\begin{description}
\sphinxlineitem{Returns}
\sphinxAtStartPar
Deep clone of this instance.

\sphinxlineitem{Return type}
\sphinxAtStartPar
State

\end{description}\end{quote}

\end{fulllineitems}


\end{fulllineitems}

\index{KalmanEstimatedState (class in qnav.estimation.state)@\spxentry{KalmanEstimatedState}\spxextra{class in qnav.estimation.state}}

\begin{fulllineitems}
\phantomsection\label{\detokenize{modules/estimation/state:qnav.estimation.state.KalmanEstimatedState}}
\pysigstartsignatures
\pysiglinewithargsret{\sphinxbfcode{\sphinxupquote{class\DUrole{w}{ }}}\sphinxcode{\sphinxupquote{qnav.estimation.state.}}\sphinxbfcode{\sphinxupquote{KalmanEstimatedState}}}{\sphinxparam{ground\_truth: \textasciitilde{}qnav.waypoints.trajectory.GroundTruth, h: \textasciitilde{}numpy.ndarray, r: \textasciitilde{}numpy.ndarray, f: \textasciitilde{}numpy.ndarray, q: \textasciitilde{}numpy.ndarray, state\_errors: \textasciitilde{}numpy.ndarray = array({[}{[}1.e\sphinxhyphen{}09, 0.e+00, 0.e+00, 0.e+00, 0.e+00, 0.e+00, 0.e+00, 0.e+00,         0.e+00, 0.e+00, 0.e+00, 0.e+00, 0.e+00, 0.e+00, 0.e+00{]},        {[}0.e+00, 1.e\sphinxhyphen{}09, 0.e+00, 0.e+00, 0.e+00, 0.e+00, 0.e+00, 0.e+00,         0.e+00, 0.e+00, 0.e+00, 0.e+00, 0.e+00, 0.e+00, 0.e+00{]},        {[}0.e+00, 0.e+00, 1.e\sphinxhyphen{}09, 0.e+00, 0.e+00, 0.e+00, 0.e+00, 0.e+00,         0.e+00, 0.e+00, 0.e+00, 0.e+00, 0.e+00, 0.e+00, 0.e+00{]},        {[}0.e+00, 0.e+00, 0.e+00, 1.e\sphinxhyphen{}09, 0.e+00, 0.e+00, 0.e+00, 0.e+00,         0.e+00, 0.e+00, 0.e+00, 0.e+00, 0.e+00, 0.e+00, 0.e+00{]},        {[}0.e+00, 0.e+00, 0.e+00, 0.e+00, 1.e\sphinxhyphen{}09, 0.e+00, 0.e+00, 0.e+00,         0.e+00, 0.e+00, 0.e+00, 0.e+00, 0.e+00, 0.e+00, 0.e+00{]},        {[}0.e+00, 0.e+00, 0.e+00, 0.e+00, 0.e+00, 1.e\sphinxhyphen{}09, 0.e+00, 0.e+00,         0.e+00, 0.e+00, 0.e+00, 0.e+00, 0.e+00, 0.e+00, 0.e+00{]},        {[}0.e+00, 0.e+00, 0.e+00, 0.e+00, 0.e+00, 0.e+00, 1.e\sphinxhyphen{}09, 0.e+00,         0.e+00, 0.e+00, 0.e+00, 0.e+00, 0.e+00, 0.e+00, 0.e+00{]},        {[}0.e+00, 0.e+00, 0.e+00, 0.e+00, 0.e+00, 0.e+00, 0.e+00, 1.e\sphinxhyphen{}09,         0.e+00, 0.e+00, 0.e+00, 0.e+00, 0.e+00, 0.e+00, 0.e+00{]},        {[}0.e+00, 0.e+00, 0.e+00, 0.e+00, 0.e+00, 0.e+00, 0.e+00, 0.e+00,         1.e\sphinxhyphen{}09, 0.e+00, 0.e+00, 0.e+00, 0.e+00, 0.e+00, 0.e+00{]},        {[}0.e+00, 0.e+00, 0.e+00, 0.e+00, 0.e+00, 0.e+00, 0.e+00, 0.e+00,         0.e+00, 1.e\sphinxhyphen{}09, 0.e+00, 0.e+00, 0.e+00, 0.e+00, 0.e+00{]},        {[}0.e+00, 0.e+00, 0.e+00, 0.e+00, 0.e+00, 0.e+00, 0.e+00, 0.e+00,         0.e+00, 0.e+00, 1.e\sphinxhyphen{}09, 0.e+00, 0.e+00, 0.e+00, 0.e+00{]},        {[}0.e+00, 0.e+00, 0.e+00, 0.e+00, 0.e+00, 0.e+00, 0.e+00, 0.e+00,         0.e+00, 0.e+00, 0.e+00, 1.e\sphinxhyphen{}09, 0.e+00, 0.e+00, 0.e+00{]},        {[}0.e+00, 0.e+00, 0.e+00, 0.e+00, 0.e+00, 0.e+00, 0.e+00, 0.e+00,         0.e+00, 0.e+00, 0.e+00, 0.e+00, 1.e\sphinxhyphen{}09, 0.e+00, 0.e+00{]},        {[}0.e+00, 0.e+00, 0.e+00, 0.e+00, 0.e+00, 0.e+00, 0.e+00, 0.e+00,         0.e+00, 0.e+00, 0.e+00, 0.e+00, 0.e+00, 1.e\sphinxhyphen{}09, 0.e+00{]},        {[}0.e+00, 0.e+00, 0.e+00, 0.e+00, 0.e+00, 0.e+00, 0.e+00, 0.e+00,         0.e+00, 0.e+00, 0.e+00, 0.e+00, 0.e+00, 0.e+00, 1.e\sphinxhyphen{}09{]}{]}), gravity\_model: \textasciitilde{}qnav.gravity.base.GravityModel = \textless{}qnav.gravity.simple.SimpleUniform object\textgreater{}}}{}
\pysigstopsignatures
\sphinxAtStartPar
Extended estimation state featuring Kalman Filter properties.
This is an extension for the standard estimated state so that Kalman
states can also be represented. These additions include shared Kalman
Filter matrices and state errors.
\index{state\_vector (qnav.estimation.state.KalmanEstimatedState property)@\spxentry{state\_vector}\spxextra{qnav.estimation.state.KalmanEstimatedState property}}

\begin{fulllineitems}
\phantomsection\label{\detokenize{modules/estimation/state:qnav.estimation.state.KalmanEstimatedState.state_vector}}
\pysigstartsignatures
\pysigline{\sphinxbfcode{\sphinxupquote{property\DUrole{w}{ }}}\sphinxbfcode{\sphinxupquote{state\_vector}}}
\pysigstopsignatures
\sphinxAtStartPar
Returns estimated states in form of state vector.
A commonly used representation of estimated states in Kalman Filters.
\begin{quote}\begin{description}
\sphinxlineitem{Returns}
\sphinxAtStartPar
Kalman State Vector

\sphinxlineitem{Return type}
\sphinxAtStartPar
numpy.ndarray (15\sphinxhyphen{}elements)

\end{description}\end{quote}

\end{fulllineitems}

\index{update\_estimates() (qnav.estimation.state.KalmanEstimatedState method)@\spxentry{update\_estimates()}\spxextra{qnav.estimation.state.KalmanEstimatedState method}}

\begin{fulllineitems}
\phantomsection\label{\detokenize{modules/estimation/state:qnav.estimation.state.KalmanEstimatedState.update_estimates}}
\pysigstartsignatures
\pysiglinewithargsret{\sphinxbfcode{\sphinxupquote{update\_estimates}}}{\sphinxparam{\DUrole{n}{timestamp}\DUrole{p}{:}\DUrole{w}{ }\DUrole{n}{float}\DUrole{w}{ }\DUrole{o}{=}\DUrole{w}{ }\DUrole{default_value}{None}}\sphinxparamcomma \sphinxparam{\DUrole{n}{position}\DUrole{p}{:}\DUrole{w}{ }\DUrole{n}{ndarray}\DUrole{w}{ }\DUrole{o}{=}\DUrole{w}{ }\DUrole{default_value}{None}}\sphinxparamcomma \sphinxparam{\DUrole{n}{velocity}\DUrole{p}{:}\DUrole{w}{ }\DUrole{n}{ndarray}\DUrole{w}{ }\DUrole{o}{=}\DUrole{w}{ }\DUrole{default_value}{None}}\sphinxparamcomma \sphinxparam{\DUrole{n}{acceleration}\DUrole{p}{:}\DUrole{w}{ }\DUrole{n}{ndarray}\DUrole{w}{ }\DUrole{o}{=}\DUrole{w}{ }\DUrole{default_value}{None}}\sphinxparamcomma \sphinxparam{\DUrole{n}{attitude}\DUrole{p}{:}\DUrole{w}{ }\DUrole{n}{ndarray}\DUrole{w}{ }\DUrole{o}{=}\DUrole{w}{ }\DUrole{default_value}{None}}\sphinxparamcomma \sphinxparam{\DUrole{n}{angle\_rates}\DUrole{p}{:}\DUrole{w}{ }\DUrole{n}{ndarray}\DUrole{w}{ }\DUrole{o}{=}\DUrole{w}{ }\DUrole{default_value}{None}}\sphinxparamcomma \sphinxparam{\DUrole{n}{state\_errors}\DUrole{p}{:}\DUrole{w}{ }\DUrole{n}{ndarray}\DUrole{w}{ }\DUrole{o}{=}\DUrole{w}{ }\DUrole{default_value}{None}}}{}
\pysigstopsignatures
\sphinxAtStartPar
Overwrites selected current estimated states.
This is namely used when estimation fusion is applied and the
internal states must be updated. To prevent accidental overwriting of
estimated states or the setting of invalid data, all changes must be
submitted through this method.
\begin{quote}\begin{description}
\sphinxlineitem{Parameters}\begin{itemize}
\item {} 
\sphinxAtStartPar
\sphinxstyleliteralstrong{\sphinxupquote{timestamp}} (\sphinxstyleliteralemphasis{\sphinxupquote{float}}) \textendash{} The updated timestamp marking the estimated current
duration of the simulation (in seconds).

\item {} 
\sphinxAtStartPar
\sphinxstyleliteralstrong{\sphinxupquote{position}} (\sphinxstyleliteralemphasis{\sphinxupquote{numpy.ndarray}}\sphinxstyleliteralemphasis{\sphinxupquote{ (}}\sphinxstyleliteralemphasis{\sphinxupquote{3\sphinxhyphen{}elements}}\sphinxstyleliteralemphasis{\sphinxupquote{)}}) \textendash{} The updated position vector containing the
latitude (in degrees), longitude (in degrees) and altitude
(in metres) respectively.

\item {} 
\sphinxAtStartPar
\sphinxstyleliteralstrong{\sphinxupquote{velocity}} (\sphinxstyleliteralemphasis{\sphinxupquote{numpy.ndarray}}\sphinxstyleliteralemphasis{\sphinxupquote{ (}}\sphinxstyleliteralemphasis{\sphinxupquote{3\sphinxhyphen{}elements}}\sphinxstyleliteralemphasis{\sphinxupquote{)}}) \textendash{} The updated velocity vector for the
North, East and Down components respectively (in m/s).

\item {} 
\sphinxAtStartPar
\sphinxstyleliteralstrong{\sphinxupquote{acceleration}} (\sphinxstyleliteralemphasis{\sphinxupquote{numpy.ndarray}}\sphinxstyleliteralemphasis{\sphinxupquote{ (}}\sphinxstyleliteralemphasis{\sphinxupquote{3\sphinxhyphen{}elements}}\sphinxstyleliteralemphasis{\sphinxupquote{)}}) \textendash{} The updated acceleration vector for
the North, East and Down components respectively (in m/s\textasciicircum{}2).

\item {} 
\sphinxAtStartPar
\sphinxstyleliteralstrong{\sphinxupquote{attitude}} (\sphinxstyleliteralemphasis{\sphinxupquote{numpy.ndarray}}\sphinxstyleliteralemphasis{\sphinxupquote{ (}}\sphinxstyleliteralemphasis{\sphinxupquote{3\sphinxhyphen{}elements}}\sphinxstyleliteralemphasis{\sphinxupquote{)}}) \textendash{} The updated attitude vector for the Heading,
Pitch and Roll components respectively (in degrees).

\item {} 
\sphinxAtStartPar
\sphinxstyleliteralstrong{\sphinxupquote{angle\_rates}} (\sphinxstyleliteralemphasis{\sphinxupquote{numpy.ndarray}}\sphinxstyleliteralemphasis{\sphinxupquote{ (}}\sphinxstyleliteralemphasis{\sphinxupquote{3\sphinxhyphen{}elements}}\sphinxstyleliteralemphasis{\sphinxupquote{)}}) \textendash{} The updated angle rate vector for the
P, Q and R body rates respectively (in deg/s).

\item {} 
\sphinxAtStartPar
\sphinxstyleliteralstrong{\sphinxupquote{state\_errors}} (\sphinxstyleliteralemphasis{\sphinxupquote{numpy.ndarray}}\sphinxstyleliteralemphasis{\sphinxupquote{ (}}\sphinxstyleliteralemphasis{\sphinxupquote{15\sphinxhyphen{}by\sphinxhyphen{}15 elements}}\sphinxstyleliteralemphasis{\sphinxupquote{)}}) \textendash{} The update Kalman state errors.

\end{itemize}

\end{description}\end{quote}

\end{fulllineitems}


\end{fulllineitems}



\subsection{Fusion}
\label{\detokenize{usage/packages:fusion}}
\sphinxAtStartPar
Modules related to standard sensor fusion methods:

\sphinxstepscope


\subsubsection{qnav.fusion.ins}
\label{\detokenize{modules/fusion/ins:module-qnav.fusion.ins}}\label{\detokenize{modules/fusion/ins:qnav-fusion-ins}}\label{\detokenize{modules/fusion/ins::doc}}\index{module@\spxentry{module}!qnav.fusion.ins@\spxentry{qnav.fusion.ins}}\index{qnav.fusion.ins@\spxentry{qnav.fusion.ins}!module@\spxentry{module}}

\paragraph{ins.py}
\label{\detokenize{modules/fusion/ins:ins-py}}\begin{quote}\begin{description}
\sphinxlineitem{summary}
\sphinxAtStartPar
Fusion method for supplying Inertial Navigation.
Provides a selection of standard sensor fusion methods for means of basic
internal navigation. Each of these takes measurements from accelerometer
and gyroscope sensors (IMU).

\sphinxlineitem{authors}
\begin{DUlineblock}{0em}
\item[] Prof. Jason Ralph \sphinxhyphen{} \sphinxhref{mailto:jfralph@liverpool.ac.uk}{jfralph@liverpool.ac.uk}
\item[] Michael Wright \sphinxhyphen{} \sphinxhref{mailto:mjwright@liverpool.ac.uk}{mjwright@liverpool.ac.uk}
\end{DUlineblock}

\sphinxlineitem{version}
\sphinxAtStartPar
1.0.0 \sphinxhyphen{} First full implementation of module.

\end{description}\end{quote}
\index{NumericalINS (class in qnav.fusion.ins)@\spxentry{NumericalINS}\spxextra{class in qnav.fusion.ins}}

\begin{fulllineitems}
\phantomsection\label{\detokenize{modules/fusion/ins:qnav.fusion.ins.NumericalINS}}
\pysigstartsignatures
\pysiglinewithargsret{\sphinxbfcode{\sphinxupquote{class\DUrole{w}{ }}}\sphinxcode{\sphinxupquote{qnav.fusion.ins.}}\sphinxbfcode{\sphinxupquote{NumericalINS}}}{\sphinxparam{\DUrole{n}{accelerometer}\DUrole{p}{:}\DUrole{w}{ }\DUrole{n}{{\hyperref[\detokenize{modules/measurement/accelerometer:qnav.measurement.accelerometer.Accelerometer}]{\sphinxcrossref{Accelerometer}}}}}\sphinxparamcomma \sphinxparam{\DUrole{n}{gyroscope}\DUrole{p}{:}\DUrole{w}{ }\DUrole{n}{{\hyperref[\detokenize{modules/measurement/gyroscope:qnav.measurement.gyroscope.Gyroscope}]{\sphinxcrossref{Gyroscope}}}}}\sphinxparamcomma \sphinxparam{\DUrole{n}{estimated\_acc\_axis}\DUrole{p}{:}\DUrole{w}{ }\DUrole{n}{{\hyperref[\detokenize{modules/measurement/platform:qnav.measurement.platform.SensorAxis}]{\sphinxcrossref{SensorAxis}}}}\DUrole{w}{ }\DUrole{o}{=}\DUrole{w}{ }\DUrole{default_value}{None}}\sphinxparamcomma \sphinxparam{\DUrole{n}{estimated\_gyro\_axis}\DUrole{p}{:}\DUrole{w}{ }\DUrole{n}{{\hyperref[\detokenize{modules/measurement/platform:qnav.measurement.platform.SensorAxis}]{\sphinxcrossref{SensorAxis}}}}\DUrole{w}{ }\DUrole{o}{=}\DUrole{w}{ }\DUrole{default_value}{None}}}{}
\pysigstopsignatures
\sphinxAtStartPar
A fairly basic simple numerical solution for Internal Navigation (INS).
This INS solution is based on simple numerical integration. While its
estimation calculations are relatively primitive, it serves as a base
model that other solutions can extend from. The simple implementation
also provides straight\sphinxhyphen{}forward means for testing calculations.
\index{perform\_fusion() (qnav.fusion.ins.NumericalINS method)@\spxentry{perform\_fusion()}\spxextra{qnav.fusion.ins.NumericalINS method}}

\begin{fulllineitems}
\phantomsection\label{\detokenize{modules/fusion/ins:qnav.fusion.ins.NumericalINS.perform_fusion}}
\pysigstartsignatures
\pysiglinewithargsret{\sphinxbfcode{\sphinxupquote{perform\_fusion}}}{\sphinxparam{\DUrole{n}{estimated\_state}\DUrole{p}{:}\DUrole{w}{ }\DUrole{n}{{\hyperref[\detokenize{modules/estimation/state:qnav.estimation.state.EstimatedState}]{\sphinxcrossref{EstimatedState}}}}}}{{ $\rightarrow$ None}}
\pysigstopsignatures
\sphinxAtStartPar
Performs fusion using latest accelerometer and gyroscope measurements.
Using the latest basic inertial measurements, updates the full current
estimated state using a simple numerical model.
\begin{quote}\begin{description}
\sphinxlineitem{Parameters}
\sphinxAtStartPar
\sphinxstyleliteralstrong{\sphinxupquote{estimated\_state}} ({\hyperref[\detokenize{modules/estimation/state:qnav.estimation.state.EstimatedState}]{\sphinxcrossref{\sphinxstyleliteralemphasis{\sphinxupquote{EstimatedState}}}}}) \textendash{} The current estimated state.

\end{description}\end{quote}

\end{fulllineitems}


\end{fulllineitems}

\index{RungeKutta (class in qnav.fusion.ins)@\spxentry{RungeKutta}\spxextra{class in qnav.fusion.ins}}

\begin{fulllineitems}
\phantomsection\label{\detokenize{modules/fusion/ins:qnav.fusion.ins.RungeKutta}}
\pysigstartsignatures
\pysiglinewithargsret{\sphinxbfcode{\sphinxupquote{class\DUrole{w}{ }}}\sphinxcode{\sphinxupquote{qnav.fusion.ins.}}\sphinxbfcode{\sphinxupquote{RungeKutta}}}{\sphinxparam{\DUrole{n}{accelerometer}\DUrole{p}{:}\DUrole{w}{ }\DUrole{n}{{\hyperref[\detokenize{modules/measurement/accelerometer:qnav.measurement.accelerometer.Accelerometer}]{\sphinxcrossref{Accelerometer}}}}}\sphinxparamcomma \sphinxparam{\DUrole{n}{gyroscope}\DUrole{p}{:}\DUrole{w}{ }\DUrole{n}{{\hyperref[\detokenize{modules/measurement/gyroscope:qnav.measurement.gyroscope.Gyroscope}]{\sphinxcrossref{Gyroscope}}}}}\sphinxparamcomma \sphinxparam{\DUrole{n}{estimated\_acc\_axis}\DUrole{p}{:}\DUrole{w}{ }\DUrole{n}{{\hyperref[\detokenize{modules/measurement/platform:qnav.measurement.platform.SensorAxis}]{\sphinxcrossref{SensorAxis}}}}\DUrole{w}{ }\DUrole{o}{=}\DUrole{w}{ }\DUrole{default_value}{None}}\sphinxparamcomma \sphinxparam{\DUrole{n}{estimated\_gyro\_axis}\DUrole{p}{:}\DUrole{w}{ }\DUrole{n}{{\hyperref[\detokenize{modules/measurement/platform:qnav.measurement.platform.SensorAxis}]{\sphinxcrossref{SensorAxis}}}}\DUrole{w}{ }\DUrole{o}{=}\DUrole{w}{ }\DUrole{default_value}{None}}}{}
\pysigstopsignatures
\sphinxAtStartPar
An extension of the numerical INS that employs Runge\sphinxhyphen{}Kutta methods.
Processes the measurement vector and updates the estimation states using
Runge Kutta (4th order) integration.
\begin{description}
\sphinxlineitem{Does include:}\begin{itemize}
\item {} 
\sphinxAtStartPar
Simple Accelerometer Measurement Model (basic \sphinxhyphen{} independent
measurement errors, giving drift term)

\item {} 
\sphinxAtStartPar
3 x Individual Accelerator Bias values (fixed)

\item {} 
\sphinxAtStartPar
3 x Individual Accelerator Scaling errors (fixed)

\item {} 
\sphinxAtStartPar
3D Accelerometer non\sphinxhyphen{}orthogonality errors (fixed) for cross\sphinxhyphen{}coupling

\item {} 
\sphinxAtStartPar
Simple (not position dependent) Gravity Compensation

\item {} 
\sphinxAtStartPar
Simple Gyroscope Measurement Model (basic \sphinxhyphen{} independent
measurement errors, giving drift term)

\item {} 
\sphinxAtStartPar
3 x Individual Gyroscope Bias values (fixed)

\item {} 
\sphinxAtStartPar
3 x Individual Gyroscope Scaling errors (fixed)

\item {} 
\sphinxAtStartPar
3D Gyroscope non\sphinxhyphen{}orthogonality errors (fixed) for cross\sphinxhyphen{}coupling

\item {} 
\sphinxAtStartPar
Sensor bias drift for accelerometers and gyroscopes

\item {} 
\sphinxAtStartPar
Kalman Filtering of measurement signals

\item {} 
\sphinxAtStartPar
WGS’84 coordinates (ellipsoidal rotating Earth)

\item {} 
\sphinxAtStartPar
Effect of Coriolis effect due to Earth’s rotation

\end{itemize}

\sphinxlineitem{Does NOT currently include:}\begin{itemize}
\item {} 
\sphinxAtStartPar
Schuler correction loop.

\item {} 
\sphinxAtStartPar
Physics\sphinxhyphen{}based accelerometer sensor model

\item {} 
\sphinxAtStartPar
Physics\sphinxhyphen{}based gyroscope sensor model

\item {} 
\sphinxAtStartPar
Lever\sphinxhyphen{}arm effects from rotations not around origin/centre of the IMU

\end{itemize}

\end{description}
\index{perform\_fusion() (qnav.fusion.ins.RungeKutta method)@\spxentry{perform\_fusion()}\spxextra{qnav.fusion.ins.RungeKutta method}}

\begin{fulllineitems}
\phantomsection\label{\detokenize{modules/fusion/ins:qnav.fusion.ins.RungeKutta.perform_fusion}}
\pysigstartsignatures
\pysiglinewithargsret{\sphinxbfcode{\sphinxupquote{perform\_fusion}}}{\sphinxparam{\DUrole{n}{estimated\_state}\DUrole{p}{:}\DUrole{w}{ }\DUrole{n}{{\hyperref[\detokenize{modules/estimation/state:qnav.estimation.state.EstimatedState}]{\sphinxcrossref{EstimatedState}}}}}}{{ $\rightarrow$ None}}
\pysigstopsignatures
\sphinxAtStartPar
Performs fusion using latest accelerometer and gyroscope measurements.
Using the latest basic inertial measurements, updates the full current
estimated state using a simple numerical model.
\begin{quote}\begin{description}
\sphinxlineitem{Parameters}
\sphinxAtStartPar
\sphinxstyleliteralstrong{\sphinxupquote{estimated\_state}} ({\hyperref[\detokenize{modules/estimation/state:qnav.estimation.state.EstimatedState}]{\sphinxcrossref{\sphinxstyleliteralemphasis{\sphinxupquote{EstimatedState}}}}}) \textendash{} The current estimated state.

\end{description}\end{quote}

\end{fulllineitems}


\end{fulllineitems}

\index{AdamsBashforth (class in qnav.fusion.ins)@\spxentry{AdamsBashforth}\spxextra{class in qnav.fusion.ins}}

\begin{fulllineitems}
\phantomsection\label{\detokenize{modules/fusion/ins:qnav.fusion.ins.AdamsBashforth}}
\pysigstartsignatures
\pysiglinewithargsret{\sphinxbfcode{\sphinxupquote{class\DUrole{w}{ }}}\sphinxcode{\sphinxupquote{qnav.fusion.ins.}}\sphinxbfcode{\sphinxupquote{AdamsBashforth}}}{\sphinxparam{\DUrole{n}{accelerometer}\DUrole{p}{:}\DUrole{w}{ }\DUrole{n}{{\hyperref[\detokenize{modules/measurement/accelerometer:qnav.measurement.accelerometer.Accelerometer}]{\sphinxcrossref{Accelerometer}}}}}\sphinxparamcomma \sphinxparam{\DUrole{n}{gyroscope}\DUrole{p}{:}\DUrole{w}{ }\DUrole{n}{{\hyperref[\detokenize{modules/measurement/gyroscope:qnav.measurement.gyroscope.Gyroscope}]{\sphinxcrossref{Gyroscope}}}}}\sphinxparamcomma \sphinxparam{\DUrole{n}{estimated\_acc\_axis}\DUrole{p}{:}\DUrole{w}{ }\DUrole{n}{{\hyperref[\detokenize{modules/measurement/platform:qnav.measurement.platform.SensorAxis}]{\sphinxcrossref{SensorAxis}}}}\DUrole{w}{ }\DUrole{o}{=}\DUrole{w}{ }\DUrole{default_value}{None}}\sphinxparamcomma \sphinxparam{\DUrole{n}{estimated\_gyro\_axis}\DUrole{p}{:}\DUrole{w}{ }\DUrole{n}{{\hyperref[\detokenize{modules/measurement/platform:qnav.measurement.platform.SensorAxis}]{\sphinxcrossref{SensorAxis}}}}\DUrole{w}{ }\DUrole{o}{=}\DUrole{w}{ }\DUrole{default_value}{None}}}{}
\pysigstopsignatures
\sphinxAtStartPar
An extension of the numerical INS that employs Adams\sphinxhyphen{}bashforth method.
Processes the measurement vector and updates the estimation states using
Adams\sphinxhyphen{}Bashforth integration.
\begin{description}
\sphinxlineitem{Does include:}\begin{itemize}
\item {} 
\sphinxAtStartPar
Simple Accelerometer Measurement Model (basic \sphinxhyphen{} independent
measurement errors, giving drift term)

\item {} 
\sphinxAtStartPar
3 x Individual Accelerator Bias values (fixed)

\item {} 
\sphinxAtStartPar
3 x Individual Accelerator Scaling errors (fixed)

\item {} 
\sphinxAtStartPar
3D Accelerometer non\sphinxhyphen{}orthogonality errors (fixed) for cross\sphinxhyphen{}coupling

\item {} 
\sphinxAtStartPar
Simple (not position dependent) Gravity Compensation

\item {} 
\sphinxAtStartPar
Simple Gyroscope Measurement Model (basic \sphinxhyphen{} independent
measurement errors, giving drift term)

\item {} 
\sphinxAtStartPar
3 x Individual Gyroscope Bias values (fixed)

\item {} 
\sphinxAtStartPar
3 x Individual Gyroscope Scaling errors (fixed)

\item {} 
\sphinxAtStartPar
3D Gyroscope non\sphinxhyphen{}orthogonality errors (fixed) for cross\sphinxhyphen{}coupling

\item {} 
\sphinxAtStartPar
Sensor bias drift for accelerometers and gyroscopes

\item {} 
\sphinxAtStartPar
Kalman Filtering of measurement signals

\item {} 
\sphinxAtStartPar
WGS’84 coordinates (ellipsoidal rotating Earth)

\item {} 
\sphinxAtStartPar
Effect of Coriolis effect due to Earth’s rotation

\end{itemize}

\sphinxlineitem{Does NOT currently include:}\begin{itemize}
\item {} 
\sphinxAtStartPar
Schuler correction loop.

\item {} 
\sphinxAtStartPar
Physics\sphinxhyphen{}based accelerometer sensor model

\item {} 
\sphinxAtStartPar
Physics\sphinxhyphen{}based gyroscope sensor model

\item {} 
\sphinxAtStartPar
Lever\sphinxhyphen{}arm effects from rotations not around origin/centre of the IMU

\end{itemize}

\end{description}
\index{perform\_fusion() (qnav.fusion.ins.AdamsBashforth method)@\spxentry{perform\_fusion()}\spxextra{qnav.fusion.ins.AdamsBashforth method}}

\begin{fulllineitems}
\phantomsection\label{\detokenize{modules/fusion/ins:qnav.fusion.ins.AdamsBashforth.perform_fusion}}
\pysigstartsignatures
\pysiglinewithargsret{\sphinxbfcode{\sphinxupquote{perform\_fusion}}}{\sphinxparam{\DUrole{n}{estimated\_state}\DUrole{p}{:}\DUrole{w}{ }\DUrole{n}{{\hyperref[\detokenize{modules/estimation/state:qnav.estimation.state.EstimatedState}]{\sphinxcrossref{EstimatedState}}}}}}{{ $\rightarrow$ None}}
\pysigstopsignatures
\sphinxAtStartPar
Performs fusion using latest accelerometer and gyroscope measurements.
Using the latest basic inertial measurements, updates the full current
estimated state using a simple numerical model.
\begin{quote}\begin{description}
\sphinxlineitem{Parameters}
\sphinxAtStartPar
\sphinxstyleliteralstrong{\sphinxupquote{estimated\_state}} ({\hyperref[\detokenize{modules/estimation/state:qnav.estimation.state.EstimatedState}]{\sphinxcrossref{\sphinxstyleliteralemphasis{\sphinxupquote{EstimatedState}}}}}) \textendash{} The current estimated state.

\end{description}\end{quote}

\end{fulllineitems}


\end{fulllineitems}


\sphinxstepscope


\subsubsection{qnav.fusion.kalman}
\label{\detokenize{modules/fusion/kalman:module-qnav.fusion.kalman}}\label{\detokenize{modules/fusion/kalman:qnav-fusion-kalman}}\label{\detokenize{modules/fusion/kalman::doc}}\index{module@\spxentry{module}!qnav.fusion.kalman@\spxentry{qnav.fusion.kalman}}\index{qnav.fusion.kalman@\spxentry{qnav.fusion.kalman}!module@\spxentry{module}}

\paragraph{kalman.py}
\label{\detokenize{modules/fusion/kalman:kalman-py}}\begin{quote}\begin{description}
\sphinxlineitem{summary}
\sphinxAtStartPar
Kalman basd fusion method for supplying Inertial Navigation.
Provides Kalman Filter fusion for the Inertial Navigation system (INS).

\sphinxlineitem{authors}
\begin{DUlineblock}{0em}
\item[] Prof. Jason Ralph \sphinxhyphen{} \sphinxhref{mailto:jfralph@liverpool.ac.uk}{jfralph@liverpool.ac.uk}
\item[] Michael Wright \sphinxhyphen{} \sphinxhref{mailto:mjwright@liverpool.ac.uk}{mjwright@liverpool.ac.uk}
\end{DUlineblock}

\sphinxlineitem{version}
\sphinxAtStartPar
1.0.0 \sphinxhyphen{} First full implementation of module.

\end{description}\end{quote}
\index{KalmanINS (class in qnav.fusion.kalman)@\spxentry{KalmanINS}\spxextra{class in qnav.fusion.kalman}}

\begin{fulllineitems}
\phantomsection\label{\detokenize{modules/fusion/kalman:qnav.fusion.kalman.KalmanINS}}
\pysigstartsignatures
\pysiglinewithargsret{\sphinxbfcode{\sphinxupquote{class\DUrole{w}{ }}}\sphinxcode{\sphinxupquote{qnav.fusion.kalman.}}\sphinxbfcode{\sphinxupquote{KalmanINS}}}{\sphinxparam{\DUrole{n}{accelerometer}\DUrole{p}{:}\DUrole{w}{ }\DUrole{n}{{\hyperref[\detokenize{modules/measurement/accelerometer:qnav.measurement.accelerometer.Accelerometer}]{\sphinxcrossref{Accelerometer}}}}}\sphinxparamcomma \sphinxparam{\DUrole{n}{gyroscope}\DUrole{p}{:}\DUrole{w}{ }\DUrole{n}{{\hyperref[\detokenize{modules/measurement/gyroscope:qnav.measurement.gyroscope.Gyroscope}]{\sphinxcrossref{Gyroscope}}}}}\sphinxparamcomma \sphinxparam{\DUrole{n}{estimated\_acc\_axis}\DUrole{p}{:}\DUrole{w}{ }\DUrole{n}{{\hyperref[\detokenize{modules/measurement/platform:qnav.measurement.platform.SensorAxis}]{\sphinxcrossref{SensorAxis}}}}\DUrole{w}{ }\DUrole{o}{=}\DUrole{w}{ }\DUrole{default_value}{None}}\sphinxparamcomma \sphinxparam{\DUrole{n}{estimated\_gyro\_axis}\DUrole{p}{:}\DUrole{w}{ }\DUrole{n}{{\hyperref[\detokenize{modules/measurement/platform:qnav.measurement.platform.SensorAxis}]{\sphinxcrossref{SensorAxis}}}}\DUrole{w}{ }\DUrole{o}{=}\DUrole{w}{ }\DUrole{default_value}{None}}}{}
\pysigstopsignatures
\sphinxAtStartPar
Kalman Filter Inertial Navigation System (INS).
This class represents a Kalman filter that uses linear quadratic
estimations to update the estimated states and estimations of
measurement errors. This solution is known to be more accurate
than the other more basic INS methods provided in the project.
\index{perform\_fusion() (qnav.fusion.kalman.KalmanINS method)@\spxentry{perform\_fusion()}\spxextra{qnav.fusion.kalman.KalmanINS method}}

\begin{fulllineitems}
\phantomsection\label{\detokenize{modules/fusion/kalman:qnav.fusion.kalman.KalmanINS.perform_fusion}}
\pysigstartsignatures
\pysiglinewithargsret{\sphinxbfcode{\sphinxupquote{perform\_fusion}}}{\sphinxparam{\DUrole{n}{kalman\_state}\DUrole{p}{:}\DUrole{w}{ }\DUrole{n}{{\hyperref[\detokenize{modules/estimation/state:qnav.estimation.state.KalmanEstimatedState}]{\sphinxcrossref{KalmanEstimatedState}}}}}}{{ $\rightarrow$ None}}
\pysigstopsignatures
\sphinxAtStartPar
Performs fusion using latest accelerometer and gyroscope measurements.
Updates the estimated states by performing Kalman Filtering.
\begin{quote}\begin{description}
\sphinxlineitem{Parameters}
\sphinxAtStartPar
\sphinxstyleliteralstrong{\sphinxupquote{kalman\_state}} ({\hyperref[\detokenize{modules/estimation/state:qnav.estimation.state.EstimatedState}]{\sphinxcrossref{\sphinxstyleliteralemphasis{\sphinxupquote{EstimatedState}}}}}) \textendash{} The current estimated state.

\end{description}\end{quote}

\end{fulllineitems}


\end{fulllineitems}

\index{state\_vector2ned() (in module qnav.fusion.kalman)@\spxentry{state\_vector2ned()}\spxextra{in module qnav.fusion.kalman}}

\begin{fulllineitems}
\phantomsection\label{\detokenize{modules/fusion/kalman:qnav.fusion.kalman.state_vector2ned}}
\pysigstartsignatures
\pysiglinewithargsret{\sphinxcode{\sphinxupquote{qnav.fusion.kalman.}}\sphinxbfcode{\sphinxupquote{state\_vector2ned}}}{\sphinxparam{\DUrole{n}{state\_vector}\DUrole{p}{:}\DUrole{w}{ }\DUrole{n}{ndarray}}\sphinxparamcomma \sphinxparam{\DUrole{n}{ref\_lla}\DUrole{p}{:}\DUrole{w}{ }\DUrole{n}{ndarray}}\sphinxparamcomma \sphinxparam{\DUrole{n}{gravity\_model}\DUrole{p}{:}\DUrole{w}{ }\DUrole{n}{{\hyperref[\detokenize{modules/gravity/base:qnav.gravity.base.GravityModel}]{\sphinxcrossref{GravityModel}}}}}}{{ $\rightarrow$ ndarray}}
\pysigstopsignatures
\sphinxAtStartPar
This function converts the state vector array and measurements into local
North\sphinxhyphen{}East\sphinxhyphen{}Down (NED) coordinates from a given reference location,
suitable for kalman filtering.
\begin{quote}\begin{description}
\sphinxlineitem{Parameters}\begin{itemize}
\item {} 
\sphinxAtStartPar
\sphinxstyleliteralstrong{\sphinxupquote{state\_vector}} (\sphinxstyleliteralemphasis{\sphinxupquote{numpy.ndarray}}\sphinxstyleliteralemphasis{\sphinxupquote{ (}}\sphinxstyleliteralemphasis{\sphinxupquote{15\sphinxhyphen{}elements}}\sphinxstyleliteralemphasis{\sphinxupquote{)}}) \textendash{} The state vector array holding the position,
velocity, acceleration, attitude and angle rate variables.

\item {} 
\sphinxAtStartPar
\sphinxstyleliteralstrong{\sphinxupquote{ref\_lla}} (\sphinxstyleliteralemphasis{\sphinxupquote{numpy.ndarray}}\sphinxstyleliteralemphasis{\sphinxupquote{ (}}\sphinxstyleliteralemphasis{\sphinxupquote{3\sphinxhyphen{}elements}}\sphinxstyleliteralemphasis{\sphinxupquote{)}}) \textendash{} The reference latitude\sphinxhyphen{}longitude\sphinxhyphen{}altitude vector (given in
degrees\sphinxhyphen{}degrees\sphinxhyphen{}metres) identifying the centre point.

\item {} 
\sphinxAtStartPar
\sphinxstyleliteralstrong{\sphinxupquote{gravity\_model}} ({\hyperref[\detokenize{modules/gravity/base:qnav.gravity.base.GravityModel}]{\sphinxcrossref{\sphinxstyleliteralemphasis{\sphinxupquote{GravityModel}}}}}) \textendash{} The model for estimating downwards acceleration.

\end{itemize}

\sphinxlineitem{Returns}
\sphinxAtStartPar
The given state vector array and measurements converted to
local NED/Earth axes for Kalman filtering.

\sphinxlineitem{Return type}
\sphinxAtStartPar
numpy.ndarray (15\sphinxhyphen{}elements)

\end{description}\end{quote}

\end{fulllineitems}

\index{ned2state\_vector() (in module qnav.fusion.kalman)@\spxentry{ned2state\_vector()}\spxextra{in module qnav.fusion.kalman}}

\begin{fulllineitems}
\phantomsection\label{\detokenize{modules/fusion/kalman:qnav.fusion.kalman.ned2state_vector}}
\pysigstartsignatures
\pysiglinewithargsret{\sphinxcode{\sphinxupquote{qnav.fusion.kalman.}}\sphinxbfcode{\sphinxupquote{ned2state\_vector}}}{\sphinxparam{\DUrole{n}{state\_vector}\DUrole{p}{:}\DUrole{w}{ }\DUrole{n}{ndarray}}\sphinxparamcomma \sphinxparam{\DUrole{n}{ref\_lla}\DUrole{p}{:}\DUrole{w}{ }\DUrole{n}{ndarray}}\sphinxparamcomma \sphinxparam{\DUrole{n}{gravity\_model}\DUrole{p}{:}\DUrole{w}{ }\DUrole{n}{{\hyperref[\detokenize{modules/gravity/base:qnav.gravity.base.GravityModel}]{\sphinxcrossref{GravityModel}}}}}}{{ $\rightarrow$ ndarray}}
\pysigstopsignatures
\sphinxAtStartPar
This function converts the state vector array back to Latitude\sphinxhyphen{}longitude\sphinxhyphen{}
altitude (LLA) coordinates (in degrees\sphinxhyphen{}degrees\sphinxhyphen{}metres). Please note that
the state error matrix remains in Local North\sphinxhyphen{}East\sphinxhyphen{}Down (NED) coordinates.
\begin{quote}\begin{description}
\sphinxlineitem{Parameters}\begin{itemize}
\item {} 
\sphinxAtStartPar
\sphinxstyleliteralstrong{\sphinxupquote{state\_vector}} (\sphinxstyleliteralemphasis{\sphinxupquote{numpy.ndarray}}\sphinxstyleliteralemphasis{\sphinxupquote{ (}}\sphinxstyleliteralemphasis{\sphinxupquote{15\sphinxhyphen{}elements}}\sphinxstyleliteralemphasis{\sphinxupquote{)}}) \textendash{} The state vector array holding the position,
velocity, acceleration, attitude and angle rate variables.

\item {} 
\sphinxAtStartPar
\sphinxstyleliteralstrong{\sphinxupquote{ref\_lla}} (\sphinxstyleliteralemphasis{\sphinxupquote{numpy.ndarray}}\sphinxstyleliteralemphasis{\sphinxupquote{ (}}\sphinxstyleliteralemphasis{\sphinxupquote{3\sphinxhyphen{}elements}}\sphinxstyleliteralemphasis{\sphinxupquote{)}}) \textendash{} The reference latitude\sphinxhyphen{}longitude\sphinxhyphen{}altitude vector (given in
degrees\sphinxhyphen{}degrees\sphinxhyphen{}metres) identifying the centre point.

\item {} 
\sphinxAtStartPar
\sphinxstyleliteralstrong{\sphinxupquote{gravity\_model}} ({\hyperref[\detokenize{modules/gravity/base:qnav.gravity.base.GravityModel}]{\sphinxcrossref{\sphinxstyleliteralemphasis{\sphinxupquote{GravityModel}}}}}) \textendash{} The model for estimating downwards acceleration.

\end{itemize}

\sphinxlineitem{Returns}
\sphinxAtStartPar
The given state vector array converted to LLA.

\sphinxlineitem{Return type}
\sphinxAtStartPar
numpy.ndarray (15\sphinxhyphen{}elements)

\end{description}\end{quote}

\end{fulllineitems}


\sphinxstepscope


\subsubsection{qnav.fusion.altimeter}
\label{\detokenize{modules/fusion/altimeter:module-qnav.fusion.altimeter}}\label{\detokenize{modules/fusion/altimeter:qnav-fusion-altimeter}}\label{\detokenize{modules/fusion/altimeter::doc}}\index{module@\spxentry{module}!qnav.fusion.altimeter@\spxentry{qnav.fusion.altimeter}}\index{qnav.fusion.altimeter@\spxentry{qnav.fusion.altimeter}!module@\spxentry{module}}

\paragraph{altitude.py}
\label{\detokenize{modules/fusion/altimeter:altitude-py}}\begin{quote}\begin{description}
\sphinxlineitem{summary}
\sphinxAtStartPar
Supplied fusion method for altitude measurements.
This module contains methods that can be used for acquiring altimeter
measurements during simulation and fusing them with the current
estimated altitude. Additional methods may be added in future updates.

\sphinxlineitem{authors}
\begin{DUlineblock}{0em}
\item[] Prof. Jason Ralph \sphinxhyphen{} \sphinxhref{mailto:jfralph@liverpool.ac.uk}{jfralph@liverpool.ac.uk}
\item[] Michael Wright \sphinxhyphen{} \sphinxhref{mailto:mjwright@liverpool.ac.uk}{mjwright@liverpool.ac.uk}
\end{DUlineblock}

\sphinxlineitem{version}
\sphinxAtStartPar
1.0.0 \sphinxhyphen{} First full implementation of module.

\end{description}\end{quote}
\index{FixedGainAltimeter (class in qnav.fusion.altimeter)@\spxentry{FixedGainAltimeter}\spxextra{class in qnav.fusion.altimeter}}

\begin{fulllineitems}
\phantomsection\label{\detokenize{modules/fusion/altimeter:qnav.fusion.altimeter.FixedGainAltimeter}}
\pysigstartsignatures
\pysiglinewithargsret{\sphinxbfcode{\sphinxupquote{class\DUrole{w}{ }}}\sphinxcode{\sphinxupquote{qnav.fusion.altimeter.}}\sphinxbfcode{\sphinxupquote{FixedGainAltimeter}}}{\sphinxparam{\DUrole{n}{altimeter}\DUrole{p}{:}\DUrole{w}{ }\DUrole{n}{{\hyperref[\detokenize{modules/measurement/altimeter:qnav.measurement.altimeter.Altimeter}]{\sphinxcrossref{Altimeter}}}}}\sphinxparamcomma \sphinxparam{\DUrole{n}{gain\_amount}\DUrole{p}{:}\DUrole{w}{ }\DUrole{n}{float}\DUrole{w}{ }\DUrole{o}{=}\DUrole{w}{ }\DUrole{default_value}{0.5}}}{}
\pysigstopsignatures
\sphinxAtStartPar
Fuses altimeter altitude measurements using fixed gain fusion.
Updates the estimated altitude by using a weighted average of the
measurement altitude and current estimated altitude. This gain
amount can be set and updated at any time.
\index{perform\_fusion() (qnav.fusion.altimeter.FixedGainAltimeter method)@\spxentry{perform\_fusion()}\spxextra{qnav.fusion.altimeter.FixedGainAltimeter method}}

\begin{fulllineitems}
\phantomsection\label{\detokenize{modules/fusion/altimeter:qnav.fusion.altimeter.FixedGainAltimeter.perform_fusion}}
\pysigstartsignatures
\pysiglinewithargsret{\sphinxbfcode{\sphinxupquote{perform\_fusion}}}{\sphinxparam{\DUrole{n}{estimated\_state}\DUrole{p}{:}\DUrole{w}{ }\DUrole{n}{{\hyperref[\detokenize{modules/estimation/state:qnav.estimation.state.EstimatedState}]{\sphinxcrossref{EstimatedState}}}}}}{{ $\rightarrow$ None}}
\pysigstopsignatures
\sphinxAtStartPar
Performs the fixed gained fusion with last altimeter measurement.
Simply fuses the last altimeter measurement with previous measurement
using a fixed\sphinxhyphen{}gain (weighted average) update.
\begin{quote}\begin{description}
\sphinxlineitem{Parameters}
\sphinxAtStartPar
\sphinxstyleliteralstrong{\sphinxupquote{estimated\_state}} ({\hyperref[\detokenize{modules/estimation/state:qnav.estimation.state.EstimatedState}]{\sphinxcrossref{\sphinxstyleliteralemphasis{\sphinxupquote{EstimatedState}}}}}) \textendash{} The current estimated state.

\end{description}\end{quote}

\end{fulllineitems}

\index{fixed\_gain (qnav.fusion.altimeter.FixedGainAltimeter property)@\spxentry{fixed\_gain}\spxextra{qnav.fusion.altimeter.FixedGainAltimeter property}}

\begin{fulllineitems}
\phantomsection\label{\detokenize{modules/fusion/altimeter:qnav.fusion.altimeter.FixedGainAltimeter.fixed_gain}}
\pysigstartsignatures
\pysigline{\sphinxbfcode{\sphinxupquote{property\DUrole{w}{ }}}\sphinxbfcode{\sphinxupquote{fixed\_gain}}\sphinxbfcode{\sphinxupquote{\DUrole{p}{:}\DUrole{w}{ }float}}}
\pysigstopsignatures
\sphinxAtStartPar
Property getter method for fixed gain amount.
\begin{quote}\begin{description}
\sphinxlineitem{Returns}
\sphinxAtStartPar
The current gain amount for the altimeter measurements.

\sphinxlineitem{Return type}
\sphinxAtStartPar
float

\end{description}\end{quote}

\end{fulllineitems}


\end{fulllineitems}

\index{AlphaBetaAltimeter (class in qnav.fusion.altimeter)@\spxentry{AlphaBetaAltimeter}\spxextra{class in qnav.fusion.altimeter}}

\begin{fulllineitems}
\phantomsection\label{\detokenize{modules/fusion/altimeter:qnav.fusion.altimeter.AlphaBetaAltimeter}}
\pysigstartsignatures
\pysiglinewithargsret{\sphinxbfcode{\sphinxupquote{class\DUrole{w}{ }}}\sphinxcode{\sphinxupquote{qnav.fusion.altimeter.}}\sphinxbfcode{\sphinxupquote{AlphaBetaAltimeter}}}{\sphinxparam{\DUrole{n}{altimeter}\DUrole{p}{:}\DUrole{w}{ }\DUrole{n}{{\hyperref[\detokenize{modules/measurement/altimeter:qnav.measurement.altimeter.Altimeter}]{\sphinxcrossref{Altimeter}}}}}\sphinxparamcomma \sphinxparam{\DUrole{n}{alpha}\DUrole{p}{:}\DUrole{w}{ }\DUrole{n}{float}\DUrole{w}{ }\DUrole{o}{=}\DUrole{w}{ }\DUrole{default_value}{0.9}}\sphinxparamcomma \sphinxparam{\DUrole{n}{beta}\DUrole{p}{:}\DUrole{w}{ }\DUrole{n}{float}\DUrole{w}{ }\DUrole{o}{=}\DUrole{w}{ }\DUrole{default_value}{None}}}{}
\pysigstopsignatures
\sphinxAtStartPar
Fuses altimeter altitude measurements using alpha\sphinxhyphen{}beta filtering.
Updates the estimated altitude by using an alpha\sphinxhyphen{}beta filter, where alpha
relates to the expected accuracy of the altimeter and beta relates to the
precision of the vertical velocity estimate. This method updates both
estimated altitude and the vertical velocity.
\index{perform\_fusion() (qnav.fusion.altimeter.AlphaBetaAltimeter method)@\spxentry{perform\_fusion()}\spxextra{qnav.fusion.altimeter.AlphaBetaAltimeter method}}

\begin{fulllineitems}
\phantomsection\label{\detokenize{modules/fusion/altimeter:qnav.fusion.altimeter.AlphaBetaAltimeter.perform_fusion}}
\pysigstartsignatures
\pysiglinewithargsret{\sphinxbfcode{\sphinxupquote{perform\_fusion}}}{\sphinxparam{\DUrole{n}{estimated\_state}\DUrole{p}{:}\DUrole{w}{ }\DUrole{n}{{\hyperref[\detokenize{modules/estimation/state:qnav.estimation.state.EstimatedState}]{\sphinxcrossref{EstimatedState}}}}}}{{ $\rightarrow$ None}}
\pysigstopsignatures
\sphinxAtStartPar
Performs alpha\sphinxhyphen{}beta fusion with latest altimeter measurement.
Updates both the estimated altitude and the vertical velocity using
the latest measurement from the altimeter. From this both the
vertical position and vertical velocity are updated.
\begin{quote}\begin{description}
\sphinxlineitem{Parameters}
\sphinxAtStartPar
\sphinxstyleliteralstrong{\sphinxupquote{estimated\_state}} ({\hyperref[\detokenize{modules/estimation/state:qnav.estimation.state.EstimatedState}]{\sphinxcrossref{\sphinxstyleliteralemphasis{\sphinxupquote{EstimatedState}}}}}) \textendash{} The current estimated state.

\end{description}\end{quote}

\end{fulllineitems}

\index{alpha (qnav.fusion.altimeter.AlphaBetaAltimeter property)@\spxentry{alpha}\spxextra{qnav.fusion.altimeter.AlphaBetaAltimeter property}}

\begin{fulllineitems}
\phantomsection\label{\detokenize{modules/fusion/altimeter:qnav.fusion.altimeter.AlphaBetaAltimeter.alpha}}
\pysigstartsignatures
\pysigline{\sphinxbfcode{\sphinxupquote{property\DUrole{w}{ }}}\sphinxbfcode{\sphinxupquote{alpha}}\sphinxbfcode{\sphinxupquote{\DUrole{p}{:}\DUrole{w}{ }float}}}
\pysigstopsignatures
\sphinxAtStartPar
Property getter method for the alpha value.
\begin{quote}\begin{description}
\sphinxlineitem{Returns}
\sphinxAtStartPar
The current alpha value to use for filtering.

\sphinxlineitem{Return type}
\sphinxAtStartPar
float

\end{description}\end{quote}

\end{fulllineitems}

\index{beta (qnav.fusion.altimeter.AlphaBetaAltimeter property)@\spxentry{beta}\spxextra{qnav.fusion.altimeter.AlphaBetaAltimeter property}}

\begin{fulllineitems}
\phantomsection\label{\detokenize{modules/fusion/altimeter:qnav.fusion.altimeter.AlphaBetaAltimeter.beta}}
\pysigstartsignatures
\pysigline{\sphinxbfcode{\sphinxupquote{property\DUrole{w}{ }}}\sphinxbfcode{\sphinxupquote{beta}}\sphinxbfcode{\sphinxupquote{\DUrole{p}{:}\DUrole{w}{ }float}}}
\pysigstopsignatures
\sphinxAtStartPar
Property getter method for the beta value.
\begin{quote}\begin{description}
\sphinxlineitem{Returns}
\sphinxAtStartPar
The current beta value to use for filtering.

\sphinxlineitem{Return type}
\sphinxAtStartPar
float

\end{description}\end{quote}

\end{fulllineitems}


\end{fulllineitems}



\subsection{GPS}
\label{\detokenize{usage/packages:gps}}
\sphinxAtStartPar
Modules relating to GPS sensor and fusion methods:

\sphinxstepscope


\subsubsection{qnav.gps.ephemeris}
\label{\detokenize{modules/gps/ephemeris:module-qnav.gps.ephemeris}}\label{\detokenize{modules/gps/ephemeris:qnav-gps-ephemeris}}\label{\detokenize{modules/gps/ephemeris::doc}}\index{module@\spxentry{module}!qnav.gps.ephemeris@\spxentry{qnav.gps.ephemeris}}\index{qnav.gps.ephemeris@\spxentry{qnav.gps.ephemeris}!module@\spxentry{module}}

\paragraph{ephemeris.py}
\label{\detokenize{modules/gps/ephemeris:ephemeris-py}}\begin{quote}\begin{description}
\sphinxlineitem{summary}
\sphinxAtStartPar
Collection of functions for extracting GPS/GNSS information.
This module is responsible for the util satellite functionality within
the toolbox. This includes functions for downloading GPS information,
parsing RINEX files, extracting satellite information and estimating
missing values.

\sphinxlineitem{authors}
\begin{DUlineblock}{0em}
\item[] Michael Wright \sphinxhyphen{} \sphinxhref{mailto:mjwright@liverpool.ac.uk}{mjwright@liverpool.ac.uk}
\item[] Dr. Niall Moroney \sphinxhyphen{} \sphinxhref{mailto:niall.moroney17@imperial.ac.uk}{niall.moroney17@imperial.ac.uk}
\end{DUlineblock}

\sphinxlineitem{version}
\sphinxAtStartPar
1.0.0 \sphinxhyphen{} First implementation of this module.

\end{description}\end{quote}
\index{SatelliteType (class in qnav.gps.ephemeris)@\spxentry{SatelliteType}\spxextra{class in qnav.gps.ephemeris}}

\begin{fulllineitems}
\phantomsection\label{\detokenize{modules/gps/ephemeris:qnav.gps.ephemeris.SatelliteType}}
\pysigstartsignatures
\pysiglinewithargsret{\sphinxbfcode{\sphinxupquote{class\DUrole{w}{ }}}\sphinxcode{\sphinxupquote{qnav.gps.ephemeris.}}\sphinxbfcode{\sphinxupquote{SatelliteType}}}{\sphinxparam{\DUrole{n}{value}}\sphinxparamcomma \sphinxparam{\DUrole{n}{names=\textless{}not given\textgreater{}}}\sphinxparamcomma \sphinxparam{\DUrole{n}{*values}}\sphinxparamcomma \sphinxparam{\DUrole{n}{module=None}}\sphinxparamcomma \sphinxparam{\DUrole{n}{qualname=None}}\sphinxparamcomma \sphinxparam{\DUrole{n}{type=None}}\sphinxparamcomma \sphinxparam{\DUrole{n}{start=1}}\sphinxparamcomma \sphinxparam{\DUrole{n}{boundary=None}}}{}
\pysigstopsignatures
\sphinxAtStartPar
Character identifiers for the different satellite types.
These characters identify the record types corresponding to the satellite
systems supported inside a Rinex File (as of version 4).

\end{fulllineitems}

\index{get\_rinex\_path() (in module qnav.gps.ephemeris)@\spxentry{get\_rinex\_path()}\spxextra{in module qnav.gps.ephemeris}}

\begin{fulllineitems}
\phantomsection\label{\detokenize{modules/gps/ephemeris:qnav.gps.ephemeris.get_rinex_path}}
\pysigstartsignatures
\pysiglinewithargsret{\sphinxcode{\sphinxupquote{qnav.gps.ephemeris.}}\sphinxbfcode{\sphinxupquote{get\_rinex\_path}}}{\sphinxparam{\DUrole{n}{gps\_time}\DUrole{p}{:}\DUrole{w}{ }\DUrole{n}{datetime}}\sphinxparamcomma \sphinxparam{\DUrole{n}{directory}\DUrole{p}{:}\DUrole{w}{ }\DUrole{n}{Path}}}{{ $\rightarrow$ Path}}
\pysigstopsignatures
\sphinxAtStartPar
Obtain the location of a local RINEX file identified by GPS date and time.
This function will return the local filepath of a RINEX file corresponding
to a given (python) datetime. The file of the returned path may or may not
exist (requiring downloading if missing).
\begin{quote}\begin{description}
\sphinxlineitem{Parameters}\begin{itemize}
\item {} 
\sphinxAtStartPar
\sphinxstyleliteralstrong{\sphinxupquote{gps\_time}} (\sphinxstyleliteralemphasis{\sphinxupquote{datetime.datetime}}) \textendash{} Python datetime object for the time required for the ephemeris.

\item {} 
\sphinxAtStartPar
\sphinxstyleliteralstrong{\sphinxupquote{directory}} (\sphinxstyleliteralemphasis{\sphinxupquote{Path}}) \textendash{} Location for the files to be downloaded to.

\end{itemize}

\sphinxlineitem{Returns}
\sphinxAtStartPar
The path of the corresponding RINEX file.

\sphinxlineitem{Return type}
\sphinxAtStartPar
Path

\end{description}\end{quote}

\end{fulllineitems}

\index{download\_rinex\_file() (in module qnav.gps.ephemeris)@\spxentry{download\_rinex\_file()}\spxextra{in module qnav.gps.ephemeris}}

\begin{fulllineitems}
\phantomsection\label{\detokenize{modules/gps/ephemeris:qnav.gps.ephemeris.download_rinex_file}}
\pysigstartsignatures
\pysiglinewithargsret{\sphinxcode{\sphinxupquote{qnav.gps.ephemeris.}}\sphinxbfcode{\sphinxupquote{download\_rinex\_file}}}{\sphinxparam{\DUrole{n}{gps\_time}\DUrole{p}{:}\DUrole{w}{ }\DUrole{n}{datetime}}\sphinxparamcomma \sphinxparam{\DUrole{n}{database\_dir}\DUrole{p}{:}\DUrole{w}{ }\DUrole{n}{Path}}}{}
\pysigstopsignatures
\sphinxAtStartPar
This function downloads and extracts an ephemeris file for a specified
day from “cddis.gsfc.nasa.gov” and saves it under the provided directory.
\begin{quote}\begin{description}
\sphinxlineitem{Parameters}\begin{itemize}
\item {} 
\sphinxAtStartPar
\sphinxstyleliteralstrong{\sphinxupquote{gps\_time}} \textendash{} The Python datetime object for the time required for the
ephemeris data.

\item {} 
\sphinxAtStartPar
\sphinxstyleliteralstrong{\sphinxupquote{database\_dir}} (\sphinxstyleliteralemphasis{\sphinxupquote{Path}}) \textendash{} The location for the files to be downloaded to.

\end{itemize}

\sphinxlineitem{Type}
\sphinxAtStartPar
datetime.datetime

\end{description}\end{quote}

\end{fulllineitems}

\index{download\_rinex\_files() (in module qnav.gps.ephemeris)@\spxentry{download\_rinex\_files()}\spxextra{in module qnav.gps.ephemeris}}

\begin{fulllineitems}
\phantomsection\label{\detokenize{modules/gps/ephemeris:qnav.gps.ephemeris.download_rinex_files}}
\pysigstartsignatures
\pysiglinewithargsret{\sphinxcode{\sphinxupquote{qnav.gps.ephemeris.}}\sphinxbfcode{\sphinxupquote{download\_rinex\_files}}}{\sphinxparam{\DUrole{n}{from\_time}\DUrole{p}{:}\DUrole{w}{ }\DUrole{n}{datetime}}\sphinxparamcomma \sphinxparam{\DUrole{n}{to\_time}\DUrole{p}{:}\DUrole{w}{ }\DUrole{n}{datetime}}\sphinxparamcomma \sphinxparam{\DUrole{n}{database\_dir}\DUrole{p}{:}\DUrole{w}{ }\DUrole{n}{Path}}\sphinxparamcomma \sphinxparam{\DUrole{n}{inc\_last}\DUrole{p}{:}\DUrole{w}{ }\DUrole{n}{bool}\DUrole{w}{ }\DUrole{o}{=}\DUrole{w}{ }\DUrole{default_value}{True}}}{}
\pysigstopsignatures
\sphinxAtStartPar
This function downloads and extracts multiple ephemeris files.
Given a date range, this function will download all the ephemeris
files covered by the range.
\begin{quote}\begin{description}
\sphinxlineitem{Parameters}\begin{itemize}
\item {} 
\sphinxAtStartPar
\sphinxstyleliteralstrong{\sphinxupquote{from\_time}} (\sphinxstyleliteralemphasis{\sphinxupquote{datetime.datetime}}) \textendash{} The first date for the download range.
This is the date for the first GPS file to be downloaded.

\item {} 
\sphinxAtStartPar
\sphinxstyleliteralstrong{\sphinxupquote{to\_time}} (\sphinxstyleliteralemphasis{\sphinxupquote{datetime.datetime}}) \textendash{} The last date for the download range.
This is the end date for the GPS files to download.

\item {} 
\sphinxAtStartPar
\sphinxstyleliteralstrong{\sphinxupquote{database\_dir}} (\sphinxstyleliteralemphasis{\sphinxupquote{Path}}) \textendash{} The location for the files to be downloaded to.

\item {} 
\sphinxAtStartPar
\sphinxstyleliteralstrong{\sphinxupquote{inc\_last}} (\sphinxstyleliteralemphasis{\sphinxupquote{bool}}) \textendash{} Should the last date of the range be included?

\end{itemize}

\end{description}\end{quote}

\end{fulllineitems}

\index{read\_rinex\_file() (in module qnav.gps.ephemeris)@\spxentry{read\_rinex\_file()}\spxextra{in module qnav.gps.ephemeris}}

\begin{fulllineitems}
\phantomsection\label{\detokenize{modules/gps/ephemeris:qnav.gps.ephemeris.read_rinex_file}}
\pysigstartsignatures
\pysiglinewithargsret{\sphinxcode{\sphinxupquote{qnav.gps.ephemeris.}}\sphinxbfcode{\sphinxupquote{read\_rinex\_file}}}{\sphinxparam{\DUrole{n}{ephemeris\_file}\DUrole{p}{:}\DUrole{w}{ }\DUrole{n}{Path}}}{{ $\rightarrow$ dict\DUrole{p}{{[}}str\DUrole{p}{,}\DUrole{w}{ }dict\DUrole{p}{{]}}}}
\pysigstopsignatures
\sphinxAtStartPar
Parses a downloaded ephemeris file and returns its data.
Reads a given ephemeris file, skipping the header contents and returning
the contained entry records in dictionaries. Each is keyed by the ID of
the satellite.
\begin{quote}\begin{description}
\sphinxlineitem{Parameters}
\sphinxAtStartPar
\sphinxstyleliteralstrong{\sphinxupquote{ephemeris\_file}} (\sphinxstyleliteralemphasis{\sphinxupquote{Path}}) \textendash{} The local ephemeris file to read.

\sphinxlineitem{Returns}
\sphinxAtStartPar
A dictionary of extracted satellite parameters.

\sphinxlineitem{Return type}
\sphinxAtStartPar
dict{[}str, dict{]}

\end{description}\end{quote}

\end{fulllineitems}

\index{rinex\_to\_satellite\_params() (in module qnav.gps.ephemeris)@\spxentry{rinex\_to\_satellite\_params()}\spxextra{in module qnav.gps.ephemeris}}

\begin{fulllineitems}
\phantomsection\label{\detokenize{modules/gps/ephemeris:qnav.gps.ephemeris.rinex_to_satellite_params}}
\pysigstartsignatures
\pysiglinewithargsret{\sphinxcode{\sphinxupquote{qnav.gps.ephemeris.}}\sphinxbfcode{\sphinxupquote{rinex\_to\_satellite\_params}}}{\sphinxparam{\DUrole{n}{rinex\_data}\DUrole{p}{:}\DUrole{w}{ }\DUrole{n}{dict}}\sphinxparamcomma \sphinxparam{\DUrole{o}{*}\DUrole{n}{satellite\_types}\DUrole{p}{:}\DUrole{w}{ }\DUrole{n}{{\hyperref[\detokenize{modules/gps/ephemeris:qnav.gps.ephemeris.SatelliteType}]{\sphinxcrossref{SatelliteType}}}}}}{{ $\rightarrow$ list\DUrole{p}{{[}}{\hyperref[\detokenize{modules/gps/satellite:qnav.gps.satellite.SatelliteParams}]{\sphinxcrossref{SatelliteParams}}}\DUrole{p}{{]}}}}
\pysigstopsignatures
\sphinxAtStartPar
Extracts given rinex data and converts into satellite parameters class.
A helper function for extracting the usable satellite parameter data
from given Rinex data, allowing for satellite instances to be created.
Optional filtering can also be supplied to limit returns to specific
satellite types.
\begin{quote}\begin{description}
\sphinxlineitem{Parameters}\begin{itemize}
\item {} 
\sphinxAtStartPar
\sphinxstyleliteralstrong{\sphinxupquote{rinex\_data}} (\sphinxstyleliteralemphasis{\sphinxupquote{dict}}\sphinxstyleliteralemphasis{\sphinxupquote{{[}}}\sphinxstyleliteralemphasis{\sphinxupquote{str}}\sphinxstyleliteralemphasis{\sphinxupquote{, }}\sphinxstyleliteralemphasis{\sphinxupquote{dict}}\sphinxstyleliteralemphasis{\sphinxupquote{{]}}}) \textendash{} The data records read from a Rinex file.

\item {} 
\sphinxAtStartPar
\sphinxstyleliteralstrong{\sphinxupquote{satellite\_types}} ({\hyperref[\detokenize{modules/gps/ephemeris:qnav.gps.ephemeris.SatelliteType}]{\sphinxcrossref{\sphinxstyleliteralemphasis{\sphinxupquote{SatelliteType}}}}}) \textendash{} (Optional) The satellite types to filter by.
By default, all types are selected and no filtering is performed.

\end{itemize}

\sphinxlineitem{Returns}
\sphinxAtStartPar
The satellite parameters for each rinex record.

\sphinxlineitem{Return type}
\sphinxAtStartPar
list{[}{\hyperref[\detokenize{modules/gps/satellite:qnav.gps.satellite.SatelliteParams}]{\sphinxcrossref{SatelliteParams}}}{]}

\end{description}\end{quote}

\end{fulllineitems}


\sphinxstepscope


\subsubsection{qnav.gps.fusion}
\label{\detokenize{modules/gps/fusion:module-qnav.gps.fusion}}\label{\detokenize{modules/gps/fusion:qnav-gps-fusion}}\label{\detokenize{modules/gps/fusion::doc}}\index{module@\spxentry{module}!qnav.gps.fusion@\spxentry{qnav.gps.fusion}}\index{qnav.gps.fusion@\spxentry{qnav.gps.fusion}!module@\spxentry{module}}

\paragraph{fusion.py}
\label{\detokenize{modules/gps/fusion:fusion-py}}\begin{quote}\begin{description}
\sphinxlineitem{summary}
\sphinxAtStartPar
A collection of GPS/GNSS methods for fusing GPS measurements with INS.
This module contains all the GPS fusion methods for merging GPS
position and velocity updates with the INS current estimation.
Simpler methods will simply merge the estimates, whereas more
sophisticated methods will adapt expected errors to compensate
for INS noise/drift.

\sphinxlineitem{authors}
\begin{DUlineblock}{0em}
\item[] Kallum O’Hara \sphinxhyphen{} \sphinxhref{mailto:sgkohara@liverpool.ac.uk}{sgkohara@liverpool.ac.uk}
\item[] Michael Wright \sphinxhyphen{} \sphinxhref{mailto:mjwright@liverpool.ac.uk}{mjwright@liverpool.ac.uk}
\end{DUlineblock}

\sphinxlineitem{version}
\sphinxAtStartPar
1.0.0 \sphinxhyphen{} First implementation of this module.

\end{description}\end{quote}
\index{GpsFusion (class in qnav.gps.fusion)@\spxentry{GpsFusion}\spxextra{class in qnav.gps.fusion}}

\begin{fulllineitems}
\phantomsection\label{\detokenize{modules/gps/fusion:qnav.gps.fusion.GpsFusion}}
\pysigstartsignatures
\pysiglinewithargsret{\sphinxbfcode{\sphinxupquote{class\DUrole{w}{ }}}\sphinxcode{\sphinxupquote{qnav.gps.fusion.}}\sphinxbfcode{\sphinxupquote{GpsFusion}}}{\sphinxparam{\DUrole{n}{gps\_sensor}\DUrole{p}{:}\DUrole{w}{ }\DUrole{n}{{\hyperref[\detokenize{modules/gps/sensor:qnav.gps.sensor.GpsSensor}]{\sphinxcrossref{GpsSensor}}}}}\sphinxparamcomma \sphinxparam{\DUrole{n}{estimated\_state}\DUrole{p}{:}\DUrole{w}{ }\DUrole{n}{{\hyperref[\detokenize{modules/estimation/state:qnav.estimation.state.EstimatedState}]{\sphinxcrossref{EstimatedState}}}}}\sphinxparamcomma \sphinxparam{\DUrole{n}{estimated\_clock\_bias}\DUrole{p}{:}\DUrole{w}{ }\DUrole{n}{float}\DUrole{w}{ }\DUrole{o}{=}\DUrole{w}{ }\DUrole{default_value}{0.0}}\sphinxparamcomma \sphinxparam{\DUrole{n}{estimated\_clock\_drift}\DUrole{p}{:}\DUrole{w}{ }\DUrole{n}{float}\DUrole{w}{ }\DUrole{o}{=}\DUrole{w}{ }\DUrole{default_value}{0.0}}\sphinxparamcomma \sphinxparam{\DUrole{n}{burn\_in\_iterations}\DUrole{p}{:}\DUrole{w}{ }\DUrole{n}{int}\DUrole{w}{ }\DUrole{o}{=}\DUrole{w}{ }\DUrole{default_value}{3}}}{}
\pysigstopsignatures
\sphinxAtStartPar
The abstract base class for fusion methods.
This class outlines the structure each fusion method must extend from.
This essentially serves as a template class containing methods and
properties that must be incorporated by each GPS fusion method.
\index{solve() (qnav.gps.fusion.GpsFusion method)@\spxentry{solve()}\spxextra{qnav.gps.fusion.GpsFusion method}}

\begin{fulllineitems}
\phantomsection\label{\detokenize{modules/gps/fusion:qnav.gps.fusion.GpsFusion.solve}}
\pysigstartsignatures
\pysiglinewithargsret{\sphinxbfcode{\sphinxupquote{solve}}}{}{{ $\rightarrow$ None}}
\pysigstopsignatures
\sphinxAtStartPar
Performs solving of GPS measurements to estimated values.
Essentially this translates the GPS satellite measurements into
internal local estimates.

\end{fulllineitems}


\end{fulllineitems}

\index{GpsFixedGainFusion (class in qnav.gps.fusion)@\spxentry{GpsFixedGainFusion}\spxextra{class in qnav.gps.fusion}}

\begin{fulllineitems}
\phantomsection\label{\detokenize{modules/gps/fusion:qnav.gps.fusion.GpsFixedGainFusion}}
\pysigstartsignatures
\pysiglinewithargsret{\sphinxbfcode{\sphinxupquote{class\DUrole{w}{ }}}\sphinxcode{\sphinxupquote{qnav.gps.fusion.}}\sphinxbfcode{\sphinxupquote{GpsFixedGainFusion}}}{\sphinxparam{\DUrole{n}{gps\_sensor}\DUrole{p}{:}\DUrole{w}{ }\DUrole{n}{{\hyperref[\detokenize{modules/gps/sensor:qnav.gps.sensor.GpsSensor}]{\sphinxcrossref{GpsSensor}}}}}\sphinxparamcomma \sphinxparam{\DUrole{n}{estimated\_state}\DUrole{p}{:}\DUrole{w}{ }\DUrole{n}{{\hyperref[\detokenize{modules/estimation/state:qnav.estimation.state.EstimatedState}]{\sphinxcrossref{EstimatedState}}}}}\sphinxparamcomma \sphinxparam{\DUrole{n}{gain\_amount}\DUrole{p}{:}\DUrole{w}{ }\DUrole{n}{float}\DUrole{w}{ }\DUrole{o}{=}\DUrole{w}{ }\DUrole{default_value}{0.1}}}{}
\pysigstopsignatures
\sphinxAtStartPar
Performs simple fixed\sphinxhyphen{}gain GPS fusion.
Updates the estimated state’s position and velocity using simple fixed
gained fusion. Can be used with any estimated state class.
\index{perform\_fusion() (qnav.gps.fusion.GpsFixedGainFusion method)@\spxentry{perform\_fusion()}\spxextra{qnav.gps.fusion.GpsFixedGainFusion method}}

\begin{fulllineitems}
\phantomsection\label{\detokenize{modules/gps/fusion:qnav.gps.fusion.GpsFixedGainFusion.perform_fusion}}
\pysigstartsignatures
\pysiglinewithargsret{\sphinxbfcode{\sphinxupquote{perform\_fusion}}}{\sphinxparam{\DUrole{n}{estimated\_state}\DUrole{p}{:}\DUrole{w}{ }\DUrole{n}{{\hyperref[\detokenize{modules/estimation/state:qnav.estimation.state.EstimatedState}]{\sphinxcrossref{EstimatedState}}}}}}{{ $\rightarrow$ None}}
\pysigstopsignatures
\sphinxAtStartPar
Performs GPS correction using simple fixed\sphinxhyphen{}gained fusion.
\begin{quote}\begin{description}
\sphinxlineitem{Parameters}
\sphinxAtStartPar
\sphinxstyleliteralstrong{\sphinxupquote{estimated\_state}} ({\hyperref[\detokenize{modules/estimation/state:qnav.estimation.state.EstimatedState}]{\sphinxcrossref{\sphinxstyleliteralemphasis{\sphinxupquote{EstimatedState}}}}}) \textendash{} The current estimated state.

\end{description}\end{quote}

\end{fulllineitems}


\end{fulllineitems}

\index{GpsLooseFusion (class in qnav.gps.fusion)@\spxentry{GpsLooseFusion}\spxextra{class in qnav.gps.fusion}}

\begin{fulllineitems}
\phantomsection\label{\detokenize{modules/gps/fusion:qnav.gps.fusion.GpsLooseFusion}}
\pysigstartsignatures
\pysiglinewithargsret{\sphinxbfcode{\sphinxupquote{class\DUrole{w}{ }}}\sphinxcode{\sphinxupquote{qnav.gps.fusion.}}\sphinxbfcode{\sphinxupquote{GpsLooseFusion}}}{\sphinxparam{\DUrole{n}{gps\_sensor}\DUrole{p}{:}\DUrole{w}{ }\DUrole{n}{{\hyperref[\detokenize{modules/gps/sensor:qnav.gps.sensor.GpsSensor}]{\sphinxcrossref{GpsSensor}}}}}\sphinxparamcomma \sphinxparam{\DUrole{n}{estimated\_state}\DUrole{p}{:}\DUrole{w}{ }\DUrole{n}{{\hyperref[\detokenize{modules/estimation/state:qnav.estimation.state.KalmanEstimatedState}]{\sphinxcrossref{KalmanEstimatedState}}}}}\sphinxparamcomma \sphinxparam{\DUrole{n}{pos\_cov\_mat\_ned}\DUrole{p}{:}\DUrole{w}{ }\DUrole{n}{ndarray}}\sphinxparamcomma \sphinxparam{\DUrole{n}{vel\_cov\_mat\_ned}\DUrole{p}{:}\DUrole{w}{ }\DUrole{n}{ndarray}}}{}
\pysigstopsignatures
\sphinxAtStartPar
Performs loosely\sphinxhyphen{}coupled based GPS fusion.
A fusion method that uses loosely\sphinxhyphen{}coupled fusion for updating the INS
with given GPS estimates. To further improve the accuracy of the
navigation solution, the error states are fed back to the INS.
Must be used with a Kalman estimated state class.
\index{perform\_fusion() (qnav.gps.fusion.GpsLooseFusion method)@\spxentry{perform\_fusion()}\spxextra{qnav.gps.fusion.GpsLooseFusion method}}

\begin{fulllineitems}
\phantomsection\label{\detokenize{modules/gps/fusion:qnav.gps.fusion.GpsLooseFusion.perform_fusion}}
\pysigstartsignatures
\pysiglinewithargsret{\sphinxbfcode{\sphinxupquote{perform\_fusion}}}{\sphinxparam{\DUrole{n}{kalman\_state}\DUrole{p}{:}\DUrole{w}{ }\DUrole{n}{{\hyperref[\detokenize{modules/estimation/state:qnav.estimation.state.KalmanEstimatedState}]{\sphinxcrossref{KalmanEstimatedState}}}}}}{{ $\rightarrow$ None}}
\pysigstopsignatures
\sphinxAtStartPar
Applies loose\sphinxhyphen{}coupled fusion using GPS estimates and INS solution.
This method updates the given INS instance’s estimates by fusing them
with the given GPS estimates. The INS states are updated internally.
\begin{quote}\begin{description}
\sphinxlineitem{Parameters}
\sphinxAtStartPar
\sphinxstyleliteralstrong{\sphinxupquote{kalman\_state}} ({\hyperref[\detokenize{modules/estimation/state:qnav.estimation.state.KalmanEstimatedState}]{\sphinxcrossref{\sphinxstyleliteralemphasis{\sphinxupquote{KalmanEstimatedState}}}}}) \textendash{} The current estimated state.

\end{description}\end{quote}

\end{fulllineitems}


\end{fulllineitems}

\index{GpsTightFusion (class in qnav.gps.fusion)@\spxentry{GpsTightFusion}\spxextra{class in qnav.gps.fusion}}

\begin{fulllineitems}
\phantomsection\label{\detokenize{modules/gps/fusion:qnav.gps.fusion.GpsTightFusion}}
\pysigstartsignatures
\pysiglinewithargsret{\sphinxbfcode{\sphinxupquote{class\DUrole{w}{ }}}\sphinxcode{\sphinxupquote{qnav.gps.fusion.}}\sphinxbfcode{\sphinxupquote{GpsTightFusion}}}{\sphinxparam{\DUrole{n}{gps\_sensor}\DUrole{p}{:}\DUrole{w}{ }\DUrole{n}{{\hyperref[\detokenize{modules/gps/sensor:qnav.gps.sensor.GpsSensor}]{\sphinxcrossref{GpsSensor}}}}}\sphinxparamcomma \sphinxparam{\DUrole{n}{estimated\_state}\DUrole{p}{:}\DUrole{w}{ }\DUrole{n}{{\hyperref[\detokenize{modules/estimation/state:qnav.estimation.state.KalmanEstimatedState}]{\sphinxcrossref{KalmanEstimatedState}}}}}\sphinxparamcomma \sphinxparam{\DUrole{n}{tight\_state\_vector}\DUrole{p}{:}\DUrole{w}{ }\DUrole{n}{ndarray}}\sphinxparamcomma \sphinxparam{\DUrole{n}{tight\_state\_errors}\DUrole{p}{:}\DUrole{w}{ }\DUrole{n}{ndarray}}}{}
\pysigstopsignatures
\sphinxAtStartPar
Tightly\sphinxhyphen{}coupled based GPS fusion.
A fusion method that uses tightly\sphinxhyphen{}coupled fusion for updating the INS
with given GPS estimates. This method uses a centralized Kalman filter
that integrates the estimated Pseudo Ranges and Doppler shift from the
GNSS receiver, along with the estimated position, velocity and attitude
from the INS. Must be used with a Kalman estimated state class.
\index{perform\_fusion() (qnav.gps.fusion.GpsTightFusion method)@\spxentry{perform\_fusion()}\spxextra{qnav.gps.fusion.GpsTightFusion method}}

\begin{fulllineitems}
\phantomsection\label{\detokenize{modules/gps/fusion:qnav.gps.fusion.GpsTightFusion.perform_fusion}}
\pysigstartsignatures
\pysiglinewithargsret{\sphinxbfcode{\sphinxupquote{perform\_fusion}}}{\sphinxparam{\DUrole{n}{estimated\_state}\DUrole{p}{:}\DUrole{w}{ }\DUrole{n}{{\hyperref[\detokenize{modules/estimation/state:qnav.estimation.state.KalmanEstimatedState}]{\sphinxcrossref{KalmanEstimatedState}}}}}}{{ $\rightarrow$ None}}
\pysigstopsignatures
\sphinxAtStartPar
Applies tightly\sphinxhyphen{}coupled fusion using GPS estimates and INS solution.
This method updates the given INS instance’s estimates by fusing them
with the given GPS estimates. The INS states are updated internally.
\begin{quote}\begin{description}
\sphinxlineitem{Parameters}
\sphinxAtStartPar
\sphinxstyleliteralstrong{\sphinxupquote{estimated\_state}} ({\hyperref[\detokenize{modules/estimation/state:qnav.estimation.state.KalmanEstimatedState}]{\sphinxcrossref{\sphinxstyleliteralemphasis{\sphinxupquote{KalmanEstimatedState}}}}}) \textendash{} The current estimated state.

\end{description}\end{quote}

\end{fulllineitems}


\end{fulllineitems}

\index{calculate\_cov\_mat() (in module qnav.gps.fusion)@\spxentry{calculate\_cov\_mat()}\spxextra{in module qnav.gps.fusion}}

\begin{fulllineitems}
\phantomsection\label{\detokenize{modules/gps/fusion:qnav.gps.fusion.calculate_cov_mat}}
\pysigstartsignatures
\pysiglinewithargsret{\sphinxcode{\sphinxupquote{qnav.gps.fusion.}}\sphinxbfcode{\sphinxupquote{calculate\_cov\_mat}}}{\sphinxparam{\DUrole{n}{satellites}\DUrole{p}{:}\DUrole{w}{ }\DUrole{n}{Collection\DUrole{p}{{[}}{\hyperref[\detokenize{modules/gps/satellite:qnav.gps.satellite.Satellite}]{\sphinxcrossref{Satellite}}}\DUrole{p}{{]}}}}\sphinxparamcomma \sphinxparam{\DUrole{n}{ant\_pos\_lla}\DUrole{p}{:}\DUrole{w}{ }\DUrole{n}{ndarray}}\sphinxparamcomma \sphinxparam{\DUrole{n}{sigma\_position}\DUrole{p}{:}\DUrole{w}{ }\DUrole{n}{float}\DUrole{w}{ }\DUrole{o}{=}\DUrole{w}{ }\DUrole{default_value}{50}}\sphinxparamcomma \sphinxparam{\DUrole{n}{sigma\_velocity}\DUrole{p}{:}\DUrole{w}{ }\DUrole{n}{float}\DUrole{w}{ }\DUrole{o}{=}\DUrole{w}{ }\DUrole{default_value}{0.3}}}{{ $\rightarrow$ tuple\DUrole{p}{{[}}ndarray\DUrole{p}{,}\DUrole{w}{ }ndarray\DUrole{p}{{]}}}}
\pysigstopsignatures
\sphinxAtStartPar
Produces the covariance matrices for loosely\sphinxhyphen{}coupled fusion.
This method produces and returns the covariance matrices for position
and velocity based on the current position. These are namely required
for initialising a loosely\sphinxhyphen{}coupled GPS fusion instance (if matrices
are not produced manually).
\begin{quote}\begin{description}
\sphinxlineitem{Parameters}\begin{itemize}
\item {} 
\sphinxAtStartPar
\sphinxstyleliteralstrong{\sphinxupquote{satellites}} (\sphinxstyleliteralemphasis{\sphinxupquote{Collection}}\sphinxstyleliteralemphasis{\sphinxupquote{{[}}}{\hyperref[\detokenize{modules/gps/satellite:qnav.gps.satellite.Satellite}]{\sphinxcrossref{\sphinxstyleliteralemphasis{\sphinxupquote{Satellite}}}}}\sphinxstyleliteralemphasis{\sphinxupquote{{]}}}) \textendash{} The GPS satellites involved in measurements.

\item {} 
\sphinxAtStartPar
\sphinxstyleliteralstrong{\sphinxupquote{ant\_pos\_lla}} (\sphinxstyleliteralemphasis{\sphinxupquote{numpy.array}}\sphinxstyleliteralemphasis{\sphinxupquote{ (}}\sphinxstyleliteralemphasis{\sphinxupquote{3\sphinxhyphen{}elements}}\sphinxstyleliteralemphasis{\sphinxupquote{)}}) \textendash{} The true antenna position (latitude, longitude
and altitude given in deg\sphinxhyphen{}deg\sphinxhyphen{}metres).

\item {} 
\sphinxAtStartPar
\sphinxstyleliteralstrong{\sphinxupquote{sigma\_position}} (\sphinxstyleliteralemphasis{\sphinxupquote{float}}) \textendash{} The pseudo\sphinxhyphen{}range position error (in metres).

\item {} 
\sphinxAtStartPar
\sphinxstyleliteralstrong{\sphinxupquote{sigma\_velocity}} (\sphinxstyleliteralemphasis{\sphinxupquote{float}}) \textendash{} The pseudo\sphinxhyphen{}range velocity error (in metres/sec).

\end{itemize}

\sphinxlineitem{Returns}
\sphinxAtStartPar
A tuple containing the constructed position and velocity
covariance matrices, respectively.

\sphinxlineitem{Return type}
\sphinxAtStartPar
tuple{[}np.ndarray, np.ndarray{]}

\end{description}\end{quote}

\end{fulllineitems}


\sphinxstepscope


\subsubsection{qnav.gps.satellite}
\label{\detokenize{modules/gps/satellite:module-qnav.gps.satellite}}\label{\detokenize{modules/gps/satellite:qnav-gps-satellite}}\label{\detokenize{modules/gps/satellite::doc}}\index{module@\spxentry{module}!qnav.gps.satellite@\spxentry{qnav.gps.satellite}}\index{qnav.gps.satellite@\spxentry{qnav.gps.satellite}!module@\spxentry{module}}

\paragraph{satellite.py}
\label{\detokenize{modules/gps/satellite:satellite-py}}\begin{quote}\begin{description}
\sphinxlineitem{summary}
\sphinxAtStartPar
All GPS satellite and solver functionalities.
This module is responsible for simulating up\sphinxhyphen{}to\sphinxhyphen{}date satellites and
supplying a solver to produce measurements (with configured errors
and noise). During a simulation these classes will be used for
acquiring GPS information.

\sphinxlineitem{authors}
\begin{DUlineblock}{0em}
\item[] Dr. Niall Moroney \sphinxhyphen{} \sphinxhref{mailto:niall.moroney17@imperial.ac.uk}{niall.moroney17@imperial.ac.uk}
\item[] Michael Wright \sphinxhyphen{} \sphinxhref{mailto:mjwright@liverpool.ac.uk}{mjwright@liverpool.ac.uk}
\end{DUlineblock}

\sphinxlineitem{version}
\sphinxAtStartPar
0.9.0 \sphinxhyphen{} Awaiting further implementations.

\end{description}\end{quote}
\index{satellite\_param\_fields (in module qnav.gps.satellite)@\spxentry{satellite\_param\_fields}\spxextra{in module qnav.gps.satellite}}

\begin{fulllineitems}
\phantomsection\label{\detokenize{modules/gps/satellite:qnav.gps.satellite.satellite_param_fields}}
\pysigstartsignatures
\pysigline{\sphinxcode{\sphinxupquote{qnav.gps.satellite.}}\sphinxbfcode{\sphinxupquote{satellite\_param\_fields}}\sphinxbfcode{\sphinxupquote{\DUrole{w}{ }\DUrole{p}{=}\DUrole{w}{ }(\textquotesingle{}type\textquotesingle{}, \textquotesingle{}id\textquotesingle{}, \textquotesingle{}year\textquotesingle{}, \textquotesingle{}month\textquotesingle{}, \textquotesingle{}day\textquotesingle{}, \textquotesingle{}hour\textquotesingle{}, \textquotesingle{}minute\textquotesingle{}, \textquotesingle{}second\textquotesingle{}, \textquotesingle{}clock\_bias\textquotesingle{}, \textquotesingle{}drift\_freq\_bias\textquotesingle{}, \textquotesingle{}drift\_rate\_message\_time\textquotesingle{}, \textquotesingle{}t\_0e\textquotesingle{}, \textquotesingle{}prn\textquotesingle{}, \textquotesingle{}delta\_n\textquotesingle{}, \textquotesingle{}i\_dot\textquotesingle{}, \textquotesingle{}omega\_dot\textquotesingle{}, \textquotesingle{}omega\_0\textquotesingle{}, \textquotesingle{}c\_us\textquotesingle{}, \textquotesingle{}c\_uc\textquotesingle{}, \textquotesingle{}c\_is\textquotesingle{}, \textquotesingle{}c\_ic\textquotesingle{}, \textquotesingle{}c\_rs\textquotesingle{}, \textquotesingle{}c\_rc\textquotesingle{}, \textquotesingle{}sqrt\_a\textquotesingle{}, \textquotesingle{}e\textquotesingle{}, \textquotesingle{}i\_0\textquotesingle{}, \textquotesingle{}w\textquotesingle{}, \textquotesingle{}m\_bar\_0\textquotesingle{})}}}
\pysigstopsignatures
\sphinxAtStartPar
The field names for satellite parameters.

\end{fulllineitems}

\index{SatelliteParams (class in qnav.gps.satellite)@\spxentry{SatelliteParams}\spxextra{class in qnav.gps.satellite}}

\begin{fulllineitems}
\phantomsection\label{\detokenize{modules/gps/satellite:qnav.gps.satellite.SatelliteParams}}
\pysigstartsignatures
\pysiglinewithargsret{\sphinxbfcode{\sphinxupquote{class\DUrole{w}{ }}}\sphinxcode{\sphinxupquote{qnav.gps.satellite.}}\sphinxbfcode{\sphinxupquote{SatelliteParams}}}{\sphinxparam{\DUrole{n}{type}}\sphinxparamcomma \sphinxparam{\DUrole{n}{id}}\sphinxparamcomma \sphinxparam{\DUrole{n}{year}}\sphinxparamcomma \sphinxparam{\DUrole{n}{month}}\sphinxparamcomma \sphinxparam{\DUrole{n}{day}}\sphinxparamcomma \sphinxparam{\DUrole{n}{hour}}\sphinxparamcomma \sphinxparam{\DUrole{n}{minute}}\sphinxparamcomma \sphinxparam{\DUrole{n}{second}}\sphinxparamcomma \sphinxparam{\DUrole{n}{clock\_bias}}\sphinxparamcomma \sphinxparam{\DUrole{n}{drift\_freq\_bias}}\sphinxparamcomma \sphinxparam{\DUrole{n}{drift\_rate\_message\_time}}\sphinxparamcomma \sphinxparam{\DUrole{n}{t\_0e}}\sphinxparamcomma \sphinxparam{\DUrole{n}{prn}}\sphinxparamcomma \sphinxparam{\DUrole{n}{delta\_n}}\sphinxparamcomma \sphinxparam{\DUrole{n}{i\_dot}}\sphinxparamcomma \sphinxparam{\DUrole{n}{omega\_dot}}\sphinxparamcomma \sphinxparam{\DUrole{n}{omega\_0}}\sphinxparamcomma \sphinxparam{\DUrole{n}{c\_us}}\sphinxparamcomma \sphinxparam{\DUrole{n}{c\_uc}}\sphinxparamcomma \sphinxparam{\DUrole{n}{c\_is}}\sphinxparamcomma \sphinxparam{\DUrole{n}{c\_ic}}\sphinxparamcomma \sphinxparam{\DUrole{n}{c\_rs}}\sphinxparamcomma \sphinxparam{\DUrole{n}{c\_rc}}\sphinxparamcomma \sphinxparam{\DUrole{n}{sqrt\_a}}\sphinxparamcomma \sphinxparam{\DUrole{n}{e}}\sphinxparamcomma \sphinxparam{\DUrole{n}{i\_0}}\sphinxparamcomma \sphinxparam{\DUrole{n}{w}}\sphinxparamcomma \sphinxparam{\DUrole{n}{m\_bar\_0}}}{}
\pysigstopsignatures
\sphinxAtStartPar
A collection of satellite parameters needed to initialise a Satellite instance.
\index{c\_ic (qnav.gps.satellite.SatelliteParams attribute)@\spxentry{c\_ic}\spxextra{qnav.gps.satellite.SatelliteParams attribute}}

\begin{fulllineitems}
\phantomsection\label{\detokenize{modules/gps/satellite:qnav.gps.satellite.SatelliteParams.c_ic}}
\pysigstartsignatures
\pysigline{\sphinxbfcode{\sphinxupquote{c\_ic}}}
\pysigstopsignatures
\sphinxAtStartPar
Alias for field number 20

\end{fulllineitems}

\index{c\_is (qnav.gps.satellite.SatelliteParams attribute)@\spxentry{c\_is}\spxextra{qnav.gps.satellite.SatelliteParams attribute}}

\begin{fulllineitems}
\phantomsection\label{\detokenize{modules/gps/satellite:qnav.gps.satellite.SatelliteParams.c_is}}
\pysigstartsignatures
\pysigline{\sphinxbfcode{\sphinxupquote{c\_is}}}
\pysigstopsignatures
\sphinxAtStartPar
Alias for field number 19

\end{fulllineitems}

\index{c\_rc (qnav.gps.satellite.SatelliteParams attribute)@\spxentry{c\_rc}\spxextra{qnav.gps.satellite.SatelliteParams attribute}}

\begin{fulllineitems}
\phantomsection\label{\detokenize{modules/gps/satellite:qnav.gps.satellite.SatelliteParams.c_rc}}
\pysigstartsignatures
\pysigline{\sphinxbfcode{\sphinxupquote{c\_rc}}}
\pysigstopsignatures
\sphinxAtStartPar
Alias for field number 22

\end{fulllineitems}

\index{c\_rs (qnav.gps.satellite.SatelliteParams attribute)@\spxentry{c\_rs}\spxextra{qnav.gps.satellite.SatelliteParams attribute}}

\begin{fulllineitems}
\phantomsection\label{\detokenize{modules/gps/satellite:qnav.gps.satellite.SatelliteParams.c_rs}}
\pysigstartsignatures
\pysigline{\sphinxbfcode{\sphinxupquote{c\_rs}}}
\pysigstopsignatures
\sphinxAtStartPar
Alias for field number 21

\end{fulllineitems}

\index{c\_uc (qnav.gps.satellite.SatelliteParams attribute)@\spxentry{c\_uc}\spxextra{qnav.gps.satellite.SatelliteParams attribute}}

\begin{fulllineitems}
\phantomsection\label{\detokenize{modules/gps/satellite:qnav.gps.satellite.SatelliteParams.c_uc}}
\pysigstartsignatures
\pysigline{\sphinxbfcode{\sphinxupquote{c\_uc}}}
\pysigstopsignatures
\sphinxAtStartPar
Alias for field number 18

\end{fulllineitems}

\index{c\_us (qnav.gps.satellite.SatelliteParams attribute)@\spxentry{c\_us}\spxextra{qnav.gps.satellite.SatelliteParams attribute}}

\begin{fulllineitems}
\phantomsection\label{\detokenize{modules/gps/satellite:qnav.gps.satellite.SatelliteParams.c_us}}
\pysigstartsignatures
\pysigline{\sphinxbfcode{\sphinxupquote{c\_us}}}
\pysigstopsignatures
\sphinxAtStartPar
Alias for field number 17

\end{fulllineitems}

\index{clock\_bias (qnav.gps.satellite.SatelliteParams attribute)@\spxentry{clock\_bias}\spxextra{qnav.gps.satellite.SatelliteParams attribute}}

\begin{fulllineitems}
\phantomsection\label{\detokenize{modules/gps/satellite:qnav.gps.satellite.SatelliteParams.clock_bias}}
\pysigstartsignatures
\pysigline{\sphinxbfcode{\sphinxupquote{clock\_bias}}}
\pysigstopsignatures
\sphinxAtStartPar
Alias for field number 8

\end{fulllineitems}

\index{day (qnav.gps.satellite.SatelliteParams attribute)@\spxentry{day}\spxextra{qnav.gps.satellite.SatelliteParams attribute}}

\begin{fulllineitems}
\phantomsection\label{\detokenize{modules/gps/satellite:qnav.gps.satellite.SatelliteParams.day}}
\pysigstartsignatures
\pysigline{\sphinxbfcode{\sphinxupquote{day}}}
\pysigstopsignatures
\sphinxAtStartPar
Alias for field number 4

\end{fulllineitems}

\index{delta\_n (qnav.gps.satellite.SatelliteParams attribute)@\spxentry{delta\_n}\spxextra{qnav.gps.satellite.SatelliteParams attribute}}

\begin{fulllineitems}
\phantomsection\label{\detokenize{modules/gps/satellite:qnav.gps.satellite.SatelliteParams.delta_n}}
\pysigstartsignatures
\pysigline{\sphinxbfcode{\sphinxupquote{delta\_n}}}
\pysigstopsignatures
\sphinxAtStartPar
Alias for field number 13

\end{fulllineitems}

\index{drift\_freq\_bias (qnav.gps.satellite.SatelliteParams attribute)@\spxentry{drift\_freq\_bias}\spxextra{qnav.gps.satellite.SatelliteParams attribute}}

\begin{fulllineitems}
\phantomsection\label{\detokenize{modules/gps/satellite:qnav.gps.satellite.SatelliteParams.drift_freq_bias}}
\pysigstartsignatures
\pysigline{\sphinxbfcode{\sphinxupquote{drift\_freq\_bias}}}
\pysigstopsignatures
\sphinxAtStartPar
Alias for field number 9

\end{fulllineitems}

\index{drift\_rate\_message\_time (qnav.gps.satellite.SatelliteParams attribute)@\spxentry{drift\_rate\_message\_time}\spxextra{qnav.gps.satellite.SatelliteParams attribute}}

\begin{fulllineitems}
\phantomsection\label{\detokenize{modules/gps/satellite:qnav.gps.satellite.SatelliteParams.drift_rate_message_time}}
\pysigstartsignatures
\pysigline{\sphinxbfcode{\sphinxupquote{drift\_rate\_message\_time}}}
\pysigstopsignatures
\sphinxAtStartPar
Alias for field number 10

\end{fulllineitems}

\index{e (qnav.gps.satellite.SatelliteParams attribute)@\spxentry{e}\spxextra{qnav.gps.satellite.SatelliteParams attribute}}

\begin{fulllineitems}
\phantomsection\label{\detokenize{modules/gps/satellite:qnav.gps.satellite.SatelliteParams.e}}
\pysigstartsignatures
\pysigline{\sphinxbfcode{\sphinxupquote{e}}}
\pysigstopsignatures
\sphinxAtStartPar
Alias for field number 24

\end{fulllineitems}

\index{hour (qnav.gps.satellite.SatelliteParams attribute)@\spxentry{hour}\spxextra{qnav.gps.satellite.SatelliteParams attribute}}

\begin{fulllineitems}
\phantomsection\label{\detokenize{modules/gps/satellite:qnav.gps.satellite.SatelliteParams.hour}}
\pysigstartsignatures
\pysigline{\sphinxbfcode{\sphinxupquote{hour}}}
\pysigstopsignatures
\sphinxAtStartPar
Alias for field number 5

\end{fulllineitems}

\index{i\_0 (qnav.gps.satellite.SatelliteParams attribute)@\spxentry{i\_0}\spxextra{qnav.gps.satellite.SatelliteParams attribute}}

\begin{fulllineitems}
\phantomsection\label{\detokenize{modules/gps/satellite:qnav.gps.satellite.SatelliteParams.i_0}}
\pysigstartsignatures
\pysigline{\sphinxbfcode{\sphinxupquote{i\_0}}}
\pysigstopsignatures
\sphinxAtStartPar
Alias for field number 25

\end{fulllineitems}

\index{i\_dot (qnav.gps.satellite.SatelliteParams attribute)@\spxentry{i\_dot}\spxextra{qnav.gps.satellite.SatelliteParams attribute}}

\begin{fulllineitems}
\phantomsection\label{\detokenize{modules/gps/satellite:qnav.gps.satellite.SatelliteParams.i_dot}}
\pysigstartsignatures
\pysigline{\sphinxbfcode{\sphinxupquote{i\_dot}}}
\pysigstopsignatures
\sphinxAtStartPar
Alias for field number 14

\end{fulllineitems}

\index{id (qnav.gps.satellite.SatelliteParams attribute)@\spxentry{id}\spxextra{qnav.gps.satellite.SatelliteParams attribute}}

\begin{fulllineitems}
\phantomsection\label{\detokenize{modules/gps/satellite:qnav.gps.satellite.SatelliteParams.id}}
\pysigstartsignatures
\pysigline{\sphinxbfcode{\sphinxupquote{id}}}
\pysigstopsignatures
\sphinxAtStartPar
Alias for field number 1

\end{fulllineitems}

\index{m\_bar\_0 (qnav.gps.satellite.SatelliteParams attribute)@\spxentry{m\_bar\_0}\spxextra{qnav.gps.satellite.SatelliteParams attribute}}

\begin{fulllineitems}
\phantomsection\label{\detokenize{modules/gps/satellite:qnav.gps.satellite.SatelliteParams.m_bar_0}}
\pysigstartsignatures
\pysigline{\sphinxbfcode{\sphinxupquote{m\_bar\_0}}}
\pysigstopsignatures
\sphinxAtStartPar
Alias for field number 27

\end{fulllineitems}

\index{minute (qnav.gps.satellite.SatelliteParams attribute)@\spxentry{minute}\spxextra{qnav.gps.satellite.SatelliteParams attribute}}

\begin{fulllineitems}
\phantomsection\label{\detokenize{modules/gps/satellite:qnav.gps.satellite.SatelliteParams.minute}}
\pysigstartsignatures
\pysigline{\sphinxbfcode{\sphinxupquote{minute}}}
\pysigstopsignatures
\sphinxAtStartPar
Alias for field number 6

\end{fulllineitems}

\index{month (qnav.gps.satellite.SatelliteParams attribute)@\spxentry{month}\spxextra{qnav.gps.satellite.SatelliteParams attribute}}

\begin{fulllineitems}
\phantomsection\label{\detokenize{modules/gps/satellite:qnav.gps.satellite.SatelliteParams.month}}
\pysigstartsignatures
\pysigline{\sphinxbfcode{\sphinxupquote{month}}}
\pysigstopsignatures
\sphinxAtStartPar
Alias for field number 3

\end{fulllineitems}

\index{omega\_0 (qnav.gps.satellite.SatelliteParams attribute)@\spxentry{omega\_0}\spxextra{qnav.gps.satellite.SatelliteParams attribute}}

\begin{fulllineitems}
\phantomsection\label{\detokenize{modules/gps/satellite:qnav.gps.satellite.SatelliteParams.omega_0}}
\pysigstartsignatures
\pysigline{\sphinxbfcode{\sphinxupquote{omega\_0}}}
\pysigstopsignatures
\sphinxAtStartPar
Alias for field number 16

\end{fulllineitems}

\index{omega\_dot (qnav.gps.satellite.SatelliteParams attribute)@\spxentry{omega\_dot}\spxextra{qnav.gps.satellite.SatelliteParams attribute}}

\begin{fulllineitems}
\phantomsection\label{\detokenize{modules/gps/satellite:qnav.gps.satellite.SatelliteParams.omega_dot}}
\pysigstartsignatures
\pysigline{\sphinxbfcode{\sphinxupquote{omega\_dot}}}
\pysigstopsignatures
\sphinxAtStartPar
Alias for field number 15

\end{fulllineitems}

\index{prn (qnav.gps.satellite.SatelliteParams attribute)@\spxentry{prn}\spxextra{qnav.gps.satellite.SatelliteParams attribute}}

\begin{fulllineitems}
\phantomsection\label{\detokenize{modules/gps/satellite:qnav.gps.satellite.SatelliteParams.prn}}
\pysigstartsignatures
\pysigline{\sphinxbfcode{\sphinxupquote{prn}}}
\pysigstopsignatures
\sphinxAtStartPar
Alias for field number 12

\end{fulllineitems}

\index{second (qnav.gps.satellite.SatelliteParams attribute)@\spxentry{second}\spxextra{qnav.gps.satellite.SatelliteParams attribute}}

\begin{fulllineitems}
\phantomsection\label{\detokenize{modules/gps/satellite:qnav.gps.satellite.SatelliteParams.second}}
\pysigstartsignatures
\pysigline{\sphinxbfcode{\sphinxupquote{second}}}
\pysigstopsignatures
\sphinxAtStartPar
Alias for field number 7

\end{fulllineitems}

\index{sqrt\_a (qnav.gps.satellite.SatelliteParams attribute)@\spxentry{sqrt\_a}\spxextra{qnav.gps.satellite.SatelliteParams attribute}}

\begin{fulllineitems}
\phantomsection\label{\detokenize{modules/gps/satellite:qnav.gps.satellite.SatelliteParams.sqrt_a}}
\pysigstartsignatures
\pysigline{\sphinxbfcode{\sphinxupquote{sqrt\_a}}}
\pysigstopsignatures
\sphinxAtStartPar
Alias for field number 23

\end{fulllineitems}

\index{t\_0e (qnav.gps.satellite.SatelliteParams attribute)@\spxentry{t\_0e}\spxextra{qnav.gps.satellite.SatelliteParams attribute}}

\begin{fulllineitems}
\phantomsection\label{\detokenize{modules/gps/satellite:qnav.gps.satellite.SatelliteParams.t_0e}}
\pysigstartsignatures
\pysigline{\sphinxbfcode{\sphinxupquote{t\_0e}}}
\pysigstopsignatures
\sphinxAtStartPar
Alias for field number 11

\end{fulllineitems}

\index{type (qnav.gps.satellite.SatelliteParams attribute)@\spxentry{type}\spxextra{qnav.gps.satellite.SatelliteParams attribute}}

\begin{fulllineitems}
\phantomsection\label{\detokenize{modules/gps/satellite:qnav.gps.satellite.SatelliteParams.type}}
\pysigstartsignatures
\pysigline{\sphinxbfcode{\sphinxupquote{type}}}
\pysigstopsignatures
\sphinxAtStartPar
Alias for field number 0

\end{fulllineitems}

\index{w (qnav.gps.satellite.SatelliteParams attribute)@\spxentry{w}\spxextra{qnav.gps.satellite.SatelliteParams attribute}}

\begin{fulllineitems}
\phantomsection\label{\detokenize{modules/gps/satellite:qnav.gps.satellite.SatelliteParams.w}}
\pysigstartsignatures
\pysigline{\sphinxbfcode{\sphinxupquote{w}}}
\pysigstopsignatures
\sphinxAtStartPar
Alias for field number 26

\end{fulllineitems}

\index{year (qnav.gps.satellite.SatelliteParams attribute)@\spxentry{year}\spxextra{qnav.gps.satellite.SatelliteParams attribute}}

\begin{fulllineitems}
\phantomsection\label{\detokenize{modules/gps/satellite:qnav.gps.satellite.SatelliteParams.year}}
\pysigstartsignatures
\pysigline{\sphinxbfcode{\sphinxupquote{year}}}
\pysigstopsignatures
\sphinxAtStartPar
Alias for field number 2

\end{fulllineitems}


\end{fulllineitems}

\index{Satellite (class in qnav.gps.satellite)@\spxentry{Satellite}\spxextra{class in qnav.gps.satellite}}

\begin{fulllineitems}
\phantomsection\label{\detokenize{modules/gps/satellite:qnav.gps.satellite.Satellite}}
\pysigstartsignatures
\pysiglinewithargsret{\sphinxbfcode{\sphinxupquote{class\DUrole{w}{ }}}\sphinxcode{\sphinxupquote{qnav.gps.satellite.}}\sphinxbfcode{\sphinxupquote{Satellite}}}{\sphinxparam{\DUrole{n}{params}\DUrole{p}{:}\DUrole{w}{ }\DUrole{n}{{\hyperref[\detokenize{modules/gps/satellite:qnav.gps.satellite.SatelliteParams}]{\sphinxcrossref{SatelliteParams}}}}}\sphinxparamcomma \sphinxparam{\DUrole{n}{gps\_time}\DUrole{p}{:}\DUrole{w}{ }\DUrole{n}{datetime}\DUrole{w}{ }\DUrole{o}{=}\DUrole{w}{ }\DUrole{default_value}{None}}\sphinxparamcomma \sphinxparam{\DUrole{n}{pseudo\_range\_noise}\DUrole{p}{:}\DUrole{w}{ }\DUrole{n}{float}\DUrole{w}{ }\DUrole{o}{=}\DUrole{w}{ }\DUrole{default_value}{0}}\sphinxparamcomma \sphinxparam{\DUrole{n}{pseudo\_range\_rate\_noise}\DUrole{p}{:}\DUrole{w}{ }\DUrole{n}{float}\DUrole{w}{ }\DUrole{o}{=}\DUrole{w}{ }\DUrole{default_value}{0}}\sphinxparamcomma \sphinxparam{\DUrole{n}{rand\_seed}\DUrole{p}{:}\DUrole{w}{ }\DUrole{n}{int}\DUrole{w}{ }\DUrole{o}{=}\DUrole{w}{ }\DUrole{default_value}{None}}}{}
\pysigstopsignatures
\sphinxAtStartPar
Represents a GPS satellite and its current states/properties.
This class is responsible for representing information about a GPS
satellite. Methods are supplied for updating these states and allowing
for measurement information to be acquired. During simulations a
collection of satellite instances will be used for generating a GPS signal
that will be used to improve navigation accuracy.
\index{update() (qnav.gps.satellite.Satellite method)@\spxentry{update()}\spxextra{qnav.gps.satellite.Satellite method}}

\begin{fulllineitems}
\phantomsection\label{\detokenize{modules/gps/satellite:qnav.gps.satellite.Satellite.update}}
\pysigstartsignatures
\pysiglinewithargsret{\sphinxbfcode{\sphinxupquote{update}}}{\sphinxparam{\DUrole{n}{dt}\DUrole{p}{:}\DUrole{w}{ }\DUrole{n}{float}}}{}
\pysigstopsignatures
\sphinxAtStartPar
Updates the position, velocity and errors given the time difference.
This method updates the state of the satellite by a given amount of
time (in seconds). The updated position and velocity of the satellite
are calculated along with the new error states for generating noisy
measurements.
\begin{quote}\begin{description}
\sphinxlineitem{Parameters}
\sphinxAtStartPar
\sphinxstyleliteralstrong{\sphinxupquote{dt}} (\sphinxstyleliteralemphasis{\sphinxupquote{float}}) \textendash{} The time difference since the last update was requested.

\end{description}\end{quote}

\end{fulllineitems}

\index{get\_location\_azimuth\_elevation() (qnav.gps.satellite.Satellite method)@\spxentry{get\_location\_azimuth\_elevation()}\spxextra{qnav.gps.satellite.Satellite method}}

\begin{fulllineitems}
\phantomsection\label{\detokenize{modules/gps/satellite:qnav.gps.satellite.Satellite.get_location_azimuth_elevation}}
\pysigstartsignatures
\pysiglinewithargsret{\sphinxbfcode{\sphinxupquote{get\_location\_azimuth\_elevation}}}{\sphinxparam{\DUrole{n}{ant\_pos\_lla}\DUrole{p}{:}\DUrole{w}{ }\DUrole{n}{ndarray}}}{{ $\rightarrow$ tuple\DUrole{p}{{[}}float\DUrole{p}{,}\DUrole{w}{ }float\DUrole{p}{{]}}}}
\pysigstopsignatures
\sphinxAtStartPar
Get the relative location of the satellite from the current position,
in azimuth and elevation (deg\sphinxhyphen{}deg).
\begin{quote}\begin{description}
\sphinxlineitem{Parameters}
\sphinxAtStartPar
\sphinxstyleliteralstrong{\sphinxupquote{ant\_pos\_lla}} (\sphinxstyleliteralemphasis{\sphinxupquote{numpy.array}}\sphinxstyleliteralemphasis{\sphinxupquote{ (}}\sphinxstyleliteralemphasis{\sphinxupquote{3\sphinxhyphen{}elements}}\sphinxstyleliteralemphasis{\sphinxupquote{)}}) \textendash{} The true antenna position (latitude, longitude
and altitude given in deg\sphinxhyphen{}deg\sphinxhyphen{}metres).

\sphinxlineitem{Returns}
\sphinxAtStartPar
Returns the corresponding azimuth and elevation (both in
degrees).

\sphinxlineitem{Return type}
\sphinxAtStartPar
tuple{[}float, float{]}

\end{description}\end{quote}

\end{fulllineitems}

\index{get\_measurements() (qnav.gps.satellite.Satellite method)@\spxentry{get\_measurements()}\spxextra{qnav.gps.satellite.Satellite method}}

\begin{fulllineitems}
\phantomsection\label{\detokenize{modules/gps/satellite:qnav.gps.satellite.Satellite.get_measurements}}
\pysigstartsignatures
\pysiglinewithargsret{\sphinxbfcode{\sphinxupquote{get\_measurements}}}{\sphinxparam{\DUrole{n}{ant\_pos\_lla}\DUrole{p}{:}\DUrole{w}{ }\DUrole{n}{ndarray}}\sphinxparamcomma \sphinxparam{\DUrole{n}{ant\_att\_deg}\DUrole{p}{:}\DUrole{w}{ }\DUrole{n}{ndarray}}\sphinxparamcomma \sphinxparam{\DUrole{n}{ant\_vel\_body\_axes}\DUrole{p}{:}\DUrole{w}{ }\DUrole{n}{ndarray}}}{{ $\rightarrow$ dict}}
\pysigstopsignatures
\sphinxAtStartPar
Calculates and returns the noisy pseudo ranges, position, velocity
and time information for the satellite.
\begin{quote}\begin{description}
\sphinxlineitem{Parameters}\begin{itemize}
\item {} 
\sphinxAtStartPar
\sphinxstyleliteralstrong{\sphinxupquote{ant\_pos\_lla}} (\sphinxstyleliteralemphasis{\sphinxupquote{numpy.array}}\sphinxstyleliteralemphasis{\sphinxupquote{ (}}\sphinxstyleliteralemphasis{\sphinxupquote{3\sphinxhyphen{}elements}}\sphinxstyleliteralemphasis{\sphinxupquote{)}}) \textendash{} The true antenna position (latitude, longitude
and altitude given in deg\sphinxhyphen{}deg\sphinxhyphen{}metres).

\item {} 
\sphinxAtStartPar
\sphinxstyleliteralstrong{\sphinxupquote{ant\_att\_deg}} (\sphinxstyleliteralemphasis{\sphinxupquote{numpy.array}}\sphinxstyleliteralemphasis{\sphinxupquote{ (}}\sphinxstyleliteralemphasis{\sphinxupquote{3\sphinxhyphen{}elements}}\sphinxstyleliteralemphasis{\sphinxupquote{)}}) \textendash{} The true antenna attitude (heading, pitch and roll
in degrees).

\item {} 
\sphinxAtStartPar
\sphinxstyleliteralstrong{\sphinxupquote{ant\_vel\_body\_axes}} (\sphinxstyleliteralemphasis{\sphinxupquote{numpy.array}}\sphinxstyleliteralemphasis{\sphinxupquote{ (}}\sphinxstyleliteralemphasis{\sphinxupquote{3\sphinxhyphen{}elements}}\sphinxstyleliteralemphasis{\sphinxupquote{)}}) \textendash{} The true antenna velocity (in body axes in
metres/second).

\end{itemize}

\sphinxlineitem{Returns}
\sphinxAtStartPar
The relevant measurement information in a dict.

\sphinxlineitem{Return type}
\sphinxAtStartPar
dict

\end{description}\end{quote}

\end{fulllineitems}

\index{sat\_id (qnav.gps.satellite.Satellite property)@\spxentry{sat\_id}\spxextra{qnav.gps.satellite.Satellite property}}

\begin{fulllineitems}
\phantomsection\label{\detokenize{modules/gps/satellite:qnav.gps.satellite.Satellite.sat_id}}
\pysigstartsignatures
\pysigline{\sphinxbfcode{\sphinxupquote{property\DUrole{w}{ }}}\sphinxbfcode{\sphinxupquote{sat\_id}}\sphinxbfcode{\sphinxupquote{\DUrole{p}{:}\DUrole{w}{ }str}}}
\pysigstopsignatures
\sphinxAtStartPar
Returns the full satellites ID from given parameters.
\begin{quote}\begin{description}
\sphinxlineitem{Returns}
\sphinxAtStartPar


\sphinxlineitem{Return type}
\sphinxAtStartPar
str

\end{description}\end{quote}

\end{fulllineitems}

\index{sat\_time (qnav.gps.satellite.Satellite property)@\spxentry{sat\_time}\spxextra{qnav.gps.satellite.Satellite property}}

\begin{fulllineitems}
\phantomsection\label{\detokenize{modules/gps/satellite:qnav.gps.satellite.Satellite.sat_time}}
\pysigstartsignatures
\pysigline{\sphinxbfcode{\sphinxupquote{property\DUrole{w}{ }}}\sphinxbfcode{\sphinxupquote{sat\_time}}\sphinxbfcode{\sphinxupquote{\DUrole{p}{:}\DUrole{w}{ }datetime}}}
\pysigstopsignatures
\sphinxAtStartPar
The satellite’s current clock time.
\begin{quote}\begin{description}
\sphinxlineitem{Returns}
\sphinxAtStartPar
Current satellite time.

\sphinxlineitem{Return type}
\sphinxAtStartPar
datetime.datetime

\end{description}\end{quote}

\end{fulllineitems}

\index{position (qnav.gps.satellite.Satellite property)@\spxentry{position}\spxextra{qnav.gps.satellite.Satellite property}}

\begin{fulllineitems}
\phantomsection\label{\detokenize{modules/gps/satellite:qnav.gps.satellite.Satellite.position}}
\pysigstartsignatures
\pysigline{\sphinxbfcode{\sphinxupquote{property\DUrole{w}{ }}}\sphinxbfcode{\sphinxupquote{position}}\sphinxbfcode{\sphinxupquote{\DUrole{p}{:}\DUrole{w}{ }ndarray}}}
\pysigstopsignatures
\sphinxAtStartPar
The satellite’s ECEF position since last update.
\begin{quote}\begin{description}
\sphinxlineitem{Returns}
\sphinxAtStartPar
Current ECEF position.

\sphinxlineitem{Return type}
\sphinxAtStartPar
numpy.ndarray

\end{description}\end{quote}

\end{fulllineitems}

\index{velocity (qnav.gps.satellite.Satellite property)@\spxentry{velocity}\spxextra{qnav.gps.satellite.Satellite property}}

\begin{fulllineitems}
\phantomsection\label{\detokenize{modules/gps/satellite:qnav.gps.satellite.Satellite.velocity}}
\pysigstartsignatures
\pysigline{\sphinxbfcode{\sphinxupquote{property\DUrole{w}{ }}}\sphinxbfcode{\sphinxupquote{velocity}}\sphinxbfcode{\sphinxupquote{\DUrole{p}{:}\DUrole{w}{ }ndarray}}}
\pysigstopsignatures
\sphinxAtStartPar
The satellite’s ECEF velocity since last update.
\begin{quote}\begin{description}
\sphinxlineitem{Returns}
\sphinxAtStartPar
Current ECEF velocity.

\sphinxlineitem{Return type}
\sphinxAtStartPar
numpy.ndarray

\end{description}\end{quote}

\end{fulllineitems}

\index{pseudo\_range\_noise (qnav.gps.satellite.Satellite property)@\spxentry{pseudo\_range\_noise}\spxextra{qnav.gps.satellite.Satellite property}}

\begin{fulllineitems}
\phantomsection\label{\detokenize{modules/gps/satellite:qnav.gps.satellite.Satellite.pseudo_range_noise}}
\pysigstartsignatures
\pysigline{\sphinxbfcode{\sphinxupquote{property\DUrole{w}{ }}}\sphinxbfcode{\sphinxupquote{pseudo\_range\_noise}}\sphinxbfcode{\sphinxupquote{\DUrole{p}{:}\DUrole{w}{ }float}}}
\pysigstopsignatures
\sphinxAtStartPar
Returns the satellite’s initial pseudo range noise.
\begin{quote}\begin{description}
\sphinxlineitem{Returns}
\sphinxAtStartPar
Pseudo range noise.

\sphinxlineitem{Return type}
\sphinxAtStartPar
float

\end{description}\end{quote}

\end{fulllineitems}

\index{pseudo\_range\_rate\_noise (qnav.gps.satellite.Satellite property)@\spxentry{pseudo\_range\_rate\_noise}\spxextra{qnav.gps.satellite.Satellite property}}

\begin{fulllineitems}
\phantomsection\label{\detokenize{modules/gps/satellite:qnav.gps.satellite.Satellite.pseudo_range_rate_noise}}
\pysigstartsignatures
\pysigline{\sphinxbfcode{\sphinxupquote{property\DUrole{w}{ }}}\sphinxbfcode{\sphinxupquote{pseudo\_range\_rate\_noise}}\sphinxbfcode{\sphinxupquote{\DUrole{p}{:}\DUrole{w}{ }float}}}
\pysigstopsignatures
\sphinxAtStartPar
Returns the satellite’s initial pseudo range rate noise.
\begin{quote}\begin{description}
\sphinxlineitem{Returns}
\sphinxAtStartPar
Pseudo range rate noise.

\sphinxlineitem{Return type}
\sphinxAtStartPar
float

\end{description}\end{quote}

\end{fulllineitems}


\end{fulllineitems}


\sphinxstepscope


\subsubsection{qnav.gps.sensor}
\label{\detokenize{modules/gps/sensor:module-qnav.gps.sensor}}\label{\detokenize{modules/gps/sensor:qnav-gps-sensor}}\label{\detokenize{modules/gps/sensor::doc}}\index{module@\spxentry{module}!qnav.gps.sensor@\spxentry{qnav.gps.sensor}}\index{qnav.gps.sensor@\spxentry{qnav.gps.sensor}!module@\spxentry{module}}

\paragraph{gps/sensor.py}
\label{\detokenize{modules/gps/sensor:gps-sensor-py}}\begin{quote}\begin{description}
\sphinxlineitem{summary}
\sphinxAtStartPar
GPS sensor activities used for collecting measurements from satellites.
This module provides sensor classes that are responsible for obtaining
and returning GPS measurements from selected satellites. These
measurements can then be passed to the appropriate GPS fusion method.

\sphinxlineitem{authors}
\begin{DUlineblock}{0em}
\item[] Kallum O’Hara \sphinxhyphen{} \sphinxhref{mailto:sgkohara@liverpool.ac.uk}{sgkohara@liverpool.ac.uk}
\item[] Michael Wright \sphinxhyphen{} \sphinxhref{mailto:mjwright@liverpool.ac.uk}{mjwright@liverpool.ac.uk}
\end{DUlineblock}

\sphinxlineitem{version}
\sphinxAtStartPar
1.0.0 \sphinxhyphen{} First implementation of this module.

\end{description}\end{quote}
\index{GpsSensor (class in qnav.gps.sensor)@\spxentry{GpsSensor}\spxextra{class in qnav.gps.sensor}}

\begin{fulllineitems}
\phantomsection\label{\detokenize{modules/gps/sensor:qnav.gps.sensor.GpsSensor}}
\pysigstartsignatures
\pysiglinewithargsret{\sphinxbfcode{\sphinxupquote{class\DUrole{w}{ }}}\sphinxcode{\sphinxupquote{qnav.gps.sensor.}}\sphinxbfcode{\sphinxupquote{GpsSensor}}}{\sphinxparam{\DUrole{n}{frequency}\DUrole{p}{:}\DUrole{w}{ }\DUrole{n}{float}}\sphinxparamcomma \sphinxparam{\DUrole{n}{satellites}\DUrole{p}{:}\DUrole{w}{ }\DUrole{n}{Collection\DUrole{p}{{[}}{\hyperref[\detokenize{modules/gps/satellite:qnav.gps.satellite.Satellite}]{\sphinxcrossref{Satellite}}}\DUrole{p}{{]}}}}\sphinxparamcomma \sphinxparam{\DUrole{n}{gps\_zones}\DUrole{p}{:}\DUrole{w}{ }\DUrole{n}{{\hyperref[\detokenize{modules/gps/zones:qnav.gps.zones.FlaggedPolygonGroup}]{\sphinxcrossref{FlaggedPolygonGroup}}}}\DUrole{w}{ }\DUrole{o}{=}\DUrole{w}{ }\DUrole{default_value}{None}}\sphinxparamcomma \sphinxparam{\DUrole{n}{start\_time}\DUrole{p}{:}\DUrole{w}{ }\DUrole{n}{float}\DUrole{w}{ }\DUrole{o}{=}\DUrole{w}{ }\DUrole{default_value}{0.0}}\sphinxparamcomma \sphinxparam{\DUrole{n}{rand\_seed}\DUrole{p}{:}\DUrole{w}{ }\DUrole{n}{int}\DUrole{w}{ }\DUrole{o}{=}\DUrole{w}{ }\DUrole{default_value}{None}}}{}
\pysigstopsignatures
\sphinxAtStartPar
A basic sensor for collecting measurements from GPS satellites.
This sensor can be used to collect measurements on demand from
presented GPS satellites at requested frequency.
\index{take\_measurement() (qnav.gps.sensor.GpsSensor method)@\spxentry{take\_measurement()}\spxextra{qnav.gps.sensor.GpsSensor method}}

\begin{fulllineitems}
\phantomsection\label{\detokenize{modules/gps/sensor:qnav.gps.sensor.GpsSensor.take_measurement}}
\pysigstartsignatures
\pysiglinewithargsret{\sphinxbfcode{\sphinxupquote{take\_measurement}}}{\sphinxparam{\DUrole{n}{est\_time}\DUrole{p}{:}\DUrole{w}{ }\DUrole{n}{float}}\sphinxparamcomma \sphinxparam{\DUrole{n}{ground\_truth}\DUrole{p}{:}\DUrole{w}{ }\DUrole{n}{{\hyperref[\detokenize{modules/waypoints/trajectory:qnav.waypoints.trajectory.GroundTruth}]{\sphinxcrossref{GroundTruth}}}}}}{}
\pysigstopsignatures
\sphinxAtStartPar
Obtains updated measurement from GPS satellites.
Obtains measurements from all GPS satellites and updates the last
measurement. Performed prior to GPS fusion.
\begin{quote}\begin{description}
\sphinxlineitem{Parameters}\begin{itemize}
\item {} 
\sphinxAtStartPar
\sphinxstyleliteralstrong{\sphinxupquote{est\_time}} (\sphinxstyleliteralemphasis{\sphinxupquote{float}}) \textendash{} The estimated time the measurement of the sensor
is expected to be captured at (in seconds).

\item {} 
\sphinxAtStartPar
\sphinxstyleliteralstrong{\sphinxupquote{ground\_truth}} ({\hyperref[\detokenize{modules/waypoints/trajectory:qnav.waypoints.trajectory.GroundTruth}]{\sphinxcrossref{\sphinxstyleliteralemphasis{\sphinxupquote{GroundTruth}}}}}) \textendash{} The corresponding ground truth data.

\end{itemize}

\end{description}\end{quote}

\end{fulllineitems}

\index{last\_measurement (qnav.gps.sensor.GpsSensor property)@\spxentry{last\_measurement}\spxextra{qnav.gps.sensor.GpsSensor property}}

\begin{fulllineitems}
\phantomsection\label{\detokenize{modules/gps/sensor:qnav.gps.sensor.GpsSensor.last_measurement}}
\pysigstartsignatures
\pysigline{\sphinxbfcode{\sphinxupquote{property\DUrole{w}{ }}}\sphinxbfcode{\sphinxupquote{last\_measurement}}\sphinxbfcode{\sphinxupquote{\DUrole{p}{:}\DUrole{w}{ }dict}}}
\pysigstopsignatures
\sphinxAtStartPar
Returns the list of latest GPS satellite measurements.
\begin{quote}\begin{description}
\sphinxlineitem{Returns}
\sphinxAtStartPar
Latest GPS satellite measurements.

\sphinxlineitem{Return type}
\sphinxAtStartPar
dict

\end{description}\end{quote}

\end{fulllineitems}

\index{update() (qnav.gps.sensor.GpsSensor method)@\spxentry{update()}\spxextra{qnav.gps.sensor.GpsSensor method}}

\begin{fulllineitems}
\phantomsection\label{\detokenize{modules/gps/sensor:qnav.gps.sensor.GpsSensor.update}}
\pysigstartsignatures
\pysiglinewithargsret{\sphinxbfcode{\sphinxupquote{update}}}{\sphinxparam{\DUrole{n}{time\_sec}\DUrole{p}{:}\DUrole{w}{ }\DUrole{n}{float}\DUrole{w}{ }\DUrole{o}{=}\DUrole{w}{ }\DUrole{default_value}{None}}}{}
\pysigstopsignatures
\sphinxAtStartPar
Updates the state of the satellites for the given amount of seconds.
This method is used to update each of the underlying satellites in
the given collection.
\begin{quote}\begin{description}
\sphinxlineitem{Parameters}
\sphinxAtStartPar
\sphinxstyleliteralstrong{\sphinxupquote{time\_sec}} (\sphinxstyleliteralemphasis{\sphinxupquote{float}}) \textendash{} The estimated current simulation time.

\end{description}\end{quote}

\end{fulllineitems}


\end{fulllineitems}


\sphinxstepscope


\subsubsection{qnav.gps.zones}
\label{\detokenize{modules/gps/zones:module-qnav.gps.zones}}\label{\detokenize{modules/gps/zones:qnav-gps-zones}}\label{\detokenize{modules/gps/zones::doc}}\index{module@\spxentry{module}!qnav.gps.zones@\spxentry{qnav.gps.zones}}\index{qnav.gps.zones@\spxentry{qnav.gps.zones}!module@\spxentry{module}}

\paragraph{zones.py}
\label{\detokenize{modules/gps/zones:zones-py}}\begin{quote}\begin{description}
\sphinxlineitem{summary}
\sphinxAtStartPar
Handles the representation and method connected with zones/areas.
This module contains a family of classes used for representing
and querying zones. This has been designed for reuse, with various
child area\sphinxhyphen{}polygon classes and their associated group handlers.

\sphinxlineitem{author}
\sphinxAtStartPar
Michael Wright \sphinxhyphen{} \sphinxhref{mailto:mjwright@liverpool.ac.uk}{mjwright@liverpool.ac.uk}

\sphinxlineitem{version}
\sphinxAtStartPar
1.0.0 \sphinxhyphen{} First implementation of this module.

\end{description}\end{quote}
\index{Polygon (class in qnav.gps.zones)@\spxentry{Polygon}\spxextra{class in qnav.gps.zones}}

\begin{fulllineitems}
\phantomsection\label{\detokenize{modules/gps/zones:qnav.gps.zones.Polygon}}
\pysigstartsignatures
\pysiglinewithargsret{\sphinxbfcode{\sphinxupquote{class\DUrole{w}{ }}}\sphinxcode{\sphinxupquote{qnav.gps.zones.}}\sphinxbfcode{\sphinxupquote{Polygon}}}{\sphinxparam{\DUrole{n}{lat\_points}\DUrole{p}{:}\DUrole{w}{ }\DUrole{n}{ndarray}}\sphinxparamcomma \sphinxparam{\DUrole{n}{lon\_points}\DUrole{p}{:}\DUrole{w}{ }\DUrole{n}{ndarray}}\sphinxparamcomma \sphinxparam{\DUrole{n}{priority}\DUrole{p}{:}\DUrole{w}{ }\DUrole{n}{int}\DUrole{w}{ }\DUrole{o}{=}\DUrole{w}{ }\DUrole{default_value}{0}}}{}
\pysigstopsignatures
\sphinxAtStartPar
Represents a simple latitude\sphinxhyphen{}longitude closed polygon.
Can be used for defining custom 2D zones on the globe.
\index{is\_point\_inside() (qnav.gps.zones.Polygon method)@\spxentry{is\_point\_inside()}\spxextra{qnav.gps.zones.Polygon method}}

\begin{fulllineitems}
\phantomsection\label{\detokenize{modules/gps/zones:qnav.gps.zones.Polygon.is_point_inside}}
\pysigstartsignatures
\pysiglinewithargsret{\sphinxbfcode{\sphinxupquote{is\_point\_inside}}}{\sphinxparam{\DUrole{n}{lat\_point}\DUrole{p}{:}\DUrole{w}{ }\DUrole{n}{float}}\sphinxparamcomma \sphinxparam{\DUrole{n}{lon\_point}\DUrole{p}{:}\DUrole{w}{ }\DUrole{n}{float}}}{{ $\rightarrow$ bool}}
\pysigstopsignatures
\sphinxAtStartPar
Determines if a given position is inside the polygon.
Given the latitude and longitude coordinates of a single point in
decimal degrees, this function will return True if the point is
inside the polygon.
\begin{quote}\begin{description}
\sphinxlineitem{Parameters}\begin{itemize}
\item {} 
\sphinxAtStartPar
\sphinxstyleliteralstrong{\sphinxupquote{lat\_point}} (\sphinxstyleliteralemphasis{\sphinxupquote{float}}) \textendash{} The latitude coordinate of the point of interest.

\item {} 
\sphinxAtStartPar
\sphinxstyleliteralstrong{\sphinxupquote{lon\_point}} (\sphinxstyleliteralemphasis{\sphinxupquote{float}}) \textendash{} The longitude coordinate of the point of interest.

\end{itemize}

\sphinxlineitem{Returns}
\sphinxAtStartPar
Returns True if given point is inside the polygon.

\sphinxlineitem{Return type}
\sphinxAtStartPar
bool

\end{description}\end{quote}

\end{fulllineitems}

\index{is\_points\_inside() (qnav.gps.zones.Polygon method)@\spxentry{is\_points\_inside()}\spxextra{qnav.gps.zones.Polygon method}}

\begin{fulllineitems}
\phantomsection\label{\detokenize{modules/gps/zones:qnav.gps.zones.Polygon.is_points_inside}}
\pysigstartsignatures
\pysiglinewithargsret{\sphinxbfcode{\sphinxupquote{is\_points\_inside}}}{\sphinxparam{\DUrole{n}{lat\_points}\DUrole{p}{:}\DUrole{w}{ }\DUrole{n}{ndarray}}\sphinxparamcomma \sphinxparam{\DUrole{n}{lon\_points}\DUrole{p}{:}\DUrole{w}{ }\DUrole{n}{ndarray}}}{{ $\rightarrow$ ndarray}}
\pysigstopsignatures
\sphinxAtStartPar
Determines if a series of positions are inside the polygon.
Given the latitude and longitude coordinates of a multiple points
in decimal degrees, this function will return True for each point
that is inside the polygon.
\begin{quote}\begin{description}
\sphinxlineitem{Parameters}\begin{itemize}
\item {} 
\sphinxAtStartPar
\sphinxstyleliteralstrong{\sphinxupquote{lat\_points}} (\sphinxstyleliteralemphasis{\sphinxupquote{float}}) \textendash{} The latitude coordinates of the points of interest.

\item {} 
\sphinxAtStartPar
\sphinxstyleliteralstrong{\sphinxupquote{lon\_points}} (\sphinxstyleliteralemphasis{\sphinxupquote{float}}) \textendash{} The longitude coordinates of the points of interest.

\end{itemize}

\sphinxlineitem{Returns}
\sphinxAtStartPar
Returns True for each given point inside the polygon.

\sphinxlineitem{Return type}
\sphinxAtStartPar
bool

\end{description}\end{quote}

\end{fulllineitems}

\index{lat\_points (qnav.gps.zones.Polygon property)@\spxentry{lat\_points}\spxextra{qnav.gps.zones.Polygon property}}

\begin{fulllineitems}
\phantomsection\label{\detokenize{modules/gps/zones:qnav.gps.zones.Polygon.lat_points}}
\pysigstartsignatures
\pysigline{\sphinxbfcode{\sphinxupquote{property\DUrole{w}{ }}}\sphinxbfcode{\sphinxupquote{lat\_points}}\sphinxbfcode{\sphinxupquote{\DUrole{p}{:}\DUrole{w}{ }ndarray}}}
\pysigstopsignatures
\sphinxAtStartPar
Returns the list of latitude coordinate points of the polygon.
This property returns a copy, preventing accidental editing.
\begin{quote}\begin{description}
\sphinxlineitem{Returns}
\sphinxAtStartPar
The latitude coordinate points of the polygon.

\sphinxlineitem{Return type}
\sphinxAtStartPar
np.ndarray

\end{description}\end{quote}

\end{fulllineitems}

\index{lon\_points (qnav.gps.zones.Polygon property)@\spxentry{lon\_points}\spxextra{qnav.gps.zones.Polygon property}}

\begin{fulllineitems}
\phantomsection\label{\detokenize{modules/gps/zones:qnav.gps.zones.Polygon.lon_points}}
\pysigstartsignatures
\pysigline{\sphinxbfcode{\sphinxupquote{property\DUrole{w}{ }}}\sphinxbfcode{\sphinxupquote{lon\_points}}\sphinxbfcode{\sphinxupquote{\DUrole{p}{:}\DUrole{w}{ }ndarray}}}
\pysigstopsignatures
\sphinxAtStartPar
Returns the list of longitude coordinate points of the polygon.
This property returns a copy, preventing accidental editing.
\begin{quote}\begin{description}
\sphinxlineitem{Returns}
\sphinxAtStartPar
The longitude coordinate points of the polygon.

\sphinxlineitem{Return type}
\sphinxAtStartPar
np.ndarray

\end{description}\end{quote}

\end{fulllineitems}

\index{priority (qnav.gps.zones.Polygon property)@\spxentry{priority}\spxextra{qnav.gps.zones.Polygon property}}

\begin{fulllineitems}
\phantomsection\label{\detokenize{modules/gps/zones:qnav.gps.zones.Polygon.priority}}
\pysigstartsignatures
\pysigline{\sphinxbfcode{\sphinxupquote{property\DUrole{w}{ }}}\sphinxbfcode{\sphinxupquote{priority}}\sphinxbfcode{\sphinxupquote{\DUrole{p}{:}\DUrole{w}{ }int}}}
\pysigstopsignatures
\sphinxAtStartPar
Returns the priority (or layer) of the polygon.
This is used for determining the outcome when multiple polygons
overlapping one another are involved. When this occurs the highest
priority polygon is used.
\begin{quote}\begin{description}
\sphinxlineitem{Returns}
\sphinxAtStartPar
The priority (or layer) of the polygon.

\sphinxlineitem{Return type}
\sphinxAtStartPar
int

\end{description}\end{quote}

\end{fulllineitems}


\end{fulllineitems}

\index{FlaggedPolygon (class in qnav.gps.zones)@\spxentry{FlaggedPolygon}\spxextra{class in qnav.gps.zones}}

\begin{fulllineitems}
\phantomsection\label{\detokenize{modules/gps/zones:qnav.gps.zones.FlaggedPolygon}}
\pysigstartsignatures
\pysiglinewithargsret{\sphinxbfcode{\sphinxupquote{class\DUrole{w}{ }}}\sphinxcode{\sphinxupquote{qnav.gps.zones.}}\sphinxbfcode{\sphinxupquote{FlaggedPolygon}}}{\sphinxparam{\DUrole{n}{lat\_points}\DUrole{p}{:}\DUrole{w}{ }\DUrole{n}{ndarray}}\sphinxparamcomma \sphinxparam{\DUrole{n}{lon\_points}\DUrole{p}{:}\DUrole{w}{ }\DUrole{n}{ndarray}}\sphinxparamcomma \sphinxparam{\DUrole{n}{is\_active\_area}\DUrole{p}{:}\DUrole{w}{ }\DUrole{n}{bool}}\sphinxparamcomma \sphinxparam{\DUrole{n}{priority}\DUrole{p}{:}\DUrole{w}{ }\DUrole{n}{int}\DUrole{w}{ }\DUrole{o}{=}\DUrole{w}{ }\DUrole{default_value}{0}}}{}
\pysigstopsignatures
\sphinxAtStartPar
Represents a latitude\sphinxhyphen{}longitude closed polygon with binary active flag.
These polygons are used for determining “active” and “inactive” areas.
For example, areas where specific sensors, such as where GPS or
quantum sensors can and cannot be used.
\index{is\_active\_area (qnav.gps.zones.FlaggedPolygon property)@\spxentry{is\_active\_area}\spxextra{qnav.gps.zones.FlaggedPolygon property}}

\begin{fulllineitems}
\phantomsection\label{\detokenize{modules/gps/zones:qnav.gps.zones.FlaggedPolygon.is_active_area}}
\pysigstartsignatures
\pysigline{\sphinxbfcode{\sphinxupquote{property\DUrole{w}{ }}}\sphinxbfcode{\sphinxupquote{is\_active\_area}}\sphinxbfcode{\sphinxupquote{\DUrole{p}{:}\DUrole{w}{ }bool}}}
\pysigstopsignatures
\sphinxAtStartPar
Returns True if the polygon is active area.
This can be used to determine if points which are found inside the
polygon are to be considered inside an active area.
\begin{quote}\begin{description}
\sphinxlineitem{Returns}
\sphinxAtStartPar
Whether the polygon is an active area.

\sphinxlineitem{Return type}
\sphinxAtStartPar
bool

\end{description}\end{quote}

\end{fulllineitems}


\end{fulllineitems}

\index{FlaggedPolygonGroup (class in qnav.gps.zones)@\spxentry{FlaggedPolygonGroup}\spxextra{class in qnav.gps.zones}}

\begin{fulllineitems}
\phantomsection\label{\detokenize{modules/gps/zones:qnav.gps.zones.FlaggedPolygonGroup}}
\pysigstartsignatures
\pysiglinewithargsret{\sphinxbfcode{\sphinxupquote{class\DUrole{w}{ }}}\sphinxcode{\sphinxupquote{qnav.gps.zones.}}\sphinxbfcode{\sphinxupquote{FlaggedPolygonGroup}}}{\sphinxparam{\DUrole{n}{polygons}\DUrole{p}{:}\DUrole{w}{ }\DUrole{n}{List\DUrole{p}{{[}}{\hyperref[\detokenize{modules/gps/zones:qnav.gps.zones.FlaggedPolygon}]{\sphinxcrossref{FlaggedPolygon}}}\DUrole{p}{{]}}}}\sphinxparamcomma \sphinxparam{\DUrole{n}{is\_active\_by\_default}\DUrole{p}{:}\DUrole{w}{ }\DUrole{n}{bool}\DUrole{w}{ }\DUrole{o}{=}\DUrole{w}{ }\DUrole{default_value}{True}}}{}
\pysigstopsignatures
\sphinxAtStartPar
A handler for a group of flagged polygons.
Allows the querying of multiple polygons, in priority order,
with a default flag.
\index{get\_point\_flag() (qnav.gps.zones.FlaggedPolygonGroup method)@\spxentry{get\_point\_flag()}\spxextra{qnav.gps.zones.FlaggedPolygonGroup method}}

\begin{fulllineitems}
\phantomsection\label{\detokenize{modules/gps/zones:qnav.gps.zones.FlaggedPolygonGroup.get_point_flag}}
\pysigstartsignatures
\pysiglinewithargsret{\sphinxbfcode{\sphinxupquote{get\_point\_flag}}}{\sphinxparam{\DUrole{n}{lat\_point}\DUrole{p}{:}\DUrole{w}{ }\DUrole{n}{float}}\sphinxparamcomma \sphinxparam{\DUrole{n}{lon\_point}\DUrole{p}{:}\DUrole{w}{ }\DUrole{n}{float}}}{{ $\rightarrow$ bool}}
\pysigstopsignatures
\sphinxAtStartPar
Returns the boolean flag for given point.
Searches through given polygons in priority order and returns the
flag for the highest priority polygon that the point is inside of.
If the point is not in any polygons, the default flag will be used.
\begin{quote}\begin{description}
\sphinxlineitem{Parameters}\begin{itemize}
\item {} 
\sphinxAtStartPar
\sphinxstyleliteralstrong{\sphinxupquote{lat\_point}} (\sphinxstyleliteralemphasis{\sphinxupquote{float}}) \textendash{} The latitude coordinate of the point of interest.

\item {} 
\sphinxAtStartPar
\sphinxstyleliteralstrong{\sphinxupquote{lon\_point}} (\sphinxstyleliteralemphasis{\sphinxupquote{float}}) \textendash{} The longitude coordinate of the point of interest.

\end{itemize}

\sphinxlineitem{Returns}
\sphinxAtStartPar
Boolean flag of the highest priority polygon the point is
inside of (or the default flag for the group).

\sphinxlineitem{Return type}
\sphinxAtStartPar
bool

\end{description}\end{quote}

\end{fulllineitems}

\index{polygons (qnav.gps.zones.FlaggedPolygonGroup property)@\spxentry{polygons}\spxextra{qnav.gps.zones.FlaggedPolygonGroup property}}

\begin{fulllineitems}
\phantomsection\label{\detokenize{modules/gps/zones:qnav.gps.zones.FlaggedPolygonGroup.polygons}}
\pysigstartsignatures
\pysigline{\sphinxbfcode{\sphinxupquote{property\DUrole{w}{ }}}\sphinxbfcode{\sphinxupquote{polygons}}\sphinxbfcode{\sphinxupquote{\DUrole{p}{:}\DUrole{w}{ }List\DUrole{p}{{[}}{\hyperref[\detokenize{modules/gps/zones:qnav.gps.zones.FlaggedPolygon}]{\sphinxcrossref{FlaggedPolygon}}}\DUrole{p}{{]}}}}}
\pysigstopsignatures
\sphinxAtStartPar
Returns list of flagged polygons used by group.
\begin{quote}\begin{description}
\sphinxlineitem{Returns}
\sphinxAtStartPar
list of flagged polygons used by group

\sphinxlineitem{Return type}
\sphinxAtStartPar
List{[}{\hyperref[\detokenize{modules/gps/zones:qnav.gps.zones.FlaggedPolygon}]{\sphinxcrossref{FlaggedPolygon}}}{]}

\end{description}\end{quote}

\end{fulllineitems}

\index{num\_polygons (qnav.gps.zones.FlaggedPolygonGroup property)@\spxentry{num\_polygons}\spxextra{qnav.gps.zones.FlaggedPolygonGroup property}}

\begin{fulllineitems}
\phantomsection\label{\detokenize{modules/gps/zones:qnav.gps.zones.FlaggedPolygonGroup.num_polygons}}
\pysigstartsignatures
\pysigline{\sphinxbfcode{\sphinxupquote{property\DUrole{w}{ }}}\sphinxbfcode{\sphinxupquote{num\_polygons}}\sphinxbfcode{\sphinxupquote{\DUrole{p}{:}\DUrole{w}{ }int}}}
\pysigstopsignatures
\sphinxAtStartPar
Returns the number of polygons covered by the group.
\begin{quote}\begin{description}
\sphinxlineitem{Returns}
\sphinxAtStartPar
The number of polygons covered by the group.

\sphinxlineitem{Return type}
\sphinxAtStartPar
int

\end{description}\end{quote}

\end{fulllineitems}

\index{get\_point\_flags() (qnav.gps.zones.FlaggedPolygonGroup method)@\spxentry{get\_point\_flags()}\spxextra{qnav.gps.zones.FlaggedPolygonGroup method}}

\begin{fulllineitems}
\phantomsection\label{\detokenize{modules/gps/zones:qnav.gps.zones.FlaggedPolygonGroup.get_point_flags}}
\pysigstartsignatures
\pysiglinewithargsret{\sphinxbfcode{\sphinxupquote{get\_point\_flags}}}{\sphinxparam{\DUrole{n}{lat\_points}\DUrole{p}{:}\DUrole{w}{ }\DUrole{n}{ndarray}}\sphinxparamcomma \sphinxparam{\DUrole{n}{lon\_points}\DUrole{p}{:}\DUrole{w}{ }\DUrole{n}{ndarray}}}{{ $\rightarrow$ ndarray}}
\pysigstopsignatures
\sphinxAtStartPar
Returns the boolean flags for multiple given points.
Searches through given polygons in priority order and returns the
flags of the highest priority polygons that each of the point is
inside of. For points that do not fall inside nay polygon’s the
default flag is used.
\begin{quote}\begin{description}
\sphinxlineitem{Parameters}\begin{itemize}
\item {} 
\sphinxAtStartPar
\sphinxstyleliteralstrong{\sphinxupquote{lat\_points}} (\sphinxstyleliteralemphasis{\sphinxupquote{np.ndarray}}) \textendash{} The latitude coordinates of the points of interest.

\item {} 
\sphinxAtStartPar
\sphinxstyleliteralstrong{\sphinxupquote{lon\_points}} (\sphinxstyleliteralemphasis{\sphinxupquote{np.ndarray}}) \textendash{} The longitude coordinates of the points of interest.

\end{itemize}

\sphinxlineitem{Returns}
\sphinxAtStartPar
Boolean flags of the highest priority polygon each point is
inside of (or the default flag for the group).

\sphinxlineitem{Return type}
\sphinxAtStartPar
np.ndarray

\end{description}\end{quote}

\end{fulllineitems}


\end{fulllineitems}

\index{is\_inside\_polygon() (in module qnav.gps.zones)@\spxentry{is\_inside\_polygon()}\spxextra{in module qnav.gps.zones}}

\begin{fulllineitems}
\phantomsection\label{\detokenize{modules/gps/zones:qnav.gps.zones.is_inside_polygon}}
\pysigstartsignatures
\pysiglinewithargsret{\sphinxcode{\sphinxupquote{qnav.gps.zones.}}\sphinxbfcode{\sphinxupquote{is\_inside\_polygon}}}{\sphinxparam{\DUrole{n}{x\_point}\DUrole{p}{:}\DUrole{w}{ }\DUrole{n}{float}}\sphinxparamcomma \sphinxparam{\DUrole{n}{y\_point}\DUrole{p}{:}\DUrole{w}{ }\DUrole{n}{float}}\sphinxparamcomma \sphinxparam{\DUrole{n}{polygon\_x}\DUrole{p}{:}\DUrole{w}{ }\DUrole{n}{ndarray}}\sphinxparamcomma \sphinxparam{\DUrole{n}{polygon\_y}\DUrole{p}{:}\DUrole{w}{ }\DUrole{n}{ndarray}}}{{ $\rightarrow$ bool}}
\pysigstopsignatures
\sphinxAtStartPar
Determines if a given position is inside given the polygon.
Tests by generating horizontal lines and assessing the number of
polygon lines that are crossed. Using this count, a point can be
determine to be inside the closed polygon or not.
\begin{quote}\begin{description}
\sphinxlineitem{Parameters}\begin{itemize}
\item {} 
\sphinxAtStartPar
\sphinxstyleliteralstrong{\sphinxupquote{x\_point}} (\sphinxstyleliteralemphasis{\sphinxupquote{float}}) \textendash{} The x coordinate of the point of interest.

\item {} 
\sphinxAtStartPar
\sphinxstyleliteralstrong{\sphinxupquote{y\_point}} (\sphinxstyleliteralemphasis{\sphinxupquote{float}}) \textendash{} The y coordinate of the point of interest.

\item {} 
\sphinxAtStartPar
\sphinxstyleliteralstrong{\sphinxupquote{polygon\_x}} (\sphinxstyleliteralemphasis{\sphinxupquote{np.ndarray}}) \textendash{} The x\sphinxhyphen{}coordinates of the polygon. Assumed to be closed.

\item {} 
\sphinxAtStartPar
\sphinxstyleliteralstrong{\sphinxupquote{polygon\_y}} (\sphinxstyleliteralemphasis{\sphinxupquote{np.ndarray}}) \textendash{} The y\sphinxhyphen{}coordinates of the polygon. Assumed to be closed.

\end{itemize}

\sphinxlineitem{Returns}
\sphinxAtStartPar
Returns True if the given position is inside the polygon.

\sphinxlineitem{Return type}
\sphinxAtStartPar
bool

\end{description}\end{quote}

\end{fulllineitems}

\index{is\_inside\_polygon\_vec() (in module qnav.gps.zones)@\spxentry{is\_inside\_polygon\_vec()}\spxextra{in module qnav.gps.zones}}

\begin{fulllineitems}
\phantomsection\label{\detokenize{modules/gps/zones:qnav.gps.zones.is_inside_polygon_vec}}
\pysigstartsignatures
\pysiglinewithargsret{\sphinxcode{\sphinxupquote{qnav.gps.zones.}}\sphinxbfcode{\sphinxupquote{is\_inside\_polygon\_vec}}}{\sphinxparam{\DUrole{n}{x\_points}\DUrole{p}{:}\DUrole{w}{ }\DUrole{n}{ndarray}}\sphinxparamcomma \sphinxparam{\DUrole{n}{y\_points}\DUrole{p}{:}\DUrole{w}{ }\DUrole{n}{ndarray}}\sphinxparamcomma \sphinxparam{\DUrole{n}{polygon\_x}\DUrole{p}{:}\DUrole{w}{ }\DUrole{n}{ndarray}}\sphinxparamcomma \sphinxparam{\DUrole{n}{polygon\_y}\DUrole{p}{:}\DUrole{w}{ }\DUrole{n}{ndarray}}}{{ $\rightarrow$ ndarray}}
\pysigstopsignatures
\sphinxAtStartPar
Determines if a given position is inside given the polygon.
Tests by generating horizontal lines and assessing the number of
polygon lines that are crossed. Using this count, points can be
determine to be inside the closed polygon or not.
\begin{quote}\begin{description}
\sphinxlineitem{Parameters}\begin{itemize}
\item {} 
\sphinxAtStartPar
\sphinxstyleliteralstrong{\sphinxupquote{x\_points}} (\sphinxstyleliteralemphasis{\sphinxupquote{float}}) \textendash{} The x coordinates of the point of interest.

\item {} 
\sphinxAtStartPar
\sphinxstyleliteralstrong{\sphinxupquote{y\_points}} (\sphinxstyleliteralemphasis{\sphinxupquote{float}}) \textendash{} The y coordinates of the point of interest.

\item {} 
\sphinxAtStartPar
\sphinxstyleliteralstrong{\sphinxupquote{polygon\_x}} (\sphinxstyleliteralemphasis{\sphinxupquote{np.ndarray}}) \textendash{} The x\sphinxhyphen{}coordinates of the polygon. Assumed to be closed.

\item {} 
\sphinxAtStartPar
\sphinxstyleliteralstrong{\sphinxupquote{polygon\_y}} (\sphinxstyleliteralemphasis{\sphinxupquote{np.ndarray}}) \textendash{} The y\sphinxhyphen{}coordinates of the polygon. Assumed to be closed.

\end{itemize}

\sphinxlineitem{Returns}
\sphinxAtStartPar
Returns True if the given position is inside the polygon.

\sphinxlineitem{Return type}
\sphinxAtStartPar
bool

\end{description}\end{quote}

\end{fulllineitems}



\subsection{Gravity}
\label{\detokenize{usage/packages:gravity}}
\sphinxAtStartPar
Modules relating to gravitational calculations:

\sphinxstepscope


\subsubsection{qnav.gravity.base}
\label{\detokenize{modules/gravity/base:module-qnav.gravity.base}}\label{\detokenize{modules/gravity/base:qnav-gravity-base}}\label{\detokenize{modules/gravity/base::doc}}\index{module@\spxentry{module}!qnav.gravity.base@\spxentry{qnav.gravity.base}}\index{qnav.gravity.base@\spxentry{qnav.gravity.base}!module@\spxentry{module}}

\paragraph{base.py}
\label{\detokenize{modules/gravity/base:base-py}}\begin{quote}\begin{description}
\sphinxlineitem{summary}
\sphinxAtStartPar
Gravity model template class interfaces that supported gravity models
must extend from. The purpose of this is to define the functions, input
arguments and outputs that gravity models must incorporate to be used
in the toolbox. This also simplifies the ability to add new gravity
models future.

\sphinxlineitem{author}
\sphinxAtStartPar
Michael Wright \sphinxhyphen{} \sphinxhref{mailto:mjwright@liverpool.ac.uk}{mjwright@liverpool.ac.uk}

\sphinxlineitem{version}
\sphinxAtStartPar
1.0.0 \sphinxhyphen{} First implementation of this module.

\end{description}\end{quote}
\index{GravityModel (class in qnav.gravity.base)@\spxentry{GravityModel}\spxextra{class in qnav.gravity.base}}

\begin{fulllineitems}
\phantomsection\label{\detokenize{modules/gravity/base:qnav.gravity.base.GravityModel}}
\pysigstartsignatures
\pysigline{\sphinxbfcode{\sphinxupquote{class\DUrole{w}{ }}}\sphinxcode{\sphinxupquote{qnav.gravity.base.}}\sphinxbfcode{\sphinxupquote{GravityModel}}}
\pysigstopsignatures
\sphinxAtStartPar
The abstract class that gravity models are to follow.
This defines the interface that all gravity models in the project must
follow, outlining the necessary functions they must supply for
calculating various gravity components at requested positions. For a
gravity model to be used by the toolbox, it must correctly extend from
this abstract base class.
\begin{description}
\sphinxlineitem{In short, each gravity model must supply the following methods:}\begin{itemize}
\item {} 
\sphinxAtStartPar
\_\_str\_\_ \sphinxhyphen{} The descriptive name for the gravity model.

\item {} 
\sphinxAtStartPar
calc\_gravity\_z \sphinxhyphen{} Calculate downwards gravitational acceleration.

\item {} 
\sphinxAtStartPar
calc\_gravity\_xyz \sphinxhyphen{} Calculate gravitational acceleration vector.

\item {} 
\sphinxAtStartPar
calc\_gravity\_grad \sphinxhyphen{} Calculate gravitational gradient tensor.

\end{itemize}

\sphinxlineitem{In addition to this, there are also optional methods they can supply:}\begin{itemize}
\item {} 
\sphinxAtStartPar
calc\_gravity\_z\_vec \sphinxhyphen{} Vectorised downwards gravitational acceleration.

\item {} 
\sphinxAtStartPar
calc\_gravity\_xyz\_vec \sphinxhyphen{} Vectorised gravitational acceleration vector.

\item {} 
\sphinxAtStartPar
calc\_gravity\_grad\_vec \sphinxhyphen{} Vectorised gravitational gradient tensor.

\end{itemize}

\end{description}

\sphinxAtStartPar
These are used to efficiently calculate gravity for a large range of
positions. To enforce compatibility, if not implemented, calling these
methods will execute the scalar corresponding methods in a loop.
\index{calc\_gravity\_z() (qnav.gravity.base.GravityModel method)@\spxentry{calc\_gravity\_z()}\spxextra{qnav.gravity.base.GravityModel method}}

\begin{fulllineitems}
\phantomsection\label{\detokenize{modules/gravity/base:qnav.gravity.base.GravityModel.calc_gravity_z}}
\pysigstartsignatures
\pysiglinewithargsret{\sphinxbfcode{\sphinxupquote{abstract\DUrole{w}{ }}}\sphinxbfcode{\sphinxupquote{calc\_gravity\_z}}}{\sphinxparam{\DUrole{n}{lat}\DUrole{p}{:}\DUrole{w}{ }\DUrole{n}{float}}\sphinxparamcomma \sphinxparam{\DUrole{n}{lon}\DUrole{p}{:}\DUrole{w}{ }\DUrole{n}{float}}\sphinxparamcomma \sphinxparam{\DUrole{n}{alt}\DUrole{p}{:}\DUrole{w}{ }\DUrole{n}{float}}}{{ $\rightarrow$ float}}
\pysigstopsignatures
\sphinxAtStartPar
Calculates and returns the total gravitational acceleration.
Computes the gravity along the Z/Down axis (in m/s\textasciicircum{}2) for
the given location and altitude offset.
\begin{quote}\begin{description}
\sphinxlineitem{Parameters}\begin{itemize}
\item {} 
\sphinxAtStartPar
\sphinxstyleliteralstrong{\sphinxupquote{lat}} (\sphinxstyleliteralemphasis{\sphinxupquote{float}}) \textendash{} The latitude coordinate for the position of interest.
The latitude position value given in decimal degrees.

\item {} 
\sphinxAtStartPar
\sphinxstyleliteralstrong{\sphinxupquote{lon}} (\sphinxstyleliteralemphasis{\sphinxupquote{float}}) \textendash{} The longitude coordinate for the position of interest.
The longitude position value given in decimal degrees.

\item {} 
\sphinxAtStartPar
\sphinxstyleliteralstrong{\sphinxupquote{alt}} (\sphinxstyleliteralemphasis{\sphinxupquote{float}}) \textendash{} The altitude offset for the position of interest.
The elevation from the WGS\sphinxhyphen{}84 reference ellipsoid in metres.

\end{itemize}

\sphinxlineitem{Returns}
\sphinxAtStartPar
The downwards gravitational acceleration for each location.

\sphinxlineitem{Return type}
\sphinxAtStartPar
float

\end{description}\end{quote}

\end{fulllineitems}

\index{calc\_gravity\_xyz() (qnav.gravity.base.GravityModel method)@\spxentry{calc\_gravity\_xyz()}\spxextra{qnav.gravity.base.GravityModel method}}

\begin{fulllineitems}
\phantomsection\label{\detokenize{modules/gravity/base:qnav.gravity.base.GravityModel.calc_gravity_xyz}}
\pysigstartsignatures
\pysiglinewithargsret{\sphinxbfcode{\sphinxupquote{calc\_gravity\_xyz}}}{\sphinxparam{\DUrole{n}{lat}\DUrole{p}{:}\DUrole{w}{ }\DUrole{n}{float}}\sphinxparamcomma \sphinxparam{\DUrole{n}{lon}\DUrole{p}{:}\DUrole{w}{ }\DUrole{n}{float}}\sphinxparamcomma \sphinxparam{\DUrole{n}{alt}\DUrole{p}{:}\DUrole{w}{ }\DUrole{n}{float}}}{{ $\rightarrow$ ndarray}}
\pysigstopsignatures
\sphinxAtStartPar
Calculates and returns the gravitational acceleration vector.
Computes the gravity acceleration X/North, Y/East and Z/Down
components (in m/s\textasciicircum{}2) for the given location and altitude offset.
\begin{quote}\begin{description}
\sphinxlineitem{Parameters}\begin{itemize}
\item {} 
\sphinxAtStartPar
\sphinxstyleliteralstrong{\sphinxupquote{lat}} (\sphinxstyleliteralemphasis{\sphinxupquote{float}}) \textendash{} The latitude coordinate for the position of interest.
The latitude position value given in decimal degrees.

\item {} 
\sphinxAtStartPar
\sphinxstyleliteralstrong{\sphinxupquote{lon}} (\sphinxstyleliteralemphasis{\sphinxupquote{float}}) \textendash{} The longitude coordinate for the position of interest.
The longitude position value given in decimal degrees.

\item {} 
\sphinxAtStartPar
\sphinxstyleliteralstrong{\sphinxupquote{alt}} (\sphinxstyleliteralemphasis{\sphinxupquote{float}}) \textendash{} The altitude offset for the position of interest.
The elevation from the WGS\sphinxhyphen{}84 reference ellipsoid in metres.

\end{itemize}

\sphinxlineitem{Returns}
\sphinxAtStartPar
The local gravity vector for the XYZ components (in m/s\textasciicircum{}2).

\sphinxlineitem{Return type}
\sphinxAtStartPar
numpy.ndarray (3\sphinxhyphen{}elements)

\end{description}\end{quote}

\end{fulllineitems}

\index{calc\_vertical\_grad() (qnav.gravity.base.GravityModel method)@\spxentry{calc\_vertical\_grad()}\spxextra{qnav.gravity.base.GravityModel method}}

\begin{fulllineitems}
\phantomsection\label{\detokenize{modules/gravity/base:qnav.gravity.base.GravityModel.calc_vertical_grad}}
\pysigstartsignatures
\pysiglinewithargsret{\sphinxbfcode{\sphinxupquote{calc\_vertical\_grad}}}{\sphinxparam{\DUrole{n}{lat}\DUrole{p}{:}\DUrole{w}{ }\DUrole{n}{float}}\sphinxparamcomma \sphinxparam{\DUrole{n}{lon}\DUrole{p}{:}\DUrole{w}{ }\DUrole{n}{float}}\sphinxparamcomma \sphinxparam{\DUrole{n}{alt}\DUrole{p}{:}\DUrole{w}{ }\DUrole{n}{float}}\sphinxparamcomma \sphinxparam{\DUrole{n}{step}\DUrole{p}{:}\DUrole{w}{ }\DUrole{n}{float}\DUrole{w}{ }\DUrole{o}{=}\DUrole{w}{ }\DUrole{default_value}{100}}}{{ $\rightarrow$ float}}
\pysigstopsignatures
\sphinxAtStartPar
Calculates the vertical gradient of gravity.
Computes the vertical change in gravity (dg/dr) of the Z/Down
component (in m/s\textasciicircum{}2). This is primary used for gravity map\sphinxhyphen{}matching
related calculations.
\begin{quote}\begin{description}
\sphinxlineitem{Parameters}\begin{itemize}
\item {} 
\sphinxAtStartPar
\sphinxstyleliteralstrong{\sphinxupquote{lat}} (\sphinxstyleliteralemphasis{\sphinxupquote{float}}) \textendash{} The latitude coordinate for the position of interest.
The latitude position value given in decimal degrees.

\item {} 
\sphinxAtStartPar
\sphinxstyleliteralstrong{\sphinxupquote{lon}} (\sphinxstyleliteralemphasis{\sphinxupquote{float}}) \textendash{} The longitude coordinate for the position of interest.
The longitude position value given in decimal degrees.

\item {} 
\sphinxAtStartPar
\sphinxstyleliteralstrong{\sphinxupquote{alt}} (\sphinxstyleliteralemphasis{\sphinxupquote{float}}) \textendash{} The altitude offset for the position of interest.
The elevation from the WGS\sphinxhyphen{}84 reference ellipsoid in metres.

\item {} 
\sphinxAtStartPar
\sphinxstyleliteralstrong{\sphinxupquote{step}} (\sphinxstyleliteralemphasis{\sphinxupquote{float}}) \textendash{} The step size (in metres) for calculating the gradient.
To be used when the gradient cannot be computed directly.

\end{itemize}

\sphinxlineitem{Returns}
\sphinxAtStartPar
The vertical gravity gradient (in m/s\textasciicircum{}2).

\sphinxlineitem{Return type}
\sphinxAtStartPar
float

\end{description}\end{quote}

\end{fulllineitems}

\index{calc\_gravity\_z\_vec() (qnav.gravity.base.GravityModel method)@\spxentry{calc\_gravity\_z\_vec()}\spxextra{qnav.gravity.base.GravityModel method}}

\begin{fulllineitems}
\phantomsection\label{\detokenize{modules/gravity/base:qnav.gravity.base.GravityModel.calc_gravity_z_vec}}
\pysigstartsignatures
\pysiglinewithargsret{\sphinxbfcode{\sphinxupquote{calc\_gravity\_z\_vec}}}{\sphinxparam{\DUrole{n}{lat}\DUrole{p}{:}\DUrole{w}{ }\DUrole{n}{ndarray}}\sphinxparamcomma \sphinxparam{\DUrole{n}{lon}\DUrole{p}{:}\DUrole{w}{ }\DUrole{n}{ndarray}}\sphinxparamcomma \sphinxparam{\DUrole{n}{alt}\DUrole{p}{:}\DUrole{w}{ }\DUrole{n}{ndarray}}}{{ $\rightarrow$ ndarray}}
\pysigstopsignatures
\sphinxAtStartPar
In a vectorised batch, calculates and returns the gravitational
acceleration along the Z/Down axis (in m/s\textasciicircum{}2) for the given location
and altitude offset. Ideal when calculating gravity for a large
number of positions.
\begin{quote}\begin{description}
\sphinxlineitem{Parameters}\begin{itemize}
\item {} 
\sphinxAtStartPar
\sphinxstyleliteralstrong{\sphinxupquote{lat}} (\sphinxstyleliteralemphasis{\sphinxupquote{np.array}}) \textendash{} The latitude coordinate for the position of interest.
The latitude position value given in decimal degrees.

\item {} 
\sphinxAtStartPar
\sphinxstyleliteralstrong{\sphinxupquote{lon}} (\sphinxstyleliteralemphasis{\sphinxupquote{np.array}}) \textendash{} The longitude coordinate for the position of interest.
The longitude position value given in decimal degrees.

\item {} 
\sphinxAtStartPar
\sphinxstyleliteralstrong{\sphinxupquote{alt}} (\sphinxstyleliteralemphasis{\sphinxupquote{np.array}}) \textendash{} The altitude offset for the position of interest.
The elevation from the WGS\sphinxhyphen{}84 reference ellipsoid in metres.

\end{itemize}

\sphinxlineitem{Returns}
\sphinxAtStartPar
The downwards gravitational acceleration for each location
(in m/s\textasciicircum{}2).

\sphinxlineitem{Return type}
\sphinxAtStartPar
np.array (N\sphinxhyphen{}elements)

\end{description}\end{quote}

\end{fulllineitems}

\index{calc\_gravity\_xyz\_vec() (qnav.gravity.base.GravityModel method)@\spxentry{calc\_gravity\_xyz\_vec()}\spxextra{qnav.gravity.base.GravityModel method}}

\begin{fulllineitems}
\phantomsection\label{\detokenize{modules/gravity/base:qnav.gravity.base.GravityModel.calc_gravity_xyz_vec}}
\pysigstartsignatures
\pysiglinewithargsret{\sphinxbfcode{\sphinxupquote{calc\_gravity\_xyz\_vec}}}{\sphinxparam{\DUrole{n}{lat}\DUrole{p}{:}\DUrole{w}{ }\DUrole{n}{ndarray}}\sphinxparamcomma \sphinxparam{\DUrole{n}{lon}\DUrole{p}{:}\DUrole{w}{ }\DUrole{n}{ndarray}}\sphinxparamcomma \sphinxparam{\DUrole{n}{alt}\DUrole{p}{:}\DUrole{w}{ }\DUrole{n}{ndarray}}}{{ $\rightarrow$ ndarray}}
\pysigstopsignatures
\sphinxAtStartPar
In a vectorised batch, calculates and returns the gravitational
acceleration for the X/North, Y/East and Z/Down axes (in m/s\textasciicircum{}2)
for the given location and altitude offset. Ideal when calculating
gravity for a large number of positions.
\begin{quote}\begin{description}
\sphinxlineitem{Parameters}\begin{itemize}
\item {} 
\sphinxAtStartPar
\sphinxstyleliteralstrong{\sphinxupquote{lat}} (\sphinxstyleliteralemphasis{\sphinxupquote{np.array}}) \textendash{} The latitude coordinate for the position of interest.
The latitude position value given in decimal degrees.

\item {} 
\sphinxAtStartPar
\sphinxstyleliteralstrong{\sphinxupquote{lon}} (\sphinxstyleliteralemphasis{\sphinxupquote{np.array}}) \textendash{} The longitude coordinate for the position of interest.
The longitude position value given in decimal degrees.

\item {} 
\sphinxAtStartPar
\sphinxstyleliteralstrong{\sphinxupquote{alt}} (\sphinxstyleliteralemphasis{\sphinxupquote{np.array}}) \textendash{} The altitude offset for the position of interest.
The elevation from the WGS\sphinxhyphen{}84 reference ellipsoid in metres.

\end{itemize}

\sphinxlineitem{Returns}
\sphinxAtStartPar
The local gravity vector for the XYZ components (in m/s\textasciicircum{}2).

\sphinxlineitem{Return type}
\sphinxAtStartPar
numpy.ndarray (N\sphinxhyphen{}by\sphinxhyphen{}3 elements)

\end{description}\end{quote}

\end{fulllineitems}

\index{calc\_vertical\_grad\_vec() (qnav.gravity.base.GravityModel method)@\spxentry{calc\_vertical\_grad\_vec()}\spxextra{qnav.gravity.base.GravityModel method}}

\begin{fulllineitems}
\phantomsection\label{\detokenize{modules/gravity/base:qnav.gravity.base.GravityModel.calc_vertical_grad_vec}}
\pysigstartsignatures
\pysiglinewithargsret{\sphinxbfcode{\sphinxupquote{calc\_vertical\_grad\_vec}}}{\sphinxparam{\DUrole{n}{lat}\DUrole{p}{:}\DUrole{w}{ }\DUrole{n}{ndarray}}\sphinxparamcomma \sphinxparam{\DUrole{n}{lon}\DUrole{p}{:}\DUrole{w}{ }\DUrole{n}{ndarray}}\sphinxparamcomma \sphinxparam{\DUrole{n}{alt}\DUrole{p}{:}\DUrole{w}{ }\DUrole{n}{ndarray}}\sphinxparamcomma \sphinxparam{\DUrole{n}{step}\DUrole{p}{:}\DUrole{w}{ }\DUrole{n}{float}\DUrole{w}{ }\DUrole{o}{=}\DUrole{w}{ }\DUrole{default_value}{10}}}{{ $\rightarrow$ ndarray}}
\pysigstopsignatures
\sphinxAtStartPar
In a vectorised batch, calculates the vertical gradient of gravity.
Computes the vertical change in gravity (dg/dr) of the Z/Down
component (in m/s\textasciicircum{}2). This is primary used for gravity map\sphinxhyphen{}matching
related calculations.
\begin{quote}\begin{description}
\sphinxlineitem{Parameters}\begin{itemize}
\item {} 
\sphinxAtStartPar
\sphinxstyleliteralstrong{\sphinxupquote{lat}} (\sphinxstyleliteralemphasis{\sphinxupquote{np.ndarray}}) \textendash{} The latitude coordinate for the position of interest.
The latitude position value given in decimal degrees.

\item {} 
\sphinxAtStartPar
\sphinxstyleliteralstrong{\sphinxupquote{lon}} (\sphinxstyleliteralemphasis{\sphinxupquote{np.ndarray}}) \textendash{} The longitude coordinate for the position of interest.
The longitude position value given in decimal degrees.

\item {} 
\sphinxAtStartPar
\sphinxstyleliteralstrong{\sphinxupquote{alt}} (\sphinxstyleliteralemphasis{\sphinxupquote{np.ndarray}}) \textendash{} The altitude offset for the position of interest.
The elevation from the WGS\sphinxhyphen{}84 reference ellipsoid in metres.

\item {} 
\sphinxAtStartPar
\sphinxstyleliteralstrong{\sphinxupquote{step}} (\sphinxstyleliteralemphasis{\sphinxupquote{float}}) \textendash{} The step size (in metres) for calculating the gradient.
To be used when the gradient cannot be computed directly.

\end{itemize}

\sphinxlineitem{Returns}
\sphinxAtStartPar
The vertical gravity gradient (in m/s\textasciicircum{}2).

\sphinxlineitem{Return type}
\sphinxAtStartPar
float

\end{description}\end{quote}

\end{fulllineitems}

\index{base\_model (qnav.gravity.base.GravityModel property)@\spxentry{base\_model}\spxextra{qnav.gravity.base.GravityModel property}}

\begin{fulllineitems}
\phantomsection\label{\detokenize{modules/gravity/base:qnav.gravity.base.GravityModel.base_model}}
\pysigstartsignatures
\pysigline{\sphinxbfcode{\sphinxupquote{property\DUrole{w}{ }}}\sphinxbfcode{\sphinxupquote{base\_model}}\sphinxbfcode{\sphinxupquote{\DUrole{p}{:}\DUrole{w}{ }Self\DUrole{w}{ }\DUrole{p}{|}\DUrole{w}{ }None}}}
\pysigstopsignatures
\sphinxAtStartPar
Returns the base gravity model the model extends from.
Allows access to the base function/model the instance is
using (if any).
\begin{quote}\begin{description}
\sphinxlineitem{Returns}
\sphinxAtStartPar
Either the base gravity model or None.

\sphinxlineitem{Return type}
\sphinxAtStartPar
{\hyperref[\detokenize{modules/gravity/base:qnav.gravity.base.GravityModel}]{\sphinxcrossref{GravityModel}}} | None

\end{description}\end{quote}

\end{fulllineitems}


\end{fulllineitems}

\index{calculate\_curvature() (in module qnav.gravity.base)@\spxentry{calculate\_curvature()}\spxextra{in module qnav.gravity.base}}

\begin{fulllineitems}
\phantomsection\label{\detokenize{modules/gravity/base:qnav.gravity.base.calculate_curvature}}
\pysigstartsignatures
\pysiglinewithargsret{\sphinxcode{\sphinxupquote{qnav.gravity.base.}}\sphinxbfcode{\sphinxupquote{calculate\_curvature}}}{\sphinxparam{\DUrole{n}{lat}\DUrole{p}{:}\DUrole{w}{ }\DUrole{n}{float\DUrole{w}{ }\DUrole{p}{|}\DUrole{w}{ }ndarray}}}{{ $\rightarrow$ float\DUrole{w}{ }\DUrole{p}{|}\DUrole{w}{ }ndarray}}
\pysigstopsignatures
\sphinxAtStartPar
Calculates the mean curvature of the ellipsoid J at given position.
Using equations from Physical Geodesy 2006, this function calculates
J for given positions.
\begin{quote}\begin{description}
\sphinxlineitem{Parameters}
\sphinxAtStartPar
\sphinxstyleliteralstrong{\sphinxupquote{lat}} (\sphinxstyleliteralemphasis{\sphinxupquote{float}}\sphinxstyleliteralemphasis{\sphinxupquote{ | }}\sphinxstyleliteralemphasis{\sphinxupquote{np.ndarray}}) \textendash{} The latitude of the point(s) of interest (in decimal degrees).

\sphinxlineitem{Returns}
\sphinxAtStartPar
The mean curvature of the ellipsoid.

\sphinxlineitem{Return type}
\sphinxAtStartPar
float | np.ndarray

\end{description}\end{quote}

\end{fulllineitems}


\sphinxstepscope


\subsubsection{qnav.gravity.simple}
\label{\detokenize{modules/gravity/simple:module-qnav.gravity.simple}}\label{\detokenize{modules/gravity/simple:qnav-gravity-simple}}\label{\detokenize{modules/gravity/simple::doc}}\index{module@\spxentry{module}!qnav.gravity.simple@\spxentry{qnav.gravity.simple}}\index{qnav.gravity.simple@\spxentry{qnav.gravity.simple}!module@\spxentry{module}}

\paragraph{simple.py}
\label{\detokenize{modules/gravity/simple:simple-py}}\begin{quote}\begin{description}
\sphinxlineitem{summary}
\sphinxAtStartPar
Implementation of simple gravity modelling classes and functions.
This class contains very simple gravity modelling functions that are
namely used for testing, replacing gravity with values that are fixed
or vary very little. Primarily used for cancelling out any impact of
gravity during experiments.

\sphinxlineitem{author}
\sphinxAtStartPar
Michael Wright \sphinxhyphen{} \sphinxhref{mailto:mjwright@liverpool.ac.uk}{mjwright@liverpool.ac.uk}

\sphinxlineitem{version}
\sphinxAtStartPar
1.0.0 \sphinxhyphen{} First implementation of this module.

\end{description}\end{quote}
\index{FixedValue (class in qnav.gravity.simple)@\spxentry{FixedValue}\spxextra{class in qnav.gravity.simple}}

\begin{fulllineitems}
\phantomsection\label{\detokenize{modules/gravity/simple:qnav.gravity.simple.FixedValue}}
\pysigstartsignatures
\pysigline{\sphinxbfcode{\sphinxupquote{class\DUrole{w}{ }}}\sphinxcode{\sphinxupquote{qnav.gravity.simple.}}\sphinxbfcode{\sphinxupquote{FixedValue}}}
\pysigstopsignatures
\sphinxAtStartPar
Sets gravity equal to fixed constant regardless of location.
Using this model, gravity will be fixed and the same constant
value will be returned.
\index{calc\_gravity\_z() (qnav.gravity.simple.FixedValue method)@\spxentry{calc\_gravity\_z()}\spxextra{qnav.gravity.simple.FixedValue method}}

\begin{fulllineitems}
\phantomsection\label{\detokenize{modules/gravity/simple:qnav.gravity.simple.FixedValue.calc_gravity_z}}
\pysigstartsignatures
\pysiglinewithargsret{\sphinxbfcode{\sphinxupquote{calc\_gravity\_z}}}{\sphinxparam{\DUrole{n}{lat}\DUrole{p}{:}\DUrole{w}{ }\DUrole{n}{float}\DUrole{w}{ }\DUrole{o}{=}\DUrole{w}{ }\DUrole{default_value}{None}}\sphinxparamcomma \sphinxparam{\DUrole{n}{lon}\DUrole{p}{:}\DUrole{w}{ }\DUrole{n}{float}\DUrole{w}{ }\DUrole{o}{=}\DUrole{w}{ }\DUrole{default_value}{None}}\sphinxparamcomma \sphinxparam{\DUrole{n}{alt}\DUrole{p}{:}\DUrole{w}{ }\DUrole{n}{float}\DUrole{w}{ }\DUrole{o}{=}\DUrole{w}{ }\DUrole{default_value}{None}}}{{ $\rightarrow$ float}}
\pysigstopsignatures
\sphinxAtStartPar
Calculates and returns the total gravitational acceleration.
Computes the gravity along the Z/Down axis (in m/s\textasciicircum{}2) for
the given location and altitude offset.
\begin{quote}\begin{description}
\sphinxlineitem{Parameters}\begin{itemize}
\item {} 
\sphinxAtStartPar
\sphinxstyleliteralstrong{\sphinxupquote{lat}} (\sphinxstyleliteralemphasis{\sphinxupquote{float}}) \textendash{} The latitude coordinate for the position of interest.
The latitude position value given in decimal degrees.

\item {} 
\sphinxAtStartPar
\sphinxstyleliteralstrong{\sphinxupquote{lon}} (\sphinxstyleliteralemphasis{\sphinxupquote{float}}) \textendash{} The longitude coordinate for the position of interest.
The longitude position value given in decimal degrees.

\item {} 
\sphinxAtStartPar
\sphinxstyleliteralstrong{\sphinxupquote{alt}} (\sphinxstyleliteralemphasis{\sphinxupquote{float}}) \textendash{} The altitude offset for the position of interest.
The elevation from the WGS\sphinxhyphen{}84 reference ellipsoid in metres.

\end{itemize}

\sphinxlineitem{Returns}
\sphinxAtStartPar
The downwards gravitational acceleration for each location.

\sphinxlineitem{Return type}
\sphinxAtStartPar
float

\end{description}\end{quote}

\end{fulllineitems}

\index{calc\_gravity\_z\_vec() (qnav.gravity.simple.FixedValue method)@\spxentry{calc\_gravity\_z\_vec()}\spxextra{qnav.gravity.simple.FixedValue method}}

\begin{fulllineitems}
\phantomsection\label{\detokenize{modules/gravity/simple:qnav.gravity.simple.FixedValue.calc_gravity_z_vec}}
\pysigstartsignatures
\pysiglinewithargsret{\sphinxbfcode{\sphinxupquote{calc\_gravity\_z\_vec}}}{\sphinxparam{\DUrole{n}{lat}\DUrole{p}{:}\DUrole{w}{ }\DUrole{n}{ndarray}}\sphinxparamcomma \sphinxparam{\DUrole{n}{lon}\DUrole{p}{:}\DUrole{w}{ }\DUrole{n}{ndarray}\DUrole{w}{ }\DUrole{o}{=}\DUrole{w}{ }\DUrole{default_value}{None}}\sphinxparamcomma \sphinxparam{\DUrole{n}{alt}\DUrole{p}{:}\DUrole{w}{ }\DUrole{n}{ndarray}\DUrole{w}{ }\DUrole{o}{=}\DUrole{w}{ }\DUrole{default_value}{None}}}{{ $\rightarrow$ ndarray}}
\pysigstopsignatures
\sphinxAtStartPar
In a vectorised batch, calculates and returns the gravitational
acceleration along the Z/Down axis (in m/s\textasciicircum{}2) for the given location
and altitude offset. Ideal when calculating gravity for a large
number of positions.
\begin{quote}\begin{description}
\sphinxlineitem{Parameters}\begin{itemize}
\item {} 
\sphinxAtStartPar
\sphinxstyleliteralstrong{\sphinxupquote{lat}} (\sphinxstyleliteralemphasis{\sphinxupquote{np.array}}) \textendash{} The latitude coordinate for the position of interest.
The latitude position value given in decimal degrees.

\item {} 
\sphinxAtStartPar
\sphinxstyleliteralstrong{\sphinxupquote{lon}} (\sphinxstyleliteralemphasis{\sphinxupquote{np.array}}) \textendash{} The longitude coordinate for the position of interest.
The longitude position value given in decimal degrees.

\item {} 
\sphinxAtStartPar
\sphinxstyleliteralstrong{\sphinxupquote{alt}} (\sphinxstyleliteralemphasis{\sphinxupquote{np.array}}) \textendash{} The altitude offset for the position of interest.
The elevation from the WGS\sphinxhyphen{}84 reference ellipsoid in metres.

\end{itemize}

\sphinxlineitem{Returns}
\sphinxAtStartPar
The downwards gravitational acceleration for each location
(in m/s\textasciicircum{}2).

\sphinxlineitem{Return type}
\sphinxAtStartPar
np.array (N\sphinxhyphen{}elements)

\end{description}\end{quote}

\end{fulllineitems}


\end{fulllineitems}

\index{SimpleUniform (class in qnav.gravity.simple)@\spxentry{SimpleUniform}\spxextra{class in qnav.gravity.simple}}

\begin{fulllineitems}
\phantomsection\label{\detokenize{modules/gravity/simple:qnav.gravity.simple.SimpleUniform}}
\pysigstartsignatures
\pysigline{\sphinxbfcode{\sphinxupquote{class\DUrole{w}{ }}}\sphinxcode{\sphinxupquote{qnav.gravity.simple.}}\sphinxbfcode{\sphinxupquote{SimpleUniform}}}
\pysigstopsignatures
\sphinxAtStartPar
Calculates gravity using uniformly across the Earth.
Produces values using weighted average of the gravity between the poles
and equator, meaning only latitude affected the calculations produced.
Changes in longitude or altitude will have no impact.
\index{calc\_gravity\_z() (qnav.gravity.simple.SimpleUniform method)@\spxentry{calc\_gravity\_z()}\spxextra{qnav.gravity.simple.SimpleUniform method}}

\begin{fulllineitems}
\phantomsection\label{\detokenize{modules/gravity/simple:qnav.gravity.simple.SimpleUniform.calc_gravity_z}}
\pysigstartsignatures
\pysiglinewithargsret{\sphinxbfcode{\sphinxupquote{calc\_gravity\_z}}}{\sphinxparam{\DUrole{n}{lat}\DUrole{p}{:}\DUrole{w}{ }\DUrole{n}{float}}\sphinxparamcomma \sphinxparam{\DUrole{n}{lon}\DUrole{p}{:}\DUrole{w}{ }\DUrole{n}{float}\DUrole{w}{ }\DUrole{o}{=}\DUrole{w}{ }\DUrole{default_value}{None}}\sphinxparamcomma \sphinxparam{\DUrole{n}{alt}\DUrole{p}{:}\DUrole{w}{ }\DUrole{n}{float}\DUrole{w}{ }\DUrole{o}{=}\DUrole{w}{ }\DUrole{default_value}{None}}}{{ $\rightarrow$ float}}
\pysigstopsignatures
\sphinxAtStartPar
Calculates and returns the total gravitational acceleration.
Computes the gravity along the Z/Down axis (in m/s\textasciicircum{}2) for
the given location and altitude offset.
\begin{quote}\begin{description}
\sphinxlineitem{Parameters}\begin{itemize}
\item {} 
\sphinxAtStartPar
\sphinxstyleliteralstrong{\sphinxupquote{lat}} (\sphinxstyleliteralemphasis{\sphinxupquote{float}}) \textendash{} The latitude coordinate for the position of interest.
The latitude position value given in decimal degrees.

\item {} 
\sphinxAtStartPar
\sphinxstyleliteralstrong{\sphinxupquote{lon}} (\sphinxstyleliteralemphasis{\sphinxupquote{float}}) \textendash{} The longitude coordinate for the position of interest.
The longitude position value given in decimal degrees.

\item {} 
\sphinxAtStartPar
\sphinxstyleliteralstrong{\sphinxupquote{alt}} (\sphinxstyleliteralemphasis{\sphinxupquote{float}}) \textendash{} The altitude offset for the position of interest.
The elevation from the WGS\sphinxhyphen{}84 reference ellipsoid in metres.

\end{itemize}

\sphinxlineitem{Returns}
\sphinxAtStartPar
The downwards gravitational acceleration for each location.

\sphinxlineitem{Return type}
\sphinxAtStartPar
float

\end{description}\end{quote}

\end{fulllineitems}

\index{calc\_gravity\_z\_vec() (qnav.gravity.simple.SimpleUniform method)@\spxentry{calc\_gravity\_z\_vec()}\spxextra{qnav.gravity.simple.SimpleUniform method}}

\begin{fulllineitems}
\phantomsection\label{\detokenize{modules/gravity/simple:qnav.gravity.simple.SimpleUniform.calc_gravity_z_vec}}
\pysigstartsignatures
\pysiglinewithargsret{\sphinxbfcode{\sphinxupquote{calc\_gravity\_z\_vec}}}{\sphinxparam{\DUrole{n}{lat}\DUrole{p}{:}\DUrole{w}{ }\DUrole{n}{ndarray}}\sphinxparamcomma \sphinxparam{\DUrole{n}{lon}\DUrole{p}{:}\DUrole{w}{ }\DUrole{n}{ndarray}\DUrole{w}{ }\DUrole{o}{=}\DUrole{w}{ }\DUrole{default_value}{None}}\sphinxparamcomma \sphinxparam{\DUrole{n}{alt}\DUrole{p}{:}\DUrole{w}{ }\DUrole{n}{ndarray}\DUrole{w}{ }\DUrole{o}{=}\DUrole{w}{ }\DUrole{default_value}{None}}}{{ $\rightarrow$ float}}
\pysigstopsignatures
\sphinxAtStartPar
In a vectorised batch, calculates and returns the gravitational
acceleration along the Z/Down axis (in m/s\textasciicircum{}2) for the given location
and altitude offset. Ideal when calculating gravity for a large
number of positions.
\begin{quote}\begin{description}
\sphinxlineitem{Parameters}\begin{itemize}
\item {} 
\sphinxAtStartPar
\sphinxstyleliteralstrong{\sphinxupquote{lat}} (\sphinxstyleliteralemphasis{\sphinxupquote{np.array}}) \textendash{} The latitude coordinate for the position of interest.
The latitude position value given in decimal degrees.

\item {} 
\sphinxAtStartPar
\sphinxstyleliteralstrong{\sphinxupquote{lon}} (\sphinxstyleliteralemphasis{\sphinxupquote{np.array}}) \textendash{} The longitude coordinate for the position of interest.
The longitude position value given in decimal degrees.

\item {} 
\sphinxAtStartPar
\sphinxstyleliteralstrong{\sphinxupquote{alt}} (\sphinxstyleliteralemphasis{\sphinxupquote{np.array}}) \textendash{} The altitude offset for the position of interest.
The elevation from the WGS\sphinxhyphen{}84 reference ellipsoid in metres.

\end{itemize}

\sphinxlineitem{Returns}
\sphinxAtStartPar
The downwards gravitational acceleration for each location
(in m/s\textasciicircum{}2).

\sphinxlineitem{Return type}
\sphinxAtStartPar
np.array (N\sphinxhyphen{}elements)

\end{description}\end{quote}

\end{fulllineitems}


\end{fulllineitems}


\sphinxstepscope


\subsubsection{qnav.gravity.somigliana}
\label{\detokenize{modules/gravity/somigliana:module-qnav.gravity.somigliana}}\label{\detokenize{modules/gravity/somigliana:qnav-gravity-somigliana}}\label{\detokenize{modules/gravity/somigliana::doc}}\index{module@\spxentry{module}!qnav.gravity.somigliana@\spxentry{qnav.gravity.somigliana}}\index{qnav.gravity.somigliana@\spxentry{qnav.gravity.somigliana}!module@\spxentry{module}}

\paragraph{somigliana.py}
\label{\detokenize{modules/gravity/somigliana:somigliana-py}}\begin{quote}\begin{description}
\sphinxlineitem{summary}
\sphinxAtStartPar
Implementation of the Somigliana using WGS\sphinxhyphen{}84 parameters.
This class contains gravity functions related to the Somigliana
formula (with height correction).

\sphinxlineitem{author}
\sphinxAtStartPar
Michael Wright \sphinxhyphen{} \sphinxhref{mailto:mjwright@liverpool.ac.uk}{mjwright@liverpool.ac.uk}

\sphinxlineitem{version}
\sphinxAtStartPar
1.0.0 \sphinxhyphen{} First implementation of this module.

\end{description}\end{quote}
\index{Somigliana (class in qnav.gravity.somigliana)@\spxentry{Somigliana}\spxextra{class in qnav.gravity.somigliana}}

\begin{fulllineitems}
\phantomsection\label{\detokenize{modules/gravity/somigliana:qnav.gravity.somigliana.Somigliana}}
\pysigstartsignatures
\pysigline{\sphinxbfcode{\sphinxupquote{class\DUrole{w}{ }}}\sphinxcode{\sphinxupquote{qnav.gravity.somigliana.}}\sphinxbfcode{\sphinxupquote{Somigliana}}}
\pysigstopsignatures
\sphinxAtStartPar
Gravity function implementation of the Somigliana formula.
This allows the calculation of the vertical gravity, gravitational
acceleration and gravity gradients using the Somigliana formula
(with altitude interpolation).
\index{calc\_gravity\_z() (qnav.gravity.somigliana.Somigliana method)@\spxentry{calc\_gravity\_z()}\spxextra{qnav.gravity.somigliana.Somigliana method}}

\begin{fulllineitems}
\phantomsection\label{\detokenize{modules/gravity/somigliana:qnav.gravity.somigliana.Somigliana.calc_gravity_z}}
\pysigstartsignatures
\pysiglinewithargsret{\sphinxbfcode{\sphinxupquote{calc\_gravity\_z}}}{\sphinxparam{\DUrole{n}{lat}\DUrole{p}{:}\DUrole{w}{ }\DUrole{n}{float}}\sphinxparamcomma \sphinxparam{\DUrole{n}{lon}\DUrole{p}{:}\DUrole{w}{ }\DUrole{n}{float}}\sphinxparamcomma \sphinxparam{\DUrole{n}{alt}\DUrole{p}{:}\DUrole{w}{ }\DUrole{n}{float}}}{{ $\rightarrow$ float}}
\pysigstopsignatures
\sphinxAtStartPar
Calculates and returns the gravitational acceleration along the Z/Down
axis (in m/s\textasciicircum{}2) for the given location and altitude offset.
\begin{quote}\begin{description}
\sphinxlineitem{Parameters}\begin{itemize}
\item {} 
\sphinxAtStartPar
\sphinxstyleliteralstrong{\sphinxupquote{lat}} (\sphinxstyleliteralemphasis{\sphinxupquote{float}}) \textendash{} The latitude coordinate for the position of interest.
The latitude position value given in decimal degrees.

\item {} 
\sphinxAtStartPar
\sphinxstyleliteralstrong{\sphinxupquote{lon}} (\sphinxstyleliteralemphasis{\sphinxupquote{float}}) \textendash{} The longitude coordinate for the position of interest.
The longitude position value given in decimal degrees.

\item {} 
\sphinxAtStartPar
\sphinxstyleliteralstrong{\sphinxupquote{alt}} (\sphinxstyleliteralemphasis{\sphinxupquote{float}}) \textendash{} The altitude offset for the position of interest.
The elevation from the WGS\sphinxhyphen{}84 reference ellipsoid in metres.

\end{itemize}

\sphinxlineitem{Returns}
\sphinxAtStartPar
The downwards gravitational acceleration for each location.

\sphinxlineitem{Return type}
\sphinxAtStartPar
float

\end{description}\end{quote}

\end{fulllineitems}

\index{calc\_gravity\_xyz() (qnav.gravity.somigliana.Somigliana method)@\spxentry{calc\_gravity\_xyz()}\spxextra{qnav.gravity.somigliana.Somigliana method}}

\begin{fulllineitems}
\phantomsection\label{\detokenize{modules/gravity/somigliana:qnav.gravity.somigliana.Somigliana.calc_gravity_xyz}}
\pysigstartsignatures
\pysiglinewithargsret{\sphinxbfcode{\sphinxupquote{calc\_gravity\_xyz}}}{\sphinxparam{\DUrole{n}{lat}\DUrole{p}{:}\DUrole{w}{ }\DUrole{n}{float}}\sphinxparamcomma \sphinxparam{\DUrole{n}{lon}\DUrole{p}{:}\DUrole{w}{ }\DUrole{n}{float}}\sphinxparamcomma \sphinxparam{\DUrole{n}{alt}\DUrole{p}{:}\DUrole{w}{ }\DUrole{n}{float}}}{{ $\rightarrow$ ndarray}}
\pysigstopsignatures
\sphinxAtStartPar
Calculates and returns the gravitational acceleration for the
X/North, Y/East and Z/Down axes (in m/s\textasciicircum{}2) for the given location
and altitude offset.
\begin{quote}\begin{description}
\sphinxlineitem{Parameters}\begin{itemize}
\item {} 
\sphinxAtStartPar
\sphinxstyleliteralstrong{\sphinxupquote{lat}} (\sphinxstyleliteralemphasis{\sphinxupquote{float}}) \textendash{} The latitude coordinate for the position of interest.
The latitude position value given in decimal degrees.

\item {} 
\sphinxAtStartPar
\sphinxstyleliteralstrong{\sphinxupquote{lon}} (\sphinxstyleliteralemphasis{\sphinxupquote{float}}) \textendash{} The longitude coordinate for the position of interest.
The longitude position value given in decimal degrees.

\item {} 
\sphinxAtStartPar
\sphinxstyleliteralstrong{\sphinxupquote{alt}} (\sphinxstyleliteralemphasis{\sphinxupquote{float}}) \textendash{} The altitude offset for the position of interest.
The elevation from the WGS\sphinxhyphen{}84 reference ellipsoid in metres.

\end{itemize}

\sphinxlineitem{Returns}
\sphinxAtStartPar
The local gravity vector for the XYZ components (in m/s\textasciicircum{}2).

\sphinxlineitem{Return type}
\sphinxAtStartPar
numpy.ndarray (3\sphinxhyphen{}elements)

\end{description}\end{quote}

\end{fulllineitems}

\index{calc\_gravity\_z\_vec() (qnav.gravity.somigliana.Somigliana method)@\spxentry{calc\_gravity\_z\_vec()}\spxextra{qnav.gravity.somigliana.Somigliana method}}

\begin{fulllineitems}
\phantomsection\label{\detokenize{modules/gravity/somigliana:qnav.gravity.somigliana.Somigliana.calc_gravity_z_vec}}
\pysigstartsignatures
\pysiglinewithargsret{\sphinxbfcode{\sphinxupquote{calc\_gravity\_z\_vec}}}{\sphinxparam{\DUrole{n}{lat}\DUrole{p}{:}\DUrole{w}{ }\DUrole{n}{ndarray}}\sphinxparamcomma \sphinxparam{\DUrole{n}{lon}\DUrole{p}{:}\DUrole{w}{ }\DUrole{n}{ndarray}}\sphinxparamcomma \sphinxparam{\DUrole{n}{alt}\DUrole{p}{:}\DUrole{w}{ }\DUrole{n}{ndarray}}}{{ $\rightarrow$ ndarray}}
\pysigstopsignatures
\sphinxAtStartPar
In a vectorised batch, calculates and returns the gravitational
acceleration along the Z/Down axis (in m/s\textasciicircum{}2) for the given location
and altitude offset. Ideal when calculating gravity for a large
number of positions.
\begin{quote}\begin{description}
\sphinxlineitem{Parameters}\begin{itemize}
\item {} 
\sphinxAtStartPar
\sphinxstyleliteralstrong{\sphinxupquote{lat}} (\sphinxstyleliteralemphasis{\sphinxupquote{np.array}}) \textendash{} The latitude coordinate for the position of interest.
The latitude position value given in decimal degrees.

\item {} 
\sphinxAtStartPar
\sphinxstyleliteralstrong{\sphinxupquote{lon}} (\sphinxstyleliteralemphasis{\sphinxupquote{np.array}}) \textendash{} The longitude coordinate for the position of interest.
The longitude position value given in decimal degrees.

\item {} 
\sphinxAtStartPar
\sphinxstyleliteralstrong{\sphinxupquote{alt}} (\sphinxstyleliteralemphasis{\sphinxupquote{np.array}}) \textendash{} The altitude offset for the position of interest.
The elevation from the WGS\sphinxhyphen{}84 reference ellipsoid in metres.

\end{itemize}

\sphinxlineitem{Returns}
\sphinxAtStartPar
The downwards gravitational acceleration for each location
(in m/s\textasciicircum{}2).

\sphinxlineitem{Return type}
\sphinxAtStartPar
np.array (N\sphinxhyphen{}elements)

\end{description}\end{quote}

\end{fulllineitems}

\index{calc\_gravity\_xyz\_vec() (qnav.gravity.somigliana.Somigliana method)@\spxentry{calc\_gravity\_xyz\_vec()}\spxextra{qnav.gravity.somigliana.Somigliana method}}

\begin{fulllineitems}
\phantomsection\label{\detokenize{modules/gravity/somigliana:qnav.gravity.somigliana.Somigliana.calc_gravity_xyz_vec}}
\pysigstartsignatures
\pysiglinewithargsret{\sphinxbfcode{\sphinxupquote{calc\_gravity\_xyz\_vec}}}{\sphinxparam{\DUrole{n}{lat}\DUrole{p}{:}\DUrole{w}{ }\DUrole{n}{ndarray}}\sphinxparamcomma \sphinxparam{\DUrole{n}{lon}\DUrole{p}{:}\DUrole{w}{ }\DUrole{n}{ndarray}}\sphinxparamcomma \sphinxparam{\DUrole{n}{alt}\DUrole{p}{:}\DUrole{w}{ }\DUrole{n}{ndarray}}}{{ $\rightarrow$ ndarray}}
\pysigstopsignatures
\sphinxAtStartPar
In a vectorised batch, calculates and returns the gravitational
acceleration for the X/North, Y/East and Z/Down axes (in m/s\textasciicircum{}2)
for the given location and altitude offset. Ideal when calculating
gravity for a large number of positions.
\begin{quote}\begin{description}
\sphinxlineitem{Parameters}\begin{itemize}
\item {} 
\sphinxAtStartPar
\sphinxstyleliteralstrong{\sphinxupquote{lat}} (\sphinxstyleliteralemphasis{\sphinxupquote{np.array}}) \textendash{} The latitude coordinate for the position of interest.
The latitude position value given in decimal degrees.

\item {} 
\sphinxAtStartPar
\sphinxstyleliteralstrong{\sphinxupquote{lon}} (\sphinxstyleliteralemphasis{\sphinxupquote{np.array}}) \textendash{} The longitude coordinate for the position of interest.
The longitude position value given in decimal degrees.

\item {} 
\sphinxAtStartPar
\sphinxstyleliteralstrong{\sphinxupquote{alt}} (\sphinxstyleliteralemphasis{\sphinxupquote{np.array}}) \textendash{} The altitude offset for the position of interest.
The elevation from the WGS\sphinxhyphen{}84 reference ellipsoid in metres.

\end{itemize}

\sphinxlineitem{Returns}
\sphinxAtStartPar
The local gravity vector for the XYZ components (in m/s\textasciicircum{}2).

\sphinxlineitem{Return type}
\sphinxAtStartPar
numpy.ndarray (N\sphinxhyphen{}by\sphinxhyphen{}3 elements)

\end{description}\end{quote}

\end{fulllineitems}


\end{fulllineitems}


\sphinxstepscope


\subsubsection{qnav.gravity.nima}
\label{\detokenize{modules/gravity/nima:module-qnav.gravity.nima}}\label{\detokenize{modules/gravity/nima:qnav-gravity-nima}}\label{\detokenize{modules/gravity/nima::doc}}\index{module@\spxentry{module}!qnav.gravity.nima@\spxentry{qnav.gravity.nima}}\index{qnav.gravity.nima@\spxentry{qnav.gravity.nima}!module@\spxentry{module}}

\paragraph{nima.py}
\label{\detokenize{modules/gravity/nima:nima-py}}\begin{quote}\begin{description}
\sphinxlineitem{summary}
\sphinxAtStartPar
National Imagery and Mapping Agency (NIMA) gravity functions.
This module contains all the gravity calculation methods published by
NIMA, specifically those stated in “Department of Defense World
Geodetic System 1984: Its Definition, and Relationship with Local
Geodetic Systems, TR8350.2, Third Ed.” Department of Defense,
Washington, DC: 1997.

\sphinxlineitem{author}
\sphinxAtStartPar
Michael Wright \sphinxhyphen{} \sphinxhref{mailto:mjwright@liverpool.ac.uk}{mjwright@liverpool.ac.uk}

\sphinxlineitem{version}
\sphinxAtStartPar
1.0.0 \sphinxhyphen{} First implementation of this module.

\end{description}\end{quote}
\index{WGS84Gravity (class in qnav.gravity.nima)@\spxentry{WGS84Gravity}\spxextra{class in qnav.gravity.nima}}

\begin{fulllineitems}
\phantomsection\label{\detokenize{modules/gravity/nima:qnav.gravity.nima.WGS84Gravity}}
\pysigstartsignatures
\pysigline{\sphinxbfcode{\sphinxupquote{class\DUrole{w}{ }}}\sphinxcode{\sphinxupquote{qnav.gravity.nima.}}\sphinxbfcode{\sphinxupquote{WGS84Gravity}}}
\pysigstopsignatures
\sphinxAtStartPar
Implementation of the WGS84 representation of Earth gravity.
This model uses the equations published by NIMA. This is known to be
a relatively accurate close approximation model for heights up to
20,000 meters above the WGS84 reference Ellipsoid.
\index{calc\_gravity\_z() (qnav.gravity.nima.WGS84Gravity method)@\spxentry{calc\_gravity\_z()}\spxextra{qnav.gravity.nima.WGS84Gravity method}}

\begin{fulllineitems}
\phantomsection\label{\detokenize{modules/gravity/nima:qnav.gravity.nima.WGS84Gravity.calc_gravity_z}}
\pysigstartsignatures
\pysiglinewithargsret{\sphinxbfcode{\sphinxupquote{calc\_gravity\_z}}}{\sphinxparam{\DUrole{n}{lat}\DUrole{p}{:}\DUrole{w}{ }\DUrole{n}{float}}\sphinxparamcomma \sphinxparam{\DUrole{n}{lon}\DUrole{p}{:}\DUrole{w}{ }\DUrole{n}{float}}\sphinxparamcomma \sphinxparam{\DUrole{n}{alt}\DUrole{p}{:}\DUrole{w}{ }\DUrole{n}{float}}}{{ $\rightarrow$ float}}
\pysigstopsignatures
\sphinxAtStartPar
Calculates and returns the gravitational acceleration along the Z/Down
axis (in m/s\textasciicircum{}2) for the given location and altitude offset.
\begin{quote}\begin{description}
\sphinxlineitem{Parameters}\begin{itemize}
\item {} 
\sphinxAtStartPar
\sphinxstyleliteralstrong{\sphinxupquote{lat}} (\sphinxstyleliteralemphasis{\sphinxupquote{float}}) \textendash{} The latitude coordinate for the position of interest.
The latitude position value given in decimal degrees.

\item {} 
\sphinxAtStartPar
\sphinxstyleliteralstrong{\sphinxupquote{lon}} (\sphinxstyleliteralemphasis{\sphinxupquote{float}}) \textendash{} The longitude coordinate for the position of interest.
The longitude position value given in decimal degrees.

\item {} 
\sphinxAtStartPar
\sphinxstyleliteralstrong{\sphinxupquote{alt}} (\sphinxstyleliteralemphasis{\sphinxupquote{float}}) \textendash{} The altitude offset for the position of interest.
The elevation from the WGS\sphinxhyphen{}84 reference ellipsoid in metres.

\end{itemize}

\sphinxlineitem{Returns}
\sphinxAtStartPar
The downwards gravitational acceleration for each location.

\sphinxlineitem{Return type}
\sphinxAtStartPar
float

\end{description}\end{quote}

\end{fulllineitems}

\index{calc\_gravity\_xyz() (qnav.gravity.nima.WGS84Gravity method)@\spxentry{calc\_gravity\_xyz()}\spxextra{qnav.gravity.nima.WGS84Gravity method}}

\begin{fulllineitems}
\phantomsection\label{\detokenize{modules/gravity/nima:qnav.gravity.nima.WGS84Gravity.calc_gravity_xyz}}
\pysigstartsignatures
\pysiglinewithargsret{\sphinxbfcode{\sphinxupquote{calc\_gravity\_xyz}}}{\sphinxparam{\DUrole{n}{lat}\DUrole{p}{:}\DUrole{w}{ }\DUrole{n}{float}}\sphinxparamcomma \sphinxparam{\DUrole{n}{lon}\DUrole{p}{:}\DUrole{w}{ }\DUrole{n}{float}}\sphinxparamcomma \sphinxparam{\DUrole{n}{alt}\DUrole{p}{:}\DUrole{w}{ }\DUrole{n}{float}}}{{ $\rightarrow$ ndarray}}
\pysigstopsignatures
\sphinxAtStartPar
Calculates and returns the gravitational acceleration for the
X/North, Y/East and Z/Down axes (in m/s\textasciicircum{}2) for the given location
and altitude offset.
\begin{quote}\begin{description}
\sphinxlineitem{Parameters}\begin{itemize}
\item {} 
\sphinxAtStartPar
\sphinxstyleliteralstrong{\sphinxupquote{lat}} (\sphinxstyleliteralemphasis{\sphinxupquote{float}}) \textendash{} The latitude coordinate for the position of interest.
The latitude position value given in decimal degrees.

\item {} 
\sphinxAtStartPar
\sphinxstyleliteralstrong{\sphinxupquote{lon}} (\sphinxstyleliteralemphasis{\sphinxupquote{float}}) \textendash{} The longitude coordinate for the position of interest.
The longitude position value given in decimal degrees.

\item {} 
\sphinxAtStartPar
\sphinxstyleliteralstrong{\sphinxupquote{alt}} (\sphinxstyleliteralemphasis{\sphinxupquote{float}}) \textendash{} The altitude offset for the position of interest.
The elevation from the WGS\sphinxhyphen{}84 reference ellipsoid in metres.

\end{itemize}

\sphinxlineitem{Returns}
\sphinxAtStartPar
The local gravity vector for the XYZ components (in m/s\textasciicircum{}2).

\sphinxlineitem{Return type}
\sphinxAtStartPar
numpy.ndarray (3\sphinxhyphen{}elements)

\end{description}\end{quote}

\end{fulllineitems}


\end{fulllineitems}


\sphinxstepscope


\subsubsection{qnav.gravity.map}
\label{\detokenize{modules/gravity/map:module-qnav.gravity.map}}\label{\detokenize{modules/gravity/map:qnav-gravity-map}}\label{\detokenize{modules/gravity/map::doc}}\index{module@\spxentry{module}!qnav.gravity.map@\spxentry{qnav.gravity.map}}\index{qnav.gravity.map@\spxentry{qnav.gravity.map}!module@\spxentry{module}}

\paragraph{map.py}
\label{\detokenize{modules/gravity/map:map-py}}\begin{quote}\begin{description}
\sphinxlineitem{summary}
\sphinxAtStartPar
Module provides template implementation for gravity correction maps.
Due to their high resolution and spatial complexity, methods from
the original gravity model base class have had to be adjusted.
The largest feature provided by this module is the method for
correctly approximating the vertical gravity gradient.

\sphinxlineitem{authors}
\begin{DUlineblock}{0em}
\item[] Prof. Jason Ralph \sphinxhyphen{} \sphinxhref{mailto:jfralph@liverpool.ac.uk}{jfralph@liverpool.ac.uk}
\item[] Michael Wright \sphinxhyphen{} \sphinxhref{mailto:mjwright@liverpool.ac.uk}{mjwright@liverpool.ac.uk}
\end{DUlineblock}

\sphinxlineitem{version}
\sphinxAtStartPar
1.0.0 \sphinxhyphen{} First implementation of this module.

\end{description}\end{quote}
\index{CircularGrid (class in qnav.gravity.map)@\spxentry{CircularGrid}\spxextra{class in qnav.gravity.map}}

\begin{fulllineitems}
\phantomsection\label{\detokenize{modules/gravity/map:qnav.gravity.map.CircularGrid}}
\pysigstartsignatures
\pysiglinewithargsret{\sphinxbfcode{\sphinxupquote{class\DUrole{w}{ }}}\sphinxcode{\sphinxupquote{qnav.gravity.map.}}\sphinxbfcode{\sphinxupquote{CircularGrid}}}{\sphinxparam{\DUrole{n}{num\_rings}\DUrole{p}{:}\DUrole{w}{ }\DUrole{n}{int}\DUrole{w}{ }\DUrole{o}{=}\DUrole{w}{ }\DUrole{default_value}{100}}\sphinxparamcomma \sphinxparam{\DUrole{n}{num\_angles}\DUrole{p}{:}\DUrole{w}{ }\DUrole{n}{int}\DUrole{w}{ }\DUrole{o}{=}\DUrole{w}{ }\DUrole{default_value}{50}}\sphinxparamcomma \sphinxparam{\DUrole{n}{max\_radius}\DUrole{p}{:}\DUrole{w}{ }\DUrole{n}{float}\DUrole{w}{ }\DUrole{o}{=}\DUrole{w}{ }\DUrole{default_value}{0.25}}\sphinxparamcomma \sphinxparam{\DUrole{n}{l\_threshold}\DUrole{p}{:}\DUrole{w}{ }\DUrole{n}{float}\DUrole{w}{ }\DUrole{o}{=}\DUrole{w}{ }\DUrole{default_value}{1.0}}}{}
\pysigstopsignatures
\sphinxAtStartPar
Definition for a circular reference grid from a centre position.
Circular grid points used for involving the surrounding data for
accurately computing the vertical gravity gradient.
\index{generate\_grid() (qnav.gravity.map.CircularGrid method)@\spxentry{generate\_grid()}\spxextra{qnav.gravity.map.CircularGrid method}}

\begin{fulllineitems}
\phantomsection\label{\detokenize{modules/gravity/map:qnav.gravity.map.CircularGrid.generate_grid}}
\pysigstartsignatures
\pysiglinewithargsret{\sphinxbfcode{\sphinxupquote{generate\_grid}}}{\sphinxparam{\DUrole{n}{lat}\DUrole{p}{:}\DUrole{w}{ }\DUrole{n}{float}}\sphinxparamcomma \sphinxparam{\DUrole{n}{lon}\DUrole{p}{:}\DUrole{w}{ }\DUrole{n}{float}}\sphinxparamcomma \sphinxparam{\DUrole{n}{alt}\DUrole{p}{:}\DUrole{w}{ }\DUrole{n}{float}}}{{ $\rightarrow$ dict\DUrole{p}{{[}}str\DUrole{p}{,}\DUrole{w}{ }ndarray\DUrole{p}{{]}}}}
\pysigstopsignatures
\sphinxAtStartPar
Generates a circular grid of query positions surrounding a point.
On succession, a dictionary containing the positional and
angular information of points surrounding a position of interest are
returned. This also includes commonly calculated values step size
values.
\begin{quote}\begin{description}
\sphinxlineitem{Parameters}\begin{itemize}
\item {} 
\sphinxAtStartPar
\sphinxstyleliteralstrong{\sphinxupquote{lat}} (\sphinxstyleliteralemphasis{\sphinxupquote{float}}) \textendash{} The latitude position of the point (in decimal degrees).

\item {} 
\sphinxAtStartPar
\sphinxstyleliteralstrong{\sphinxupquote{lon}} (\sphinxstyleliteralemphasis{\sphinxupquote{float}}) \textendash{} The longitude position of the point (in decimal degrees).

\item {} 
\sphinxAtStartPar
\sphinxstyleliteralstrong{\sphinxupquote{alt}} (\sphinxstyleliteralemphasis{\sphinxupquote{float}}) \textendash{} The altitude position of the point (in metres).

\end{itemize}

\sphinxlineitem{Returns}
\sphinxAtStartPar

\sphinxAtStartPar
A dictionary holding the calculated positions (lla),
NED offsets (ned), azimuth angles (alpha), distances from centre
\begin{quote}

\sphinxAtStartPar
(psi) and step sizes (delta\_psi and delta\_alpha).
\end{quote}


\sphinxlineitem{Return type}
\sphinxAtStartPar
dict{[}str, np.ndarray{]}

\end{description}\end{quote}

\end{fulllineitems}


\end{fulllineitems}

\index{set\_default\_grid() (in module qnav.gravity.map)@\spxentry{set\_default\_grid()}\spxextra{in module qnav.gravity.map}}

\begin{fulllineitems}
\phantomsection\label{\detokenize{modules/gravity/map:qnav.gravity.map.set_default_grid}}
\pysigstartsignatures
\pysiglinewithargsret{\sphinxcode{\sphinxupquote{qnav.gravity.map.}}\sphinxbfcode{\sphinxupquote{set\_default\_grid}}}{\sphinxparam{\DUrole{n}{num\_rings}\DUrole{p}{:}\DUrole{w}{ }\DUrole{n}{int}}\sphinxparamcomma \sphinxparam{\DUrole{n}{num\_angles}\DUrole{p}{:}\DUrole{w}{ }\DUrole{n}{int}}\sphinxparamcomma \sphinxparam{\DUrole{n}{max\_radius}\DUrole{p}{:}\DUrole{w}{ }\DUrole{n}{float}}\sphinxparamcomma \sphinxparam{\DUrole{n}{l\_threshold}\DUrole{p}{:}\DUrole{w}{ }\DUrole{n}{float}}}{}
\pysigstopsignatures
\sphinxAtStartPar
Overwrites the settings to use for the default circular reference grid.
This is the grid that will be used by default for calculating the
vertical gravity gradient.
\begin{quote}\begin{description}
\sphinxlineitem{Parameters}\begin{itemize}
\item {} 
\sphinxAtStartPar
\sphinxstyleliteralstrong{\sphinxupquote{num\_rings}} (\sphinxstyleliteralemphasis{\sphinxupquote{int}}) \textendash{} Number of radial distance steps.

\item {} 
\sphinxAtStartPar
\sphinxstyleliteralstrong{\sphinxupquote{num\_angles}} (\sphinxstyleliteralemphasis{\sphinxupquote{int}}) \textendash{} Number of azimuth angle steps.

\item {} 
\sphinxAtStartPar
\sphinxstyleliteralstrong{\sphinxupquote{max\_radius}} (\sphinxstyleliteralemphasis{\sphinxupquote{float}}) \textendash{} Maximum radial distance (in degrees).

\item {} 
\sphinxAtStartPar
\sphinxstyleliteralstrong{\sphinxupquote{l\_threshold}} (\sphinxstyleliteralemphasis{\sphinxupquote{float}}) \textendash{} Threshold for internal points (in metres).

\end{itemize}

\end{description}\end{quote}

\end{fulllineitems}

\index{GravityMap (class in qnav.gravity.map)@\spxentry{GravityMap}\spxextra{class in qnav.gravity.map}}

\begin{fulllineitems}
\phantomsection\label{\detokenize{modules/gravity/map:qnav.gravity.map.GravityMap}}
\pysigstartsignatures
\pysiglinewithargsret{\sphinxbfcode{\sphinxupquote{class\DUrole{w}{ }}}\sphinxcode{\sphinxupquote{qnav.gravity.map.}}\sphinxbfcode{\sphinxupquote{GravityMap}}}{\sphinxparam{\DUrole{n}{base\_model}\DUrole{p}{:}\DUrole{w}{ }\DUrole{n}{{\hyperref[\detokenize{modules/gravity/base:qnav.gravity.base.GravityModel}]{\sphinxcrossref{GravityModel}}}}}\sphinxparamcomma \sphinxparam{\DUrole{n}{geoid\_model}\DUrole{p}{:}\DUrole{w}{ }\DUrole{n}{{\hyperref[\detokenize{modules/earth/geoid:qnav.earth.geoid.GeoidModel}]{\sphinxcrossref{GeoidModel}}}\DUrole{w}{ }\DUrole{p}{|}\DUrole{w}{ }None}}\sphinxparamcomma \sphinxparam{\DUrole{n}{map\_dir}\DUrole{p}{:}\DUrole{w}{ }\DUrole{n}{Path\DUrole{w}{ }\DUrole{p}{|}\DUrole{w}{ }None}}\sphinxparamcomma \sphinxparam{\DUrole{n}{auto\_download}\DUrole{p}{:}\DUrole{w}{ }\DUrole{n}{bool}\DUrole{w}{ }\DUrole{o}{=}\DUrole{w}{ }\DUrole{default_value}{True}}}{}
\pysigstopsignatures
\sphinxAtStartPar
An extension to the abstract base class GravityModel that provides
re\sphinxhyphen{}worked methods for handling high\sphinxhyphen{}resolution lookup maps.
\index{get\_anomaly() (qnav.gravity.map.GravityMap method)@\spxentry{get\_anomaly()}\spxextra{qnav.gravity.map.GravityMap method}}

\begin{fulllineitems}
\phantomsection\label{\detokenize{modules/gravity/map:qnav.gravity.map.GravityMap.get_anomaly}}
\pysigstartsignatures
\pysiglinewithargsret{\sphinxbfcode{\sphinxupquote{abstract\DUrole{w}{ }}}\sphinxbfcode{\sphinxupquote{get\_anomaly}}}{\sphinxparam{\DUrole{n}{lat}\DUrole{p}{:}\DUrole{w}{ }\DUrole{n}{float}}\sphinxparamcomma \sphinxparam{\DUrole{n}{lon}\DUrole{p}{:}\DUrole{w}{ }\DUrole{n}{float}}}{{ $\rightarrow$ float}}
\pysigstopsignatures
\sphinxAtStartPar
Calculates the gravity anomaly for given position.
Retrieves the gravity anomaly value for the requested position
(in m/s\textasciicircum{}2) at the base surface. For gravity maps this value is
interpolated from its dataset, following possible conversion.
\begin{quote}\begin{description}
\sphinxlineitem{Parameters}\begin{itemize}
\item {} 
\sphinxAtStartPar
\sphinxstyleliteralstrong{\sphinxupquote{lat}} (\sphinxstyleliteralemphasis{\sphinxupquote{float}}) \textendash{} The latitude coordinate for the position of interest.
The latitude position value given in decimal degrees.

\item {} 
\sphinxAtStartPar
\sphinxstyleliteralstrong{\sphinxupquote{lon}} (\sphinxstyleliteralemphasis{\sphinxupquote{float}}) \textendash{} The longitude coordinate for the position of interest.
The longitude position value given in decimal degrees.

\end{itemize}

\sphinxlineitem{Returns}
\sphinxAtStartPar
The gravity anomaly value (in m/s\textasciicircum{}2).

\sphinxlineitem{Return type}
\sphinxAtStartPar
float

\end{description}\end{quote}

\end{fulllineitems}

\index{get\_disturbance() (qnav.gravity.map.GravityMap method)@\spxentry{get\_disturbance()}\spxextra{qnav.gravity.map.GravityMap method}}

\begin{fulllineitems}
\phantomsection\label{\detokenize{modules/gravity/map:qnav.gravity.map.GravityMap.get_disturbance}}
\pysigstartsignatures
\pysiglinewithargsret{\sphinxbfcode{\sphinxupquote{abstract\DUrole{w}{ }}}\sphinxbfcode{\sphinxupquote{get\_disturbance}}}{\sphinxparam{\DUrole{n}{lat}\DUrole{p}{:}\DUrole{w}{ }\DUrole{n}{float}}\sphinxparamcomma \sphinxparam{\DUrole{n}{lon}\DUrole{p}{:}\DUrole{w}{ }\DUrole{n}{float}}}{{ $\rightarrow$ float}}
\pysigstopsignatures
\sphinxAtStartPar
Calculates the gravity disturbance for given position.
Retrieves the gravity disturbance value for the requested position
(in m/s\textasciicircum{}2) at the base surface. For gravity maps this value is
interpolated from its dataset, following possible conversion.
\begin{quote}\begin{description}
\sphinxlineitem{Parameters}\begin{itemize}
\item {} 
\sphinxAtStartPar
\sphinxstyleliteralstrong{\sphinxupquote{lat}} (\sphinxstyleliteralemphasis{\sphinxupquote{float}}) \textendash{} The latitude coordinate for the position of interest.
The latitude position value given in decimal degrees.

\item {} 
\sphinxAtStartPar
\sphinxstyleliteralstrong{\sphinxupquote{lon}} (\sphinxstyleliteralemphasis{\sphinxupquote{float}}) \textendash{} The longitude coordinate for the position of interest.
The longitude position value given in decimal degrees.

\end{itemize}

\sphinxlineitem{Returns}
\sphinxAtStartPar
The gravity disturbance value (in m/s\textasciicircum{}2).

\sphinxlineitem{Return type}
\sphinxAtStartPar
float

\end{description}\end{quote}

\end{fulllineitems}

\index{get\_anomaly\_vec() (qnav.gravity.map.GravityMap method)@\spxentry{get\_anomaly\_vec()}\spxextra{qnav.gravity.map.GravityMap method}}

\begin{fulllineitems}
\phantomsection\label{\detokenize{modules/gravity/map:qnav.gravity.map.GravityMap.get_anomaly_vec}}
\pysigstartsignatures
\pysiglinewithargsret{\sphinxbfcode{\sphinxupquote{abstract\DUrole{w}{ }}}\sphinxbfcode{\sphinxupquote{get\_anomaly\_vec}}}{\sphinxparam{\DUrole{n}{lat}\DUrole{p}{:}\DUrole{w}{ }\DUrole{n}{ndarray}}\sphinxparamcomma \sphinxparam{\DUrole{n}{lon}\DUrole{p}{:}\DUrole{w}{ }\DUrole{n}{ndarray}}}{{ $\rightarrow$ ndarray}}
\pysigstopsignatures
\sphinxAtStartPar
Calculates the gravity anomaly for multiple given positions.
Retrieves the gravity anomaly value for the requested position
(in m/s\textasciicircum{}2) at the base surface. For gravity maps this value is
interpolated from its dataset, following possible conversion.
\begin{quote}\begin{description}
\sphinxlineitem{Parameters}\begin{itemize}
\item {} 
\sphinxAtStartPar
\sphinxstyleliteralstrong{\sphinxupquote{lat}} (\sphinxstyleliteralemphasis{\sphinxupquote{float}}) \textendash{} The latitude coordinate for the position of interest.
The latitude position value given in decimal degrees.

\item {} 
\sphinxAtStartPar
\sphinxstyleliteralstrong{\sphinxupquote{lon}} (\sphinxstyleliteralemphasis{\sphinxupquote{float}}) \textendash{} The longitude coordinate for the position of interest.
The longitude position value given in decimal degrees.

\end{itemize}

\sphinxlineitem{Returns}
\sphinxAtStartPar
The gravity anomaly value (in m/s\textasciicircum{}2).

\sphinxlineitem{Return type}
\sphinxAtStartPar
float

\end{description}\end{quote}

\end{fulllineitems}

\index{get\_disturbance\_vec() (qnav.gravity.map.GravityMap method)@\spxentry{get\_disturbance\_vec()}\spxextra{qnav.gravity.map.GravityMap method}}

\begin{fulllineitems}
\phantomsection\label{\detokenize{modules/gravity/map:qnav.gravity.map.GravityMap.get_disturbance_vec}}
\pysigstartsignatures
\pysiglinewithargsret{\sphinxbfcode{\sphinxupquote{abstract\DUrole{w}{ }}}\sphinxbfcode{\sphinxupquote{get\_disturbance\_vec}}}{\sphinxparam{\DUrole{n}{lat}\DUrole{p}{:}\DUrole{w}{ }\DUrole{n}{ndarray}}\sphinxparamcomma \sphinxparam{\DUrole{n}{lon}\DUrole{p}{:}\DUrole{w}{ }\DUrole{n}{ndarray}}}{{ $\rightarrow$ ndarray}}
\pysigstopsignatures
\sphinxAtStartPar
Calculates the gravity disturbance for multiple given positions.
Retrieves the gravity disturbance value for the requested position
(in m/s\textasciicircum{}2) at the base surface. For gravity maps this value is
interpolated from its dataset, following possible conversion.
\begin{quote}\begin{description}
\sphinxlineitem{Parameters}\begin{itemize}
\item {} 
\sphinxAtStartPar
\sphinxstyleliteralstrong{\sphinxupquote{lat}} (\sphinxstyleliteralemphasis{\sphinxupquote{float}}) \textendash{} The latitude coordinate for the position of interest.
The latitude position value given in decimal degrees.

\item {} 
\sphinxAtStartPar
\sphinxstyleliteralstrong{\sphinxupquote{lon}} (\sphinxstyleliteralemphasis{\sphinxupquote{float}}) \textendash{} The longitude coordinate for the position of interest.
The longitude position value given in decimal degrees.

\end{itemize}

\sphinxlineitem{Returns}
\sphinxAtStartPar
The gravity disturbance value (in m/s\textasciicircum{}2).

\sphinxlineitem{Return type}
\sphinxAtStartPar
float

\end{description}\end{quote}

\end{fulllineitems}

\index{calc\_gravity\_xyz() (qnav.gravity.map.GravityMap method)@\spxentry{calc\_gravity\_xyz()}\spxextra{qnav.gravity.map.GravityMap method}}

\begin{fulllineitems}
\phantomsection\label{\detokenize{modules/gravity/map:qnav.gravity.map.GravityMap.calc_gravity_xyz}}
\pysigstartsignatures
\pysiglinewithargsret{\sphinxbfcode{\sphinxupquote{calc\_gravity\_xyz}}}{\sphinxparam{\DUrole{n}{lat}\DUrole{p}{:}\DUrole{w}{ }\DUrole{n}{float}}\sphinxparamcomma \sphinxparam{\DUrole{n}{lon}\DUrole{p}{:}\DUrole{w}{ }\DUrole{n}{float}}\sphinxparamcomma \sphinxparam{\DUrole{n}{alt}\DUrole{p}{:}\DUrole{w}{ }\DUrole{n}{float}}}{{ $\rightarrow$ ndarray}}
\pysigstopsignatures
\sphinxAtStartPar
Calculates and returns the gravitational acceleration vector.
Computes the gravity acceleration X/North, Y/East and Z/Down
components (in m/s\textasciicircum{}2) for the given location and altitude offset.
\begin{quote}\begin{description}
\sphinxlineitem{Parameters}\begin{itemize}
\item {} 
\sphinxAtStartPar
\sphinxstyleliteralstrong{\sphinxupquote{lat}} (\sphinxstyleliteralemphasis{\sphinxupquote{float}}) \textendash{} The latitude coordinate for the position of interest.
The latitude position value given in decimal degrees.

\item {} 
\sphinxAtStartPar
\sphinxstyleliteralstrong{\sphinxupquote{lon}} (\sphinxstyleliteralemphasis{\sphinxupquote{float}}) \textendash{} The longitude coordinate for the position of interest.
The longitude position value given in decimal degrees.

\item {} 
\sphinxAtStartPar
\sphinxstyleliteralstrong{\sphinxupquote{alt}} (\sphinxstyleliteralemphasis{\sphinxupquote{float}}) \textendash{} The altitude offset for the position of interest.
The elevation from the WGS\sphinxhyphen{}84 reference ellipsoid in metres.

\end{itemize}

\sphinxlineitem{Returns}
\sphinxAtStartPar
The local gravity vector for the XYZ components (in m/s\textasciicircum{}2).

\sphinxlineitem{Return type}
\sphinxAtStartPar
numpy.ndarray (3\sphinxhyphen{}elements)

\end{description}\end{quote}

\end{fulllineitems}

\index{calc\_gravity\_xyz\_vec() (qnav.gravity.map.GravityMap method)@\spxentry{calc\_gravity\_xyz\_vec()}\spxextra{qnav.gravity.map.GravityMap method}}

\begin{fulllineitems}
\phantomsection\label{\detokenize{modules/gravity/map:qnav.gravity.map.GravityMap.calc_gravity_xyz_vec}}
\pysigstartsignatures
\pysiglinewithargsret{\sphinxbfcode{\sphinxupquote{calc\_gravity\_xyz\_vec}}}{\sphinxparam{\DUrole{n}{lat}\DUrole{p}{:}\DUrole{w}{ }\DUrole{n}{ndarray}}\sphinxparamcomma \sphinxparam{\DUrole{n}{lon}\DUrole{p}{:}\DUrole{w}{ }\DUrole{n}{ndarray}}\sphinxparamcomma \sphinxparam{\DUrole{n}{alt}\DUrole{p}{:}\DUrole{w}{ }\DUrole{n}{ndarray}}}{{ $\rightarrow$ ndarray}}
\pysigstopsignatures
\sphinxAtStartPar
In a vectorised batch, calculates and returns the gravitational
acceleration for the X/North, Y/East and Z/Down axes (in m/s\textasciicircum{}2)
for the given location and altitude offset. Ideal when calculating
gravity for a large number of positions.
\begin{quote}\begin{description}
\sphinxlineitem{Parameters}\begin{itemize}
\item {} 
\sphinxAtStartPar
\sphinxstyleliteralstrong{\sphinxupquote{lat}} (\sphinxstyleliteralemphasis{\sphinxupquote{np.array}}) \textendash{} The latitude coordinate for the position of interest.
The latitude position value given in decimal degrees.

\item {} 
\sphinxAtStartPar
\sphinxstyleliteralstrong{\sphinxupquote{lon}} (\sphinxstyleliteralemphasis{\sphinxupquote{np.array}}) \textendash{} The longitude coordinate for the position of interest.
The longitude position value given in decimal degrees.

\item {} 
\sphinxAtStartPar
\sphinxstyleliteralstrong{\sphinxupquote{alt}} (\sphinxstyleliteralemphasis{\sphinxupquote{np.array}}) \textendash{} The altitude offset for the position of interest.
The elevation from the WGS\sphinxhyphen{}84 reference ellipsoid in metres.

\end{itemize}

\sphinxlineitem{Returns}
\sphinxAtStartPar
The local gravity vector for the XYZ components (in m/s\textasciicircum{}2).

\sphinxlineitem{Return type}
\sphinxAtStartPar
numpy.ndarray (N\sphinxhyphen{}by\sphinxhyphen{}3 elements)

\end{description}\end{quote}

\end{fulllineitems}

\index{calc\_vertical\_grad() (qnav.gravity.map.GravityMap method)@\spxentry{calc\_vertical\_grad()}\spxextra{qnav.gravity.map.GravityMap method}}

\begin{fulllineitems}
\phantomsection\label{\detokenize{modules/gravity/map:qnav.gravity.map.GravityMap.calc_vertical_grad}}
\pysigstartsignatures
\pysiglinewithargsret{\sphinxbfcode{\sphinxupquote{calc\_vertical\_grad}}}{\sphinxparam{\DUrole{n}{lat}\DUrole{p}{:}\DUrole{w}{ }\DUrole{n}{float}}\sphinxparamcomma \sphinxparam{\DUrole{n}{lon}\DUrole{p}{:}\DUrole{w}{ }\DUrole{n}{float}}\sphinxparamcomma \sphinxparam{\DUrole{n}{alt}\DUrole{p}{:}\DUrole{w}{ }\DUrole{n}{float}}\sphinxparamcomma \sphinxparam{\DUrole{n}{step}\DUrole{p}{:}\DUrole{w}{ }\DUrole{n}{float}\DUrole{w}{ }\DUrole{o}{=}\DUrole{w}{ }\DUrole{default_value}{100}}\sphinxparamcomma \sphinxparam{\DUrole{n}{circular\_grid}\DUrole{p}{:}\DUrole{w}{ }\DUrole{n}{{\hyperref[\detokenize{modules/gravity/map:qnav.gravity.map.CircularGrid}]{\sphinxcrossref{CircularGrid}}}}\DUrole{w}{ }\DUrole{o}{=}\DUrole{w}{ }\DUrole{default_value}{None}}}{{ $\rightarrow$ float}}
\pysigstopsignatures
\sphinxAtStartPar
Calculates the vertical gradient of gravity.
Computes the vertical change in gravity (dg/dr) of the Z/Down
component (in m/s\textasciicircum{}2). This is primary used for gravity map\sphinxhyphen{}matching
related calculations.
\begin{quote}\begin{description}
\sphinxlineitem{Parameters}\begin{itemize}
\item {} 
\sphinxAtStartPar
\sphinxstyleliteralstrong{\sphinxupquote{lat}} (\sphinxstyleliteralemphasis{\sphinxupquote{float}}) \textendash{} The latitude coordinate for the position of interest.
The latitude position value given in decimal degrees.

\item {} 
\sphinxAtStartPar
\sphinxstyleliteralstrong{\sphinxupquote{lon}} (\sphinxstyleliteralemphasis{\sphinxupquote{float}}) \textendash{} The longitude coordinate for the position of interest.
The longitude position value given in decimal degrees.

\item {} 
\sphinxAtStartPar
\sphinxstyleliteralstrong{\sphinxupquote{alt}} (\sphinxstyleliteralemphasis{\sphinxupquote{float}}) \textendash{} The altitude offset for the position of interest.
The elevation from the WGS\sphinxhyphen{}84 reference ellipsoid in metres.

\item {} 
\sphinxAtStartPar
\sphinxstyleliteralstrong{\sphinxupquote{step}} (\sphinxstyleliteralemphasis{\sphinxupquote{float}}) \textendash{} The step size (in metres) for calculating the gradient.
To be used when the gradient cannot be computed directly.

\end{itemize}

\sphinxlineitem{Returns}
\sphinxAtStartPar
The vertical gravity gradient (in m/s\textasciicircum{}2).

\sphinxlineitem{Return type}
\sphinxAtStartPar
float

\end{description}\end{quote}

\end{fulllineitems}

\index{calc\_vertical\_grad\_vec() (qnav.gravity.map.GravityMap method)@\spxentry{calc\_vertical\_grad\_vec()}\spxextra{qnav.gravity.map.GravityMap method}}

\begin{fulllineitems}
\phantomsection\label{\detokenize{modules/gravity/map:qnav.gravity.map.GravityMap.calc_vertical_grad_vec}}
\pysigstartsignatures
\pysiglinewithargsret{\sphinxbfcode{\sphinxupquote{calc\_vertical\_grad\_vec}}}{\sphinxparam{\DUrole{n}{lat}\DUrole{p}{:}\DUrole{w}{ }\DUrole{n}{ndarray}}\sphinxparamcomma \sphinxparam{\DUrole{n}{lon}\DUrole{p}{:}\DUrole{w}{ }\DUrole{n}{ndarray}}\sphinxparamcomma \sphinxparam{\DUrole{n}{alt}\DUrole{p}{:}\DUrole{w}{ }\DUrole{n}{ndarray}}\sphinxparamcomma \sphinxparam{\DUrole{n}{step}\DUrole{p}{:}\DUrole{w}{ }\DUrole{n}{float}\DUrole{w}{ }\DUrole{o}{=}\DUrole{w}{ }\DUrole{default_value}{100}}\sphinxparamcomma \sphinxparam{\DUrole{n}{circular\_grid}\DUrole{p}{:}\DUrole{w}{ }\DUrole{n}{{\hyperref[\detokenize{modules/gravity/map:qnav.gravity.map.CircularGrid}]{\sphinxcrossref{CircularGrid}}}}\DUrole{w}{ }\DUrole{o}{=}\DUrole{w}{ }\DUrole{default_value}{None}}}{{ $\rightarrow$ ndarray}}
\pysigstopsignatures
\sphinxAtStartPar
In a vectorised batch, calculates the vertical gradient of gravity.
Computes the vertical change in gravity (dg/dr) of the Z/Down
component (in m/s\textasciicircum{}2). This is primary used for gravity map\sphinxhyphen{}matching
related calculations.
\begin{quote}\begin{description}
\sphinxlineitem{Parameters}\begin{itemize}
\item {} 
\sphinxAtStartPar
\sphinxstyleliteralstrong{\sphinxupquote{lat}} (\sphinxstyleliteralemphasis{\sphinxupquote{np.ndarray}}) \textendash{} The latitude coordinate for the position of interest.
The latitude position value given in decimal degrees.

\item {} 
\sphinxAtStartPar
\sphinxstyleliteralstrong{\sphinxupquote{lon}} (\sphinxstyleliteralemphasis{\sphinxupquote{np.ndarray}}) \textendash{} The longitude coordinate for the position of interest.
The longitude position value given in decimal degrees.

\item {} 
\sphinxAtStartPar
\sphinxstyleliteralstrong{\sphinxupquote{alt}} (\sphinxstyleliteralemphasis{\sphinxupquote{np.ndarray}}) \textendash{} The altitude offset for the position of interest.
The elevation from the WGS\sphinxhyphen{}84 reference ellipsoid in metres.

\item {} 
\sphinxAtStartPar
\sphinxstyleliteralstrong{\sphinxupquote{step}} (\sphinxstyleliteralemphasis{\sphinxupquote{float}}) \textendash{} The step size (in metres) for calculating the gradient.
To be used when the gradient cannot be computed directly.

\end{itemize}

\sphinxlineitem{Returns}
\sphinxAtStartPar
The vertical gravity gradient (in m/s\textasciicircum{}2).

\sphinxlineitem{Return type}
\sphinxAtStartPar
float

\end{description}\end{quote}

\end{fulllineitems}

\index{base\_model (qnav.gravity.map.GravityMap property)@\spxentry{base\_model}\spxextra{qnav.gravity.map.GravityMap property}}

\begin{fulllineitems}
\phantomsection\label{\detokenize{modules/gravity/map:qnav.gravity.map.GravityMap.base_model}}
\pysigstartsignatures
\pysigline{\sphinxbfcode{\sphinxupquote{property\DUrole{w}{ }}}\sphinxbfcode{\sphinxupquote{base\_model}}\sphinxbfcode{\sphinxupquote{\DUrole{p}{:}\DUrole{w}{ }Self\DUrole{w}{ }\DUrole{p}{|}\DUrole{w}{ }None}}}
\pysigstopsignatures
\sphinxAtStartPar
Returns the base gravity model the model extends from.
Allows access to the base function/model the instance is using.
\begin{quote}\begin{description}
\sphinxlineitem{Returns}
\sphinxAtStartPar
The base gravity model.

\sphinxlineitem{Return type}
\sphinxAtStartPar
{\hyperref[\detokenize{modules/gravity/base:qnav.gravity.base.GravityModel}]{\sphinxcrossref{GravityModel}}} | None

\end{description}\end{quote}

\end{fulllineitems}


\end{fulllineitems}


\sphinxstepscope


\subsubsection{qnav.gravity.egm}
\label{\detokenize{modules/gravity/egm:module-qnav.gravity.egm}}\label{\detokenize{modules/gravity/egm:qnav-gravity-egm}}\label{\detokenize{modules/gravity/egm::doc}}\index{module@\spxentry{module}!qnav.gravity.egm@\spxentry{qnav.gravity.egm}}\index{qnav.gravity.egm@\spxentry{qnav.gravity.egm}!module@\spxentry{module}}

\paragraph{egm.py}
\label{\detokenize{modules/gravity/egm:egm-py}}\begin{quote}\begin{description}
\sphinxlineitem{summary}
\sphinxAtStartPar
A module for supplying Earth Gravitational Models (EGMs).
This provides implementations for EGM gravity models, using a standard
‘base’ gravity function and a geoid model for height and surface
deflection correction.

\sphinxlineitem{authors}
\begin{DUlineblock}{0em}
\item[] Michael Wright \sphinxhyphen{} \sphinxhref{mailto:mjwright@liverpool.ac.uk}{mjwright@liverpool.ac.uk}
\end{DUlineblock}

\sphinxlineitem{version}
\sphinxAtStartPar
1.0.0 \sphinxhyphen{} First full implementation of this module.

\end{description}\end{quote}
\index{GeoidCorrection (class in qnav.gravity.egm)@\spxentry{GeoidCorrection}\spxextra{class in qnav.gravity.egm}}

\begin{fulllineitems}
\phantomsection\label{\detokenize{modules/gravity/egm:qnav.gravity.egm.GeoidCorrection}}
\pysigstartsignatures
\pysiglinewithargsret{\sphinxbfcode{\sphinxupquote{class\DUrole{w}{ }}}\sphinxcode{\sphinxupquote{qnav.gravity.egm.}}\sphinxbfcode{\sphinxupquote{GeoidCorrection}}}{\sphinxparam{\DUrole{n}{base\_model}\DUrole{p}{:}\DUrole{w}{ }\DUrole{n}{{\hyperref[\detokenize{modules/gravity/base:qnav.gravity.base.GravityModel}]{\sphinxcrossref{GravityModel}}}}}\sphinxparamcomma \sphinxparam{\DUrole{n}{geoid\_model}\DUrole{p}{:}\DUrole{w}{ }\DUrole{n}{{\hyperref[\detokenize{modules/earth/geoid:qnav.earth.geoid.GeoidModel}]{\sphinxcrossref{GeoidModel}}}}}}{}
\pysigstopsignatures
\sphinxAtStartPar
Gravity Map using Geoid Model for applying acceleration corrections.
Using a given geoid model to applying corrections to the gravity
accelerations, using the height offsets and deflections of vertical.
\index{calc\_gravity\_z() (qnav.gravity.egm.GeoidCorrection method)@\spxentry{calc\_gravity\_z()}\spxextra{qnav.gravity.egm.GeoidCorrection method}}

\begin{fulllineitems}
\phantomsection\label{\detokenize{modules/gravity/egm:qnav.gravity.egm.GeoidCorrection.calc_gravity_z}}
\pysigstartsignatures
\pysiglinewithargsret{\sphinxbfcode{\sphinxupquote{calc\_gravity\_z}}}{\sphinxparam{\DUrole{n}{lat}\DUrole{p}{:}\DUrole{w}{ }\DUrole{n}{float}}\sphinxparamcomma \sphinxparam{\DUrole{n}{lon}\DUrole{p}{:}\DUrole{w}{ }\DUrole{n}{float}}\sphinxparamcomma \sphinxparam{\DUrole{n}{alt}\DUrole{p}{:}\DUrole{w}{ }\DUrole{n}{float}}}{{ $\rightarrow$ float}}
\pysigstopsignatures
\sphinxAtStartPar
Calculates and returns the gravitational acceleration along the Z/Down
axis (in m/s\textasciicircum{}2) for the given location and altitude offset.
\begin{quote}\begin{description}
\sphinxlineitem{Parameters}\begin{itemize}
\item {} 
\sphinxAtStartPar
\sphinxstyleliteralstrong{\sphinxupquote{lat}} (\sphinxstyleliteralemphasis{\sphinxupquote{float}}) \textendash{} The latitude coordinate for the position of interest.
The latitude position value given in decimal degrees.

\item {} 
\sphinxAtStartPar
\sphinxstyleliteralstrong{\sphinxupquote{lon}} (\sphinxstyleliteralemphasis{\sphinxupquote{float}}) \textendash{} The longitude coordinate for the position of interest.
The longitude position value given in decimal degrees.

\item {} 
\sphinxAtStartPar
\sphinxstyleliteralstrong{\sphinxupquote{alt}} (\sphinxstyleliteralemphasis{\sphinxupquote{float}}) \textendash{} The altitude offset for the position of interest.
The elevation from the WGS\sphinxhyphen{}84 reference ellipsoid in metres.

\end{itemize}

\sphinxlineitem{Returns}
\sphinxAtStartPar
The downwards gravitational acceleration for each location.

\sphinxlineitem{Return type}
\sphinxAtStartPar
float

\end{description}\end{quote}

\end{fulllineitems}

\index{calc\_gravity\_z\_vec() (qnav.gravity.egm.GeoidCorrection method)@\spxentry{calc\_gravity\_z\_vec()}\spxextra{qnav.gravity.egm.GeoidCorrection method}}

\begin{fulllineitems}
\phantomsection\label{\detokenize{modules/gravity/egm:qnav.gravity.egm.GeoidCorrection.calc_gravity_z_vec}}
\pysigstartsignatures
\pysiglinewithargsret{\sphinxbfcode{\sphinxupquote{calc\_gravity\_z\_vec}}}{\sphinxparam{\DUrole{n}{lat}\DUrole{p}{:}\DUrole{w}{ }\DUrole{n}{ndarray}}\sphinxparamcomma \sphinxparam{\DUrole{n}{lon}\DUrole{p}{:}\DUrole{w}{ }\DUrole{n}{ndarray}}\sphinxparamcomma \sphinxparam{\DUrole{n}{alt}\DUrole{p}{:}\DUrole{w}{ }\DUrole{n}{ndarray}}}{{ $\rightarrow$ ndarray}}
\pysigstopsignatures
\sphinxAtStartPar
In a vectorised batch, calculates and returns the gravitational
acceleration along the Z/Down axis (in m/s\textasciicircum{}2) for the given location
and altitude offset. Ideal when calculating gravity for a large
number of positions.
\begin{quote}\begin{description}
\sphinxlineitem{Parameters}\begin{itemize}
\item {} 
\sphinxAtStartPar
\sphinxstyleliteralstrong{\sphinxupquote{lat}} (\sphinxstyleliteralemphasis{\sphinxupquote{np.array}}) \textendash{} The latitude coordinate for the position of interest.
The latitude position value given in decimal degrees.

\item {} 
\sphinxAtStartPar
\sphinxstyleliteralstrong{\sphinxupquote{lon}} (\sphinxstyleliteralemphasis{\sphinxupquote{np.array}}) \textendash{} The longitude coordinate for the position of interest.
The longitude position value given in decimal degrees.

\item {} 
\sphinxAtStartPar
\sphinxstyleliteralstrong{\sphinxupquote{alt}} (\sphinxstyleliteralemphasis{\sphinxupquote{np.array}}) \textendash{} The altitude offset for the position of interest.
The elevation from the WGS\sphinxhyphen{}84 reference ellipsoid in metres.

\end{itemize}

\sphinxlineitem{Returns}
\sphinxAtStartPar
The downwards gravitational acceleration for each location
(in m/s\textasciicircum{}2).

\sphinxlineitem{Return type}
\sphinxAtStartPar
np.array (N\sphinxhyphen{}elements)

\end{description}\end{quote}

\end{fulllineitems}

\index{get\_anomaly() (qnav.gravity.egm.GeoidCorrection method)@\spxentry{get\_anomaly()}\spxextra{qnav.gravity.egm.GeoidCorrection method}}

\begin{fulllineitems}
\phantomsection\label{\detokenize{modules/gravity/egm:qnav.gravity.egm.GeoidCorrection.get_anomaly}}
\pysigstartsignatures
\pysiglinewithargsret{\sphinxbfcode{\sphinxupquote{get\_anomaly}}}{\sphinxparam{\DUrole{n}{lat}\DUrole{p}{:}\DUrole{w}{ }\DUrole{n}{float}}\sphinxparamcomma \sphinxparam{\DUrole{n}{lon}\DUrole{p}{:}\DUrole{w}{ }\DUrole{n}{float}}}{{ $\rightarrow$ float}}
\pysigstopsignatures
\sphinxAtStartPar
Calculates the gravity anomaly for given position.
Retrieves the gravity anomaly value for the requested position
(in m/s\textasciicircum{}2) at the base surface. For gravity maps this value is
interpolated from its dataset, following possible conversion.
\begin{quote}\begin{description}
\sphinxlineitem{Parameters}\begin{itemize}
\item {} 
\sphinxAtStartPar
\sphinxstyleliteralstrong{\sphinxupquote{lat}} (\sphinxstyleliteralemphasis{\sphinxupquote{float}}) \textendash{} The latitude coordinate for the position of interest.
The latitude position value given in decimal degrees.

\item {} 
\sphinxAtStartPar
\sphinxstyleliteralstrong{\sphinxupquote{lon}} (\sphinxstyleliteralemphasis{\sphinxupquote{float}}) \textendash{} The longitude coordinate for the position of interest.
The longitude position value given in decimal degrees.

\end{itemize}

\sphinxlineitem{Returns}
\sphinxAtStartPar
The gravity anomaly value (in m/s\textasciicircum{}2).

\sphinxlineitem{Return type}
\sphinxAtStartPar
float

\end{description}\end{quote}

\end{fulllineitems}

\index{get\_anomaly\_vec() (qnav.gravity.egm.GeoidCorrection method)@\spxentry{get\_anomaly\_vec()}\spxextra{qnav.gravity.egm.GeoidCorrection method}}

\begin{fulllineitems}
\phantomsection\label{\detokenize{modules/gravity/egm:qnav.gravity.egm.GeoidCorrection.get_anomaly_vec}}
\pysigstartsignatures
\pysiglinewithargsret{\sphinxbfcode{\sphinxupquote{get\_anomaly\_vec}}}{\sphinxparam{\DUrole{n}{lat}\DUrole{p}{:}\DUrole{w}{ }\DUrole{n}{ndarray}}\sphinxparamcomma \sphinxparam{\DUrole{n}{lon}\DUrole{p}{:}\DUrole{w}{ }\DUrole{n}{ndarray}}}{{ $\rightarrow$ ndarray}}
\pysigstopsignatures
\sphinxAtStartPar
Calculates the gravity anomaly for multiple given positions.
Retrieves the gravity anomaly value for the requested position
(in m/s\textasciicircum{}2) at the base surface. For gravity maps this value is
interpolated from its dataset, following possible conversion.
\begin{quote}\begin{description}
\sphinxlineitem{Parameters}\begin{itemize}
\item {} 
\sphinxAtStartPar
\sphinxstyleliteralstrong{\sphinxupquote{lat}} (\sphinxstyleliteralemphasis{\sphinxupquote{float}}) \textendash{} The latitude coordinate for the position of interest.
The latitude position value given in decimal degrees.

\item {} 
\sphinxAtStartPar
\sphinxstyleliteralstrong{\sphinxupquote{lon}} (\sphinxstyleliteralemphasis{\sphinxupquote{float}}) \textendash{} The longitude coordinate for the position of interest.
The longitude position value given in decimal degrees.

\end{itemize}

\sphinxlineitem{Returns}
\sphinxAtStartPar
The gravity anomaly value (in m/s\textasciicircum{}2).

\sphinxlineitem{Return type}
\sphinxAtStartPar
float

\end{description}\end{quote}

\end{fulllineitems}

\index{get\_disturbance() (qnav.gravity.egm.GeoidCorrection method)@\spxentry{get\_disturbance()}\spxextra{qnav.gravity.egm.GeoidCorrection method}}

\begin{fulllineitems}
\phantomsection\label{\detokenize{modules/gravity/egm:qnav.gravity.egm.GeoidCorrection.get_disturbance}}
\pysigstartsignatures
\pysiglinewithargsret{\sphinxbfcode{\sphinxupquote{get\_disturbance}}}{\sphinxparam{\DUrole{n}{lat}\DUrole{p}{:}\DUrole{w}{ }\DUrole{n}{float}}\sphinxparamcomma \sphinxparam{\DUrole{n}{lon}\DUrole{p}{:}\DUrole{w}{ }\DUrole{n}{float}}}{{ $\rightarrow$ float}}
\pysigstopsignatures
\sphinxAtStartPar
Calculates the gravity disturbance for given position.
Retrieves the gravity disturbance value for the requested position
(in m/s\textasciicircum{}2) at the base surface. For gravity maps this value is
interpolated from its dataset, following possible conversion.
\begin{quote}\begin{description}
\sphinxlineitem{Parameters}\begin{itemize}
\item {} 
\sphinxAtStartPar
\sphinxstyleliteralstrong{\sphinxupquote{lat}} (\sphinxstyleliteralemphasis{\sphinxupquote{float}}) \textendash{} The latitude coordinate for the position of interest.
The latitude position value given in decimal degrees.

\item {} 
\sphinxAtStartPar
\sphinxstyleliteralstrong{\sphinxupquote{lon}} (\sphinxstyleliteralemphasis{\sphinxupquote{float}}) \textendash{} The longitude coordinate for the position of interest.
The longitude position value given in decimal degrees.

\end{itemize}

\sphinxlineitem{Returns}
\sphinxAtStartPar
The gravity disturbance value (in m/s\textasciicircum{}2).

\sphinxlineitem{Return type}
\sphinxAtStartPar
float

\end{description}\end{quote}

\end{fulllineitems}

\index{get\_disturbance\_vec() (qnav.gravity.egm.GeoidCorrection method)@\spxentry{get\_disturbance\_vec()}\spxextra{qnav.gravity.egm.GeoidCorrection method}}

\begin{fulllineitems}
\phantomsection\label{\detokenize{modules/gravity/egm:qnav.gravity.egm.GeoidCorrection.get_disturbance_vec}}
\pysigstartsignatures
\pysiglinewithargsret{\sphinxbfcode{\sphinxupquote{get\_disturbance\_vec}}}{\sphinxparam{\DUrole{n}{lat}\DUrole{p}{:}\DUrole{w}{ }\DUrole{n}{ndarray}}\sphinxparamcomma \sphinxparam{\DUrole{n}{lon}\DUrole{p}{:}\DUrole{w}{ }\DUrole{n}{ndarray}}}{{ $\rightarrow$ ndarray}}
\pysigstopsignatures
\sphinxAtStartPar
Calculates the gravity disturbance for multiple given positions.
Retrieves the gravity disturbance value for the requested position
(in m/s\textasciicircum{}2) at the base surface. For gravity maps this value is
interpolated from its dataset, following possible conversion.
\begin{quote}\begin{description}
\sphinxlineitem{Parameters}\begin{itemize}
\item {} 
\sphinxAtStartPar
\sphinxstyleliteralstrong{\sphinxupquote{lat}} (\sphinxstyleliteralemphasis{\sphinxupquote{float}}) \textendash{} The latitude coordinate for the position of interest.
The latitude position value given in decimal degrees.

\item {} 
\sphinxAtStartPar
\sphinxstyleliteralstrong{\sphinxupquote{lon}} (\sphinxstyleliteralemphasis{\sphinxupquote{float}}) \textendash{} The longitude coordinate for the position of interest.
The longitude position value given in decimal degrees.

\end{itemize}

\sphinxlineitem{Returns}
\sphinxAtStartPar
The gravity disturbance value (in m/s\textasciicircum{}2).

\sphinxlineitem{Return type}
\sphinxAtStartPar
float

\end{description}\end{quote}

\end{fulllineitems}


\end{fulllineitems}


\sphinxstepscope


\subsubsection{qnav.gravity.geosat}
\label{\detokenize{modules/gravity/geosat:module-qnav.gravity.geosat}}\label{\detokenize{modules/gravity/geosat:qnav-gravity-geosat}}\label{\detokenize{modules/gravity/geosat::doc}}\index{module@\spxentry{module}!qnav.gravity.geosat@\spxentry{qnav.gravity.geosat}}\index{qnav.gravity.geosat@\spxentry{qnav.gravity.geosat}!module@\spxentry{module}}

\paragraph{geosat.py}
\label{\detokenize{modules/gravity/geosat:geosat-py}}\begin{quote}\begin{description}
\sphinxlineitem{summary}
\sphinxAtStartPar
Geosat\sphinxhyphen{}44 reading, interpolation and gravity functions.
This module contains all the gravity calculation functionalities related
to the geosat\sphinxhyphen{}44 dataset. The Geosat44 database contains records of
time\sphinxhyphen{}offsets, positions, geoid heights, gravity anomalies and
calculation uncertainties.

\sphinxlineitem{author}
\sphinxAtStartPar
Michael Wright \sphinxhyphen{} \sphinxhref{mailto:mjwright@liverpool.ac.uk}{mjwright@liverpool.ac.uk}

\sphinxlineitem{version}
\sphinxAtStartPar
1.0.0 \sphinxhyphen{} First implementation of this module.

\end{description}\end{quote}
\index{Geosat44 (class in qnav.gravity.geosat)@\spxentry{Geosat44}\spxextra{class in qnav.gravity.geosat}}

\begin{fulllineitems}
\phantomsection\label{\detokenize{modules/gravity/geosat:qnav.gravity.geosat.Geosat44}}
\pysigstartsignatures
\pysiglinewithargsret{\sphinxbfcode{\sphinxupquote{class\DUrole{w}{ }}}\sphinxcode{\sphinxupquote{qnav.gravity.geosat.}}\sphinxbfcode{\sphinxupquote{Geosat44}}}{\sphinxparam{\DUrole{n}{base\_model}\DUrole{p}{:}\DUrole{w}{ }\DUrole{n}{{\hyperref[\detokenize{modules/gravity/base:qnav.gravity.base.GravityModel}]{\sphinxcrossref{GravityModel}}}}}\sphinxparamcomma \sphinxparam{\DUrole{n}{geoid\_model}\DUrole{p}{:}\DUrole{w}{ }\DUrole{n}{{\hyperref[\detokenize{modules/earth/geoid:qnav.earth.geoid.GeoidModel}]{\sphinxcrossref{GeoidModel}}}\DUrole{w}{ }\DUrole{p}{|}\DUrole{w}{ }None}}\sphinxparamcomma \sphinxparam{\DUrole{n}{map\_dir}\DUrole{p}{:}\DUrole{w}{ }\DUrole{n}{Path}}\sphinxparamcomma \sphinxparam{\DUrole{n}{auto\_download}\DUrole{p}{:}\DUrole{w}{ }\DUrole{n}{bool}\DUrole{w}{ }\DUrole{o}{=}\DUrole{w}{ }\DUrole{default_value}{True}}\sphinxparamcomma \sphinxparam{\DUrole{n}{auto\_load}\DUrole{p}{:}\DUrole{w}{ }\DUrole{n}{bool}\DUrole{w}{ }\DUrole{o}{=}\DUrole{w}{ }\DUrole{default_value}{True}}}{}
\pysigstopsignatures
\sphinxAtStartPar
Gravity Map implementation for the Geosat\sphinxhyphen{}44 dataset.
Provides corrections for a given base gravity function, using the
non\sphinxhyphen{}gridded geosat data values. Performs linear interpolation by
producing a Delaunay triangulation mesh with fixed zero value for
extrapolation.
\index{calc\_gravity\_z() (qnav.gravity.geosat.Geosat44 method)@\spxentry{calc\_gravity\_z()}\spxextra{qnav.gravity.geosat.Geosat44 method}}

\begin{fulllineitems}
\phantomsection\label{\detokenize{modules/gravity/geosat:qnav.gravity.geosat.Geosat44.calc_gravity_z}}
\pysigstartsignatures
\pysiglinewithargsret{\sphinxbfcode{\sphinxupquote{calc\_gravity\_z}}}{\sphinxparam{\DUrole{n}{lat}\DUrole{p}{:}\DUrole{w}{ }\DUrole{n}{float}}\sphinxparamcomma \sphinxparam{\DUrole{n}{lon}\DUrole{p}{:}\DUrole{w}{ }\DUrole{n}{float}}\sphinxparamcomma \sphinxparam{\DUrole{n}{alt}\DUrole{p}{:}\DUrole{w}{ }\DUrole{n}{float}}}{{ $\rightarrow$ float}}
\pysigstopsignatures
\sphinxAtStartPar
Calculates and returns the total gravitational acceleration.
Computes the gravity along the Z/Down axis (in m/s\textasciicircum{}2) for
the given location and altitude offset.
\begin{quote}\begin{description}
\sphinxlineitem{Parameters}\begin{itemize}
\item {} 
\sphinxAtStartPar
\sphinxstyleliteralstrong{\sphinxupquote{lat}} (\sphinxstyleliteralemphasis{\sphinxupquote{float}}) \textendash{} The latitude coordinate for the position of interest.
The latitude position value given in decimal degrees.

\item {} 
\sphinxAtStartPar
\sphinxstyleliteralstrong{\sphinxupquote{lon}} (\sphinxstyleliteralemphasis{\sphinxupquote{float}}) \textendash{} The longitude coordinate for the position of interest.
The longitude position value given in decimal degrees.

\item {} 
\sphinxAtStartPar
\sphinxstyleliteralstrong{\sphinxupquote{alt}} (\sphinxstyleliteralemphasis{\sphinxupquote{float}}) \textendash{} The altitude offset for the position of interest.
The elevation from the WGS\sphinxhyphen{}84 reference ellipsoid in metres.

\end{itemize}

\sphinxlineitem{Returns}
\sphinxAtStartPar
The downwards gravitational acceleration for each location.

\sphinxlineitem{Return type}
\sphinxAtStartPar
float

\end{description}\end{quote}

\end{fulllineitems}

\index{calc\_gravity\_z\_vec() (qnav.gravity.geosat.Geosat44 method)@\spxentry{calc\_gravity\_z\_vec()}\spxextra{qnav.gravity.geosat.Geosat44 method}}

\begin{fulllineitems}
\phantomsection\label{\detokenize{modules/gravity/geosat:qnav.gravity.geosat.Geosat44.calc_gravity_z_vec}}
\pysigstartsignatures
\pysiglinewithargsret{\sphinxbfcode{\sphinxupquote{calc\_gravity\_z\_vec}}}{\sphinxparam{\DUrole{n}{lat}\DUrole{p}{:}\DUrole{w}{ }\DUrole{n}{ndarray}}\sphinxparamcomma \sphinxparam{\DUrole{n}{lon}\DUrole{p}{:}\DUrole{w}{ }\DUrole{n}{ndarray}}\sphinxparamcomma \sphinxparam{\DUrole{n}{alt}\DUrole{p}{:}\DUrole{w}{ }\DUrole{n}{ndarray}}}{{ $\rightarrow$ ndarray}}
\pysigstopsignatures
\sphinxAtStartPar
In a vectorised batch, calculates and returns the gravitational
acceleration along the Z/Down axis (in m/s\textasciicircum{}2) for the given location
and altitude offset. Ideal when calculating gravity for a large
number of positions.
\begin{quote}\begin{description}
\sphinxlineitem{Parameters}\begin{itemize}
\item {} 
\sphinxAtStartPar
\sphinxstyleliteralstrong{\sphinxupquote{lat}} (\sphinxstyleliteralemphasis{\sphinxupquote{np.array}}) \textendash{} The latitude coordinate for the position of interest.
The latitude position value given in decimal degrees.

\item {} 
\sphinxAtStartPar
\sphinxstyleliteralstrong{\sphinxupquote{lon}} (\sphinxstyleliteralemphasis{\sphinxupquote{np.array}}) \textendash{} The longitude coordinate for the position of interest.
The longitude position value given in decimal degrees.

\item {} 
\sphinxAtStartPar
\sphinxstyleliteralstrong{\sphinxupquote{alt}} (\sphinxstyleliteralemphasis{\sphinxupquote{np.array}}) \textendash{} The altitude offset for the position of interest.
The elevation from the WGS\sphinxhyphen{}84 reference ellipsoid in metres.

\end{itemize}

\sphinxlineitem{Returns}
\sphinxAtStartPar
The downwards gravitational acceleration for each location
(in m/s\textasciicircum{}2).

\sphinxlineitem{Return type}
\sphinxAtStartPar
np.array (N\sphinxhyphen{}elements)

\end{description}\end{quote}

\end{fulllineitems}

\index{get\_map\_value() (qnav.gravity.geosat.Geosat44 method)@\spxentry{get\_map\_value()}\spxextra{qnav.gravity.geosat.Geosat44 method}}

\begin{fulllineitems}
\phantomsection\label{\detokenize{modules/gravity/geosat:qnav.gravity.geosat.Geosat44.get_map_value}}
\pysigstartsignatures
\pysiglinewithargsret{\sphinxbfcode{\sphinxupquote{get\_map\_value}}}{\sphinxparam{\DUrole{n}{lat}\DUrole{p}{:}\DUrole{w}{ }\DUrole{n}{float\DUrole{w}{ }\DUrole{p}{|}\DUrole{w}{ }ndarray}}\sphinxparamcomma \sphinxparam{\DUrole{n}{lon}\DUrole{p}{:}\DUrole{w}{ }\DUrole{n}{float\DUrole{w}{ }\DUrole{p}{|}\DUrole{w}{ }ndarray}}}{{ $\rightarrow$ float\DUrole{w}{ }\DUrole{p}{|}\DUrole{w}{ }ndarray}}
\pysigstopsignatures
\sphinxAtStartPar
Interpolates the gravity anomaly for requested position(s).
Looks\sphinxhyphen{}up and returns the gravity anomaly value for requested latitude
and longitude positions. Supports both scalar and array input.
\begin{quote}\begin{description}
\sphinxlineitem{Parameters}\begin{itemize}
\item {} 
\sphinxAtStartPar
\sphinxstyleliteralstrong{\sphinxupquote{lat}} (\sphinxstyleliteralemphasis{\sphinxupquote{float}}\sphinxstyleliteralemphasis{\sphinxupquote{ | }}\sphinxstyleliteralemphasis{\sphinxupquote{np.ndarray}}) \textendash{} The latitude coordinate(s) in decimal degrees.

\item {} 
\sphinxAtStartPar
\sphinxstyleliteralstrong{\sphinxupquote{lon}} (\sphinxstyleliteralemphasis{\sphinxupquote{float}}\sphinxstyleliteralemphasis{\sphinxupquote{ | }}\sphinxstyleliteralemphasis{\sphinxupquote{np.ndarray}}) \textendash{} The longitude coordinate(s) in decimal degrees.

\end{itemize}

\sphinxlineitem{Returns}
\sphinxAtStartPar
The interpolated gravity anomaly value(s) in m/s\textasciicircum{}2.

\sphinxlineitem{Return type}
\sphinxAtStartPar
float | np.ndarray

\end{description}\end{quote}

\end{fulllineitems}

\index{get\_anomaly() (qnav.gravity.geosat.Geosat44 method)@\spxentry{get\_anomaly()}\spxextra{qnav.gravity.geosat.Geosat44 method}}

\begin{fulllineitems}
\phantomsection\label{\detokenize{modules/gravity/geosat:qnav.gravity.geosat.Geosat44.get_anomaly}}
\pysigstartsignatures
\pysiglinewithargsret{\sphinxbfcode{\sphinxupquote{get\_anomaly}}}{\sphinxparam{\DUrole{n}{lat}\DUrole{p}{:}\DUrole{w}{ }\DUrole{n}{float}}\sphinxparamcomma \sphinxparam{\DUrole{n}{lon}\DUrole{p}{:}\DUrole{w}{ }\DUrole{n}{float}}}{{ $\rightarrow$ float}}
\pysigstopsignatures
\sphinxAtStartPar
Calculates the gravity anomaly for given position.
Retrieves the gravity anomaly value for the requested position
(in m/s\textasciicircum{}2) at the base surface. For gravity maps this value is
interpolated from its dataset, following possible conversion.
\begin{quote}\begin{description}
\sphinxlineitem{Parameters}\begin{itemize}
\item {} 
\sphinxAtStartPar
\sphinxstyleliteralstrong{\sphinxupquote{lat}} (\sphinxstyleliteralemphasis{\sphinxupquote{float}}) \textendash{} The latitude coordinate for the position of interest.
The latitude position value given in decimal degrees.

\item {} 
\sphinxAtStartPar
\sphinxstyleliteralstrong{\sphinxupquote{lon}} (\sphinxstyleliteralemphasis{\sphinxupquote{float}}) \textendash{} The longitude coordinate for the position of interest.
The longitude position value given in decimal degrees.

\end{itemize}

\sphinxlineitem{Returns}
\sphinxAtStartPar
The gravity anomaly value (in m/s\textasciicircum{}2).

\sphinxlineitem{Return type}
\sphinxAtStartPar
float

\end{description}\end{quote}

\end{fulllineitems}

\index{get\_anomaly\_vec() (qnav.gravity.geosat.Geosat44 method)@\spxentry{get\_anomaly\_vec()}\spxextra{qnav.gravity.geosat.Geosat44 method}}

\begin{fulllineitems}
\phantomsection\label{\detokenize{modules/gravity/geosat:qnav.gravity.geosat.Geosat44.get_anomaly_vec}}
\pysigstartsignatures
\pysiglinewithargsret{\sphinxbfcode{\sphinxupquote{get\_anomaly\_vec}}}{\sphinxparam{\DUrole{n}{lat}\DUrole{p}{:}\DUrole{w}{ }\DUrole{n}{ndarray}}\sphinxparamcomma \sphinxparam{\DUrole{n}{lon}\DUrole{p}{:}\DUrole{w}{ }\DUrole{n}{ndarray}}}{{ $\rightarrow$ ndarray}}
\pysigstopsignatures
\sphinxAtStartPar
Calculates the gravity anomaly for multiple given positions.
Retrieves the gravity anomaly value for the requested position
(in m/s\textasciicircum{}2) at the base surface. For gravity maps this value is
interpolated from its dataset, following possible conversion.
\begin{quote}\begin{description}
\sphinxlineitem{Parameters}\begin{itemize}
\item {} 
\sphinxAtStartPar
\sphinxstyleliteralstrong{\sphinxupquote{lat}} (\sphinxstyleliteralemphasis{\sphinxupquote{float}}) \textendash{} The latitude coordinate for the position of interest.
The latitude position value given in decimal degrees.

\item {} 
\sphinxAtStartPar
\sphinxstyleliteralstrong{\sphinxupquote{lon}} (\sphinxstyleliteralemphasis{\sphinxupquote{float}}) \textendash{} The longitude coordinate for the position of interest.
The longitude position value given in decimal degrees.

\end{itemize}

\sphinxlineitem{Returns}
\sphinxAtStartPar
The gravity anomaly value (in m/s\textasciicircum{}2).

\sphinxlineitem{Return type}
\sphinxAtStartPar
float

\end{description}\end{quote}

\end{fulllineitems}

\index{get\_disturbance() (qnav.gravity.geosat.Geosat44 method)@\spxentry{get\_disturbance()}\spxextra{qnav.gravity.geosat.Geosat44 method}}

\begin{fulllineitems}
\phantomsection\label{\detokenize{modules/gravity/geosat:qnav.gravity.geosat.Geosat44.get_disturbance}}
\pysigstartsignatures
\pysiglinewithargsret{\sphinxbfcode{\sphinxupquote{get\_disturbance}}}{\sphinxparam{\DUrole{n}{lat}\DUrole{p}{:}\DUrole{w}{ }\DUrole{n}{float}}\sphinxparamcomma \sphinxparam{\DUrole{n}{lon}\DUrole{p}{:}\DUrole{w}{ }\DUrole{n}{float}}}{{ $\rightarrow$ float}}
\pysigstopsignatures
\sphinxAtStartPar
Calculates the gravity disturbance for given position.
Retrieves the gravity disturbance value for the requested position
(in m/s\textasciicircum{}2) at the base surface. For gravity maps this value is
interpolated from its dataset, following possible conversion.
\begin{quote}\begin{description}
\sphinxlineitem{Parameters}\begin{itemize}
\item {} 
\sphinxAtStartPar
\sphinxstyleliteralstrong{\sphinxupquote{lat}} (\sphinxstyleliteralemphasis{\sphinxupquote{float}}) \textendash{} The latitude coordinate for the position of interest.
The latitude position value given in decimal degrees.

\item {} 
\sphinxAtStartPar
\sphinxstyleliteralstrong{\sphinxupquote{lon}} (\sphinxstyleliteralemphasis{\sphinxupquote{float}}) \textendash{} The longitude coordinate for the position of interest.
The longitude position value given in decimal degrees.

\end{itemize}

\sphinxlineitem{Returns}
\sphinxAtStartPar
The gravity disturbance value (in m/s\textasciicircum{}2).

\sphinxlineitem{Return type}
\sphinxAtStartPar
float

\end{description}\end{quote}

\end{fulllineitems}

\index{get\_disturbance\_vec() (qnav.gravity.geosat.Geosat44 method)@\spxentry{get\_disturbance\_vec()}\spxextra{qnav.gravity.geosat.Geosat44 method}}

\begin{fulllineitems}
\phantomsection\label{\detokenize{modules/gravity/geosat:qnav.gravity.geosat.Geosat44.get_disturbance_vec}}
\pysigstartsignatures
\pysiglinewithargsret{\sphinxbfcode{\sphinxupquote{get\_disturbance\_vec}}}{\sphinxparam{\DUrole{n}{lat}\DUrole{p}{:}\DUrole{w}{ }\DUrole{n}{ndarray}}\sphinxparamcomma \sphinxparam{\DUrole{n}{lon}\DUrole{p}{:}\DUrole{w}{ }\DUrole{n}{ndarray}}}{{ $\rightarrow$ ndarray}}
\pysigstopsignatures
\sphinxAtStartPar
Calculates the gravity disturbance for multiple given positions.
Retrieves the gravity disturbance value for the requested position
(in m/s\textasciicircum{}2) at the base surface. For gravity maps this value is
interpolated from its dataset, following possible conversion.
\begin{quote}\begin{description}
\sphinxlineitem{Parameters}\begin{itemize}
\item {} 
\sphinxAtStartPar
\sphinxstyleliteralstrong{\sphinxupquote{lat}} (\sphinxstyleliteralemphasis{\sphinxupquote{float}}) \textendash{} The latitude coordinate for the position of interest.
The latitude position value given in decimal degrees.

\item {} 
\sphinxAtStartPar
\sphinxstyleliteralstrong{\sphinxupquote{lon}} (\sphinxstyleliteralemphasis{\sphinxupquote{float}}) \textendash{} The longitude coordinate for the position of interest.
The longitude position value given in decimal degrees.

\end{itemize}

\sphinxlineitem{Returns}
\sphinxAtStartPar
The gravity disturbance value (in m/s\textasciicircum{}2).

\sphinxlineitem{Return type}
\sphinxAtStartPar
float

\end{description}\end{quote}

\end{fulllineitems}

\index{load\_data() (qnav.gravity.geosat.Geosat44 method)@\spxentry{load\_data()}\spxextra{qnav.gravity.geosat.Geosat44 method}}

\begin{fulllineitems}
\phantomsection\label{\detokenize{modules/gravity/geosat:qnav.gravity.geosat.Geosat44.load_data}}
\pysigstartsignatures
\pysiglinewithargsret{\sphinxbfcode{\sphinxupquote{load\_data}}}{}{}
\pysigstopsignatures
\sphinxAtStartPar
Loads the Geosat data and generates internal interpolator for look\sphinxhyphen{}up.
Called on initialisation to load the required data to allow for
correction interpolation. Due to spacing between geosat records,
linear interpolation is used for stability.

\end{fulllineitems}

\index{download\_data() (qnav.gravity.geosat.Geosat44 method)@\spxentry{download\_data()}\spxextra{qnav.gravity.geosat.Geosat44 method}}

\begin{fulllineitems}
\phantomsection\label{\detokenize{modules/gravity/geosat:qnav.gravity.geosat.Geosat44.download_data}}
\pysigstartsignatures
\pysiglinewithargsret{\sphinxbfcode{\sphinxupquote{download\_data}}}{\sphinxparam{\DUrole{n}{overwrite}\DUrole{p}{:}\DUrole{w}{ }\DUrole{n}{bool}\DUrole{w}{ }\DUrole{o}{=}\DUrole{w}{ }\DUrole{default_value}{False}}}{}
\pysigstopsignatures
\sphinxAtStartPar
Downloads the required Geosat\sphinxhyphen{}44 files from the source.
If allowed and files are not present, will download the required
geosat files to the set directory. Optionally, files can be
overwritten.
\begin{quote}\begin{description}
\sphinxlineitem{Parameters}
\sphinxAtStartPar
\sphinxstyleliteralstrong{\sphinxupquote{overwrite}} (\sphinxstyleliteralemphasis{\sphinxupquote{bool}}) \textendash{} Should downloading replace existing files.
By default this is set to false and only missing files
are downloaded.

\end{description}\end{quote}

\end{fulllineitems}


\end{fulllineitems}

\index{read\_geosat\_file() (in module qnav.gravity.geosat)@\spxentry{read\_geosat\_file()}\spxextra{in module qnav.gravity.geosat}}

\begin{fulllineitems}
\phantomsection\label{\detokenize{modules/gravity/geosat:qnav.gravity.geosat.read_geosat_file}}
\pysigstartsignatures
\pysiglinewithargsret{\sphinxcode{\sphinxupquote{qnav.gravity.geosat.}}\sphinxbfcode{\sphinxupquote{read\_geosat\_file}}}{\sphinxparam{\DUrole{n}{file\_path}\DUrole{p}{:}\DUrole{w}{ }\DUrole{n}{Path}}}{{ $\rightarrow$ dict\DUrole{p}{{[}}str\DUrole{p}{,}\DUrole{w}{ }ndarray\DUrole{p}{{]}}}}
\pysigstopsignatures
\sphinxAtStartPar
Reads and extracts the content of a given geosat file.
Parses a given binary (.bin) geosat file and returns its records in a
Python dictionary. The records include:
• time\_offset:  The time since the beginning of each pass (seconds)
• latitude:     The latitude position of measurement (degrees)
• longitude:    The longitude position of measurement (degrees)
• geoid:        The geoid height from the ellipse (metres)
• anomaly:      The calculated gravity anomaly (m/s\textasciicircum{}2)
• uncertainty:  The measurement uncertainty (m/s\textasciicircum{}2)
\begin{quote}\begin{description}
\sphinxlineitem{Parameters}
\sphinxAtStartPar
\sphinxstyleliteralstrong{\sphinxupquote{file\_path}} (\sphinxstyleliteralemphasis{\sphinxupquote{Path}}) \textendash{} The path for the binary geosat file to read.

\sphinxlineitem{Returns}
\sphinxAtStartPar
A dictionary containing the read data records.

\sphinxlineitem{Return type}
\sphinxAtStartPar
dict{[}str, np.ndarray{]}

\end{description}\end{quote}

\end{fulllineitems}

\index{create\_interp\_grid() (in module qnav.gravity.geosat)@\spxentry{create\_interp\_grid()}\spxextra{in module qnav.gravity.geosat}}

\begin{fulllineitems}
\phantomsection\label{\detokenize{modules/gravity/geosat:qnav.gravity.geosat.create_interp_grid}}
\pysigstartsignatures
\pysiglinewithargsret{\sphinxcode{\sphinxupquote{qnav.gravity.geosat.}}\sphinxbfcode{\sphinxupquote{create\_interp\_grid}}}{\sphinxparam{\DUrole{n}{geosat\_data}\DUrole{p}{:}\DUrole{w}{ }\DUrole{n}{dict\DUrole{p}{{[}}str\DUrole{p}{,}\DUrole{w}{ }ndarray\DUrole{p}{{]}}}}\sphinxparamcomma \sphinxparam{\DUrole{n}{interp\_method}\DUrole{p}{:}\DUrole{w}{ }\DUrole{n}{str}\DUrole{w}{ }\DUrole{o}{=}\DUrole{w}{ }\DUrole{default_value}{\textquotesingle{}linear\textquotesingle{}}}\sphinxparamcomma \sphinxparam{\DUrole{n}{resolution}\DUrole{p}{:}\DUrole{w}{ }\DUrole{n}{float}\DUrole{w}{ }\DUrole{o}{=}\DUrole{w}{ }\DUrole{default_value}{0.1}}}{{ $\rightarrow$ tuple\DUrole{p}{{[}}ndarray\DUrole{p}{,}\DUrole{w}{ }ndarray\DUrole{p}{,}\DUrole{w}{ }ndarray\DUrole{p}{{]}}}}
\pysigstopsignatures
\sphinxAtStartPar
Generates interpolated grid points for geosat data.
For given geosat data, a set of grid points are generated using the
requested interpolation method and resolution. This is provided to
produced uniformly gridded maps for faster look\sphinxhyphen{}up and interpolation.
\begin{quote}\begin{description}
\sphinxlineitem{Parameters}\begin{itemize}
\item {} 
\sphinxAtStartPar
\sphinxstyleliteralstrong{\sphinxupquote{geosat\_data}} (\sphinxstyleliteralemphasis{\sphinxupquote{dict}}\sphinxstyleliteralemphasis{\sphinxupquote{{[}}}\sphinxstyleliteralemphasis{\sphinxupquote{str}}\sphinxstyleliteralemphasis{\sphinxupquote{, }}\sphinxstyleliteralemphasis{\sphinxupquote{np.ndarray}}\sphinxstyleliteralemphasis{\sphinxupquote{{]}}}) \textendash{} The contents read from the geosat file(s).

\item {} 
\sphinxAtStartPar
\sphinxstyleliteralstrong{\sphinxupquote{interp\_method}} (\sphinxstyleliteralemphasis{\sphinxupquote{str}}) \textendash{} (Optional) The interpolation method to use.
Supported methods include ‘nearest’, ‘linear’ and ‘cubic’.
By default, the ‘linear’ method is used.

\item {} 
\sphinxAtStartPar
\sphinxstyleliteralstrong{\sphinxupquote{resolution}} \textendash{} (Optional) The requested step\sphinxhyphen{}size for the
interpolation in decimal degrees. By default 0.1 is used.

\end{itemize}

\sphinxlineitem{Returns}
\sphinxAtStartPar
A tuple containing the produced latitude, longitude and gravity
anomaly values, respectively.

\sphinxlineitem{Return type}
\sphinxAtStartPar
tuple{[}np.ndarray, np.ndarray, np.ndarray{]}

\end{description}\end{quote}

\end{fulllineitems}


\sphinxstepscope


\subsubsection{qnav.gravity.marine}
\label{\detokenize{modules/gravity/marine:module-qnav.gravity.marine}}\label{\detokenize{modules/gravity/marine:qnav-gravity-marine}}\label{\detokenize{modules/gravity/marine::doc}}\index{module@\spxentry{module}!qnav.gravity.marine@\spxentry{qnav.gravity.marine}}\index{qnav.gravity.marine@\spxentry{qnav.gravity.marine}!module@\spxentry{module}}

\paragraph{marine.py}
\label{\detokenize{modules/gravity/marine:marine-py}}\begin{quote}\begin{description}
\sphinxlineitem{summary}
\sphinxAtStartPar
Marine Gravity Map reading, interpolation and gravity functions.
This module contains all the gravity calculation functionalities related
to the marine gravity dataset. The Marine Gravity Map database contains
records of positions and free\sphinxhyphen{}air gravity anomalies, uncertainties in
anomalies and vertical gravity gradients. Data was last updated
August 2022.

\sphinxlineitem{author}
\sphinxAtStartPar
Michael Wright \sphinxhyphen{} \sphinxhref{mailto:mjwright@liverpool.ac.uk}{mjwright@liverpool.ac.uk}

\sphinxlineitem{version}
\sphinxAtStartPar
1.0.0 \sphinxhyphen{} First implementation of this module.

\end{description}\end{quote}
\index{MarineGravity (class in qnav.gravity.marine)@\spxentry{MarineGravity}\spxextra{class in qnav.gravity.marine}}

\begin{fulllineitems}
\phantomsection\label{\detokenize{modules/gravity/marine:qnav.gravity.marine.MarineGravity}}
\pysigstartsignatures
\pysiglinewithargsret{\sphinxbfcode{\sphinxupquote{class\DUrole{w}{ }}}\sphinxcode{\sphinxupquote{qnav.gravity.marine.}}\sphinxbfcode{\sphinxupquote{MarineGravity}}}{\sphinxparam{\DUrole{n}{base\_model}\DUrole{p}{:}\DUrole{w}{ }\DUrole{n}{{\hyperref[\detokenize{modules/gravity/base:qnav.gravity.base.GravityModel}]{\sphinxcrossref{GravityModel}}}}}\sphinxparamcomma \sphinxparam{\DUrole{n}{geoid\_model}\DUrole{p}{:}\DUrole{w}{ }\DUrole{n}{{\hyperref[\detokenize{modules/earth/geoid:qnav.earth.geoid.GeoidModel}]{\sphinxcrossref{GeoidModel}}}\DUrole{w}{ }\DUrole{p}{|}\DUrole{w}{ }None}}\sphinxparamcomma \sphinxparam{\DUrole{n}{map\_dir}\DUrole{p}{:}\DUrole{w}{ }\DUrole{n}{Path}}\sphinxparamcomma \sphinxparam{\DUrole{n}{auto\_download}\DUrole{p}{:}\DUrole{w}{ }\DUrole{n}{bool}\DUrole{w}{ }\DUrole{o}{=}\DUrole{w}{ }\DUrole{default_value}{True}}\sphinxparamcomma \sphinxparam{\DUrole{n}{auto\_load}\DUrole{p}{:}\DUrole{w}{ }\DUrole{n}{bool}\DUrole{w}{ }\DUrole{o}{=}\DUrole{w}{ }\DUrole{default_value}{True}}}{}
\pysigstopsignatures
\sphinxAtStartPar
Gravity Map implementation for the Marine Gravity dataset.
Provides corrections for a given base gravity function, using the
marine gravity map anomaly values.
\index{calc\_gravity\_z() (qnav.gravity.marine.MarineGravity method)@\spxentry{calc\_gravity\_z()}\spxextra{qnav.gravity.marine.MarineGravity method}}

\begin{fulllineitems}
\phantomsection\label{\detokenize{modules/gravity/marine:qnav.gravity.marine.MarineGravity.calc_gravity_z}}
\pysigstartsignatures
\pysiglinewithargsret{\sphinxbfcode{\sphinxupquote{calc\_gravity\_z}}}{\sphinxparam{\DUrole{n}{lat}\DUrole{p}{:}\DUrole{w}{ }\DUrole{n}{float}}\sphinxparamcomma \sphinxparam{\DUrole{n}{lon}\DUrole{p}{:}\DUrole{w}{ }\DUrole{n}{float}}\sphinxparamcomma \sphinxparam{\DUrole{n}{alt}\DUrole{p}{:}\DUrole{w}{ }\DUrole{n}{float}}}{{ $\rightarrow$ float}}
\pysigstopsignatures
\sphinxAtStartPar
Calculates and returns the total gravitational acceleration.
Computes the gravity along the Z/Down axis (in m/s\textasciicircum{}2) for
the given location and altitude offset.
\begin{quote}\begin{description}
\sphinxlineitem{Parameters}\begin{itemize}
\item {} 
\sphinxAtStartPar
\sphinxstyleliteralstrong{\sphinxupquote{lat}} (\sphinxstyleliteralemphasis{\sphinxupquote{float}}) \textendash{} The latitude coordinate for the position of interest.
The latitude position value given in decimal degrees.

\item {} 
\sphinxAtStartPar
\sphinxstyleliteralstrong{\sphinxupquote{lon}} (\sphinxstyleliteralemphasis{\sphinxupquote{float}}) \textendash{} The longitude coordinate for the position of interest.
The longitude position value given in decimal degrees.

\item {} 
\sphinxAtStartPar
\sphinxstyleliteralstrong{\sphinxupquote{alt}} (\sphinxstyleliteralemphasis{\sphinxupquote{float}}) \textendash{} The altitude offset for the position of interest.
The elevation from the WGS\sphinxhyphen{}84 reference ellipsoid in metres.

\end{itemize}

\sphinxlineitem{Returns}
\sphinxAtStartPar
The downwards gravitational acceleration for each location.

\sphinxlineitem{Return type}
\sphinxAtStartPar
float

\end{description}\end{quote}

\end{fulllineitems}

\index{calc\_gravity\_z\_vec() (qnav.gravity.marine.MarineGravity method)@\spxentry{calc\_gravity\_z\_vec()}\spxextra{qnav.gravity.marine.MarineGravity method}}

\begin{fulllineitems}
\phantomsection\label{\detokenize{modules/gravity/marine:qnav.gravity.marine.MarineGravity.calc_gravity_z_vec}}
\pysigstartsignatures
\pysiglinewithargsret{\sphinxbfcode{\sphinxupquote{calc\_gravity\_z\_vec}}}{\sphinxparam{\DUrole{n}{lat}\DUrole{p}{:}\DUrole{w}{ }\DUrole{n}{ndarray}}\sphinxparamcomma \sphinxparam{\DUrole{n}{lon}\DUrole{p}{:}\DUrole{w}{ }\DUrole{n}{ndarray}}\sphinxparamcomma \sphinxparam{\DUrole{n}{alt}\DUrole{p}{:}\DUrole{w}{ }\DUrole{n}{ndarray}}}{{ $\rightarrow$ ndarray}}
\pysigstopsignatures
\sphinxAtStartPar
In a vectorised batch, calculates and returns the gravitational
acceleration along the Z/Down axis (in m/s\textasciicircum{}2) for the given location
and altitude offset. Ideal when calculating gravity for a large
number of positions.
\begin{quote}\begin{description}
\sphinxlineitem{Parameters}\begin{itemize}
\item {} 
\sphinxAtStartPar
\sphinxstyleliteralstrong{\sphinxupquote{lat}} (\sphinxstyleliteralemphasis{\sphinxupquote{np.array}}) \textendash{} The latitude coordinate for the position of interest.
The latitude position value given in decimal degrees.

\item {} 
\sphinxAtStartPar
\sphinxstyleliteralstrong{\sphinxupquote{lon}} (\sphinxstyleliteralemphasis{\sphinxupquote{np.array}}) \textendash{} The longitude coordinate for the position of interest.
The longitude position value given in decimal degrees.

\item {} 
\sphinxAtStartPar
\sphinxstyleliteralstrong{\sphinxupquote{alt}} (\sphinxstyleliteralemphasis{\sphinxupquote{np.array}}) \textendash{} The altitude offset for the position of interest.
The elevation from the WGS\sphinxhyphen{}84 reference ellipsoid in metres.

\end{itemize}

\sphinxlineitem{Returns}
\sphinxAtStartPar
The downwards gravitational acceleration for each location
(in m/s\textasciicircum{}2).

\sphinxlineitem{Return type}
\sphinxAtStartPar
np.array (N\sphinxhyphen{}elements)

\end{description}\end{quote}

\end{fulllineitems}

\index{get\_map\_value() (qnav.gravity.marine.MarineGravity method)@\spxentry{get\_map\_value()}\spxextra{qnav.gravity.marine.MarineGravity method}}

\begin{fulllineitems}
\phantomsection\label{\detokenize{modules/gravity/marine:qnav.gravity.marine.MarineGravity.get_map_value}}
\pysigstartsignatures
\pysiglinewithargsret{\sphinxbfcode{\sphinxupquote{get\_map\_value}}}{\sphinxparam{\DUrole{n}{lat}\DUrole{p}{:}\DUrole{w}{ }\DUrole{n}{float\DUrole{w}{ }\DUrole{p}{|}\DUrole{w}{ }ndarray}}\sphinxparamcomma \sphinxparam{\DUrole{n}{lon}\DUrole{p}{:}\DUrole{w}{ }\DUrole{n}{float\DUrole{w}{ }\DUrole{p}{|}\DUrole{w}{ }ndarray}}}{{ $\rightarrow$ float\DUrole{w}{ }\DUrole{p}{|}\DUrole{w}{ }ndarray}}
\pysigstopsignatures
\sphinxAtStartPar
Interpolates the gravity anomaly for requested position(s).
Looks\sphinxhyphen{}up and returns the gravity anomaly value for requested latitude
and longitude positions. Supports both scalar and array input.
\begin{quote}\begin{description}
\sphinxlineitem{Parameters}\begin{itemize}
\item {} 
\sphinxAtStartPar
\sphinxstyleliteralstrong{\sphinxupquote{lat}} (\sphinxstyleliteralemphasis{\sphinxupquote{float}}\sphinxstyleliteralemphasis{\sphinxupquote{ | }}\sphinxstyleliteralemphasis{\sphinxupquote{np.ndarray}}) \textendash{} The latitude coordinate(s) in decimal degrees.

\item {} 
\sphinxAtStartPar
\sphinxstyleliteralstrong{\sphinxupquote{lon}} (\sphinxstyleliteralemphasis{\sphinxupquote{float}}\sphinxstyleliteralemphasis{\sphinxupquote{ | }}\sphinxstyleliteralemphasis{\sphinxupquote{np.ndarray}}) \textendash{} The longitude coordinate(s) in decimal degrees.

\end{itemize}

\sphinxlineitem{Returns}
\sphinxAtStartPar
The interpolated gravity anomaly value(s) in m/s\textasciicircum{}2.

\sphinxlineitem{Return type}
\sphinxAtStartPar
float | np.ndarray

\end{description}\end{quote}

\end{fulllineitems}

\index{get\_anomaly() (qnav.gravity.marine.MarineGravity method)@\spxentry{get\_anomaly()}\spxextra{qnav.gravity.marine.MarineGravity method}}

\begin{fulllineitems}
\phantomsection\label{\detokenize{modules/gravity/marine:qnav.gravity.marine.MarineGravity.get_anomaly}}
\pysigstartsignatures
\pysiglinewithargsret{\sphinxbfcode{\sphinxupquote{get\_anomaly}}}{\sphinxparam{\DUrole{n}{lat}\DUrole{p}{:}\DUrole{w}{ }\DUrole{n}{float}}\sphinxparamcomma \sphinxparam{\DUrole{n}{lon}\DUrole{p}{:}\DUrole{w}{ }\DUrole{n}{float}}}{{ $\rightarrow$ float}}
\pysigstopsignatures
\sphinxAtStartPar
Calculates the gravity anomaly for given position.
Retrieves the gravity anomaly value for the requested position
(in m/s\textasciicircum{}2) at the base surface. For gravity maps this value is
interpolated from its dataset, following possible conversion.
\begin{quote}\begin{description}
\sphinxlineitem{Parameters}\begin{itemize}
\item {} 
\sphinxAtStartPar
\sphinxstyleliteralstrong{\sphinxupquote{lat}} (\sphinxstyleliteralemphasis{\sphinxupquote{float}}) \textendash{} The latitude coordinate for the position of interest.
The latitude position value given in decimal degrees.

\item {} 
\sphinxAtStartPar
\sphinxstyleliteralstrong{\sphinxupquote{lon}} (\sphinxstyleliteralemphasis{\sphinxupquote{float}}) \textendash{} The longitude coordinate for the position of interest.
The longitude position value given in decimal degrees.

\end{itemize}

\sphinxlineitem{Returns}
\sphinxAtStartPar
The gravity anomaly value (in m/s\textasciicircum{}2).

\sphinxlineitem{Return type}
\sphinxAtStartPar
float

\end{description}\end{quote}

\end{fulllineitems}

\index{get\_anomaly\_vec() (qnav.gravity.marine.MarineGravity method)@\spxentry{get\_anomaly\_vec()}\spxextra{qnav.gravity.marine.MarineGravity method}}

\begin{fulllineitems}
\phantomsection\label{\detokenize{modules/gravity/marine:qnav.gravity.marine.MarineGravity.get_anomaly_vec}}
\pysigstartsignatures
\pysiglinewithargsret{\sphinxbfcode{\sphinxupquote{get\_anomaly\_vec}}}{\sphinxparam{\DUrole{n}{lat}\DUrole{p}{:}\DUrole{w}{ }\DUrole{n}{ndarray}}\sphinxparamcomma \sphinxparam{\DUrole{n}{lon}\DUrole{p}{:}\DUrole{w}{ }\DUrole{n}{ndarray}}}{{ $\rightarrow$ ndarray}}
\pysigstopsignatures
\sphinxAtStartPar
Calculates the gravity anomaly for multiple given positions.
Retrieves the gravity anomaly value for the requested position
(in m/s\textasciicircum{}2) at the base surface. For gravity maps this value is
interpolated from its dataset, following possible conversion.
\begin{quote}\begin{description}
\sphinxlineitem{Parameters}\begin{itemize}
\item {} 
\sphinxAtStartPar
\sphinxstyleliteralstrong{\sphinxupquote{lat}} (\sphinxstyleliteralemphasis{\sphinxupquote{float}}) \textendash{} The latitude coordinate for the position of interest.
The latitude position value given in decimal degrees.

\item {} 
\sphinxAtStartPar
\sphinxstyleliteralstrong{\sphinxupquote{lon}} (\sphinxstyleliteralemphasis{\sphinxupquote{float}}) \textendash{} The longitude coordinate for the position of interest.
The longitude position value given in decimal degrees.

\end{itemize}

\sphinxlineitem{Returns}
\sphinxAtStartPar
The gravity anomaly value (in m/s\textasciicircum{}2).

\sphinxlineitem{Return type}
\sphinxAtStartPar
float

\end{description}\end{quote}

\end{fulllineitems}

\index{get\_disturbance() (qnav.gravity.marine.MarineGravity method)@\spxentry{get\_disturbance()}\spxextra{qnav.gravity.marine.MarineGravity method}}

\begin{fulllineitems}
\phantomsection\label{\detokenize{modules/gravity/marine:qnav.gravity.marine.MarineGravity.get_disturbance}}
\pysigstartsignatures
\pysiglinewithargsret{\sphinxbfcode{\sphinxupquote{get\_disturbance}}}{\sphinxparam{\DUrole{n}{lat}\DUrole{p}{:}\DUrole{w}{ }\DUrole{n}{float}}\sphinxparamcomma \sphinxparam{\DUrole{n}{lon}\DUrole{p}{:}\DUrole{w}{ }\DUrole{n}{float}}}{{ $\rightarrow$ float}}
\pysigstopsignatures
\sphinxAtStartPar
Calculates the gravity disturbance for given position.
Retrieves the gravity disturbance value for the requested position
(in m/s\textasciicircum{}2) at the base surface. For gravity maps this value is
interpolated from its dataset, following possible conversion.
\begin{quote}\begin{description}
\sphinxlineitem{Parameters}\begin{itemize}
\item {} 
\sphinxAtStartPar
\sphinxstyleliteralstrong{\sphinxupquote{lat}} (\sphinxstyleliteralemphasis{\sphinxupquote{float}}) \textendash{} The latitude coordinate for the position of interest.
The latitude position value given in decimal degrees.

\item {} 
\sphinxAtStartPar
\sphinxstyleliteralstrong{\sphinxupquote{lon}} (\sphinxstyleliteralemphasis{\sphinxupquote{float}}) \textendash{} The longitude coordinate for the position of interest.
The longitude position value given in decimal degrees.

\end{itemize}

\sphinxlineitem{Returns}
\sphinxAtStartPar
The gravity disturbance value (in m/s\textasciicircum{}2).

\sphinxlineitem{Return type}
\sphinxAtStartPar
float

\end{description}\end{quote}

\end{fulllineitems}

\index{get\_disturbance\_vec() (qnav.gravity.marine.MarineGravity method)@\spxentry{get\_disturbance\_vec()}\spxextra{qnav.gravity.marine.MarineGravity method}}

\begin{fulllineitems}
\phantomsection\label{\detokenize{modules/gravity/marine:qnav.gravity.marine.MarineGravity.get_disturbance_vec}}
\pysigstartsignatures
\pysiglinewithargsret{\sphinxbfcode{\sphinxupquote{get\_disturbance\_vec}}}{\sphinxparam{\DUrole{n}{lat}\DUrole{p}{:}\DUrole{w}{ }\DUrole{n}{ndarray}}\sphinxparamcomma \sphinxparam{\DUrole{n}{lon}\DUrole{p}{:}\DUrole{w}{ }\DUrole{n}{ndarray}}}{{ $\rightarrow$ ndarray}}
\pysigstopsignatures
\sphinxAtStartPar
Calculates the gravity disturbance for multiple given positions.
Retrieves the gravity disturbance value for the requested position
(in m/s\textasciicircum{}2) at the base surface. For gravity maps this value is
interpolated from its dataset, following possible conversion.
\begin{quote}\begin{description}
\sphinxlineitem{Parameters}\begin{itemize}
\item {} 
\sphinxAtStartPar
\sphinxstyleliteralstrong{\sphinxupquote{lat}} (\sphinxstyleliteralemphasis{\sphinxupquote{float}}) \textendash{} The latitude coordinate for the position of interest.
The latitude position value given in decimal degrees.

\item {} 
\sphinxAtStartPar
\sphinxstyleliteralstrong{\sphinxupquote{lon}} (\sphinxstyleliteralemphasis{\sphinxupquote{float}}) \textendash{} The longitude coordinate for the position of interest.
The longitude position value given in decimal degrees.

\end{itemize}

\sphinxlineitem{Returns}
\sphinxAtStartPar
The gravity disturbance value (in m/s\textasciicircum{}2).

\sphinxlineitem{Return type}
\sphinxAtStartPar
float

\end{description}\end{quote}

\end{fulllineitems}

\index{load\_data() (qnav.gravity.marine.MarineGravity method)@\spxentry{load\_data()}\spxextra{qnav.gravity.marine.MarineGravity method}}

\begin{fulllineitems}
\phantomsection\label{\detokenize{modules/gravity/marine:qnav.gravity.marine.MarineGravity.load_data}}
\pysigstartsignatures
\pysiglinewithargsret{\sphinxbfcode{\sphinxupquote{load\_data}}}{}{}
\pysigstopsignatures
\sphinxAtStartPar
Loads the Marine Gravity data and generates interpolator for look\sphinxhyphen{}up.
Called on initialisation to load the required data to allow for
correction interpolation.

\end{fulllineitems}

\index{download\_data() (qnav.gravity.marine.MarineGravity method)@\spxentry{download\_data()}\spxextra{qnav.gravity.marine.MarineGravity method}}

\begin{fulllineitems}
\phantomsection\label{\detokenize{modules/gravity/marine:qnav.gravity.marine.MarineGravity.download_data}}
\pysigstartsignatures
\pysiglinewithargsret{\sphinxbfcode{\sphinxupquote{download\_data}}}{\sphinxparam{\DUrole{n}{overwrite}\DUrole{p}{:}\DUrole{w}{ }\DUrole{n}{bool}\DUrole{w}{ }\DUrole{o}{=}\DUrole{w}{ }\DUrole{default_value}{False}}}{}
\pysigstopsignatures
\sphinxAtStartPar
Downloads the required marine gravity files from the source.
If allowed and files are not present, will download the required
geosat files to the set directory. Optionally, files can be
overwritten.
\begin{quote}\begin{description}
\sphinxlineitem{Parameters}
\sphinxAtStartPar
\sphinxstyleliteralstrong{\sphinxupquote{overwrite}} (\sphinxstyleliteralemphasis{\sphinxupquote{bool}}) \textendash{} Should downloading replace existing files.
By default this is set to false and only missing files
are downloaded.

\end{description}\end{quote}

\end{fulllineitems}


\end{fulllineitems}

\index{read\_marine\_file() (in module qnav.gravity.marine)@\spxentry{read\_marine\_file()}\spxextra{in module qnav.gravity.marine}}

\begin{fulllineitems}
\phantomsection\label{\detokenize{modules/gravity/marine:qnav.gravity.marine.read_marine_file}}
\pysigstartsignatures
\pysiglinewithargsret{\sphinxcode{\sphinxupquote{qnav.gravity.marine.}}\sphinxbfcode{\sphinxupquote{read\_marine\_file}}}{\sphinxparam{\DUrole{n}{file\_path}\DUrole{p}{:}\DUrole{w}{ }\DUrole{n}{Path}}}{{ $\rightarrow$ dict\DUrole{p}{{[}}str\DUrole{p}{,}\DUrole{w}{ }ndarray\DUrole{p}{{]}}}}
\pysigstopsignatures
\sphinxAtStartPar
This function extracts the data from a Marine gravity dataset. Marine
gravity files are encoded in HDF5 scientific data format. Once
successfully read, the gravity anomaly matrix, along with the latitude
and longitude positions will be returned. The gravity values are in
metres/seconds\textasciicircum{}2.
\begin{quote}\begin{description}
\sphinxlineitem{Parameters}
\sphinxAtStartPar
\sphinxstyleliteralstrong{\sphinxupquote{file\_path}} (\sphinxstyleliteralemphasis{\sphinxupquote{Path}}) \textendash{} The location of the Marine gravity dataset to read.

\sphinxlineitem{Returns}
\sphinxAtStartPar
A dictionary containing the marine gravity map data, consisting
of gravity anomaly, latitude and longitude values.

\sphinxlineitem{Return type}
\sphinxAtStartPar
dict{[}str, np.ndarray{]}

\end{description}\end{quote}

\end{fulllineitems}


\sphinxstepscope


\subsubsection{qnav.gravity.ggm\_plus}
\label{\detokenize{modules/gravity/ggm_plus:module-qnav.gravity.ggm_plus}}\label{\detokenize{modules/gravity/ggm_plus:qnav-gravity-ggm-plus}}\label{\detokenize{modules/gravity/ggm_plus::doc}}\index{module@\spxentry{module}!qnav.gravity.ggm\_plus@\spxentry{qnav.gravity.ggm\_plus}}\index{qnav.gravity.ggm\_plus@\spxentry{qnav.gravity.ggm\_plus}!module@\spxentry{module}}

\paragraph{ggm\_plus.py}
\label{\detokenize{modules/gravity/ggm_plus:ggm-plus-py}}\begin{quote}\begin{description}
\sphinxlineitem{summary}
\sphinxAtStartPar
Global Gravity Map Plus (GGMPlus) gravity map implementations.
This module provides support for the GGMPlus gravity models. Specifically,
two implementations for the GGPlus gravity acceleration and gravity
disturbance corrections.

\sphinxlineitem{author}
\sphinxAtStartPar
Michael Wright \sphinxhyphen{} \sphinxhref{mailto:mjwright@liverpool.ac.uk}{mjwright@liverpool.ac.uk}

\sphinxlineitem{version}
\sphinxAtStartPar
1.0.0 \sphinxhyphen{} First implementation of this module.

\end{description}\end{quote}
\index{GGMPlusAcc (class in qnav.gravity.ggm\_plus)@\spxentry{GGMPlusAcc}\spxextra{class in qnav.gravity.ggm\_plus}}

\begin{fulllineitems}
\phantomsection\label{\detokenize{modules/gravity/ggm_plus:qnav.gravity.ggm_plus.GGMPlusAcc}}
\pysigstartsignatures
\pysiglinewithargsret{\sphinxbfcode{\sphinxupquote{class\DUrole{w}{ }}}\sphinxcode{\sphinxupquote{qnav.gravity.ggm\_plus.}}\sphinxbfcode{\sphinxupquote{GGMPlusAcc}}}{\sphinxparam{\DUrole{n}{base\_model}\DUrole{p}{:}\DUrole{w}{ }\DUrole{n}{{\hyperref[\detokenize{modules/gravity/base:qnav.gravity.base.GravityModel}]{\sphinxcrossref{GravityModel}}}}}\sphinxparamcomma \sphinxparam{\DUrole{n}{geoid\_model}\DUrole{p}{:}\DUrole{w}{ }\DUrole{n}{{\hyperref[\detokenize{modules/earth/geoid:qnav.earth.geoid.GeoidModel}]{\sphinxcrossref{GeoidModel}}}\DUrole{w}{ }\DUrole{p}{|}\DUrole{w}{ }None}}\sphinxparamcomma \sphinxparam{\DUrole{n}{map\_dir}\DUrole{p}{:}\DUrole{w}{ }\DUrole{n}{Path}}\sphinxparamcomma \sphinxparam{\DUrole{n}{auto\_download}\DUrole{p}{:}\DUrole{w}{ }\DUrole{n}{bool}\DUrole{w}{ }\DUrole{o}{=}\DUrole{w}{ }\DUrole{default_value}{True}}\sphinxparamcomma \sphinxparam{\DUrole{n}{margin\_size}\DUrole{p}{:}\DUrole{w}{ }\DUrole{n}{int}\DUrole{w}{ }\DUrole{o}{=}\DUrole{w}{ }\DUrole{default_value}{1}}}{}
\pysigstopsignatures
\sphinxAtStartPar
Global Gravity Map Plus for Acceleration Correction.
An implementation of GGMPlus, using the provided base gravity
acceleration values. This overwrites the base values zero\sphinxhyphen{}altitude
gravity estimation, but uses its altitude extrapolation.
\index{calc\_gravity\_z() (qnav.gravity.ggm\_plus.GGMPlusAcc method)@\spxentry{calc\_gravity\_z()}\spxextra{qnav.gravity.ggm\_plus.GGMPlusAcc method}}

\begin{fulllineitems}
\phantomsection\label{\detokenize{modules/gravity/ggm_plus:qnav.gravity.ggm_plus.GGMPlusAcc.calc_gravity_z}}
\pysigstartsignatures
\pysiglinewithargsret{\sphinxbfcode{\sphinxupquote{calc\_gravity\_z}}}{\sphinxparam{\DUrole{n}{lat}\DUrole{p}{:}\DUrole{w}{ }\DUrole{n}{float}}\sphinxparamcomma \sphinxparam{\DUrole{n}{lon}\DUrole{p}{:}\DUrole{w}{ }\DUrole{n}{float}}\sphinxparamcomma \sphinxparam{\DUrole{n}{alt}\DUrole{p}{:}\DUrole{w}{ }\DUrole{n}{float}}}{{ $\rightarrow$ float}}
\pysigstopsignatures
\sphinxAtStartPar
Calculates and returns the total gravitational acceleration.
Computes the gravity along the Z/Down axis (in m/s\textasciicircum{}2) for
the given location and altitude offset.
\begin{quote}\begin{description}
\sphinxlineitem{Parameters}\begin{itemize}
\item {} 
\sphinxAtStartPar
\sphinxstyleliteralstrong{\sphinxupquote{lat}} (\sphinxstyleliteralemphasis{\sphinxupquote{float}}) \textendash{} The latitude coordinate for the position of interest.
The latitude position value given in decimal degrees.

\item {} 
\sphinxAtStartPar
\sphinxstyleliteralstrong{\sphinxupquote{lon}} (\sphinxstyleliteralemphasis{\sphinxupquote{float}}) \textendash{} The longitude coordinate for the position of interest.
The longitude position value given in decimal degrees.

\item {} 
\sphinxAtStartPar
\sphinxstyleliteralstrong{\sphinxupquote{alt}} (\sphinxstyleliteralemphasis{\sphinxupquote{float}}) \textendash{} The altitude offset for the position of interest.
The elevation from the WGS\sphinxhyphen{}84 reference ellipsoid in metres.

\end{itemize}

\sphinxlineitem{Returns}
\sphinxAtStartPar
The downwards gravitational acceleration for each location.

\sphinxlineitem{Return type}
\sphinxAtStartPar
float

\end{description}\end{quote}

\end{fulllineitems}

\index{calc\_gravity\_z\_vec() (qnav.gravity.ggm\_plus.GGMPlusAcc method)@\spxentry{calc\_gravity\_z\_vec()}\spxextra{qnav.gravity.ggm\_plus.GGMPlusAcc method}}

\begin{fulllineitems}
\phantomsection\label{\detokenize{modules/gravity/ggm_plus:qnav.gravity.ggm_plus.GGMPlusAcc.calc_gravity_z_vec}}
\pysigstartsignatures
\pysiglinewithargsret{\sphinxbfcode{\sphinxupquote{calc\_gravity\_z\_vec}}}{\sphinxparam{\DUrole{n}{lat}\DUrole{p}{:}\DUrole{w}{ }\DUrole{n}{ndarray}}\sphinxparamcomma \sphinxparam{\DUrole{n}{lon}\DUrole{p}{:}\DUrole{w}{ }\DUrole{n}{ndarray}}\sphinxparamcomma \sphinxparam{\DUrole{n}{alt}\DUrole{p}{:}\DUrole{w}{ }\DUrole{n}{ndarray}}}{{ $\rightarrow$ ndarray}}
\pysigstopsignatures
\sphinxAtStartPar
In a vectorised batch, calculates and returns the gravitational
acceleration along the Z/Down axis (in m/s\textasciicircum{}2) for the given location
and altitude offset. Ideal when calculating gravity for a large
number of positions.
\begin{quote}\begin{description}
\sphinxlineitem{Parameters}\begin{itemize}
\item {} 
\sphinxAtStartPar
\sphinxstyleliteralstrong{\sphinxupquote{lat}} (\sphinxstyleliteralemphasis{\sphinxupquote{np.array}}) \textendash{} The latitude coordinate for the position of interest.
The latitude position value given in decimal degrees.

\item {} 
\sphinxAtStartPar
\sphinxstyleliteralstrong{\sphinxupquote{lon}} (\sphinxstyleliteralemphasis{\sphinxupquote{np.array}}) \textendash{} The longitude coordinate for the position of interest.
The longitude position value given in decimal degrees.

\item {} 
\sphinxAtStartPar
\sphinxstyleliteralstrong{\sphinxupquote{alt}} (\sphinxstyleliteralemphasis{\sphinxupquote{np.array}}) \textendash{} The altitude offset for the position of interest.
The elevation from the WGS\sphinxhyphen{}84 reference ellipsoid in metres.

\end{itemize}

\sphinxlineitem{Returns}
\sphinxAtStartPar
The downwards gravitational acceleration for each location
(in m/s\textasciicircum{}2).

\sphinxlineitem{Return type}
\sphinxAtStartPar
np.array (N\sphinxhyphen{}elements)

\end{description}\end{quote}

\end{fulllineitems}

\index{get\_anomaly() (qnav.gravity.ggm\_plus.GGMPlusAcc method)@\spxentry{get\_anomaly()}\spxextra{qnav.gravity.ggm\_plus.GGMPlusAcc method}}

\begin{fulllineitems}
\phantomsection\label{\detokenize{modules/gravity/ggm_plus:qnav.gravity.ggm_plus.GGMPlusAcc.get_anomaly}}
\pysigstartsignatures
\pysiglinewithargsret{\sphinxbfcode{\sphinxupquote{get\_anomaly}}}{\sphinxparam{\DUrole{n}{lat}\DUrole{p}{:}\DUrole{w}{ }\DUrole{n}{float}}\sphinxparamcomma \sphinxparam{\DUrole{n}{lon}\DUrole{p}{:}\DUrole{w}{ }\DUrole{n}{float}}}{{ $\rightarrow$ float}}
\pysigstopsignatures
\sphinxAtStartPar
Calculates the gravity anomaly for given position.
Retrieves the gravity anomaly value for the requested position
(in m/s\textasciicircum{}2) at the base surface. For gravity maps this value is
interpolated from its dataset, following possible conversion.
\begin{quote}\begin{description}
\sphinxlineitem{Parameters}\begin{itemize}
\item {} 
\sphinxAtStartPar
\sphinxstyleliteralstrong{\sphinxupquote{lat}} (\sphinxstyleliteralemphasis{\sphinxupquote{float}}) \textendash{} The latitude coordinate for the position of interest.
The latitude position value given in decimal degrees.

\item {} 
\sphinxAtStartPar
\sphinxstyleliteralstrong{\sphinxupquote{lon}} (\sphinxstyleliteralemphasis{\sphinxupquote{float}}) \textendash{} The longitude coordinate for the position of interest.
The longitude position value given in decimal degrees.

\end{itemize}

\sphinxlineitem{Returns}
\sphinxAtStartPar
The gravity anomaly value (in m/s\textasciicircum{}2).

\sphinxlineitem{Return type}
\sphinxAtStartPar
float

\end{description}\end{quote}

\end{fulllineitems}

\index{get\_anomaly\_vec() (qnav.gravity.ggm\_plus.GGMPlusAcc method)@\spxentry{get\_anomaly\_vec()}\spxextra{qnav.gravity.ggm\_plus.GGMPlusAcc method}}

\begin{fulllineitems}
\phantomsection\label{\detokenize{modules/gravity/ggm_plus:qnav.gravity.ggm_plus.GGMPlusAcc.get_anomaly_vec}}
\pysigstartsignatures
\pysiglinewithargsret{\sphinxbfcode{\sphinxupquote{get\_anomaly\_vec}}}{\sphinxparam{\DUrole{n}{lat}\DUrole{p}{:}\DUrole{w}{ }\DUrole{n}{ndarray}}\sphinxparamcomma \sphinxparam{\DUrole{n}{lon}\DUrole{p}{:}\DUrole{w}{ }\DUrole{n}{ndarray}}}{{ $\rightarrow$ ndarray}}
\pysigstopsignatures
\sphinxAtStartPar
Calculates the gravity anomaly for multiple given positions.
Retrieves the gravity anomaly value for the requested position
(in m/s\textasciicircum{}2) at the base surface. For gravity maps this value is
interpolated from its dataset, following possible conversion.
\begin{quote}\begin{description}
\sphinxlineitem{Parameters}\begin{itemize}
\item {} 
\sphinxAtStartPar
\sphinxstyleliteralstrong{\sphinxupquote{lat}} (\sphinxstyleliteralemphasis{\sphinxupquote{float}}) \textendash{} The latitude coordinate for the position of interest.
The latitude position value given in decimal degrees.

\item {} 
\sphinxAtStartPar
\sphinxstyleliteralstrong{\sphinxupquote{lon}} (\sphinxstyleliteralemphasis{\sphinxupquote{float}}) \textendash{} The longitude coordinate for the position of interest.
The longitude position value given in decimal degrees.

\end{itemize}

\sphinxlineitem{Returns}
\sphinxAtStartPar
The gravity anomaly value (in m/s\textasciicircum{}2).

\sphinxlineitem{Return type}
\sphinxAtStartPar
float

\end{description}\end{quote}

\end{fulllineitems}

\index{get\_disturbance() (qnav.gravity.ggm\_plus.GGMPlusAcc method)@\spxentry{get\_disturbance()}\spxextra{qnav.gravity.ggm\_plus.GGMPlusAcc method}}

\begin{fulllineitems}
\phantomsection\label{\detokenize{modules/gravity/ggm_plus:qnav.gravity.ggm_plus.GGMPlusAcc.get_disturbance}}
\pysigstartsignatures
\pysiglinewithargsret{\sphinxbfcode{\sphinxupquote{get\_disturbance}}}{\sphinxparam{\DUrole{n}{lat}\DUrole{p}{:}\DUrole{w}{ }\DUrole{n}{float}}\sphinxparamcomma \sphinxparam{\DUrole{n}{lon}\DUrole{p}{:}\DUrole{w}{ }\DUrole{n}{float}}}{{ $\rightarrow$ float}}
\pysigstopsignatures
\sphinxAtStartPar
Calculates the gravity disturbance for given position.
Retrieves the gravity disturbance value for the requested position
(in m/s\textasciicircum{}2) at the base surface. For gravity maps this value is
interpolated from its dataset, following possible conversion.
\begin{quote}\begin{description}
\sphinxlineitem{Parameters}\begin{itemize}
\item {} 
\sphinxAtStartPar
\sphinxstyleliteralstrong{\sphinxupquote{lat}} (\sphinxstyleliteralemphasis{\sphinxupquote{float}}) \textendash{} The latitude coordinate for the position of interest.
The latitude position value given in decimal degrees.

\item {} 
\sphinxAtStartPar
\sphinxstyleliteralstrong{\sphinxupquote{lon}} (\sphinxstyleliteralemphasis{\sphinxupquote{float}}) \textendash{} The longitude coordinate for the position of interest.
The longitude position value given in decimal degrees.

\end{itemize}

\sphinxlineitem{Returns}
\sphinxAtStartPar
The gravity disturbance value (in m/s\textasciicircum{}2).

\sphinxlineitem{Return type}
\sphinxAtStartPar
float

\end{description}\end{quote}

\end{fulllineitems}

\index{get\_disturbance\_vec() (qnav.gravity.ggm\_plus.GGMPlusAcc method)@\spxentry{get\_disturbance\_vec()}\spxextra{qnav.gravity.ggm\_plus.GGMPlusAcc method}}

\begin{fulllineitems}
\phantomsection\label{\detokenize{modules/gravity/ggm_plus:qnav.gravity.ggm_plus.GGMPlusAcc.get_disturbance_vec}}
\pysigstartsignatures
\pysiglinewithargsret{\sphinxbfcode{\sphinxupquote{get\_disturbance\_vec}}}{\sphinxparam{\DUrole{n}{lat}\DUrole{p}{:}\DUrole{w}{ }\DUrole{n}{ndarray}}\sphinxparamcomma \sphinxparam{\DUrole{n}{lon}\DUrole{p}{:}\DUrole{w}{ }\DUrole{n}{ndarray}}}{{ $\rightarrow$ ndarray}}
\pysigstopsignatures
\sphinxAtStartPar
Calculates the gravity disturbance for multiple given positions.
Retrieves the gravity disturbance value for the requested position
(in m/s\textasciicircum{}2) at the base surface. For gravity maps this value is
interpolated from its dataset, following possible conversion.
\begin{quote}\begin{description}
\sphinxlineitem{Parameters}\begin{itemize}
\item {} 
\sphinxAtStartPar
\sphinxstyleliteralstrong{\sphinxupquote{lat}} (\sphinxstyleliteralemphasis{\sphinxupquote{float}}) \textendash{} The latitude coordinate for the position of interest.
The latitude position value given in decimal degrees.

\item {} 
\sphinxAtStartPar
\sphinxstyleliteralstrong{\sphinxupquote{lon}} (\sphinxstyleliteralemphasis{\sphinxupquote{float}}) \textendash{} The longitude coordinate for the position of interest.
The longitude position value given in decimal degrees.

\end{itemize}

\sphinxlineitem{Returns}
\sphinxAtStartPar
The gravity disturbance value (in m/s\textasciicircum{}2).

\sphinxlineitem{Return type}
\sphinxAtStartPar
float

\end{description}\end{quote}

\end{fulllineitems}

\index{get\_map\_value() (qnav.gravity.ggm\_plus.GGMPlusAcc method)@\spxentry{get\_map\_value()}\spxextra{qnav.gravity.ggm\_plus.GGMPlusAcc method}}

\begin{fulllineitems}
\phantomsection\label{\detokenize{modules/gravity/ggm_plus:qnav.gravity.ggm_plus.GGMPlusAcc.get_map_value}}
\pysigstartsignatures
\pysiglinewithargsret{\sphinxbfcode{\sphinxupquote{get\_map\_value}}}{\sphinxparam{\DUrole{n}{lat}\DUrole{p}{:}\DUrole{w}{ }\DUrole{n}{float}}\sphinxparamcomma \sphinxparam{\DUrole{n}{lon}\DUrole{p}{:}\DUrole{w}{ }\DUrole{n}{float}}}{{ $\rightarrow$ float}}
\pysigstopsignatures
\sphinxAtStartPar
Interpolates the gravity acceleration for requested position(.
Looks\sphinxhyphen{}up and returns the gravity anomaly value for requested latitude
and longitude position.
\begin{quote}\begin{description}
\sphinxlineitem{Parameters}\begin{itemize}
\item {} 
\sphinxAtStartPar
\sphinxstyleliteralstrong{\sphinxupquote{lat}} (\sphinxstyleliteralemphasis{\sphinxupquote{float}}) \textendash{} The latitude coordinate in decimal degrees.

\item {} 
\sphinxAtStartPar
\sphinxstyleliteralstrong{\sphinxupquote{lon}} (\sphinxstyleliteralemphasis{\sphinxupquote{float}}) \textendash{} The longitude coordinate in decimal degrees.

\end{itemize}

\sphinxlineitem{Returns}
\sphinxAtStartPar
The interpolated gravity map value in m/s\textasciicircum{}2.

\sphinxlineitem{Return type}
\sphinxAtStartPar
float

\end{description}\end{quote}

\end{fulllineitems}

\index{get\_map\_values() (qnav.gravity.ggm\_plus.GGMPlusAcc method)@\spxentry{get\_map\_values()}\spxextra{qnav.gravity.ggm\_plus.GGMPlusAcc method}}

\begin{fulllineitems}
\phantomsection\label{\detokenize{modules/gravity/ggm_plus:qnav.gravity.ggm_plus.GGMPlusAcc.get_map_values}}
\pysigstartsignatures
\pysiglinewithargsret{\sphinxbfcode{\sphinxupquote{get\_map\_values}}}{\sphinxparam{\DUrole{n}{lat}\DUrole{p}{:}\DUrole{w}{ }\DUrole{n}{ndarray}}\sphinxparamcomma \sphinxparam{\DUrole{n}{lon}\DUrole{p}{:}\DUrole{w}{ }\DUrole{n}{ndarray}}}{{ $\rightarrow$ ndarray}}
\pysigstopsignatures
\sphinxAtStartPar
Looks\sphinxhyphen{}up the correction values using the map.
This will dynamically load data to cover the range, keeping the
maximum number of loaded tiles under the margin value, to conserve
memory constraints.
\begin{quote}\begin{description}
\sphinxlineitem{Parameters}\begin{itemize}
\item {} 
\sphinxAtStartPar
\sphinxstyleliteralstrong{\sphinxupquote{lat}} (\sphinxstyleliteralemphasis{\sphinxupquote{np.ndarray}}) \textendash{} The latitude coordinate in decimal degrees.

\item {} 
\sphinxAtStartPar
\sphinxstyleliteralstrong{\sphinxupquote{lon}} (\sphinxstyleliteralemphasis{\sphinxupquote{np.ndarray}}) \textendash{} The longitude coordinate in decimal degrees.

\end{itemize}

\sphinxlineitem{Returns}
\sphinxAtStartPar
The corresponding values interpolated from the map.

\sphinxlineitem{Return type}
\sphinxAtStartPar
np.ndarray

\end{description}\end{quote}

\end{fulllineitems}

\index{download\_data() (qnav.gravity.ggm\_plus.GGMPlusAcc method)@\spxentry{download\_data()}\spxextra{qnav.gravity.ggm\_plus.GGMPlusAcc method}}

\begin{fulllineitems}
\phantomsection\label{\detokenize{modules/gravity/ggm_plus:qnav.gravity.ggm_plus.GGMPlusAcc.download_data}}
\pysigstartsignatures
\pysiglinewithargsret{\sphinxbfcode{\sphinxupquote{download\_data}}}{\sphinxparam{\DUrole{n}{origin\_lat}\DUrole{p}{:}\DUrole{w}{ }\DUrole{n}{int}}\sphinxparamcomma \sphinxparam{\DUrole{n}{origin\_lon}\DUrole{p}{:}\DUrole{w}{ }\DUrole{n}{int}}\sphinxparamcomma \sphinxparam{\DUrole{n}{overwrite}\DUrole{p}{:}\DUrole{w}{ }\DUrole{n}{bool}\DUrole{w}{ }\DUrole{o}{=}\DUrole{w}{ }\DUrole{default_value}{False}}}{}
\pysigstopsignatures
\sphinxAtStartPar
Downloads the required marine gravity files from the source.
If allowed and files are not present, will download the required
geosat files to the set directory. Optionally, files can be
overwritten.
\begin{quote}\begin{description}
\sphinxlineitem{Parameters}\begin{itemize}
\item {} 
\sphinxAtStartPar
\sphinxstyleliteralstrong{\sphinxupquote{origin\_lat}} (\sphinxstyleliteralemphasis{\sphinxupquote{int}}) \textendash{} Origin latitude of tile to download.

\item {} 
\sphinxAtStartPar
\sphinxstyleliteralstrong{\sphinxupquote{origin\_lon}} (\sphinxstyleliteralemphasis{\sphinxupquote{int}}) \textendash{} Origin longitude of tile to download.

\item {} 
\sphinxAtStartPar
\sphinxstyleliteralstrong{\sphinxupquote{overwrite}} (\sphinxstyleliteralemphasis{\sphinxupquote{bool}}) \textendash{} Should downloading replace existing files.
By default this is set to false and only missing files
are downloaded.

\end{itemize}

\end{description}\end{quote}

\end{fulllineitems}


\end{fulllineitems}

\index{GGMPlusDist (class in qnav.gravity.ggm\_plus)@\spxentry{GGMPlusDist}\spxextra{class in qnav.gravity.ggm\_plus}}

\begin{fulllineitems}
\phantomsection\label{\detokenize{modules/gravity/ggm_plus:qnav.gravity.ggm_plus.GGMPlusDist}}
\pysigstartsignatures
\pysiglinewithargsret{\sphinxbfcode{\sphinxupquote{class\DUrole{w}{ }}}\sphinxcode{\sphinxupquote{qnav.gravity.ggm\_plus.}}\sphinxbfcode{\sphinxupquote{GGMPlusDist}}}{\sphinxparam{\DUrole{n}{base\_model}\DUrole{p}{:}\DUrole{w}{ }\DUrole{n}{{\hyperref[\detokenize{modules/gravity/base:qnav.gravity.base.GravityModel}]{\sphinxcrossref{GravityModel}}}}}\sphinxparamcomma \sphinxparam{\DUrole{n}{geoid\_model}\DUrole{p}{:}\DUrole{w}{ }\DUrole{n}{{\hyperref[\detokenize{modules/earth/geoid:qnav.earth.geoid.GeoidModel}]{\sphinxcrossref{GeoidModel}}}\DUrole{w}{ }\DUrole{p}{|}\DUrole{w}{ }None}}\sphinxparamcomma \sphinxparam{\DUrole{n}{map\_dir}\DUrole{p}{:}\DUrole{w}{ }\DUrole{n}{Path}}\sphinxparamcomma \sphinxparam{\DUrole{n}{auto\_download}\DUrole{p}{:}\DUrole{w}{ }\DUrole{n}{bool}\DUrole{w}{ }\DUrole{o}{=}\DUrole{w}{ }\DUrole{default_value}{True}}\sphinxparamcomma \sphinxparam{\DUrole{n}{margin\_size}\DUrole{p}{:}\DUrole{w}{ }\DUrole{n}{int}\DUrole{w}{ }\DUrole{o}{=}\DUrole{w}{ }\DUrole{default_value}{1}}}{}
\pysigstopsignatures
\sphinxAtStartPar
Global Gravity Map Plus for Gravity Disturbance Correction.
An implementation of GGMPlus, using the acquired gravity disturbance
values values.
\index{get\_map\_value() (qnav.gravity.ggm\_plus.GGMPlusDist method)@\spxentry{get\_map\_value()}\spxextra{qnav.gravity.ggm\_plus.GGMPlusDist method}}

\begin{fulllineitems}
\phantomsection\label{\detokenize{modules/gravity/ggm_plus:qnav.gravity.ggm_plus.GGMPlusDist.get_map_value}}
\pysigstartsignatures
\pysiglinewithargsret{\sphinxbfcode{\sphinxupquote{get\_map\_value}}}{\sphinxparam{\DUrole{n}{lat}\DUrole{p}{:}\DUrole{w}{ }\DUrole{n}{float}}\sphinxparamcomma \sphinxparam{\DUrole{n}{lon}\DUrole{p}{:}\DUrole{w}{ }\DUrole{n}{float}}}{{ $\rightarrow$ float}}
\pysigstopsignatures
\sphinxAtStartPar
Interpolates the gravity acceleration for requested position(.
Looks\sphinxhyphen{}up and returns the gravity anomaly value for requested latitude
and longitude position.
\begin{quote}\begin{description}
\sphinxlineitem{Parameters}\begin{itemize}
\item {} 
\sphinxAtStartPar
\sphinxstyleliteralstrong{\sphinxupquote{lat}} (\sphinxstyleliteralemphasis{\sphinxupquote{float}}) \textendash{} The latitude coordinate in decimal degrees.

\item {} 
\sphinxAtStartPar
\sphinxstyleliteralstrong{\sphinxupquote{lon}} (\sphinxstyleliteralemphasis{\sphinxupquote{float}}) \textendash{} The longitude coordinate in decimal degrees.

\end{itemize}

\sphinxlineitem{Returns}
\sphinxAtStartPar
The interpolated gravity map value in m/s\textasciicircum{}2.

\sphinxlineitem{Return type}
\sphinxAtStartPar
float

\end{description}\end{quote}

\end{fulllineitems}

\index{get\_map\_values() (qnav.gravity.ggm\_plus.GGMPlusDist method)@\spxentry{get\_map\_values()}\spxextra{qnav.gravity.ggm\_plus.GGMPlusDist method}}

\begin{fulllineitems}
\phantomsection\label{\detokenize{modules/gravity/ggm_plus:qnav.gravity.ggm_plus.GGMPlusDist.get_map_values}}
\pysigstartsignatures
\pysiglinewithargsret{\sphinxbfcode{\sphinxupquote{get\_map\_values}}}{\sphinxparam{\DUrole{n}{lat}\DUrole{p}{:}\DUrole{w}{ }\DUrole{n}{ndarray}}\sphinxparamcomma \sphinxparam{\DUrole{n}{lon}\DUrole{p}{:}\DUrole{w}{ }\DUrole{n}{ndarray}}}{{ $\rightarrow$ ndarray}}
\pysigstopsignatures
\sphinxAtStartPar
Looks\sphinxhyphen{}up the correction values using the map.
This will dynamically load data to cover the range, keeping the
maximum number of loaded tiles under the margin value, to conserve
memory constraints.
\begin{quote}\begin{description}
\sphinxlineitem{Parameters}\begin{itemize}
\item {} 
\sphinxAtStartPar
\sphinxstyleliteralstrong{\sphinxupquote{lat}} (\sphinxstyleliteralemphasis{\sphinxupquote{np.ndarray}}) \textendash{} The latitude coordinate in decimal degrees.

\item {} 
\sphinxAtStartPar
\sphinxstyleliteralstrong{\sphinxupquote{lon}} (\sphinxstyleliteralemphasis{\sphinxupquote{np.ndarray}}) \textendash{} The longitude coordinate in decimal degrees.

\end{itemize}

\sphinxlineitem{Returns}
\sphinxAtStartPar
The corresponding values interpolated from the map.

\sphinxlineitem{Return type}
\sphinxAtStartPar
np.ndarray

\end{description}\end{quote}

\end{fulllineitems}

\index{calc\_gravity\_z() (qnav.gravity.ggm\_plus.GGMPlusDist method)@\spxentry{calc\_gravity\_z()}\spxextra{qnav.gravity.ggm\_plus.GGMPlusDist method}}

\begin{fulllineitems}
\phantomsection\label{\detokenize{modules/gravity/ggm_plus:qnav.gravity.ggm_plus.GGMPlusDist.calc_gravity_z}}
\pysigstartsignatures
\pysiglinewithargsret{\sphinxbfcode{\sphinxupquote{calc\_gravity\_z}}}{\sphinxparam{\DUrole{n}{lat}\DUrole{p}{:}\DUrole{w}{ }\DUrole{n}{float}}\sphinxparamcomma \sphinxparam{\DUrole{n}{lon}\DUrole{p}{:}\DUrole{w}{ }\DUrole{n}{float}}\sphinxparamcomma \sphinxparam{\DUrole{n}{alt}\DUrole{p}{:}\DUrole{w}{ }\DUrole{n}{float}}}{{ $\rightarrow$ float}}
\pysigstopsignatures
\sphinxAtStartPar
Calculates and returns the total gravitational acceleration.
Computes the gravity along the Z/Down axis (in m/s\textasciicircum{}2) for
the given location and altitude offset.
\begin{quote}\begin{description}
\sphinxlineitem{Parameters}\begin{itemize}
\item {} 
\sphinxAtStartPar
\sphinxstyleliteralstrong{\sphinxupquote{lat}} (\sphinxstyleliteralemphasis{\sphinxupquote{float}}) \textendash{} The latitude coordinate for the position of interest.
The latitude position value given in decimal degrees.

\item {} 
\sphinxAtStartPar
\sphinxstyleliteralstrong{\sphinxupquote{lon}} (\sphinxstyleliteralemphasis{\sphinxupquote{float}}) \textendash{} The longitude coordinate for the position of interest.
The longitude position value given in decimal degrees.

\item {} 
\sphinxAtStartPar
\sphinxstyleliteralstrong{\sphinxupquote{alt}} (\sphinxstyleliteralemphasis{\sphinxupquote{float}}) \textendash{} The altitude offset for the position of interest.
The elevation from the WGS\sphinxhyphen{}84 reference ellipsoid in metres.

\end{itemize}

\sphinxlineitem{Returns}
\sphinxAtStartPar
The downwards gravitational acceleration for each location.

\sphinxlineitem{Return type}
\sphinxAtStartPar
float

\end{description}\end{quote}

\end{fulllineitems}

\index{calc\_gravity\_z\_vec() (qnav.gravity.ggm\_plus.GGMPlusDist method)@\spxentry{calc\_gravity\_z\_vec()}\spxextra{qnav.gravity.ggm\_plus.GGMPlusDist method}}

\begin{fulllineitems}
\phantomsection\label{\detokenize{modules/gravity/ggm_plus:qnav.gravity.ggm_plus.GGMPlusDist.calc_gravity_z_vec}}
\pysigstartsignatures
\pysiglinewithargsret{\sphinxbfcode{\sphinxupquote{calc\_gravity\_z\_vec}}}{\sphinxparam{\DUrole{n}{lat}\DUrole{p}{:}\DUrole{w}{ }\DUrole{n}{ndarray}}\sphinxparamcomma \sphinxparam{\DUrole{n}{lon}\DUrole{p}{:}\DUrole{w}{ }\DUrole{n}{ndarray}}\sphinxparamcomma \sphinxparam{\DUrole{n}{alt}\DUrole{p}{:}\DUrole{w}{ }\DUrole{n}{ndarray}}}{{ $\rightarrow$ ndarray}}
\pysigstopsignatures
\sphinxAtStartPar
In a vectorised batch, calculates and returns the gravitational
acceleration along the Z/Down axis (in m/s\textasciicircum{}2) for the given location
and altitude offset. Ideal when calculating gravity for a large
number of positions.
\begin{quote}\begin{description}
\sphinxlineitem{Parameters}\begin{itemize}
\item {} 
\sphinxAtStartPar
\sphinxstyleliteralstrong{\sphinxupquote{lat}} (\sphinxstyleliteralemphasis{\sphinxupquote{np.array}}) \textendash{} The latitude coordinate for the position of interest.
The latitude position value given in decimal degrees.

\item {} 
\sphinxAtStartPar
\sphinxstyleliteralstrong{\sphinxupquote{lon}} (\sphinxstyleliteralemphasis{\sphinxupquote{np.array}}) \textendash{} The longitude coordinate for the position of interest.
The longitude position value given in decimal degrees.

\item {} 
\sphinxAtStartPar
\sphinxstyleliteralstrong{\sphinxupquote{alt}} (\sphinxstyleliteralemphasis{\sphinxupquote{np.array}}) \textendash{} The altitude offset for the position of interest.
The elevation from the WGS\sphinxhyphen{}84 reference ellipsoid in metres.

\end{itemize}

\sphinxlineitem{Returns}
\sphinxAtStartPar
The downwards gravitational acceleration for each location
(in m/s\textasciicircum{}2).

\sphinxlineitem{Return type}
\sphinxAtStartPar
np.array (N\sphinxhyphen{}elements)

\end{description}\end{quote}

\end{fulllineitems}

\index{get\_disturbance() (qnav.gravity.ggm\_plus.GGMPlusDist method)@\spxentry{get\_disturbance()}\spxextra{qnav.gravity.ggm\_plus.GGMPlusDist method}}

\begin{fulllineitems}
\phantomsection\label{\detokenize{modules/gravity/ggm_plus:qnav.gravity.ggm_plus.GGMPlusDist.get_disturbance}}
\pysigstartsignatures
\pysiglinewithargsret{\sphinxbfcode{\sphinxupquote{get\_disturbance}}}{\sphinxparam{\DUrole{n}{lat}\DUrole{p}{:}\DUrole{w}{ }\DUrole{n}{float}}\sphinxparamcomma \sphinxparam{\DUrole{n}{lon}\DUrole{p}{:}\DUrole{w}{ }\DUrole{n}{float}}}{{ $\rightarrow$ float}}
\pysigstopsignatures
\sphinxAtStartPar
Calculates the gravity disturbance for given position.
Retrieves the gravity disturbance value for the requested position
(in m/s\textasciicircum{}2) at the base surface. For gravity maps this value is
interpolated from its dataset, following possible conversion.
\begin{quote}\begin{description}
\sphinxlineitem{Parameters}\begin{itemize}
\item {} 
\sphinxAtStartPar
\sphinxstyleliteralstrong{\sphinxupquote{lat}} (\sphinxstyleliteralemphasis{\sphinxupquote{float}}) \textendash{} The latitude coordinate for the position of interest.
The latitude position value given in decimal degrees.

\item {} 
\sphinxAtStartPar
\sphinxstyleliteralstrong{\sphinxupquote{lon}} (\sphinxstyleliteralemphasis{\sphinxupquote{float}}) \textendash{} The longitude coordinate for the position of interest.
The longitude position value given in decimal degrees.

\end{itemize}

\sphinxlineitem{Returns}
\sphinxAtStartPar
The gravity disturbance value (in m/s\textasciicircum{}2).

\sphinxlineitem{Return type}
\sphinxAtStartPar
float

\end{description}\end{quote}

\end{fulllineitems}

\index{get\_disturbance\_vec() (qnav.gravity.ggm\_plus.GGMPlusDist method)@\spxentry{get\_disturbance\_vec()}\spxextra{qnav.gravity.ggm\_plus.GGMPlusDist method}}

\begin{fulllineitems}
\phantomsection\label{\detokenize{modules/gravity/ggm_plus:qnav.gravity.ggm_plus.GGMPlusDist.get_disturbance_vec}}
\pysigstartsignatures
\pysiglinewithargsret{\sphinxbfcode{\sphinxupquote{get\_disturbance\_vec}}}{\sphinxparam{\DUrole{n}{lat}\DUrole{p}{:}\DUrole{w}{ }\DUrole{n}{ndarray}}\sphinxparamcomma \sphinxparam{\DUrole{n}{lon}\DUrole{p}{:}\DUrole{w}{ }\DUrole{n}{ndarray}}}{{ $\rightarrow$ ndarray}}
\pysigstopsignatures
\sphinxAtStartPar
Calculates the gravity disturbance for multiple given positions.
Retrieves the gravity disturbance value for the requested position
(in m/s\textasciicircum{}2) at the base surface. For gravity maps this value is
interpolated from its dataset, following possible conversion.
\begin{quote}\begin{description}
\sphinxlineitem{Parameters}\begin{itemize}
\item {} 
\sphinxAtStartPar
\sphinxstyleliteralstrong{\sphinxupquote{lat}} (\sphinxstyleliteralemphasis{\sphinxupquote{float}}) \textendash{} The latitude coordinate for the position of interest.
The latitude position value given in decimal degrees.

\item {} 
\sphinxAtStartPar
\sphinxstyleliteralstrong{\sphinxupquote{lon}} (\sphinxstyleliteralemphasis{\sphinxupquote{float}}) \textendash{} The longitude coordinate for the position of interest.
The longitude position value given in decimal degrees.

\end{itemize}

\sphinxlineitem{Returns}
\sphinxAtStartPar
The gravity disturbance value (in m/s\textasciicircum{}2).

\sphinxlineitem{Return type}
\sphinxAtStartPar
float

\end{description}\end{quote}

\end{fulllineitems}


\end{fulllineitems}

\index{read\_ggm\_file() (in module qnav.gravity.ggm\_plus)@\spxentry{read\_ggm\_file()}\spxextra{in module qnav.gravity.ggm\_plus}}

\begin{fulllineitems}
\phantomsection\label{\detokenize{modules/gravity/ggm_plus:qnav.gravity.ggm_plus.read_ggm_file}}
\pysigstartsignatures
\pysiglinewithargsret{\sphinxcode{\sphinxupquote{qnav.gravity.ggm\_plus.}}\sphinxbfcode{\sphinxupquote{read\_ggm\_file}}}{\sphinxparam{\DUrole{n}{file\_path}\DUrole{p}{:}\DUrole{w}{ }\DUrole{n}{Path}}}{{ $\rightarrow$ ndarray}}
\pysigstopsignatures
\sphinxAtStartPar
Reads the data from a given GGM\_plus file. Once read, the returned
2500\sphinxhyphen{}by\sphinxhyphen{}2500 matrix data will be in \(metres/second^2\). No latitude
or longitude coordinates are contained within any GGM\_plus file.
\begin{quote}\begin{description}
\sphinxlineitem{Parameters}
\sphinxAtStartPar
\sphinxstyleliteralstrong{\sphinxupquote{file\_path}} (\sphinxstyleliteralemphasis{\sphinxupquote{Path}}\sphinxstyleliteralemphasis{\sphinxupquote{ or }}\sphinxstyleliteralemphasis{\sphinxupquote{str}}) \textendash{} The location of the GGM\_plus file to read.

\sphinxlineitem{Returns}
\sphinxAtStartPar
The gravitational acceleration matrix read from the file,
measured in \(metres/second^2\)

\sphinxlineitem{Return type}
\sphinxAtStartPar
NumPy Matrix (2500\sphinxhyphen{}by\sphinxhyphen{}2500 elements)

\end{description}\end{quote}

\end{fulllineitems}

\index{get\_empty\_ggm\_tile() (in module qnav.gravity.ggm\_plus)@\spxentry{get\_empty\_ggm\_tile()}\spxextra{in module qnav.gravity.ggm\_plus}}

\begin{fulllineitems}
\phantomsection\label{\detokenize{modules/gravity/ggm_plus:qnav.gravity.ggm_plus.get_empty_ggm_tile}}
\pysigstartsignatures
\pysiglinewithargsret{\sphinxcode{\sphinxupquote{qnav.gravity.ggm\_plus.}}\sphinxbfcode{\sphinxupquote{get\_empty\_ggm\_tile}}}{}{{ $\rightarrow$ ndarray}}
\pysigstopsignatures
\sphinxAtStartPar
Produces an empty tile of dimensions of GGM Plus data file.
Used when data is unavailable and cannot be downloaded.
\begin{quote}\begin{description}
\sphinxlineitem{Returns}
\sphinxAtStartPar
Empty GGM Plus tile data for missing files.

\sphinxlineitem{Return type}
\sphinxAtStartPar
NumPy Matrix (2500\sphinxhyphen{}by\sphinxhyphen{}2500 elements)

\end{description}\end{quote}

\end{fulllineitems}


\sphinxstepscope


\subsubsection{qnav.gravity.srtm2gravity}
\label{\detokenize{modules/gravity/srtm2gravity:module-qnav.gravity.srtm2gravity}}\label{\detokenize{modules/gravity/srtm2gravity:qnav-gravity-srtm2gravity}}\label{\detokenize{modules/gravity/srtm2gravity::doc}}\index{module@\spxentry{module}!qnav.gravity.srtm2gravity@\spxentry{qnav.gravity.srtm2gravity}}\index{qnav.gravity.srtm2gravity@\spxentry{qnav.gravity.srtm2gravity}!module@\spxentry{module}}

\paragraph{srtm2gravity.py}
\label{\detokenize{modules/gravity/srtm2gravity:srtm2gravity-py}}\begin{quote}\begin{description}
\sphinxlineitem{summary}
\sphinxAtStartPar
SRTM2Gravity gravity map implementations.
This module provides support for the SRTM2Gravity gravity models.
Specifically, the full\sphinxhyphen{}scale and residual gravity maps.

\sphinxlineitem{author}
\sphinxAtStartPar
Michael Wright \sphinxhyphen{} \sphinxhref{mailto:mjwright@liverpool.ac.uk}{mjwright@liverpool.ac.uk}

\sphinxlineitem{version}
\sphinxAtStartPar
1.0.0 \sphinxhyphen{} First implementation of this module.

\end{description}\end{quote}
\index{SRTM2GravityFull (class in qnav.gravity.srtm2gravity)@\spxentry{SRTM2GravityFull}\spxextra{class in qnav.gravity.srtm2gravity}}

\begin{fulllineitems}
\phantomsection\label{\detokenize{modules/gravity/srtm2gravity:qnav.gravity.srtm2gravity.SRTM2GravityFull}}
\pysigstartsignatures
\pysiglinewithargsret{\sphinxbfcode{\sphinxupquote{class\DUrole{w}{ }}}\sphinxcode{\sphinxupquote{qnav.gravity.srtm2gravity.}}\sphinxbfcode{\sphinxupquote{SRTM2GravityFull}}}{\sphinxparam{\DUrole{n}{base\_model}\DUrole{p}{:}\DUrole{w}{ }\DUrole{n}{{\hyperref[\detokenize{modules/gravity/base:qnav.gravity.base.GravityModel}]{\sphinxcrossref{GravityModel}}}}}\sphinxparamcomma \sphinxparam{\DUrole{n}{geoid\_model}\DUrole{p}{:}\DUrole{w}{ }\DUrole{n}{{\hyperref[\detokenize{modules/earth/geoid:qnav.earth.geoid.GeoidModel}]{\sphinxcrossref{GeoidModel}}}\DUrole{w}{ }\DUrole{p}{|}\DUrole{w}{ }None}}\sphinxparamcomma \sphinxparam{\DUrole{n}{map\_dir}\DUrole{p}{:}\DUrole{w}{ }\DUrole{n}{Path}}\sphinxparamcomma \sphinxparam{\DUrole{n}{auto\_download}\DUrole{p}{:}\DUrole{w}{ }\DUrole{n}{bool}\DUrole{w}{ }\DUrole{o}{=}\DUrole{w}{ }\DUrole{default_value}{True}}\sphinxparamcomma \sphinxparam{\DUrole{n}{margin\_size}\DUrole{p}{:}\DUrole{w}{ }\DUrole{n}{int}\DUrole{w}{ }\DUrole{o}{=}\DUrole{w}{ }\DUrole{default_value}{1}}}{}
\pysigstopsignatures
\sphinxAtStartPar
Implementation of SRTM2Gravity Full\sphinxhyphen{}Scale gravity.
Provides a gravity correction map that uses the SRTM2Gravity Full\sphinxhyphen{}Scale
dataset to provide gravity disturbance values.

\end{fulllineitems}

\index{SRTM2GravityRes (class in qnav.gravity.srtm2gravity)@\spxentry{SRTM2GravityRes}\spxextra{class in qnav.gravity.srtm2gravity}}

\begin{fulllineitems}
\phantomsection\label{\detokenize{modules/gravity/srtm2gravity:qnav.gravity.srtm2gravity.SRTM2GravityRes}}
\pysigstartsignatures
\pysiglinewithargsret{\sphinxbfcode{\sphinxupquote{class\DUrole{w}{ }}}\sphinxcode{\sphinxupquote{qnav.gravity.srtm2gravity.}}\sphinxbfcode{\sphinxupquote{SRTM2GravityRes}}}{\sphinxparam{\DUrole{n}{base\_model}\DUrole{p}{:}\DUrole{w}{ }\DUrole{n}{{\hyperref[\detokenize{modules/gravity/base:qnav.gravity.base.GravityModel}]{\sphinxcrossref{GravityModel}}}}}\sphinxparamcomma \sphinxparam{\DUrole{n}{geoid\_model}\DUrole{p}{:}\DUrole{w}{ }\DUrole{n}{{\hyperref[\detokenize{modules/earth/geoid:qnav.earth.geoid.GeoidModel}]{\sphinxcrossref{GeoidModel}}}\DUrole{w}{ }\DUrole{p}{|}\DUrole{w}{ }None}}\sphinxparamcomma \sphinxparam{\DUrole{n}{map\_dir}\DUrole{p}{:}\DUrole{w}{ }\DUrole{n}{Path}}\sphinxparamcomma \sphinxparam{\DUrole{n}{auto\_download}\DUrole{p}{:}\DUrole{w}{ }\DUrole{n}{bool}\DUrole{w}{ }\DUrole{o}{=}\DUrole{w}{ }\DUrole{default_value}{True}}\sphinxparamcomma \sphinxparam{\DUrole{n}{margin\_size}\DUrole{p}{:}\DUrole{w}{ }\DUrole{n}{int}\DUrole{w}{ }\DUrole{o}{=}\DUrole{w}{ }\DUrole{default_value}{1}}}{}
\pysigstopsignatures
\end{fulllineitems}

\index{read\_srtm\_file() (in module qnav.gravity.srtm2gravity)@\spxentry{read\_srtm\_file()}\spxextra{in module qnav.gravity.srtm2gravity}}

\begin{fulllineitems}
\phantomsection\label{\detokenize{modules/gravity/srtm2gravity:qnav.gravity.srtm2gravity.read_srtm_file}}
\pysigstartsignatures
\pysiglinewithargsret{\sphinxcode{\sphinxupquote{qnav.gravity.srtm2gravity.}}\sphinxbfcode{\sphinxupquote{read\_srtm\_file}}}{\sphinxparam{\DUrole{n}{file\_path}\DUrole{p}{:}\DUrole{w}{ }\DUrole{n}{Path}}}{{ $\rightarrow$ ndarray}}
\pysigstopsignatures
\sphinxAtStartPar
Reads the data from a given SRTM2Gravity file. Once read, the returned
1200\sphinxhyphen{}by\sphinxhyphen{}1200 matrix data will be in \(metres/second^2\). No latitude
or longitude coordinates are contained within any SRTM2Gravity file.
\begin{quote}\begin{description}
\sphinxlineitem{Parameters}
\sphinxAtStartPar
\sphinxstyleliteralstrong{\sphinxupquote{file\_path}} (\sphinxstyleliteralemphasis{\sphinxupquote{Path}}\sphinxstyleliteralemphasis{\sphinxupquote{ or }}\sphinxstyleliteralemphasis{\sphinxupquote{str}}) \textendash{} The location of the GGM\_plus file to read.

\sphinxlineitem{Returns}
\sphinxAtStartPar
The gravity disturbance matrix read from the file,
measured in \(metres/second^2\)

\sphinxlineitem{Return type}
\sphinxAtStartPar
NumPy Matrix (1200\sphinxhyphen{}by\sphinxhyphen{}1200 elements)

\end{description}\end{quote}

\end{fulllineitems}

\index{get\_empty\_srtm\_tile() (in module qnav.gravity.srtm2gravity)@\spxentry{get\_empty\_srtm\_tile()}\spxextra{in module qnav.gravity.srtm2gravity}}

\begin{fulllineitems}
\phantomsection\label{\detokenize{modules/gravity/srtm2gravity:qnav.gravity.srtm2gravity.get_empty_srtm_tile}}
\pysigstartsignatures
\pysiglinewithargsret{\sphinxcode{\sphinxupquote{qnav.gravity.srtm2gravity.}}\sphinxbfcode{\sphinxupquote{get\_empty\_srtm\_tile}}}{}{{ $\rightarrow$ ndarray}}
\pysigstopsignatures
\sphinxAtStartPar
Produces an empty tile of dimensions of GGM Plus data file.
Used when data is unavailable and cannot be downloaded.
\begin{quote}\begin{description}
\sphinxlineitem{Returns}
\sphinxAtStartPar
Empty GGM Plus tile data for missing files.

\sphinxlineitem{Return type}
\sphinxAtStartPar
NumPy Matrix (1200\sphinxhyphen{}by\sphinxhyphen{}1200 elements)

\end{description}\end{quote}

\end{fulllineitems}


\sphinxstepscope


\subsubsection{qnav.gravity.irish}
\label{\detokenize{modules/gravity/irish:module-qnav.gravity.irish}}\label{\detokenize{modules/gravity/irish:qnav-gravity-irish}}\label{\detokenize{modules/gravity/irish::doc}}\index{module@\spxentry{module}!qnav.gravity.irish@\spxentry{qnav.gravity.irish}}\index{qnav.gravity.irish@\spxentry{qnav.gravity.irish}!module@\spxentry{module}}\index{IrishSeaGravity (class in qnav.gravity.irish)@\spxentry{IrishSeaGravity}\spxextra{class in qnav.gravity.irish}}

\begin{fulllineitems}
\phantomsection\label{\detokenize{modules/gravity/irish:qnav.gravity.irish.IrishSeaGravity}}
\pysigstartsignatures
\pysiglinewithargsret{\sphinxbfcode{\sphinxupquote{class\DUrole{w}{ }}}\sphinxcode{\sphinxupquote{qnav.gravity.irish.}}\sphinxbfcode{\sphinxupquote{IrishSeaGravity}}}{\sphinxparam{\DUrole{n}{base\_model}\DUrole{p}{:}\DUrole{w}{ }\DUrole{n}{{\hyperref[\detokenize{modules/gravity/base:qnav.gravity.base.GravityModel}]{\sphinxcrossref{GravityModel}}}}}\sphinxparamcomma \sphinxparam{\DUrole{n}{geoid\_model}\DUrole{p}{:}\DUrole{w}{ }\DUrole{n}{{\hyperref[\detokenize{modules/earth/geoid:qnav.earth.geoid.GeoidModel}]{\sphinxcrossref{GeoidModel}}}\DUrole{w}{ }\DUrole{p}{|}\DUrole{w}{ }None}}\sphinxparamcomma \sphinxparam{\DUrole{n}{map\_dir}\DUrole{p}{:}\DUrole{w}{ }\DUrole{n}{Path\DUrole{w}{ }\DUrole{p}{|}\DUrole{w}{ }None}}}{}
\pysigstopsignatures\index{load\_data() (qnav.gravity.irish.IrishSeaGravity method)@\spxentry{load\_data()}\spxextra{qnav.gravity.irish.IrishSeaGravity method}}

\begin{fulllineitems}
\phantomsection\label{\detokenize{modules/gravity/irish:qnav.gravity.irish.IrishSeaGravity.load_data}}
\pysigstartsignatures
\pysiglinewithargsret{\sphinxbfcode{\sphinxupquote{load\_data}}}{}{}
\pysigstopsignatures
\sphinxAtStartPar
Loads the Marine Gravity data and generates interpolator for look\sphinxhyphen{}up.
Called on initialisation to load the required data to allow for
correction interpolation.

\end{fulllineitems}

\index{get\_map\_value() (qnav.gravity.irish.IrishSeaGravity method)@\spxentry{get\_map\_value()}\spxextra{qnav.gravity.irish.IrishSeaGravity method}}

\begin{fulllineitems}
\phantomsection\label{\detokenize{modules/gravity/irish:qnav.gravity.irish.IrishSeaGravity.get_map_value}}
\pysigstartsignatures
\pysiglinewithargsret{\sphinxbfcode{\sphinxupquote{get\_map\_value}}}{\sphinxparam{\DUrole{n}{lat}\DUrole{p}{:}\DUrole{w}{ }\DUrole{n}{float\DUrole{w}{ }\DUrole{p}{|}\DUrole{w}{ }ndarray}}\sphinxparamcomma \sphinxparam{\DUrole{n}{lon}\DUrole{p}{:}\DUrole{w}{ }\DUrole{n}{float\DUrole{w}{ }\DUrole{p}{|}\DUrole{w}{ }ndarray}}}{{ $\rightarrow$ float\DUrole{w}{ }\DUrole{p}{|}\DUrole{w}{ }ndarray}}
\pysigstopsignatures
\sphinxAtStartPar
Interpolates the gravity anomaly for requested position(s).
Looks\sphinxhyphen{}up and returns the gravity anomaly value for requested latitude
and longitude positions. Supports both scalar and array input.
\begin{quote}\begin{description}
\sphinxlineitem{Parameters}\begin{itemize}
\item {} 
\sphinxAtStartPar
\sphinxstyleliteralstrong{\sphinxupquote{lat}} (\sphinxstyleliteralemphasis{\sphinxupquote{float}}\sphinxstyleliteralemphasis{\sphinxupquote{ | }}\sphinxstyleliteralemphasis{\sphinxupquote{np.ndarray}}) \textendash{} The latitude coordinate(s) in decimal degrees.

\item {} 
\sphinxAtStartPar
\sphinxstyleliteralstrong{\sphinxupquote{lon}} (\sphinxstyleliteralemphasis{\sphinxupquote{float}}\sphinxstyleliteralemphasis{\sphinxupquote{ | }}\sphinxstyleliteralemphasis{\sphinxupquote{np.ndarray}}) \textendash{} The longitude coordinate(s) in decimal degrees.

\end{itemize}

\sphinxlineitem{Returns}
\sphinxAtStartPar
The interpolated gravity anomaly value(s) in m/s\textasciicircum{}2.

\sphinxlineitem{Return type}
\sphinxAtStartPar
float | np.ndarray

\end{description}\end{quote}

\end{fulllineitems}

\index{calc\_gravity\_z() (qnav.gravity.irish.IrishSeaGravity method)@\spxentry{calc\_gravity\_z()}\spxextra{qnav.gravity.irish.IrishSeaGravity method}}

\begin{fulllineitems}
\phantomsection\label{\detokenize{modules/gravity/irish:qnav.gravity.irish.IrishSeaGravity.calc_gravity_z}}
\pysigstartsignatures
\pysiglinewithargsret{\sphinxbfcode{\sphinxupquote{calc\_gravity\_z}}}{\sphinxparam{\DUrole{n}{lat}\DUrole{p}{:}\DUrole{w}{ }\DUrole{n}{float}}\sphinxparamcomma \sphinxparam{\DUrole{n}{lon}\DUrole{p}{:}\DUrole{w}{ }\DUrole{n}{float}}\sphinxparamcomma \sphinxparam{\DUrole{n}{alt}\DUrole{p}{:}\DUrole{w}{ }\DUrole{n}{float}}}{{ $\rightarrow$ float}}
\pysigstopsignatures
\sphinxAtStartPar
Calculates and returns the total gravitational acceleration.
Computes the gravity along the Z/Down axis (in m/s\textasciicircum{}2) for
the given location and altitude offset.
\begin{quote}\begin{description}
\sphinxlineitem{Parameters}\begin{itemize}
\item {} 
\sphinxAtStartPar
\sphinxstyleliteralstrong{\sphinxupquote{lat}} (\sphinxstyleliteralemphasis{\sphinxupquote{float}}) \textendash{} The latitude coordinate for the position of interest.
The latitude position value given in decimal degrees.

\item {} 
\sphinxAtStartPar
\sphinxstyleliteralstrong{\sphinxupquote{lon}} (\sphinxstyleliteralemphasis{\sphinxupquote{float}}) \textendash{} The longitude coordinate for the position of interest.
The longitude position value given in decimal degrees.

\item {} 
\sphinxAtStartPar
\sphinxstyleliteralstrong{\sphinxupquote{alt}} (\sphinxstyleliteralemphasis{\sphinxupquote{float}}) \textendash{} The altitude offset for the position of interest.
The elevation from the WGS\sphinxhyphen{}84 reference ellipsoid in metres.

\end{itemize}

\sphinxlineitem{Returns}
\sphinxAtStartPar
The downwards gravitational acceleration for each location.

\sphinxlineitem{Return type}
\sphinxAtStartPar
float

\end{description}\end{quote}

\end{fulllineitems}

\index{calc\_gravity\_z\_vec() (qnav.gravity.irish.IrishSeaGravity method)@\spxentry{calc\_gravity\_z\_vec()}\spxextra{qnav.gravity.irish.IrishSeaGravity method}}

\begin{fulllineitems}
\phantomsection\label{\detokenize{modules/gravity/irish:qnav.gravity.irish.IrishSeaGravity.calc_gravity_z_vec}}
\pysigstartsignatures
\pysiglinewithargsret{\sphinxbfcode{\sphinxupquote{calc\_gravity\_z\_vec}}}{\sphinxparam{\DUrole{n}{lat}\DUrole{p}{:}\DUrole{w}{ }\DUrole{n}{ndarray}}\sphinxparamcomma \sphinxparam{\DUrole{n}{lon}\DUrole{p}{:}\DUrole{w}{ }\DUrole{n}{ndarray}}\sphinxparamcomma \sphinxparam{\DUrole{n}{alt}\DUrole{p}{:}\DUrole{w}{ }\DUrole{n}{ndarray}}}{{ $\rightarrow$ ndarray}}
\pysigstopsignatures
\sphinxAtStartPar
In a vectorised batch, calculates and returns the gravitational
acceleration along the Z/Down axis (in m/s\textasciicircum{}2) for the given location
and altitude offset. Ideal when calculating gravity for a large
number of positions.
\begin{quote}\begin{description}
\sphinxlineitem{Parameters}\begin{itemize}
\item {} 
\sphinxAtStartPar
\sphinxstyleliteralstrong{\sphinxupquote{lat}} (\sphinxstyleliteralemphasis{\sphinxupquote{np.array}}) \textendash{} The latitude coordinate for the position of interest.
The latitude position value given in decimal degrees.

\item {} 
\sphinxAtStartPar
\sphinxstyleliteralstrong{\sphinxupquote{lon}} (\sphinxstyleliteralemphasis{\sphinxupquote{np.array}}) \textendash{} The longitude coordinate for the position of interest.
The longitude position value given in decimal degrees.

\item {} 
\sphinxAtStartPar
\sphinxstyleliteralstrong{\sphinxupquote{alt}} (\sphinxstyleliteralemphasis{\sphinxupquote{np.array}}) \textendash{} The altitude offset for the position of interest.
The elevation from the WGS\sphinxhyphen{}84 reference ellipsoid in metres.

\end{itemize}

\sphinxlineitem{Returns}
\sphinxAtStartPar
The downwards gravitational acceleration for each location
(in m/s\textasciicircum{}2).

\sphinxlineitem{Return type}
\sphinxAtStartPar
np.array (N\sphinxhyphen{}elements)

\end{description}\end{quote}

\end{fulllineitems}

\index{get\_anomaly() (qnav.gravity.irish.IrishSeaGravity method)@\spxentry{get\_anomaly()}\spxextra{qnav.gravity.irish.IrishSeaGravity method}}

\begin{fulllineitems}
\phantomsection\label{\detokenize{modules/gravity/irish:qnav.gravity.irish.IrishSeaGravity.get_anomaly}}
\pysigstartsignatures
\pysiglinewithargsret{\sphinxbfcode{\sphinxupquote{get\_anomaly}}}{\sphinxparam{\DUrole{n}{lat}\DUrole{p}{:}\DUrole{w}{ }\DUrole{n}{float}}\sphinxparamcomma \sphinxparam{\DUrole{n}{lon}\DUrole{p}{:}\DUrole{w}{ }\DUrole{n}{float}}}{{ $\rightarrow$ float}}
\pysigstopsignatures
\sphinxAtStartPar
Calculates the gravity anomaly for given position.
Retrieves the gravity anomaly value for the requested position
(in m/s\textasciicircum{}2) at the base surface. For gravity maps this value is
interpolated from its dataset, following possible conversion.
\begin{quote}\begin{description}
\sphinxlineitem{Parameters}\begin{itemize}
\item {} 
\sphinxAtStartPar
\sphinxstyleliteralstrong{\sphinxupquote{lat}} (\sphinxstyleliteralemphasis{\sphinxupquote{float}}) \textendash{} The latitude coordinate for the position of interest.
The latitude position value given in decimal degrees.

\item {} 
\sphinxAtStartPar
\sphinxstyleliteralstrong{\sphinxupquote{lon}} (\sphinxstyleliteralemphasis{\sphinxupquote{float}}) \textendash{} The longitude coordinate for the position of interest.
The longitude position value given in decimal degrees.

\end{itemize}

\sphinxlineitem{Returns}
\sphinxAtStartPar
The gravity anomaly value (in m/s\textasciicircum{}2).

\sphinxlineitem{Return type}
\sphinxAtStartPar
float

\end{description}\end{quote}

\end{fulllineitems}

\index{get\_anomaly\_vec() (qnav.gravity.irish.IrishSeaGravity method)@\spxentry{get\_anomaly\_vec()}\spxextra{qnav.gravity.irish.IrishSeaGravity method}}

\begin{fulllineitems}
\phantomsection\label{\detokenize{modules/gravity/irish:qnav.gravity.irish.IrishSeaGravity.get_anomaly_vec}}
\pysigstartsignatures
\pysiglinewithargsret{\sphinxbfcode{\sphinxupquote{get\_anomaly\_vec}}}{\sphinxparam{\DUrole{n}{lat}\DUrole{p}{:}\DUrole{w}{ }\DUrole{n}{ndarray}}\sphinxparamcomma \sphinxparam{\DUrole{n}{lon}\DUrole{p}{:}\DUrole{w}{ }\DUrole{n}{ndarray}}}{{ $\rightarrow$ ndarray}}
\pysigstopsignatures
\sphinxAtStartPar
Calculates the gravity anomaly for multiple given positions.
Retrieves the gravity anomaly value for the requested position
(in m/s\textasciicircum{}2) at the base surface. For gravity maps this value is
interpolated from its dataset, following possible conversion.
\begin{quote}\begin{description}
\sphinxlineitem{Parameters}\begin{itemize}
\item {} 
\sphinxAtStartPar
\sphinxstyleliteralstrong{\sphinxupquote{lat}} (\sphinxstyleliteralemphasis{\sphinxupquote{float}}) \textendash{} The latitude coordinate for the position of interest.
The latitude position value given in decimal degrees.

\item {} 
\sphinxAtStartPar
\sphinxstyleliteralstrong{\sphinxupquote{lon}} (\sphinxstyleliteralemphasis{\sphinxupquote{float}}) \textendash{} The longitude coordinate for the position of interest.
The longitude position value given in decimal degrees.

\end{itemize}

\sphinxlineitem{Returns}
\sphinxAtStartPar
The gravity anomaly value (in m/s\textasciicircum{}2).

\sphinxlineitem{Return type}
\sphinxAtStartPar
float

\end{description}\end{quote}

\end{fulllineitems}

\index{get\_disturbance() (qnav.gravity.irish.IrishSeaGravity method)@\spxentry{get\_disturbance()}\spxextra{qnav.gravity.irish.IrishSeaGravity method}}

\begin{fulllineitems}
\phantomsection\label{\detokenize{modules/gravity/irish:qnav.gravity.irish.IrishSeaGravity.get_disturbance}}
\pysigstartsignatures
\pysiglinewithargsret{\sphinxbfcode{\sphinxupquote{get\_disturbance}}}{\sphinxparam{\DUrole{n}{lat}\DUrole{p}{:}\DUrole{w}{ }\DUrole{n}{float}}\sphinxparamcomma \sphinxparam{\DUrole{n}{lon}\DUrole{p}{:}\DUrole{w}{ }\DUrole{n}{float}}}{{ $\rightarrow$ float}}
\pysigstopsignatures
\sphinxAtStartPar
Calculates the gravity disturbance for given position.
Retrieves the gravity disturbance value for the requested position
(in m/s\textasciicircum{}2) at the base surface. For gravity maps this value is
interpolated from its dataset, following possible conversion.
\begin{quote}\begin{description}
\sphinxlineitem{Parameters}\begin{itemize}
\item {} 
\sphinxAtStartPar
\sphinxstyleliteralstrong{\sphinxupquote{lat}} (\sphinxstyleliteralemphasis{\sphinxupquote{float}}) \textendash{} The latitude coordinate for the position of interest.
The latitude position value given in decimal degrees.

\item {} 
\sphinxAtStartPar
\sphinxstyleliteralstrong{\sphinxupquote{lon}} (\sphinxstyleliteralemphasis{\sphinxupquote{float}}) \textendash{} The longitude coordinate for the position of interest.
The longitude position value given in decimal degrees.

\end{itemize}

\sphinxlineitem{Returns}
\sphinxAtStartPar
The gravity disturbance value (in m/s\textasciicircum{}2).

\sphinxlineitem{Return type}
\sphinxAtStartPar
float

\end{description}\end{quote}

\end{fulllineitems}

\index{get\_disturbance\_vec() (qnav.gravity.irish.IrishSeaGravity method)@\spxentry{get\_disturbance\_vec()}\spxextra{qnav.gravity.irish.IrishSeaGravity method}}

\begin{fulllineitems}
\phantomsection\label{\detokenize{modules/gravity/irish:qnav.gravity.irish.IrishSeaGravity.get_disturbance_vec}}
\pysigstartsignatures
\pysiglinewithargsret{\sphinxbfcode{\sphinxupquote{get\_disturbance\_vec}}}{\sphinxparam{\DUrole{n}{lat}\DUrole{p}{:}\DUrole{w}{ }\DUrole{n}{ndarray}}\sphinxparamcomma \sphinxparam{\DUrole{n}{lon}\DUrole{p}{:}\DUrole{w}{ }\DUrole{n}{ndarray}}}{{ $\rightarrow$ ndarray}}
\pysigstopsignatures
\sphinxAtStartPar
Calculates the gravity disturbance for multiple given positions.
Retrieves the gravity disturbance value for the requested position
(in m/s\textasciicircum{}2) at the base surface. For gravity maps this value is
interpolated from its dataset, following possible conversion.
\begin{quote}\begin{description}
\sphinxlineitem{Parameters}\begin{itemize}
\item {} 
\sphinxAtStartPar
\sphinxstyleliteralstrong{\sphinxupquote{lat}} (\sphinxstyleliteralemphasis{\sphinxupquote{float}}) \textendash{} The latitude coordinate for the position of interest.
The latitude position value given in decimal degrees.

\item {} 
\sphinxAtStartPar
\sphinxstyleliteralstrong{\sphinxupquote{lon}} (\sphinxstyleliteralemphasis{\sphinxupquote{float}}) \textendash{} The longitude coordinate for the position of interest.
The longitude position value given in decimal degrees.

\end{itemize}

\sphinxlineitem{Returns}
\sphinxAtStartPar
The gravity disturbance value (in m/s\textasciicircum{}2).

\sphinxlineitem{Return type}
\sphinxAtStartPar
float

\end{description}\end{quote}

\end{fulllineitems}


\end{fulllineitems}

\index{read\_irish\_gravity\_file() (in module qnav.gravity.irish)@\spxentry{read\_irish\_gravity\_file()}\spxextra{in module qnav.gravity.irish}}

\begin{fulllineitems}
\phantomsection\label{\detokenize{modules/gravity/irish:qnav.gravity.irish.read_irish_gravity_file}}
\pysigstartsignatures
\pysiglinewithargsret{\sphinxcode{\sphinxupquote{qnav.gravity.irish.}}\sphinxbfcode{\sphinxupquote{read\_irish\_gravity\_file}}}{\sphinxparam{\DUrole{n}{data\_file}\DUrole{p}{:}\DUrole{w}{ }\DUrole{n}{Path}}\sphinxparamcomma \sphinxparam{\DUrole{n}{null\_value}\DUrole{p}{:}\DUrole{w}{ }\DUrole{n}{float}\DUrole{w}{ }\DUrole{o}{=}\DUrole{w}{ }\DUrole{default_value}{\sphinxhyphen{}99999}}}{{ $\rightarrow$ ndarray}}
\pysigstopsignatures
\sphinxAtStartPar
Extracts the data from a given Irish Sea Gravity Map data file, using
values read from its corresponding header file. On succession, the read
gravitational anomaly matrix will be returned (measured in m/s\textasciicircum{}2). No axis
information is contained within data files.
\begin{quote}\begin{description}
\sphinxlineitem{Parameters}\begin{itemize}
\item {} 
\sphinxAtStartPar
\sphinxstyleliteralstrong{\sphinxupquote{data\_file}} (\sphinxstyleliteralemphasis{\sphinxupquote{str}}\sphinxstyleliteralemphasis{\sphinxupquote{ or }}\sphinxstyleliteralemphasis{\sphinxupquote{Path}}) \textendash{} The location of the Irish Sea Gravity Map data file.

\item {} 
\sphinxAtStartPar
\sphinxstyleliteralstrong{\sphinxupquote{null\_value}} (\sphinxstyleliteralemphasis{\sphinxupquote{float}}) \textendash{} Data values less than or equal to this value will be
converted to NaN values (default=\sphinxhyphen{}99999).

\end{itemize}

\sphinxlineitem{Returns}
\sphinxAtStartPar
The extracted gravitational anomaly matrix (measured in m/s\textasciicircum{}2).

\sphinxlineitem{Return type}
\sphinxAtStartPar
2D NumPy Array

\end{description}\end{quote}

\end{fulllineitems}

\index{read\_irish\_ers\_file() (in module qnav.gravity.irish)@\spxentry{read\_irish\_ers\_file()}\spxextra{in module qnav.gravity.irish}}

\begin{fulllineitems}
\phantomsection\label{\detokenize{modules/gravity/irish:qnav.gravity.irish.read_irish_ers_file}}
\pysigstartsignatures
\pysiglinewithargsret{\sphinxcode{\sphinxupquote{qnav.gravity.irish.}}\sphinxbfcode{\sphinxupquote{read\_irish\_ers\_file}}}{\sphinxparam{\DUrole{n}{header\_file}}}{{ $\rightarrow$ dict\DUrole{p}{{[}}str\DUrole{p}{,}\DUrole{w}{ }any\DUrole{p}{{]}}}}
\pysigstopsignatures
\sphinxAtStartPar
This function will read a given Irish Sea Gravity Map header file (a file
normally with the extension ‘.ers’) and extract the relevant information
from it.
\begin{quote}\begin{description}
\sphinxlineitem{Parameters}
\sphinxAtStartPar
\sphinxstyleliteralstrong{\sphinxupquote{header\_file}} (\sphinxstyleliteralemphasis{\sphinxupquote{str}}\sphinxstyleliteralemphasis{\sphinxupquote{ or }}\sphinxstyleliteralemphasis{\sphinxupquote{Path}}) \textendash{} The location of the header corresponding file.

\sphinxlineitem{Returns}
\sphinxAtStartPar
A dictionary holding the extracted northings and eastings
origin coordinates, data dimensions, and cell spacing information.

\end{description}\end{quote}

\end{fulllineitems}



\subsection{Input}
\label{\detokenize{usage/packages:input}}
\sphinxAtStartPar
Modules related to reading in input data:

\sphinxstepscope


\subsubsection{qnav.input.config}
\label{\detokenize{modules/input/config:module-qnav.input.config}}\label{\detokenize{modules/input/config:qnav-input-config}}\label{\detokenize{modules/input/config::doc}}\index{module@\spxentry{module}!qnav.input.config@\spxentry{qnav.input.config}}\index{qnav.input.config@\spxentry{qnav.input.config}!module@\spxentry{module}}

\paragraph{config.py}
\label{\detokenize{modules/input/config:config-py}}\begin{quote}\begin{description}
\sphinxlineitem{summary}
\sphinxAtStartPar
Obtains primary simulation components from configuration file.
This module is used for initialising and obtaining the main
components used in simulations and for logging details to
the console.

\sphinxlineitem{author}
\sphinxAtStartPar
Michael Wright \sphinxhyphen{} \sphinxhref{mailto:mjwright@liverpool.ac.uk}{mjwright@liverpool.ac.uk}

\sphinxlineitem{version}
\sphinxAtStartPar
1.0.0 \sphinxhyphen{} First full implemented version.

\end{description}\end{quote}
\index{get\_configuration() (in module qnav.input.config)@\spxentry{get\_configuration()}\spxextra{in module qnav.input.config}}

\begin{fulllineitems}
\phantomsection\label{\detokenize{modules/input/config:qnav.input.config.get_configuration}}
\pysigstartsignatures
\pysiglinewithargsret{\sphinxcode{\sphinxupquote{qnav.input.config.}}\sphinxbfcode{\sphinxupquote{get\_configuration}}}{\sphinxparam{\DUrole{n}{config\_file}\DUrole{p}{:}\DUrole{w}{ }\DUrole{n}{Path}\DUrole{w}{ }\DUrole{o}{=}\DUrole{w}{ }\DUrole{default_value}{None}}}{{ $\rightarrow$ {\hyperref[\detokenize{modules/input/config_handler:qnav.input.config_handler.ConfigHandler}]{\sphinxcrossref{ConfigHandler}}}}}
\pysigstopsignatures
\sphinxAtStartPar
Initialises and returns a \sphinxcode{\sphinxupquote{qnav.input.config.ConfigHandler}}
instance using the contents of the given file. On succession, the
file’s variables, along with those in any sub\sphinxhyphen{}configurations,
will be parsed and loaded. If not file is given, a default
settings file will be used.
\begin{quote}\begin{description}
\sphinxlineitem{Parameters}
\sphinxAtStartPar
\sphinxstyleliteralstrong{\sphinxupquote{config\_file}} (\sphinxstyleliteralemphasis{\sphinxupquote{Path}}) \textendash{} A path for a configuration (.ini) file to use.

\sphinxlineitem{Returns}
\sphinxAtStartPar
An initialised ConfigHandler instance.

\sphinxlineitem{Return type}
\sphinxAtStartPar
{\hyperref[\detokenize{modules/input/config_handler:qnav.input.config_handler.ConfigHandler}]{\sphinxcrossref{ConfigHandler}}}

\end{description}\end{quote}

\end{fulllineitems}

\index{get\_trajectory() (in module qnav.input.config)@\spxentry{get\_trajectory()}\spxextra{in module qnav.input.config}}

\begin{fulllineitems}
\phantomsection\label{\detokenize{modules/input/config:qnav.input.config.get_trajectory}}
\pysigstartsignatures
\pysiglinewithargsret{\sphinxcode{\sphinxupquote{qnav.input.config.}}\sphinxbfcode{\sphinxupquote{get\_trajectory}}}{\sphinxparam{\DUrole{n}{config}\DUrole{p}{:}\DUrole{w}{ }\DUrole{n}{{\hyperref[\detokenize{modules/input/config_handler:qnav.input.config_handler.ConfigHandler}]{\sphinxcrossref{ConfigHandler}}}}}}{{ $\rightarrow$ {\hyperref[\detokenize{modules/waypoints/trajectory:qnav.waypoints.trajectory.Trajectory}]{\sphinxcrossref{Trajectory}}}}}
\pysigstopsignatures
\sphinxAtStartPar
Extracts values and produces ground truth trajectories.
When called, this function reads the required variables from the
given configuration and uses them to generate a trajectory. On
succession, this is return ready for use in a simulation.
\begin{quote}\begin{description}
\sphinxlineitem{Parameters}
\sphinxAtStartPar
\sphinxstyleliteralstrong{\sphinxupquote{config}} ({\hyperref[\detokenize{modules/input/config_handler:qnav.input.config_handler.ConfigHandler}]{\sphinxcrossref{\sphinxstyleliteralemphasis{\sphinxupquote{ConfigHandler}}}}}) \textendash{} The ConfigHandler instance to use for simulation settings.

\sphinxlineitem{Returns}
\sphinxAtStartPar
The trajectory ready for use in a simulation.

\sphinxlineitem{Return type}
\sphinxAtStartPar
{\hyperref[\detokenize{modules/waypoints/trajectory:qnav.waypoints.trajectory.Trajectory}]{\sphinxcrossref{Trajectory}}}

\end{description}\end{quote}

\end{fulllineitems}

\index{get\_estimation() (in module qnav.input.config)@\spxentry{get\_estimation()}\spxextra{in module qnav.input.config}}

\begin{fulllineitems}
\phantomsection\label{\detokenize{modules/input/config:qnav.input.config.get_estimation}}
\pysigstartsignatures
\pysiglinewithargsret{\sphinxcode{\sphinxupquote{qnav.input.config.}}\sphinxbfcode{\sphinxupquote{get\_estimation}}}{\sphinxparam{\DUrole{n}{config}\DUrole{p}{:}\DUrole{w}{ }\DUrole{n}{{\hyperref[\detokenize{modules/input/config_handler:qnav.input.config_handler.ConfigHandler}]{\sphinxcrossref{ConfigHandler}}}}}\sphinxparamcomma \sphinxparam{\DUrole{n}{trajectory}\DUrole{p}{:}\DUrole{w}{ }\DUrole{n}{{\hyperref[\detokenize{modules/waypoints/trajectory:qnav.waypoints.trajectory.Trajectory}]{\sphinxcrossref{Trajectory}}}}}}{{ $\rightarrow$ {\hyperref[\detokenize{modules/estimation/state:qnav.estimation.state.EstimatedState}]{\sphinxcrossref{EstimatedState}}}}}
\pysigstopsignatures
\sphinxAtStartPar
Extracts values and produces the initial estimation state.
When called, this function reads the required variables from the
given configuration and uses them to initialise the starting
estimation information. This uses the first record from the
previously generated trajectory and induces any requested errors.
On succession, an initial estimated state is return.
\begin{quote}\begin{description}
\sphinxlineitem{Parameters}\begin{itemize}
\item {} 
\sphinxAtStartPar
\sphinxstyleliteralstrong{\sphinxupquote{config}} ({\hyperref[\detokenize{modules/input/config_handler:qnav.input.config_handler.ConfigHandler}]{\sphinxcrossref{\sphinxstyleliteralemphasis{\sphinxupquote{ConfigHandler}}}}}) \textendash{} The ConfigHandler instance to use for simulation settings.

\item {} 
\sphinxAtStartPar
\sphinxstyleliteralstrong{\sphinxupquote{trajectory}} ({\hyperref[\detokenize{modules/waypoints/trajectory:qnav.waypoints.trajectory.Trajectory}]{\sphinxcrossref{\sphinxstyleliteralemphasis{\sphinxupquote{Trajectory}}}}}) \textendash{} The trajectory to be used during the simulation.
Required to obtain the true state at the first time step.

\end{itemize}

\sphinxlineitem{Returns}
\sphinxAtStartPar
The initial estimated state to use in the simulation.

\sphinxlineitem{Return type}
\sphinxAtStartPar
{\hyperref[\detokenize{modules/estimation/state:qnav.estimation.state.EstimatedState}]{\sphinxcrossref{EstimatedState}}}

\end{description}\end{quote}

\end{fulllineitems}

\index{get\_platform\_hardware() (in module qnav.input.config)@\spxentry{get\_platform\_hardware()}\spxextra{in module qnav.input.config}}

\begin{fulllineitems}
\phantomsection\label{\detokenize{modules/input/config:qnav.input.config.get_platform_hardware}}
\pysigstartsignatures
\pysiglinewithargsret{\sphinxcode{\sphinxupquote{qnav.input.config.}}\sphinxbfcode{\sphinxupquote{get\_platform\_hardware}}}{\sphinxparam{\DUrole{n}{config}\DUrole{p}{:}\DUrole{w}{ }\DUrole{n}{{\hyperref[\detokenize{modules/input/config_handler:qnav.input.config_handler.ConfigHandler}]{\sphinxcrossref{ConfigHandler}}}}}\sphinxparamcomma \sphinxparam{\DUrole{n}{estimated\_state}\DUrole{p}{:}\DUrole{w}{ }\DUrole{n}{{\hyperref[\detokenize{modules/estimation/state:qnav.estimation.state.EstimatedState}]{\sphinxcrossref{EstimatedState}}}}}}{{ $\rightarrow$ list\DUrole{p}{{[}}{\hyperref[\detokenize{modules/measurement/sensor:qnav.measurement.sensor.SensorFusion}]{\sphinxcrossref{SensorFusion}}}\DUrole{p}{{]}}}}
\pysigstopsignatures
\sphinxAtStartPar
Extracts values and produces the virtual platform hardware instances.
When called, this function reads the required variables from the
given configuration and uses them initialize and return the requested
Sensor\sphinxhyphen{}Fusion instance to use during the simulation. On succession,
a list of SensorFusion classes to use in a simulation are returned.
\begin{quote}\begin{description}
\sphinxlineitem{Parameters}\begin{itemize}
\item {} 
\sphinxAtStartPar
\sphinxstyleliteralstrong{\sphinxupquote{config}} ({\hyperref[\detokenize{modules/input/config_handler:qnav.input.config_handler.ConfigHandler}]{\sphinxcrossref{\sphinxstyleliteralemphasis{\sphinxupquote{ConfigHandler}}}}}) \textendash{} The ConfigHandler instance to use for simulation settings.

\item {} 
\sphinxAtStartPar
\sphinxstyleliteralstrong{\sphinxupquote{estimated\_state}} ({\hyperref[\detokenize{modules/estimation/state:qnav.estimation.state.EstimatedState}]{\sphinxcrossref{\sphinxstyleliteralemphasis{\sphinxupquote{EstimatedState}}}}}) \textendash{} The initial estimates before start of simulation.
Required as some fusion methods need an initial estimate to start with.

\end{itemize}

\sphinxlineitem{Returns}
\sphinxAtStartPar
The list of SensorFusion instances to use in a simulation.

\sphinxlineitem{Return type}
\sphinxAtStartPar
list{[}{\hyperref[\detokenize{modules/measurement/sensor:qnav.measurement.sensor.SensorFusion}]{\sphinxcrossref{SensorFusion}}}{]}

\end{description}\end{quote}

\end{fulllineitems}

\index{get\_platform\_clock() (in module qnav.input.config)@\spxentry{get\_platform\_clock()}\spxextra{in module qnav.input.config}}

\begin{fulllineitems}
\phantomsection\label{\detokenize{modules/input/config:qnav.input.config.get_platform_clock}}
\pysigstartsignatures
\pysiglinewithargsret{\sphinxcode{\sphinxupquote{qnav.input.config.}}\sphinxbfcode{\sphinxupquote{get\_platform\_clock}}}{\sphinxparam{\DUrole{n}{config}\DUrole{p}{:}\DUrole{w}{ }\DUrole{n}{{\hyperref[\detokenize{modules/input/config_handler:qnav.input.config_handler.ConfigHandler}]{\sphinxcrossref{ConfigHandler}}}}}}{{ $\rightarrow$ {\hyperref[\detokenize{modules/measurement/clock:qnav.measurement.clock.Clock}]{\sphinxcrossref{Clock}}}}}
\pysigstopsignatures
\sphinxAtStartPar
Extracts values and produces a virtual platform clock instances.
When called, this function reads the required variables from the
given configuration and uses them to initialise and return the
platform clock to use.
\begin{quote}\begin{description}
\sphinxlineitem{Parameters}
\sphinxAtStartPar
\sphinxstyleliteralstrong{\sphinxupquote{config}} ({\hyperref[\detokenize{modules/input/config_handler:qnav.input.config_handler.ConfigHandler}]{\sphinxcrossref{\sphinxstyleliteralemphasis{\sphinxupquote{ConfigHandler}}}}}) \textendash{} The ConfigHandler instance to use for simulation settings.

\sphinxlineitem{Returns}
\sphinxAtStartPar
The Clock instance to use in a simulation.

\sphinxlineitem{Return type}
\sphinxAtStartPar
{\hyperref[\detokenize{modules/measurement/clock:qnav.measurement.clock.Clock}]{\sphinxcrossref{Clock}}}

\end{description}\end{quote}

\end{fulllineitems}

\index{get\_collector() (in module qnav.input.config)@\spxentry{get\_collector()}\spxextra{in module qnav.input.config}}

\begin{fulllineitems}
\phantomsection\label{\detokenize{modules/input/config:qnav.input.config.get_collector}}
\pysigstartsignatures
\pysiglinewithargsret{\sphinxcode{\sphinxupquote{qnav.input.config.}}\sphinxbfcode{\sphinxupquote{get\_collector}}}{\sphinxparam{\DUrole{n}{config}\DUrole{p}{:}\DUrole{w}{ }\DUrole{n}{{\hyperref[\detokenize{modules/input/config_handler:qnav.input.config_handler.ConfigHandler}]{\sphinxcrossref{ConfigHandler}}}}}}{{ $\rightarrow$ ResultsCapture}}
\pysigstopsignatures
\sphinxAtStartPar
Extracts values and produces the results collection instance to use.
When called, this function reads the required variables from the
given configuration and uses them initialise and return a results
collection instance, responsible for filtering and collecting
simulation results.
\begin{quote}\begin{description}
\sphinxlineitem{Parameters}
\sphinxAtStartPar
\sphinxstyleliteralstrong{\sphinxupquote{config}} ({\hyperref[\detokenize{modules/input/config_handler:qnav.input.config_handler.ConfigHandler}]{\sphinxcrossref{\sphinxstyleliteralemphasis{\sphinxupquote{ConfigHandler}}}}}) \textendash{} The ConfigHandler instance to use for simulation settings.

\sphinxlineitem{Returns}
\sphinxAtStartPar
The ResultsCapture instance to use in a simulation.

\sphinxlineitem{Return type}
\sphinxAtStartPar
ResultsCapture

\end{description}\end{quote}

\end{fulllineitems}

\index{get\_displayer() (in module qnav.input.config)@\spxentry{get\_displayer()}\spxextra{in module qnav.input.config}}

\begin{fulllineitems}
\phantomsection\label{\detokenize{modules/input/config:qnav.input.config.get_displayer}}
\pysigstartsignatures
\pysiglinewithargsret{\sphinxcode{\sphinxupquote{qnav.input.config.}}\sphinxbfcode{\sphinxupquote{get\_displayer}}}{\sphinxparam{\DUrole{n}{config}\DUrole{p}{:}\DUrole{w}{ }\DUrole{n}{{\hyperref[\detokenize{modules/input/config_handler:qnav.input.config_handler.ConfigHandler}]{\sphinxcrossref{ConfigHandler}}}}}}{{ $\rightarrow$ {\hyperref[\detokenize{modules/simulation/display:qnav.simulation.display.Display}]{\sphinxcrossref{Display}}}}}
\pysigstopsignatures
\sphinxAtStartPar
Extracts values and produces the displaying instance to use.
When called, this function reads the required variables from the
given configuration and uses them initialise and return a display
instance, responsible for feeding back real\sphinxhyphen{}time information about
the simulation. Can also be used for stopping the simulation
prematurely when criteria is met.
\begin{quote}\begin{description}
\sphinxlineitem{Parameters}
\sphinxAtStartPar
\sphinxstyleliteralstrong{\sphinxupquote{config}} ({\hyperref[\detokenize{modules/input/config_handler:qnav.input.config_handler.ConfigHandler}]{\sphinxcrossref{\sphinxstyleliteralemphasis{\sphinxupquote{ConfigHandler}}}}}) \textendash{} The ConfigHandler instance to use for simulation settings.

\sphinxlineitem{Returns}
\sphinxAtStartPar
The Display instance to use in a simulation.

\sphinxlineitem{Return type}
\sphinxAtStartPar
{\hyperref[\detokenize{modules/simulation/display:qnav.simulation.display.Display}]{\sphinxcrossref{Display}}}

\end{description}\end{quote}

\end{fulllineitems}

\index{get\_results\_settings() (in module qnav.input.config)@\spxentry{get\_results\_settings()}\spxextra{in module qnav.input.config}}

\begin{fulllineitems}
\phantomsection\label{\detokenize{modules/input/config:qnav.input.config.get_results_settings}}
\pysigstartsignatures
\pysiglinewithargsret{\sphinxcode{\sphinxupquote{qnav.input.config.}}\sphinxbfcode{\sphinxupquote{get\_results\_settings}}}{\sphinxparam{\DUrole{n}{config}\DUrole{p}{:}\DUrole{w}{ }\DUrole{n}{{\hyperref[\detokenize{modules/input/config_handler:qnav.input.config_handler.ConfigHandler}]{\sphinxcrossref{ConfigHandler}}}}}}{{ $\rightarrow$ dict}}
\pysigstopsignatures
\sphinxAtStartPar
Extracts values and produces the results generation settings.
When called, this function reads the required variables from the
given configuration and uses them to define how results will be
gathered after the simulation and exported.
\begin{quote}\begin{description}
\sphinxlineitem{Parameters}
\sphinxAtStartPar
\sphinxstyleliteralstrong{\sphinxupquote{config}} ({\hyperref[\detokenize{modules/input/config_handler:qnav.input.config_handler.ConfigHandler}]{\sphinxcrossref{\sphinxstyleliteralemphasis{\sphinxupquote{ConfigHandler}}}}}) \textendash{} The ConfigHandler instance to use for simulation settings.

\sphinxlineitem{Returns}
\sphinxAtStartPar
A dictionary containing settings for exporting results.

\sphinxlineitem{Return type}
\sphinxAtStartPar
dict

\end{description}\end{quote}

\end{fulllineitems}

\index{get\_figure\_settings() (in module qnav.input.config)@\spxentry{get\_figure\_settings()}\spxextra{in module qnav.input.config}}

\begin{fulllineitems}
\phantomsection\label{\detokenize{modules/input/config:qnav.input.config.get_figure_settings}}
\pysigstartsignatures
\pysiglinewithargsret{\sphinxcode{\sphinxupquote{qnav.input.config.}}\sphinxbfcode{\sphinxupquote{get\_figure\_settings}}}{\sphinxparam{\DUrole{n}{config}\DUrole{p}{:}\DUrole{w}{ }\DUrole{n}{{\hyperref[\detokenize{modules/input/config_handler:qnav.input.config_handler.ConfigHandler}]{\sphinxcrossref{ConfigHandler}}}}}}{{ $\rightarrow$ dict}}
\pysigstopsignatures
\sphinxAtStartPar
Extracts values to use for results figure and report generation.
When called, this function reads the required variables from the
given configuration and returns them in a dictionary. These can
then be used later for generating the relevant after\sphinxhyphen{}simulation
graphics and report results.
\begin{quote}\begin{description}
\sphinxlineitem{Parameters}
\sphinxAtStartPar
\sphinxstyleliteralstrong{\sphinxupquote{config}} ({\hyperref[\detokenize{modules/input/config_handler:qnav.input.config_handler.ConfigHandler}]{\sphinxcrossref{\sphinxstyleliteralemphasis{\sphinxupquote{ConfigHandler}}}}}) \textendash{} The ConfigHandler instance to use for simulation settings.

\sphinxlineitem{Returns}
\sphinxAtStartPar
Dictionary containing figure generation settings.

\sphinxlineitem{Return type}
\sphinxAtStartPar
dict

\end{description}\end{quote}

\end{fulllineitems}

\index{print\_banner() (in module qnav.input.config)@\spxentry{print\_banner()}\spxextra{in module qnav.input.config}}

\begin{fulllineitems}
\phantomsection\label{\detokenize{modules/input/config:qnav.input.config.print_banner}}
\pysigstartsignatures
\pysiglinewithargsret{\sphinxcode{\sphinxupquote{qnav.input.config.}}\sphinxbfcode{\sphinxupquote{print\_banner}}}{\sphinxparam{\DUrole{n}{text}\DUrole{p}{:}\DUrole{w}{ }\DUrole{n}{str}\DUrole{w}{ }\DUrole{o}{=}\DUrole{w}{ }\DUrole{default_value}{None}}}{}
\pysigstopsignatures
\sphinxAtStartPar
Used to print formatted heading in the output console.
Simply used to group list points together cleanly.
\begin{quote}\begin{description}
\sphinxlineitem{Parameters}
\sphinxAtStartPar
\sphinxstyleliteralstrong{\sphinxupquote{text}} (\sphinxstyleliteralemphasis{\sphinxupquote{str}}) \textendash{} (Optional) The text to include in the heading.
If not provided a dividing line will be used.

\end{description}\end{quote}

\end{fulllineitems}

\index{print\_line() (in module qnav.input.config)@\spxentry{print\_line()}\spxextra{in module qnav.input.config}}

\begin{fulllineitems}
\phantomsection\label{\detokenize{modules/input/config:qnav.input.config.print_line}}
\pysigstartsignatures
\pysiglinewithargsret{\sphinxcode{\sphinxupquote{qnav.input.config.}}\sphinxbfcode{\sphinxupquote{print\_line}}}{\sphinxparam{\DUrole{n}{field}\DUrole{p}{:}\DUrole{w}{ }\DUrole{n}{str}}\sphinxparamcomma \sphinxparam{\DUrole{n}{value}\DUrole{p}{:}\DUrole{w}{ }\DUrole{n}{any}}}{}
\pysigstopsignatures
\sphinxAtStartPar
Used to print formated lines in the output console.
Simply used to produce clean uniformly spaced list points.
\begin{quote}\begin{description}
\sphinxlineitem{Parameters}\begin{itemize}
\item {} 
\sphinxAtStartPar
\sphinxstyleliteralstrong{\sphinxupquote{field}} (\sphinxstyleliteralemphasis{\sphinxupquote{str}}) \textendash{} The text to place on the left\sphinxhyphen{}hand side.

\item {} 
\sphinxAtStartPar
\sphinxstyleliteralstrong{\sphinxupquote{value}} (\sphinxstyleliteralemphasis{\sphinxupquote{any}}) \textendash{} The value to place on the right\sphinxhyphen{}hand side.

\end{itemize}

\end{description}\end{quote}

\end{fulllineitems}


\sphinxstepscope


\subsubsection{qnav.input.config\_handler}
\label{\detokenize{modules/input/config_handler:module-qnav.input.config_handler}}\label{\detokenize{modules/input/config_handler:qnav-input-config-handler}}\label{\detokenize{modules/input/config_handler::doc}}\index{module@\spxentry{module}!qnav.input.config\_handler@\spxentry{qnav.input.config\_handler}}\index{qnav.input.config\_handler@\spxentry{qnav.input.config\_handler}!module@\spxentry{module}}

\paragraph{config\_handler.py}
\label{\detokenize{modules/input/config_handler:config-handler-py}}\begin{quote}\begin{description}
\sphinxlineitem{summary}
\sphinxAtStartPar
This class provides a wrapper for Python’s ConfigHandler class.
This supplies additional functionality such as automatically reading from
sub\sphinxhyphen{}configuration files and providing platform independent means for
extracting configuration values in usable formats.

\sphinxlineitem{authors}
\sphinxAtStartPar
Michael Wright \sphinxhyphen{} \sphinxhref{mailto:mjwright@liverpool.ac.uk}{mjwright@liverpool.ac.uk}

\sphinxlineitem{version}
\sphinxAtStartPar
1.0.0 \sphinxhyphen{} First implementation of class.

\end{description}\end{quote}
\index{ConfigHandler (class in qnav.input.config\_handler)@\spxentry{ConfigHandler}\spxextra{class in qnav.input.config\_handler}}

\begin{fulllineitems}
\phantomsection\label{\detokenize{modules/input/config_handler:qnav.input.config_handler.ConfigHandler}}
\pysigstartsignatures
\pysiglinewithargsret{\sphinxbfcode{\sphinxupquote{class\DUrole{w}{ }}}\sphinxcode{\sphinxupquote{qnav.input.config\_handler.}}\sphinxbfcode{\sphinxupquote{ConfigHandler}}}{\sphinxparam{\DUrole{n}{config\_file}\DUrole{p}{:}\DUrole{w}{ }\DUrole{n}{Path}\DUrole{w}{ }\DUrole{o}{=}\DUrole{w}{ }\DUrole{default_value}{PosixPath(\textquotesingle{}default\_settings.ini\textquotesingle{})}}}{}
\pysigstopsignatures
\sphinxAtStartPar
This class provides means for reading the user’s configuration preferences
from a given configuration file. Once read, values can be extracted within
various formats (strings, booleans, floats, arrays, etc..). This class is
essentially a wrapper for Python’s ConfigHandler class, providing
additional functionality such as the ability to split configuration
across multiple files.
\index{get\_config\_name() (qnav.input.config\_handler.ConfigHandler method)@\spxentry{get\_config\_name()}\spxextra{qnav.input.config\_handler.ConfigHandler method}}

\begin{fulllineitems}
\phantomsection\label{\detokenize{modules/input/config_handler:qnav.input.config_handler.ConfigHandler.get_config_name}}
\pysigstartsignatures
\pysiglinewithargsret{\sphinxbfcode{\sphinxupquote{get\_config\_name}}}{}{{ $\rightarrow$ str}}
\pysigstopsignatures
\sphinxAtStartPar
This simple function returns the name of the configuration file being
used. This function will automatically remove the project database\_dir
from the returned name (if the selected config is located within the
project database\_dir).
\begin{quote}\begin{description}
\sphinxlineitem{Returns}
\sphinxAtStartPar
The short name of the selected configuration file that is
being used.

\sphinxlineitem{Return type}
\sphinxAtStartPar
str

\end{description}\end{quote}

\end{fulllineitems}

\index{load\_configuration() (qnav.input.config\_handler.ConfigHandler method)@\spxentry{load\_configuration()}\spxextra{qnav.input.config\_handler.ConfigHandler method}}

\begin{fulllineitems}
\phantomsection\label{\detokenize{modules/input/config_handler:qnav.input.config_handler.ConfigHandler.load_configuration}}
\pysigstartsignatures
\pysiglinewithargsret{\sphinxbfcode{\sphinxupquote{load\_configuration}}}{}{}
\pysigstopsignatures
\sphinxAtStartPar
This function reads and records the content of the given configuration
file (and expected sub\sphinxhyphen{}files). Expected sections/sensors are
automatically searched and their sub\sphinxhyphen{}files are also included within
the loading process. This function can be used to re\sphinxhyphen{}load/re\sphinxhyphen{}read the
information from configuration files (if changes are made during
runtime).

\end{fulllineitems}

\index{append\_config() (qnav.input.config\_handler.ConfigHandler method)@\spxentry{append\_config()}\spxextra{qnav.input.config\_handler.ConfigHandler method}}

\begin{fulllineitems}
\phantomsection\label{\detokenize{modules/input/config_handler:qnav.input.config_handler.ConfigHandler.append_config}}
\pysigstartsignatures
\pysiglinewithargsret{\sphinxbfcode{\sphinxupquote{append\_config}}}{\sphinxparam{\DUrole{n}{config\_file}\DUrole{p}{:}\DUrole{w}{ }\DUrole{n}{str}}\sphinxparamcomma \sphinxparam{\DUrole{n}{from\_section}\DUrole{p}{:}\DUrole{w}{ }\DUrole{n}{str}}\sphinxparamcomma \sphinxparam{\DUrole{n}{to\_section}\DUrole{p}{:}\DUrole{w}{ }\DUrole{n}{str}}\sphinxparamcomma \sphinxparam{\DUrole{n}{to\_overwrite}\DUrole{p}{:}\DUrole{w}{ }\DUrole{n}{bool}\DUrole{w}{ }\DUrole{o}{=}\DUrole{w}{ }\DUrole{default_value}{False}}}{}
\pysigstopsignatures
\sphinxAtStartPar
This function will read and record the values of an external
configuration files (i.e. a sub\sphinxhyphen{}configuration) and merge the read
information with the current ConfigHandler instance.
\begin{quote}\begin{description}
\sphinxlineitem{Parameters}\begin{itemize}
\item {} 
\sphinxAtStartPar
\sphinxstyleliteralstrong{\sphinxupquote{config\_file}} (\sphinxstyleliteralemphasis{\sphinxupquote{str}}) \textendash{} The sub\sphinxhyphen{}configuration file to read.

\item {} 
\sphinxAtStartPar
\sphinxstyleliteralstrong{\sphinxupquote{from\_section}} (\sphinxstyleliteralemphasis{\sphinxupquote{str}}) \textendash{} The section within the sub\sphinxhyphen{}configuration
file to read and record.

\item {} 
\sphinxAtStartPar
\sphinxstyleliteralstrong{\sphinxupquote{to\_section}} (\sphinxstyleliteralemphasis{\sphinxupquote{str}}) \textendash{} The section within the current instance to append
the read information to (can be a new section).

\item {} 
\sphinxAtStartPar
\sphinxstyleliteralstrong{\sphinxupquote{to\_overwrite}} (\sphinxstyleliteralemphasis{\sphinxupquote{bool}}) \textendash{} If set to true, previous existing values from
the main configuration will be overwritten by those read from
sub\sphinxhyphen{}configuration files (default=False).

\end{itemize}

\end{description}\end{quote}

\end{fulllineitems}

\index{expand\_section() (qnav.input.config\_handler.ConfigHandler method)@\spxentry{expand\_section()}\spxextra{qnav.input.config\_handler.ConfigHandler method}}

\begin{fulllineitems}
\phantomsection\label{\detokenize{modules/input/config_handler:qnav.input.config_handler.ConfigHandler.expand_section}}
\pysigstartsignatures
\pysiglinewithargsret{\sphinxbfcode{\sphinxupquote{expand\_section}}}{\sphinxparam{\DUrole{n}{section}\DUrole{p}{:}\DUrole{w}{ }\DUrole{n}{str}}\sphinxparamcomma \sphinxparam{\DUrole{n}{path\_property}\DUrole{p}{:}\DUrole{w}{ }\DUrole{n}{str}}\sphinxparamcomma \sphinxparam{\DUrole{n}{sub\_file\_section}\DUrole{p}{:}\DUrole{w}{ }\DUrole{n}{str}\DUrole{w}{ }\DUrole{o}{=}\DUrole{w}{ }\DUrole{default_value}{None}}\sphinxparamcomma \sphinxparam{\DUrole{n}{to\_overwrite}\DUrole{p}{:}\DUrole{w}{ }\DUrole{n}{bool}\DUrole{w}{ }\DUrole{o}{=}\DUrole{w}{ }\DUrole{default_value}{False}}}{}
\pysigstopsignatures
\sphinxAtStartPar
Expand section by importing the referenced sub\sphinxhyphen{}configuration path.
Imports additional properties into section, for a given section and
property referencing the sub\sphinxhyphen{}configuration file.
\begin{quote}\begin{description}
\sphinxlineitem{Parameters}\begin{itemize}
\item {} 
\sphinxAtStartPar
\sphinxstyleliteralstrong{\sphinxupquote{section}} (\sphinxstyleliteralemphasis{\sphinxupquote{str}}) \textendash{} The section containing the property with the
sub\sphinxhyphen{}configuration path.

\item {} 
\sphinxAtStartPar
\sphinxstyleliteralstrong{\sphinxupquote{path\_property}} (\sphinxstyleliteralemphasis{\sphinxupquote{str}}) \textendash{} The property containing the sub\sphinxhyphen{}
configuration path.

\item {} 
\sphinxAtStartPar
\sphinxstyleliteralstrong{\sphinxupquote{sub\_file\_section}} (\sphinxstyleliteralemphasis{\sphinxupquote{str}}) \textendash{} (Optional) The section to import from the
sub\sphinxhyphen{}configuration file.

\item {} 
\sphinxAtStartPar
\sphinxstyleliteralstrong{\sphinxupquote{to\_overwrite}} (\sphinxstyleliteralemphasis{\sphinxupquote{bool}}) \textendash{} Should this import overwrite existing settings?

\end{itemize}

\end{description}\end{quote}

\end{fulllineitems}

\index{quick\_append\_config() (qnav.input.config\_handler.ConfigHandler method)@\spxentry{quick\_append\_config()}\spxextra{qnav.input.config\_handler.ConfigHandler method}}

\begin{fulllineitems}
\phantomsection\label{\detokenize{modules/input/config_handler:qnav.input.config_handler.ConfigHandler.quick_append_config}}
\pysigstartsignatures
\pysiglinewithargsret{\sphinxbfcode{\sphinxupquote{quick\_append\_config}}}{\sphinxparam{\DUrole{n}{from\_section}\DUrole{p}{:}\DUrole{w}{ }\DUrole{n}{Buffer\DUrole{w}{ }\DUrole{p}{|}\DUrole{w}{ }\_SupportsArray\DUrole{p}{{[}}dtype\DUrole{p}{{[}}Any\DUrole{p}{{]}}\DUrole{p}{{]}}\DUrole{w}{ }\DUrole{p}{|}\DUrole{w}{ }\_NestedSequence\DUrole{p}{{[}}\_SupportsArray\DUrole{p}{{[}}dtype\DUrole{p}{{[}}Any\DUrole{p}{{]}}\DUrole{p}{{]}}\DUrole{p}{{]}}\DUrole{w}{ }\DUrole{p}{|}\DUrole{w}{ }bool\DUrole{w}{ }\DUrole{p}{|}\DUrole{w}{ }int\DUrole{w}{ }\DUrole{p}{|}\DUrole{w}{ }float\DUrole{w}{ }\DUrole{p}{|}\DUrole{w}{ }complex\DUrole{w}{ }\DUrole{p}{|}\DUrole{w}{ }str\DUrole{w}{ }\DUrole{p}{|}\DUrole{w}{ }bytes\DUrole{w}{ }\DUrole{p}{|}\DUrole{w}{ }\_NestedSequence\DUrole{p}{{[}}bool\DUrole{w}{ }\DUrole{p}{|}\DUrole{w}{ }int\DUrole{w}{ }\DUrole{p}{|}\DUrole{w}{ }float\DUrole{w}{ }\DUrole{p}{|}\DUrole{w}{ }complex\DUrole{w}{ }\DUrole{p}{|}\DUrole{w}{ }str\DUrole{w}{ }\DUrole{p}{|}\DUrole{w}{ }bytes\DUrole{p}{{]}}}}\sphinxparamcomma \sphinxparam{\DUrole{n}{from\_variable}\DUrole{p}{:}\DUrole{w}{ }\DUrole{n}{Buffer\DUrole{w}{ }\DUrole{p}{|}\DUrole{w}{ }\_SupportsArray\DUrole{p}{{[}}dtype\DUrole{p}{{[}}Any\DUrole{p}{{]}}\DUrole{p}{{]}}\DUrole{w}{ }\DUrole{p}{|}\DUrole{w}{ }\_NestedSequence\DUrole{p}{{[}}\_SupportsArray\DUrole{p}{{[}}dtype\DUrole{p}{{[}}Any\DUrole{p}{{]}}\DUrole{p}{{]}}\DUrole{p}{{]}}\DUrole{w}{ }\DUrole{p}{|}\DUrole{w}{ }bool\DUrole{w}{ }\DUrole{p}{|}\DUrole{w}{ }int\DUrole{w}{ }\DUrole{p}{|}\DUrole{w}{ }float\DUrole{w}{ }\DUrole{p}{|}\DUrole{w}{ }complex\DUrole{w}{ }\DUrole{p}{|}\DUrole{w}{ }str\DUrole{w}{ }\DUrole{p}{|}\DUrole{w}{ }bytes\DUrole{w}{ }\DUrole{p}{|}\DUrole{w}{ }\_NestedSequence\DUrole{p}{{[}}bool\DUrole{w}{ }\DUrole{p}{|}\DUrole{w}{ }int\DUrole{w}{ }\DUrole{p}{|}\DUrole{w}{ }float\DUrole{w}{ }\DUrole{p}{|}\DUrole{w}{ }complex\DUrole{w}{ }\DUrole{p}{|}\DUrole{w}{ }str\DUrole{w}{ }\DUrole{p}{|}\DUrole{w}{ }bytes\DUrole{p}{{]}}}}\sphinxparamcomma \sphinxparam{\DUrole{n}{external\_section}\DUrole{p}{:}\DUrole{w}{ }\DUrole{n}{Buffer\DUrole{w}{ }\DUrole{p}{|}\DUrole{w}{ }\_SupportsArray\DUrole{p}{{[}}dtype\DUrole{p}{{[}}Any\DUrole{p}{{]}}\DUrole{p}{{]}}\DUrole{w}{ }\DUrole{p}{|}\DUrole{w}{ }\_NestedSequence\DUrole{p}{{[}}\_SupportsArray\DUrole{p}{{[}}dtype\DUrole{p}{{[}}Any\DUrole{p}{{]}}\DUrole{p}{{]}}\DUrole{p}{{]}}\DUrole{w}{ }\DUrole{p}{|}\DUrole{w}{ }bool\DUrole{w}{ }\DUrole{p}{|}\DUrole{w}{ }int\DUrole{w}{ }\DUrole{p}{|}\DUrole{w}{ }float\DUrole{w}{ }\DUrole{p}{|}\DUrole{w}{ }complex\DUrole{w}{ }\DUrole{p}{|}\DUrole{w}{ }str\DUrole{w}{ }\DUrole{p}{|}\DUrole{w}{ }bytes\DUrole{w}{ }\DUrole{p}{|}\DUrole{w}{ }\_NestedSequence\DUrole{p}{{[}}bool\DUrole{w}{ }\DUrole{p}{|}\DUrole{w}{ }int\DUrole{w}{ }\DUrole{p}{|}\DUrole{w}{ }float\DUrole{w}{ }\DUrole{p}{|}\DUrole{w}{ }complex\DUrole{w}{ }\DUrole{p}{|}\DUrole{w}{ }str\DUrole{w}{ }\DUrole{p}{|}\DUrole{w}{ }bytes\DUrole{p}{{]}}}}\sphinxparamcomma \sphinxparam{\DUrole{n}{to\_overwrite}\DUrole{p}{:}\DUrole{w}{ }\DUrole{n}{bool}\DUrole{w}{ }\DUrole{o}{=}\DUrole{w}{ }\DUrole{default_value}{False}}}{}
\pysigstopsignatures
\sphinxAtStartPar
This function simplifies merging multiple configuration file sections
to the current configuration.
\begin{quote}\begin{description}
\sphinxlineitem{Parameters}\begin{itemize}
\item {} 
\sphinxAtStartPar
\sphinxstyleliteralstrong{\sphinxupquote{from\_section}} (\sphinxstyleliteralemphasis{\sphinxupquote{numpy.typing.ArrayLike}}) \textendash{} The sections within the main configuration file
that contain the variables to read from.

\item {} 
\sphinxAtStartPar
\sphinxstyleliteralstrong{\sphinxupquote{from\_variable}} (\sphinxstyleliteralemphasis{\sphinxupquote{numpy.typing.ArrayLike}}) \textendash{} The filepath variable names from the given
section stating the external configuration files to read from.

\item {} 
\sphinxAtStartPar
\sphinxstyleliteralstrong{\sphinxupquote{external\_section}} (\sphinxstyleliteralemphasis{\sphinxupquote{numpy.typing.ArrayLike}}) \textendash{} The external sub\sphinxhyphen{}configuration file sections
to merge with the main configuration.

\item {} 
\sphinxAtStartPar
\sphinxstyleliteralstrong{\sphinxupquote{to\_overwrite}} (\sphinxstyleliteralemphasis{\sphinxupquote{bool}}) \textendash{} If set to true, previous existing values from
the main configuration will be overwritten by those read from
sub\sphinxhyphen{}configuration files (default=False).

\end{itemize}

\end{description}\end{quote}

\end{fulllineitems}

\index{get\_all\_values() (qnav.input.config\_handler.ConfigHandler method)@\spxentry{get\_all\_values()}\spxextra{qnav.input.config\_handler.ConfigHandler method}}

\begin{fulllineitems}
\phantomsection\label{\detokenize{modules/input/config_handler:qnav.input.config_handler.ConfigHandler.get_all_values}}
\pysigstartsignatures
\pysiglinewithargsret{\sphinxbfcode{\sphinxupquote{get\_all\_values}}}{}{{ $\rightarrow$ ndarray}}
\pysigstopsignatures
\sphinxAtStartPar
This function returns a matrix containing all the values read from the
given configuration file(s). This function is namely for debugging
purposes. Values are displayed as raw strings.
\begin{quote}\begin{description}
\sphinxlineitem{Returns}
\sphinxAtStartPar
An N\sphinxhyphen{}by\sphinxhyphen{}3 matrix of values containing the section, variable
name and value for each option read.

\sphinxlineitem{Return type}
\sphinxAtStartPar
numpy.ndarray

\end{description}\end{quote}

\end{fulllineitems}

\index{get\_sections() (qnav.input.config\_handler.ConfigHandler method)@\spxentry{get\_sections()}\spxextra{qnav.input.config\_handler.ConfigHandler method}}

\begin{fulllineitems}
\phantomsection\label{\detokenize{modules/input/config_handler:qnav.input.config_handler.ConfigHandler.get_sections}}
\pysigstartsignatures
\pysiglinewithargsret{\sphinxbfcode{\sphinxupquote{get\_sections}}}{}{}
\pysigstopsignatures
\sphinxAtStartPar
This function returns a list of all the sections read from the given
configuration file(s). This is useful for debugging and ensure
required content is being correctly read.
\begin{quote}\begin{description}
\sphinxlineitem{Returns}
\sphinxAtStartPar
A list containing all section names.

\sphinxlineitem{Return type}
\sphinxAtStartPar
list

\end{description}\end{quote}

\end{fulllineitems}

\index{get\_str() (qnav.input.config\_handler.ConfigHandler method)@\spxentry{get\_str()}\spxextra{qnav.input.config\_handler.ConfigHandler method}}

\begin{fulllineitems}
\phantomsection\label{\detokenize{modules/input/config_handler:qnav.input.config_handler.ConfigHandler.get_str}}
\pysigstartsignatures
\pysiglinewithargsret{\sphinxbfcode{\sphinxupquote{get\_str}}}{\sphinxparam{\DUrole{n}{section}\DUrole{p}{:}\DUrole{w}{ }\DUrole{n}{str}}\sphinxparamcomma \sphinxparam{\DUrole{n}{name}\DUrole{p}{:}\DUrole{w}{ }\DUrole{n}{str}}\sphinxparamcomma \sphinxparam{\DUrole{n}{fallback}\DUrole{p}{:}\DUrole{w}{ }\DUrole{n}{str}\DUrole{w}{ }\DUrole{o}{=}\DUrole{w}{ }\DUrole{default_value}{None}}}{{ $\rightarrow$ str}}
\pysigstopsignatures
\sphinxAtStartPar
This function will return the string value for a parameter in the
selected configuration file, identified by its section and name. If
the value is missing, an optional default value can be used instead.
\begin{quote}\begin{description}
\sphinxlineitem{Parameters}\begin{itemize}
\item {} 
\sphinxAtStartPar
\sphinxstyleliteralstrong{\sphinxupquote{section}} (\sphinxstyleliteralemphasis{\sphinxupquote{str}}) \textendash{} The section of the configuration file.

\item {} 
\sphinxAtStartPar
\sphinxstyleliteralstrong{\sphinxupquote{name}} (\sphinxstyleliteralemphasis{\sphinxupquote{str}}) \textendash{} The name of the variable to read from the given section.

\item {} 
\sphinxAtStartPar
\sphinxstyleliteralstrong{\sphinxupquote{fallback}} (\sphinxstyleliteralemphasis{\sphinxupquote{str}}) \textendash{} (Optional) The default value to use if not found.

\end{itemize}

\sphinxlineitem{Returns}
\sphinxAtStartPar
Returns either the string found within the config file
for the given section and name, or the default value if not
found.

\sphinxlineitem{Return type}
\sphinxAtStartPar
str

\end{description}\end{quote}

\end{fulllineitems}

\index{get\_str\_alpha() (qnav.input.config\_handler.ConfigHandler method)@\spxentry{get\_str\_alpha()}\spxextra{qnav.input.config\_handler.ConfigHandler method}}

\begin{fulllineitems}
\phantomsection\label{\detokenize{modules/input/config_handler:qnav.input.config_handler.ConfigHandler.get_str_alpha}}
\pysigstartsignatures
\pysiglinewithargsret{\sphinxbfcode{\sphinxupquote{get\_str\_alpha}}}{\sphinxparam{\DUrole{n}{section}\DUrole{p}{:}\DUrole{w}{ }\DUrole{n}{str}}\sphinxparamcomma \sphinxparam{\DUrole{n}{name}\DUrole{p}{:}\DUrole{w}{ }\DUrole{n}{str}}\sphinxparamcomma \sphinxparam{\DUrole{n}{fallback}\DUrole{p}{:}\DUrole{w}{ }\DUrole{n}{str}\DUrole{w}{ }\DUrole{o}{=}\DUrole{w}{ }\DUrole{default_value}{None}}}{{ $\rightarrow$ str}}
\pysigstopsignatures
\sphinxAtStartPar
This function will return the alphanumeric string value for a
parameter in the selected configuration file, identified by its section
and name. If the value is missing, an optional default value can be
used instead.
\begin{quote}\begin{description}
\sphinxlineitem{Parameters}\begin{itemize}
\item {} 
\sphinxAtStartPar
\sphinxstyleliteralstrong{\sphinxupquote{section}} (\sphinxstyleliteralemphasis{\sphinxupquote{str}}) \textendash{} The section of the configuration file.

\item {} 
\sphinxAtStartPar
\sphinxstyleliteralstrong{\sphinxupquote{name}} (\sphinxstyleliteralemphasis{\sphinxupquote{str}}) \textendash{} The name of the variable to read from the given section.

\item {} 
\sphinxAtStartPar
\sphinxstyleliteralstrong{\sphinxupquote{fallback}} (\sphinxstyleliteralemphasis{\sphinxupquote{str}}) \textendash{} (Optional) The default value to use if not found.

\end{itemize}

\sphinxlineitem{Returns}
\sphinxAtStartPar
Returns either the alphanumeric parts of the string found
within the config file for the given section and name, or the
default value if not found.

\sphinxlineitem{Return type}
\sphinxAtStartPar
str

\end{description}\end{quote}

\end{fulllineitems}

\index{get\_int() (qnav.input.config\_handler.ConfigHandler method)@\spxentry{get\_int()}\spxextra{qnav.input.config\_handler.ConfigHandler method}}

\begin{fulllineitems}
\phantomsection\label{\detokenize{modules/input/config_handler:qnav.input.config_handler.ConfigHandler.get_int}}
\pysigstartsignatures
\pysiglinewithargsret{\sphinxbfcode{\sphinxupquote{get\_int}}}{\sphinxparam{\DUrole{n}{section}\DUrole{p}{:}\DUrole{w}{ }\DUrole{n}{str}}\sphinxparamcomma \sphinxparam{\DUrole{n}{name}\DUrole{p}{:}\DUrole{w}{ }\DUrole{n}{str}}\sphinxparamcomma \sphinxparam{\DUrole{n}{fallback}\DUrole{p}{:}\DUrole{w}{ }\DUrole{n}{int}\DUrole{w}{ }\DUrole{o}{=}\DUrole{w}{ }\DUrole{default_value}{None}}}{{ $\rightarrow$ int}}
\pysigstopsignatures
\sphinxAtStartPar
This function will return the integer value for a parameter in the
selected configuration file, identified by its section and name.
If the value is missing, an optional default value can be used instead.
\begin{quote}\begin{description}
\sphinxlineitem{Parameters}\begin{itemize}
\item {} 
\sphinxAtStartPar
\sphinxstyleliteralstrong{\sphinxupquote{section}} (\sphinxstyleliteralemphasis{\sphinxupquote{str}}) \textendash{} The section of the configuration file.

\item {} 
\sphinxAtStartPar
\sphinxstyleliteralstrong{\sphinxupquote{name}} (\sphinxstyleliteralemphasis{\sphinxupquote{str}}) \textendash{} The name of the variable to read from the given section.

\item {} 
\sphinxAtStartPar
\sphinxstyleliteralstrong{\sphinxupquote{fallback}} (\sphinxstyleliteralemphasis{\sphinxupquote{int}}) \textendash{} (Optional) The default value to use if not found.

\end{itemize}

\sphinxlineitem{Returns}
\sphinxAtStartPar
Returns either the integer value found within the config file
for the given section and name, or the default value if not found.

\sphinxlineitem{Return type}
\sphinxAtStartPar
int

\end{description}\end{quote}

\end{fulllineitems}

\index{get\_float() (qnav.input.config\_handler.ConfigHandler method)@\spxentry{get\_float()}\spxextra{qnav.input.config\_handler.ConfigHandler method}}

\begin{fulllineitems}
\phantomsection\label{\detokenize{modules/input/config_handler:qnav.input.config_handler.ConfigHandler.get_float}}
\pysigstartsignatures
\pysiglinewithargsret{\sphinxbfcode{\sphinxupquote{get\_float}}}{\sphinxparam{\DUrole{n}{section}\DUrole{p}{:}\DUrole{w}{ }\DUrole{n}{str}}\sphinxparamcomma \sphinxparam{\DUrole{n}{name}\DUrole{p}{:}\DUrole{w}{ }\DUrole{n}{str}}\sphinxparamcomma \sphinxparam{\DUrole{n}{fallback}\DUrole{p}{:}\DUrole{w}{ }\DUrole{n}{float}\DUrole{w}{ }\DUrole{o}{=}\DUrole{w}{ }\DUrole{default_value}{None}}}{{ $\rightarrow$ float}}
\pysigstopsignatures
\sphinxAtStartPar
This function will return the float value for a parameter in the
selected configuration file, identified by its section and name.
This function returns a numpy 64\sphinxhyphen{}bit float instead of a normal float.
If the value is missing, an optional default value can be used instead.
\begin{quote}\begin{description}
\sphinxlineitem{Parameters}\begin{itemize}
\item {} 
\sphinxAtStartPar
\sphinxstyleliteralstrong{\sphinxupquote{section}} (\sphinxstyleliteralemphasis{\sphinxupquote{str}}) \textendash{} The section of the configuration file.

\item {} 
\sphinxAtStartPar
\sphinxstyleliteralstrong{\sphinxupquote{name}} (\sphinxstyleliteralemphasis{\sphinxupquote{str}}) \textendash{} The name of the variable to read from the given section.

\item {} 
\sphinxAtStartPar
\sphinxstyleliteralstrong{\sphinxupquote{fallback}} (\sphinxstyleliteralemphasis{\sphinxupquote{float}}) \textendash{} (Optional) The default value to use if not found.

\end{itemize}

\sphinxlineitem{Returns}
\sphinxAtStartPar
Returns either the float value found within the config file
for the given section and name, or the default value if not found.

\sphinxlineitem{Return type}
\sphinxAtStartPar
float

\end{description}\end{quote}

\end{fulllineitems}

\index{get\_bool() (qnav.input.config\_handler.ConfigHandler method)@\spxentry{get\_bool()}\spxextra{qnav.input.config\_handler.ConfigHandler method}}

\begin{fulllineitems}
\phantomsection\label{\detokenize{modules/input/config_handler:qnav.input.config_handler.ConfigHandler.get_bool}}
\pysigstartsignatures
\pysiglinewithargsret{\sphinxbfcode{\sphinxupquote{get\_bool}}}{\sphinxparam{\DUrole{n}{section}\DUrole{p}{:}\DUrole{w}{ }\DUrole{n}{str}}\sphinxparamcomma \sphinxparam{\DUrole{n}{name}\DUrole{p}{:}\DUrole{w}{ }\DUrole{n}{str}}\sphinxparamcomma \sphinxparam{\DUrole{n}{fallback}\DUrole{p}{:}\DUrole{w}{ }\DUrole{n}{bool}\DUrole{w}{ }\DUrole{o}{=}\DUrole{w}{ }\DUrole{default_value}{None}}}{{ $\rightarrow$ bool}}
\pysigstopsignatures
\sphinxAtStartPar
This function will return the boolean value for a parameter in the
selected configuration file, identified by its section and name.
Acceptable True values are ‘yes’, ‘true’, ‘on’ and ‘1’. Acceptable
False values are ‘no’, ‘false’, ‘off’ and ‘0’. If the value is missing,
an optional default value can be used instead.
\begin{quote}\begin{description}
\sphinxlineitem{Parameters}\begin{itemize}
\item {} 
\sphinxAtStartPar
\sphinxstyleliteralstrong{\sphinxupquote{section}} (\sphinxstyleliteralemphasis{\sphinxupquote{str}}) \textendash{} The section of the configuration file.

\item {} 
\sphinxAtStartPar
\sphinxstyleliteralstrong{\sphinxupquote{name}} (\sphinxstyleliteralemphasis{\sphinxupquote{str}}) \textendash{} The name of the variable to read from the given section.

\item {} 
\sphinxAtStartPar
\sphinxstyleliteralstrong{\sphinxupquote{fallback}} (\sphinxstyleliteralemphasis{\sphinxupquote{bool}}) \textendash{} (Optional) The default value to use if not found.

\end{itemize}

\sphinxlineitem{Returns}
\sphinxAtStartPar
Returns either the boolean found within the config file for
the given section and name, or the default value if not found.

\sphinxlineitem{Return type}
\sphinxAtStartPar
bool

\end{description}\end{quote}

\end{fulllineitems}

\index{get\_csv\_numeric() (qnav.input.config\_handler.ConfigHandler method)@\spxentry{get\_csv\_numeric()}\spxextra{qnav.input.config\_handler.ConfigHandler method}}

\begin{fulllineitems}
\phantomsection\label{\detokenize{modules/input/config_handler:qnav.input.config_handler.ConfigHandler.get_csv_numeric}}
\pysigstartsignatures
\pysiglinewithargsret{\sphinxbfcode{\sphinxupquote{get\_csv\_numeric}}}{\sphinxparam{\DUrole{n}{section}\DUrole{p}{:}\DUrole{w}{ }\DUrole{n}{str}}\sphinxparamcomma \sphinxparam{\DUrole{n}{name}\DUrole{p}{:}\DUrole{w}{ }\DUrole{n}{str}}\sphinxparamcomma \sphinxparam{\DUrole{n}{fallback}\DUrole{p}{:}\DUrole{w}{ }\DUrole{n}{Buffer\DUrole{w}{ }\DUrole{p}{|}\DUrole{w}{ }\_SupportsArray\DUrole{p}{{[}}dtype\DUrole{p}{{[}}Any\DUrole{p}{{]}}\DUrole{p}{{]}}\DUrole{w}{ }\DUrole{p}{|}\DUrole{w}{ }\_NestedSequence\DUrole{p}{{[}}\_SupportsArray\DUrole{p}{{[}}dtype\DUrole{p}{{[}}Any\DUrole{p}{{]}}\DUrole{p}{{]}}\DUrole{p}{{]}}\DUrole{w}{ }\DUrole{p}{|}\DUrole{w}{ }bool\DUrole{w}{ }\DUrole{p}{|}\DUrole{w}{ }int\DUrole{w}{ }\DUrole{p}{|}\DUrole{w}{ }float\DUrole{w}{ }\DUrole{p}{|}\DUrole{w}{ }complex\DUrole{w}{ }\DUrole{p}{|}\DUrole{w}{ }str\DUrole{w}{ }\DUrole{p}{|}\DUrole{w}{ }bytes\DUrole{w}{ }\DUrole{p}{|}\DUrole{w}{ }\_NestedSequence\DUrole{p}{{[}}bool\DUrole{w}{ }\DUrole{p}{|}\DUrole{w}{ }int\DUrole{w}{ }\DUrole{p}{|}\DUrole{w}{ }float\DUrole{w}{ }\DUrole{p}{|}\DUrole{w}{ }complex\DUrole{w}{ }\DUrole{p}{|}\DUrole{w}{ }str\DUrole{w}{ }\DUrole{p}{|}\DUrole{w}{ }bytes\DUrole{p}{{]}}}\DUrole{w}{ }\DUrole{o}{=}\DUrole{w}{ }\DUrole{default_value}{None}}\sphinxparamcomma \sphinxparam{\DUrole{n}{sep}\DUrole{p}{:}\DUrole{w}{ }\DUrole{n}{str}\DUrole{w}{ }\DUrole{o}{=}\DUrole{w}{ }\DUrole{default_value}{\textquotesingle{},\textquotesingle{}}}}{{ $\rightarrow$ Buffer\DUrole{w}{ }\DUrole{p}{|}\DUrole{w}{ }\_SupportsArray\DUrole{p}{{[}}dtype\DUrole{p}{{[}}Any\DUrole{p}{{]}}\DUrole{p}{{]}}\DUrole{w}{ }\DUrole{p}{|}\DUrole{w}{ }\_NestedSequence\DUrole{p}{{[}}\_SupportsArray\DUrole{p}{{[}}dtype\DUrole{p}{{[}}Any\DUrole{p}{{]}}\DUrole{p}{{]}}\DUrole{p}{{]}}\DUrole{w}{ }\DUrole{p}{|}\DUrole{w}{ }bool\DUrole{w}{ }\DUrole{p}{|}\DUrole{w}{ }int\DUrole{w}{ }\DUrole{p}{|}\DUrole{w}{ }float\DUrole{w}{ }\DUrole{p}{|}\DUrole{w}{ }complex\DUrole{w}{ }\DUrole{p}{|}\DUrole{w}{ }str\DUrole{w}{ }\DUrole{p}{|}\DUrole{w}{ }bytes\DUrole{w}{ }\DUrole{p}{|}\DUrole{w}{ }\_NestedSequence\DUrole{p}{{[}}bool\DUrole{w}{ }\DUrole{p}{|}\DUrole{w}{ }int\DUrole{w}{ }\DUrole{p}{|}\DUrole{w}{ }float\DUrole{w}{ }\DUrole{p}{|}\DUrole{w}{ }complex\DUrole{w}{ }\DUrole{p}{|}\DUrole{w}{ }str\DUrole{w}{ }\DUrole{p}{|}\DUrole{w}{ }bytes\DUrole{p}{{]}}\DUrole{w}{ }\DUrole{p}{|}\DUrole{w}{ }None}}
\pysigstopsignatures
\sphinxAtStartPar
This function will return an array of values that are separated by a
common character (comma by default) from the selected configuration
file, identified by its section and name. This function returns a
numpy array. If the value is missing, an optional default value can
be used instead.
\begin{quote}\begin{description}
\sphinxlineitem{Parameters}\begin{itemize}
\item {} 
\sphinxAtStartPar
\sphinxstyleliteralstrong{\sphinxupquote{section}} (\sphinxstyleliteralemphasis{\sphinxupquote{str}}) \textendash{} The section of the configuration file.

\item {} 
\sphinxAtStartPar
\sphinxstyleliteralstrong{\sphinxupquote{name}} (\sphinxstyleliteralemphasis{\sphinxupquote{str}}) \textendash{} The name of the variable to read from the given section.

\item {} 
\sphinxAtStartPar
\sphinxstyleliteralstrong{\sphinxupquote{fallback}} (\sphinxstyleliteralemphasis{\sphinxupquote{list}}) \textendash{} (Optional) The default value(s) to use if not found.

\item {} 
\sphinxAtStartPar
\sphinxstyleliteralstrong{\sphinxupquote{sep}} \textendash{} (Optional) The item separator character (default=’,’).

\end{itemize}

\sphinxlineitem{Type}
\sphinxAtStartPar
str

\sphinxlineitem{Returns}
\sphinxAtStartPar
A numpy array containing the written items.

\sphinxlineitem{Return type}
\sphinxAtStartPar
list

\end{description}\end{quote}

\end{fulllineitems}

\index{get\_csv\_str() (qnav.input.config\_handler.ConfigHandler method)@\spxentry{get\_csv\_str()}\spxextra{qnav.input.config\_handler.ConfigHandler method}}

\begin{fulllineitems}
\phantomsection\label{\detokenize{modules/input/config_handler:qnav.input.config_handler.ConfigHandler.get_csv_str}}
\pysigstartsignatures
\pysiglinewithargsret{\sphinxbfcode{\sphinxupquote{get\_csv\_str}}}{\sphinxparam{\DUrole{n}{section}\DUrole{p}{:}\DUrole{w}{ }\DUrole{n}{str}}\sphinxparamcomma \sphinxparam{\DUrole{n}{name}\DUrole{p}{:}\DUrole{w}{ }\DUrole{n}{str}}\sphinxparamcomma \sphinxparam{\DUrole{n}{fallback}\DUrole{p}{:}\DUrole{w}{ }\DUrole{n}{list}\DUrole{w}{ }\DUrole{o}{=}\DUrole{w}{ }\DUrole{default_value}{None}}\sphinxparamcomma \sphinxparam{\DUrole{n}{sep}\DUrole{p}{:}\DUrole{w}{ }\DUrole{n}{str}\DUrole{w}{ }\DUrole{o}{=}\DUrole{w}{ }\DUrole{default_value}{\textquotesingle{},\textquotesingle{}}}}{{ $\rightarrow$ list\DUrole{w}{ }\DUrole{p}{|}\DUrole{w}{ }None}}
\pysigstopsignatures
\sphinxAtStartPar
This function will return an array of string values that are
separated by a common character (comma by default) from the selected
configuration file, identified by its section and name. This function
returns a list of one or more character values. If the value is
missing, an optional default value can be used instead.
\begin{quote}\begin{description}
\sphinxlineitem{Parameters}\begin{itemize}
\item {} 
\sphinxAtStartPar
\sphinxstyleliteralstrong{\sphinxupquote{section}} (\sphinxstyleliteralemphasis{\sphinxupquote{str}}) \textendash{} The section of the configuration file.

\item {} 
\sphinxAtStartPar
\sphinxstyleliteralstrong{\sphinxupquote{name}} (\sphinxstyleliteralemphasis{\sphinxupquote{str}}) \textendash{} The name of the variable to read from the given section.

\item {} 
\sphinxAtStartPar
\sphinxstyleliteralstrong{\sphinxupquote{fallback}} (\sphinxstyleliteralemphasis{\sphinxupquote{list}}\sphinxstyleliteralemphasis{\sphinxupquote{ | }}\sphinxstyleliteralemphasis{\sphinxupquote{None}}) \textendash{} (Optional) The default value(s) to use if not found.

\item {} 
\sphinxAtStartPar
\sphinxstyleliteralstrong{\sphinxupquote{sep}} \textendash{} (Optional) The item separator character (default=’,’).

\end{itemize}

\sphinxlineitem{Type}
\sphinxAtStartPar
str

\sphinxlineitem{Returns}
\sphinxAtStartPar
A Python list containing the written items.

\sphinxlineitem{Return type}
\sphinxAtStartPar
list | None

\end{description}\end{quote}

\end{fulllineitems}

\index{get\_value\_filepath() (qnav.input.config\_handler.ConfigHandler method)@\spxentry{get\_value\_filepath()}\spxextra{qnav.input.config\_handler.ConfigHandler method}}

\begin{fulllineitems}
\phantomsection\label{\detokenize{modules/input/config_handler:qnav.input.config_handler.ConfigHandler.get_value_filepath}}
\pysigstartsignatures
\pysiglinewithargsret{\sphinxbfcode{\sphinxupquote{get\_value\_filepath}}}{\sphinxparam{\DUrole{n}{section}\DUrole{p}{:}\DUrole{w}{ }\DUrole{n}{str}}\sphinxparamcomma \sphinxparam{\DUrole{n}{name}\DUrole{p}{:}\DUrole{w}{ }\DUrole{n}{str}}\sphinxparamcomma \sphinxparam{\DUrole{n}{fallback}\DUrole{p}{:}\DUrole{w}{ }\DUrole{n}{str\DUrole{w}{ }\DUrole{p}{|}\DUrole{w}{ }Path}\DUrole{w}{ }\DUrole{o}{=}\DUrole{w}{ }\DUrole{default_value}{None}}\sphinxparamcomma \sphinxparam{\DUrole{n}{must\_exist}\DUrole{p}{:}\DUrole{w}{ }\DUrole{n}{bool}\DUrole{w}{ }\DUrole{o}{=}\DUrole{w}{ }\DUrole{default_value}{False}}}{{ $\rightarrow$ Path}}
\pysigstopsignatures
\sphinxAtStartPar
This function will return a local filepath read for a parameter in the
selected configuration file, identified by its section and name. This
function will handle read paths in a platform independent manner,
automatically deciding and converting when needed. This allows paths
read from the configuration file to be written in either Linux/Unix or
Windows/NT format without affecting behaviour on different platforms.
\begin{quote}\begin{description}
\sphinxlineitem{Parameters}\begin{itemize}
\item {} 
\sphinxAtStartPar
\sphinxstyleliteralstrong{\sphinxupquote{section}} (\sphinxstyleliteralemphasis{\sphinxupquote{str}}) \textendash{} The section of the configuration file.

\item {} 
\sphinxAtStartPar
\sphinxstyleliteralstrong{\sphinxupquote{name}} (\sphinxstyleliteralemphasis{\sphinxupquote{str}}) \textendash{} The name of the variable to read from the given section.

\item {} 
\sphinxAtStartPar
\sphinxstyleliteralstrong{\sphinxupquote{fallback}} (\sphinxstyleliteralemphasis{\sphinxupquote{str}}) \textendash{} (Optional) The default value to use if not found.

\item {} 
\sphinxAtStartPar
\sphinxstyleliteralstrong{\sphinxupquote{must\_exist}} \textendash{} (Optional) Should the path exist on the filesystem?
If enabled an error will raise if the given path is not found on
the system (default=True).

\end{itemize}

\sphinxlineitem{Returns}
\sphinxAtStartPar
A valid filepath for the current system to the given
location/resources.

\sphinxlineitem{Return type}
\sphinxAtStartPar
Path

\end{description}\end{quote}

\end{fulllineitems}

\index{get\_value\_filepath\_str() (qnav.input.config\_handler.ConfigHandler method)@\spxentry{get\_value\_filepath\_str()}\spxextra{qnav.input.config\_handler.ConfigHandler method}}

\begin{fulllineitems}
\phantomsection\label{\detokenize{modules/input/config_handler:qnav.input.config_handler.ConfigHandler.get_value_filepath_str}}
\pysigstartsignatures
\pysiglinewithargsret{\sphinxbfcode{\sphinxupquote{get\_value\_filepath\_str}}}{\sphinxparam{\DUrole{n}{section}\DUrole{p}{:}\DUrole{w}{ }\DUrole{n}{str}}\sphinxparamcomma \sphinxparam{\DUrole{n}{name}\DUrole{p}{:}\DUrole{w}{ }\DUrole{n}{str}}\sphinxparamcomma \sphinxparam{\DUrole{n}{fallback}\DUrole{p}{:}\DUrole{w}{ }\DUrole{n}{str}\DUrole{w}{ }\DUrole{o}{=}\DUrole{w}{ }\DUrole{default_value}{None}}}{}
\pysigstopsignatures
\sphinxAtStartPar
This function will return a local filepath read for a parameter in the
selected configuration file, identified by its section and name. This
function will handle read paths in a platform independent manner,
automatically deciding and converting when needed. This allows paths
read from the configuration file to be written in either Linux/Unix or
Windows/NT format without affecting behaviour on different platforms.

\sphinxAtStartPar
Unlike ‘get\_value\_filepath’ this function does not enforce the
existence of a requested file/database\_dir.
\begin{quote}\begin{description}
\sphinxlineitem{Parameters}\begin{itemize}
\item {} 
\sphinxAtStartPar
\sphinxstyleliteralstrong{\sphinxupquote{section}} (\sphinxstyleliteralemphasis{\sphinxupquote{str}}) \textendash{} The section of the configuration file.

\item {} 
\sphinxAtStartPar
\sphinxstyleliteralstrong{\sphinxupquote{name}} (\sphinxstyleliteralemphasis{\sphinxupquote{str}}) \textendash{} The name of the variable to read from the given section.

\item {} 
\sphinxAtStartPar
\sphinxstyleliteralstrong{\sphinxupquote{fallback}} (\sphinxstyleliteralemphasis{\sphinxupquote{str}}) \textendash{} (Optional) The default value to use if not found.

\end{itemize}

\sphinxlineitem{Returns}
\sphinxAtStartPar
A filepath for the current system to the given location/
resources. This resource does not have to exist.

\sphinxlineitem{Return type}
\sphinxAtStartPar
Path

\end{description}\end{quote}

\end{fulllineitems}

\index{is\_value\_set() (qnav.input.config\_handler.ConfigHandler method)@\spxentry{is\_value\_set()}\spxextra{qnav.input.config\_handler.ConfigHandler method}}

\begin{fulllineitems}
\phantomsection\label{\detokenize{modules/input/config_handler:qnav.input.config_handler.ConfigHandler.is_value_set}}
\pysigstartsignatures
\pysiglinewithargsret{\sphinxbfcode{\sphinxupquote{is\_value\_set}}}{\sphinxparam{\DUrole{n}{section}\DUrole{p}{:}\DUrole{w}{ }\DUrole{n}{list}}\sphinxparamcomma \sphinxparam{\DUrole{n}{value}\DUrole{p}{:}\DUrole{w}{ }\DUrole{n}{list}}}{{ $\rightarrow$ bool}}
\pysigstopsignatures
\sphinxAtStartPar
Returns true if there is an entry for (all of) the given sections and
values. This function works for both single values and lists of
values. For multiple values, both lists must be of the same length.
\begin{quote}\begin{description}
\sphinxlineitem{Parameters}\begin{itemize}
\item {} 
\sphinxAtStartPar
\sphinxstyleliteralstrong{\sphinxupquote{section}} (\sphinxstyleliteralemphasis{\sphinxupquote{list}}\sphinxstyleliteralemphasis{\sphinxupquote{ or }}\sphinxstyleliteralemphasis{\sphinxupquote{str}}) \textendash{} A single or list of sections names.

\item {} 
\sphinxAtStartPar
\sphinxstyleliteralstrong{\sphinxupquote{value}} \textendash{} A single or list of value names.

\end{itemize}

\sphinxlineitem{Returns}
\sphinxAtStartPar
True if an entry could be found in the provided configuration
for all the given section\sphinxhyphen{}list pairs. False if one or more entries
could not be found.

\sphinxlineitem{Return type}
\sphinxAtStartPar
bool

\end{description}\end{quote}

\end{fulllineitems}

\index{overwrite\_value() (qnav.input.config\_handler.ConfigHandler method)@\spxentry{overwrite\_value()}\spxextra{qnav.input.config\_handler.ConfigHandler method}}

\begin{fulllineitems}
\phantomsection\label{\detokenize{modules/input/config_handler:qnav.input.config_handler.ConfigHandler.overwrite_value}}
\pysigstartsignatures
\pysiglinewithargsret{\sphinxbfcode{\sphinxupquote{overwrite\_value}}}{\sphinxparam{\DUrole{n}{section}\DUrole{p}{:}\DUrole{w}{ }\DUrole{n}{str}}\sphinxparamcomma \sphinxparam{\DUrole{n}{name}\DUrole{p}{:}\DUrole{w}{ }\DUrole{n}{str}}\sphinxparamcomma \sphinxparam{\DUrole{n}{new\_value}}}{}
\pysigstopsignatures
\sphinxAtStartPar
This simple function can be used to overwrite or insert any value
extracted from the given configuration file. This is only to be used
for automatically generated values and is not an alternative to
fallbacks.
\begin{quote}\begin{description}
\sphinxlineitem{Parameters}\begin{itemize}
\item {} 
\sphinxAtStartPar
\sphinxstyleliteralstrong{\sphinxupquote{section}} (\sphinxstyleliteralemphasis{\sphinxupquote{str}}) \textendash{} The section name of the variable to override

\item {} 
\sphinxAtStartPar
\sphinxstyleliteralstrong{\sphinxupquote{name}} (\sphinxstyleliteralemphasis{\sphinxupquote{str}}) \textendash{} The name of the variable to read from the given section.

\item {} 
\sphinxAtStartPar
\sphinxstyleliteralstrong{\sphinxupquote{new\_value}} (\sphinxstyleliteralemphasis{\sphinxupquote{str}}\sphinxstyleliteralemphasis{\sphinxupquote{, }}\sphinxstyleliteralemphasis{\sphinxupquote{int}}\sphinxstyleliteralemphasis{\sphinxupquote{, }}\sphinxstyleliteralemphasis{\sphinxupquote{float}}) \textendash{} The new value to use for the given section and name.

\end{itemize}

\end{description}\end{quote}

\end{fulllineitems}

\index{copy\_file() (qnav.input.config\_handler.ConfigHandler method)@\spxentry{copy\_file()}\spxextra{qnav.input.config\_handler.ConfigHandler method}}

\begin{fulllineitems}
\phantomsection\label{\detokenize{modules/input/config_handler:qnav.input.config_handler.ConfigHandler.copy_file}}
\pysigstartsignatures
\pysiglinewithargsret{\sphinxbfcode{\sphinxupquote{copy\_file}}}{\sphinxparam{\DUrole{n}{dest}\DUrole{p}{:}\DUrole{w}{ }\DUrole{n}{Path}}}{}
\pysigstopsignatures
\sphinxAtStartPar
Copies the configuration file to a destination.
This method produces a copy of the configuration file at the
requested destination.
\begin{quote}\begin{description}
\sphinxlineitem{Parameters}
\sphinxAtStartPar
\sphinxstyleliteralstrong{\sphinxupquote{dest}} (\sphinxstyleliteralemphasis{\sphinxupquote{Path}}) \textendash{} The location to write the file to.

\end{description}\end{quote}

\end{fulllineitems}

\index{write\_file() (qnav.input.config\_handler.ConfigHandler method)@\spxentry{write\_file()}\spxextra{qnav.input.config\_handler.ConfigHandler method}}

\begin{fulllineitems}
\phantomsection\label{\detokenize{modules/input/config_handler:qnav.input.config_handler.ConfigHandler.write_file}}
\pysigstartsignatures
\pysiglinewithargsret{\sphinxbfcode{\sphinxupquote{write\_file}}}{\sphinxparam{\DUrole{n}{dest}\DUrole{p}{:}\DUrole{w}{ }\DUrole{n}{Path}}}{}
\pysigstopsignatures
\sphinxAtStartPar
Writes the configuration values to a file.
This method writes the configuration values in their current state
(including modifications) to a requested file. The produced file can
then be re\sphinxhyphen{}used as a settings file to repeat experimentation.
\begin{quote}\begin{description}
\sphinxlineitem{Parameters}
\sphinxAtStartPar
\sphinxstyleliteralstrong{\sphinxupquote{dest}} (\sphinxstyleliteralemphasis{\sphinxupquote{Path}}) \textendash{} The location to write the file to.

\end{description}\end{quote}

\end{fulllineitems}


\end{fulllineitems}

\index{NavConfigWarning@\spxentry{NavConfigWarning}}

\begin{fulllineitems}
\phantomsection\label{\detokenize{modules/input/config_handler:qnav.input.config_handler.NavConfigWarning}}
\pysigstartsignatures
\pysigline{\sphinxbfcode{\sphinxupquote{exception\DUrole{w}{ }}}\sphinxcode{\sphinxupquote{qnav.input.config\_handler.}}\sphinxbfcode{\sphinxupquote{NavConfigWarning}}}
\pysigstopsignatures
\sphinxAtStartPar
A custom warning type used to inform the user about invalid, unsuitable
or missing values present in the configuration file. If these warnings
generated can be safely ignored, this warning type can be suppressed
to hide all configuration file warnings for the session. Suppressing
this warning type will not affect the displaying of other warnings.

\end{fulllineitems}

\index{NavConfigError@\spxentry{NavConfigError}}

\begin{fulllineitems}
\phantomsection\label{\detokenize{modules/input/config_handler:qnav.input.config_handler.NavConfigError}}
\pysigstartsignatures
\pysiglinewithargsret{\sphinxbfcode{\sphinxupquote{exception\DUrole{w}{ }}}\sphinxcode{\sphinxupquote{qnav.input.config\_handler.}}\sphinxbfcode{\sphinxupquote{NavConfigError}}}{\sphinxparam{\DUrole{n}{section\_name}\DUrole{p}{:}\DUrole{w}{ }\DUrole{n}{str}}\sphinxparamcomma \sphinxparam{\DUrole{n}{variable\_name}\DUrole{p}{:}\DUrole{w}{ }\DUrole{n}{str}}\sphinxparamcomma \sphinxparam{\DUrole{n}{error\_type}\DUrole{p}{:}\DUrole{w}{ }\DUrole{n}{str}}\sphinxparamcomma \sphinxparam{\DUrole{n}{error\_desc}\DUrole{p}{:}\DUrole{w}{ }\DUrole{n}{str}}}{}
\pysigstopsignatures
\sphinxAtStartPar
A custom exception type used for inform the user about erroneous values
present in the provided configuration file that cannot be ignored. This
includes syntax, value type and conflicting value errors. This exception
type is primarily used for formatting the errors being displayed. This
custom class also allows configuration based errors to be distinguished
from serious errors such faults not detected until after delivery.

\end{fulllineitems}



\subsection{Measurement}
\label{\detokenize{usage/packages:measurement}}
\sphinxAtStartPar
Modules related to standard sensors for measurements:

\sphinxstepscope


\subsubsection{qnav.measurement.sensor}
\label{\detokenize{modules/measurement/sensor:module-qnav.measurement.sensor}}\label{\detokenize{modules/measurement/sensor:qnav-measurement-sensor}}\label{\detokenize{modules/measurement/sensor::doc}}\index{module@\spxentry{module}!qnav.measurement.sensor@\spxentry{qnav.measurement.sensor}}\index{qnav.measurement.sensor@\spxentry{qnav.measurement.sensor}!module@\spxentry{module}}

\paragraph{sensor.py}
\label{\detokenize{modules/measurement/sensor:sensor-py}}\begin{quote}\begin{description}
\sphinxlineitem{summary}
\sphinxAtStartPar
Defines the abstract base classes for sensor usage in simulations.
This module contains the base classes for both sensors (used for taking
GroundTruth records and producing measurements) and sensor fusion (methods
for applying measurement information to the current estimated states).

\sphinxlineitem{author}
\sphinxAtStartPar
Michael Wright \sphinxhyphen{} \sphinxhref{mailto:mjwright@liverpool.ac.uk}{mjwright@liverpool.ac.uk}

\sphinxlineitem{version}
\sphinxAtStartPar
1.0.0 \sphinxhyphen{} First full implementation of class.

\end{description}\end{quote}
\index{Sensor (class in qnav.measurement.sensor)@\spxentry{Sensor}\spxextra{class in qnav.measurement.sensor}}

\begin{fulllineitems}
\phantomsection\label{\detokenize{modules/measurement/sensor:qnav.measurement.sensor.Sensor}}
\pysigstartsignatures
\pysiglinewithargsret{\sphinxbfcode{\sphinxupquote{class\DUrole{w}{ }}}\sphinxcode{\sphinxupquote{qnav.measurement.sensor.}}\sphinxbfcode{\sphinxupquote{Sensor}}}{\sphinxparam{\DUrole{n}{frequency}\DUrole{p}{:}\DUrole{w}{ }\DUrole{n}{float}}\sphinxparamcomma \sphinxparam{\DUrole{n}{start\_time}\DUrole{p}{:}\DUrole{w}{ }\DUrole{n}{float}\DUrole{w}{ }\DUrole{o}{=}\DUrole{w}{ }\DUrole{default_value}{0.0}}\sphinxparamcomma \sphinxparam{\DUrole{n}{rand\_seed}\DUrole{p}{:}\DUrole{w}{ }\DUrole{n}{int}\DUrole{w}{ }\DUrole{o}{=}\DUrole{w}{ }\DUrole{default_value}{None}}}{}
\pysigstopsignatures
\sphinxAtStartPar
Base class for sensors used in navigation.
Defines the format of sensors used in the toolbox, specifically the
methods, inputs and outputs expected for fusion and other classes.
\index{take\_measurement() (qnav.measurement.sensor.Sensor method)@\spxentry{take\_measurement()}\spxextra{qnav.measurement.sensor.Sensor method}}

\begin{fulllineitems}
\phantomsection\label{\detokenize{modules/measurement/sensor:qnav.measurement.sensor.Sensor.take_measurement}}
\pysigstartsignatures
\pysiglinewithargsret{\sphinxbfcode{\sphinxupquote{abstract\DUrole{w}{ }}}\sphinxbfcode{\sphinxupquote{take\_measurement}}}{\sphinxparam{\DUrole{n}{est\_time}\DUrole{p}{:}\DUrole{w}{ }\DUrole{n}{float}}\sphinxparamcomma \sphinxparam{\DUrole{n}{ground\_truth}\DUrole{p}{:}\DUrole{w}{ }\DUrole{n}{{\hyperref[\detokenize{modules/waypoints/trajectory:qnav.waypoints.trajectory.GroundTruth}]{\sphinxcrossref{GroundTruth}}}}}}{}
\pysigstopsignatures
\sphinxAtStartPar
Produces the measurement captured by the sensor.
Given the current estimated time and the ground truth, this method
shall return the measurement captured at that time. Alternatively, a
None value can be returned to signify a measurement is not ready.
\begin{quote}\begin{description}
\sphinxlineitem{Parameters}\begin{itemize}
\item {} 
\sphinxAtStartPar
\sphinxstyleliteralstrong{\sphinxupquote{est\_time}} (\sphinxstyleliteralemphasis{\sphinxupquote{float}}) \textendash{} The estimated time the measurement of the sensor
is expected to be captured at (in seconds).

\item {} 
\sphinxAtStartPar
\sphinxstyleliteralstrong{\sphinxupquote{ground\_truth}} ({\hyperref[\detokenize{modules/waypoints/trajectory:qnav.waypoints.trajectory.GroundTruth}]{\sphinxcrossref{\sphinxstyleliteralemphasis{\sphinxupquote{GroundTruth}}}}}) \textendash{} The corresponding ground truth data.

\end{itemize}

\end{description}\end{quote}

\end{fulllineitems}

\index{last\_measurement (qnav.measurement.sensor.Sensor property)@\spxentry{last\_measurement}\spxextra{qnav.measurement.sensor.Sensor property}}

\begin{fulllineitems}
\phantomsection\label{\detokenize{modules/measurement/sensor:qnav.measurement.sensor.Sensor.last_measurement}}
\pysigstartsignatures
\pysigline{\sphinxbfcode{\sphinxupquote{abstract\DUrole{w}{ }property\DUrole{w}{ }}}\sphinxbfcode{\sphinxupquote{last\_measurement}}\sphinxbfcode{\sphinxupquote{\DUrole{p}{:}\DUrole{w}{ }any}}}
\pysigstopsignatures
\sphinxAtStartPar
Returns the latest recorded measurement produced by sensor.
\begin{quote}\begin{description}
\sphinxlineitem{Returns}
\sphinxAtStartPar
Latest measurement produced by sensor.

\sphinxlineitem{Return type}
\sphinxAtStartPar
any

\end{description}\end{quote}

\end{fulllineitems}

\index{update() (qnav.measurement.sensor.Sensor method)@\spxentry{update()}\spxextra{qnav.measurement.sensor.Sensor method}}

\begin{fulllineitems}
\phantomsection\label{\detokenize{modules/measurement/sensor:qnav.measurement.sensor.Sensor.update}}
\pysigstartsignatures
\pysiglinewithargsret{\sphinxbfcode{\sphinxupquote{update}}}{\sphinxparam{\DUrole{n}{time\_sec}\DUrole{p}{:}\DUrole{w}{ }\DUrole{n}{float}\DUrole{w}{ }\DUrole{o}{=}\DUrole{w}{ }\DUrole{default_value}{None}}}{}
\pysigstopsignatures
\sphinxAtStartPar
Called on demand to update the sensor’s state.
Typically called after measurements have been taken. When called
this increases the next measurement time and propagates internal
states, to the time difference.
\begin{quote}\begin{description}
\sphinxlineitem{Parameters}
\sphinxAtStartPar
\sphinxstyleliteralstrong{\sphinxupquote{time\_sec}} (\sphinxstyleliteralemphasis{\sphinxupquote{float}}) \textendash{} The current simulation time in seconds.

\end{description}\end{quote}

\end{fulllineitems}

\index{last\_update (qnav.measurement.sensor.Sensor property)@\spxentry{last\_update}\spxextra{qnav.measurement.sensor.Sensor property}}

\begin{fulllineitems}
\phantomsection\label{\detokenize{modules/measurement/sensor:qnav.measurement.sensor.Sensor.last_update}}
\pysigstartsignatures
\pysigline{\sphinxbfcode{\sphinxupquote{property\DUrole{w}{ }}}\sphinxbfcode{\sphinxupquote{last\_update}}\sphinxbfcode{\sphinxupquote{\DUrole{p}{:}\DUrole{w}{ }float}}}
\pysigstopsignatures
\sphinxAtStartPar
Returns the time (in seconds) of the last measurement.
\begin{quote}\begin{description}
\sphinxlineitem{Returns}
\sphinxAtStartPar
Time of last measurement (in seconds).

\sphinxlineitem{Return type}
\sphinxAtStartPar
float

\end{description}\end{quote}

\end{fulllineitems}

\index{next\_update (qnav.measurement.sensor.Sensor property)@\spxentry{next\_update}\spxextra{qnav.measurement.sensor.Sensor property}}

\begin{fulllineitems}
\phantomsection\label{\detokenize{modules/measurement/sensor:qnav.measurement.sensor.Sensor.next_update}}
\pysigstartsignatures
\pysigline{\sphinxbfcode{\sphinxupquote{property\DUrole{w}{ }}}\sphinxbfcode{\sphinxupquote{next\_update}}\sphinxbfcode{\sphinxupquote{\DUrole{p}{:}\DUrole{w}{ }float}}}
\pysigstopsignatures
\sphinxAtStartPar
Returns the due time (in seconds) of the next measurement.
\begin{quote}\begin{description}
\sphinxlineitem{Returns}
\sphinxAtStartPar
Time of next measurement (in seconds).

\sphinxlineitem{Return type}
\sphinxAtStartPar
float

\end{description}\end{quote}

\end{fulllineitems}

\index{frequency (qnav.measurement.sensor.Sensor property)@\spxentry{frequency}\spxextra{qnav.measurement.sensor.Sensor property}}

\begin{fulllineitems}
\phantomsection\label{\detokenize{modules/measurement/sensor:qnav.measurement.sensor.Sensor.frequency}}
\pysigstartsignatures
\pysigline{\sphinxbfcode{\sphinxupquote{property\DUrole{w}{ }}}\sphinxbfcode{\sphinxupquote{frequency}}\sphinxbfcode{\sphinxupquote{\DUrole{p}{:}\DUrole{w}{ }float}}}
\pysigstopsignatures
\sphinxAtStartPar
Returns the sensor’s measurement frequency in Hz.
\begin{quote}\begin{description}
\sphinxlineitem{Returns}
\sphinxAtStartPar
The frequency in Hz.

\sphinxlineitem{Return type}
\sphinxAtStartPar
float

\end{description}\end{quote}

\end{fulllineitems}

\index{time\_step (qnav.measurement.sensor.Sensor property)@\spxentry{time\_step}\spxextra{qnav.measurement.sensor.Sensor property}}

\begin{fulllineitems}
\phantomsection\label{\detokenize{modules/measurement/sensor:qnav.measurement.sensor.Sensor.time_step}}
\pysigstartsignatures
\pysigline{\sphinxbfcode{\sphinxupquote{property\DUrole{w}{ }}}\sphinxbfcode{\sphinxupquote{time\_step}}\sphinxbfcode{\sphinxupquote{\DUrole{p}{:}\DUrole{w}{ }float}}}
\pysigstopsignatures
\sphinxAtStartPar
Returns the sensor’s measurement time step in seconds.
\begin{quote}\begin{description}
\sphinxlineitem{Returns}
\sphinxAtStartPar
The measurement interval in seconds.

\sphinxlineitem{Return type}
\sphinxAtStartPar
float

\end{description}\end{quote}

\end{fulllineitems}


\end{fulllineitems}

\index{FusionTrigger (class in qnav.measurement.sensor)@\spxentry{FusionTrigger}\spxextra{class in qnav.measurement.sensor}}

\begin{fulllineitems}
\phantomsection\label{\detokenize{modules/measurement/sensor:qnav.measurement.sensor.FusionTrigger}}
\pysigstartsignatures
\pysiglinewithargsret{\sphinxbfcode{\sphinxupquote{class\DUrole{w}{ }}}\sphinxcode{\sphinxupquote{qnav.measurement.sensor.}}\sphinxbfcode{\sphinxupquote{FusionTrigger}}}{\sphinxparam{\DUrole{n}{value}}\sphinxparamcomma \sphinxparam{\DUrole{n}{names=\textless{}not given\textgreater{}}}\sphinxparamcomma \sphinxparam{\DUrole{n}{*values}}\sphinxparamcomma \sphinxparam{\DUrole{n}{module=None}}\sphinxparamcomma \sphinxparam{\DUrole{n}{qualname=None}}\sphinxparamcomma \sphinxparam{\DUrole{n}{type=None}}\sphinxparamcomma \sphinxparam{\DUrole{n}{start=1}}\sphinxparamcomma \sphinxparam{\DUrole{n}{boundary=None}}}{}
\pysigstopsignatures
\sphinxAtStartPar
Flags used to indicate the triggering mode for fusion.
These essentially define when fusion is triggered, subject to the sensor
measurements that ara available.

\end{fulllineitems}

\index{SensorFusion (class in qnav.measurement.sensor)@\spxentry{SensorFusion}\spxextra{class in qnav.measurement.sensor}}

\begin{fulllineitems}
\phantomsection\label{\detokenize{modules/measurement/sensor:qnav.measurement.sensor.SensorFusion}}
\pysigstartsignatures
\pysiglinewithargsret{\sphinxbfcode{\sphinxupquote{class\DUrole{w}{ }}}\sphinxcode{\sphinxupquote{qnav.measurement.sensor.}}\sphinxbfcode{\sphinxupquote{SensorFusion}}}{\sphinxparam{\DUrole{n}{trigger}\DUrole{p}{:}\DUrole{w}{ }\DUrole{n}{{\hyperref[\detokenize{modules/measurement/sensor:qnav.measurement.sensor.FusionTrigger}]{\sphinxcrossref{FusionTrigger}}}}}\sphinxparamcomma \sphinxparam{\DUrole{o}{*}\DUrole{n}{sensors}\DUrole{p}{:}\DUrole{w}{ }\DUrole{n}{{\hyperref[\detokenize{modules/measurement/sensor:qnav.measurement.sensor.Sensor}]{\sphinxcrossref{Sensor}}}}}}{}
\pysigstopsignatures
\sphinxAtStartPar
Base class for a sensor fusion method used in navigation.
Defines the format of data fusion used in the toolbox, specifically the
methods, inputs and outputs expected.
\index{trigger (qnav.measurement.sensor.SensorFusion property)@\spxentry{trigger}\spxextra{qnav.measurement.sensor.SensorFusion property}}

\begin{fulllineitems}
\phantomsection\label{\detokenize{modules/measurement/sensor:qnav.measurement.sensor.SensorFusion.trigger}}
\pysigstartsignatures
\pysigline{\sphinxbfcode{\sphinxupquote{property\DUrole{w}{ }}}\sphinxbfcode{\sphinxupquote{trigger}}\sphinxbfcode{\sphinxupquote{\DUrole{p}{:}\DUrole{w}{ }{\hyperref[\detokenize{modules/measurement/sensor:qnav.measurement.sensor.FusionTrigger}]{\sphinxcrossref{FusionTrigger}}}}}}
\pysigstopsignatures
\sphinxAtStartPar
The selected trigger method for deciding when to apply fusion.
This decides when fusion can be applied, subject to the available
sensor measurements.
\begin{quote}\begin{description}
\sphinxlineitem{Returns}
\sphinxAtStartPar
The selected fusion trigger method.

\sphinxlineitem{Return type}
\sphinxAtStartPar
{\hyperref[\detokenize{modules/measurement/sensor:qnav.measurement.sensor.FusionTrigger}]{\sphinxcrossref{FusionTrigger}}}

\end{description}\end{quote}

\end{fulllineitems}

\index{sensors (qnav.measurement.sensor.SensorFusion property)@\spxentry{sensors}\spxextra{qnav.measurement.sensor.SensorFusion property}}

\begin{fulllineitems}
\phantomsection\label{\detokenize{modules/measurement/sensor:qnav.measurement.sensor.SensorFusion.sensors}}
\pysigstartsignatures
\pysigline{\sphinxbfcode{\sphinxupquote{property\DUrole{w}{ }}}\sphinxbfcode{\sphinxupquote{sensors}}\sphinxbfcode{\sphinxupquote{\DUrole{p}{:}\DUrole{w}{ }Iterable\DUrole{p}{{[}}{\hyperref[\detokenize{modules/measurement/sensor:qnav.measurement.sensor.Sensor}]{\sphinxcrossref{Sensor}}}\DUrole{p}{{]}}}}}
\pysigstopsignatures
\sphinxAtStartPar
The provided sensors associated with the fusion method.
This is a collection of all shared sensors whose measurements will be
used by the fusion method.
\begin{quote}\begin{description}
\sphinxlineitem{Returns}
\sphinxAtStartPar
The shared sensors used by the fusion method.

\sphinxlineitem{Return type}
\sphinxAtStartPar
Iterable{[}{\hyperref[\detokenize{modules/measurement/sensor:qnav.measurement.sensor.Sensor}]{\sphinxcrossref{Sensor}}}{]}

\end{description}\end{quote}

\end{fulllineitems}

\index{perform\_fusion() (qnav.measurement.sensor.SensorFusion method)@\spxentry{perform\_fusion()}\spxextra{qnav.measurement.sensor.SensorFusion method}}

\begin{fulllineitems}
\phantomsection\label{\detokenize{modules/measurement/sensor:qnav.measurement.sensor.SensorFusion.perform_fusion}}
\pysigstartsignatures
\pysiglinewithargsret{\sphinxbfcode{\sphinxupquote{abstract\DUrole{w}{ }}}\sphinxbfcode{\sphinxupquote{perform\_fusion}}}{\sphinxparam{\DUrole{n}{estimated\_state}\DUrole{p}{:}\DUrole{w}{ }\DUrole{n}{{\hyperref[\detokenize{modules/estimation/state:qnav.estimation.state.EstimatedState}]{\sphinxcrossref{EstimatedState}}}}}}{{ $\rightarrow$ None}}
\pysigstopsignatures
\sphinxAtStartPar
Performs the fusion with sensors and updates the estimated state.
When called at scheduled time during a simulation loop, this method
performs the data fusion by extracting last measurements from sensors
and using them with the fusion algorithm to update the given current
estimated state.
\begin{quote}\begin{description}
\sphinxlineitem{Parameters}
\sphinxAtStartPar
\sphinxstyleliteralstrong{\sphinxupquote{estimated\_state}} ({\hyperref[\detokenize{modules/estimation/state:qnav.estimation.state.EstimatedState}]{\sphinxcrossref{\sphinxstyleliteralemphasis{\sphinxupquote{EstimatedState}}}}}) \textendash{} The current estimated states of the platform.
This is what the fusion method will update.

\end{description}\end{quote}

\end{fulllineitems}


\end{fulllineitems}


\sphinxstepscope


\subsubsection{qnav.measurement.platform}
\label{\detokenize{modules/measurement/platform:module-qnav.measurement.platform}}\label{\detokenize{modules/measurement/platform:qnav-measurement-platform}}\label{\detokenize{modules/measurement/platform::doc}}\index{module@\spxentry{module}!qnav.measurement.platform@\spxentry{qnav.measurement.platform}}\index{qnav.measurement.platform@\spxentry{qnav.measurement.platform}!module@\spxentry{module}}

\paragraph{platform.py}
\label{\detokenize{modules/measurement/platform:platform-py}}\begin{quote}\begin{description}
\sphinxlineitem{summary}
\sphinxAtStartPar
Provides conversion to and from sensor reference frame.
This module is responsible for representing the orientation and location
of sensor’s on a platform body, providing means for common conversion
between the reference frames.

\sphinxlineitem{authors}
\begin{DUlineblock}{0em}
\item[] Prof. Jason Ralph \sphinxhyphen{} \sphinxhref{mailto:jfralph@liverpool.ac.uk}{jfralph@liverpool.ac.uk}
\item[] Michael Wright \sphinxhyphen{} \sphinxhref{mailto:mjwright@liverpool.ac.uk}{mjwright@liverpool.ac.uk}
\end{DUlineblock}

\sphinxlineitem{version}
\sphinxAtStartPar
1.0.0 \sphinxhyphen{} First full implementation of module.

\end{description}\end{quote}
\index{SensorAxis (class in qnav.measurement.platform)@\spxentry{SensorAxis}\spxextra{class in qnav.measurement.platform}}

\begin{fulllineitems}
\phantomsection\label{\detokenize{modules/measurement/platform:qnav.measurement.platform.SensorAxis}}
\pysigstartsignatures
\pysiglinewithargsret{\sphinxbfcode{\sphinxupquote{class\DUrole{w}{ }}}\sphinxcode{\sphinxupquote{qnav.measurement.platform.}}\sphinxbfcode{\sphinxupquote{SensorAxis}}}{\sphinxparam{\DUrole{n}{sensor\_angles}\DUrole{p}{:}\DUrole{w}{ }\DUrole{n}{Buffer\DUrole{w}{ }\DUrole{p}{|}\DUrole{w}{ }\_SupportsArray\DUrole{p}{{[}}dtype\DUrole{p}{{[}}Any\DUrole{p}{{]}}\DUrole{p}{{]}}\DUrole{w}{ }\DUrole{p}{|}\DUrole{w}{ }\_NestedSequence\DUrole{p}{{[}}\_SupportsArray\DUrole{p}{{[}}dtype\DUrole{p}{{[}}Any\DUrole{p}{{]}}\DUrole{p}{{]}}\DUrole{p}{{]}}\DUrole{w}{ }\DUrole{p}{|}\DUrole{w}{ }bool\DUrole{w}{ }\DUrole{p}{|}\DUrole{w}{ }int\DUrole{w}{ }\DUrole{p}{|}\DUrole{w}{ }float\DUrole{w}{ }\DUrole{p}{|}\DUrole{w}{ }complex\DUrole{w}{ }\DUrole{p}{|}\DUrole{w}{ }str\DUrole{w}{ }\DUrole{p}{|}\DUrole{w}{ }bytes\DUrole{w}{ }\DUrole{p}{|}\DUrole{w}{ }\_NestedSequence\DUrole{p}{{[}}bool\DUrole{w}{ }\DUrole{p}{|}\DUrole{w}{ }int\DUrole{w}{ }\DUrole{p}{|}\DUrole{w}{ }float\DUrole{w}{ }\DUrole{p}{|}\DUrole{w}{ }complex\DUrole{w}{ }\DUrole{p}{|}\DUrole{w}{ }str\DUrole{w}{ }\DUrole{p}{|}\DUrole{w}{ }bytes\DUrole{p}{{]}}}\DUrole{w}{ }\DUrole{o}{=}\DUrole{w}{ }\DUrole{default_value}{array({[}0., 0., 0.{]})}}\sphinxparamcomma \sphinxparam{\DUrole{n}{lever\_arm}\DUrole{p}{:}\DUrole{w}{ }\DUrole{n}{Buffer\DUrole{w}{ }\DUrole{p}{|}\DUrole{w}{ }\_SupportsArray\DUrole{p}{{[}}dtype\DUrole{p}{{[}}Any\DUrole{p}{{]}}\DUrole{p}{{]}}\DUrole{w}{ }\DUrole{p}{|}\DUrole{w}{ }\_NestedSequence\DUrole{p}{{[}}\_SupportsArray\DUrole{p}{{[}}dtype\DUrole{p}{{[}}Any\DUrole{p}{{]}}\DUrole{p}{{]}}\DUrole{p}{{]}}\DUrole{w}{ }\DUrole{p}{|}\DUrole{w}{ }bool\DUrole{w}{ }\DUrole{p}{|}\DUrole{w}{ }int\DUrole{w}{ }\DUrole{p}{|}\DUrole{w}{ }float\DUrole{w}{ }\DUrole{p}{|}\DUrole{w}{ }complex\DUrole{w}{ }\DUrole{p}{|}\DUrole{w}{ }str\DUrole{w}{ }\DUrole{p}{|}\DUrole{w}{ }bytes\DUrole{w}{ }\DUrole{p}{|}\DUrole{w}{ }\_NestedSequence\DUrole{p}{{[}}bool\DUrole{w}{ }\DUrole{p}{|}\DUrole{w}{ }int\DUrole{w}{ }\DUrole{p}{|}\DUrole{w}{ }float\DUrole{w}{ }\DUrole{p}{|}\DUrole{w}{ }complex\DUrole{w}{ }\DUrole{p}{|}\DUrole{w}{ }str\DUrole{w}{ }\DUrole{p}{|}\DUrole{w}{ }bytes\DUrole{p}{{]}}}\DUrole{w}{ }\DUrole{o}{=}\DUrole{w}{ }\DUrole{default_value}{array({[}0., 0., 0.{]})}}}{}
\pysigstopsignatures
\sphinxAtStartPar
Specifies the orientation and offset of sensors relative to platform.
For a corresponding sensor, this class specifies the angle orientation
and offset location of the sensor on the platform body. This supplies
a consistent practice for local conversion between sensor and
platform reference frames.
\index{sensor\_angles\_deg (qnav.measurement.platform.SensorAxis property)@\spxentry{sensor\_angles\_deg}\spxextra{qnav.measurement.platform.SensorAxis property}}

\begin{fulllineitems}
\phantomsection\label{\detokenize{modules/measurement/platform:qnav.measurement.platform.SensorAxis.sensor_angles_deg}}
\pysigstartsignatures
\pysigline{\sphinxbfcode{\sphinxupquote{property\DUrole{w}{ }}}\sphinxbfcode{\sphinxupquote{sensor\_angles\_deg}}\sphinxbfcode{\sphinxupquote{\DUrole{p}{:}\DUrole{w}{ }ndarray}}}
\pysigstopsignatures
\sphinxAtStartPar
Returns the sensor orientation angles in degrees.
\begin{quote}\begin{description}
\sphinxlineitem{Returns}
\sphinxAtStartPar
Sensor angles (in degrees).

\sphinxlineitem{Return type}
\sphinxAtStartPar
np.ndarray (3\sphinxhyphen{}elements)

\end{description}\end{quote}

\end{fulllineitems}

\index{sensor\_angles\_rad (qnav.measurement.platform.SensorAxis property)@\spxentry{sensor\_angles\_rad}\spxextra{qnav.measurement.platform.SensorAxis property}}

\begin{fulllineitems}
\phantomsection\label{\detokenize{modules/measurement/platform:qnav.measurement.platform.SensorAxis.sensor_angles_rad}}
\pysigstartsignatures
\pysigline{\sphinxbfcode{\sphinxupquote{property\DUrole{w}{ }}}\sphinxbfcode{\sphinxupquote{sensor\_angles\_rad}}\sphinxbfcode{\sphinxupquote{\DUrole{p}{:}\DUrole{w}{ }ndarray}}}
\pysigstopsignatures
\sphinxAtStartPar
Returns the sensor orientation angles in radians.
\begin{quote}\begin{description}
\sphinxlineitem{Returns}
\sphinxAtStartPar
Sensor angles (in radians).

\sphinxlineitem{Return type}
\sphinxAtStartPar
np.ndarray (3\sphinxhyphen{}elements)

\end{description}\end{quote}

\end{fulllineitems}

\index{lever\_arm (qnav.measurement.platform.SensorAxis property)@\spxentry{lever\_arm}\spxextra{qnav.measurement.platform.SensorAxis property}}

\begin{fulllineitems}
\phantomsection\label{\detokenize{modules/measurement/platform:qnav.measurement.platform.SensorAxis.lever_arm}}
\pysigstartsignatures
\pysigline{\sphinxbfcode{\sphinxupquote{property\DUrole{w}{ }}}\sphinxbfcode{\sphinxupquote{lever\_arm}}\sphinxbfcode{\sphinxupquote{\DUrole{p}{:}\DUrole{w}{ }ndarray}}}
\pysigstopsignatures
\sphinxAtStartPar
The lever arm position offset in metres.
\begin{quote}\begin{description}
\sphinxlineitem{Returns}
\sphinxAtStartPar
Lever arm offset (in metres).

\sphinxlineitem{Return type}
\sphinxAtStartPar
np.ndarray (3\sphinxhyphen{}elements)

\end{description}\end{quote}

\end{fulllineitems}

\index{body2sensor\_mat (qnav.measurement.platform.SensorAxis property)@\spxentry{body2sensor\_mat}\spxextra{qnav.measurement.platform.SensorAxis property}}

\begin{fulllineitems}
\phantomsection\label{\detokenize{modules/measurement/platform:qnav.measurement.platform.SensorAxis.body2sensor_mat}}
\pysigstartsignatures
\pysigline{\sphinxbfcode{\sphinxupquote{property\DUrole{w}{ }}}\sphinxbfcode{\sphinxupquote{body2sensor\_mat}}\sphinxbfcode{\sphinxupquote{\DUrole{p}{:}\DUrole{w}{ }ndarray}}}
\pysigstopsignatures
\sphinxAtStartPar
The rotation matrix for body to sensor conversion.
\begin{quote}\begin{description}
\sphinxlineitem{Returns}
\sphinxAtStartPar
The rotation matrix for reference frame conversion.

\sphinxlineitem{Return type}
\sphinxAtStartPar
np.ndarray (3\sphinxhyphen{}by\sphinxhyphen{}3 elements)

\end{description}\end{quote}

\end{fulllineitems}


\end{fulllineitems}


\sphinxstepscope


\subsubsection{qnav.measurement.error\_properties}
\label{\detokenize{modules/measurement/error_properties:module-qnav.measurement.error_properties}}\label{\detokenize{modules/measurement/error_properties:qnav-measurement-error-properties}}\label{\detokenize{modules/measurement/error_properties::doc}}\index{module@\spxentry{module}!qnav.measurement.error\_properties@\spxentry{qnav.measurement.error\_properties}}\index{qnav.measurement.error\_properties@\spxentry{qnav.measurement.error\_properties}!module@\spxentry{module}}

\paragraph{error\_properties.py}
\label{\detokenize{modules/measurement/error_properties:error-properties-py}}\begin{quote}\begin{description}
\sphinxlineitem{summary}
\sphinxAtStartPar
Contains representations for sensor error profiles.
Module contains classes that are used for profiling the common error
types represented by various sensor types. Namely used to reduce
repeated common class arguments and methods.

\sphinxlineitem{authors}
\begin{DUlineblock}{0em}
\item[] Prof. Jason Ralph \sphinxhyphen{} \sphinxhref{mailto:jfralph@liverpool.ac.uk}{jfralph@liverpool.ac.uk}
\item[] Michael Wright \sphinxhyphen{} \sphinxhref{mailto:mjwright@liverpool.ac.uk}{mjwright@liverpool.ac.uk}
\end{DUlineblock}

\sphinxlineitem{version}
\sphinxAtStartPar
1.0.0 \sphinxhyphen{} First full implementation of module.

\end{description}\end{quote}
\index{ErrorProperties (class in qnav.measurement.error\_properties)@\spxentry{ErrorProperties}\spxextra{class in qnav.measurement.error\_properties}}

\begin{fulllineitems}
\phantomsection\label{\detokenize{modules/measurement/error_properties:qnav.measurement.error_properties.ErrorProperties}}
\pysigstartsignatures
\pysiglinewithargsret{\sphinxbfcode{\sphinxupquote{class\DUrole{w}{ }}}\sphinxcode{\sphinxupquote{qnav.measurement.error\_properties.}}\sphinxbfcode{\sphinxupquote{ErrorProperties}}}{\sphinxparam{\DUrole{n}{bias\_error}\DUrole{p}{:}\DUrole{w}{ }\DUrole{n}{float\DUrole{w}{ }\DUrole{p}{|}\DUrole{w}{ }Buffer\DUrole{w}{ }\DUrole{p}{|}\DUrole{w}{ }\_SupportsArray\DUrole{p}{{[}}dtype\DUrole{p}{{[}}Any\DUrole{p}{{]}}\DUrole{p}{{]}}\DUrole{w}{ }\DUrole{p}{|}\DUrole{w}{ }\_NestedSequence\DUrole{p}{{[}}\_SupportsArray\DUrole{p}{{[}}dtype\DUrole{p}{{[}}Any\DUrole{p}{{]}}\DUrole{p}{{]}}\DUrole{p}{{]}}\DUrole{w}{ }\DUrole{p}{|}\DUrole{w}{ }bool\DUrole{w}{ }\DUrole{p}{|}\DUrole{w}{ }int\DUrole{w}{ }\DUrole{p}{|}\DUrole{w}{ }complex\DUrole{w}{ }\DUrole{p}{|}\DUrole{w}{ }str\DUrole{w}{ }\DUrole{p}{|}\DUrole{w}{ }bytes\DUrole{w}{ }\DUrole{p}{|}\DUrole{w}{ }\_NestedSequence\DUrole{p}{{[}}bool\DUrole{w}{ }\DUrole{p}{|}\DUrole{w}{ }int\DUrole{w}{ }\DUrole{p}{|}\DUrole{w}{ }float\DUrole{w}{ }\DUrole{p}{|}\DUrole{w}{ }complex\DUrole{w}{ }\DUrole{p}{|}\DUrole{w}{ }str\DUrole{w}{ }\DUrole{p}{|}\DUrole{w}{ }bytes\DUrole{p}{{]}}}\DUrole{w}{ }\DUrole{o}{=}\DUrole{w}{ }\DUrole{default_value}{0}}\sphinxparamcomma \sphinxparam{\DUrole{n}{bias\_drift\_rate}\DUrole{p}{:}\DUrole{w}{ }\DUrole{n}{float\DUrole{w}{ }\DUrole{p}{|}\DUrole{w}{ }Buffer\DUrole{w}{ }\DUrole{p}{|}\DUrole{w}{ }\_SupportsArray\DUrole{p}{{[}}dtype\DUrole{p}{{[}}Any\DUrole{p}{{]}}\DUrole{p}{{]}}\DUrole{w}{ }\DUrole{p}{|}\DUrole{w}{ }\_NestedSequence\DUrole{p}{{[}}\_SupportsArray\DUrole{p}{{[}}dtype\DUrole{p}{{[}}Any\DUrole{p}{{]}}\DUrole{p}{{]}}\DUrole{p}{{]}}\DUrole{w}{ }\DUrole{p}{|}\DUrole{w}{ }bool\DUrole{w}{ }\DUrole{p}{|}\DUrole{w}{ }int\DUrole{w}{ }\DUrole{p}{|}\DUrole{w}{ }complex\DUrole{w}{ }\DUrole{p}{|}\DUrole{w}{ }str\DUrole{w}{ }\DUrole{p}{|}\DUrole{w}{ }bytes\DUrole{w}{ }\DUrole{p}{|}\DUrole{w}{ }\_NestedSequence\DUrole{p}{{[}}bool\DUrole{w}{ }\DUrole{p}{|}\DUrole{w}{ }int\DUrole{w}{ }\DUrole{p}{|}\DUrole{w}{ }float\DUrole{w}{ }\DUrole{p}{|}\DUrole{w}{ }complex\DUrole{w}{ }\DUrole{p}{|}\DUrole{w}{ }str\DUrole{w}{ }\DUrole{p}{|}\DUrole{w}{ }bytes\DUrole{p}{{]}}}\DUrole{w}{ }\DUrole{o}{=}\DUrole{w}{ }\DUrole{default_value}{0}}\sphinxparamcomma \sphinxparam{\DUrole{n}{scale\_error}\DUrole{p}{:}\DUrole{w}{ }\DUrole{n}{float\DUrole{w}{ }\DUrole{p}{|}\DUrole{w}{ }Buffer\DUrole{w}{ }\DUrole{p}{|}\DUrole{w}{ }\_SupportsArray\DUrole{p}{{[}}dtype\DUrole{p}{{[}}Any\DUrole{p}{{]}}\DUrole{p}{{]}}\DUrole{w}{ }\DUrole{p}{|}\DUrole{w}{ }\_NestedSequence\DUrole{p}{{[}}\_SupportsArray\DUrole{p}{{[}}dtype\DUrole{p}{{[}}Any\DUrole{p}{{]}}\DUrole{p}{{]}}\DUrole{p}{{]}}\DUrole{w}{ }\DUrole{p}{|}\DUrole{w}{ }bool\DUrole{w}{ }\DUrole{p}{|}\DUrole{w}{ }int\DUrole{w}{ }\DUrole{p}{|}\DUrole{w}{ }complex\DUrole{w}{ }\DUrole{p}{|}\DUrole{w}{ }str\DUrole{w}{ }\DUrole{p}{|}\DUrole{w}{ }bytes\DUrole{w}{ }\DUrole{p}{|}\DUrole{w}{ }\_NestedSequence\DUrole{p}{{[}}bool\DUrole{w}{ }\DUrole{p}{|}\DUrole{w}{ }int\DUrole{w}{ }\DUrole{p}{|}\DUrole{w}{ }float\DUrole{w}{ }\DUrole{p}{|}\DUrole{w}{ }complex\DUrole{w}{ }\DUrole{p}{|}\DUrole{w}{ }str\DUrole{w}{ }\DUrole{p}{|}\DUrole{w}{ }bytes\DUrole{p}{{]}}}\DUrole{w}{ }\DUrole{o}{=}\DUrole{w}{ }\DUrole{default_value}{0}}\sphinxparamcomma \sphinxparam{\DUrole{n}{non\_orth\_error}\DUrole{p}{:}\DUrole{w}{ }\DUrole{n}{float\DUrole{w}{ }\DUrole{p}{|}\DUrole{w}{ }Buffer\DUrole{w}{ }\DUrole{p}{|}\DUrole{w}{ }\_SupportsArray\DUrole{p}{{[}}dtype\DUrole{p}{{[}}Any\DUrole{p}{{]}}\DUrole{p}{{]}}\DUrole{w}{ }\DUrole{p}{|}\DUrole{w}{ }\_NestedSequence\DUrole{p}{{[}}\_SupportsArray\DUrole{p}{{[}}dtype\DUrole{p}{{[}}Any\DUrole{p}{{]}}\DUrole{p}{{]}}\DUrole{p}{{]}}\DUrole{w}{ }\DUrole{p}{|}\DUrole{w}{ }bool\DUrole{w}{ }\DUrole{p}{|}\DUrole{w}{ }int\DUrole{w}{ }\DUrole{p}{|}\DUrole{w}{ }complex\DUrole{w}{ }\DUrole{p}{|}\DUrole{w}{ }str\DUrole{w}{ }\DUrole{p}{|}\DUrole{w}{ }bytes\DUrole{w}{ }\DUrole{p}{|}\DUrole{w}{ }\_NestedSequence\DUrole{p}{{[}}bool\DUrole{w}{ }\DUrole{p}{|}\DUrole{w}{ }int\DUrole{w}{ }\DUrole{p}{|}\DUrole{w}{ }float\DUrole{w}{ }\DUrole{p}{|}\DUrole{w}{ }complex\DUrole{w}{ }\DUrole{p}{|}\DUrole{w}{ }str\DUrole{w}{ }\DUrole{p}{|}\DUrole{w}{ }bytes\DUrole{p}{{]}}}\DUrole{w}{ }\DUrole{o}{=}\DUrole{w}{ }\DUrole{default_value}{0}}\sphinxparamcomma \sphinxparam{\DUrole{n}{avg\_meas\_noise}\DUrole{p}{:}\DUrole{w}{ }\DUrole{n}{float\DUrole{w}{ }\DUrole{p}{|}\DUrole{w}{ }Buffer\DUrole{w}{ }\DUrole{p}{|}\DUrole{w}{ }\_SupportsArray\DUrole{p}{{[}}dtype\DUrole{p}{{[}}Any\DUrole{p}{{]}}\DUrole{p}{{]}}\DUrole{w}{ }\DUrole{p}{|}\DUrole{w}{ }\_NestedSequence\DUrole{p}{{[}}\_SupportsArray\DUrole{p}{{[}}dtype\DUrole{p}{{[}}Any\DUrole{p}{{]}}\DUrole{p}{{]}}\DUrole{p}{{]}}\DUrole{w}{ }\DUrole{p}{|}\DUrole{w}{ }bool\DUrole{w}{ }\DUrole{p}{|}\DUrole{w}{ }int\DUrole{w}{ }\DUrole{p}{|}\DUrole{w}{ }complex\DUrole{w}{ }\DUrole{p}{|}\DUrole{w}{ }str\DUrole{w}{ }\DUrole{p}{|}\DUrole{w}{ }bytes\DUrole{w}{ }\DUrole{p}{|}\DUrole{w}{ }\_NestedSequence\DUrole{p}{{[}}bool\DUrole{w}{ }\DUrole{p}{|}\DUrole{w}{ }int\DUrole{w}{ }\DUrole{p}{|}\DUrole{w}{ }float\DUrole{w}{ }\DUrole{p}{|}\DUrole{w}{ }complex\DUrole{w}{ }\DUrole{p}{|}\DUrole{w}{ }str\DUrole{w}{ }\DUrole{p}{|}\DUrole{w}{ }bytes\DUrole{p}{{]}}}\DUrole{w}{ }\DUrole{o}{=}\DUrole{w}{ }\DUrole{default_value}{0}}}{}
\pysigstopsignatures
\sphinxAtStartPar
Defines common error properties used for 3D sensors.
This data class contains the errors associated with 3D sensors;
specifically, the bias error, bias drift rate, scaling errors,
non\sphinxhyphen{}orthogonality and measurement noise properties. These are commonly
used when defining sensors such as accelerometer and gyroscope, along
with their quantum counterparts.
\index{bias\_error (qnav.measurement.error\_properties.ErrorProperties property)@\spxentry{bias\_error}\spxextra{qnav.measurement.error\_properties.ErrorProperties property}}

\begin{fulllineitems}
\phantomsection\label{\detokenize{modules/measurement/error_properties:qnav.measurement.error_properties.ErrorProperties.bias_error}}
\pysigstartsignatures
\pysigline{\sphinxbfcode{\sphinxupquote{property\DUrole{w}{ }}}\sphinxbfcode{\sphinxupquote{bias\_error}}\sphinxbfcode{\sphinxupquote{\DUrole{p}{:}\DUrole{w}{ }ndarray}}}
\pysigstopsignatures
\sphinxAtStartPar
The bias errors for each axis (in micro\sphinxhyphen{}g).
\begin{quote}\begin{description}
\sphinxlineitem{Returns}
\sphinxAtStartPar
Sensor bias errors (in micro\sphinxhyphen{}g).

\sphinxlineitem{Return type}
\sphinxAtStartPar
np.ndarray (3\sphinxhyphen{}elements)

\end{description}\end{quote}

\end{fulllineitems}

\index{bias\_drift\_rate (qnav.measurement.error\_properties.ErrorProperties property)@\spxentry{bias\_drift\_rate}\spxextra{qnav.measurement.error\_properties.ErrorProperties property}}

\begin{fulllineitems}
\phantomsection\label{\detokenize{modules/measurement/error_properties:qnav.measurement.error_properties.ErrorProperties.bias_drift_rate}}
\pysigstartsignatures
\pysigline{\sphinxbfcode{\sphinxupquote{property\DUrole{w}{ }}}\sphinxbfcode{\sphinxupquote{bias\_drift\_rate}}\sphinxbfcode{\sphinxupquote{\DUrole{p}{:}\DUrole{w}{ }ndarray}}}
\pysigstopsignatures
\sphinxAtStartPar
The bias drift rate for each axis (in micro\sphinxhyphen{}g/root sec).
\begin{quote}\begin{description}
\sphinxlineitem{Returns}
\sphinxAtStartPar
Bias drift rate (in micro\sphinxhyphen{}g/root sec).

\sphinxlineitem{Return type}
\sphinxAtStartPar
np.ndarray (3\sphinxhyphen{}elements)

\end{description}\end{quote}

\end{fulllineitems}

\index{scale\_error (qnav.measurement.error\_properties.ErrorProperties property)@\spxentry{scale\_error}\spxextra{qnav.measurement.error\_properties.ErrorProperties property}}

\begin{fulllineitems}
\phantomsection\label{\detokenize{modules/measurement/error_properties:qnav.measurement.error_properties.ErrorProperties.scale_error}}
\pysigstartsignatures
\pysigline{\sphinxbfcode{\sphinxupquote{property\DUrole{w}{ }}}\sphinxbfcode{\sphinxupquote{scale\_error}}\sphinxbfcode{\sphinxupquote{\DUrole{p}{:}\DUrole{w}{ }ndarray}}}
\pysigstopsignatures
\sphinxAtStartPar
The scaling errors for each axis (in ppm).
\begin{quote}\begin{description}
\sphinxlineitem{Returns}
\sphinxAtStartPar
Scaling errors (in ppm).

\sphinxlineitem{Return type}
\sphinxAtStartPar
np.ndarray (3\sphinxhyphen{}elements)

\end{description}\end{quote}

\end{fulllineitems}

\index{non\_orth\_error (qnav.measurement.error\_properties.ErrorProperties property)@\spxentry{non\_orth\_error}\spxextra{qnav.measurement.error\_properties.ErrorProperties property}}

\begin{fulllineitems}
\phantomsection\label{\detokenize{modules/measurement/error_properties:qnav.measurement.error_properties.ErrorProperties.non_orth_error}}
\pysigstartsignatures
\pysigline{\sphinxbfcode{\sphinxupquote{property\DUrole{w}{ }}}\sphinxbfcode{\sphinxupquote{non\_orth\_error}}\sphinxbfcode{\sphinxupquote{\DUrole{p}{:}\DUrole{w}{ }ndarray}}}
\pysigstopsignatures
\sphinxAtStartPar
The non\sphinxhyphen{}orthogonal “cross\sphinxhyphen{}coupling” errors (in micro\sphinxhyphen{}g).
Specifically the alignment errors for axis xy, xz, yx, yz,
zx and zy respectively.
\begin{quote}\begin{description}
\sphinxlineitem{Returns}
\sphinxAtStartPar
non\sphinxhyphen{}orthogonal errors (in micro\sphinxhyphen{}g).

\sphinxlineitem{Return type}
\sphinxAtStartPar
np.ndarray (6\sphinxhyphen{}elements)

\end{description}\end{quote}

\end{fulllineitems}

\index{avg\_meas\_noise (qnav.measurement.error\_properties.ErrorProperties property)@\spxentry{avg\_meas\_noise}\spxextra{qnav.measurement.error\_properties.ErrorProperties property}}

\begin{fulllineitems}
\phantomsection\label{\detokenize{modules/measurement/error_properties:qnav.measurement.error_properties.ErrorProperties.avg_meas_noise}}
\pysigstartsignatures
\pysigline{\sphinxbfcode{\sphinxupquote{property\DUrole{w}{ }}}\sphinxbfcode{\sphinxupquote{avg\_meas\_noise}}\sphinxbfcode{\sphinxupquote{\DUrole{p}{:}\DUrole{w}{ }ndarray}}}
\pysigstopsignatures
\sphinxAtStartPar
The average measurement error for each axis (in micro\sphinxhyphen{}g.root Hz).
\begin{quote}\begin{description}
\sphinxlineitem{Returns}
\sphinxAtStartPar
average measurement error (in micro\sphinxhyphen{}g.root Hz).

\sphinxlineitem{Return type}
\sphinxAtStartPar
np.ndarray (3\sphinxhyphen{}elements)

\end{description}\end{quote}

\end{fulllineitems}

\index{error\_matrix (qnav.measurement.error\_properties.ErrorProperties property)@\spxentry{error\_matrix}\spxextra{qnav.measurement.error\_properties.ErrorProperties property}}

\begin{fulllineitems}
\phantomsection\label{\detokenize{modules/measurement/error_properties:qnav.measurement.error_properties.ErrorProperties.error_matrix}}
\pysigstartsignatures
\pysigline{\sphinxbfcode{\sphinxupquote{property\DUrole{w}{ }}}\sphinxbfcode{\sphinxupquote{error\_matrix}}\sphinxbfcode{\sphinxupquote{\DUrole{p}{:}\DUrole{w}{ }ndarray}}}
\pysigstopsignatures
\sphinxAtStartPar
A matrix containing the scale factor and cross\sphinxhyphen{}coupling errors for a
nominally orthogonal sensor. This is a commonly used representation.
For scale errors S and non\sphinxhyphen{}orthogonal errors M, this is formatted as:
{[}{[} Sx, Mxy, Mxz {]},
\begin{quote}

\sphinxAtStartPar
{[} Myx, Sy, Myz {]},
{[} Mzx, Mzy, Sz {]}{]}
\end{quote}
\begin{quote}\begin{description}
\sphinxlineitem{Returns}
\sphinxAtStartPar
Error matrix for scale factor and cross\sphinxhyphen{}coupling errors.

\sphinxlineitem{Return type}
\sphinxAtStartPar
np.ndarray (3\sphinxhyphen{}by\sphinxhyphen{}3 elements)

\end{description}\end{quote}

\end{fulllineitems}


\end{fulllineitems}


\sphinxstepscope


\subsubsection{qnav.measurement.accelerometer}
\label{\detokenize{modules/measurement/accelerometer:module-qnav.measurement.accelerometer}}\label{\detokenize{modules/measurement/accelerometer:qnav-measurement-accelerometer}}\label{\detokenize{modules/measurement/accelerometer::doc}}\index{module@\spxentry{module}!qnav.measurement.accelerometer@\spxentry{qnav.measurement.accelerometer}}\index{qnav.measurement.accelerometer@\spxentry{qnav.measurement.accelerometer}!module@\spxentry{module}}

\paragraph{accelerometer.py}
\label{\detokenize{modules/measurement/accelerometer:accelerometer-py}}\begin{quote}\begin{description}
\sphinxlineitem{summary}
\sphinxAtStartPar
Contains the virtual accelerometer classes for producing measurements.
This module contains virtual accelerometer classes that can be used for
generating accelerometer measurements, from a simulated environment,
with a configured level of noise/errors.

\sphinxlineitem{authors}
\begin{DUlineblock}{0em}
\item[] Prof. Jason Ralph \sphinxhyphen{} \sphinxhref{mailto:jfralph@liverpool.ac.uk}{jfralph@liverpool.ac.uk}
\item[] Michael Wright \sphinxhyphen{} \sphinxhref{mailto:mjwright@liverpool.ac.uk}{mjwright@liverpool.ac.uk}
\end{DUlineblock}

\sphinxlineitem{version}
\sphinxAtStartPar
1.0.0 \sphinxhyphen{} First full implementation of module.

\end{description}\end{quote}
\index{Accelerometer (class in qnav.measurement.accelerometer)@\spxentry{Accelerometer}\spxextra{class in qnav.measurement.accelerometer}}

\begin{fulllineitems}
\phantomsection\label{\detokenize{modules/measurement/accelerometer:qnav.measurement.accelerometer.Accelerometer}}
\pysigstartsignatures
\pysiglinewithargsret{\sphinxbfcode{\sphinxupquote{class\DUrole{w}{ }}}\sphinxcode{\sphinxupquote{qnav.measurement.accelerometer.}}\sphinxbfcode{\sphinxupquote{Accelerometer}}}{\sphinxparam{\DUrole{n}{frequency: float}}\sphinxparamcomma \sphinxparam{\DUrole{n}{error\_profile: \textasciitilde{}qnav.measurement.error\_properties.ErrorProperties}}\sphinxparamcomma \sphinxparam{\DUrole{n}{sensor\_axis: \textasciitilde{}qnav.measurement.platform.SensorAxis = \textless{}qnav.measurement.platform.SensorAxis object\textgreater{}}}\sphinxparamcomma \sphinxparam{\DUrole{n}{start\_time: float = 0.0}}\sphinxparamcomma \sphinxparam{\DUrole{n}{rand\_seed: int = None}}}{}
\pysigstopsignatures
\sphinxAtStartPar
Provides the functionality of an accelerometer in a simulated environment.
This class essentially represents a virtual accelerometer converting
waypoint accelerations to body rate axis and applying noise/errors.
The output of this can then be passed to a selected INS fusion class.
\index{take\_measurement() (qnav.measurement.accelerometer.Accelerometer method)@\spxentry{take\_measurement()}\spxextra{qnav.measurement.accelerometer.Accelerometer method}}

\begin{fulllineitems}
\phantomsection\label{\detokenize{modules/measurement/accelerometer:qnav.measurement.accelerometer.Accelerometer.take_measurement}}
\pysigstartsignatures
\pysiglinewithargsret{\sphinxbfcode{\sphinxupquote{take\_measurement}}}{\sphinxparam{\DUrole{n}{est\_time}\DUrole{p}{:}\DUrole{w}{ }\DUrole{n}{float}}\sphinxparamcomma \sphinxparam{\DUrole{n}{ground\_truth}\DUrole{p}{:}\DUrole{w}{ }\DUrole{n}{{\hyperref[\detokenize{modules/waypoints/trajectory:qnav.waypoints.trajectory.GroundTruth}]{\sphinxcrossref{GroundTruth}}}}}}{}
\pysigstopsignatures
\sphinxAtStartPar
Produces an accelerometer measurement from given truth data.
This method converts the given waypoint data to body axes and
applies the configured errors/noise. The output can then be
fed to the selected INS solution.
\begin{quote}\begin{description}
\sphinxlineitem{Parameters}\begin{itemize}
\item {} 
\sphinxAtStartPar
\sphinxstyleliteralstrong{\sphinxupquote{est\_time}} (\sphinxstyleliteralemphasis{\sphinxupquote{float}}) \textendash{} The estimated time the measurement of the sensor
is expected to be captured at (in seconds).

\item {} 
\sphinxAtStartPar
\sphinxstyleliteralstrong{\sphinxupquote{ground\_truth}} ({\hyperref[\detokenize{modules/waypoints/trajectory:qnav.waypoints.trajectory.GroundTruth}]{\sphinxcrossref{\sphinxstyleliteralemphasis{\sphinxupquote{GroundTruth}}}}}) \textendash{} The corresponding ground truth data.

\end{itemize}

\end{description}\end{quote}

\end{fulllineitems}

\index{last\_measurement (qnav.measurement.accelerometer.Accelerometer property)@\spxentry{last\_measurement}\spxextra{qnav.measurement.accelerometer.Accelerometer property}}

\begin{fulllineitems}
\phantomsection\label{\detokenize{modules/measurement/accelerometer:qnav.measurement.accelerometer.Accelerometer.last_measurement}}
\pysigstartsignatures
\pysigline{\sphinxbfcode{\sphinxupquote{property\DUrole{w}{ }}}\sphinxbfcode{\sphinxupquote{last\_measurement}}\sphinxbfcode{\sphinxupquote{\DUrole{p}{:}\DUrole{w}{ }ndarray}}}
\pysigstopsignatures
\sphinxAtStartPar
Returns the latest recorded acceleration produced by sensor.
The measured acceleration of the accelerometer (North, East
and Down components in metres/second).
\begin{quote}\begin{description}
\sphinxlineitem{Returns}
\sphinxAtStartPar
Latest acceleration measurement produced by sensor.

\sphinxlineitem{Return type}
\sphinxAtStartPar
np.ndarray (3\sphinxhyphen{}elements)

\end{description}\end{quote}

\end{fulllineitems}

\index{update() (qnav.measurement.accelerometer.Accelerometer method)@\spxentry{update()}\spxextra{qnav.measurement.accelerometer.Accelerometer method}}

\begin{fulllineitems}
\phantomsection\label{\detokenize{modules/measurement/accelerometer:qnav.measurement.accelerometer.Accelerometer.update}}
\pysigstartsignatures
\pysiglinewithargsret{\sphinxbfcode{\sphinxupquote{update}}}{\sphinxparam{\DUrole{n}{time\_sec}\DUrole{p}{:}\DUrole{w}{ }\DUrole{n}{float}\DUrole{w}{ }\DUrole{o}{=}\DUrole{w}{ }\DUrole{default_value}{None}}}{}
\pysigstopsignatures
\sphinxAtStartPar
Called every iteration to update simulated sensor errors.
As sensor errors propagate with time, it is important that this
method is called prior to every measurement update.
\begin{quote}\begin{description}
\sphinxlineitem{Parameters}
\sphinxAtStartPar
\sphinxstyleliteralstrong{\sphinxupquote{time\_sec}} (\sphinxstyleliteralemphasis{\sphinxupquote{float}}) \textendash{} The current simulation time in seconds.

\end{description}\end{quote}

\end{fulllineitems}

\index{sensor\_axis (qnav.measurement.accelerometer.Accelerometer property)@\spxentry{sensor\_axis}\spxextra{qnav.measurement.accelerometer.Accelerometer property}}

\begin{fulllineitems}
\phantomsection\label{\detokenize{modules/measurement/accelerometer:qnav.measurement.accelerometer.Accelerometer.sensor_axis}}
\pysigstartsignatures
\pysigline{\sphinxbfcode{\sphinxupquote{property\DUrole{w}{ }}}\sphinxbfcode{\sphinxupquote{sensor\_axis}}\sphinxbfcode{\sphinxupquote{\DUrole{p}{:}\DUrole{w}{ }{\hyperref[\detokenize{modules/measurement/platform:qnav.measurement.platform.SensorAxis}]{\sphinxcrossref{SensorAxis}}}}}}
\pysigstopsignatures
\sphinxAtStartPar
Returns the true axis for the accelerometer.
\begin{quote}\begin{description}
\sphinxlineitem{Returns}
\sphinxAtStartPar
The sensor axis of the accelerometer.

\sphinxlineitem{Return type}
\sphinxAtStartPar
{\hyperref[\detokenize{modules/measurement/platform:qnav.measurement.platform.SensorAxis}]{\sphinxcrossref{SensorAxis}}}

\end{description}\end{quote}

\end{fulllineitems}


\end{fulllineitems}


\sphinxstepscope


\subsubsection{qnav.measurement.gyroscope}
\label{\detokenize{modules/measurement/gyroscope:module-qnav.measurement.gyroscope}}\label{\detokenize{modules/measurement/gyroscope:qnav-measurement-gyroscope}}\label{\detokenize{modules/measurement/gyroscope::doc}}\index{module@\spxentry{module}!qnav.measurement.gyroscope@\spxentry{qnav.measurement.gyroscope}}\index{qnav.measurement.gyroscope@\spxentry{qnav.measurement.gyroscope}!module@\spxentry{module}}

\paragraph{gyroscope.py}
\label{\detokenize{modules/measurement/gyroscope:gyroscope-py}}\begin{quote}\begin{description}
\sphinxlineitem{summary}
\sphinxAtStartPar
Contains the virtual gyroscope classes for producing measurements.
This module contains virtual gyroscope classes that can be used for
generating gyroscope measurements, from a simulated environment,
with a configured level of noise/errors.

\sphinxlineitem{authors}
\begin{DUlineblock}{0em}
\item[] Prof. Jason Ralph \sphinxhyphen{} \sphinxhref{mailto:jfralph@liverpool.ac.uk}{jfralph@liverpool.ac.uk}
\item[] Michael Wright \sphinxhyphen{} \sphinxhref{mailto:mjwright@liverpool.ac.uk}{mjwright@liverpool.ac.uk}
\end{DUlineblock}

\sphinxlineitem{version}
\sphinxAtStartPar
1.0.0 \sphinxhyphen{} First full implementation of module.

\end{description}\end{quote}
\index{Gyroscope (class in qnav.measurement.gyroscope)@\spxentry{Gyroscope}\spxextra{class in qnav.measurement.gyroscope}}

\begin{fulllineitems}
\phantomsection\label{\detokenize{modules/measurement/gyroscope:qnav.measurement.gyroscope.Gyroscope}}
\pysigstartsignatures
\pysiglinewithargsret{\sphinxbfcode{\sphinxupquote{class\DUrole{w}{ }}}\sphinxcode{\sphinxupquote{qnav.measurement.gyroscope.}}\sphinxbfcode{\sphinxupquote{Gyroscope}}}{\sphinxparam{\DUrole{n}{frequency: float}}\sphinxparamcomma \sphinxparam{\DUrole{n}{error\_profile: \textasciitilde{}qnav.measurement.error\_properties.ErrorProperties}}\sphinxparamcomma \sphinxparam{\DUrole{n}{sensor\_axis: \textasciitilde{}qnav.measurement.platform.SensorAxis = \textless{}qnav.measurement.platform.SensorAxis object\textgreater{}}}\sphinxparamcomma \sphinxparam{\DUrole{n}{start\_time: float = 0.0}}\sphinxparamcomma \sphinxparam{\DUrole{n}{rand\_seed: int = None}}}{}
\pysigstopsignatures
\sphinxAtStartPar
Provides the functionality of a gyroscope in a simulated environment.
This class essentially represents a virtual gyroscope converting
waypoint angle rates to body rate axis and applying noise/errors.
The output of this can then be passed to a selected INS solution.
\index{take\_measurement() (qnav.measurement.gyroscope.Gyroscope method)@\spxentry{take\_measurement()}\spxextra{qnav.measurement.gyroscope.Gyroscope method}}

\begin{fulllineitems}
\phantomsection\label{\detokenize{modules/measurement/gyroscope:qnav.measurement.gyroscope.Gyroscope.take_measurement}}
\pysigstartsignatures
\pysiglinewithargsret{\sphinxbfcode{\sphinxupquote{take\_measurement}}}{\sphinxparam{\DUrole{n}{est\_time}\DUrole{p}{:}\DUrole{w}{ }\DUrole{n}{float}}\sphinxparamcomma \sphinxparam{\DUrole{n}{ground\_truth}\DUrole{p}{:}\DUrole{w}{ }\DUrole{n}{{\hyperref[\detokenize{modules/waypoints/trajectory:qnav.waypoints.trajectory.GroundTruth}]{\sphinxcrossref{GroundTruth}}}}}}{}
\pysigstopsignatures
\sphinxAtStartPar
Produces a gyroscope measurement from given truth data.
This method converts the given waypoint data to body axes and
applies the configured errors/noise. The output can then be
fed to the selected INS solution.
\begin{quote}\begin{description}
\sphinxlineitem{Parameters}\begin{itemize}
\item {} 
\sphinxAtStartPar
\sphinxstyleliteralstrong{\sphinxupquote{est\_time}} (\sphinxstyleliteralemphasis{\sphinxupquote{float}}) \textendash{} The estimated time the measurement of the sensor
is expected to be captured at (in seconds).

\item {} 
\sphinxAtStartPar
\sphinxstyleliteralstrong{\sphinxupquote{ground\_truth}} ({\hyperref[\detokenize{modules/waypoints/trajectory:qnav.waypoints.trajectory.GroundTruth}]{\sphinxcrossref{\sphinxstyleliteralemphasis{\sphinxupquote{GroundTruth}}}}}) \textendash{} The corresponding ground truth data.

\end{itemize}

\end{description}\end{quote}

\end{fulllineitems}

\index{last\_measurement (qnav.measurement.gyroscope.Gyroscope property)@\spxentry{last\_measurement}\spxextra{qnav.measurement.gyroscope.Gyroscope property}}

\begin{fulllineitems}
\phantomsection\label{\detokenize{modules/measurement/gyroscope:qnav.measurement.gyroscope.Gyroscope.last_measurement}}
\pysigstartsignatures
\pysigline{\sphinxbfcode{\sphinxupquote{property\DUrole{w}{ }}}\sphinxbfcode{\sphinxupquote{last\_measurement}}\sphinxbfcode{\sphinxupquote{\DUrole{p}{:}\DUrole{w}{ }ndarray}}}
\pysigstopsignatures
\sphinxAtStartPar
Returns the latest recorded gyroscope produced by sensor.
The measured angle rates of the gyroscope (P, Q and R
angles in degrees/second).
\begin{quote}\begin{description}
\sphinxlineitem{Returns}
\sphinxAtStartPar
Latest angle rates measurement produced by sensor.

\sphinxlineitem{Return type}
\sphinxAtStartPar
np.ndarray (3\sphinxhyphen{}elements)

\end{description}\end{quote}

\end{fulllineitems}

\index{update() (qnav.measurement.gyroscope.Gyroscope method)@\spxentry{update()}\spxextra{qnav.measurement.gyroscope.Gyroscope method}}

\begin{fulllineitems}
\phantomsection\label{\detokenize{modules/measurement/gyroscope:qnav.measurement.gyroscope.Gyroscope.update}}
\pysigstartsignatures
\pysiglinewithargsret{\sphinxbfcode{\sphinxupquote{update}}}{\sphinxparam{\DUrole{n}{time\_sec}\DUrole{p}{:}\DUrole{w}{ }\DUrole{n}{float}\DUrole{w}{ }\DUrole{o}{=}\DUrole{w}{ }\DUrole{default_value}{None}}}{}
\pysigstopsignatures
\sphinxAtStartPar
Called every iteration to update simulated sensor errors.
As sensor errors propagate with time, it is important that this
method is called prior to every measurement update.
\begin{quote}\begin{description}
\sphinxlineitem{Parameters}
\sphinxAtStartPar
\sphinxstyleliteralstrong{\sphinxupquote{time\_sec}} (\sphinxstyleliteralemphasis{\sphinxupquote{float}}) \textendash{} The current simulation time in seconds.

\end{description}\end{quote}

\end{fulllineitems}

\index{sensor\_axis (qnav.measurement.gyroscope.Gyroscope property)@\spxentry{sensor\_axis}\spxextra{qnav.measurement.gyroscope.Gyroscope property}}

\begin{fulllineitems}
\phantomsection\label{\detokenize{modules/measurement/gyroscope:qnav.measurement.gyroscope.Gyroscope.sensor_axis}}
\pysigstartsignatures
\pysigline{\sphinxbfcode{\sphinxupquote{property\DUrole{w}{ }}}\sphinxbfcode{\sphinxupquote{sensor\_axis}}\sphinxbfcode{\sphinxupquote{\DUrole{p}{:}\DUrole{w}{ }{\hyperref[\detokenize{modules/measurement/platform:qnav.measurement.platform.SensorAxis}]{\sphinxcrossref{SensorAxis}}}}}}
\pysigstopsignatures
\sphinxAtStartPar
Returns the true axis for the gyroscope.
\begin{quote}\begin{description}
\sphinxlineitem{Returns}
\sphinxAtStartPar
The sensor axis of the gyroscope.

\sphinxlineitem{Return type}
\sphinxAtStartPar
{\hyperref[\detokenize{modules/measurement/platform:qnav.measurement.platform.SensorAxis}]{\sphinxcrossref{SensorAxis}}}

\end{description}\end{quote}

\end{fulllineitems}


\end{fulllineitems}


\sphinxstepscope


\subsubsection{qnav.measurement.altimeter}
\label{\detokenize{modules/measurement/altimeter:module-qnav.measurement.altimeter}}\label{\detokenize{modules/measurement/altimeter:qnav-measurement-altimeter}}\label{\detokenize{modules/measurement/altimeter::doc}}\index{module@\spxentry{module}!qnav.measurement.altimeter@\spxentry{qnav.measurement.altimeter}}\index{qnav.measurement.altimeter@\spxentry{qnav.measurement.altimeter}!module@\spxentry{module}}

\paragraph{altimeter.py}
\label{\detokenize{modules/measurement/altimeter:altimeter-py}}\begin{quote}\begin{description}
\sphinxlineitem{summary}
\sphinxAtStartPar
Contains the virtual altimeter classes for producing measurements.
This module contains virtual altimeter classes that can be used for
generating altimeter measurements, from a simulated environment,
with a configured level of noise/errors.

\sphinxlineitem{authors}
\begin{DUlineblock}{0em}
\item[] Prof. Jason Ralph \sphinxhyphen{} \sphinxhref{mailto:jfralph@liverpool.ac.uk}{jfralph@liverpool.ac.uk}
\item[] Michael Wright \sphinxhyphen{} \sphinxhref{mailto:mjwright@liverpool.ac.uk}{mjwright@liverpool.ac.uk}
\end{DUlineblock}

\sphinxlineitem{version}
\sphinxAtStartPar
1.0.0 \sphinxhyphen{} First full implementation of module.

\end{description}\end{quote}
\index{Altimeter (class in qnav.measurement.altimeter)@\spxentry{Altimeter}\spxextra{class in qnav.measurement.altimeter}}

\begin{fulllineitems}
\phantomsection\label{\detokenize{modules/measurement/altimeter:qnav.measurement.altimeter.Altimeter}}
\pysigstartsignatures
\pysiglinewithargsret{\sphinxbfcode{\sphinxupquote{class\DUrole{w}{ }}}\sphinxcode{\sphinxupquote{qnav.measurement.altimeter.}}\sphinxbfcode{\sphinxupquote{Altimeter}}}{\sphinxparam{\DUrole{n}{frequency}\DUrole{p}{:}\DUrole{w}{ }\DUrole{n}{float}}\sphinxparamcomma \sphinxparam{\DUrole{n}{bias\_error}\DUrole{p}{:}\DUrole{w}{ }\DUrole{n}{float}\DUrole{w}{ }\DUrole{o}{=}\DUrole{w}{ }\DUrole{default_value}{0.0}}\sphinxparamcomma \sphinxparam{\DUrole{n}{bias\_drift\_rate}\DUrole{p}{:}\DUrole{w}{ }\DUrole{n}{float}\DUrole{w}{ }\DUrole{o}{=}\DUrole{w}{ }\DUrole{default_value}{0.0}}\sphinxparamcomma \sphinxparam{\DUrole{n}{start\_time}\DUrole{p}{:}\DUrole{w}{ }\DUrole{n}{float}\DUrole{w}{ }\DUrole{o}{=}\DUrole{w}{ }\DUrole{default_value}{0.0}}\sphinxparamcomma \sphinxparam{\DUrole{n}{rand\_seed}\DUrole{p}{:}\DUrole{w}{ }\DUrole{n}{int}\DUrole{w}{ }\DUrole{o}{=}\DUrole{w}{ }\DUrole{default_value}{None}}}{}
\pysigstopsignatures
\sphinxAtStartPar
Provides the functionality of an altimeter in a simulated environment.
This class essentially represents a virtual altimeter that takes the
altitude from the waypoint data and adds a configured noise/error.
The output of this can then be passed to a selected INS solution.
\index{take\_measurement() (qnav.measurement.altimeter.Altimeter method)@\spxentry{take\_measurement()}\spxextra{qnav.measurement.altimeter.Altimeter method}}

\begin{fulllineitems}
\phantomsection\label{\detokenize{modules/measurement/altimeter:qnav.measurement.altimeter.Altimeter.take_measurement}}
\pysigstartsignatures
\pysiglinewithargsret{\sphinxbfcode{\sphinxupquote{take\_measurement}}}{\sphinxparam{\DUrole{n}{est\_time}\DUrole{p}{:}\DUrole{w}{ }\DUrole{n}{float}}\sphinxparamcomma \sphinxparam{\DUrole{n}{ground\_truth}\DUrole{p}{:}\DUrole{w}{ }\DUrole{n}{{\hyperref[\detokenize{modules/waypoints/trajectory:qnav.waypoints.trajectory.GroundTruth}]{\sphinxcrossref{GroundTruth}}}}}}{}
\pysigstopsignatures
\sphinxAtStartPar
Produces an altimeter measurement from given truth data.
This method takes the true altitude (above sea level) and
applies the configured errors/noise. The output can then be
fed to the selected INS solution.
\begin{quote}\begin{description}
\sphinxlineitem{Parameters}\begin{itemize}
\item {} 
\sphinxAtStartPar
\sphinxstyleliteralstrong{\sphinxupquote{est\_time}} (\sphinxstyleliteralemphasis{\sphinxupquote{float}}) \textendash{} The estimated time the measurement of the sensor
is expected to be captured at (in seconds).

\item {} 
\sphinxAtStartPar
\sphinxstyleliteralstrong{\sphinxupquote{ground\_truth}} ({\hyperref[\detokenize{modules/waypoints/trajectory:qnav.waypoints.trajectory.GroundTruth}]{\sphinxcrossref{\sphinxstyleliteralemphasis{\sphinxupquote{GroundTruth}}}}}) \textendash{} The corresponding ground truth data.

\end{itemize}

\end{description}\end{quote}

\end{fulllineitems}

\index{last\_measurement (qnav.measurement.altimeter.Altimeter property)@\spxentry{last\_measurement}\spxextra{qnav.measurement.altimeter.Altimeter property}}

\begin{fulllineitems}
\phantomsection\label{\detokenize{modules/measurement/altimeter:qnav.measurement.altimeter.Altimeter.last_measurement}}
\pysigstartsignatures
\pysigline{\sphinxbfcode{\sphinxupquote{property\DUrole{w}{ }}}\sphinxbfcode{\sphinxupquote{last\_measurement}}\sphinxbfcode{\sphinxupquote{\DUrole{p}{:}\DUrole{w}{ }float}}}
\pysigstopsignatures
\sphinxAtStartPar
Returns the latest recorded altitude produced by the altimeter.
\begin{quote}\begin{description}
\sphinxlineitem{Returns}
\sphinxAtStartPar
Latest altitude measurement produced.

\sphinxlineitem{Return type}
\sphinxAtStartPar
float

\end{description}\end{quote}

\end{fulllineitems}

\index{update() (qnav.measurement.altimeter.Altimeter method)@\spxentry{update()}\spxextra{qnav.measurement.altimeter.Altimeter method}}

\begin{fulllineitems}
\phantomsection\label{\detokenize{modules/measurement/altimeter:qnav.measurement.altimeter.Altimeter.update}}
\pysigstartsignatures
\pysiglinewithargsret{\sphinxbfcode{\sphinxupquote{update}}}{\sphinxparam{\DUrole{n}{time\_sec}\DUrole{p}{:}\DUrole{w}{ }\DUrole{n}{float}\DUrole{w}{ }\DUrole{o}{=}\DUrole{w}{ }\DUrole{default_value}{None}}}{}
\pysigstopsignatures
\sphinxAtStartPar
Called on demand to update the altimeter’s state.
Typically called after measurements have been taken. When called
this increases the next measurement time and propagates internal
states, to the time difference.
\begin{quote}\begin{description}
\sphinxlineitem{Parameters}
\sphinxAtStartPar
\sphinxstyleliteralstrong{\sphinxupquote{time\_sec}} (\sphinxstyleliteralemphasis{\sphinxupquote{float}}) \textendash{} The current simulation time in seconds.

\end{description}\end{quote}

\end{fulllineitems}


\end{fulllineitems}


\sphinxstepscope


\subsubsection{qnav.measurement.clock}
\label{\detokenize{modules/measurement/clock:module-qnav.measurement.clock}}\label{\detokenize{modules/measurement/clock:qnav-measurement-clock}}\label{\detokenize{modules/measurement/clock::doc}}\index{module@\spxentry{module}!qnav.measurement.clock@\spxentry{qnav.measurement.clock}}\index{qnav.measurement.clock@\spxentry{qnav.measurement.clock}!module@\spxentry{module}}

\paragraph{clock.py}
\label{\detokenize{modules/measurement/clock:clock-py}}\begin{quote}\begin{description}
\sphinxlineitem{summary}
\sphinxAtStartPar
Contains the virtual clock/timer classes for estimating elapsed time.
This module contains virtual clock classes that can be used for simulating
recorded time (with errors and drift) within a simulated environment. This
is namely to simulate clock errors within simulations.

\sphinxlineitem{authors}
\begin{DUlineblock}{0em}
\item[] Prof. Jason Ralph \sphinxhyphen{} \sphinxhref{mailto:jfralph@liverpool.ac.uk}{jfralph@liverpool.ac.uk}
\item[] Michael Wright \sphinxhyphen{} \sphinxhref{mailto:mjwright@liverpool.ac.uk}{mjwright@liverpool.ac.uk}
\end{DUlineblock}

\sphinxlineitem{version}
\sphinxAtStartPar
1.0.0 \sphinxhyphen{} First full implementation of module.

\end{description}\end{quote}
\index{Clock (class in qnav.measurement.clock)@\spxentry{Clock}\spxextra{class in qnav.measurement.clock}}

\begin{fulllineitems}
\phantomsection\label{\detokenize{modules/measurement/clock:qnav.measurement.clock.Clock}}
\pysigstartsignatures
\pysiglinewithargsret{\sphinxbfcode{\sphinxupquote{class\DUrole{w}{ }}}\sphinxcode{\sphinxupquote{qnav.measurement.clock.}}\sphinxbfcode{\sphinxupquote{Clock}}}{\sphinxparam{\DUrole{n}{clock\_bias}\DUrole{p}{:}\DUrole{w}{ }\DUrole{n}{float}\DUrole{w}{ }\DUrole{o}{=}\DUrole{w}{ }\DUrole{default_value}{0}}\sphinxparamcomma \sphinxparam{\DUrole{n}{bias\_drift\_rate}\DUrole{p}{:}\DUrole{w}{ }\DUrole{n}{float}\DUrole{w}{ }\DUrole{o}{=}\DUrole{w}{ }\DUrole{default_value}{0}}\sphinxparamcomma \sphinxparam{\DUrole{n}{start\_time}\DUrole{p}{:}\DUrole{w}{ }\DUrole{n}{float}\DUrole{w}{ }\DUrole{o}{=}\DUrole{w}{ }\DUrole{default_value}{0.0}}\sphinxparamcomma \sphinxparam{\DUrole{n}{rand\_seed}\DUrole{p}{:}\DUrole{w}{ }\DUrole{n}{int}\DUrole{w}{ }\DUrole{o}{=}\DUrole{w}{ }\DUrole{default_value}{None}}}{}
\pysigstopsignatures
\sphinxAtStartPar
Provides the functionality of a system clock in a simulated environment.
Distinct from typical sensor hardware, the clock represents a virtual
clock that simulated the recording of elapsed time with configured error drift.
\index{update() (qnav.measurement.clock.Clock method)@\spxentry{update()}\spxextra{qnav.measurement.clock.Clock method}}

\begin{fulllineitems}
\phantomsection\label{\detokenize{modules/measurement/clock:qnav.measurement.clock.Clock.update}}
\pysigstartsignatures
\pysiglinewithargsret{\sphinxbfcode{\sphinxupquote{update}}}{\sphinxparam{\DUrole{n}{time\_sec}\DUrole{p}{:}\DUrole{w}{ }\DUrole{n}{float}}}{}
\pysigstopsignatures
\sphinxAtStartPar
Called every iteration to update clock measurement errors.
As clock errors propagate with time, it is important that this
method is called prior to every measurement update.
\begin{quote}\begin{description}
\sphinxlineitem{Parameters}
\sphinxAtStartPar
\sphinxstyleliteralstrong{\sphinxupquote{time\_sec}} (\sphinxstyleliteralemphasis{\sphinxupquote{float}}) \textendash{} The current simulation time in seconds.

\end{description}\end{quote}

\end{fulllineitems}

\index{get\_time() (qnav.measurement.clock.Clock method)@\spxentry{get\_time()}\spxextra{qnav.measurement.clock.Clock method}}

\begin{fulllineitems}
\phantomsection\label{\detokenize{modules/measurement/clock:qnav.measurement.clock.Clock.get_time}}
\pysigstartsignatures
\pysiglinewithargsret{\sphinxbfcode{\sphinxupquote{get\_time}}}{\sphinxparam{\DUrole{n}{current\_time}\DUrole{p}{:}\DUrole{w}{ }\DUrole{n}{float}}}{{ $\rightarrow$ float}}
\pysigstopsignatures
\sphinxAtStartPar
Returns the believed elapsed time since the start of the simulation.
This is given in seconds (with considerations for rounding errors).
\begin{quote}\begin{description}
\sphinxlineitem{Parameters}
\sphinxAtStartPar
\sphinxstyleliteralstrong{\sphinxupquote{current\_time}} (\sphinxstyleliteralemphasis{\sphinxupquote{float}}) \textendash{} The true current time of the simulation.

\sphinxlineitem{Returns}
\sphinxAtStartPar
The current time with clock errors applied (in seconds).

\sphinxlineitem{Return type}
\sphinxAtStartPar
float

\end{description}\end{quote}

\end{fulllineitems}


\end{fulllineitems}

\index{DummyClock (class in qnav.measurement.clock)@\spxentry{DummyClock}\spxextra{class in qnav.measurement.clock}}

\begin{fulllineitems}
\phantomsection\label{\detokenize{modules/measurement/clock:qnav.measurement.clock.DummyClock}}
\pysigstartsignatures
\pysiglinewithargsret{\sphinxbfcode{\sphinxupquote{class\DUrole{w}{ }}}\sphinxcode{\sphinxupquote{qnav.measurement.clock.}}\sphinxbfcode{\sphinxupquote{DummyClock}}}{\sphinxparam{\DUrole{n}{start\_time}\DUrole{p}{:}\DUrole{w}{ }\DUrole{n}{float}\DUrole{w}{ }\DUrole{o}{=}\DUrole{w}{ }\DUrole{default_value}{0.0}}}{}
\pysigstopsignatures\index{update() (qnav.measurement.clock.DummyClock method)@\spxentry{update()}\spxextra{qnav.measurement.clock.DummyClock method}}

\begin{fulllineitems}
\phantomsection\label{\detokenize{modules/measurement/clock:qnav.measurement.clock.DummyClock.update}}
\pysigstartsignatures
\pysiglinewithargsret{\sphinxbfcode{\sphinxupquote{update}}}{\sphinxparam{\DUrole{n}{time\_sec}\DUrole{p}{:}\DUrole{w}{ }\DUrole{n}{float}}}{}
\pysigstopsignatures
\sphinxAtStartPar
Called every iteration to update clock measurement errors.
As clock errors propagate with time, it is important that this
method is called prior to every measurement update.
\begin{quote}\begin{description}
\sphinxlineitem{Parameters}
\sphinxAtStartPar
\sphinxstyleliteralstrong{\sphinxupquote{time\_sec}} (\sphinxstyleliteralemphasis{\sphinxupquote{float}}) \textendash{} The current simulation time in seconds.

\end{description}\end{quote}

\end{fulllineitems}

\index{get\_time() (qnav.measurement.clock.DummyClock method)@\spxentry{get\_time()}\spxextra{qnav.measurement.clock.DummyClock method}}

\begin{fulllineitems}
\phantomsection\label{\detokenize{modules/measurement/clock:qnav.measurement.clock.DummyClock.get_time}}
\pysigstartsignatures
\pysiglinewithargsret{\sphinxbfcode{\sphinxupquote{get\_time}}}{\sphinxparam{\DUrole{n}{current\_time}\DUrole{p}{:}\DUrole{w}{ }\DUrole{n}{float}}}{{ $\rightarrow$ float}}
\pysigstopsignatures
\sphinxAtStartPar
Returns the believed elapsed time since the start of the simulation.
This is given in seconds (with considerations for rounding errors).
\begin{quote}\begin{description}
\sphinxlineitem{Parameters}
\sphinxAtStartPar
\sphinxstyleliteralstrong{\sphinxupquote{current\_time}} (\sphinxstyleliteralemphasis{\sphinxupquote{float}}) \textendash{} The true current time of the simulation.

\sphinxlineitem{Returns}
\sphinxAtStartPar
The current time with clock errors applied (in seconds).

\sphinxlineitem{Return type}
\sphinxAtStartPar
float

\end{description}\end{quote}

\end{fulllineitems}


\end{fulllineitems}



\subsection{Output}
\label{\detokenize{usage/packages:output}}
\sphinxAtStartPar
Modules related to exporting output data:

\sphinxstepscope


\subsubsection{qnav.output.data}
\label{\detokenize{modules/output/data:module-qnav.output.data}}\label{\detokenize{modules/output/data:qnav-output-data}}\label{\detokenize{modules/output/data::doc}}\index{module@\spxentry{module}!qnav.output.data@\spxentry{qnav.output.data}}\index{qnav.output.data@\spxentry{qnav.output.data}!module@\spxentry{module}}\index{ResultsTable (class in qnav.output.data)@\spxentry{ResultsTable}\spxextra{class in qnav.output.data}}

\begin{fulllineitems}
\phantomsection\label{\detokenize{modules/output/data:qnav.output.data.ResultsTable}}
\pysigstartsignatures
\pysiglinewithargsret{\sphinxbfcode{\sphinxupquote{class\DUrole{w}{ }}}\sphinxcode{\sphinxupquote{qnav.output.data.}}\sphinxbfcode{\sphinxupquote{ResultsTable}}}{\sphinxparam{\DUrole{n}{file\_path}\DUrole{p}{:}\DUrole{w}{ }\DUrole{n}{Path}}\sphinxparamcomma \sphinxparam{\DUrole{n}{read\_only}\DUrole{p}{:}\DUrole{w}{ }\DUrole{n}{bool}\DUrole{w}{ }\DUrole{o}{=}\DUrole{w}{ }\DUrole{default_value}{False}}\sphinxparamcomma \sphinxparam{\DUrole{n}{auto\_load}\DUrole{p}{:}\DUrole{w}{ }\DUrole{n}{bool}\DUrole{w}{ }\DUrole{o}{=}\DUrole{w}{ }\DUrole{default_value}{True}}}{}
\pysigstopsignatures
\sphinxAtStartPar
Dynamic on\sphinxhyphen{}disk data container for results records.
Provides a container that uses virtual\sphinxhyphen{}memory, storing the results on
disk but still allowing fast querying and retrieval. This can be useful
for a number of tasks, including conversion and report generation,
without loading all data into memory.
\index{is\_data\_loaded (qnav.output.data.ResultsTable property)@\spxentry{is\_data\_loaded}\spxextra{qnav.output.data.ResultsTable property}}

\begin{fulllineitems}
\phantomsection\label{\detokenize{modules/output/data:qnav.output.data.ResultsTable.is_data_loaded}}
\pysigstartsignatures
\pysigline{\sphinxbfcode{\sphinxupquote{property\DUrole{w}{ }}}\sphinxbfcode{\sphinxupquote{is\_data\_loaded}}\sphinxbfcode{\sphinxupquote{\DUrole{p}{:}\DUrole{w}{ }bool}}}
\pysigstopsignatures
\sphinxAtStartPar
States data has already been loaded.
\begin{quote}\begin{description}
\sphinxlineitem{Returns}
\sphinxAtStartPar
Returns True if data is loaded.

\sphinxlineitem{Return type}
\sphinxAtStartPar
bool

\end{description}\end{quote}

\end{fulllineitems}


\end{fulllineitems}


\sphinxstepscope


\subsubsection{qnav.output.formats}
\label{\detokenize{modules/output/formats:module-qnav.output.formats}}\label{\detokenize{modules/output/formats:qnav-output-formats}}\label{\detokenize{modules/output/formats::doc}}\index{module@\spxentry{module}!qnav.output.formats@\spxentry{qnav.output.formats}}\index{qnav.output.formats@\spxentry{qnav.output.formats}!module@\spxentry{module}}\index{FileFormat (class in qnav.output.formats)@\spxentry{FileFormat}\spxextra{class in qnav.output.formats}}

\begin{fulllineitems}
\phantomsection\label{\detokenize{modules/output/formats:qnav.output.formats.FileFormat}}
\pysigstartsignatures
\pysiglinewithargsret{\sphinxbfcode{\sphinxupquote{class\DUrole{w}{ }}}\sphinxcode{\sphinxupquote{qnav.output.formats.}}\sphinxbfcode{\sphinxupquote{FileFormat}}}{\sphinxparam{\DUrole{n}{match\_id}\DUrole{p}{:}\DUrole{w}{ }\DUrole{n}{str}}\sphinxparamcomma \sphinxparam{\DUrole{n}{file\_ext}\DUrole{p}{:}\DUrole{w}{ }\DUrole{n}{str}}}{}
\pysigstopsignatures
\sphinxAtStartPar
Defines the fields for each supported file format.
Included as a simpler way to represent the SaveFormat enum class.

\end{fulllineitems}

\index{SaveFormat (class in qnav.output.formats)@\spxentry{SaveFormat}\spxextra{class in qnav.output.formats}}

\begin{fulllineitems}
\phantomsection\label{\detokenize{modules/output/formats:qnav.output.formats.SaveFormat}}
\pysigstartsignatures
\pysiglinewithargsret{\sphinxbfcode{\sphinxupquote{class\DUrole{w}{ }}}\sphinxcode{\sphinxupquote{qnav.output.formats.}}\sphinxbfcode{\sphinxupquote{SaveFormat}}}{\sphinxparam{\DUrole{n}{value}}\sphinxparamcomma \sphinxparam{\DUrole{n}{names=\textless{}not given\textgreater{}}}\sphinxparamcomma \sphinxparam{\DUrole{n}{*values}}\sphinxparamcomma \sphinxparam{\DUrole{n}{module=None}}\sphinxparamcomma \sphinxparam{\DUrole{n}{qualname=None}}\sphinxparamcomma \sphinxparam{\DUrole{n}{type=None}}\sphinxparamcomma \sphinxparam{\DUrole{n}{start=1}}\sphinxparamcomma \sphinxparam{\DUrole{n}{boundary=None}}}{}
\pysigstopsignatures
\sphinxAtStartPar
A collection of the supported file formats for results table.
Controllable flags for selecting the use of the supported file formats
when exporting data. A useful way to efficiently export data in
multiple formats, without employing multiple functions.

\end{fulllineitems}

\index{save\_as\_csv() (in module qnav.output.formats)@\spxentry{save\_as\_csv()}\spxextra{in module qnav.output.formats}}

\begin{fulllineitems}
\phantomsection\label{\detokenize{modules/output/formats:qnav.output.formats.save_as_csv}}
\pysigstartsignatures
\pysiglinewithargsret{\sphinxcode{\sphinxupquote{qnav.output.formats.}}\sphinxbfcode{\sphinxupquote{save\_as\_csv}}}{\sphinxparam{\DUrole{n}{file\_path}\DUrole{p}{:}\DUrole{w}{ }\DUrole{n}{Path}}\sphinxparamcomma \sphinxparam{\DUrole{n}{results}\DUrole{p}{:}\DUrole{w}{ }\DUrole{n}{{\hyperref[\detokenize{modules/output/data:qnav.output.data.ResultsTable}]{\sphinxcrossref{ResultsTable}}}}}\sphinxparamcomma \sphinxparam{\DUrole{n}{down\_sample}\DUrole{p}{:}\DUrole{w}{ }\DUrole{n}{int}\DUrole{w}{ }\DUrole{o}{=}\DUrole{w}{ }\DUrole{default_value}{1}}}{}
\pysigstopsignatures
\sphinxAtStartPar
Exports the given results table as a CSV file.
\begin{quote}\begin{description}
\sphinxlineitem{Parameters}\begin{itemize}
\item {} 
\sphinxAtStartPar
\sphinxstyleliteralstrong{\sphinxupquote{file\_path}} (\sphinxstyleliteralemphasis{\sphinxupquote{pathlib.Path}}) \textendash{} The file path to save the CSV data to.

\item {} 
\sphinxAtStartPar
\sphinxstyleliteralstrong{\sphinxupquote{results}} ({\hyperref[\detokenize{modules/output/data:qnav.output.data.ResultsTable}]{\sphinxcrossref{\sphinxstyleliteralemphasis{\sphinxupquote{ResultsTable}}}}}) \textendash{} The results to convert to CSV.

\item {} 
\sphinxAtStartPar
\sphinxstyleliteralstrong{\sphinxupquote{down\_sample}} (\sphinxstyleliteralemphasis{\sphinxupquote{int}}) \textendash{} Optional down\sphinxhyphen{}sampling factor to apply.

\end{itemize}

\end{description}\end{quote}

\end{fulllineitems}

\index{save\_as\_csv\_comp() (in module qnav.output.formats)@\spxentry{save\_as\_csv\_comp()}\spxextra{in module qnav.output.formats}}

\begin{fulllineitems}
\phantomsection\label{\detokenize{modules/output/formats:qnav.output.formats.save_as_csv_comp}}
\pysigstartsignatures
\pysiglinewithargsret{\sphinxcode{\sphinxupquote{qnav.output.formats.}}\sphinxbfcode{\sphinxupquote{save\_as\_csv\_comp}}}{\sphinxparam{\DUrole{n}{file\_path}\DUrole{p}{:}\DUrole{w}{ }\DUrole{n}{Path}}\sphinxparamcomma \sphinxparam{\DUrole{n}{results}\DUrole{p}{:}\DUrole{w}{ }\DUrole{n}{{\hyperref[\detokenize{modules/output/data:qnav.output.data.ResultsTable}]{\sphinxcrossref{ResultsTable}}}}}\sphinxparamcomma \sphinxparam{\DUrole{n}{down\_sample}\DUrole{p}{:}\DUrole{w}{ }\DUrole{n}{int}\DUrole{w}{ }\DUrole{o}{=}\DUrole{w}{ }\DUrole{default_value}{1}}}{}
\pysigstopsignatures
\sphinxAtStartPar
Exports the given results table as a compressed CSV file.
\begin{quote}\begin{description}
\sphinxlineitem{Parameters}\begin{itemize}
\item {} 
\sphinxAtStartPar
\sphinxstyleliteralstrong{\sphinxupquote{file\_path}} (\sphinxstyleliteralemphasis{\sphinxupquote{pathlib.Path}}) \textendash{} The file path to save the compressed CSV data to.

\item {} 
\sphinxAtStartPar
\sphinxstyleliteralstrong{\sphinxupquote{results}} ({\hyperref[\detokenize{modules/output/data:qnav.output.data.ResultsTable}]{\sphinxcrossref{\sphinxstyleliteralemphasis{\sphinxupquote{ResultsTable}}}}}) \textendash{} The results to convert to compressed CSV.

\item {} 
\sphinxAtStartPar
\sphinxstyleliteralstrong{\sphinxupquote{down\_sample}} (\sphinxstyleliteralemphasis{\sphinxupquote{int}}) \textendash{} Optional down\sphinxhyphen{}sampling factor to apply.

\end{itemize}

\end{description}\end{quote}

\end{fulllineitems}

\index{save\_as\_numpy() (in module qnav.output.formats)@\spxentry{save\_as\_numpy()}\spxextra{in module qnav.output.formats}}

\begin{fulllineitems}
\phantomsection\label{\detokenize{modules/output/formats:qnav.output.formats.save_as_numpy}}
\pysigstartsignatures
\pysiglinewithargsret{\sphinxcode{\sphinxupquote{qnav.output.formats.}}\sphinxbfcode{\sphinxupquote{save\_as\_numpy}}}{\sphinxparam{\DUrole{n}{file\_path}\DUrole{p}{:}\DUrole{w}{ }\DUrole{n}{Path}}\sphinxparamcomma \sphinxparam{\DUrole{n}{results}\DUrole{p}{:}\DUrole{w}{ }\DUrole{n}{{\hyperref[\detokenize{modules/output/data:qnav.output.data.ResultsTable}]{\sphinxcrossref{ResultsTable}}}}}\sphinxparamcomma \sphinxparam{\DUrole{n}{down\_sample}\DUrole{p}{:}\DUrole{w}{ }\DUrole{n}{int}\DUrole{w}{ }\DUrole{o}{=}\DUrole{w}{ }\DUrole{default_value}{1}}}{}
\pysigstopsignatures
\sphinxAtStartPar
Exports the given results table as a numpy data file.
\begin{quote}\begin{description}
\sphinxlineitem{Parameters}\begin{itemize}
\item {} 
\sphinxAtStartPar
\sphinxstyleliteralstrong{\sphinxupquote{file\_path}} (\sphinxstyleliteralemphasis{\sphinxupquote{pathlib.Path}}) \textendash{} The file path to save the numpy data to.

\item {} 
\sphinxAtStartPar
\sphinxstyleliteralstrong{\sphinxupquote{results}} ({\hyperref[\detokenize{modules/output/data:qnav.output.data.ResultsTable}]{\sphinxcrossref{\sphinxstyleliteralemphasis{\sphinxupquote{ResultsTable}}}}}) \textendash{} The results to convert to numpy save file.

\item {} 
\sphinxAtStartPar
\sphinxstyleliteralstrong{\sphinxupquote{down\_sample}} (\sphinxstyleliteralemphasis{\sphinxupquote{int}}) \textendash{} Optional down\sphinxhyphen{}sampling factor to apply.

\end{itemize}

\end{description}\end{quote}

\end{fulllineitems}

\index{save\_as\_numpy\_comp() (in module qnav.output.formats)@\spxentry{save\_as\_numpy\_comp()}\spxextra{in module qnav.output.formats}}

\begin{fulllineitems}
\phantomsection\label{\detokenize{modules/output/formats:qnav.output.formats.save_as_numpy_comp}}
\pysigstartsignatures
\pysiglinewithargsret{\sphinxcode{\sphinxupquote{qnav.output.formats.}}\sphinxbfcode{\sphinxupquote{save\_as\_numpy\_comp}}}{\sphinxparam{\DUrole{n}{file\_path}\DUrole{p}{:}\DUrole{w}{ }\DUrole{n}{Path}}\sphinxparamcomma \sphinxparam{\DUrole{n}{results}\DUrole{p}{:}\DUrole{w}{ }\DUrole{n}{{\hyperref[\detokenize{modules/output/data:qnav.output.data.ResultsTable}]{\sphinxcrossref{ResultsTable}}}}}\sphinxparamcomma \sphinxparam{\DUrole{n}{down\_sample}\DUrole{p}{:}\DUrole{w}{ }\DUrole{n}{int}\DUrole{w}{ }\DUrole{o}{=}\DUrole{w}{ }\DUrole{default_value}{1}}}{}
\pysigstopsignatures
\sphinxAtStartPar
Exports the given results table as a compressed numpy data file.
\begin{quote}\begin{description}
\sphinxlineitem{Parameters}\begin{itemize}
\item {} 
\sphinxAtStartPar
\sphinxstyleliteralstrong{\sphinxupquote{file\_path}} (\sphinxstyleliteralemphasis{\sphinxupquote{pathlib.Path}}) \textendash{} The file path to save the compressed numpy data to.

\item {} 
\sphinxAtStartPar
\sphinxstyleliteralstrong{\sphinxupquote{results}} ({\hyperref[\detokenize{modules/output/data:qnav.output.data.ResultsTable}]{\sphinxcrossref{\sphinxstyleliteralemphasis{\sphinxupquote{ResultsTable}}}}}) \textendash{} The results to convert to compressed numpy save file.

\item {} 
\sphinxAtStartPar
\sphinxstyleliteralstrong{\sphinxupquote{down\_sample}} (\sphinxstyleliteralemphasis{\sphinxupquote{int}}) \textendash{} Optional down\sphinxhyphen{}sampling factor to apply.

\end{itemize}

\end{description}\end{quote}

\end{fulllineitems}

\index{save\_as\_binary() (in module qnav.output.formats)@\spxentry{save\_as\_binary()}\spxextra{in module qnav.output.formats}}

\begin{fulllineitems}
\phantomsection\label{\detokenize{modules/output/formats:qnav.output.formats.save_as_binary}}
\pysigstartsignatures
\pysiglinewithargsret{\sphinxcode{\sphinxupquote{qnav.output.formats.}}\sphinxbfcode{\sphinxupquote{save\_as\_binary}}}{\sphinxparam{\DUrole{n}{file\_path}\DUrole{p}{:}\DUrole{w}{ }\DUrole{n}{Path}}\sphinxparamcomma \sphinxparam{\DUrole{n}{results}\DUrole{p}{:}\DUrole{w}{ }\DUrole{n}{{\hyperref[\detokenize{modules/output/data:qnav.output.data.ResultsTable}]{\sphinxcrossref{ResultsTable}}}}}\sphinxparamcomma \sphinxparam{\DUrole{n}{down\_sample}\DUrole{p}{:}\DUrole{w}{ }\DUrole{n}{int}\DUrole{w}{ }\DUrole{o}{=}\DUrole{w}{ }\DUrole{default_value}{1}}}{}
\pysigstopsignatures
\sphinxAtStartPar
Exports the given results table as a compressed numpy data file.
\begin{quote}\begin{description}
\sphinxlineitem{Parameters}\begin{itemize}
\item {} 
\sphinxAtStartPar
\sphinxstyleliteralstrong{\sphinxupquote{file\_path}} (\sphinxstyleliteralemphasis{\sphinxupquote{pathlib.Path}}) \textendash{} The file path to save the compressed numpy data to.

\item {} 
\sphinxAtStartPar
\sphinxstyleliteralstrong{\sphinxupquote{results}} ({\hyperref[\detokenize{modules/output/data:qnav.output.data.ResultsTable}]{\sphinxcrossref{\sphinxstyleliteralemphasis{\sphinxupquote{ResultsTable}}}}}) \textendash{} The results to convert to compressed numpy save file.

\item {} 
\sphinxAtStartPar
\sphinxstyleliteralstrong{\sphinxupquote{down\_sample}} (\sphinxstyleliteralemphasis{\sphinxupquote{int}}) \textendash{} Optional down\sphinxhyphen{}sampling factor to apply.

\end{itemize}

\end{description}\end{quote}

\end{fulllineitems}

\index{save\_as\_matlab() (in module qnav.output.formats)@\spxentry{save\_as\_matlab()}\spxextra{in module qnav.output.formats}}

\begin{fulllineitems}
\phantomsection\label{\detokenize{modules/output/formats:qnav.output.formats.save_as_matlab}}
\pysigstartsignatures
\pysiglinewithargsret{\sphinxcode{\sphinxupquote{qnav.output.formats.}}\sphinxbfcode{\sphinxupquote{save\_as\_matlab}}}{\sphinxparam{\DUrole{n}{file\_path}\DUrole{p}{:}\DUrole{w}{ }\DUrole{n}{Path}}\sphinxparamcomma \sphinxparam{\DUrole{n}{results}\DUrole{p}{:}\DUrole{w}{ }\DUrole{n}{{\hyperref[\detokenize{modules/output/data:qnav.output.data.ResultsTable}]{\sphinxcrossref{ResultsTable}}}}}\sphinxparamcomma \sphinxparam{\DUrole{n}{down\_sample}\DUrole{p}{:}\DUrole{w}{ }\DUrole{n}{int}\DUrole{w}{ }\DUrole{o}{=}\DUrole{w}{ }\DUrole{default_value}{1}}}{}
\pysigstopsignatures
\sphinxAtStartPar
Exports the given results table as a (non\sphinxhyphen{}compressed) MATLAB data file.
\begin{quote}\begin{description}
\sphinxlineitem{Parameters}\begin{itemize}
\item {} 
\sphinxAtStartPar
\sphinxstyleliteralstrong{\sphinxupquote{file\_path}} (\sphinxstyleliteralemphasis{\sphinxupquote{pathlib.Path}}) \textendash{} The file path to save the MATLAB data to.

\item {} 
\sphinxAtStartPar
\sphinxstyleliteralstrong{\sphinxupquote{results}} ({\hyperref[\detokenize{modules/output/data:qnav.output.data.ResultsTable}]{\sphinxcrossref{\sphinxstyleliteralemphasis{\sphinxupquote{ResultsTable}}}}}) \textendash{} The results to convert to MATLAB file.

\item {} 
\sphinxAtStartPar
\sphinxstyleliteralstrong{\sphinxupquote{down\_sample}} (\sphinxstyleliteralemphasis{\sphinxupquote{int}}) \textendash{} Optional down\sphinxhyphen{}sampling factor to apply.

\end{itemize}

\end{description}\end{quote}

\end{fulllineitems}

\index{save\_as\_matlab\_comp() (in module qnav.output.formats)@\spxentry{save\_as\_matlab\_comp()}\spxextra{in module qnav.output.formats}}

\begin{fulllineitems}
\phantomsection\label{\detokenize{modules/output/formats:qnav.output.formats.save_as_matlab_comp}}
\pysigstartsignatures
\pysiglinewithargsret{\sphinxcode{\sphinxupquote{qnav.output.formats.}}\sphinxbfcode{\sphinxupquote{save\_as\_matlab\_comp}}}{\sphinxparam{\DUrole{n}{file\_path}\DUrole{p}{:}\DUrole{w}{ }\DUrole{n}{Path}}\sphinxparamcomma \sphinxparam{\DUrole{n}{results}\DUrole{p}{:}\DUrole{w}{ }\DUrole{n}{{\hyperref[\detokenize{modules/output/data:qnav.output.data.ResultsTable}]{\sphinxcrossref{ResultsTable}}}}}\sphinxparamcomma \sphinxparam{\DUrole{n}{down\_sample}\DUrole{p}{:}\DUrole{w}{ }\DUrole{n}{int}\DUrole{w}{ }\DUrole{o}{=}\DUrole{w}{ }\DUrole{default_value}{1}}}{}
\pysigstopsignatures
\sphinxAtStartPar
Exports the given results table as a MATLAB data file (with compression).
\begin{quote}\begin{description}
\sphinxlineitem{Parameters}\begin{itemize}
\item {} 
\sphinxAtStartPar
\sphinxstyleliteralstrong{\sphinxupquote{file\_path}} (\sphinxstyleliteralemphasis{\sphinxupquote{pathlib.Path}}) \textendash{} The file path to save the MATLAB data to.

\item {} 
\sphinxAtStartPar
\sphinxstyleliteralstrong{\sphinxupquote{results}} ({\hyperref[\detokenize{modules/output/data:qnav.output.data.ResultsTable}]{\sphinxcrossref{\sphinxstyleliteralemphasis{\sphinxupquote{ResultsTable}}}}}) \textendash{} The results to convert to MATLAB file.

\item {} 
\sphinxAtStartPar
\sphinxstyleliteralstrong{\sphinxupquote{down\_sample}} (\sphinxstyleliteralemphasis{\sphinxupquote{int}}) \textendash{} Optional down\sphinxhyphen{}sampling factor to apply.

\end{itemize}

\end{description}\end{quote}

\end{fulllineitems}

\index{save\_as() (in module qnav.output.formats)@\spxentry{save\_as()}\spxextra{in module qnav.output.formats}}

\begin{fulllineitems}
\phantomsection\label{\detokenize{modules/output/formats:qnav.output.formats.save_as}}
\pysigstartsignatures
\pysiglinewithargsret{\sphinxcode{\sphinxupquote{qnav.output.formats.}}\sphinxbfcode{\sphinxupquote{save\_as}}}{\sphinxparam{\DUrole{n}{file\_dir}\DUrole{p}{:}\DUrole{w}{ }\DUrole{n}{Path}}\sphinxparamcomma \sphinxparam{\DUrole{n}{file\_name}\DUrole{p}{:}\DUrole{w}{ }\DUrole{n}{str}}\sphinxparamcomma \sphinxparam{\DUrole{n}{results}\DUrole{p}{:}\DUrole{w}{ }\DUrole{n}{{\hyperref[\detokenize{modules/output/data:qnav.output.data.ResultsTable}]{\sphinxcrossref{ResultsTable}}}}}\sphinxparamcomma \sphinxparam{\DUrole{n}{down\_sample}\DUrole{p}{:}\DUrole{w}{ }\DUrole{n}{int}\DUrole{w}{ }\DUrole{o}{=}\DUrole{w}{ }\DUrole{default_value}{1}}\sphinxparamcomma \sphinxparam{\DUrole{o}{*}\DUrole{n}{formats}\DUrole{p}{:}\DUrole{w}{ }\DUrole{n}{{\hyperref[\detokenize{modules/output/formats:qnav.output.formats.SaveFormat}]{\sphinxcrossref{SaveFormat}}}}}}{}
\pysigstopsignatures
\sphinxAtStartPar
Saves the given results as one or format types.
\begin{quote}\begin{description}
\sphinxlineitem{Parameters}\begin{itemize}
\item {} 
\sphinxAtStartPar
\sphinxstyleliteralstrong{\sphinxupquote{file\_dir}} (\sphinxstyleliteralemphasis{\sphinxupquote{pathlib.Path}}) \textendash{} The location to write results to.

\item {} 
\sphinxAtStartPar
\sphinxstyleliteralstrong{\sphinxupquote{file\_name}} (\sphinxstyleliteralemphasis{\sphinxupquote{str}}) \textendash{} The filename prefix to use.

\item {} 
\sphinxAtStartPar
\sphinxstyleliteralstrong{\sphinxupquote{results}} ({\hyperref[\detokenize{modules/output/data:qnav.output.data.ResultsTable}]{\sphinxcrossref{\sphinxstyleliteralemphasis{\sphinxupquote{ResultsTable}}}}}) \textendash{} The results to convert and write.

\item {} 
\sphinxAtStartPar
\sphinxstyleliteralstrong{\sphinxupquote{down\_sample}} (\sphinxstyleliteralemphasis{\sphinxupquote{int}}) \textendash{} Optional down\sphinxhyphen{}sampling factor to apply.

\item {} 
\sphinxAtStartPar
\sphinxstyleliteralstrong{\sphinxupquote{formats}} ({\hyperref[\detokenize{modules/output/formats:qnav.output.formats.SaveFormat}]{\sphinxcrossref{\sphinxstyleliteralemphasis{\sphinxupquote{SaveFormat}}}}}) \textendash{} The output formats to use.

\end{itemize}

\end{description}\end{quote}

\end{fulllineitems}


\sphinxstepscope


\subsubsection{qnav.output.figures}
\label{\detokenize{modules/output/figures:module-qnav.output.figures}}\label{\detokenize{modules/output/figures:qnav-output-figures}}\label{\detokenize{modules/output/figures::doc}}\index{module@\spxentry{module}!qnav.output.figures@\spxentry{qnav.output.figures}}\index{qnav.output.figures@\spxentry{qnav.output.figures}!module@\spxentry{module}}
\sphinxstepscope


\subsubsection{qnav.output.summary}
\label{\detokenize{modules/output/summary:module-qnav.output.summary}}\label{\detokenize{modules/output/summary:qnav-output-summary}}\label{\detokenize{modules/output/summary::doc}}\index{module@\spxentry{module}!qnav.output.summary@\spxentry{qnav.output.summary}}\index{qnav.output.summary@\spxentry{qnav.output.summary}!module@\spxentry{module}}

\subsection{Quantum}
\label{\detokenize{usage/packages:quantum}}
\sphinxAtStartPar
Modules relating to quantum sensor and fusion methods:

\sphinxstepscope


\subsubsection{qnav.quantum.base}
\label{\detokenize{modules/quantum/base:module-qnav.quantum.base}}\label{\detokenize{modules/quantum/base:qnav-quantum-base}}\label{\detokenize{modules/quantum/base::doc}}\index{module@\spxentry{module}!qnav.quantum.base@\spxentry{qnav.quantum.base}}\index{qnav.quantum.base@\spxentry{qnav.quantum.base}!module@\spxentry{module}}

\paragraph{base.py}
\label{\detokenize{modules/quantum/base:base-py}}\begin{quote}\begin{description}
\sphinxlineitem{summary}
\sphinxAtStartPar
Provides basic concept quantum sensors and fusion methods.
These serve as base models that more advanced and realistic
implementations can extend from.

\sphinxlineitem{author}
\sphinxAtStartPar
Michael Wright \sphinxhyphen{} \sphinxhref{mailto:mjwright@liverpool.ac.uk}{mjwright@liverpool.ac.uk}

\sphinxlineitem{version}
\sphinxAtStartPar
1.0.0 \sphinxhyphen{} First full implementation of class.

\end{description}\end{quote}
\index{ConceptQuantumImu (class in qnav.quantum.base)@\spxentry{ConceptQuantumImu}\spxextra{class in qnav.quantum.base}}

\begin{fulllineitems}
\phantomsection\label{\detokenize{modules/quantum/base:qnav.quantum.base.ConceptQuantumImu}}
\pysigstartsignatures
\pysiglinewithargsret{\sphinxbfcode{\sphinxupquote{class\DUrole{w}{ }}}\sphinxcode{\sphinxupquote{qnav.quantum.base.}}\sphinxbfcode{\sphinxupquote{ConceptQuantumImu}}}{\sphinxparam{\DUrole{n}{full\_frequency: float}}\sphinxparamcomma \sphinxparam{\DUrole{n}{imu\_frequency: float}}\sphinxparamcomma \sphinxparam{\DUrole{n}{accelerometer\_errors: \textasciitilde{}qnav.measurement.error\_properties.ErrorProperties}}\sphinxparamcomma \sphinxparam{\DUrole{n}{gyroscope\_errors: \textasciitilde{}qnav.measurement.error\_properties.ErrorProperties}}\sphinxparamcomma \sphinxparam{\DUrole{n}{sensor\_axis: \textasciitilde{}qnav.measurement.platform.SensorAxis = \textless{}qnav.measurement.platform.SensorAxis object\textgreater{}}}\sphinxparamcomma \sphinxparam{\DUrole{n}{duty\_cycle: float = 0.5}}\sphinxparamcomma \sphinxparam{\DUrole{n}{start\_time: float = 0.0}}\sphinxparamcomma \sphinxparam{\DUrole{n}{rand\_seed: int = None}}}{}
\pysigstopsignatures
\sphinxAtStartPar
Concept Quantum IMU featuring use of quantum accelerometer and gyroscope.
While not realistic in implementation, this base class produces
measurements mimicking a quantum accelerometer and gyroscope,
requiring multiple samples of the ground truth.
\index{take\_measurement() (qnav.quantum.base.ConceptQuantumImu method)@\spxentry{take\_measurement()}\spxextra{qnav.quantum.base.ConceptQuantumImu method}}

\begin{fulllineitems}
\phantomsection\label{\detokenize{modules/quantum/base:qnav.quantum.base.ConceptQuantumImu.take_measurement}}
\pysigstartsignatures
\pysiglinewithargsret{\sphinxbfcode{\sphinxupquote{take\_measurement}}}{\sphinxparam{\DUrole{n}{est\_time}\DUrole{p}{:}\DUrole{w}{ }\DUrole{n}{float}}\sphinxparamcomma \sphinxparam{\DUrole{n}{ground\_truth}\DUrole{p}{:}\DUrole{w}{ }\DUrole{n}{{\hyperref[\detokenize{modules/waypoints/trajectory:qnav.waypoints.trajectory.GroundTruth}]{\sphinxcrossref{GroundTruth}}}}}}{}
\pysigstopsignatures
\sphinxAtStartPar
Samples from the ground truth to collect data to produce measurements.
Using the ground truth, captures the current acceleration and angles
rates at the time. A series of these are required to produce a
measurement before fusion.
\begin{quote}\begin{description}
\sphinxlineitem{Parameters}\begin{itemize}
\item {} 
\sphinxAtStartPar
\sphinxstyleliteralstrong{\sphinxupquote{est\_time}} (\sphinxstyleliteralemphasis{\sphinxupquote{float}}) \textendash{} The estimated time the measurement of the sensor
is expected to be captured at (in seconds).

\item {} 
\sphinxAtStartPar
\sphinxstyleliteralstrong{\sphinxupquote{ground\_truth}} ({\hyperref[\detokenize{modules/waypoints/trajectory:qnav.waypoints.trajectory.GroundTruth}]{\sphinxcrossref{\sphinxstyleliteralemphasis{\sphinxupquote{GroundTruth}}}}}) \textendash{} The corresponding ground truth data.

\end{itemize}

\end{description}\end{quote}

\end{fulllineitems}

\index{update() (qnav.quantum.base.ConceptQuantumImu method)@\spxentry{update()}\spxextra{qnav.quantum.base.ConceptQuantumImu method}}

\begin{fulllineitems}
\phantomsection\label{\detokenize{modules/quantum/base:qnav.quantum.base.ConceptQuantumImu.update}}
\pysigstartsignatures
\pysiglinewithargsret{\sphinxbfcode{\sphinxupquote{update}}}{\sphinxparam{\DUrole{n}{time\_sec}\DUrole{p}{:}\DUrole{w}{ }\DUrole{n}{float}\DUrole{w}{ }\DUrole{o}{=}\DUrole{w}{ }\DUrole{default_value}{None}}}{}
\pysigstopsignatures
\sphinxAtStartPar
Called on demand to update the sensor’s state.
Typically called after measurements have been taken. When called
this increases the next measurement time and propagates internal
states, to the time difference.
\begin{quote}\begin{description}
\sphinxlineitem{Parameters}
\sphinxAtStartPar
\sphinxstyleliteralstrong{\sphinxupquote{time\_sec}} (\sphinxstyleliteralemphasis{\sphinxupquote{float}}) \textendash{} The current simulation time in seconds.

\end{description}\end{quote}

\end{fulllineitems}

\index{reset() (qnav.quantum.base.ConceptQuantumImu method)@\spxentry{reset()}\spxextra{qnav.quantum.base.ConceptQuantumImu method}}

\begin{fulllineitems}
\phantomsection\label{\detokenize{modules/quantum/base:qnav.quantum.base.ConceptQuantumImu.reset}}
\pysigstartsignatures
\pysiglinewithargsret{\sphinxbfcode{\sphinxupquote{reset}}}{}{}
\pysigstopsignatures
\sphinxAtStartPar
Resets the sensor, clearing all previous measurements.
Normally used directly after fusion.

\end{fulllineitems}

\index{last\_measurement (qnav.quantum.base.ConceptQuantumImu property)@\spxentry{last\_measurement}\spxextra{qnav.quantum.base.ConceptQuantumImu property}}

\begin{fulllineitems}
\phantomsection\label{\detokenize{modules/quantum/base:qnav.quantum.base.ConceptQuantumImu.last_measurement}}
\pysigstartsignatures
\pysigline{\sphinxbfcode{\sphinxupquote{property\DUrole{w}{ }}}\sphinxbfcode{\sphinxupquote{last\_measurement}}\sphinxbfcode{\sphinxupquote{\DUrole{p}{:}\DUrole{w}{ }tuple\DUrole{p}{{[}}ndarray\DUrole{p}{,}\DUrole{w}{ }ndarray\DUrole{p}{{]}}}}}
\pysigstopsignatures
\sphinxAtStartPar
Returns the last acceleration and angle rate measurements captured
by the quantum sensor, respectively.
\begin{quote}\begin{description}
\sphinxlineitem{Returns}
\sphinxAtStartPar
The last measurements of the sensor.

\sphinxlineitem{Return type}
\sphinxAtStartPar
tuple{[}np.ndarray, np.ndarray{]}

\end{description}\end{quote}

\end{fulllineitems}

\index{num\_active\_steps (qnav.quantum.base.ConceptQuantumImu property)@\spxentry{num\_active\_steps}\spxextra{qnav.quantum.base.ConceptQuantumImu property}}

\begin{fulllineitems}
\phantomsection\label{\detokenize{modules/quantum/base:qnav.quantum.base.ConceptQuantumImu.num_active_steps}}
\pysigstartsignatures
\pysigline{\sphinxbfcode{\sphinxupquote{property\DUrole{w}{ }}}\sphinxbfcode{\sphinxupquote{num\_active\_steps}}\sphinxbfcode{\sphinxupquote{\DUrole{p}{:}\DUrole{w}{ }int}}}
\pysigstopsignatures
\sphinxAtStartPar
Returns the number of measurement steps required,
\begin{quote}\begin{description}
\sphinxlineitem{Returns}
\sphinxAtStartPar
Number of measurements required.

\sphinxlineitem{Return type}
\sphinxAtStartPar
int

\end{description}\end{quote}

\end{fulllineitems}

\index{num\_steps (qnav.quantum.base.ConceptQuantumImu property)@\spxentry{num\_steps}\spxextra{qnav.quantum.base.ConceptQuantumImu property}}

\begin{fulllineitems}
\phantomsection\label{\detokenize{modules/quantum/base:qnav.quantum.base.ConceptQuantumImu.num_steps}}
\pysigstartsignatures
\pysigline{\sphinxbfcode{\sphinxupquote{property\DUrole{w}{ }}}\sphinxbfcode{\sphinxupquote{num\_steps}}\sphinxbfcode{\sphinxupquote{\DUrole{p}{:}\DUrole{w}{ }int}}}
\pysigstopsignatures
\sphinxAtStartPar
Returns the number of total steps required,
\begin{quote}\begin{description}
\sphinxlineitem{Returns}
\sphinxAtStartPar
Number of steps required for fusion.

\sphinxlineitem{Return type}
\sphinxAtStartPar
int

\end{description}\end{quote}

\end{fulllineitems}

\index{sensor\_axis (qnav.quantum.base.ConceptQuantumImu property)@\spxentry{sensor\_axis}\spxextra{qnav.quantum.base.ConceptQuantumImu property}}

\begin{fulllineitems}
\phantomsection\label{\detokenize{modules/quantum/base:qnav.quantum.base.ConceptQuantumImu.sensor_axis}}
\pysigstartsignatures
\pysigline{\sphinxbfcode{\sphinxupquote{property\DUrole{w}{ }}}\sphinxbfcode{\sphinxupquote{sensor\_axis}}\sphinxbfcode{\sphinxupquote{\DUrole{p}{:}\DUrole{w}{ }{\hyperref[\detokenize{modules/measurement/platform:qnav.measurement.platform.SensorAxis}]{\sphinxcrossref{SensorAxis}}}}}}
\pysigstopsignatures
\sphinxAtStartPar
Returns the true sensor axis for the sensor.
\begin{quote}\begin{description}
\sphinxlineitem{Returns}
\sphinxAtStartPar
The sensor axis used by the sensor.

\sphinxlineitem{Return type}
\sphinxAtStartPar
{\hyperref[\detokenize{modules/measurement/platform:qnav.measurement.platform.SensorAxis}]{\sphinxcrossref{SensorAxis}}}

\end{description}\end{quote}

\end{fulllineitems}


\end{fulllineitems}

\index{ConceptQuantumFusion (class in qnav.quantum.base)@\spxentry{ConceptQuantumFusion}\spxextra{class in qnav.quantum.base}}

\begin{fulllineitems}
\phantomsection\label{\detokenize{modules/quantum/base:qnav.quantum.base.ConceptQuantumFusion}}
\pysigstartsignatures
\pysiglinewithargsret{\sphinxbfcode{\sphinxupquote{class\DUrole{w}{ }}}\sphinxcode{\sphinxupquote{qnav.quantum.base.}}\sphinxbfcode{\sphinxupquote{ConceptQuantumFusion}}}{\sphinxparam{\DUrole{n}{qs\_imu}\DUrole{p}{:}\DUrole{w}{ }\DUrole{n}{{\hyperref[\detokenize{modules/quantum/base:qnav.quantum.base.ConceptQuantumImu}]{\sphinxcrossref{ConceptQuantumImu}}}}}\sphinxparamcomma \sphinxparam{\DUrole{n}{shared\_accelerometer}\DUrole{p}{:}\DUrole{w}{ }\DUrole{n}{{\hyperref[\detokenize{modules/measurement/accelerometer:qnav.measurement.accelerometer.Accelerometer}]{\sphinxcrossref{Accelerometer}}}}}\sphinxparamcomma \sphinxparam{\DUrole{n}{shared\_gyroscope}\DUrole{p}{:}\DUrole{w}{ }\DUrole{n}{{\hyperref[\detokenize{modules/measurement/gyroscope:qnav.measurement.gyroscope.Gyroscope}]{\sphinxcrossref{Gyroscope}}}}}\sphinxparamcomma \sphinxparam{\DUrole{n}{est\_qs\_imu\_axis}\DUrole{p}{:}\DUrole{w}{ }\DUrole{n}{{\hyperref[\detokenize{modules/measurement/platform:qnav.measurement.platform.SensorAxis}]{\sphinxcrossref{SensorAxis}}}}\DUrole{w}{ }\DUrole{o}{=}\DUrole{w}{ }\DUrole{default_value}{None}}\sphinxparamcomma \sphinxparam{\DUrole{n}{est\_acc\_axis}\DUrole{p}{:}\DUrole{w}{ }\DUrole{n}{{\hyperref[\detokenize{modules/measurement/platform:qnav.measurement.platform.SensorAxis}]{\sphinxcrossref{SensorAxis}}}}\DUrole{w}{ }\DUrole{o}{=}\DUrole{w}{ }\DUrole{default_value}{None}}\sphinxparamcomma \sphinxparam{\DUrole{n}{est\_gyro\_axis}\DUrole{p}{:}\DUrole{w}{ }\DUrole{n}{{\hyperref[\detokenize{modules/measurement/platform:qnav.measurement.platform.SensorAxis}]{\sphinxcrossref{SensorAxis}}}}\DUrole{w}{ }\DUrole{o}{=}\DUrole{w}{ }\DUrole{default_value}{None}}}{}
\pysigstopsignatures
\sphinxAtStartPar
A basic version of quantum fusion, namely used as a base class.
This is the simplest implementation of position correction fusion for
quantum IMU. Many of its methods can be extended and swapped out for
more advanced implementations.
\index{perform\_fusion() (qnav.quantum.base.ConceptQuantumFusion method)@\spxentry{perform\_fusion()}\spxextra{qnav.quantum.base.ConceptQuantumFusion method}}

\begin{fulllineitems}
\phantomsection\label{\detokenize{modules/quantum/base:qnav.quantum.base.ConceptQuantumFusion.perform_fusion}}
\pysigstartsignatures
\pysiglinewithargsret{\sphinxbfcode{\sphinxupquote{perform\_fusion}}}{\sphinxparam{\DUrole{n}{estimated\_state}\DUrole{p}{:}\DUrole{w}{ }\DUrole{n}{{\hyperref[\detokenize{modules/estimation/state:qnav.estimation.state.EstimatedState}]{\sphinxcrossref{EstimatedState}}}}}}{{ $\rightarrow$ None}}
\pysigstopsignatures
\sphinxAtStartPar
Collects measurements and applies position fixing.
Each cycle collects measurements from the typical IMU sensors and
quantum sensor. After a set number of calls, this uses the single
quantum measurement to calculate an approximate bias correction for
the IMU. If successful, the position fixing is applied.
\begin{quote}\begin{description}
\sphinxlineitem{Parameters}
\sphinxAtStartPar
\sphinxstyleliteralstrong{\sphinxupquote{estimated\_state}} ({\hyperref[\detokenize{modules/estimation/state:qnav.estimation.state.EstimatedState}]{\sphinxcrossref{\sphinxstyleliteralemphasis{\sphinxupquote{EstimatedState}}}}}) \textendash{} The current estimated state to update.

\end{description}\end{quote}

\end{fulllineitems}

\index{reset() (qnav.quantum.base.ConceptQuantumFusion method)@\spxentry{reset()}\spxextra{qnav.quantum.base.ConceptQuantumFusion method}}

\begin{fulllineitems}
\phantomsection\label{\detokenize{modules/quantum/base:qnav.quantum.base.ConceptQuantumFusion.reset}}
\pysigstartsignatures
\pysiglinewithargsret{\sphinxbfcode{\sphinxupquote{reset}}}{}{}
\pysigstopsignatures
\sphinxAtStartPar
Resets the state of the fusion, clearing internal temporary states.
When call the fusion will return to expecting the first measurement,
then continuing from there.

\end{fulllineitems}

\index{sensors (qnav.quantum.base.ConceptQuantumFusion property)@\spxentry{sensors}\spxextra{qnav.quantum.base.ConceptQuantumFusion property}}

\begin{fulllineitems}
\phantomsection\label{\detokenize{modules/quantum/base:qnav.quantum.base.ConceptQuantumFusion.sensors}}
\pysigstartsignatures
\pysigline{\sphinxbfcode{\sphinxupquote{property\DUrole{w}{ }}}\sphinxbfcode{\sphinxupquote{sensors}}\sphinxbfcode{\sphinxupquote{\DUrole{p}{:}\DUrole{w}{ }(\textless{}class \textquotesingle{}qnav.quantum.base.ConceptQuantumImu\textquotesingle{}\textgreater{}, \textless{}class \textquotesingle{}qnav.measurement.accelerometer.Accelerometer\textquotesingle{}\textgreater{}, \textless{}class \textquotesingle{}qnav.measurement.gyroscope.Gyroscope\textquotesingle{}\textgreater{})}}}
\pysigstopsignatures
\sphinxAtStartPar
The provided sensors associated with the fusion method.
This is a collection of all shared sensors whose measurements will be
used by the fusion method.
\begin{quote}\begin{description}
\sphinxlineitem{Returns}
\sphinxAtStartPar
The shared sensors used by the fusion method.

\sphinxlineitem{Return type}
\sphinxAtStartPar
Iterable{[}{\hyperref[\detokenize{modules/measurement/sensor:qnav.measurement.sensor.Sensor}]{\sphinxcrossref{Sensor}}}{]}

\end{description}\end{quote}

\end{fulllineitems}


\end{fulllineitems}


\sphinxstepscope


\subsubsection{qnav.quantum.realistic}
\label{\detokenize{modules/quantum/realistic:module-qnav.quantum.realistic}}\label{\detokenize{modules/quantum/realistic:qnav-quantum-realistic}}\label{\detokenize{modules/quantum/realistic::doc}}\index{module@\spxentry{module}!qnav.quantum.realistic@\spxentry{qnav.quantum.realistic}}\index{qnav.quantum.realistic@\spxentry{qnav.quantum.realistic}!module@\spxentry{module}}

\paragraph{realistic.py}
\label{\detokenize{modules/quantum/realistic:realistic-py}}\begin{quote}\begin{description}
\sphinxlineitem{summary}
\sphinxAtStartPar
Provides quantum sensors with realistic implementations.
Extends from the base quantum sensor and includes realistic
physics for capturing measurements.

\sphinxlineitem{author}
\sphinxAtStartPar
Michael Wright \sphinxhyphen{} \sphinxhref{mailto:mjwright@liverpool.ac.uk}{mjwright@liverpool.ac.uk}

\sphinxlineitem{version}
\sphinxAtStartPar
1.0.0 \sphinxhyphen{} First full implementation of class.

\end{description}\end{quote}
\index{ColdAtomInterferometer (class in qnav.quantum.realistic)@\spxentry{ColdAtomInterferometer}\spxextra{class in qnav.quantum.realistic}}

\begin{fulllineitems}
\phantomsection\label{\detokenize{modules/quantum/realistic:qnav.quantum.realistic.ColdAtomInterferometer}}
\pysigstartsignatures
\pysiglinewithargsret{\sphinxbfcode{\sphinxupquote{class\DUrole{w}{ }}}\sphinxcode{\sphinxupquote{qnav.quantum.realistic.}}\sphinxbfcode{\sphinxupquote{ColdAtomInterferometer}}}{\sphinxparam{\DUrole{n}{num\_of\_atoms}\DUrole{p}{:}\DUrole{w}{ }\DUrole{n}{int}\DUrole{w}{ }\DUrole{o}{=}\DUrole{w}{ }\DUrole{default_value}{1000000.0}}\sphinxparamcomma \sphinxparam{\DUrole{n}{atom\_mass}\DUrole{p}{:}\DUrole{w}{ }\DUrole{n}{float}\DUrole{w}{ }\DUrole{o}{=}\DUrole{w}{ }\DUrole{default_value}{2.20693925e\sphinxhyphen{}25}}\sphinxparamcomma \sphinxparam{\DUrole{n}{sensor\_length}\DUrole{p}{:}\DUrole{w}{ }\DUrole{n}{float}\DUrole{w}{ }\DUrole{o}{=}\DUrole{w}{ }\DUrole{default_value}{0.5}}\sphinxparamcomma \sphinxparam{\DUrole{n}{beam\_width}\DUrole{p}{:}\DUrole{w}{ }\DUrole{n}{float}\DUrole{w}{ }\DUrole{o}{=}\DUrole{w}{ }\DUrole{default_value}{0.01}}\sphinxparamcomma \sphinxparam{\DUrole{n}{eta}\DUrole{p}{:}\DUrole{w}{ }\DUrole{n}{float}\DUrole{w}{ }\DUrole{o}{=}\DUrole{w}{ }\DUrole{default_value}{0.3}}\sphinxparamcomma \sphinxparam{\DUrole{n}{time\_pulse}\DUrole{p}{:}\DUrole{w}{ }\DUrole{n}{float}\DUrole{w}{ }\DUrole{o}{=}\DUrole{w}{ }\DUrole{default_value}{0.16}}\sphinxparamcomma \sphinxparam{\DUrole{n}{recoil\_velocity}\DUrole{p}{:}\DUrole{w}{ }\DUrole{n}{float}\DUrole{w}{ }\DUrole{o}{=}\DUrole{w}{ }\DUrole{default_value}{0.007}}}{}
\pysigstopsignatures
\sphinxAtStartPar
Parameters used for defining a Cold Atom Interferometer.
These properties describe the qualities of an ultra\sphinxhyphen{}cold atom
interferometer, which is used for capturing quantum
accelerations and angle rates.

\end{fulllineitems}

\index{QuantumIMU (class in qnav.quantum.realistic)@\spxentry{QuantumIMU}\spxextra{class in qnav.quantum.realistic}}

\begin{fulllineitems}
\phantomsection\label{\detokenize{modules/quantum/realistic:qnav.quantum.realistic.QuantumIMU}}
\pysigstartsignatures
\pysiglinewithargsret{\sphinxbfcode{\sphinxupquote{class\DUrole{w}{ }}}\sphinxcode{\sphinxupquote{qnav.quantum.realistic.}}\sphinxbfcode{\sphinxupquote{QuantumIMU}}}{\sphinxparam{\DUrole{n}{full\_frequency: float}}\sphinxparamcomma \sphinxparam{\DUrole{n}{imu\_frequency: float}}\sphinxparamcomma \sphinxparam{\DUrole{n}{accelerometer\_errors: \textasciitilde{}qnav.measurement.error\_properties.ErrorProperties}}\sphinxparamcomma \sphinxparam{\DUrole{n}{gyroscope\_errors: \textasciitilde{}qnav.measurement.error\_properties.ErrorProperties}}\sphinxparamcomma \sphinxparam{\DUrole{n}{interferometer: \textasciitilde{}qnav.quantum.realistic.ColdAtomInterferometer}}\sphinxparamcomma \sphinxparam{\DUrole{n}{sensor\_axis: \textasciitilde{}qnav.measurement.platform.SensorAxis = \textless{}qnav.measurement.platform.SensorAxis object\textgreater{}}}\sphinxparamcomma \sphinxparam{\DUrole{n}{duty\_cycle: float = 0.5}}\sphinxparamcomma \sphinxparam{\DUrole{n}{start\_time: float = 0.0}}\sphinxparamcomma \sphinxparam{\DUrole{n}{rand\_seed: int = None}}}{}
\pysigstopsignatures
\sphinxAtStartPar
Quantum IMU that uses an ultra\sphinxhyphen{}cold atom interferometer for measurements.
This extends from the base concept quantum IMU by replacing the acquiring
of measurement methods with realistic implementations. These feature
\index{take\_measurement() (qnav.quantum.realistic.QuantumIMU method)@\spxentry{take\_measurement()}\spxextra{qnav.quantum.realistic.QuantumIMU method}}

\begin{fulllineitems}
\phantomsection\label{\detokenize{modules/quantum/realistic:qnav.quantum.realistic.QuantumIMU.take_measurement}}
\pysigstartsignatures
\pysiglinewithargsret{\sphinxbfcode{\sphinxupquote{take\_measurement}}}{\sphinxparam{\DUrole{n}{est\_time}\DUrole{p}{:}\DUrole{w}{ }\DUrole{n}{float}}\sphinxparamcomma \sphinxparam{\DUrole{n}{ground\_truth}\DUrole{p}{:}\DUrole{w}{ }\DUrole{n}{{\hyperref[\detokenize{modules/waypoints/trajectory:qnav.waypoints.trajectory.GroundTruth}]{\sphinxcrossref{GroundTruth}}}}}}{}
\pysigstopsignatures
\sphinxAtStartPar
Samples from the ground truth to collect data to produce measurements.
Using the ground truth, captures the current acceleration and angles
rates at the time. A series of these are required to produce a
measurement before fusion.
\begin{quote}\begin{description}
\sphinxlineitem{Parameters}\begin{itemize}
\item {} 
\sphinxAtStartPar
\sphinxstyleliteralstrong{\sphinxupquote{est\_time}} (\sphinxstyleliteralemphasis{\sphinxupquote{float}}) \textendash{} The estimated time the measurement of the sensor
is expected to be captured at (in seconds).

\item {} 
\sphinxAtStartPar
\sphinxstyleliteralstrong{\sphinxupquote{ground\_truth}} ({\hyperref[\detokenize{modules/waypoints/trajectory:qnav.waypoints.trajectory.GroundTruth}]{\sphinxcrossref{\sphinxstyleliteralemphasis{\sphinxupquote{GroundTruth}}}}}) \textendash{} The corresponding ground truth data.

\end{itemize}

\end{description}\end{quote}

\end{fulllineitems}

\index{reset() (qnav.quantum.realistic.QuantumIMU method)@\spxentry{reset()}\spxextra{qnav.quantum.realistic.QuantumIMU method}}

\begin{fulllineitems}
\phantomsection\label{\detokenize{modules/quantum/realistic:qnav.quantum.realistic.QuantumIMU.reset}}
\pysigstartsignatures
\pysiglinewithargsret{\sphinxbfcode{\sphinxupquote{reset}}}{}{}
\pysigstopsignatures
\sphinxAtStartPar
Resets the sensor, clearing all previous measurements.
Normally used directly after fusion.

\end{fulllineitems}


\end{fulllineitems}

\index{ang\_diff() (in module qnav.quantum.realistic)@\spxentry{ang\_diff()}\spxextra{in module qnav.quantum.realistic}}

\begin{fulllineitems}
\phantomsection\label{\detokenize{modules/quantum/realistic:qnav.quantum.realistic.ang_diff}}
\pysigstartsignatures
\pysiglinewithargsret{\sphinxcode{\sphinxupquote{qnav.quantum.realistic.}}\sphinxbfcode{\sphinxupquote{ang\_diff}}}{\sphinxparam{\DUrole{n}{a}\DUrole{p}{:}\DUrole{w}{ }\DUrole{n}{float}}\sphinxparamcomma \sphinxparam{\DUrole{n}{b}\DUrole{p}{:}\DUrole{w}{ }\DUrole{n}{float}}}{{ $\rightarrow$ float}}
\pysigstopsignatures
\sphinxAtStartPar
Calculate the difference between two angles (in radians).
\begin{quote}\begin{description}
\sphinxlineitem{Parameters}\begin{itemize}
\item {} 
\sphinxAtStartPar
\sphinxstyleliteralstrong{\sphinxupquote{a}} (\sphinxstyleliteralemphasis{\sphinxupquote{float}}) \textendash{} The first angle value.

\item {} 
\sphinxAtStartPar
\sphinxstyleliteralstrong{\sphinxupquote{b}} (\sphinxstyleliteralemphasis{\sphinxupquote{float}}) \textendash{} The second angle value.

\end{itemize}

\sphinxlineitem{Returns}
\sphinxAtStartPar
The difference between the given angles.

\sphinxlineitem{Return type}
\sphinxAtStartPar
float

\end{description}\end{quote}

\end{fulllineitems}


\sphinxstepscope


\subsubsection{qnav.quantum.particle\_filter}
\label{\detokenize{modules/quantum/particle_filter:module-qnav.quantum.particle_filter}}\label{\detokenize{modules/quantum/particle_filter:qnav-quantum-particle-filter}}\label{\detokenize{modules/quantum/particle_filter::doc}}\index{module@\spxentry{module}!qnav.quantum.particle\_filter@\spxentry{qnav.quantum.particle\_filter}}\index{qnav.quantum.particle\_filter@\spxentry{qnav.quantum.particle\_filter}!module@\spxentry{module}}

\paragraph{particle\_filter.py}
\label{\detokenize{modules/quantum/particle_filter:particle-filter-py}}\begin{quote}\begin{description}
\sphinxlineitem{summary}
\sphinxAtStartPar
Particle filter implementations for quantum sensor fusion.
This module contains fusion methods for quantum sensors that employ
particle filter estimation.

\sphinxlineitem{author}
\sphinxAtStartPar
Michael Wright \sphinxhyphen{} \sphinxhref{mailto:mjwright@liverpool.ac.uk}{mjwright@liverpool.ac.uk}

\sphinxlineitem{version}
\sphinxAtStartPar
1.0.0 \sphinxhyphen{} First full implementation of class.

\end{description}\end{quote}
\index{ParticleFilterParams (class in qnav.quantum.particle\_filter)@\spxentry{ParticleFilterParams}\spxextra{class in qnav.quantum.particle\_filter}}

\begin{fulllineitems}
\phantomsection\label{\detokenize{modules/quantum/particle_filter:qnav.quantum.particle_filter.ParticleFilterParams}}
\pysigstartsignatures
\pysiglinewithargsret{\sphinxbfcode{\sphinxupquote{class\DUrole{w}{ }}}\sphinxcode{\sphinxupquote{qnav.quantum.particle\_filter.}}\sphinxbfcode{\sphinxupquote{ParticleFilterParams}}}{\sphinxparam{\DUrole{n}{num\_particles}\DUrole{p}{:}\DUrole{w}{ }\DUrole{n}{int}\DUrole{w}{ }\DUrole{o}{=}\DUrole{w}{ }\DUrole{default_value}{200}}\sphinxparamcomma \sphinxparam{\DUrole{n}{sigma\_accelerometer}\DUrole{p}{:}\DUrole{w}{ }\DUrole{n}{float}\DUrole{w}{ }\DUrole{o}{=}\DUrole{w}{ }\DUrole{default_value}{0.01}}\sphinxparamcomma \sphinxparam{\DUrole{n}{sigma\_gyroscope}\DUrole{p}{:}\DUrole{w}{ }\DUrole{n}{float}\DUrole{w}{ }\DUrole{o}{=}\DUrole{w}{ }\DUrole{default_value}{0.01}}\sphinxparamcomma \sphinxparam{\DUrole{n}{sigma\_misalignment}\DUrole{p}{:}\DUrole{w}{ }\DUrole{n}{float}\DUrole{w}{ }\DUrole{o}{=}\DUrole{w}{ }\DUrole{default_value}{0.01}}\sphinxparamcomma \sphinxparam{\DUrole{n}{rng\_seed}\DUrole{p}{:}\DUrole{w}{ }\DUrole{n}{int}\DUrole{w}{ }\DUrole{o}{=}\DUrole{w}{ }\DUrole{default_value}{None}}\sphinxparamcomma \sphinxparam{\DUrole{n}{low\_noise\_prob}\DUrole{p}{:}\DUrole{w}{ }\DUrole{n}{float}\DUrole{w}{ }\DUrole{o}{=}\DUrole{w}{ }\DUrole{default_value}{0.9}}\sphinxparamcomma \sphinxparam{\DUrole{n}{low\_noise}\DUrole{p}{:}\DUrole{w}{ }\DUrole{n}{float}\DUrole{w}{ }\DUrole{o}{=}\DUrole{w}{ }\DUrole{default_value}{0.1}}\sphinxparamcomma \sphinxparam{\DUrole{n}{high\_noise}\DUrole{p}{:}\DUrole{w}{ }\DUrole{n}{float}\DUrole{w}{ }\DUrole{o}{=}\DUrole{w}{ }\DUrole{default_value}{2.0}}}{}
\pysigstopsignatures
\sphinxAtStartPar
Parameters for a particle filter fusion for quantum sensors.
This contains the arguments used for applying particle filtering for
position fixing with IMU quantum sensors. This namely contains the
number of particles to use and the expected errors.

\end{fulllineitems}

\index{QuantumPFFusion (class in qnav.quantum.particle\_filter)@\spxentry{QuantumPFFusion}\spxextra{class in qnav.quantum.particle\_filter}}

\begin{fulllineitems}
\phantomsection\label{\detokenize{modules/quantum/particle_filter:qnav.quantum.particle_filter.QuantumPFFusion}}
\pysigstartsignatures
\pysiglinewithargsret{\sphinxbfcode{\sphinxupquote{class\DUrole{w}{ }}}\sphinxcode{\sphinxupquote{qnav.quantum.particle\_filter.}}\sphinxbfcode{\sphinxupquote{QuantumPFFusion}}}{\sphinxparam{\DUrole{n}{qs\_imu}\DUrole{p}{:}\DUrole{w}{ }\DUrole{n}{{\hyperref[\detokenize{modules/quantum/base:qnav.quantum.base.ConceptQuantumImu}]{\sphinxcrossref{ConceptQuantumImu}}}}}\sphinxparamcomma \sphinxparam{\DUrole{n}{shared\_accelerometer}\DUrole{p}{:}\DUrole{w}{ }\DUrole{n}{{\hyperref[\detokenize{modules/measurement/accelerometer:qnav.measurement.accelerometer.Accelerometer}]{\sphinxcrossref{Accelerometer}}}}}\sphinxparamcomma \sphinxparam{\DUrole{n}{shared\_gyroscope}\DUrole{p}{:}\DUrole{w}{ }\DUrole{n}{{\hyperref[\detokenize{modules/measurement/gyroscope:qnav.measurement.gyroscope.Gyroscope}]{\sphinxcrossref{Gyroscope}}}}}\sphinxparamcomma \sphinxparam{\DUrole{n}{est\_qs\_imu\_axis}\DUrole{p}{:}\DUrole{w}{ }\DUrole{n}{{\hyperref[\detokenize{modules/measurement/platform:qnav.measurement.platform.SensorAxis}]{\sphinxcrossref{SensorAxis}}}}\DUrole{w}{ }\DUrole{o}{=}\DUrole{w}{ }\DUrole{default_value}{None}}\sphinxparamcomma \sphinxparam{\DUrole{n}{est\_acc\_axis}\DUrole{p}{:}\DUrole{w}{ }\DUrole{n}{{\hyperref[\detokenize{modules/measurement/platform:qnav.measurement.platform.SensorAxis}]{\sphinxcrossref{SensorAxis}}}}\DUrole{w}{ }\DUrole{o}{=}\DUrole{w}{ }\DUrole{default_value}{None}}\sphinxparamcomma \sphinxparam{\DUrole{n}{est\_gyro\_axis}\DUrole{p}{:}\DUrole{w}{ }\DUrole{n}{{\hyperref[\detokenize{modules/measurement/platform:qnav.measurement.platform.SensorAxis}]{\sphinxcrossref{SensorAxis}}}}\DUrole{w}{ }\DUrole{o}{=}\DUrole{w}{ }\DUrole{default_value}{None}}\sphinxparamcomma \sphinxparam{\DUrole{n}{filter\_params}\DUrole{p}{:}\DUrole{w}{ }\DUrole{n}{{\hyperref[\detokenize{modules/quantum/particle_filter:qnav.quantum.particle_filter.ParticleFilterParams}]{\sphinxcrossref{ParticleFilterParams}}}}\DUrole{w}{ }\DUrole{o}{=}\DUrole{w}{ }\DUrole{default_value}{ParticleFilterParams(num\_particles=200, sigma\_accelerometer=0.01, sigma\_gyroscope=0.01, sigma\_misalignment=0.01, rng\_seed=None, low\_noise\_prob=0.9, low\_noise=0.1, high\_noise=2.0)}}}{}
\pysigstopsignatures
\sphinxAtStartPar
Quantum sensor fusion with particle filter estimation.
This provides position correction by using a particle filter approach for
estimating the bias errors within the class IMU sensors. This extends
from the base class, overriding the relevant methods.

\end{fulllineitems}

\index{weighted\_cov() (in module qnav.quantum.particle\_filter)@\spxentry{weighted\_cov()}\spxextra{in module qnav.quantum.particle\_filter}}

\begin{fulllineitems}
\phantomsection\label{\detokenize{modules/quantum/particle_filter:qnav.quantum.particle_filter.weighted_cov}}
\pysigstartsignatures
\pysiglinewithargsret{\sphinxcode{\sphinxupquote{qnav.quantum.particle\_filter.}}\sphinxbfcode{\sphinxupquote{weighted\_cov}}}{\sphinxparam{\DUrole{n}{y}\DUrole{p}{:}\DUrole{w}{ }\DUrole{n}{ndarray}}\sphinxparamcomma \sphinxparam{\DUrole{n}{w}\DUrole{p}{:}\DUrole{w}{ }\DUrole{n}{ndarray}}}{}
\pysigstopsignatures
\sphinxAtStartPar
This function calculates a Weighted Covariance Matrix.
Given a matrix of observations (rows) and variables (columns) along with
an array of corresponding weights, this function returns a symmetric
matrix of weighted covariances (a positive semi\sphinxhyphen{}definite matrix i.e. all
its eigenvalues are non\sphinxhyphen{}negative).
\begin{quote}\begin{description}
\sphinxlineitem{Parameters}\begin{itemize}
\item {} 
\sphinxAtStartPar
\sphinxstyleliteralstrong{\sphinxupquote{y}} (\sphinxstyleliteralemphasis{\sphinxupquote{NumPy Array}}\sphinxstyleliteralemphasis{\sphinxupquote{ (}}\sphinxstyleliteralemphasis{\sphinxupquote{T\sphinxhyphen{}by\sphinxhyphen{}N elements}}\sphinxstyleliteralemphasis{\sphinxupquote{)}}) \textendash{} A T\sphinxhyphen{}by\sphinxhyphen{}N matrix of observations (rows) and variables (columns).

\item {} 
\sphinxAtStartPar
\sphinxstyleliteralstrong{\sphinxupquote{w}} (\sphinxstyleliteralemphasis{\sphinxupquote{NumPy Array}}\sphinxstyleliteralemphasis{\sphinxupquote{ (}}\sphinxstyleliteralemphasis{\sphinxupquote{T elements}}\sphinxstyleliteralemphasis{\sphinxupquote{)}}) \textendash{} A T\sphinxhyphen{}element array of weights for the observations.

\end{itemize}

\sphinxlineitem{Returns}
\sphinxAtStartPar
A symmetric matrix of weighted covariances.

\sphinxlineitem{Return type}
\sphinxAtStartPar
NumPy Array

\end{description}\end{quote}

\end{fulllineitems}


\sphinxstepscope


\subsubsection{qnav.quantum.misalignment}
\label{\detokenize{modules/quantum/misalignment:module-qnav.quantum.misalignment}}\label{\detokenize{modules/quantum/misalignment:qnav-quantum-misalignment}}\label{\detokenize{modules/quantum/misalignment::doc}}\index{module@\spxentry{module}!qnav.quantum.misalignment@\spxentry{qnav.quantum.misalignment}}\index{qnav.quantum.misalignment@\spxentry{qnav.quantum.misalignment}!module@\spxentry{module}}

\paragraph{misalignment.py}
\label{\detokenize{modules/quantum/misalignment:misalignment-py}}\begin{quote}\begin{description}
\sphinxlineitem{summary}
\sphinxAtStartPar
Extended implementations of Quantum Sensor Fusion for axis misalignment.
This module contains fusion methods for quantum sensors that, in addition
to position fixing, also correct axis misalignment, resulting in errors
due to non\sphinxhyphen{}matching actual and estimated sensor axis.

\sphinxlineitem{author}
\sphinxAtStartPar
Michael Wright \sphinxhyphen{} \sphinxhref{mailto:mjwright@liverpool.ac.uk}{mjwright@liverpool.ac.uk}

\sphinxlineitem{version}
\sphinxAtStartPar
1.0.0 \sphinxhyphen{} First full implementation of class.

\end{description}\end{quote}
\index{QuantumPFMFusion (class in qnav.quantum.misalignment)@\spxentry{QuantumPFMFusion}\spxextra{class in qnav.quantum.misalignment}}

\begin{fulllineitems}
\phantomsection\label{\detokenize{modules/quantum/misalignment:qnav.quantum.misalignment.QuantumPFMFusion}}
\pysigstartsignatures
\pysiglinewithargsret{\sphinxbfcode{\sphinxupquote{class\DUrole{w}{ }}}\sphinxcode{\sphinxupquote{qnav.quantum.misalignment.}}\sphinxbfcode{\sphinxupquote{QuantumPFMFusion}}}{\sphinxparam{\DUrole{n}{qs\_imu}\DUrole{p}{:}\DUrole{w}{ }\DUrole{n}{{\hyperref[\detokenize{modules/quantum/base:qnav.quantum.base.ConceptQuantumImu}]{\sphinxcrossref{ConceptQuantumImu}}}}}\sphinxparamcomma \sphinxparam{\DUrole{n}{shared\_accelerometer}\DUrole{p}{:}\DUrole{w}{ }\DUrole{n}{{\hyperref[\detokenize{modules/measurement/accelerometer:qnav.measurement.accelerometer.Accelerometer}]{\sphinxcrossref{Accelerometer}}}}}\sphinxparamcomma \sphinxparam{\DUrole{n}{shared\_gyroscope}\DUrole{p}{:}\DUrole{w}{ }\DUrole{n}{{\hyperref[\detokenize{modules/measurement/gyroscope:qnav.measurement.gyroscope.Gyroscope}]{\sphinxcrossref{Gyroscope}}}}}\sphinxparamcomma \sphinxparam{\DUrole{n}{est\_qs\_imu\_axis}\DUrole{p}{:}\DUrole{w}{ }\DUrole{n}{{\hyperref[\detokenize{modules/measurement/platform:qnav.measurement.platform.SensorAxis}]{\sphinxcrossref{SensorAxis}}}}\DUrole{w}{ }\DUrole{o}{=}\DUrole{w}{ }\DUrole{default_value}{None}}\sphinxparamcomma \sphinxparam{\DUrole{n}{est\_acc\_axis}\DUrole{p}{:}\DUrole{w}{ }\DUrole{n}{{\hyperref[\detokenize{modules/measurement/platform:qnav.measurement.platform.SensorAxis}]{\sphinxcrossref{SensorAxis}}}}\DUrole{w}{ }\DUrole{o}{=}\DUrole{w}{ }\DUrole{default_value}{None}}\sphinxparamcomma \sphinxparam{\DUrole{n}{est\_gyro\_axis}\DUrole{p}{:}\DUrole{w}{ }\DUrole{n}{{\hyperref[\detokenize{modules/measurement/platform:qnav.measurement.platform.SensorAxis}]{\sphinxcrossref{SensorAxis}}}}\DUrole{w}{ }\DUrole{o}{=}\DUrole{w}{ }\DUrole{default_value}{None}}\sphinxparamcomma \sphinxparam{\DUrole{n}{filter\_params}\DUrole{p}{:}\DUrole{w}{ }\DUrole{n}{{\hyperref[\detokenize{modules/quantum/particle_filter:qnav.quantum.particle_filter.ParticleFilterParams}]{\sphinxcrossref{ParticleFilterParams}}}}\DUrole{w}{ }\DUrole{o}{=}\DUrole{w}{ }\DUrole{default_value}{ParticleFilterParams(num\_particles=200, sigma\_accelerometer=0.01, sigma\_gyroscope=0.01, sigma\_misalignment=0.01, rng\_seed=None, low\_noise\_prob=0.9, low\_noise=0.1, high\_noise=2.0)}}}{}
\pysigstopsignatures
\sphinxAtStartPar
Quantum sensor fusion with particle filter estimation for axis misalignment.
This provides position and axis misalignment correction by using a particle
filter approach for estimating the bias errors within the class IMU sensors.
This extends from the base class, overriding the relevant methods.

\end{fulllineitems}


\sphinxstepscope


\subsubsection{qnav.quantum.gravity\_gradient}
\label{\detokenize{modules/quantum/gravity_gradient:module-qnav.quantum.gravity_gradient}}\label{\detokenize{modules/quantum/gravity_gradient:qnav-quantum-gravity-gradient}}\label{\detokenize{modules/quantum/gravity_gradient::doc}}\index{module@\spxentry{module}!qnav.quantum.gravity\_gradient@\spxentry{qnav.quantum.gravity\_gradient}}\index{qnav.quantum.gravity\_gradient@\spxentry{qnav.quantum.gravity\_gradient}!module@\spxentry{module}}

\paragraph{gravity\_gradient.py}
\label{\detokenize{modules/quantum/gravity_gradient:gravity-gradient-py}}\begin{quote}\begin{description}
\sphinxlineitem{summary}
\sphinxAtStartPar
Provides implementation for Quantum gravity gradient sensing and fusion.
Implementation is mainly based on Peters et al, ‘High precision gravity
measurements using atom interferometry’, Metrologia, Vol.38,
pp.25\sphinxhyphen{}61 (2001).

\sphinxlineitem{authors}
\begin{DUlineblock}{0em}
\item[] Prof. Jason Ralph \sphinxhyphen{} \sphinxhref{mailto:jfralph@liverpool.ac.uk}{jfralph@liverpool.ac.uk}
\item[] Michael Wright \sphinxhyphen{} \sphinxhref{mailto:mjwright@liverpool.ac.uk}{mjwright@liverpool.ac.uk}
\end{DUlineblock}

\sphinxlineitem{version}
\sphinxAtStartPar
1.0.0 \sphinxhyphen{} First full implementation of class.

\end{description}\end{quote}
\index{GravityGradInterferometer (class in qnav.quantum.gravity\_gradient)@\spxentry{GravityGradInterferometer}\spxextra{class in qnav.quantum.gravity\_gradient}}

\begin{fulllineitems}
\phantomsection\label{\detokenize{modules/quantum/gravity_gradient:qnav.quantum.gravity_gradient.GravityGradInterferometer}}
\pysigstartsignatures
\pysiglinewithargsret{\sphinxbfcode{\sphinxupquote{class\DUrole{w}{ }}}\sphinxcode{\sphinxupquote{qnav.quantum.gravity\_gradient.}}\sphinxbfcode{\sphinxupquote{GravityGradInterferometer}}}{\sphinxparam{\DUrole{n}{num\_of\_atoms}\DUrole{p}{:}\DUrole{w}{ }\DUrole{n}{int}\DUrole{w}{ }\DUrole{o}{=}\DUrole{w}{ }\DUrole{default_value}{200000}}\sphinxparamcomma \sphinxparam{\DUrole{n}{atom\_mass}\DUrole{p}{:}\DUrole{w}{ }\DUrole{n}{float}\DUrole{w}{ }\DUrole{o}{=}\DUrole{w}{ }\DUrole{default_value}{1.4431608262e\sphinxhyphen{}25}}\sphinxparamcomma \sphinxparam{\DUrole{n}{recoil\_velocity}\DUrole{p}{:}\DUrole{w}{ }\DUrole{n}{float}\DUrole{w}{ }\DUrole{o}{=}\DUrole{w}{ }\DUrole{default_value}{0.0058845}}\sphinxparamcomma \sphinxparam{\DUrole{n}{stationary\_velocity}\DUrole{p}{:}\DUrole{w}{ }\DUrole{n}{float}\DUrole{w}{ }\DUrole{o}{=}\DUrole{w}{ }\DUrole{default_value}{1.57}}\sphinxparamcomma \sphinxparam{\DUrole{n}{sensor\_z\_top}\DUrole{p}{:}\DUrole{w}{ }\DUrole{n}{float}\DUrole{w}{ }\DUrole{o}{=}\DUrole{w}{ }\DUrole{default_value}{0.0}}\sphinxparamcomma \sphinxparam{\DUrole{n}{sensor\_z\_bottom}\DUrole{p}{:}\DUrole{w}{ }\DUrole{n}{float}\DUrole{w}{ }\DUrole{o}{=}\DUrole{w}{ }\DUrole{default_value}{\sphinxhyphen{}1.0}}\sphinxparamcomma \sphinxparam{\DUrole{n}{eta}\DUrole{p}{:}\DUrole{w}{ }\DUrole{n}{float}\DUrole{w}{ }\DUrole{o}{=}\DUrole{w}{ }\DUrole{default_value}{0.25}}\sphinxparamcomma \sphinxparam{\DUrole{n}{time\_pulse}\DUrole{p}{:}\DUrole{w}{ }\DUrole{n}{float}\DUrole{w}{ }\DUrole{o}{=}\DUrole{w}{ }\DUrole{default_value}{0.2}}\sphinxparamcomma \sphinxparam{\DUrole{n}{sigma\_phi}\DUrole{p}{:}\DUrole{w}{ }\DUrole{n}{float}\DUrole{w}{ }\DUrole{o}{=}\DUrole{w}{ }\DUrole{default_value}{0.025}}\sphinxparamcomma \sphinxparam{\DUrole{n}{sigma\_n}\DUrole{p}{:}\DUrole{w}{ }\DUrole{n}{float}\DUrole{w}{ }\DUrole{o}{=}\DUrole{w}{ }\DUrole{default_value}{1.0}}\sphinxparamcomma \sphinxparam{\DUrole{n}{beam\_width}\DUrole{p}{:}\DUrole{w}{ }\DUrole{n}{float}\DUrole{w}{ }\DUrole{o}{=}\DUrole{w}{ }\DUrole{default_value}{0.01}}\sphinxparamcomma \sphinxparam{\DUrole{n}{sensor\_length}\DUrole{p}{:}\DUrole{w}{ }\DUrole{n}{float}\DUrole{w}{ }\DUrole{o}{=}\DUrole{w}{ }\DUrole{default_value}{1.5}}\sphinxparamcomma \sphinxparam{\DUrole{n}{rand\_failure\_prob}\DUrole{p}{:}\DUrole{w}{ }\DUrole{n}{float}\DUrole{w}{ }\DUrole{o}{=}\DUrole{w}{ }\DUrole{default_value}{0.0}}}{}
\pysigstopsignatures
\sphinxAtStartPar
Parameters used by the Cold Atom Gravity Gradient Interferometer.
These properties describe the qualities of an ultra\sphinxhyphen{}cold atom
interferometer, which is used for capturing gravity gradient
estimations from the local environment.
\index{k\_eff (qnav.quantum.gravity\_gradient.GravityGradInterferometer property)@\spxentry{k\_eff}\spxextra{qnav.quantum.gravity\_gradient.GravityGradInterferometer property}}

\begin{fulllineitems}
\phantomsection\label{\detokenize{modules/quantum/gravity_gradient:qnav.quantum.gravity_gradient.GravityGradInterferometer.k_eff}}
\pysigstartsignatures
\pysigline{\sphinxbfcode{\sphinxupquote{property\DUrole{w}{ }}}\sphinxbfcode{\sphinxupquote{k\_eff}}\sphinxbfcode{\sphinxupquote{\DUrole{p}{:}\DUrole{w}{ }float}}}
\pysigstopsignatures
\sphinxAtStartPar
Effective wave number, the number of wave cycles per unit distance.
A property provided for a value commonly calculated.
\begin{quote}\begin{description}
\sphinxlineitem{Returns}
\sphinxAtStartPar
A measure of the spatial frequency of a wave.

\sphinxlineitem{Return type}
\sphinxAtStartPar
float

\end{description}\end{quote}

\end{fulllineitems}


\end{fulllineitems}

\index{GravityGradiometer (class in qnav.quantum.gravity\_gradient)@\spxentry{GravityGradiometer}\spxextra{class in qnav.quantum.gravity\_gradient}}

\begin{fulllineitems}
\phantomsection\label{\detokenize{modules/quantum/gravity_gradient:qnav.quantum.gravity_gradient.GravityGradiometer}}
\pysigstartsignatures
\pysiglinewithargsret{\sphinxbfcode{\sphinxupquote{class\DUrole{w}{ }}}\sphinxcode{\sphinxupquote{qnav.quantum.gravity\_gradient.}}\sphinxbfcode{\sphinxupquote{GravityGradiometer}}}{\sphinxparam{\DUrole{n}{full\_frequency: float}}\sphinxparamcomma \sphinxparam{\DUrole{n}{imu\_frequency: float}}\sphinxparamcomma \sphinxparam{\DUrole{n}{accelerometer\_errors: \textasciitilde{}qnav.measurement.error\_properties.ErrorProperties}}\sphinxparamcomma \sphinxparam{\DUrole{n}{interferometer: \textasciitilde{}qnav.quantum.gravity\_gradient.GravityGradInterferometer}}\sphinxparamcomma \sphinxparam{\DUrole{n}{gravity\_model: \textasciitilde{}qnav.gravity.base.GravityModel}}\sphinxparamcomma \sphinxparam{\DUrole{n}{sensor\_axis: \textasciitilde{}qnav.measurement.platform.SensorAxis = \textless{}qnav.measurement.platform.SensorAxis object\textgreater{}}}\sphinxparamcomma \sphinxparam{\DUrole{n}{duty\_cycle: float = 0.5}}\sphinxparamcomma \sphinxparam{\DUrole{n}{start\_time: float = 0.0}}\sphinxparamcomma \sphinxparam{\DUrole{n}{rand\_seed: int = None}}}{}
\pysigstopsignatures
\sphinxAtStartPar
Quantum Gravity Gradiometer Sensor for estimating vertical gravity gradient.
A Quantum sensor that features two ultra\sphinxhyphen{}cold atom accelerometers, measuring
the change in the vertical gravity gradient.
\index{take\_measurement() (qnav.quantum.gravity\_gradient.GravityGradiometer method)@\spxentry{take\_measurement()}\spxextra{qnav.quantum.gravity\_gradient.GravityGradiometer method}}

\begin{fulllineitems}
\phantomsection\label{\detokenize{modules/quantum/gravity_gradient:qnav.quantum.gravity_gradient.GravityGradiometer.take_measurement}}
\pysigstartsignatures
\pysiglinewithargsret{\sphinxbfcode{\sphinxupquote{take\_measurement}}}{\sphinxparam{\DUrole{n}{est\_time}\DUrole{p}{:}\DUrole{w}{ }\DUrole{n}{float}}\sphinxparamcomma \sphinxparam{\DUrole{n}{ground\_truth}\DUrole{p}{:}\DUrole{w}{ }\DUrole{n}{{\hyperref[\detokenize{modules/waypoints/trajectory:qnav.waypoints.trajectory.GroundTruth}]{\sphinxcrossref{GroundTruth}}}}}}{}
\pysigstopsignatures
\sphinxAtStartPar
Samples from the ground truth to collect data to produce measurements.
Using the current ground truth, captures the
\begin{quote}\begin{description}
\sphinxlineitem{Parameters}\begin{itemize}
\item {} 
\sphinxAtStartPar
\sphinxstyleliteralstrong{\sphinxupquote{est\_time}} (\sphinxstyleliteralemphasis{\sphinxupquote{float}}) \textendash{} The estimated time the measurement of the sensor
is expected to be captured at (in seconds).

\item {} 
\sphinxAtStartPar
\sphinxstyleliteralstrong{\sphinxupquote{ground\_truth}} ({\hyperref[\detokenize{modules/waypoints/trajectory:qnav.waypoints.trajectory.GroundTruth}]{\sphinxcrossref{\sphinxstyleliteralemphasis{\sphinxupquote{GroundTruth}}}}}) \textendash{} The corresponding ground truth data.

\end{itemize}

\end{description}\end{quote}

\end{fulllineitems}

\index{reset() (qnav.quantum.gravity\_gradient.GravityGradiometer method)@\spxentry{reset()}\spxextra{qnav.quantum.gravity\_gradient.GravityGradiometer method}}

\begin{fulllineitems}
\phantomsection\label{\detokenize{modules/quantum/gravity_gradient:qnav.quantum.gravity_gradient.GravityGradiometer.reset}}
\pysigstartsignatures
\pysiglinewithargsret{\sphinxbfcode{\sphinxupquote{reset}}}{}{}
\pysigstopsignatures
\sphinxAtStartPar
Resets the sensor, clearing all previous measurements.
Normally used directly after fusion.

\end{fulllineitems}

\index{update() (qnav.quantum.gravity\_gradient.GravityGradiometer method)@\spxentry{update()}\spxextra{qnav.quantum.gravity\_gradient.GravityGradiometer method}}

\begin{fulllineitems}
\phantomsection\label{\detokenize{modules/quantum/gravity_gradient:qnav.quantum.gravity_gradient.GravityGradiometer.update}}
\pysigstartsignatures
\pysiglinewithargsret{\sphinxbfcode{\sphinxupquote{update}}}{\sphinxparam{\DUrole{n}{time\_sec}\DUrole{p}{:}\DUrole{w}{ }\DUrole{n}{float}\DUrole{w}{ }\DUrole{o}{=}\DUrole{w}{ }\DUrole{default_value}{None}}}{}
\pysigstopsignatures
\sphinxAtStartPar
Called on demand to update the sensor’s state.
Typically called after measurements have been taken. When called
this increases the next measurement time and propagates internal
states, to the time difference.
\begin{quote}\begin{description}
\sphinxlineitem{Parameters}
\sphinxAtStartPar
\sphinxstyleliteralstrong{\sphinxupquote{time\_sec}} (\sphinxstyleliteralemphasis{\sphinxupquote{float}}) \textendash{} The current simulation time in seconds.

\end{description}\end{quote}

\end{fulllineitems}

\index{last\_measurement (qnav.quantum.gravity\_gradient.GravityGradiometer property)@\spxentry{last\_measurement}\spxextra{qnav.quantum.gravity\_gradient.GravityGradiometer property}}

\begin{fulllineitems}
\phantomsection\label{\detokenize{modules/quantum/gravity_gradient:qnav.quantum.gravity_gradient.GravityGradiometer.last_measurement}}
\pysigstartsignatures
\pysigline{\sphinxbfcode{\sphinxupquote{property\DUrole{w}{ }}}\sphinxbfcode{\sphinxupquote{last\_measurement}}\sphinxbfcode{\sphinxupquote{\DUrole{p}{:}\DUrole{w}{ }tuple\DUrole{p}{{[}}float\DUrole{p}{,}\DUrole{w}{ }float\DUrole{p}{{]}}}}}
\pysigstopsignatures
\sphinxAtStartPar
Returns the latest measurement produced from the sensor.
\begin{quote}\begin{description}
\sphinxlineitem{Returns}
\sphinxAtStartPar
The last produced gravity gradient signal.

\sphinxlineitem{Return type}
\sphinxAtStartPar
tuple{[}float, float{]}

\end{description}\end{quote}

\end{fulllineitems}

\index{num\_steps (qnav.quantum.gravity\_gradient.GravityGradiometer property)@\spxentry{num\_steps}\spxextra{qnav.quantum.gravity\_gradient.GravityGradiometer property}}

\begin{fulllineitems}
\phantomsection\label{\detokenize{modules/quantum/gravity_gradient:qnav.quantum.gravity_gradient.GravityGradiometer.num_steps}}
\pysigstartsignatures
\pysigline{\sphinxbfcode{\sphinxupquote{property\DUrole{w}{ }}}\sphinxbfcode{\sphinxupquote{num\_steps}}\sphinxbfcode{\sphinxupquote{\DUrole{p}{:}\DUrole{w}{ }int}}}
\pysigstopsignatures
\sphinxAtStartPar
Returns the total number of internal measurement steps required.
Used namely for sharing array sizes.
\begin{quote}\begin{description}
\sphinxlineitem{Returns}
\sphinxAtStartPar
The number of internal measurement steps used.

\sphinxlineitem{Return type}
\sphinxAtStartPar
int

\end{description}\end{quote}

\end{fulllineitems}

\index{interferometer (qnav.quantum.gravity\_gradient.GravityGradiometer property)@\spxentry{interferometer}\spxextra{qnav.quantum.gravity\_gradient.GravityGradiometer property}}

\begin{fulllineitems}
\phantomsection\label{\detokenize{modules/quantum/gravity_gradient:qnav.quantum.gravity_gradient.GravityGradiometer.interferometer}}
\pysigstartsignatures
\pysigline{\sphinxbfcode{\sphinxupquote{property\DUrole{w}{ }}}\sphinxbfcode{\sphinxupquote{interferometer}}\sphinxbfcode{\sphinxupquote{\DUrole{p}{:}\DUrole{w}{ }{\hyperref[\detokenize{modules/quantum/gravity_gradient:qnav.quantum.gravity_gradient.GravityGradInterferometer}]{\sphinxcrossref{GravityGradInterferometer}}}}}}
\pysigstopsignatures
\sphinxAtStartPar
Returns the interferometer properties used by the sensor.
\begin{quote}\begin{description}
\sphinxlineitem{Returns}
\sphinxAtStartPar
Interferometer properties.

\sphinxlineitem{Return type}
\sphinxAtStartPar
{\hyperref[\detokenize{modules/quantum/gravity_gradient:qnav.quantum.gravity_gradient.GravityGradInterferometer}]{\sphinxcrossref{GravityGradInterferometer}}}

\end{description}\end{quote}

\end{fulllineitems}

\index{gravity\_model (qnav.quantum.gravity\_gradient.GravityGradiometer property)@\spxentry{gravity\_model}\spxextra{qnav.quantum.gravity\_gradient.GravityGradiometer property}}

\begin{fulllineitems}
\phantomsection\label{\detokenize{modules/quantum/gravity_gradient:qnav.quantum.gravity_gradient.GravityGradiometer.gravity_model}}
\pysigstartsignatures
\pysigline{\sphinxbfcode{\sphinxupquote{property\DUrole{w}{ }}}\sphinxbfcode{\sphinxupquote{gravity\_model}}\sphinxbfcode{\sphinxupquote{\DUrole{p}{:}\DUrole{w}{ }{\hyperref[\detokenize{modules/gravity/base:qnav.gravity.base.GravityModel}]{\sphinxcrossref{GravityModel}}}}}}
\pysigstopsignatures
\sphinxAtStartPar
Returns the internal gravity model used by the sensor.
Can be accessed for testing or producing similar measurements.
\begin{quote}\begin{description}
\sphinxlineitem{Returns}
\sphinxAtStartPar
Gravity Model instance used by the sensor.

\sphinxlineitem{Return type}
\sphinxAtStartPar
{\hyperref[\detokenize{modules/gravity/base:qnav.gravity.base.GravityModel}]{\sphinxcrossref{GravityModel}}}

\end{description}\end{quote}

\end{fulllineitems}


\end{fulllineitems}

\index{ParticleFilterParams (class in qnav.quantum.gravity\_gradient)@\spxentry{ParticleFilterParams}\spxextra{class in qnav.quantum.gravity\_gradient}}

\begin{fulllineitems}
\phantomsection\label{\detokenize{modules/quantum/gravity_gradient:qnav.quantum.gravity_gradient.ParticleFilterParams}}
\pysigstartsignatures
\pysiglinewithargsret{\sphinxbfcode{\sphinxupquote{class\DUrole{w}{ }}}\sphinxcode{\sphinxupquote{qnav.quantum.gravity\_gradient.}}\sphinxbfcode{\sphinxupquote{ParticleFilterParams}}}{\sphinxparam{\DUrole{n}{num\_particles}\DUrole{p}{:}\DUrole{w}{ }\DUrole{n}{int}\DUrole{w}{ }\DUrole{o}{=}\DUrole{w}{ }\DUrole{default_value}{1000}}\sphinxparamcomma \sphinxparam{\DUrole{n}{init\_position\_sigma}\DUrole{p}{:}\DUrole{w}{ }\DUrole{n}{float}\DUrole{w}{ }\DUrole{o}{=}\DUrole{w}{ }\DUrole{default_value}{20.0}}\sphinxparamcomma \sphinxparam{\DUrole{n}{init\_velocity\_sigma}\DUrole{p}{:}\DUrole{w}{ }\DUrole{n}{float}\DUrole{w}{ }\DUrole{o}{=}\DUrole{w}{ }\DUrole{default_value}{1e\sphinxhyphen{}05}}\sphinxparamcomma \sphinxparam{\DUrole{n}{init\_attitude\_sigma}\DUrole{p}{:}\DUrole{w}{ }\DUrole{n}{float}\DUrole{w}{ }\DUrole{o}{=}\DUrole{w}{ }\DUrole{default_value}{1e\sphinxhyphen{}05}}\sphinxparamcomma \sphinxparam{\DUrole{n}{weight\_sigma}\DUrole{p}{:}\DUrole{w}{ }\DUrole{n}{float}\DUrole{w}{ }\DUrole{o}{=}\DUrole{w}{ }\DUrole{default_value}{0.00126}}\sphinxparamcomma \sphinxparam{\DUrole{n}{alpha}\DUrole{p}{:}\DUrole{w}{ }\DUrole{n}{float}\DUrole{w}{ }\DUrole{o}{=}\DUrole{w}{ }\DUrole{default_value}{0.05}}\sphinxparamcomma \sphinxparam{\DUrole{n}{beta}\DUrole{p}{:}\DUrole{w}{ }\DUrole{n}{float}\DUrole{w}{ }\DUrole{o}{=}\DUrole{w}{ }\DUrole{default_value}{0.05}}\sphinxparamcomma \sphinxparam{\DUrole{n}{gamma}\DUrole{p}{:}\DUrole{w}{ }\DUrole{n}{float}\DUrole{w}{ }\DUrole{o}{=}\DUrole{w}{ }\DUrole{default_value}{0.05}}\sphinxparamcomma \sphinxparam{\DUrole{n}{position\_sigma}\DUrole{p}{:}\DUrole{w}{ }\DUrole{n}{float}\DUrole{w}{ }\DUrole{o}{=}\DUrole{w}{ }\DUrole{default_value}{2.0}}\sphinxparamcomma \sphinxparam{\DUrole{n}{velocity\_sigma}\DUrole{p}{:}\DUrole{w}{ }\DUrole{n}{float}\DUrole{w}{ }\DUrole{o}{=}\DUrole{w}{ }\DUrole{default_value}{0.005}}\sphinxparamcomma \sphinxparam{\DUrole{n}{attitude\_sigma}\DUrole{p}{:}\DUrole{w}{ }\DUrole{n}{float}\DUrole{w}{ }\DUrole{o}{=}\DUrole{w}{ }\DUrole{default_value}{0.001}}\sphinxparamcomma \sphinxparam{\DUrole{n}{ellipse\_intervals}\DUrole{p}{:}\DUrole{w}{ }\DUrole{n}{int}\DUrole{w}{ }\DUrole{o}{=}\DUrole{w}{ }\DUrole{default_value}{361}}\sphinxparamcomma \sphinxparam{\DUrole{n}{ellipse\_iterations}\DUrole{p}{:}\DUrole{w}{ }\DUrole{n}{int}\DUrole{w}{ }\DUrole{o}{=}\DUrole{w}{ }\DUrole{default_value}{5}}}{}
\pysigstopsignatures
\sphinxAtStartPar
Particle Filter Properties for the gravity gradient fusion methods.
A list of properties relating to the particle filter processing.

\end{fulllineitems}

\index{GravityGradientPF (class in qnav.quantum.gravity\_gradient)@\spxentry{GravityGradientPF}\spxextra{class in qnav.quantum.gravity\_gradient}}

\begin{fulllineitems}
\phantomsection\label{\detokenize{modules/quantum/gravity_gradient:qnav.quantum.gravity_gradient.GravityGradientPF}}
\pysigstartsignatures
\pysiglinewithargsret{\sphinxbfcode{\sphinxupquote{class\DUrole{w}{ }}}\sphinxcode{\sphinxupquote{qnav.quantum.gravity\_gradient.}}\sphinxbfcode{\sphinxupquote{GravityGradientPF}}}{\sphinxparam{\DUrole{n}{qs\_gravity\_grad}\DUrole{p}{:}\DUrole{w}{ }\DUrole{n}{{\hyperref[\detokenize{modules/quantum/gravity_gradient:qnav.quantum.gravity_gradient.GravityGradiometer}]{\sphinxcrossref{GravityGradiometer}}}}}\sphinxparamcomma \sphinxparam{\DUrole{n}{filter\_params}\DUrole{p}{:}\DUrole{w}{ }\DUrole{n}{{\hyperref[\detokenize{modules/quantum/gravity_gradient:qnav.quantum.gravity_gradient.ParticleFilterParams}]{\sphinxcrossref{ParticleFilterParams}}}}\DUrole{w}{ }\DUrole{o}{=}\DUrole{w}{ }\DUrole{default_value}{ParticleFilterParams(num\_particles=1000, init\_position\_sigma=20.0, init\_velocity\_sigma=1e\sphinxhyphen{}05, init\_attitude\_sigma=1e\sphinxhyphen{}05, weight\_sigma=0.00126, alpha=0.05, beta=0.05, gamma=0.05, position\_sigma=2.0, velocity\_sigma=0.005, attitude\_sigma=0.001, ellipse\_intervals=361, ellipse\_iterations=5)}}\sphinxparamcomma \sphinxparam{\DUrole{n}{est\_interferometer}\DUrole{p}{:}\DUrole{w}{ }\DUrole{n}{{\hyperref[\detokenize{modules/quantum/gravity_gradient:qnav.quantum.gravity_gradient.GravityGradInterferometer}]{\sphinxcrossref{GravityGradInterferometer}}}}\DUrole{w}{ }\DUrole{o}{=}\DUrole{w}{ }\DUrole{default_value}{None}}\sphinxparamcomma \sphinxparam{\DUrole{n}{rng\_seed}\DUrole{p}{:}\DUrole{w}{ }\DUrole{n}{int}\DUrole{w}{ }\DUrole{o}{=}\DUrole{w}{ }\DUrole{default_value}{None}}}{}
\pysigstopsignatures
\sphinxAtStartPar
Gravity Gradient Particle Filter Fusion.
Applies position correction using gravity gradient signals,
via a particle filter approach.
\index{perform\_fusion() (qnav.quantum.gravity\_gradient.GravityGradientPF method)@\spxentry{perform\_fusion()}\spxextra{qnav.quantum.gravity\_gradient.GravityGradientPF method}}

\begin{fulllineitems}
\phantomsection\label{\detokenize{modules/quantum/gravity_gradient:qnav.quantum.gravity_gradient.GravityGradientPF.perform_fusion}}
\pysigstartsignatures
\pysiglinewithargsret{\sphinxbfcode{\sphinxupquote{perform\_fusion}}}{\sphinxparam{\DUrole{n}{estimated\_state}\DUrole{p}{:}\DUrole{w}{ }\DUrole{n}{{\hyperref[\detokenize{modules/estimation/state:qnav.estimation.state.EstimatedState}]{\sphinxcrossref{EstimatedState}}}}}}{{ $\rightarrow$ None}}
\pysigstopsignatures
\sphinxAtStartPar
Performs the fusion with sensors and updates the estimated state.
When called performs fusion when gravity gradient sensor
measurements are ready. When so, the particle filter algorithms is
employed and the estimated state is updated.
\begin{quote}\begin{description}
\sphinxlineitem{Parameters}
\sphinxAtStartPar
\sphinxstyleliteralstrong{\sphinxupquote{estimated\_state}} ({\hyperref[\detokenize{modules/estimation/state:qnav.estimation.state.EstimatedState}]{\sphinxcrossref{\sphinxstyleliteralemphasis{\sphinxupquote{EstimatedState}}}}}) \textendash{} The current estimated states of the platform.
This is what the fusion method will update.

\end{description}\end{quote}

\end{fulllineitems}

\index{init\_particle\_filter() (qnav.quantum.gravity\_gradient.GravityGradientPF method)@\spxentry{init\_particle\_filter()}\spxextra{qnav.quantum.gravity\_gradient.GravityGradientPF method}}

\begin{fulllineitems}
\phantomsection\label{\detokenize{modules/quantum/gravity_gradient:qnav.quantum.gravity_gradient.GravityGradientPF.init_particle_filter}}
\pysigstartsignatures
\pysiglinewithargsret{\sphinxbfcode{\sphinxupquote{init\_particle\_filter}}}{\sphinxparam{\DUrole{n}{estimated\_state}\DUrole{p}{:}\DUrole{w}{ }\DUrole{n}{{\hyperref[\detokenize{modules/estimation/state:qnav.estimation.state.EstimatedState}]{\sphinxcrossref{EstimatedState}}}}}}{{ $\rightarrow$ None}}
\pysigstopsignatures
\sphinxAtStartPar
Initialises the particle filter using current estimate.
This is called during the first iteration, setting the particles’
estimated position, velocity and attitude values based around the
current estimated state. The range of estimation is based on the
initial sigma values set in the filter properties for each.
\begin{quote}\begin{description}
\sphinxlineitem{Parameters}
\sphinxAtStartPar
\sphinxstyleliteralstrong{\sphinxupquote{estimated\_state}} ({\hyperref[\detokenize{modules/estimation/state:qnav.estimation.state.EstimatedState}]{\sphinxcrossref{\sphinxstyleliteralemphasis{\sphinxupquote{EstimatedState}}}}}) \textendash{} The current estimated states of the platform.

\end{description}\end{quote}

\end{fulllineitems}

\index{reset() (qnav.quantum.gravity\_gradient.GravityGradientPF method)@\spxentry{reset()}\spxextra{qnav.quantum.gravity\_gradient.GravityGradientPF method}}

\begin{fulllineitems}
\phantomsection\label{\detokenize{modules/quantum/gravity_gradient:qnav.quantum.gravity_gradient.GravityGradientPF.reset}}
\pysigstartsignatures
\pysiglinewithargsret{\sphinxbfcode{\sphinxupquote{reset}}}{}{}
\pysigstopsignatures
\sphinxAtStartPar
Resets the state of the fusion, clearing internal temporary states.
When call the fusion will return to expecting the first measurement,
then continuing from there.

\end{fulllineitems}


\end{fulllineitems}


\sphinxstepscope


\subsubsection{qnav.quantum.dummy}
\label{\detokenize{modules/quantum/dummy:module-qnav.quantum.dummy}}\label{\detokenize{modules/quantum/dummy:qnav-quantum-dummy}}\label{\detokenize{modules/quantum/dummy::doc}}\index{module@\spxentry{module}!qnav.quantum.dummy@\spxentry{qnav.quantum.dummy}}\index{qnav.quantum.dummy@\spxentry{qnav.quantum.dummy}!module@\spxentry{module}}

\paragraph{dummy\_sensors.py}
\label{\detokenize{modules/quantum/dummy:dummy-sensors-py}}\begin{quote}\begin{description}
\sphinxlineitem{summary}
\sphinxAtStartPar
Presents dummy sensors used for re\sphinxhyphen{}applying fusion algorithms.
This module contains dummy sensor models used for re\sphinxhyphen{}applying previous
fusion algorithms with corrected measurements. These dummy sensors allow
internal states to be exposed and edited. These do not feature involvement
of real ground truth data, errors states or “real” measurements.

\sphinxlineitem{authors}
\begin{DUlineblock}{0em}
\item[] Prof. Jason Ralph \sphinxhyphen{} \sphinxhref{mailto:jfralph@liverpool.ac.uk}{jfralph@liverpool.ac.uk}
\item[] Michael Wright \sphinxhyphen{} \sphinxhref{mailto:mjwright@liverpool.ac.uk}{mjwright@liverpool.ac.uk}
\end{DUlineblock}

\sphinxlineitem{version}
\sphinxAtStartPar
1.0.0 \sphinxhyphen{} First full implementation of module.

\end{description}\end{quote}
\index{DummySensor (class in qnav.quantum.dummy)@\spxentry{DummySensor}\spxextra{class in qnav.quantum.dummy}}

\begin{fulllineitems}
\phantomsection\label{\detokenize{modules/quantum/dummy:qnav.quantum.dummy.DummySensor}}
\pysigstartsignatures
\pysiglinewithargsret{\sphinxbfcode{\sphinxupquote{class\DUrole{w}{ }}}\sphinxcode{\sphinxupquote{qnav.quantum.dummy.}}\sphinxbfcode{\sphinxupquote{DummySensor}}}{\sphinxparam{\DUrole{n}{frequency}\DUrole{p}{:}\DUrole{w}{ }\DUrole{n}{float}}\sphinxparamcomma \sphinxparam{\DUrole{n}{start\_time}\DUrole{p}{:}\DUrole{w}{ }\DUrole{n}{float}\DUrole{w}{ }\DUrole{o}{=}\DUrole{w}{ }\DUrole{default_value}{0.0}}\sphinxparamcomma \sphinxparam{\DUrole{n}{rand\_seed}\DUrole{p}{:}\DUrole{w}{ }\DUrole{n}{int}\DUrole{w}{ }\DUrole{o}{=}\DUrole{w}{ }\DUrole{default_value}{None}}}{}
\pysigstopsignatures
\sphinxAtStartPar
Abstract class for fake sensor to handle synthetic measurements.
Provides an abstract class for dummy sensor, which allows its
measurements to be fabricated. This is solely used for solutions
that require the re\sphinxhyphen{}use of fusion methods with updates/corrections
applied to previous measurements.
\index{last\_measurement (qnav.quantum.dummy.DummySensor property)@\spxentry{last\_measurement}\spxextra{qnav.quantum.dummy.DummySensor property}}

\begin{fulllineitems}
\phantomsection\label{\detokenize{modules/quantum/dummy:qnav.quantum.dummy.DummySensor.last_measurement}}
\pysigstartsignatures
\pysigline{\sphinxbfcode{\sphinxupquote{property\DUrole{w}{ }}}\sphinxbfcode{\sphinxupquote{last\_measurement}}\sphinxbfcode{\sphinxupquote{\DUrole{p}{:}\DUrole{w}{ }any}}}
\pysigstopsignatures
\sphinxAtStartPar
Returns the last measurement produced by the sensor.
\begin{quote}\begin{description}
\sphinxlineitem{Returns}
\sphinxAtStartPar
Last measurement value produced.

\sphinxlineitem{Return type}
\sphinxAtStartPar
any

\end{description}\end{quote}

\end{fulllineitems}

\index{last\_update (qnav.quantum.dummy.DummySensor property)@\spxentry{last\_update}\spxextra{qnav.quantum.dummy.DummySensor property}}

\begin{fulllineitems}
\phantomsection\label{\detokenize{modules/quantum/dummy:qnav.quantum.dummy.DummySensor.last_update}}
\pysigstartsignatures
\pysigline{\sphinxbfcode{\sphinxupquote{property\DUrole{w}{ }}}\sphinxbfcode{\sphinxupquote{last\_update}}\sphinxbfcode{\sphinxupquote{\DUrole{p}{:}\DUrole{w}{ }float}}}
\pysigstopsignatures
\sphinxAtStartPar
The time of the last measurement (in seconds).
\begin{quote}\begin{description}
\sphinxlineitem{Returns}
\sphinxAtStartPar
Time of last measurement in seconds.

\sphinxlineitem{Return type}
\sphinxAtStartPar
float

\end{description}\end{quote}

\end{fulllineitems}

\index{next\_update (qnav.quantum.dummy.DummySensor property)@\spxentry{next\_update}\spxextra{qnav.quantum.dummy.DummySensor property}}

\begin{fulllineitems}
\phantomsection\label{\detokenize{modules/quantum/dummy:qnav.quantum.dummy.DummySensor.next_update}}
\pysigstartsignatures
\pysigline{\sphinxbfcode{\sphinxupquote{property\DUrole{w}{ }}}\sphinxbfcode{\sphinxupquote{next\_update}}\sphinxbfcode{\sphinxupquote{\DUrole{p}{:}\DUrole{w}{ }float}}}
\pysigstopsignatures
\sphinxAtStartPar
The time of the next measurement (in seconds).
\begin{quote}\begin{description}
\sphinxlineitem{Returns}
\sphinxAtStartPar
Time of next measurement in seconds.

\sphinxlineitem{Return type}
\sphinxAtStartPar
float

\end{description}\end{quote}

\end{fulllineitems}


\end{fulllineitems}

\index{DummyAccelerometer (class in qnav.quantum.dummy)@\spxentry{DummyAccelerometer}\spxextra{class in qnav.quantum.dummy}}

\begin{fulllineitems}
\phantomsection\label{\detokenize{modules/quantum/dummy:qnav.quantum.dummy.DummyAccelerometer}}
\pysigstartsignatures
\pysiglinewithargsret{\sphinxbfcode{\sphinxupquote{class\DUrole{w}{ }}}\sphinxcode{\sphinxupquote{qnav.quantum.dummy.}}\sphinxbfcode{\sphinxupquote{DummyAccelerometer}}}{\sphinxparam{\DUrole{n}{accelerometer}\DUrole{p}{:}\DUrole{w}{ }\DUrole{n}{{\hyperref[\detokenize{modules/measurement/accelerometer:qnav.measurement.accelerometer.Accelerometer}]{\sphinxcrossref{Accelerometer}}}}}}{}
\pysigstopsignatures
\sphinxAtStartPar
Dummy Accelerometer sensor used to mimic accelerometer measurements.
This class is a dummy sensor that inherits functionalities for the
Accelerometer class. Essentially supports synthetic accelerometer
measurements to be used in fusion methods.

\end{fulllineitems}

\index{DummyGyroscope (class in qnav.quantum.dummy)@\spxentry{DummyGyroscope}\spxextra{class in qnav.quantum.dummy}}

\begin{fulllineitems}
\phantomsection\label{\detokenize{modules/quantum/dummy:qnav.quantum.dummy.DummyGyroscope}}
\pysigstartsignatures
\pysiglinewithargsret{\sphinxbfcode{\sphinxupquote{class\DUrole{w}{ }}}\sphinxcode{\sphinxupquote{qnav.quantum.dummy.}}\sphinxbfcode{\sphinxupquote{DummyGyroscope}}}{\sphinxparam{\DUrole{n}{gyroscope}\DUrole{p}{:}\DUrole{w}{ }\DUrole{n}{{\hyperref[\detokenize{modules/measurement/gyroscope:qnav.measurement.gyroscope.Gyroscope}]{\sphinxcrossref{Gyroscope}}}}}}{}
\pysigstopsignatures
\sphinxAtStartPar
Dummy Gyroscope sensor used to mimic gyroscope measurements.
This class is a dummy sensor that inherits functionalities for the
Gyroscope class. Essentially supports synthetic gyroscope measurements
to be used in fusion methods.

\end{fulllineitems}



\subsection{Simulation}
\label{\detokenize{usage/packages:simulation}}
\sphinxAtStartPar
Modules relating to simulation loop handling:

\sphinxstepscope


\subsubsection{qnav.simulation.simulation}
\label{\detokenize{modules/simulation/simulation:module-qnav.simulation.simulation}}\label{\detokenize{modules/simulation/simulation:qnav-simulation-simulation}}\label{\detokenize{modules/simulation/simulation::doc}}\index{module@\spxentry{module}!qnav.simulation.simulation@\spxentry{qnav.simulation.simulation}}\index{qnav.simulation.simulation@\spxentry{qnav.simulation.simulation}!module@\spxentry{module}}

\paragraph{simulation.py}
\label{\detokenize{modules/simulation/simulation:simulation-py}}\begin{quote}\begin{description}
\sphinxlineitem{summary}
\sphinxAtStartPar
The core harness responsible for coordinating simulations.
This module contains the core functionalities for managing simulations.
This mostly revolves around a core loop and the scheduling of sensor
time intervals, sensor measurements and data fusion.

\sphinxlineitem{author}
\sphinxAtStartPar
Michael Wright \sphinxhyphen{} \sphinxhref{mailto:mjwright@liverpool.ac.uk}{mjwright@liverpool.ac.uk}

\sphinxlineitem{version}
\sphinxAtStartPar
1.0.0 \sphinxhyphen{} First full implementation.

\end{description}\end{quote}
\index{run\_simulation\_loop() (in module qnav.simulation.simulation)@\spxentry{run\_simulation\_loop()}\spxextra{in module qnav.simulation.simulation}}

\begin{fulllineitems}
\phantomsection\label{\detokenize{modules/simulation/simulation:qnav.simulation.simulation.run_simulation_loop}}
\pysigstartsignatures
\pysiglinewithargsret{\sphinxcode{\sphinxupquote{qnav.simulation.simulation.}}\sphinxbfcode{\sphinxupquote{run\_simulation\_loop}}}{\sphinxparam{trajectory: \textasciitilde{}qnav.waypoints.trajectory.Trajectory, estimated\_state: \textasciitilde{}qnav.estimation.state.EstimatedState, hardware: \textasciitilde{}typing.Collection{[}\textasciitilde{}qnav.measurement.sensor.SensorFusion{]}, clock: \textasciitilde{}qnav.measurement.clock.Clock = \textless{}qnav.measurement.clock.DummyClock object\textgreater{}, results: \textasciitilde{}qnav.simulation.results.ResultsCapture = \textless{}qnav.simulation.results.CaptureAll object\textgreater{}, displayer: \textasciitilde{}qnav.simulation.display.Display = \textless{}qnav.simulation.display.ConsoleDisplay object\textgreater{}}}{}
\pysigstopsignatures
\sphinxAtStartPar
Executes the full simulation loop from start to finish.
On succession returns states and results in requested format.
\begin{quote}\begin{description}
\sphinxlineitem{Parameters}\begin{itemize}
\item {} 
\sphinxAtStartPar
\sphinxstyleliteralstrong{\sphinxupquote{trajectory}} ({\hyperref[\detokenize{modules/waypoints/trajectory:qnav.waypoints.trajectory.Trajectory}]{\sphinxcrossref{\sphinxstyleliteralemphasis{\sphinxupquote{Trajectory}}}}}) \textendash{} The trajectory to draw ground truth records from.
This is what the virtual platform is to follow.

\item {} 
\sphinxAtStartPar
\sphinxstyleliteralstrong{\sphinxupquote{estimated\_state}} ({\hyperref[\detokenize{modules/estimation/state:qnav.estimation.state.EstimatedState}]{\sphinxcrossref{\sphinxstyleliteralemphasis{\sphinxupquote{EstimatedState}}}}}) \textendash{} The initial estimated state. This will be
updated throughout the simulation via fusion.

\item {} 
\sphinxAtStartPar
\sphinxstyleliteralstrong{\sphinxupquote{hardware}} (\sphinxstyleliteralemphasis{\sphinxupquote{Collection}}\sphinxstyleliteralemphasis{\sphinxupquote{{[}}}{\hyperref[\detokenize{modules/measurement/sensor:qnav.measurement.sensor.SensorFusion}]{\sphinxcrossref{\sphinxstyleliteralemphasis{\sphinxupquote{SensorFusion}}}}}\sphinxstyleliteralemphasis{\sphinxupquote{{]}}}) \textendash{} A collection of SensorFusion instances. This represents
the hardware and fusion algorithms attached to the virtual platform.

\item {} 
\sphinxAtStartPar
\sphinxstyleliteralstrong{\sphinxupquote{clock}} \textendash{} A virtual clock added to simulate clock errors.
This is used when virtual clock drift is required.

\item {} 
\sphinxAtStartPar
\sphinxstyleliteralstrong{\sphinxupquote{results}} \textendash{} The results capture instance used to manage result data.
This controls how and what data is collected throughout the simulation.

\item {} 
\sphinxAtStartPar
\sphinxstyleliteralstrong{\sphinxupquote{displayer}} \textendash{} The displayer instance used for displaying progress.
This presents a display of the ongoing simulation allowing for
immediate feedback of its progress.

\end{itemize}

\end{description}\end{quote}

\end{fulllineitems}


\sphinxstepscope


\subsubsection{qnav.simulation.scheduling}
\label{\detokenize{modules/simulation/scheduling:module-qnav.simulation.scheduling}}\label{\detokenize{modules/simulation/scheduling:qnav-simulation-scheduling}}\label{\detokenize{modules/simulation/scheduling::doc}}\index{module@\spxentry{module}!qnav.simulation.scheduling@\spxentry{qnav.simulation.scheduling}}\index{qnav.simulation.scheduling@\spxentry{qnav.simulation.scheduling}!module@\spxentry{module}}

\paragraph{scheduling.py}
\label{\detokenize{modules/simulation/scheduling:scheduling-py}}\begin{quote}\begin{description}
\sphinxlineitem{summary}
\sphinxAtStartPar
Coordinates the scheduling of sensors and fusion methods.
This module provides classes that are used for organising the timing of
sensor measurements and fusion methods during simulations. Essentially
this is used to efficiently get the next non\sphinxhyphen{}inform time increment
along with what is to occur at that time.

\sphinxlineitem{author}
\sphinxAtStartPar
Michael Wright \sphinxhyphen{} \sphinxhref{mailto:mjwright@liverpool.ac.uk}{mjwright@liverpool.ac.uk}

\sphinxlineitem{version}
\sphinxAtStartPar
1.0.0 \sphinxhyphen{} First full implementation of module.

\end{description}\end{quote}
\index{TimingHandler (class in qnav.simulation.scheduling)@\spxentry{TimingHandler}\spxextra{class in qnav.simulation.scheduling}}

\begin{fulllineitems}
\phantomsection\label{\detokenize{modules/simulation/scheduling:qnav.simulation.scheduling.TimingHandler}}
\pysigstartsignatures
\pysiglinewithargsret{\sphinxbfcode{\sphinxupquote{class\DUrole{w}{ }}}\sphinxcode{\sphinxupquote{qnav.simulation.scheduling.}}\sphinxbfcode{\sphinxupquote{TimingHandler}}}{\sphinxparam{\DUrole{n}{hardware}\DUrole{p}{:}\DUrole{w}{ }\DUrole{n}{Iterable\DUrole{p}{{[}}{\hyperref[\detokenize{modules/measurement/sensor:qnav.measurement.sensor.SensorFusion}]{\sphinxcrossref{SensorFusion}}}\DUrole{p}{{]}}}}}{}
\pysigstopsignatures
\sphinxAtStartPar
A handler for managing the timings for sensors and fusion methods.
On creation this creates a mapping between shared sensors and fusion
methods. Through this the sensor with the earliest due time can be
obtained. Once this sensor is used (i.e a measurement has been produced)
a flag can be set for the fusion methods associate with it. When
sufficient sensor flags for a fusion method are set, the fusion method
can then be also obtained and returned.
\index{update\_sensors\_past\_time() (qnav.simulation.scheduling.TimingHandler method)@\spxentry{update\_sensors\_past\_time()}\spxextra{qnav.simulation.scheduling.TimingHandler method}}

\begin{fulllineitems}
\phantomsection\label{\detokenize{modules/simulation/scheduling:qnav.simulation.scheduling.TimingHandler.update_sensors_past_time}}
\pysigstartsignatures
\pysiglinewithargsret{\sphinxbfcode{\sphinxupquote{update\_sensors\_past\_time}}}{\sphinxparam{\DUrole{n}{delay\_time}\DUrole{p}{:}\DUrole{w}{ }\DUrole{n}{float}}}{{ $\rightarrow$ None}}
\pysigstopsignatures
\sphinxAtStartPar
Repeatedly updates sensors until all due time are past the given time.
Applies updates and skips measurements of all sensors until their due
times for next measurements are past the given time. This is used to
correctly induce a delay for measurements from the start of the
simulation. Note this does not collect measurements or change any
of the connected statuses used for fusion.
\begin{quote}\begin{description}
\sphinxlineitem{Parameters}
\sphinxAtStartPar
\sphinxstyleliteralstrong{\sphinxupquote{delay\_time}} (\sphinxstyleliteralemphasis{\sphinxupquote{float}}) \textendash{} The time to update all sensors past (in seconds).
Essentially the delay before measurements are to be processed.

\end{description}\end{quote}

\end{fulllineitems}

\index{set\_sensor\_used() (qnav.simulation.scheduling.TimingHandler method)@\spxentry{set\_sensor\_used()}\spxextra{qnav.simulation.scheduling.TimingHandler method}}

\begin{fulllineitems}
\phantomsection\label{\detokenize{modules/simulation/scheduling:qnav.simulation.scheduling.TimingHandler.set_sensor_used}}
\pysigstartsignatures
\pysiglinewithargsret{\sphinxbfcode{\sphinxupquote{set\_sensor\_used}}}{\sphinxparam{\DUrole{n}{sensor}\DUrole{p}{:}\DUrole{w}{ }\DUrole{n}{{\hyperref[\detokenize{modules/measurement/sensor:qnav.measurement.sensor.Sensor}]{\sphinxcrossref{Sensor}}}}}}{{ $\rightarrow$ None}}
\pysigstopsignatures
\sphinxAtStartPar
Flags that the given sensor has been used.
Marks the given sensor instance as used for all fusion methods that
requires this. If conditions are met, this trigger fusion methods
that use this sensor to be ready.
\begin{quote}\begin{description}
\sphinxlineitem{Parameters}
\sphinxAtStartPar
\sphinxstyleliteralstrong{\sphinxupquote{sensor}} ({\hyperref[\detokenize{modules/measurement/sensor:qnav.measurement.sensor.Sensor}]{\sphinxcrossref{\sphinxstyleliteralemphasis{\sphinxupquote{Sensor}}}}}) \textendash{} The sensor instance that has been used recently.

\end{description}\end{quote}

\end{fulllineitems}

\index{get\_next\_sensor() (qnav.simulation.scheduling.TimingHandler method)@\spxentry{get\_next\_sensor()}\spxextra{qnav.simulation.scheduling.TimingHandler method}}

\begin{fulllineitems}
\phantomsection\label{\detokenize{modules/simulation/scheduling:qnav.simulation.scheduling.TimingHandler.get_next_sensor}}
\pysigstartsignatures
\pysiglinewithargsret{\sphinxbfcode{\sphinxupquote{get\_next\_sensor}}}{}{{ $\rightarrow$ {\hyperref[\detokenize{modules/measurement/sensor:qnav.measurement.sensor.Sensor}]{\sphinxcrossref{Sensor}}}}}
\pysigstopsignatures
\sphinxAtStartPar
Returns the sensor with the earliest due time for measurement.
Searches through shared sensors and returns the one with the
earlier due time.
\begin{quote}\begin{description}
\sphinxlineitem{Returns}
\sphinxAtStartPar
Sensor with the earliest due time.

\sphinxlineitem{Return type}
\sphinxAtStartPar
{\hyperref[\detokenize{modules/measurement/sensor:qnav.measurement.sensor.Sensor}]{\sphinxcrossref{Sensor}}}

\end{description}\end{quote}

\end{fulllineitems}

\index{sensor\_iterator() (qnav.simulation.scheduling.TimingHandler method)@\spxentry{sensor\_iterator()}\spxextra{qnav.simulation.scheduling.TimingHandler method}}

\begin{fulllineitems}
\phantomsection\label{\detokenize{modules/simulation/scheduling:qnav.simulation.scheduling.TimingHandler.sensor_iterator}}
\pysigstartsignatures
\pysiglinewithargsret{\sphinxbfcode{\sphinxupquote{sensor\_iterator}}}{\sphinxparam{\DUrole{n}{max\_time}\DUrole{p}{:}\DUrole{w}{ }\DUrole{n}{float}}}{{ $\rightarrow$ Iterator\DUrole{p}{{[}}{\hyperref[\detokenize{modules/measurement/sensor:qnav.measurement.sensor.Sensor}]{\sphinxcrossref{Sensor}}}\DUrole{p}{{]}}}}
\pysigstopsignatures
\sphinxAtStartPar
Generates an iterator for continually getting the next sensor.
The produced iterator will be return the sensor with the earliest
measurement due time. The iteration stops when the next earliest
due time exceeds the given max time.
\begin{quote}\begin{description}
\sphinxlineitem{Parameters}
\sphinxAtStartPar
\sphinxstyleliteralstrong{\sphinxupquote{max\_time}} (\sphinxstyleliteralemphasis{\sphinxupquote{float}}) \textendash{} The maximum time due for measurements (in seconds).

\sphinxlineitem{Returns}
\sphinxAtStartPar
An iterator returning the sensor with the earliest due time.

\sphinxlineitem{Return type}
\sphinxAtStartPar
Iterator{[}{\hyperref[\detokenize{modules/measurement/sensor:qnav.measurement.sensor.Sensor}]{\sphinxcrossref{Sensor}}}{]}

\end{description}\end{quote}

\end{fulllineitems}

\index{get\_next\_fusion() (qnav.simulation.scheduling.TimingHandler method)@\spxentry{get\_next\_fusion()}\spxextra{qnav.simulation.scheduling.TimingHandler method}}

\begin{fulllineitems}
\phantomsection\label{\detokenize{modules/simulation/scheduling:qnav.simulation.scheduling.TimingHandler.get_next_fusion}}
\pysigstartsignatures
\pysiglinewithargsret{\sphinxbfcode{\sphinxupquote{get\_next\_fusion}}}{\sphinxparam{\DUrole{n}{sensor}\DUrole{p}{:}\DUrole{w}{ }\DUrole{n}{{\hyperref[\detokenize{modules/measurement/sensor:qnav.measurement.sensor.Sensor}]{\sphinxcrossref{Sensor}}}}\DUrole{w}{ }\DUrole{o}{=}\DUrole{w}{ }\DUrole{default_value}{None}}}{{ $\rightarrow$ list\DUrole{p}{{[}}{\hyperref[\detokenize{modules/measurement/sensor:qnav.measurement.sensor.SensorFusion}]{\sphinxcrossref{SensorFusion}}}\DUrole{p}{{]}}}}
\pysigstopsignatures
\sphinxAtStartPar
Returns any fusion methods that are to be used.
Searches through fusion methods and returns a list collection of those
that are ready to be used. These are automatically marked as used on
return. If no fusion methods are ready, then the list will be empty.
For faster retrieval, the search of fusion methods can be filtered
to only those assigned with a given sensor instance.
\begin{quote}\begin{description}
\sphinxlineitem{Parameters}
\sphinxAtStartPar
\sphinxstyleliteralstrong{\sphinxupquote{sensor}} ({\hyperref[\detokenize{modules/measurement/sensor:qnav.measurement.sensor.Sensor}]{\sphinxcrossref{\sphinxstyleliteralemphasis{\sphinxupquote{Sensor}}}}}) \textendash{} Sensor instance to filter fusion methods by (Optional).

\sphinxlineitem{Returns}
\sphinxAtStartPar
A list of fusion methods ready to be applied.

\sphinxlineitem{Return type}
\sphinxAtStartPar
list{[}{\hyperref[\detokenize{modules/measurement/sensor:qnav.measurement.sensor.SensorFusion}]{\sphinxcrossref{SensorFusion}}}{]}

\end{description}\end{quote}

\end{fulllineitems}

\index{print\_mappings() (qnav.simulation.scheduling.TimingHandler method)@\spxentry{print\_mappings()}\spxextra{qnav.simulation.scheduling.TimingHandler method}}

\begin{fulllineitems}
\phantomsection\label{\detokenize{modules/simulation/scheduling:qnav.simulation.scheduling.TimingHandler.print_mappings}}
\pysigstartsignatures
\pysiglinewithargsret{\sphinxbfcode{\sphinxupquote{print\_mappings}}}{}{}
\pysigstopsignatures
\sphinxAtStartPar
Displays the internal scheduling mappings in the output console.
When called, prints the fusion methods attached to each sensor,
along with their ready status. Useful for debugging states.

\end{fulllineitems}


\end{fulllineitems}


\sphinxstepscope


\subsubsection{qnav.simulation.holonomic}
\label{\detokenize{modules/simulation/holonomic:module-qnav.simulation.holonomic}}\label{\detokenize{modules/simulation/holonomic:qnav-simulation-holonomic}}\label{\detokenize{modules/simulation/holonomic::doc}}\index{module@\spxentry{module}!qnav.simulation.holonomic@\spxentry{qnav.simulation.holonomic}}\index{qnav.simulation.holonomic@\spxentry{qnav.simulation.holonomic}!module@\spxentry{module}}

\paragraph{holonomic.py}
\label{\detokenize{modules/simulation/holonomic:holonomic-py}}\begin{quote}\begin{description}
\sphinxlineitem{summary}
\sphinxAtStartPar
Supplies means of applying configured holonomic constraints.
Virtual simulations are vulnerable to numerical instabilities, causing
unintended drift between the ground truth and estimated state over time.
Unlike with artificial noise that is purposely added through configured
sensor error properties, this is undesirable and can cause unintended
behavior. To combat this, this holonomic module can be used to apply
a variety of constraints to restrict the estimated state of the system
to values within the expected behaviour. In other words, this is used
to slightly nudge the estimated state toward that of the ground truth
at a scale to overcome numerical instabilities (such as rounding errors)
but not to affect simulated errors.

\sphinxlineitem{authors}
\begin{DUlineblock}{0em}
\item[] Prof. Jason Ralph \sphinxhyphen{} \sphinxhref{mailto:jfralph@liverpool.ac.uk}{jfralph@liverpool.ac.uk}
\item[] Michael Wright \sphinxhyphen{} \sphinxhref{mailto:mjwright@liverpool.ac.uk}{mjwright@liverpool.ac.uk}
\end{DUlineblock}

\sphinxlineitem{version}
\sphinxAtStartPar
1.0.0 \sphinxhyphen{} First full implementation.

\end{description}\end{quote}
\index{HolonomicConstraint (class in qnav.simulation.holonomic)@\spxentry{HolonomicConstraint}\spxextra{class in qnav.simulation.holonomic}}

\begin{fulllineitems}
\phantomsection\label{\detokenize{modules/simulation/holonomic:qnav.simulation.holonomic.HolonomicConstraint}}
\pysigstartsignatures
\pysiglinewithargsret{\sphinxbfcode{\sphinxupquote{class\DUrole{w}{ }}}\sphinxcode{\sphinxupquote{qnav.simulation.holonomic.}}\sphinxbfcode{\sphinxupquote{HolonomicConstraint}}}{\sphinxparam{\DUrole{n}{frequency}\DUrole{p}{:}\DUrole{w}{ }\DUrole{n}{float}}\sphinxparamcomma \sphinxparam{\DUrole{n}{gt\_filter}\DUrole{p}{:}\DUrole{w}{ }\DUrole{n}{str}}\sphinxparamcomma \sphinxparam{\DUrole{n}{fixed\_gain}\DUrole{p}{:}\DUrole{w}{ }\DUrole{n}{float}}\sphinxparamcomma \sphinxparam{\DUrole{n}{axis\_filter}\DUrole{p}{:}\DUrole{w}{ }\DUrole{n}{ndarray}\DUrole{w}{ }\DUrole{o}{=}\DUrole{w}{ }\DUrole{default_value}{None}}}{}
\pysigstopsignatures
\sphinxAtStartPar
Obtains holonomic correction values during simulation.
For a given ground truth record, gently shifts the estimated state towards
the ground truth state at a rate to overcome unintended drift, but not
affect simulation accuracy (i.e. intended behaviour from
artificial/configured errors).
\index{filter\_id (qnav.simulation.holonomic.HolonomicConstraint property)@\spxentry{filter\_id}\spxextra{qnav.simulation.holonomic.HolonomicConstraint property}}

\begin{fulllineitems}
\phantomsection\label{\detokenize{modules/simulation/holonomic:qnav.simulation.holonomic.HolonomicConstraint.filter_id}}
\pysigstartsignatures
\pysigline{\sphinxbfcode{\sphinxupquote{property\DUrole{w}{ }}}\sphinxbfcode{\sphinxupquote{filter\_id}}\sphinxbfcode{\sphinxupquote{\DUrole{p}{:}\DUrole{w}{ }str}}}
\pysigstopsignatures
\sphinxAtStartPar
Displays the tag used for filtering the ground truth.
Simply used for checking filter doesn’t already exist.
\begin{quote}\begin{description}
\sphinxlineitem{Returns}
\sphinxAtStartPar
The tag ID for the ground truth record to filter by.

\sphinxlineitem{Return type}
\sphinxAtStartPar
str

\end{description}\end{quote}

\end{fulllineitems}

\index{filter\_gain (qnav.simulation.holonomic.HolonomicConstraint property)@\spxentry{filter\_gain}\spxextra{qnav.simulation.holonomic.HolonomicConstraint property}}

\begin{fulllineitems}
\phantomsection\label{\detokenize{modules/simulation/holonomic:qnav.simulation.holonomic.HolonomicConstraint.filter_gain}}
\pysigstartsignatures
\pysigline{\sphinxbfcode{\sphinxupquote{property\DUrole{w}{ }}}\sphinxbfcode{\sphinxupquote{filter\_gain}}\sphinxbfcode{\sphinxupquote{\DUrole{p}{:}\DUrole{w}{ }float\DUrole{w}{ }\DUrole{p}{|}\DUrole{w}{ }ndarray}}}
\pysigstopsignatures
\sphinxAtStartPar
Displays the gain amount (percentage) being applied.
Simply used for checking the filter’s value.
\begin{quote}\begin{description}
\sphinxlineitem{Returns}
\sphinxAtStartPar
The gain amount to apply to each axis

\sphinxlineitem{Return type}
\sphinxAtStartPar
float | np.ndarray

\end{description}\end{quote}

\end{fulllineitems}

\index{take\_measurement() (qnav.simulation.holonomic.HolonomicConstraint method)@\spxentry{take\_measurement()}\spxextra{qnav.simulation.holonomic.HolonomicConstraint method}}

\begin{fulllineitems}
\phantomsection\label{\detokenize{modules/simulation/holonomic:qnav.simulation.holonomic.HolonomicConstraint.take_measurement}}
\pysigstartsignatures
\pysiglinewithargsret{\sphinxbfcode{\sphinxupquote{take\_measurement}}}{\sphinxparam{\DUrole{n}{est\_time}\DUrole{p}{:}\DUrole{w}{ }\DUrole{n}{float}}\sphinxparamcomma \sphinxparam{\DUrole{n}{ground\_truth}\DUrole{p}{:}\DUrole{w}{ }\DUrole{n}{{\hyperref[\detokenize{modules/waypoints/trajectory:qnav.waypoints.trajectory.GroundTruth}]{\sphinxcrossref{GroundTruth}}}}}}{}
\pysigstopsignatures
\sphinxAtStartPar
Extracts data from ground truth to obtain correction amount.
When called during simulation loop, this checks the relevant data can
be extracted from the ground truth and uses it to update the
correction value to apply at that time.
\begin{quote}\begin{description}
\sphinxlineitem{Parameters}\begin{itemize}
\item {} 
\sphinxAtStartPar
\sphinxstyleliteralstrong{\sphinxupquote{est\_time}} (\sphinxstyleliteralemphasis{\sphinxupquote{float}}) \textendash{} The estimate time of the measurement (not used).

\item {} 
\sphinxAtStartPar
\sphinxstyleliteralstrong{\sphinxupquote{ground\_truth}} ({\hyperref[\detokenize{modules/waypoints/trajectory:qnav.waypoints.trajectory.GroundTruth}]{\sphinxcrossref{\sphinxstyleliteralemphasis{\sphinxupquote{GroundTruth}}}}}) \textendash{} The ground truth record at the given time.

\end{itemize}

\end{description}\end{quote}

\end{fulllineitems}

\index{clear\_measurement() (qnav.simulation.holonomic.HolonomicConstraint method)@\spxentry{clear\_measurement()}\spxextra{qnav.simulation.holonomic.HolonomicConstraint method}}

\begin{fulllineitems}
\phantomsection\label{\detokenize{modules/simulation/holonomic:qnav.simulation.holonomic.HolonomicConstraint.clear_measurement}}
\pysigstartsignatures
\pysiglinewithargsret{\sphinxbfcode{\sphinxupquote{clear\_measurement}}}{}{}
\pysigstopsignatures
\sphinxAtStartPar
Clears the last measurement value obtained.
Useful in preventing the accidental reuse of measurements when using
multiple holonomic constraints asynchronously.

\end{fulllineitems}

\index{last\_measurement (qnav.simulation.holonomic.HolonomicConstraint property)@\spxentry{last\_measurement}\spxextra{qnav.simulation.holonomic.HolonomicConstraint property}}

\begin{fulllineitems}
\phantomsection\label{\detokenize{modules/simulation/holonomic:qnav.simulation.holonomic.HolonomicConstraint.last_measurement}}
\pysigstartsignatures
\pysigline{\sphinxbfcode{\sphinxupquote{property\DUrole{w}{ }}}\sphinxbfcode{\sphinxupquote{last\_measurement}}\sphinxbfcode{\sphinxupquote{\DUrole{p}{:}\DUrole{w}{ }any}}}
\pysigstopsignatures
\sphinxAtStartPar
Returns the last measurement captured (i.e. correction to apply).
Upon error this will return None.
\begin{quote}\begin{description}
\sphinxlineitem{Returns}
\sphinxAtStartPar
The holonomic correction to apply.

\sphinxlineitem{Return type}
\sphinxAtStartPar
float | np.ndarray

\end{description}\end{quote}

\end{fulllineitems}


\end{fulllineitems}

\index{HolonomicFusion (class in qnav.simulation.holonomic)@\spxentry{HolonomicFusion}\spxextra{class in qnav.simulation.holonomic}}

\begin{fulllineitems}
\phantomsection\label{\detokenize{modules/simulation/holonomic:qnav.simulation.holonomic.HolonomicFusion}}
\pysigstartsignatures
\pysiglinewithargsret{\sphinxbfcode{\sphinxupquote{class\DUrole{w}{ }}}\sphinxcode{\sphinxupquote{qnav.simulation.holonomic.}}\sphinxbfcode{\sphinxupquote{HolonomicFusion}}}{\sphinxparam{\DUrole{n}{trigger}\DUrole{p}{:}\DUrole{w}{ }\DUrole{n}{{\hyperref[\detokenize{modules/measurement/sensor:qnav.measurement.sensor.FusionTrigger}]{\sphinxcrossref{FusionTrigger}}}}}\sphinxparamcomma \sphinxparam{\DUrole{o}{*}\DUrole{n}{holonomics}\DUrole{p}{:}\DUrole{w}{ }\DUrole{n}{{\hyperref[\detokenize{modules/simulation/holonomic:qnav.simulation.holonomic.HolonomicConstraint}]{\sphinxcrossref{HolonomicConstraint}}}}}}{}
\pysigstopsignatures
\sphinxAtStartPar
Applies the holonomic constraint corrections during simulation.
Manages a correction of holonomic constraints and applying them
to correct the current estimated state.
\index{perform\_fusion() (qnav.simulation.holonomic.HolonomicFusion method)@\spxentry{perform\_fusion()}\spxextra{qnav.simulation.holonomic.HolonomicFusion method}}

\begin{fulllineitems}
\phantomsection\label{\detokenize{modules/simulation/holonomic:qnav.simulation.holonomic.HolonomicFusion.perform_fusion}}
\pysigstartsignatures
\pysiglinewithargsret{\sphinxbfcode{\sphinxupquote{perform\_fusion}}}{\sphinxparam{\DUrole{n}{estimated\_state}\DUrole{p}{:}\DUrole{w}{ }\DUrole{n}{{\hyperref[\detokenize{modules/estimation/state:qnav.estimation.state.EstimatedState}]{\sphinxcrossref{EstimatedState}}}}}}{{ $\rightarrow$ None}}
\pysigstopsignatures
\sphinxAtStartPar
Apply the holonomic constraints and corrections during simulation.
When called, performs fusion by overwriting current estimates.
\begin{quote}\begin{description}
\sphinxlineitem{Parameters}
\sphinxAtStartPar
\sphinxstyleliteralstrong{\sphinxupquote{estimated\_state}} ({\hyperref[\detokenize{modules/estimation/state:qnav.estimation.state.EstimatedState}]{\sphinxcrossref{\sphinxstyleliteralemphasis{\sphinxupquote{EstimatedState}}}}}) \textendash{} The current estimated states to update.

\end{description}\end{quote}

\end{fulllineitems}


\end{fulllineitems}


\sphinxstepscope


\subsubsection{qnav.simulation.simulation}
\label{\detokenize{modules/simulation/results:module-qnav.simulation.simulation}}\label{\detokenize{modules/simulation/results:qnav-simulation-simulation}}\label{\detokenize{modules/simulation/results::doc}}\index{module@\spxentry{module}!qnav.simulation.simulation@\spxentry{qnav.simulation.simulation}}\index{qnav.simulation.simulation@\spxentry{qnav.simulation.simulation}!module@\spxentry{module}}

\paragraph{simulation.py}
\label{\detokenize{modules/simulation/results:simulation-py}}\begin{quote}\begin{description}
\sphinxlineitem{summary}
\sphinxAtStartPar
The core harness responsible for coordinating simulations.
This module contains the core functionalities for managing simulations.
This mostly revolves around a core loop and the scheduling of sensor
time intervals, sensor measurements and data fusion.

\sphinxlineitem{author}
\sphinxAtStartPar
Michael Wright \sphinxhyphen{} \sphinxhref{mailto:mjwright@liverpool.ac.uk}{mjwright@liverpool.ac.uk}

\sphinxlineitem{version}
\sphinxAtStartPar
1.0.0 \sphinxhyphen{} First full implementation.

\end{description}\end{quote}
\index{run\_simulation\_loop() (in module qnav.simulation.simulation)@\spxentry{run\_simulation\_loop()}\spxextra{in module qnav.simulation.simulation}}

\begin{fulllineitems}
\phantomsection\label{\detokenize{modules/simulation/results:qnav.simulation.simulation.run_simulation_loop}}
\pysigstartsignatures
\pysiglinewithargsret{\sphinxcode{\sphinxupquote{qnav.simulation.simulation.}}\sphinxbfcode{\sphinxupquote{run\_simulation\_loop}}}{\sphinxparam{trajectory: \textasciitilde{}qnav.waypoints.trajectory.Trajectory, estimated\_state: \textasciitilde{}qnav.estimation.state.EstimatedState, hardware: \textasciitilde{}typing.Collection{[}\textasciitilde{}qnav.measurement.sensor.SensorFusion{]}, clock: \textasciitilde{}qnav.measurement.clock.Clock = \textless{}qnav.measurement.clock.DummyClock object\textgreater{}, results: \textasciitilde{}qnav.simulation.results.ResultsCapture = \textless{}qnav.simulation.results.CaptureAll object\textgreater{}, displayer: \textasciitilde{}qnav.simulation.display.Display = \textless{}qnav.simulation.display.ConsoleDisplay object\textgreater{}}}{}
\pysigstopsignatures
\sphinxAtStartPar
Executes the full simulation loop from start to finish.
On succession returns states and results in requested format.
\begin{quote}\begin{description}
\sphinxlineitem{Parameters}\begin{itemize}
\item {} 
\sphinxAtStartPar
\sphinxstyleliteralstrong{\sphinxupquote{trajectory}} ({\hyperref[\detokenize{modules/waypoints/trajectory:qnav.waypoints.trajectory.Trajectory}]{\sphinxcrossref{\sphinxstyleliteralemphasis{\sphinxupquote{Trajectory}}}}}) \textendash{} The trajectory to draw ground truth records from.
This is what the virtual platform is to follow.

\item {} 
\sphinxAtStartPar
\sphinxstyleliteralstrong{\sphinxupquote{estimated\_state}} ({\hyperref[\detokenize{modules/estimation/state:qnav.estimation.state.EstimatedState}]{\sphinxcrossref{\sphinxstyleliteralemphasis{\sphinxupquote{EstimatedState}}}}}) \textendash{} The initial estimated state. This will be
updated throughout the simulation via fusion.

\item {} 
\sphinxAtStartPar
\sphinxstyleliteralstrong{\sphinxupquote{hardware}} (\sphinxstyleliteralemphasis{\sphinxupquote{Collection}}\sphinxstyleliteralemphasis{\sphinxupquote{{[}}}{\hyperref[\detokenize{modules/measurement/sensor:qnav.measurement.sensor.SensorFusion}]{\sphinxcrossref{\sphinxstyleliteralemphasis{\sphinxupquote{SensorFusion}}}}}\sphinxstyleliteralemphasis{\sphinxupquote{{]}}}) \textendash{} A collection of SensorFusion instances. This represents
the hardware and fusion algorithms attached to the virtual platform.

\item {} 
\sphinxAtStartPar
\sphinxstyleliteralstrong{\sphinxupquote{clock}} \textendash{} A virtual clock added to simulate clock errors.
This is used when virtual clock drift is required.

\item {} 
\sphinxAtStartPar
\sphinxstyleliteralstrong{\sphinxupquote{results}} \textendash{} The results capture instance used to manage result data.
This controls how and what data is collected throughout the simulation.

\item {} 
\sphinxAtStartPar
\sphinxstyleliteralstrong{\sphinxupquote{displayer}} \textendash{} The displayer instance used for displaying progress.
This presents a display of the ongoing simulation allowing for
immediate feedback of its progress.

\end{itemize}

\end{description}\end{quote}

\end{fulllineitems}


\sphinxstepscope


\subsubsection{qnav.simulation.display}
\label{\detokenize{modules/simulation/display:module-qnav.simulation.display}}\label{\detokenize{modules/simulation/display:qnav-simulation-display}}\label{\detokenize{modules/simulation/display::doc}}\index{module@\spxentry{module}!qnav.simulation.display@\spxentry{qnav.simulation.display}}\index{qnav.simulation.display@\spxentry{qnav.simulation.display}!module@\spxentry{module}}

\paragraph{display.py}
\label{\detokenize{modules/simulation/display:display-py}}\begin{quote}\begin{description}
\sphinxlineitem{summary}
\sphinxAtStartPar
Allows means for displaying simulation progress in real\sphinxhyphen{}time.
This module is responsible for providing an interface for allowing the
displaying of simulation data so that progress can be presented to the
user in real\sphinxhyphen{}time.

\sphinxlineitem{author}
\sphinxAtStartPar
Michael Wright \sphinxhyphen{} \sphinxhref{mailto:mjwright@liverpool.ac.uk}{mjwright@liverpool.ac.uk}

\sphinxlineitem{version}
\sphinxAtStartPar
1.0.0 \sphinxhyphen{} First full implementation.

\end{description}\end{quote}
\index{is\_using\_terminal() (in module qnav.simulation.display)@\spxentry{is\_using\_terminal()}\spxextra{in module qnav.simulation.display}}

\begin{fulllineitems}
\phantomsection\label{\detokenize{modules/simulation/display:qnav.simulation.display.is_using_terminal}}
\pysigstartsignatures
\pysiglinewithargsret{\sphinxcode{\sphinxupquote{qnav.simulation.display.}}\sphinxbfcode{\sphinxupquote{is\_using\_terminal}}}{}{{ $\rightarrow$ bool}}
\pysigstopsignatures
\sphinxAtStartPar
Checks if the application is being executed via terminal.
Used to determine if the toolbox is being run and displayed through a
command\sphinxhyphen{}line terminal, and not some interactive shell (or IDE). Used to
obtain an idea of how output data is to be displayed.
\begin{quote}\begin{description}
\sphinxlineitem{Returns}
\sphinxAtStartPar
True if standard output is being displayed inside terminal.

\sphinxlineitem{Return type}
\sphinxAtStartPar
bool

\end{description}\end{quote}

\end{fulllineitems}

\index{clear\_terminal\_line() (in module qnav.simulation.display)@\spxentry{clear\_terminal\_line()}\spxextra{in module qnav.simulation.display}}

\begin{fulllineitems}
\phantomsection\label{\detokenize{modules/simulation/display:qnav.simulation.display.clear_terminal_line}}
\pysigstartsignatures
\pysiglinewithargsret{\sphinxcode{\sphinxupquote{qnav.simulation.display.}}\sphinxbfcode{\sphinxupquote{clear\_terminal\_line}}}{\sphinxparam{\DUrole{n}{num\_lines}\DUrole{p}{:}\DUrole{w}{ }\DUrole{n}{int}\DUrole{w}{ }\DUrole{o}{=}\DUrole{w}{ }\DUrole{default_value}{1}}}{{ $\rightarrow$ None}}
\pysigstopsignatures
\sphinxAtStartPar
Clears a previously printed line in the console.
Only works when executing application from command line.
\begin{quote}\begin{description}
\sphinxlineitem{Parameters}
\sphinxAtStartPar
\sphinxstyleliteralstrong{\sphinxupquote{num\_lines}} (\sphinxstyleliteralemphasis{\sphinxupquote{int}}) \textendash{} Number of lines to clear (default = 1).

\end{description}\end{quote}

\end{fulllineitems}

\index{sec\_to\_str() (in module qnav.simulation.display)@\spxentry{sec\_to\_str()}\spxextra{in module qnav.simulation.display}}

\begin{fulllineitems}
\phantomsection\label{\detokenize{modules/simulation/display:qnav.simulation.display.sec_to_str}}
\pysigstartsignatures
\pysiglinewithargsret{\sphinxcode{\sphinxupquote{qnav.simulation.display.}}\sphinxbfcode{\sphinxupquote{sec\_to\_str}}}{\sphinxparam{\DUrole{n}{time\_sec}\DUrole{p}{:}\DUrole{w}{ }\DUrole{n}{int}}}{{ $\rightarrow$ str}}
\pysigstopsignatures
\sphinxAtStartPar
Converts seconds to printed hours, minutes and seconds.
Used to make printed time more human readable.
\begin{quote}\begin{description}
\sphinxlineitem{Parameters}
\sphinxAtStartPar
\sphinxstyleliteralstrong{\sphinxupquote{time\_sec}} (\sphinxstyleliteralemphasis{\sphinxupquote{int}}) \textendash{} The time to display in seconds.

\sphinxlineitem{Returns}
\sphinxAtStartPar
The given time in HH:MM:SS string format.

\sphinxlineitem{Return type}
\sphinxAtStartPar
str

\end{description}\end{quote}

\end{fulllineitems}

\index{Display (class in qnav.simulation.display)@\spxentry{Display}\spxextra{class in qnav.simulation.display}}

\begin{fulllineitems}
\phantomsection\label{\detokenize{modules/simulation/display:qnav.simulation.display.Display}}
\pysigstartsignatures
\pysiglinewithargsret{\sphinxbfcode{\sphinxupquote{class\DUrole{w}{ }}}\sphinxcode{\sphinxupquote{qnav.simulation.display.}}\sphinxbfcode{\sphinxupquote{Display}}}{\sphinxparam{\DUrole{n}{update\_interval}\DUrole{p}{:}\DUrole{w}{ }\DUrole{n}{float}}}{}
\pysigstopsignatures
\sphinxAtStartPar
A template used to define classes for displaying simulation progress.
This provides the definition to be followed for classes that can be
used for displaying simulation data during runtime. This is useful
for users who wish to see verbose information or plot data during
simulations, without having to wait until the end.
\index{on\_start() (qnav.simulation.display.Display method)@\spxentry{on\_start()}\spxextra{qnav.simulation.display.Display method}}

\begin{fulllineitems}
\phantomsection\label{\detokenize{modules/simulation/display:qnav.simulation.display.Display.on_start}}
\pysigstartsignatures
\pysiglinewithargsret{\sphinxbfcode{\sphinxupquote{abstract\DUrole{w}{ }}}\sphinxbfcode{\sphinxupquote{on\_start}}}{\sphinxparam{\DUrole{n}{trajectory}\DUrole{p}{:}\DUrole{w}{ }\DUrole{n}{{\hyperref[\detokenize{modules/waypoints/trajectory:qnav.waypoints.trajectory.Trajectory}]{\sphinxcrossref{Trajectory}}}}}}{}
\pysigstopsignatures
\sphinxAtStartPar
Called at the start of the simulation before first iteration.
Can be used to initialise any displaying components by extracting
necessary data from the trajectory.
\begin{quote}\begin{description}
\sphinxlineitem{Parameters}
\sphinxAtStartPar
\sphinxstyleliteralstrong{\sphinxupquote{trajectory}} ({\hyperref[\detokenize{modules/waypoints/trajectory:qnav.waypoints.trajectory.Trajectory}]{\sphinxcrossref{\sphinxstyleliteralemphasis{\sphinxupquote{Trajectory}}}}}) \textendash{} The ground truth trajectory instance.

\end{description}\end{quote}

\end{fulllineitems}

\index{on\_update() (qnav.simulation.display.Display method)@\spxentry{on\_update()}\spxextra{qnav.simulation.display.Display method}}

\begin{fulllineitems}
\phantomsection\label{\detokenize{modules/simulation/display:qnav.simulation.display.Display.on_update}}
\pysigstartsignatures
\pysiglinewithargsret{\sphinxbfcode{\sphinxupquote{abstract\DUrole{w}{ }}}\sphinxbfcode{\sphinxupquote{on\_update}}}{\sphinxparam{\DUrole{n}{timestamp}\DUrole{p}{:}\DUrole{w}{ }\DUrole{n}{float}}\sphinxparamcomma \sphinxparam{\DUrole{n}{estimated\_state}\DUrole{p}{:}\DUrole{w}{ }\DUrole{n}{{\hyperref[\detokenize{modules/estimation/state:qnav.estimation.state.EstimatedState}]{\sphinxcrossref{EstimatedState}}}}}\sphinxparamcomma \sphinxparam{\DUrole{n}{ground\_truth}\DUrole{p}{:}\DUrole{w}{ }\DUrole{n}{{\hyperref[\detokenize{modules/waypoints/trajectory:qnav.waypoints.trajectory.GroundTruth}]{\sphinxcrossref{GroundTruth}}}}}}{}
\pysigstopsignatures
\sphinxAtStartPar
Called during the simulation after each fusion update.
Presents the current (true) timestep, estimated states and ground
truth record.
\begin{quote}\begin{description}
\sphinxlineitem{Parameters}\begin{itemize}
\item {} 
\sphinxAtStartPar
\sphinxstyleliteralstrong{\sphinxupquote{timestamp}} (\sphinxstyleliteralemphasis{\sphinxupquote{float}}) \textendash{} The current (true) simulation time in seconds.

\item {} 
\sphinxAtStartPar
\sphinxstyleliteralstrong{\sphinxupquote{estimated\_state}} ({\hyperref[\detokenize{modules/estimation/state:qnav.estimation.state.EstimatedState}]{\sphinxcrossref{\sphinxstyleliteralemphasis{\sphinxupquote{EstimatedState}}}}}) \textendash{} The current estimated state of the simulation.

\item {} 
\sphinxAtStartPar
\sphinxstyleliteralstrong{\sphinxupquote{ground\_truth}} ({\hyperref[\detokenize{modules/waypoints/trajectory:qnav.waypoints.trajectory.GroundTruth}]{\sphinxcrossref{\sphinxstyleliteralemphasis{\sphinxupquote{GroundTruth}}}}}) \textendash{} The current ground truth state of the simulation.

\end{itemize}

\end{description}\end{quote}

\end{fulllineitems}

\index{on\_finish() (qnav.simulation.display.Display method)@\spxentry{on\_finish()}\spxextra{qnav.simulation.display.Display method}}

\begin{fulllineitems}
\phantomsection\label{\detokenize{modules/simulation/display:qnav.simulation.display.Display.on_finish}}
\pysigstartsignatures
\pysiglinewithargsret{\sphinxbfcode{\sphinxupquote{abstract\DUrole{w}{ }}}\sphinxbfcode{\sphinxupquote{on\_finish}}}{\sphinxparam{\DUrole{n}{timestamp}\DUrole{p}{:}\DUrole{w}{ }\DUrole{n}{float}}\sphinxparamcomma \sphinxparam{\DUrole{n}{estimated\_state}\DUrole{p}{:}\DUrole{w}{ }\DUrole{n}{{\hyperref[\detokenize{modules/estimation/state:qnav.estimation.state.EstimatedState}]{\sphinxcrossref{EstimatedState}}}}}\sphinxparamcomma \sphinxparam{\DUrole{n}{ground\_truth}\DUrole{p}{:}\DUrole{w}{ }\DUrole{n}{{\hyperref[\detokenize{modules/waypoints/trajectory:qnav.waypoints.trajectory.GroundTruth}]{\sphinxcrossref{GroundTruth}}}}}}{}
\pysigstopsignatures
\sphinxAtStartPar
Called at the end of the simulation after final iteration.
Can be used to finish up any displaying of information. Presents the
final (true) timestep, estimated states and ground truth record.
\begin{quote}\begin{description}
\sphinxlineitem{Parameters}\begin{itemize}
\item {} 
\sphinxAtStartPar
\sphinxstyleliteralstrong{\sphinxupquote{timestamp}} (\sphinxstyleliteralemphasis{\sphinxupquote{float}}) \textendash{} The final (true) simulation time in seconds.

\item {} 
\sphinxAtStartPar
\sphinxstyleliteralstrong{\sphinxupquote{estimated\_state}} ({\hyperref[\detokenize{modules/estimation/state:qnav.estimation.state.EstimatedState}]{\sphinxcrossref{\sphinxstyleliteralemphasis{\sphinxupquote{EstimatedState}}}}}) \textendash{} The final estimated state of the simulation.

\item {} 
\sphinxAtStartPar
\sphinxstyleliteralstrong{\sphinxupquote{ground\_truth}} ({\hyperref[\detokenize{modules/waypoints/trajectory:qnav.waypoints.trajectory.GroundTruth}]{\sphinxcrossref{\sphinxstyleliteralemphasis{\sphinxupquote{GroundTruth}}}}}) \textendash{} The final ground truth state of the simulation.

\end{itemize}

\end{description}\end{quote}

\end{fulllineitems}


\end{fulllineitems}

\index{ConsoleDisplay (class in qnav.simulation.display)@\spxentry{ConsoleDisplay}\spxextra{class in qnav.simulation.display}}

\begin{fulllineitems}
\phantomsection\label{\detokenize{modules/simulation/display:qnav.simulation.display.ConsoleDisplay}}
\pysigstartsignatures
\pysiglinewithargsret{\sphinxbfcode{\sphinxupquote{class\DUrole{w}{ }}}\sphinxcode{\sphinxupquote{qnav.simulation.display.}}\sphinxbfcode{\sphinxupquote{ConsoleDisplay}}}{\sphinxparam{\DUrole{n}{update\_interval}\DUrole{p}{:}\DUrole{w}{ }\DUrole{n}{float}}}{}
\pysigstopsignatures
\sphinxAtStartPar
A simple display for printing simulation progress in console.
Used to display that periodically prints the simulation’s runtime progress
in the output console. Namely used as an example display class.
\index{on\_start() (qnav.simulation.display.ConsoleDisplay method)@\spxentry{on\_start()}\spxextra{qnav.simulation.display.ConsoleDisplay method}}

\begin{fulllineitems}
\phantomsection\label{\detokenize{modules/simulation/display:qnav.simulation.display.ConsoleDisplay.on_start}}
\pysigstartsignatures
\pysiglinewithargsret{\sphinxbfcode{\sphinxupquote{on\_start}}}{\sphinxparam{\DUrole{n}{trajectory}\DUrole{p}{:}\DUrole{w}{ }\DUrole{n}{{\hyperref[\detokenize{modules/waypoints/trajectory:qnav.waypoints.trajectory.Trajectory}]{\sphinxcrossref{Trajectory}}}}}}{}
\pysigstopsignatures
\sphinxAtStartPar
Called at the start of the simulation before first iteration.
Used to obtain the start and end times of the trajectory.
\begin{quote}\begin{description}
\sphinxlineitem{Parameters}
\sphinxAtStartPar
\sphinxstyleliteralstrong{\sphinxupquote{trajectory}} ({\hyperref[\detokenize{modules/waypoints/trajectory:qnav.waypoints.trajectory.Trajectory}]{\sphinxcrossref{\sphinxstyleliteralemphasis{\sphinxupquote{Trajectory}}}}}) \textendash{} The ground truth trajectory instance.

\end{description}\end{quote}

\end{fulllineitems}

\index{on\_update() (qnav.simulation.display.ConsoleDisplay method)@\spxentry{on\_update()}\spxextra{qnav.simulation.display.ConsoleDisplay method}}

\begin{fulllineitems}
\phantomsection\label{\detokenize{modules/simulation/display:qnav.simulation.display.ConsoleDisplay.on_update}}
\pysigstartsignatures
\pysiglinewithargsret{\sphinxbfcode{\sphinxupquote{on\_update}}}{\sphinxparam{\DUrole{n}{timestamp}\DUrole{p}{:}\DUrole{w}{ }\DUrole{n}{float}}\sphinxparamcomma \sphinxparam{\DUrole{n}{estimated\_state}\DUrole{p}{:}\DUrole{w}{ }\DUrole{n}{{\hyperref[\detokenize{modules/estimation/state:qnav.estimation.state.EstimatedState}]{\sphinxcrossref{EstimatedState}}}}}\sphinxparamcomma \sphinxparam{\DUrole{n}{ground\_truth}\DUrole{p}{:}\DUrole{w}{ }\DUrole{n}{{\hyperref[\detokenize{modules/waypoints/trajectory:qnav.waypoints.trajectory.GroundTruth}]{\sphinxcrossref{GroundTruth}}}}}}{}
\pysigstopsignatures
\sphinxAtStartPar
Called during the simulation after each fusion update.
Updates the runtime progress display of the simulation.
\begin{quote}\begin{description}
\sphinxlineitem{Parameters}\begin{itemize}
\item {} 
\sphinxAtStartPar
\sphinxstyleliteralstrong{\sphinxupquote{timestamp}} (\sphinxstyleliteralemphasis{\sphinxupquote{float}}) \textendash{} The current (true) simulation time in seconds.

\item {} 
\sphinxAtStartPar
\sphinxstyleliteralstrong{\sphinxupquote{estimated\_state}} ({\hyperref[\detokenize{modules/estimation/state:qnav.estimation.state.EstimatedState}]{\sphinxcrossref{\sphinxstyleliteralemphasis{\sphinxupquote{EstimatedState}}}}}) \textendash{} The current estimated state of the simulation.

\item {} 
\sphinxAtStartPar
\sphinxstyleliteralstrong{\sphinxupquote{ground\_truth}} ({\hyperref[\detokenize{modules/waypoints/trajectory:qnav.waypoints.trajectory.GroundTruth}]{\sphinxcrossref{\sphinxstyleliteralemphasis{\sphinxupquote{GroundTruth}}}}}) \textendash{} The current ground truth state of the simulation.

\end{itemize}

\end{description}\end{quote}

\end{fulllineitems}

\index{on\_finish() (qnav.simulation.display.ConsoleDisplay method)@\spxentry{on\_finish()}\spxextra{qnav.simulation.display.ConsoleDisplay method}}

\begin{fulllineitems}
\phantomsection\label{\detokenize{modules/simulation/display:qnav.simulation.display.ConsoleDisplay.on_finish}}
\pysigstartsignatures
\pysiglinewithargsret{\sphinxbfcode{\sphinxupquote{on\_finish}}}{\sphinxparam{\DUrole{n}{timestamp}\DUrole{p}{:}\DUrole{w}{ }\DUrole{n}{float}}\sphinxparamcomma \sphinxparam{\DUrole{n}{estimated\_state}\DUrole{p}{:}\DUrole{w}{ }\DUrole{n}{{\hyperref[\detokenize{modules/estimation/state:qnav.estimation.state.EstimatedState}]{\sphinxcrossref{EstimatedState}}}}}\sphinxparamcomma \sphinxparam{\DUrole{n}{ground\_truth}\DUrole{p}{:}\DUrole{w}{ }\DUrole{n}{{\hyperref[\detokenize{modules/waypoints/trajectory:qnav.waypoints.trajectory.GroundTruth}]{\sphinxcrossref{GroundTruth}}}}}}{}
\pysigstopsignatures
\sphinxAtStartPar
Called at the end of the simulation after final iteration.
Finishes the runtime progress display of the simulation.
\begin{quote}\begin{description}
\sphinxlineitem{Parameters}\begin{itemize}
\item {} 
\sphinxAtStartPar
\sphinxstyleliteralstrong{\sphinxupquote{timestamp}} (\sphinxstyleliteralemphasis{\sphinxupquote{float}}) \textendash{} The final (true) simulation time in seconds.

\item {} 
\sphinxAtStartPar
\sphinxstyleliteralstrong{\sphinxupquote{estimated\_state}} ({\hyperref[\detokenize{modules/estimation/state:qnav.estimation.state.EstimatedState}]{\sphinxcrossref{\sphinxstyleliteralemphasis{\sphinxupquote{EstimatedState}}}}}) \textendash{} The final estimated state of the simulation.

\item {} 
\sphinxAtStartPar
\sphinxstyleliteralstrong{\sphinxupquote{ground\_truth}} ({\hyperref[\detokenize{modules/waypoints/trajectory:qnav.waypoints.trajectory.GroundTruth}]{\sphinxcrossref{\sphinxstyleliteralemphasis{\sphinxupquote{GroundTruth}}}}}) \textendash{} The final ground truth state of the simulation.

\end{itemize}

\end{description}\end{quote}

\end{fulllineitems}


\end{fulllineitems}

\index{ConsoleDisplayPlus (class in qnav.simulation.display)@\spxentry{ConsoleDisplayPlus}\spxextra{class in qnav.simulation.display}}

\begin{fulllineitems}
\phantomsection\label{\detokenize{modules/simulation/display:qnav.simulation.display.ConsoleDisplayPlus}}
\pysigstartsignatures
\pysiglinewithargsret{\sphinxbfcode{\sphinxupquote{class\DUrole{w}{ }}}\sphinxcode{\sphinxupquote{qnav.simulation.display.}}\sphinxbfcode{\sphinxupquote{ConsoleDisplayPlus}}}{\sphinxparam{\DUrole{n}{update\_interval}\DUrole{p}{:}\DUrole{w}{ }\DUrole{n}{float}}\sphinxparamcomma \sphinxparam{\DUrole{n}{max\_dist\_error}\DUrole{p}{:}\DUrole{w}{ }\DUrole{n}{float}\DUrole{w}{ }\DUrole{o}{=}\DUrole{w}{ }\DUrole{default_value}{None}}}{}
\pysigstopsignatures
\sphinxAtStartPar
Extension to the basic displaying class that includes estimation errors.
This subclass also calculates the estimated position error at each
interval and displays it in the console. In addition to this, an
optional threshold can be set for the maximum error distance for
prematurely stopping the simulation.
\index{on\_update() (qnav.simulation.display.ConsoleDisplayPlus method)@\spxentry{on\_update()}\spxextra{qnav.simulation.display.ConsoleDisplayPlus method}}

\begin{fulllineitems}
\phantomsection\label{\detokenize{modules/simulation/display:qnav.simulation.display.ConsoleDisplayPlus.on_update}}
\pysigstartsignatures
\pysiglinewithargsret{\sphinxbfcode{\sphinxupquote{on\_update}}}{\sphinxparam{\DUrole{n}{timestamp}\DUrole{p}{:}\DUrole{w}{ }\DUrole{n}{float}}\sphinxparamcomma \sphinxparam{\DUrole{n}{estimated\_state}\DUrole{p}{:}\DUrole{w}{ }\DUrole{n}{{\hyperref[\detokenize{modules/estimation/state:qnav.estimation.state.EstimatedState}]{\sphinxcrossref{EstimatedState}}}}}\sphinxparamcomma \sphinxparam{\DUrole{n}{ground\_truth}\DUrole{p}{:}\DUrole{w}{ }\DUrole{n}{{\hyperref[\detokenize{modules/waypoints/trajectory:qnav.waypoints.trajectory.GroundTruth}]{\sphinxcrossref{GroundTruth}}}}}}{}
\pysigstopsignatures
\sphinxAtStartPar
Called during the simulation after each fusion update.
Updates the runtime progress display of the simulation.
\begin{quote}\begin{description}
\sphinxlineitem{Parameters}\begin{itemize}
\item {} 
\sphinxAtStartPar
\sphinxstyleliteralstrong{\sphinxupquote{timestamp}} (\sphinxstyleliteralemphasis{\sphinxupquote{float}}) \textendash{} The current (true) simulation time in seconds.

\item {} 
\sphinxAtStartPar
\sphinxstyleliteralstrong{\sphinxupquote{estimated\_state}} ({\hyperref[\detokenize{modules/estimation/state:qnav.estimation.state.EstimatedState}]{\sphinxcrossref{\sphinxstyleliteralemphasis{\sphinxupquote{EstimatedState}}}}}) \textendash{} The current estimated state of the simulation.

\item {} 
\sphinxAtStartPar
\sphinxstyleliteralstrong{\sphinxupquote{ground\_truth}} ({\hyperref[\detokenize{modules/waypoints/trajectory:qnav.waypoints.trajectory.GroundTruth}]{\sphinxcrossref{\sphinxstyleliteralemphasis{\sphinxupquote{GroundTruth}}}}}) \textendash{} The current ground truth state of the simulation.

\end{itemize}

\end{description}\end{quote}

\end{fulllineitems}

\index{on\_finish() (qnav.simulation.display.ConsoleDisplayPlus method)@\spxentry{on\_finish()}\spxextra{qnav.simulation.display.ConsoleDisplayPlus method}}

\begin{fulllineitems}
\phantomsection\label{\detokenize{modules/simulation/display:qnav.simulation.display.ConsoleDisplayPlus.on_finish}}
\pysigstartsignatures
\pysiglinewithargsret{\sphinxbfcode{\sphinxupquote{on\_finish}}}{\sphinxparam{\DUrole{n}{timestamp}\DUrole{p}{:}\DUrole{w}{ }\DUrole{n}{float}}\sphinxparamcomma \sphinxparam{\DUrole{n}{estimated\_state}\DUrole{p}{:}\DUrole{w}{ }\DUrole{n}{{\hyperref[\detokenize{modules/estimation/state:qnav.estimation.state.EstimatedState}]{\sphinxcrossref{EstimatedState}}}}}\sphinxparamcomma \sphinxparam{\DUrole{n}{ground\_truth}\DUrole{p}{:}\DUrole{w}{ }\DUrole{n}{{\hyperref[\detokenize{modules/waypoints/trajectory:qnav.waypoints.trajectory.GroundTruth}]{\sphinxcrossref{GroundTruth}}}}}}{}
\pysigstopsignatures
\sphinxAtStartPar
Called at the end of the simulation after final iteration.
Finishes the runtime progress display of the simulation.
\begin{quote}\begin{description}
\sphinxlineitem{Parameters}\begin{itemize}
\item {} 
\sphinxAtStartPar
\sphinxstyleliteralstrong{\sphinxupquote{timestamp}} (\sphinxstyleliteralemphasis{\sphinxupquote{float}}) \textendash{} The final (true) simulation time in seconds.

\item {} 
\sphinxAtStartPar
\sphinxstyleliteralstrong{\sphinxupquote{estimated\_state}} ({\hyperref[\detokenize{modules/estimation/state:qnav.estimation.state.EstimatedState}]{\sphinxcrossref{\sphinxstyleliteralemphasis{\sphinxupquote{EstimatedState}}}}}) \textendash{} The final estimated state of the simulation.

\item {} 
\sphinxAtStartPar
\sphinxstyleliteralstrong{\sphinxupquote{ground\_truth}} ({\hyperref[\detokenize{modules/waypoints/trajectory:qnav.waypoints.trajectory.GroundTruth}]{\sphinxcrossref{\sphinxstyleliteralemphasis{\sphinxupquote{GroundTruth}}}}}) \textendash{} The final ground truth state of the simulation.

\end{itemize}

\end{description}\end{quote}

\end{fulllineitems}


\end{fulllineitems}

\index{MultilineDisplay (class in qnav.simulation.display)@\spxentry{MultilineDisplay}\spxextra{class in qnav.simulation.display}}

\begin{fulllineitems}
\phantomsection\label{\detokenize{modules/simulation/display:qnav.simulation.display.MultilineDisplay}}
\pysigstartsignatures
\pysiglinewithargsret{\sphinxbfcode{\sphinxupquote{class\DUrole{w}{ }}}\sphinxcode{\sphinxupquote{qnav.simulation.display.}}\sphinxbfcode{\sphinxupquote{MultilineDisplay}}}{\sphinxparam{\DUrole{n}{update\_interval}\DUrole{p}{:}\DUrole{w}{ }\DUrole{n}{float}}\sphinxparamcomma \sphinxparam{\DUrole{n}{max\_dist\_error}\DUrole{p}{:}\DUrole{w}{ }\DUrole{n}{float}\DUrole{w}{ }\DUrole{o}{=}\DUrole{w}{ }\DUrole{default_value}{None}}}{}
\pysigstopsignatures\index{on\_start() (qnav.simulation.display.MultilineDisplay method)@\spxentry{on\_start()}\spxextra{qnav.simulation.display.MultilineDisplay method}}

\begin{fulllineitems}
\phantomsection\label{\detokenize{modules/simulation/display:qnav.simulation.display.MultilineDisplay.on_start}}
\pysigstartsignatures
\pysiglinewithargsret{\sphinxbfcode{\sphinxupquote{on\_start}}}{\sphinxparam{\DUrole{n}{trajectory}\DUrole{p}{:}\DUrole{w}{ }\DUrole{n}{{\hyperref[\detokenize{modules/waypoints/trajectory:qnav.waypoints.trajectory.Trajectory}]{\sphinxcrossref{Trajectory}}}}}}{}
\pysigstopsignatures
\sphinxAtStartPar
Called at the start of the simulation before first iteration.
Used to obtain the start and end times of the trajectory.
\begin{quote}\begin{description}
\sphinxlineitem{Parameters}
\sphinxAtStartPar
\sphinxstyleliteralstrong{\sphinxupquote{trajectory}} ({\hyperref[\detokenize{modules/waypoints/trajectory:qnav.waypoints.trajectory.Trajectory}]{\sphinxcrossref{\sphinxstyleliteralemphasis{\sphinxupquote{Trajectory}}}}}) \textendash{} The ground truth trajectory instance.

\end{description}\end{quote}

\end{fulllineitems}

\index{on\_update() (qnav.simulation.display.MultilineDisplay method)@\spxentry{on\_update()}\spxextra{qnav.simulation.display.MultilineDisplay method}}

\begin{fulllineitems}
\phantomsection\label{\detokenize{modules/simulation/display:qnav.simulation.display.MultilineDisplay.on_update}}
\pysigstartsignatures
\pysiglinewithargsret{\sphinxbfcode{\sphinxupquote{on\_update}}}{\sphinxparam{\DUrole{n}{timestamp}\DUrole{p}{:}\DUrole{w}{ }\DUrole{n}{float}}\sphinxparamcomma \sphinxparam{\DUrole{n}{estimated\_state}\DUrole{p}{:}\DUrole{w}{ }\DUrole{n}{{\hyperref[\detokenize{modules/estimation/state:qnav.estimation.state.EstimatedState}]{\sphinxcrossref{EstimatedState}}}}}\sphinxparamcomma \sphinxparam{\DUrole{n}{ground\_truth}\DUrole{p}{:}\DUrole{w}{ }\DUrole{n}{{\hyperref[\detokenize{modules/waypoints/trajectory:qnav.waypoints.trajectory.GroundTruth}]{\sphinxcrossref{GroundTruth}}}}}}{}
\pysigstopsignatures
\sphinxAtStartPar
Called during the simulation after each fusion update.
Updates the runtime progress display of the simulation.
\begin{quote}\begin{description}
\sphinxlineitem{Parameters}\begin{itemize}
\item {} 
\sphinxAtStartPar
\sphinxstyleliteralstrong{\sphinxupquote{timestamp}} (\sphinxstyleliteralemphasis{\sphinxupquote{float}}) \textendash{} The current (true) simulation time in seconds.

\item {} 
\sphinxAtStartPar
\sphinxstyleliteralstrong{\sphinxupquote{estimated\_state}} ({\hyperref[\detokenize{modules/estimation/state:qnav.estimation.state.EstimatedState}]{\sphinxcrossref{\sphinxstyleliteralemphasis{\sphinxupquote{EstimatedState}}}}}) \textendash{} The current estimated state of the simulation.

\item {} 
\sphinxAtStartPar
\sphinxstyleliteralstrong{\sphinxupquote{ground\_truth}} ({\hyperref[\detokenize{modules/waypoints/trajectory:qnav.waypoints.trajectory.GroundTruth}]{\sphinxcrossref{\sphinxstyleliteralemphasis{\sphinxupquote{GroundTruth}}}}}) \textendash{} The current ground truth state of the simulation.

\end{itemize}

\end{description}\end{quote}

\end{fulllineitems}

\index{on\_finish() (qnav.simulation.display.MultilineDisplay method)@\spxentry{on\_finish()}\spxextra{qnav.simulation.display.MultilineDisplay method}}

\begin{fulllineitems}
\phantomsection\label{\detokenize{modules/simulation/display:qnav.simulation.display.MultilineDisplay.on_finish}}
\pysigstartsignatures
\pysiglinewithargsret{\sphinxbfcode{\sphinxupquote{on\_finish}}}{\sphinxparam{\DUrole{n}{timestamp}\DUrole{p}{:}\DUrole{w}{ }\DUrole{n}{float}}\sphinxparamcomma \sphinxparam{\DUrole{n}{estimated\_state}\DUrole{p}{:}\DUrole{w}{ }\DUrole{n}{{\hyperref[\detokenize{modules/estimation/state:qnav.estimation.state.EstimatedState}]{\sphinxcrossref{EstimatedState}}}}}\sphinxparamcomma \sphinxparam{\DUrole{n}{ground\_truth}\DUrole{p}{:}\DUrole{w}{ }\DUrole{n}{{\hyperref[\detokenize{modules/waypoints/trajectory:qnav.waypoints.trajectory.GroundTruth}]{\sphinxcrossref{GroundTruth}}}}}}{}
\pysigstopsignatures
\sphinxAtStartPar
Called at the end of the simulation after final iteration.
Finishes the runtime progress display of the simulation.
\begin{quote}\begin{description}
\sphinxlineitem{Parameters}\begin{itemize}
\item {} 
\sphinxAtStartPar
\sphinxstyleliteralstrong{\sphinxupquote{timestamp}} (\sphinxstyleliteralemphasis{\sphinxupquote{float}}) \textendash{} The final (true) simulation time in seconds.

\item {} 
\sphinxAtStartPar
\sphinxstyleliteralstrong{\sphinxupquote{estimated\_state}} ({\hyperref[\detokenize{modules/estimation/state:qnav.estimation.state.EstimatedState}]{\sphinxcrossref{\sphinxstyleliteralemphasis{\sphinxupquote{EstimatedState}}}}}) \textendash{} The final estimated state of the simulation.

\item {} 
\sphinxAtStartPar
\sphinxstyleliteralstrong{\sphinxupquote{ground\_truth}} ({\hyperref[\detokenize{modules/waypoints/trajectory:qnav.waypoints.trajectory.GroundTruth}]{\sphinxcrossref{\sphinxstyleliteralemphasis{\sphinxupquote{GroundTruth}}}}}) \textendash{} The final ground truth state of the simulation.

\end{itemize}

\end{description}\end{quote}

\end{fulllineitems}


\end{fulllineitems}



\subsection{Util}
\label{\detokenize{usage/packages:util}}
\sphinxAtStartPar
Modules relating to commonly used utility functions:

\sphinxstepscope


\subsubsection{qnav.util.constants}
\label{\detokenize{modules/util/constants:module-qnav.util.constants}}\label{\detokenize{modules/util/constants:qnav-util-constants}}\label{\detokenize{modules/util/constants::doc}}\index{module@\spxentry{module}!qnav.util.constants@\spxentry{qnav.util.constants}}\index{qnav.util.constants@\spxentry{qnav.util.constants}!module@\spxentry{module}}

\paragraph{constants.py}
\label{\detokenize{modules/util/constants:constants-py}}\begin{quote}\begin{description}
\sphinxlineitem{summary}
\sphinxAtStartPar
A list of calculation constants that are across the toolbox.
This shared module contains a collection of geological constants that are
repeatedly used by multiple functions within the toolbox

\sphinxlineitem{author}
\sphinxAtStartPar
Michael Wright \sphinxhyphen{} \sphinxhref{mailto:mjwright@liverpool.ac.uk}{mjwright@liverpool.ac.uk}

\sphinxlineitem{version}
\sphinxAtStartPar
1.0.0 \sphinxhyphen{} First implementation of this module.

\end{description}\end{quote}
\index{G (in module qnav.util.constants)@\spxentry{G}\spxextra{in module qnav.util.constants}}

\begin{fulllineitems}
\phantomsection\label{\detokenize{modules/util/constants:qnav.util.constants.G}}
\pysigstartsignatures
\pysigline{\sphinxcode{\sphinxupquote{qnav.util.constants.}}\sphinxbfcode{\sphinxupquote{G}}\sphinxbfcode{\sphinxupquote{\DUrole{w}{ }\DUrole{p}{=}\DUrole{w}{ }9.81}}}
\pysigstopsignatures
\sphinxAtStartPar
\(g\) \sphinxhyphen{} The average gravitation acceleration.

\end{fulllineitems}

\index{F (in module qnav.util.constants)@\spxentry{F}\spxextra{in module qnav.util.constants}}

\begin{fulllineitems}
\phantomsection\label{\detokenize{modules/util/constants:qnav.util.constants.F}}
\pysigstartsignatures
\pysigline{\sphinxcode{\sphinxupquote{qnav.util.constants.}}\sphinxbfcode{\sphinxupquote{F}}\sphinxbfcode{\sphinxupquote{\DUrole{w}{ }\DUrole{p}{=}\DUrole{w}{ }0.0033528106647474805}}}
\pysigstopsignatures
\sphinxAtStartPar
\(f\) \sphinxhyphen{} Flattening measure of the WGS84 spheroid.

\end{fulllineitems}

\index{POLAR\_AXIS\_A (in module qnav.util.constants)@\spxentry{POLAR\_AXIS\_A}\spxextra{in module qnav.util.constants}}

\begin{fulllineitems}
\phantomsection\label{\detokenize{modules/util/constants:qnav.util.constants.POLAR_AXIS_A}}
\pysigstartsignatures
\pysigline{\sphinxcode{\sphinxupquote{qnav.util.constants.}}\sphinxbfcode{\sphinxupquote{POLAR\_AXIS\_A}}\sphinxbfcode{\sphinxupquote{\DUrole{w}{ }\DUrole{p}{=}\DUrole{w}{ }6378137.0}}}
\pysigstopsignatures
\sphinxAtStartPar
\(a\) \sphinxhyphen{} Semi\sphinxhyphen{}major axis equator radius (m).

\end{fulllineitems}

\index{R\_E (in module qnav.util.constants)@\spxentry{R\_E}\spxextra{in module qnav.util.constants}}

\begin{fulllineitems}
\phantomsection\label{\detokenize{modules/util/constants:qnav.util.constants.R_E}}
\pysigstartsignatures
\pysigline{\sphinxcode{\sphinxupquote{qnav.util.constants.}}\sphinxbfcode{\sphinxupquote{R\_E}}\sphinxbfcode{\sphinxupquote{\DUrole{w}{ }\DUrole{p}{=}\DUrole{w}{ }6367444.65712259}}}
\pysigstopsignatures
\sphinxAtStartPar
\(R_e\) \sphinxhyphen{} The average radius of the Earth (m).

\end{fulllineitems}

\index{ECC\_SQ (in module qnav.util.constants)@\spxentry{ECC\_SQ}\spxextra{in module qnav.util.constants}}

\begin{fulllineitems}
\phantomsection\label{\detokenize{modules/util/constants:qnav.util.constants.ECC_SQ}}
\pysigstartsignatures
\pysigline{\sphinxcode{\sphinxupquote{qnav.util.constants.}}\sphinxbfcode{\sphinxupquote{ECC\_SQ}}\sphinxbfcode{\sphinxupquote{\DUrole{w}{ }\DUrole{p}{=}\DUrole{w}{ }0.006694379990141316}}}
\pysigstopsignatures
\sphinxAtStartPar
\(e^2\) \sphinxhyphen{} The squared first eccentricity value for the Earth.
\begin{quote}\begin{description}
\sphinxlineitem{Type}
\sphinxAtStartPar
\#

\end{description}\end{quote}

\end{fulllineitems}

\index{ECC (in module qnav.util.constants)@\spxentry{ECC}\spxextra{in module qnav.util.constants}}

\begin{fulllineitems}
\phantomsection\label{\detokenize{modules/util/constants:qnav.util.constants.ECC}}
\pysigstartsignatures
\pysigline{\sphinxcode{\sphinxupquote{qnav.util.constants.}}\sphinxbfcode{\sphinxupquote{ECC}}\sphinxbfcode{\sphinxupquote{\DUrole{w}{ }\DUrole{p}{=}\DUrole{w}{ }0.08181919084262149}}}
\pysigstopsignatures
\sphinxAtStartPar
\(e\) \sphinxhyphen{} The first eccentricity value for the Earth.

\end{fulllineitems}

\index{ECC\_PRIME (in module qnav.util.constants)@\spxentry{ECC\_PRIME}\spxextra{in module qnav.util.constants}}

\begin{fulllineitems}
\phantomsection\label{\detokenize{modules/util/constants:qnav.util.constants.ECC_PRIME}}
\pysigstartsignatures
\pysigline{\sphinxcode{\sphinxupquote{qnav.util.constants.}}\sphinxbfcode{\sphinxupquote{ECC\_PRIME}}\sphinxbfcode{\sphinxupquote{\DUrole{w}{ }\DUrole{p}{=}\DUrole{w}{ }0.08209443794969568}}}
\pysigstopsignatures
\sphinxAtStartPar
\(e'\) \sphinxhyphen{} The second eccentricity value for the Earth.

\end{fulllineitems}

\index{G\_E (in module qnav.util.constants)@\spxentry{G\_E}\spxextra{in module qnav.util.constants}}

\begin{fulllineitems}
\phantomsection\label{\detokenize{modules/util/constants:qnav.util.constants.G_E}}
\pysigstartsignatures
\pysigline{\sphinxcode{\sphinxupquote{qnav.util.constants.}}\sphinxbfcode{\sphinxupquote{G\_E}}\sphinxbfcode{\sphinxupquote{\DUrole{w}{ }\DUrole{p}{=}\DUrole{w}{ }9.78032533590389}}}
\pysigstopsignatures
\sphinxAtStartPar
\({\gamma}_e\) \sphinxhyphen{} Normal gravity at Equator (m/s**2).

\end{fulllineitems}

\index{G\_P (in module qnav.util.constants)@\spxentry{G\_P}\spxextra{in module qnav.util.constants}}

\begin{fulllineitems}
\phantomsection\label{\detokenize{modules/util/constants:qnav.util.constants.G_P}}
\pysigstartsignatures
\pysigline{\sphinxcode{\sphinxupquote{qnav.util.constants.}}\sphinxbfcode{\sphinxupquote{G\_P}}\sphinxbfcode{\sphinxupquote{\DUrole{w}{ }\DUrole{p}{=}\DUrole{w}{ }9.832184937863401}}}
\pysigstopsignatures
\sphinxAtStartPar
\({\gamma}_b\) \sphinxhyphen{} Normal gravity at Poles (m/s**2).

\end{fulllineitems}

\index{OMEGA\_E (in module qnav.util.constants)@\spxentry{OMEGA\_E}\spxextra{in module qnav.util.constants}}

\begin{fulllineitems}
\phantomsection\label{\detokenize{modules/util/constants:qnav.util.constants.OMEGA_E}}
\pysigstartsignatures
\pysigline{\sphinxcode{\sphinxupquote{qnav.util.constants.}}\sphinxbfcode{\sphinxupquote{OMEGA\_E}}\sphinxbfcode{\sphinxupquote{\DUrole{w}{ }\DUrole{p}{=}\DUrole{w}{ }7.2921159e\sphinxhyphen{}05}}}
\pysigstopsignatures
\sphinxAtStartPar
\({\omega}_e\) \sphinxhyphen{} Approximate rotation rate of the Earth (m/s**2).

\end{fulllineitems}

\index{GM (in module qnav.util.constants)@\spxentry{GM}\spxextra{in module qnav.util.constants}}

\begin{fulllineitems}
\phantomsection\label{\detokenize{modules/util/constants:qnav.util.constants.GM}}
\pysigstartsignatures
\pysigline{\sphinxcode{\sphinxupquote{qnav.util.constants.}}\sphinxbfcode{\sphinxupquote{GM}}\sphinxbfcode{\sphinxupquote{\DUrole{w}{ }\DUrole{p}{=}\DUrole{w}{ }398600441880000.0}}}
\pysigstopsignatures
\sphinxAtStartPar
\(GM\) \sphinxhyphen{} The Geocentric gravitational constant (m**3/s**2)

\end{fulllineitems}

\index{M (in module qnav.util.constants)@\spxentry{M}\spxextra{in module qnav.util.constants}}

\begin{fulllineitems}
\phantomsection\label{\detokenize{modules/util/constants:qnav.util.constants.M}}
\pysigstartsignatures
\pysigline{\sphinxcode{\sphinxupquote{qnav.util.constants.}}\sphinxbfcode{\sphinxupquote{M}}\sphinxbfcode{\sphinxupquote{\DUrole{w}{ }\DUrole{p}{=}\DUrole{w}{ }0.0034497873577006006}}}
\pysigstopsignatures
\sphinxAtStartPar
\(M\) \sphinxhyphen{} Somigliana constant, used for calculation normal gravity.

\end{fulllineitems}


\sphinxstepscope


\subsubsection{qnav.util.connections}
\label{\detokenize{modules/util/connections:module-qnav.util.connections}}\label{\detokenize{modules/util/connections:qnav-util-connections}}\label{\detokenize{modules/util/connections::doc}}\index{module@\spxentry{module}!qnav.util.connections@\spxentry{qnav.util.connections}}\index{qnav.util.connections@\spxentry{qnav.util.connections}!module@\spxentry{module}}\index{download\_file() (in module qnav.util.connections)@\spxentry{download\_file()}\spxextra{in module qnav.util.connections}}

\begin{fulllineitems}
\phantomsection\label{\detokenize{modules/util/connections:qnav.util.connections.download_file}}
\pysigstartsignatures
\pysiglinewithargsret{\sphinxcode{\sphinxupquote{qnav.util.connections.}}\sphinxbfcode{\sphinxupquote{download\_file}}}{\sphinxparam{\DUrole{n}{url}\DUrole{p}{:}\DUrole{w}{ }\DUrole{n}{str}}\sphinxparamcomma \sphinxparam{\DUrole{n}{dest}\DUrole{p}{:}\DUrole{w}{ }\DUrole{n}{Path}}}{}
\pysigstopsignatures
\sphinxAtStartPar
Downloads a file from a URL and writes it to destination file.
Commonly used by the toolbox for downloading database files from
selected web\sphinxhyphen{}servers.
\begin{quote}\begin{description}
\sphinxlineitem{Parameters}\begin{itemize}
\item {} 
\sphinxAtStartPar
\sphinxstyleliteralstrong{\sphinxupquote{url}} (\sphinxstyleliteralemphasis{\sphinxupquote{str}}) \textendash{} The URL of the file to be downloaded.

\item {} 
\sphinxAtStartPar
\sphinxstyleliteralstrong{\sphinxupquote{dest}} (\sphinxstyleliteralemphasis{\sphinxupquote{Path}}) \textendash{} The file to save the downloaded file as.

\end{itemize}

\end{description}\end{quote}

\end{fulllineitems}


\sphinxstepscope


\subsubsection{qnav.util.system}
\label{\detokenize{modules/util/system:module-qnav.util.system}}\label{\detokenize{modules/util/system:qnav-util-system}}\label{\detokenize{modules/util/system::doc}}\index{module@\spxentry{module}!qnav.util.system@\spxentry{qnav.util.system}}\index{qnav.util.system@\spxentry{qnav.util.system}!module@\spxentry{module}}

\paragraph{common.py}
\label{\detokenize{modules/util/system:common-py}}\begin{quote}\begin{description}
\sphinxlineitem{summary}
\sphinxAtStartPar
Provides commonly used project functions.
This module contains commonly used (non\sphinxhyphen{}calculation) functions that are
called throughout the project. These are namely for providing information
about the project and the current platform.

\sphinxlineitem{author}
\sphinxAtStartPar
Michael Wright \sphinxhyphen{} \sphinxhref{mailto:mjwright@liverpool.ac.uk}{mjwright@liverpool.ac.uk}

\sphinxlineitem{version}
\sphinxAtStartPar
1.0.0 \sphinxhyphen{} First full implementation of module

\end{description}\end{quote}
\index{get\_toolbox\_version() (in module qnav.util.system)@\spxentry{get\_toolbox\_version()}\spxextra{in module qnav.util.system}}

\begin{fulllineitems}
\phantomsection\label{\detokenize{modules/util/system:qnav.util.system.get_toolbox_version}}
\pysigstartsignatures
\pysiglinewithargsret{\sphinxcode{\sphinxupquote{qnav.util.system.}}\sphinxbfcode{\sphinxupquote{get\_toolbox\_version}}}{}{{ $\rightarrow$ str}}
\pysigstopsignatures
\sphinxAtStartPar
A simple function to return the current version number of the toolbox.
This is useful for reporting information that may dependent on the build
of the software, such as generated results or reported errors.
\begin{quote}\begin{description}
\sphinxlineitem{Returns}
\sphinxAtStartPar
The decimal version number for the toolbox.

\sphinxlineitem{Return type}
\sphinxAtStartPar
str

\end{description}\end{quote}

\end{fulllineitems}

\index{get\_system\_name() (in module qnav.util.system)@\spxentry{get\_system\_name()}\spxextra{in module qnav.util.system}}

\begin{fulllineitems}
\phantomsection\label{\detokenize{modules/util/system:qnav.util.system.get_system_name}}
\pysigstartsignatures
\pysiglinewithargsret{\sphinxcode{\sphinxupquote{qnav.util.system.}}\sphinxbfcode{\sphinxupquote{get\_system\_name}}}{\sphinxparam{\DUrole{n}{return\_min}\DUrole{p}{:}\DUrole{w}{ }\DUrole{n}{bool}\DUrole{w}{ }\DUrole{o}{=}\DUrole{w}{ }\DUrole{default_value}{False}}}{{ $\rightarrow$ str}}
\pysigstopsignatures
\sphinxAtStartPar
Return a short string for identifying the host operating system. This is
namely used for identifying the operating system in a human\sphinxhyphen{}readable
format. This is useful for identifying the type of system used when
running an experiment.
\begin{quote}\begin{description}
\sphinxlineitem{Parameters}
\sphinxAtStartPar
\sphinxstyleliteralstrong{\sphinxupquote{return\_min}} (\sphinxstyleliteralemphasis{\sphinxupquote{bool}}) \textendash{} If set to true, only the operating system name
and major version number will be returned.

\sphinxlineitem{Returns}
\sphinxAtStartPar
A short string used for identifying the host operating system.

\sphinxlineitem{Return type}
\sphinxAtStartPar
str

\end{description}\end{quote}

\end{fulllineitems}

\index{get\_platform\_details() (in module qnav.util.system)@\spxentry{get\_platform\_details()}\spxextra{in module qnav.util.system}}

\begin{fulllineitems}
\phantomsection\label{\detokenize{modules/util/system:qnav.util.system.get_platform_details}}
\pysigstartsignatures
\pysiglinewithargsret{\sphinxcode{\sphinxupquote{qnav.util.system.}}\sphinxbfcode{\sphinxupquote{get\_platform\_details}}}{}{{ $\rightarrow$ str}}
\pysigstopsignatures
\sphinxAtStartPar
A simple function that returns a description of the platform.
This function returns a string including the python version, processor
platform architecture of the host system.
\begin{quote}\begin{description}
\sphinxlineitem{Returns}
\sphinxAtStartPar
A string describing the platform.

\sphinxlineitem{Return type}
\sphinxAtStartPar
str

\end{description}\end{quote}

\end{fulllineitems}

\index{is\_window\_platform() (in module qnav.util.system)@\spxentry{is\_window\_platform()}\spxextra{in module qnav.util.system}}

\begin{fulllineitems}
\phantomsection\label{\detokenize{modules/util/system:qnav.util.system.is_window_platform}}
\pysigstartsignatures
\pysiglinewithargsret{\sphinxcode{\sphinxupquote{qnav.util.system.}}\sphinxbfcode{\sphinxupquote{is\_window\_platform}}}{}{{ $\rightarrow$ bool}}
\pysigstopsignatures
\sphinxAtStartPar
A simple function which will return true if the system’s operating system
is Microsoft Windows based. Otherwise, this function will return false.
\begin{quote}\begin{description}
\sphinxlineitem{Returns}
\sphinxAtStartPar
Returns true if the host operating system is Windows based.

\sphinxlineitem{Return type}
\sphinxAtStartPar
bool

\end{description}\end{quote}

\end{fulllineitems}

\index{is\_mac\_platform() (in module qnav.util.system)@\spxentry{is\_mac\_platform()}\spxextra{in module qnav.util.system}}

\begin{fulllineitems}
\phantomsection\label{\detokenize{modules/util/system:qnav.util.system.is_mac_platform}}
\pysigstartsignatures
\pysiglinewithargsret{\sphinxcode{\sphinxupquote{qnav.util.system.}}\sphinxbfcode{\sphinxupquote{is\_mac\_platform}}}{}{{ $\rightarrow$ bool}}
\pysigstopsignatures
\sphinxAtStartPar
A simple function which will return true if the system’s operating system
is Mac\sphinxhyphen{}OS based. Otherwise, this function will return false.
\begin{quote}\begin{description}
\sphinxlineitem{Returns}
\sphinxAtStartPar
Returns true if the host operating system is Mac\sphinxhyphen{}OS based.

\sphinxlineitem{Return type}
\sphinxAtStartPar
bool

\end{description}\end{quote}

\end{fulllineitems}

\index{is\_linux\_platform() (in module qnav.util.system)@\spxentry{is\_linux\_platform()}\spxextra{in module qnav.util.system}}

\begin{fulllineitems}
\phantomsection\label{\detokenize{modules/util/system:qnav.util.system.is_linux_platform}}
\pysigstartsignatures
\pysiglinewithargsret{\sphinxcode{\sphinxupquote{qnav.util.system.}}\sphinxbfcode{\sphinxupquote{is\_linux\_platform}}}{}{{ $\rightarrow$ bool}}
\pysigstopsignatures
\sphinxAtStartPar
A simple function which will return true if the system’s operating system
is Linux/Unix based. Otherwise, this function will return false.
\begin{quote}\begin{description}
\sphinxlineitem{Returns}
\sphinxAtStartPar
Returns true if the host operating system is Linux based.

\sphinxlineitem{Return type}
\sphinxAtStartPar
bool

\end{description}\end{quote}

\end{fulllineitems}


\sphinxstepscope


\subsubsection{qnav.util.rng}
\label{\detokenize{modules/util/rng:module-qnav.util.rng}}\label{\detokenize{modules/util/rng:qnav-util-rng}}\label{\detokenize{modules/util/rng::doc}}\index{module@\spxentry{module}!qnav.util.rng@\spxentry{qnav.util.rng}}\index{qnav.util.rng@\spxentry{qnav.util.rng}!module@\spxentry{module}}

\paragraph{rng.py}
\label{\detokenize{modules/util/rng:rng-py}}\begin{quote}\begin{description}
\sphinxlineitem{summary}
\sphinxAtStartPar
Provides alternative means for pseudo\sphinxhyphen{}random number generation.
This module contains implementations for controlled random number
generation, subject to realistic factors. Currently, this only
supplies gaussian\sphinxhyphen{}markov random walk based number generation.

\sphinxlineitem{authors}
\begin{DUlineblock}{0em}
\item[] Michael Wright \sphinxhyphen{} \sphinxhref{mailto:mjwright@liverpool.ac.uk}{mjwright@liverpool.ac.uk}
\end{DUlineblock}

\sphinxlineitem{version}
\sphinxAtStartPar
1.0.0 \sphinxhyphen{} First implementation of module.

\end{description}\end{quote}
\index{GaussianMarkov (class in qnav.util.rng)@\spxentry{GaussianMarkov}\spxextra{class in qnav.util.rng}}

\begin{fulllineitems}
\phantomsection\label{\detokenize{modules/util/rng:qnav.util.rng.GaussianMarkov}}
\pysigstartsignatures
\pysiglinewithargsret{\sphinxbfcode{\sphinxupquote{class\DUrole{w}{ }}}\sphinxcode{\sphinxupquote{qnav.util.rng.}}\sphinxbfcode{\sphinxupquote{GaussianMarkov}}}{\sphinxparam{\DUrole{n}{corr\_time}\DUrole{p}{:}\DUrole{w}{ }\DUrole{n}{float}}\sphinxparamcomma \sphinxparam{\DUrole{n}{psd}\DUrole{p}{:}\DUrole{w}{ }\DUrole{n}{float}}\sphinxparamcomma \sphinxparam{\DUrole{n}{output\_shape}\DUrole{p}{:}\DUrole{w}{ }\DUrole{n}{int\DUrole{w}{ }\DUrole{p}{|}\DUrole{w}{ }tuple\DUrole{p}{{[}}int\DUrole{p}{{]}}}\DUrole{w}{ }\DUrole{o}{=}\DUrole{w}{ }\DUrole{default_value}{1}}\sphinxparamcomma \sphinxparam{\DUrole{n}{start\_time}\DUrole{p}{:}\DUrole{w}{ }\DUrole{n}{float}\DUrole{w}{ }\DUrole{o}{=}\DUrole{w}{ }\DUrole{default_value}{0}}\sphinxparamcomma \sphinxparam{\DUrole{n}{seed}\DUrole{p}{:}\DUrole{w}{ }\DUrole{n}{int}\DUrole{w}{ }\DUrole{o}{=}\DUrole{w}{ }\DUrole{default_value}{None}}}{}
\pysigstopsignatures
\sphinxAtStartPar
Gaussian Markov First\sphinxhyphen{}Order Random Walk Model.
Generates controlled random numbers following a Gaussian Markov random
walk. Can be used to replace typical normal random number generation.
\index{get\_random() (qnav.util.rng.GaussianMarkov method)@\spxentry{get\_random()}\spxextra{qnav.util.rng.GaussianMarkov method}}

\begin{fulllineitems}
\phantomsection\label{\detokenize{modules/util/rng:qnav.util.rng.GaussianMarkov.get_random}}
\pysigstartsignatures
\pysiglinewithargsret{\sphinxbfcode{\sphinxupquote{get\_random}}}{\sphinxparam{\DUrole{n}{clock\_time}\DUrole{p}{:}\DUrole{w}{ }\DUrole{n}{float}}}{{ $\rightarrow$ float\DUrole{w}{ }\DUrole{p}{|}\DUrole{w}{ }ndarray}}
\pysigstopsignatures
\sphinxAtStartPar
Generates random number(s) for requested time and shape.
Calling this method continues the random walk to the given time in
seconds. Calling with a time earlier than that of the last
measurement will result in a value error due to illegal state.
\begin{quote}\begin{description}
\sphinxlineitem{Parameters}
\sphinxAtStartPar
\sphinxstyleliteralstrong{\sphinxupquote{clock\_time}} (\sphinxstyleliteralemphasis{\sphinxupquote{float}}) \textendash{} The current simulation clock time.

\sphinxlineitem{Returns}
\sphinxAtStartPar
Randomly generated number(s).

\sphinxlineitem{Return type}
\sphinxAtStartPar
float | np.ndarray

\end{description}\end{quote}

\end{fulllineitems}

\index{reset() (qnav.util.rng.GaussianMarkov method)@\spxentry{reset()}\spxextra{qnav.util.rng.GaussianMarkov method}}

\begin{fulllineitems}
\phantomsection\label{\detokenize{modules/util/rng:qnav.util.rng.GaussianMarkov.reset}}
\pysigstartsignatures
\pysiglinewithargsret{\sphinxbfcode{\sphinxupquote{reset}}}{\sphinxparam{\DUrole{n}{clock\_time}\DUrole{p}{:}\DUrole{w}{ }\DUrole{n}{float}\DUrole{w}{ }\DUrole{o}{=}\DUrole{w}{ }\DUrole{default_value}{0}}}{{ $\rightarrow$ None}}
\pysigstopsignatures
\sphinxAtStartPar
Resets the internal state of the random walk, clearing previous
history and information.
\begin{quote}\begin{description}
\sphinxlineitem{Parameters}
\sphinxAtStartPar
\sphinxstyleliteralstrong{\sphinxupquote{clock\_time}} (\sphinxstyleliteralemphasis{\sphinxupquote{float}}) \textendash{} Optional new initialisation clock time.
By default this is 0 seconds.

\end{description}\end{quote}

\end{fulllineitems}

\index{output\_shape (qnav.util.rng.GaussianMarkov property)@\spxentry{output\_shape}\spxextra{qnav.util.rng.GaussianMarkov property}}

\begin{fulllineitems}
\phantomsection\label{\detokenize{modules/util/rng:qnav.util.rng.GaussianMarkov.output_shape}}
\pysigstartsignatures
\pysigline{\sphinxbfcode{\sphinxupquote{property\DUrole{w}{ }}}\sphinxbfcode{\sphinxupquote{output\_shape}}\sphinxbfcode{\sphinxupquote{\DUrole{p}{:}\DUrole{w}{ }int\DUrole{w}{ }\DUrole{p}{|}\DUrole{w}{ }tuple\DUrole{p}{{[}}int\DUrole{p}{{]}}}}}
\pysigstopsignatures
\sphinxAtStartPar
Returns the output shape numbers are being generated for.
This is the dimensions of the output of the number generator.
\begin{quote}\begin{description}
\sphinxlineitem{Returns}
\sphinxAtStartPar
Output shape for random number generation.

\sphinxlineitem{Return type}
\sphinxAtStartPar
int | tuple{[}int{]}

\end{description}\end{quote}

\end{fulllineitems}

\index{last\_value (qnav.util.rng.GaussianMarkov property)@\spxentry{last\_value}\spxextra{qnav.util.rng.GaussianMarkov property}}

\begin{fulllineitems}
\phantomsection\label{\detokenize{modules/util/rng:qnav.util.rng.GaussianMarkov.last_value}}
\pysigstartsignatures
\pysigline{\sphinxbfcode{\sphinxupquote{property\DUrole{w}{ }}}\sphinxbfcode{\sphinxupquote{last\_value}}\sphinxbfcode{\sphinxupquote{\DUrole{p}{:}\DUrole{w}{ }float\DUrole{w}{ }\DUrole{p}{|}\DUrole{w}{ }ndarray\DUrole{w}{ }\DUrole{p}{|}\DUrole{w}{ }None}}}
\pysigstopsignatures
\sphinxAtStartPar
Returns the last random value generated.
\begin{quote}\begin{description}
\sphinxlineitem{Returns}
\sphinxAtStartPar
Last value returned by generator.

\sphinxlineitem{Return type}
\sphinxAtStartPar
float | np.ndarray | None

\end{description}\end{quote}

\end{fulllineitems}

\index{last\_time (qnav.util.rng.GaussianMarkov property)@\spxentry{last\_time}\spxextra{qnav.util.rng.GaussianMarkov property}}

\begin{fulllineitems}
\phantomsection\label{\detokenize{modules/util/rng:qnav.util.rng.GaussianMarkov.last_time}}
\pysigstartsignatures
\pysigline{\sphinxbfcode{\sphinxupquote{property\DUrole{w}{ }}}\sphinxbfcode{\sphinxupquote{last\_time}}\sphinxbfcode{\sphinxupquote{\DUrole{p}{:}\DUrole{w}{ }float}}}
\pysigstopsignatures
\sphinxAtStartPar
Returns time of last number(s) generated in seconds.
\begin{quote}\begin{description}
\sphinxlineitem{Returns}
\sphinxAtStartPar
Time of last number(s) generated.

\sphinxlineitem{Return type}
\sphinxAtStartPar
float

\end{description}\end{quote}

\end{fulllineitems}


\end{fulllineitems}



\subsection{Waypoints}
\label{\detokenize{usage/packages:waypoints}}
\sphinxAtStartPar
Modules related to trajectory generation:

\sphinxstepscope


\subsubsection{qnav.waypoints.generation}
\label{\detokenize{modules/waypoints/generation:module-qnav.waypoints.generation}}\label{\detokenize{modules/waypoints/generation:qnav-waypoints-generation}}\label{\detokenize{modules/waypoints/generation::doc}}\index{module@\spxentry{module}!qnav.waypoints.generation@\spxentry{qnav.waypoints.generation}}\index{qnav.waypoints.generation@\spxentry{qnav.waypoints.generation}!module@\spxentry{module}}

\paragraph{generation.py}
\label{\detokenize{modules/waypoints/generation:generation-py}}\begin{quote}\begin{description}
\sphinxlineitem{summary}
\sphinxAtStartPar
Handles all interpolation and generation of waypoint trajectory data.
This module contains a collection of functions for generating waypoint
trajectories for various vehicle models. In addition to this, waypoint
data can be imported from CSV files or binary numpy files.

\sphinxlineitem{authors}
\begin{DUlineblock}{0em}
\item[] Prof. Jason Ralph \sphinxhyphen{} \sphinxhref{mailto:jfralph@liverpool.ac.uk}{jfralph@liverpool.ac.uk}
\item[] Michael Wright \sphinxhyphen{} \sphinxhref{mailto:mjwright@liverpool.ac.uk}{mjwright@liverpool.ac.uk}
\end{DUlineblock}

\sphinxlineitem{version}
\sphinxAtStartPar
1.1.0 \sphinxhyphen{} Includes fixes and improvements.

\end{description}\end{quote}
\index{get\_positions() (in module qnav.waypoints.generation)@\spxentry{get\_positions()}\spxextra{in module qnav.waypoints.generation}}

\begin{fulllineitems}
\phantomsection\label{\detokenize{modules/waypoints/generation:qnav.waypoints.generation.get_positions}}
\pysigstartsignatures
\pysiglinewithargsret{\sphinxcode{\sphinxupquote{qnav.waypoints.generation.}}\sphinxbfcode{\sphinxupquote{get\_positions}}}{\sphinxparam{\DUrole{n}{lat\_points}\DUrole{p}{:}\DUrole{w}{ }\DUrole{n}{ndarray}}\sphinxparamcomma \sphinxparam{\DUrole{n}{lon\_points}\DUrole{p}{:}\DUrole{w}{ }\DUrole{n}{ndarray}}\sphinxparamcomma \sphinxparam{\DUrole{n}{alt\_points}\DUrole{p}{:}\DUrole{w}{ }\DUrole{n}{ndarray}}\sphinxparamcomma \sphinxparam{\DUrole{n}{vehicle\_model}\DUrole{p}{:}\DUrole{w}{ }\DUrole{n}{{\hyperref[\detokenize{modules/waypoints/vehicle:qnav.waypoints.vehicle.Vehicle}]{\sphinxcrossref{Vehicle}}}}}\sphinxparamcomma \sphinxparam{\DUrole{n}{frequency}\DUrole{p}{:}\DUrole{w}{ }\DUrole{n}{float}\DUrole{w}{ }\DUrole{o}{=}\DUrole{w}{ }\DUrole{default_value}{200.0}}\sphinxparamcomma \sphinxparam{\DUrole{n}{avg\_speed}\DUrole{p}{:}\DUrole{w}{ }\DUrole{n}{float}\DUrole{w}{ }\DUrole{o}{=}\DUrole{w}{ }\DUrole{default_value}{100.0}}\sphinxparamcomma \sphinxparam{\DUrole{n}{travel\_times}\DUrole{p}{:}\DUrole{w}{ }\DUrole{n}{ndarray}\DUrole{w}{ }\DUrole{o}{=}\DUrole{w}{ }\DUrole{default_value}{None}}}{}
\pysigstopsignatures
\sphinxAtStartPar
This function interpolates positions (latitude, longitude and altitude)
between provided waypoints that are separated by either stated distance or
time (not both). The generated route is subject to the limitations of the
vehicle. The produced route can then be used for calculating the
required trajectory for simulations.
\begin{quote}\begin{description}
\sphinxlineitem{Parameters}\begin{itemize}
\item {} 
\sphinxAtStartPar
\sphinxstyleliteralstrong{\sphinxupquote{lat\_points}} (\sphinxstyleliteralemphasis{\sphinxupquote{numpy.ndarray}}) \textendash{} The decimal latitude points for the given waypoints.
Use negative values for southern coordinates (in degrees).

\item {} 
\sphinxAtStartPar
\sphinxstyleliteralstrong{\sphinxupquote{lon\_points}} (\sphinxstyleliteralemphasis{\sphinxupquote{numpy.ndarray}}) \textendash{} The decimal longitude points for the given waypoints.
Use negative values for western coordinates (in degrees).

\item {} 
\sphinxAtStartPar
\sphinxstyleliteralstrong{\sphinxupquote{alt\_points}} (\sphinxstyleliteralemphasis{\sphinxupquote{numpy.ndarray}}) \textendash{} The altitude for the given waypoints (in metres).

\item {} 
\sphinxAtStartPar
\sphinxstyleliteralstrong{\sphinxupquote{vehicle\_model}} ({\hyperref[\detokenize{modules/waypoints/vehicle:qnav.waypoints.vehicle.Vehicle}]{\sphinxcrossref{\sphinxstyleliteralemphasis{\sphinxupquote{Vehicle}}}}}) \textendash{} An initialised vehicle model containing the
manoeuvring limitations to consider when generating the waypoint
positions.

\item {} 
\sphinxAtStartPar
\sphinxstyleliteralstrong{\sphinxupquote{frequency}} (\sphinxstyleliteralemphasis{\sphinxupquote{float}}) \textendash{} The simulation’s update frequency. This is used for
calculating the number of intermediate points and total number of
waypoints to return.

\item {} 
\sphinxAtStartPar
\sphinxstyleliteralstrong{\sphinxupquote{avg\_speed}} (\sphinxstyleliteralemphasis{\sphinxupquote{float}}) \textendash{} The average vehicle speed for the given waypoints. This
is used to calculate the travel times between waypoints. This value is
replaced by ‘travel\_times’. Therefore, a value for this is not
required if the travel times are known.

\item {} 
\sphinxAtStartPar
\sphinxstyleliteralstrong{\sphinxupquote{travel\_times}} (\sphinxstyleliteralemphasis{\sphinxupquote{numpy.ndarray}}) \textendash{} The time (in seconds) when each waypoint is reached.
This is required for interpolating the position between points.

\end{itemize}

\sphinxlineitem{Returns}
\sphinxAtStartPar
An array containing column for latitude, longitude and altitude
of the vehicle at each time step.

\sphinxlineitem{Return type}
\sphinxAtStartPar
numpy.ndarray (n\sphinxhyphen{}by\sphinxhyphen{}4 elements)

\end{description}\end{quote}

\end{fulllineitems}

\index{interpolate\_positions() (in module qnav.waypoints.generation)@\spxentry{interpolate\_positions()}\spxextra{in module qnav.waypoints.generation}}

\begin{fulllineitems}
\phantomsection\label{\detokenize{modules/waypoints/generation:qnav.waypoints.generation.interpolate_positions}}
\pysigstartsignatures
\pysiglinewithargsret{\sphinxcode{\sphinxupquote{qnav.waypoints.generation.}}\sphinxbfcode{\sphinxupquote{interpolate\_positions}}}{\sphinxparam{\DUrole{n}{base\_points}\DUrole{p}{:}\DUrole{w}{ }\DUrole{n}{ndarray}}\sphinxparamcomma \sphinxparam{\DUrole{n}{frequency}\DUrole{p}{:}\DUrole{w}{ }\DUrole{n}{float}}}{{ $\rightarrow$ ndarray}}
\pysigstopsignatures
\sphinxAtStartPar
Interpolates the base waypoints at the new higher frequency.
This function can be used to speed up and stabilise the calculation of
the waypoint positions, by interpolating the intermediate steps.
\begin{quote}\begin{description}
\sphinxlineitem{Parameters}\begin{itemize}
\item {} 
\sphinxAtStartPar
\sphinxstyleliteralstrong{\sphinxupquote{base\_points}} (\sphinxstyleliteralemphasis{\sphinxupquote{np.ndarray}}) \textendash{} The base waypoint data to interpolate.

\item {} 
\sphinxAtStartPar
\sphinxstyleliteralstrong{\sphinxupquote{frequency}} (\sphinxstyleliteralemphasis{\sphinxupquote{float}}) \textendash{} The frequency to interpolate the waypoint data by.

\end{itemize}

\sphinxlineitem{Returns}
\sphinxAtStartPar
The interpolated version of the given waypoint positions.

\sphinxlineitem{Return type}
\sphinxAtStartPar
np.ndarray

\end{description}\end{quote}

\end{fulllineitems}

\index{produce\_noise() (in module qnav.waypoints.generation)@\spxentry{produce\_noise()}\spxextra{in module qnav.waypoints.generation}}

\begin{fulllineitems}
\phantomsection\label{\detokenize{modules/waypoints/generation:qnav.waypoints.generation.produce_noise}}
\pysigstartsignatures
\pysiglinewithargsret{\sphinxcode{\sphinxupquote{qnav.waypoints.generation.}}\sphinxbfcode{\sphinxupquote{produce\_noise}}}{\sphinxparam{\DUrole{n}{num\_rows}\DUrole{p}{:}\DUrole{w}{ }\DUrole{n}{int}}\sphinxparamcomma \sphinxparam{\DUrole{n}{first\_val}\DUrole{p}{:}\DUrole{w}{ }\DUrole{n}{float}}\sphinxparamcomma \sphinxparam{\DUrole{n}{noise\_val}\DUrole{p}{:}\DUrole{w}{ }\DUrole{n}{float}}\sphinxparamcomma \sphinxparam{\DUrole{n}{dt\_vdp}\DUrole{p}{:}\DUrole{w}{ }\DUrole{n}{float}}}{}
\pysigstopsignatures
\sphinxAtStartPar
This function generated random noise that can be added to various aspects.
For instance, this function can be used to produce noise for position,
attitude and oscillation displacement. This function is required so that
the noise generated is done so in a self\sphinxhyphen{}balancing manner such that the
displacements do not cause permanent or increasing drift.
\begin{quote}\begin{description}
\sphinxlineitem{Parameters}\begin{itemize}
\item {} 
\sphinxAtStartPar
\sphinxstyleliteralstrong{\sphinxupquote{num\_rows}} (\sphinxstyleliteralemphasis{\sphinxupquote{int}}) \textendash{} The number of noise values to produce.

\item {} 
\sphinxAtStartPar
\sphinxstyleliteralstrong{\sphinxupquote{first\_val}} (\sphinxstyleliteralemphasis{\sphinxupquote{float}}\sphinxstyleliteralemphasis{\sphinxupquote{ or }}\sphinxstyleliteralemphasis{\sphinxupquote{numpy.array}}) \textendash{} The first generated noise value(s).

\item {} 
\sphinxAtStartPar
\sphinxstyleliteralstrong{\sphinxupquote{noise\_val}} (\sphinxstyleliteralemphasis{\sphinxupquote{float}}) \textendash{} The scale of the noise values.

\item {} 
\sphinxAtStartPar
\sphinxstyleliteralstrong{\sphinxupquote{dt\_vdp}} \textendash{} The dampening factor to apply. This is used to prevent the
noise from causing permanent or excessive drift.

\end{itemize}

\sphinxlineitem{Returns}
\sphinxAtStartPar
The generated noise subject to the given parameters.

\sphinxlineitem{Return type}
\sphinxAtStartPar
numpy.array

\end{description}\end{quote}

\end{fulllineitems}

\index{moving\_average() (in module qnav.waypoints.generation)@\spxentry{moving\_average()}\spxextra{in module qnav.waypoints.generation}}

\begin{fulllineitems}
\phantomsection\label{\detokenize{modules/waypoints/generation:qnav.waypoints.generation.moving_average}}
\pysigstartsignatures
\pysiglinewithargsret{\sphinxcode{\sphinxupquote{qnav.waypoints.generation.}}\sphinxbfcode{\sphinxupquote{moving\_average}}}{\sphinxparam{\DUrole{n}{input\_array}}\sphinxparamcomma \sphinxparam{\DUrole{n}{window\_size}\DUrole{p}{:}\DUrole{w}{ }\DUrole{n}{int}\DUrole{w}{ }\DUrole{o}{=}\DUrole{w}{ }\DUrole{default_value}{5}}}{}
\pysigstopsignatures
\sphinxAtStartPar
Obtains the “sliding window average” for values of an array.
This function calculates the moving average of a given 1D array using the
same approach as MATLAB. A sliding window is passed over the given array.
For the centre (right) element the moving average is calculated using the
values covered by the window. This means that fewer elements are covered
at the start and end of the array. The returned values are the same size
as the given array.
\begin{quote}\begin{description}
\sphinxlineitem{Parameters}\begin{itemize}
\item {} 
\sphinxAtStartPar
\sphinxstyleliteralstrong{\sphinxupquote{input\_array}} (\sphinxstyleliteralemphasis{\sphinxupquote{NumPy array}}) \textendash{} The values which the moving average is going to be
calculated for. This must be a 1D array.

\item {} 
\sphinxAtStartPar
\sphinxstyleliteralstrong{\sphinxupquote{window\_size}} (\sphinxstyleliteralemphasis{\sphinxupquote{int}}) \textendash{} The size of the sliding window. Must be an integer.

\end{itemize}

\sphinxlineitem{Returns}
\sphinxAtStartPar
The calculated moving averages for the given array.

\sphinxlineitem{Return type}
\sphinxAtStartPar
NumPy array

\end{description}\end{quote}

\end{fulllineitems}

\index{rescale\_values() (in module qnav.waypoints.generation)@\spxentry{rescale\_values()}\spxextra{in module qnav.waypoints.generation}}

\begin{fulllineitems}
\phantomsection\label{\detokenize{modules/waypoints/generation:qnav.waypoints.generation.rescale_values}}
\pysigstartsignatures
\pysiglinewithargsret{\sphinxcode{\sphinxupquote{qnav.waypoints.generation.}}\sphinxbfcode{\sphinxupquote{rescale\_values}}}{\sphinxparam{\DUrole{n}{values}}}{}
\pysigstopsignatures
\sphinxAtStartPar
A simple function that allows angles to be re\sphinxhyphen{}scaled to prevent inaccuracy
calculations between angles differences.
\begin{quote}\begin{description}
\sphinxlineitem{Parameters}
\sphinxAtStartPar
\sphinxstyleliteralstrong{\sphinxupquote{values}} (\sphinxstyleliteralemphasis{\sphinxupquote{numpy.array}}) \textendash{} The angle values to be rescaled.

\sphinxlineitem{Returns}
\sphinxAtStartPar
The given values re\sphinxhyphen{}scaled between \sphinxhyphen{}pi/2 and pi/2.

\sphinxlineitem{Return type}
\sphinxAtStartPar
numpy.array

\end{description}\end{quote}

\end{fulllineitems}

\index{get\_simple\_trajectory\_values() (in module qnav.waypoints.generation)@\spxentry{get\_simple\_trajectory\_values()}\spxextra{in module qnav.waypoints.generation}}

\begin{fulllineitems}
\phantomsection\label{\detokenize{modules/waypoints/generation:qnav.waypoints.generation.get_simple_trajectory_values}}
\pysigstartsignatures
\pysiglinewithargsret{\sphinxcode{\sphinxupquote{qnav.waypoints.generation.}}\sphinxbfcode{\sphinxupquote{get\_simple\_trajectory\_values}}}{\sphinxparam{\DUrole{n}{waypoints}\DUrole{p}{:}\DUrole{w}{ }\DUrole{n}{ndarray}}\sphinxparamcomma \sphinxparam{\DUrole{n}{gravity\_model}\DUrole{p}{:}\DUrole{w}{ }\DUrole{n}{{\hyperref[\detokenize{modules/gravity/base:qnav.gravity.base.GravityModel}]{\sphinxcrossref{GravityModel}}}}}\sphinxparamcomma \sphinxparam{\DUrole{n}{given\_attitudes}\DUrole{p}{:}\DUrole{w}{ }\DUrole{n}{ndarray}\DUrole{w}{ }\DUrole{o}{=}\DUrole{w}{ }\DUrole{default_value}{None}}\sphinxparamcomma \sphinxparam{\DUrole{n}{attitude\_times}\DUrole{p}{:}\DUrole{w}{ }\DUrole{n}{ndarray}\DUrole{w}{ }\DUrole{o}{=}\DUrole{w}{ }\DUrole{default_value}{None}}\sphinxparamcomma \sphinxparam{\DUrole{n}{smoothing\_window\_time}\DUrole{p}{:}\DUrole{w}{ }\DUrole{n}{float}\DUrole{w}{ }\DUrole{o}{=}\DUrole{w}{ }\DUrole{default_value}{None}}\sphinxparamcomma \sphinxparam{\DUrole{n}{smoothing\_passes}\DUrole{p}{:}\DUrole{w}{ }\DUrole{n}{int}\DUrole{w}{ }\DUrole{o}{=}\DUrole{w}{ }\DUrole{default_value}{None}}\sphinxparamcomma \sphinxparam{\DUrole{n}{step\_size}\DUrole{p}{:}\DUrole{w}{ }\DUrole{n}{int}\DUrole{w}{ }\DUrole{o}{=}\DUrole{w}{ }\DUrole{default_value}{2}}\sphinxparamcomma \sphinxparam{\DUrole{n}{virtual\_mem\_file}\DUrole{p}{:}\DUrole{w}{ }\DUrole{n}{Path}\DUrole{w}{ }\DUrole{o}{=}\DUrole{w}{ }\DUrole{default_value}{None}}}{}
\pysigstopsignatures
\sphinxAtStartPar
Calculates the trajectory information and properties from given waypoints.
This function calculates the trajectory properties for each position for
a given set of waypoints. Only simple properties are generated, without
any added noise, vibrations or oscillations of the platform. In previous
releases this function was refereed to as legacy waypoint generation.
\begin{quote}\begin{description}
\sphinxlineitem{Parameters}\begin{itemize}
\item {} 
\sphinxAtStartPar
\sphinxstyleliteralstrong{\sphinxupquote{waypoints}} (\sphinxstyleliteralemphasis{\sphinxupquote{NumPy Matrix}}\sphinxstyleliteralemphasis{\sphinxupquote{ (}}\sphinxstyleliteralemphasis{\sphinxupquote{n\sphinxhyphen{}by\sphinxhyphen{}4 elements}}\sphinxstyleliteralemphasis{\sphinxupquote{)}}) \textendash{} The interpolated waypoint array previously generated. It
is important any changes/updates to the waypoints (such as altitude
changes) are applied to the waypoints before using this function.

\item {} 
\sphinxAtStartPar
\sphinxstyleliteralstrong{\sphinxupquote{gravity\_model}} ({\hyperref[\detokenize{modules/gravity/base:qnav.gravity.base.GravityModel}]{\sphinxcrossref{\sphinxstyleliteralemphasis{\sphinxupquote{GravityModel}}}}}) \textendash{} The gravity model used to calculate local gravity.
This is required for calculating the downwards gravitation
acceleration and its impact on the platform for each axis.

\item {} 
\sphinxAtStartPar
\sphinxstyleliteralstrong{\sphinxupquote{given\_attitudes}} (\sphinxstyleliteralemphasis{\sphinxupquote{numpy.array}}\sphinxstyleliteralemphasis{\sphinxupquote{ (}}\sphinxstyleliteralemphasis{\sphinxupquote{n\sphinxhyphen{}elements}}\sphinxstyleliteralemphasis{\sphinxupquote{)}}) \textendash{} (Optional) The requested attitude of the vehicle at
given times. Intermediate attitude values will be interpolated.

\item {} 
\sphinxAtStartPar
\sphinxstyleliteralstrong{\sphinxupquote{attitude\_times}} (\sphinxstyleliteralemphasis{\sphinxupquote{numpy.array}}\sphinxstyleliteralemphasis{\sphinxupquote{ (}}\sphinxstyleliteralemphasis{\sphinxupquote{n\sphinxhyphen{}elements}}\sphinxstyleliteralemphasis{\sphinxupquote{)}}) \textendash{} (Optional) The requested corresponding times for
the given attitudes. These timestamps are used to interpolate
intermediate attitude values.

\item {} 
\sphinxAtStartPar
\sphinxstyleliteralstrong{\sphinxupquote{smoothing\_window\_time}} (\sphinxstyleliteralemphasis{\sphinxupquote{float}}) \textendash{} (Optional) The size of the smoothing
(moving average) window to use for smoothing the calculated
waypoint velocities. This is useful for removing noise.

\item {} 
\sphinxAtStartPar
\sphinxstyleliteralstrong{\sphinxupquote{smoothing\_passes}} (\sphinxstyleliteralemphasis{\sphinxupquote{int}}) \textendash{} (Optional) The number of smoothing passes
(iterations) to perform on the calculated waypoint velocities.

\item {} 
\sphinxAtStartPar
\sphinxstyleliteralstrong{\sphinxupquote{step\_size}} (\sphinxstyleliteralemphasis{\sphinxupquote{int}}) \textendash{} (Optional) The step size used for calculating the
velocity between upcoming waypoints (default=1).

\item {} 
\sphinxAtStartPar
\sphinxstyleliteralstrong{\sphinxupquote{virtual\_mem\_file}} (\sphinxstyleliteralemphasis{\sphinxupquote{bool}}) \textendash{} The on\sphinxhyphen{}disk file to use for virtual memory.
Will be created or overwritten

\end{itemize}

\sphinxlineitem{Returns}
\sphinxAtStartPar
A matrix containing the position, velocity, acceleration,
attitude, angle rates and gravitational acceleration (and gradient)
for each time step.

\sphinxlineitem{Return type}
\sphinxAtStartPar
NumPy Matrix (n\sphinxhyphen{}by\sphinxhyphen{}16 elements)

\end{description}\end{quote}

\end{fulllineitems}

\index{get\_trajectory\_values() (in module qnav.waypoints.generation)@\spxentry{get\_trajectory\_values()}\spxextra{in module qnav.waypoints.generation}}

\begin{fulllineitems}
\phantomsection\label{\detokenize{modules/waypoints/generation:qnav.waypoints.generation.get_trajectory_values}}
\pysigstartsignatures
\pysiglinewithargsret{\sphinxcode{\sphinxupquote{qnav.waypoints.generation.}}\sphinxbfcode{\sphinxupquote{get\_trajectory\_values}}}{\sphinxparam{\DUrole{n}{waypoints}\DUrole{p}{:}\DUrole{w}{ }\DUrole{n}{ndarray}}\sphinxparamcomma \sphinxparam{\DUrole{n}{gravity\_model}\DUrole{p}{:}\DUrole{w}{ }\DUrole{n}{{\hyperref[\detokenize{modules/gravity/base:qnav.gravity.base.GravityModel}]{\sphinxcrossref{GravityModel}}}}}\sphinxparamcomma \sphinxparam{\DUrole{n}{vehicle\_model}\DUrole{p}{:}\DUrole{w}{ }\DUrole{n}{{\hyperref[\detokenize{modules/waypoints/vehicle:qnav.waypoints.vehicle.Vehicle}]{\sphinxcrossref{Vehicle}}}}}\sphinxparamcomma \sphinxparam{\DUrole{n}{given\_attitudes}\DUrole{p}{:}\DUrole{w}{ }\DUrole{n}{ndarray}\DUrole{w}{ }\DUrole{o}{=}\DUrole{w}{ }\DUrole{default_value}{None}}\sphinxparamcomma \sphinxparam{\DUrole{n}{attitude\_times}\DUrole{p}{:}\DUrole{w}{ }\DUrole{n}{ndarray}\DUrole{w}{ }\DUrole{o}{=}\DUrole{w}{ }\DUrole{default_value}{None}}\sphinxparamcomma \sphinxparam{\DUrole{n}{smoothing\_window\_time}\DUrole{p}{:}\DUrole{w}{ }\DUrole{n}{float}\DUrole{w}{ }\DUrole{o}{=}\DUrole{w}{ }\DUrole{default_value}{None}}\sphinxparamcomma \sphinxparam{\DUrole{n}{smoothing\_passes}\DUrole{p}{:}\DUrole{w}{ }\DUrole{n}{int}\DUrole{w}{ }\DUrole{o}{=}\DUrole{w}{ }\DUrole{default_value}{None}}\sphinxparamcomma \sphinxparam{\DUrole{n}{step\_size}\DUrole{p}{:}\DUrole{w}{ }\DUrole{n}{int}\DUrole{w}{ }\DUrole{o}{=}\DUrole{w}{ }\DUrole{default_value}{2}}\sphinxparamcomma \sphinxparam{\DUrole{n}{virtual\_mem\_file}\DUrole{p}{:}\DUrole{w}{ }\DUrole{n}{Path}\DUrole{w}{ }\DUrole{o}{=}\DUrole{w}{ }\DUrole{default_value}{None}}}{}
\pysigstopsignatures
\sphinxAtStartPar
Using waypoints route information generated prior, this function will
calculate the velocity, acceleration, attitude, angle rates and
gravitational acceleration (and gradient) for each time step. This
function calculates said values in an efficient vectorised manner
while ensuring they abide by the given vehicle model’s limitations.
\begin{quote}\begin{description}
\sphinxlineitem{Parameters}\begin{itemize}
\item {} 
\sphinxAtStartPar
\sphinxstyleliteralstrong{\sphinxupquote{waypoints}} (\sphinxstyleliteralemphasis{\sphinxupquote{NumPy Matrix}}\sphinxstyleliteralemphasis{\sphinxupquote{ (}}\sphinxstyleliteralemphasis{\sphinxupquote{n\sphinxhyphen{}by\sphinxhyphen{}4 elements}}\sphinxstyleliteralemphasis{\sphinxupquote{)}}) \textendash{} The interpolated waypoint array previously generated. It
is important any changes/updates to the waypoints (such as altitude
changes) are applied to the waypoints before using this function.

\item {} 
\sphinxAtStartPar
\sphinxstyleliteralstrong{\sphinxupquote{gravity\_model}} ({\hyperref[\detokenize{modules/gravity/base:qnav.gravity.base.GravityModel}]{\sphinxcrossref{\sphinxstyleliteralemphasis{\sphinxupquote{GravityModel}}}}}) \textendash{} The gravity model used to calculate local gravity.
This is required for calculating the downwards gravitation
acceleration and its impact on the platform for each axis.

\item {} 
\sphinxAtStartPar
\sphinxstyleliteralstrong{\sphinxupquote{vehicle\_model}} ({\hyperref[\detokenize{modules/waypoints/vehicle:qnav.waypoints.vehicle.Vehicle}]{\sphinxcrossref{\sphinxstyleliteralemphasis{\sphinxupquote{Vehicle}}}}}) \textendash{} An initialised vehicle model containing the
manoeuvring limitations to consider when generating the waypoint
positions.

\item {} 
\sphinxAtStartPar
\sphinxstyleliteralstrong{\sphinxupquote{given\_attitudes}} (\sphinxstyleliteralemphasis{\sphinxupquote{numpy.array}}\sphinxstyleliteralemphasis{\sphinxupquote{ (}}\sphinxstyleliteralemphasis{\sphinxupquote{n\sphinxhyphen{}elements}}\sphinxstyleliteralemphasis{\sphinxupquote{)}}) \textendash{} (Optional) The requested attitude of the vehicle at
given times. Intermediate attitude values will be interpolated.

\item {} 
\sphinxAtStartPar
\sphinxstyleliteralstrong{\sphinxupquote{attitude\_times}} (\sphinxstyleliteralemphasis{\sphinxupquote{numpy.array}}\sphinxstyleliteralemphasis{\sphinxupquote{ (}}\sphinxstyleliteralemphasis{\sphinxupquote{n\sphinxhyphen{}elements}}\sphinxstyleliteralemphasis{\sphinxupquote{)}}) \textendash{} (Optional) The requested corresponding times for
the given attitudes. These timestamps are used to interpolate
intermediate attitude values.

\item {} 
\sphinxAtStartPar
\sphinxstyleliteralstrong{\sphinxupquote{smoothing\_window\_time}} (\sphinxstyleliteralemphasis{\sphinxupquote{float}}) \textendash{} (Optional) The size of the smoothing
(moving average) window to use for smoothing the calculated
waypoint velocities. This is useful for removing noise.

\item {} 
\sphinxAtStartPar
\sphinxstyleliteralstrong{\sphinxupquote{smoothing\_passes}} (\sphinxstyleliteralemphasis{\sphinxupquote{int}}) \textendash{} (Optional) The number of smoothing passes
(iterations) to perform on the calculated waypoint velocities.

\item {} 
\sphinxAtStartPar
\sphinxstyleliteralstrong{\sphinxupquote{step\_size}} (\sphinxstyleliteralemphasis{\sphinxupquote{int}}) \textendash{} (Optional) The step size used for calculating the
velocity between upcoming waypoints (default=1).

\item {} 
\sphinxAtStartPar
\sphinxstyleliteralstrong{\sphinxupquote{virtual\_mem\_file}} (\sphinxstyleliteralemphasis{\sphinxupquote{bool}}) \textendash{} The on\sphinxhyphen{}disk file to use for virtual memory.
Will be created or overwritten

\end{itemize}

\sphinxlineitem{Returns}
\sphinxAtStartPar
A matrix containing the position, velocity, acceleration,
attitude, angle rates and gravitational acceleration (and gradient)
for each time step.

\sphinxlineitem{Return type}
\sphinxAtStartPar
NumPy Matrix (n\sphinxhyphen{}by\sphinxhyphen{}16 elements)

\end{description}\end{quote}

\end{fulllineitems}

\index{plot\_waypoints() (in module qnav.waypoints.generation)@\spxentry{plot\_waypoints()}\spxextra{in module qnav.waypoints.generation}}

\begin{fulllineitems}
\phantomsection\label{\detokenize{modules/waypoints/generation:qnav.waypoints.generation.plot_waypoints}}
\pysigstartsignatures
\pysiglinewithargsret{\sphinxcode{\sphinxupquote{qnav.waypoints.generation.}}\sphinxbfcode{\sphinxupquote{plot\_waypoints}}}{\sphinxparam{\DUrole{n}{data\_values}\DUrole{p}{:}\DUrole{w}{ }\DUrole{n}{ndarray}}\sphinxparamcomma \sphinxparam{\DUrole{n}{filename}\DUrole{o}{=}\DUrole{default_value}{None}}}{}
\pysigstopsignatures
\sphinxAtStartPar
A simple function to plot the generated waypoints.
\begin{quote}\begin{description}
\sphinxlineitem{Parameters}\begin{itemize}
\item {} 
\sphinxAtStartPar
\sphinxstyleliteralstrong{\sphinxupquote{data\_values}} (\sphinxstyleliteralemphasis{\sphinxupquote{numpy.ndarray}}\sphinxstyleliteralemphasis{\sphinxupquote{ (}}\sphinxstyleliteralemphasis{\sphinxupquote{n\sphinxhyphen{}by\sphinxhyphen{}3 elements}}\sphinxstyleliteralemphasis{\sphinxupquote{)}}) \textendash{} The array of values returned by the generator
function.

\item {} 
\sphinxAtStartPar
\sphinxstyleliteralstrong{\sphinxupquote{filename}} (\sphinxstyleliteralemphasis{\sphinxupquote{str}}) \textendash{} (Optional) Saves the produced plot under the given file.

\end{itemize}

\end{description}\end{quote}

\end{fulllineitems}

\index{plot\_figures() (in module qnav.waypoints.generation)@\spxentry{plot\_figures()}\spxextra{in module qnav.waypoints.generation}}

\begin{fulllineitems}
\phantomsection\label{\detokenize{modules/waypoints/generation:qnav.waypoints.generation.plot_figures}}
\pysigstartsignatures
\pysiglinewithargsret{\sphinxcode{\sphinxupquote{qnav.waypoints.generation.}}\sphinxbfcode{\sphinxupquote{plot\_figures}}}{\sphinxparam{\DUrole{n}{data\_values}\DUrole{p}{:}\DUrole{w}{ }\DUrole{n}{ndarray}}\sphinxparamcomma \sphinxparam{\DUrole{n}{filename}\DUrole{o}{=}\DUrole{default_value}{None}}}{}
\pysigstopsignatures
\sphinxAtStartPar
A simple function to plot the change in variables during the waypoint
generation. Each of these is presented within a sub\sphinxhyphen{}figure.
\begin{quote}\begin{description}
\sphinxlineitem{Parameters}\begin{itemize}
\item {} 
\sphinxAtStartPar
\sphinxstyleliteralstrong{\sphinxupquote{data\_values}} (\sphinxstyleliteralemphasis{\sphinxupquote{numpy.ndarray}}\sphinxstyleliteralemphasis{\sphinxupquote{ (}}\sphinxstyleliteralemphasis{\sphinxupquote{n\sphinxhyphen{}by\sphinxhyphen{}16 elements}}\sphinxstyleliteralemphasis{\sphinxupquote{)}}) \textendash{} The array of values returned by the generator
function.

\item {} 
\sphinxAtStartPar
\sphinxstyleliteralstrong{\sphinxupquote{filename}} (\sphinxstyleliteralemphasis{\sphinxupquote{str}}) \textendash{} (Optional) Saves the produced plot under the given file.

\end{itemize}

\end{description}\end{quote}

\end{fulllineitems}

\index{load\_waypoint\_csv() (in module qnav.waypoints.generation)@\spxentry{load\_waypoint\_csv()}\spxextra{in module qnav.waypoints.generation}}

\begin{fulllineitems}
\phantomsection\label{\detokenize{modules/waypoints/generation:qnav.waypoints.generation.load_waypoint_csv}}
\pysigstartsignatures
\pysiglinewithargsret{\sphinxcode{\sphinxupquote{qnav.waypoints.generation.}}\sphinxbfcode{\sphinxupquote{load\_waypoint\_csv}}}{\sphinxparam{\DUrole{n}{filepath}\DUrole{p}{:}\DUrole{w}{ }\DUrole{n}{Path}}\sphinxparamcomma \sphinxparam{\DUrole{n}{down\_sample\_rate}\DUrole{p}{:}\DUrole{w}{ }\DUrole{n}{int}\DUrole{w}{ }\DUrole{o}{=}\DUrole{w}{ }\DUrole{default_value}{1}}}{}
\pysigstopsignatures
\sphinxAtStartPar
This function will extract the values from a given trajectory file,
formatted as comma\sphinxhyphen{}separated\sphinxhyphen{}values. Once the relevant information is
extracted, 7 rows will be returned containing information of the extracted
timestamps, latitude points, longitude points, altitude points and three
attitude angles. These rows are returned within a list.
\begin{quote}\begin{description}
\sphinxlineitem{Parameters}\begin{itemize}
\item {} 
\sphinxAtStartPar
\sphinxstyleliteralstrong{\sphinxupquote{filepath}} (\sphinxstyleliteralemphasis{\sphinxupquote{str}}\sphinxstyleliteralemphasis{\sphinxupquote{ or }}\sphinxstyleliteralemphasis{\sphinxupquote{Path}}) \textendash{} The location of the CSV trajectory file to read.

\item {} 
\sphinxAtStartPar
\sphinxstyleliteralstrong{\sphinxupquote{down\_sample\_rate}} (\sphinxstyleliteralemphasis{\sphinxupquote{int}}) \textendash{} The down sampling multiplier to apply.
This can be used to skip rows for cases where CSV waypoints are
given too frequently or spaced too closely together.

\end{itemize}

\sphinxlineitem{Returns}
\sphinxAtStartPar
A list containing 7 arrays holding the extracted timestamps,
latitude points, longitude points, altitude points and attitude
angles from the given CSV trajectory file.

\sphinxlineitem{Return type}
\sphinxAtStartPar
NumPy array

\end{description}\end{quote}

\end{fulllineitems}

\index{crop\_waypoints() (in module qnav.waypoints.generation)@\spxentry{crop\_waypoints()}\spxextra{in module qnav.waypoints.generation}}

\begin{fulllineitems}
\phantomsection\label{\detokenize{modules/waypoints/generation:qnav.waypoints.generation.crop_waypoints}}
\pysigstartsignatures
\pysiglinewithargsret{\sphinxcode{\sphinxupquote{qnav.waypoints.generation.}}\sphinxbfcode{\sphinxupquote{crop\_waypoints}}}{\sphinxparam{\DUrole{n}{waypoints}\DUrole{p}{:}\DUrole{w}{ }\DUrole{n}{ndarray}}\sphinxparamcomma \sphinxparam{\DUrole{n}{max\_time}\DUrole{p}{:}\DUrole{w}{ }\DUrole{n}{float}}}{}
\pysigstopsignatures
\sphinxAtStartPar
Crops waypoint data between time ranges (in seconds).
\begin{quote}\begin{description}
\sphinxlineitem{Parameters}\begin{itemize}
\item {} 
\sphinxAtStartPar
\sphinxstyleliteralstrong{\sphinxupquote{waypoints}} (\sphinxstyleliteralemphasis{\sphinxupquote{numpy.ndarray}}) \textendash{} The waypoint data to crop.

\item {} 
\sphinxAtStartPar
\sphinxstyleliteralstrong{\sphinxupquote{max\_time}} (\sphinxstyleliteralemphasis{\sphinxupquote{float}}) \textendash{} The maximum time for the return waypoint data.

\end{itemize}

\sphinxlineitem{Returns}
\sphinxAtStartPar
The cropped waypoint data

\sphinxlineitem{Return type}
\sphinxAtStartPar
numpy.ndarray

\end{description}\end{quote}

\end{fulllineitems}

\index{save\_waypoints() (in module qnav.waypoints.generation)@\spxentry{save\_waypoints()}\spxextra{in module qnav.waypoints.generation}}

\begin{fulllineitems}
\phantomsection\label{\detokenize{modules/waypoints/generation:qnav.waypoints.generation.save_waypoints}}
\pysigstartsignatures
\pysiglinewithargsret{\sphinxcode{\sphinxupquote{qnav.waypoints.generation.}}\sphinxbfcode{\sphinxupquote{save\_waypoints}}}{\sphinxparam{\DUrole{n}{waypoints}\DUrole{p}{:}\DUrole{w}{ }\DUrole{n}{ndarray}}\sphinxparamcomma \sphinxparam{\DUrole{n}{file\_path}\DUrole{p}{:}\DUrole{w}{ }\DUrole{n}{Path}}\sphinxparamcomma \sphinxparam{\DUrole{n}{allow\_overwrite}\DUrole{p}{:}\DUrole{w}{ }\DUrole{n}{bool}\DUrole{w}{ }\DUrole{o}{=}\DUrole{w}{ }\DUrole{default_value}{False}}}{}
\pysigstopsignatures
\sphinxAtStartPar
Exports waypoint data to disk, allowing for reuse.
\begin{quote}\begin{description}
\sphinxlineitem{Parameters}\begin{itemize}
\item {} 
\sphinxAtStartPar
\sphinxstyleliteralstrong{\sphinxupquote{waypoints}} (\sphinxstyleliteralemphasis{\sphinxupquote{numpy.array}}) \textendash{} The waypoint data to save.

\item {} 
\sphinxAtStartPar
\sphinxstyleliteralstrong{\sphinxupquote{file\_path}} (\sphinxstyleliteralemphasis{\sphinxupquote{pathlib.Path}}) \textendash{} The file location to save the waypoint data to.

\item {} 
\sphinxAtStartPar
\sphinxstyleliteralstrong{\sphinxupquote{allow\_overwrite}} (\sphinxstyleliteralemphasis{\sphinxupquote{bool}}) \textendash{} Should file overwriting be permitted? This
safeguards against accidental overwriting (default=False).

\end{itemize}

\end{description}\end{quote}

\end{fulllineitems}

\index{load\_waypoints() (in module qnav.waypoints.generation)@\spxentry{load\_waypoints()}\spxextra{in module qnav.waypoints.generation}}

\begin{fulllineitems}
\phantomsection\label{\detokenize{modules/waypoints/generation:qnav.waypoints.generation.load_waypoints}}
\pysigstartsignatures
\pysiglinewithargsret{\sphinxcode{\sphinxupquote{qnav.waypoints.generation.}}\sphinxbfcode{\sphinxupquote{load\_waypoints}}}{\sphinxparam{\DUrole{n}{file\_path}\DUrole{p}{:}\DUrole{w}{ }\DUrole{n}{Path}}\sphinxparamcomma \sphinxparam{\DUrole{n}{as\_virtual}\DUrole{p}{:}\DUrole{w}{ }\DUrole{n}{bool}\DUrole{w}{ }\DUrole{o}{=}\DUrole{w}{ }\DUrole{default_value}{False}}}{}
\pysigstopsignatures
\sphinxAtStartPar
Imports waypoint data from disk, reusing saved values.
\begin{quote}\begin{description}
\sphinxlineitem{Parameters}\begin{itemize}
\item {} 
\sphinxAtStartPar
\sphinxstyleliteralstrong{\sphinxupquote{file\_path}} (\sphinxstyleliteralemphasis{\sphinxupquote{pathlib.Path}}) \textendash{} The file location to load the waypoint data from.

\item {} 
\sphinxAtStartPar
\sphinxstyleliteralstrong{\sphinxupquote{as\_virtual}} (\sphinxstyleliteralemphasis{\sphinxupquote{bool}}) \textendash{} If set to True prevent all the data being loaded at
once and use virtual memory. This reduces RAM usage, but impacts
runtime (default=False).

\end{itemize}

\sphinxlineitem{Returns}
\sphinxAtStartPar
The loaded waypoint data.

\sphinxlineitem{Return type}
\sphinxAtStartPar
numpy.array (n\sphinxhyphen{}by\sphinxhyphen{}16 elements)

\end{description}\end{quote}

\end{fulllineitems}

\index{get\_trajectory() (in module qnav.waypoints.generation)@\spxentry{get\_trajectory()}\spxextra{in module qnav.waypoints.generation}}

\begin{fulllineitems}
\phantomsection\label{\detokenize{modules/waypoints/generation:qnav.waypoints.generation.get_trajectory}}
\pysigstartsignatures
\pysiglinewithargsret{\sphinxcode{\sphinxupquote{qnav.waypoints.generation.}}\sphinxbfcode{\sphinxupquote{get\_trajectory}}}{\sphinxparam{\DUrole{n}{waypoints}\DUrole{p}{:}\DUrole{w}{ }\DUrole{n}{ndarray}}\sphinxparamcomma \sphinxparam{\DUrole{n}{gravity\_model}\DUrole{p}{:}\DUrole{w}{ }\DUrole{n}{{\hyperref[\detokenize{modules/gravity/base:qnav.gravity.base.GravityModel}]{\sphinxcrossref{GravityModel}}}}}}{}
\pysigstopsignatures
\sphinxAtStartPar
TODO: Add documentation
:param waypoints:
:param gravity\_model:
:return:

\end{fulllineitems}

\index{read\_csv\_waypoints() (in module qnav.waypoints.generation)@\spxentry{read\_csv\_waypoints()}\spxextra{in module qnav.waypoints.generation}}

\begin{fulllineitems}
\phantomsection\label{\detokenize{modules/waypoints/generation:qnav.waypoints.generation.read_csv_waypoints}}
\pysigstartsignatures
\pysiglinewithargsret{\sphinxcode{\sphinxupquote{qnav.waypoints.generation.}}\sphinxbfcode{\sphinxupquote{read\_csv\_waypoints}}}{\sphinxparam{\DUrole{n}{file\_path}\DUrole{p}{:}\DUrole{w}{ }\DUrole{n}{Path}}\sphinxparamcomma \sphinxparam{\DUrole{n}{delimiter}\DUrole{p}{:}\DUrole{w}{ }\DUrole{n}{str}\DUrole{w}{ }\DUrole{o}{=}\DUrole{w}{ }\DUrole{default_value}{\textquotesingle{},\textquotesingle{}}}}{{ $\rightarrow$ dict\DUrole{p}{{[}}str\DUrole{p}{,}\DUrole{w}{ }ndarray\DUrole{w}{ }\DUrole{p}{|}\DUrole{w}{ }None\DUrole{p}{{]}}}}
\pysigstopsignatures
\sphinxAtStartPar
Reads given CSV waypoint file and returns the records contained within it.
For a given file path, attempts to read its content and return it’s
extracted records inside a dictionary.
\begin{quote}\begin{description}
\sphinxlineitem{Parameters}\begin{itemize}
\item {} 
\sphinxAtStartPar
\sphinxstyleliteralstrong{\sphinxupquote{file\_path}} (\sphinxstyleliteralemphasis{\sphinxupquote{Path}}) \textendash{} The CSV file to attempt to read waypoint data from.

\item {} 
\sphinxAtStartPar
\sphinxstyleliteralstrong{\sphinxupquote{delimiter}} (\sphinxstyleliteralemphasis{\sphinxupquote{str}}) \textendash{} The delimiter used to represent column breaks.
By default, this is a comma, but other characters can be used.

\end{itemize}

\sphinxlineitem{Returns}
\sphinxAtStartPar
A dictionary holding the extracted data.

\sphinxlineitem{Return type}
\sphinxAtStartPar
dict{[}str, np.ndarray | None{]}

\end{description}\end{quote}

\end{fulllineitems}


\sphinxstepscope


\subsubsection{qnav.waypoints.trajectory}
\label{\detokenize{modules/waypoints/trajectory:module-qnav.waypoints.trajectory}}\label{\detokenize{modules/waypoints/trajectory:qnav-waypoints-trajectory}}\label{\detokenize{modules/waypoints/trajectory::doc}}\index{module@\spxentry{module}!qnav.waypoints.trajectory@\spxentry{qnav.waypoints.trajectory}}\index{qnav.waypoints.trajectory@\spxentry{qnav.waypoints.trajectory}!module@\spxentry{module}}

\paragraph{trajectory.py}
\label{\detokenize{modules/waypoints/trajectory:trajectory-py}}\begin{quote}\begin{description}
\sphinxlineitem{summary}
\sphinxAtStartPar
Used for the core representation of ground truth trajectory information.
This module is responsible for handling ground truth waypoint
trajectories, including the managing of input data, position and value
calculations and look\sphinxhyphen{}up interpolation of ground truth records.

\sphinxlineitem{author}
\sphinxAtStartPar
Michael Wright \sphinxhyphen{} \sphinxhref{mailto:mjwright@liverpool.ac.uk}{mjwright@liverpool.ac.uk}

\sphinxlineitem{version}
\sphinxAtStartPar
1.0.0 \sphinxhyphen{} First full implementation of this module.

\end{description}\end{quote}
\index{GroundTruth (class in qnav.waypoints.trajectory)@\spxentry{GroundTruth}\spxextra{class in qnav.waypoints.trajectory}}

\begin{fulllineitems}
\phantomsection\label{\detokenize{modules/waypoints/trajectory:qnav.waypoints.trajectory.GroundTruth}}
\pysigstartsignatures
\pysiglinewithargsret{\sphinxbfcode{\sphinxupquote{class\DUrole{w}{ }}}\sphinxcode{\sphinxupquote{qnav.waypoints.trajectory.}}\sphinxbfcode{\sphinxupquote{GroundTruth}}}{\sphinxparam{\DUrole{n}{timestamp}\DUrole{p}{:}\DUrole{w}{ }\DUrole{n}{float}}\sphinxparamcomma \sphinxparam{\DUrole{n}{position}\DUrole{p}{:}\DUrole{w}{ }\DUrole{n}{ndarray}}\sphinxparamcomma \sphinxparam{\DUrole{n}{velocity}\DUrole{p}{:}\DUrole{w}{ }\DUrole{n}{ndarray}}\sphinxparamcomma \sphinxparam{\DUrole{n}{acceleration}\DUrole{p}{:}\DUrole{w}{ }\DUrole{n}{ndarray}}\sphinxparamcomma \sphinxparam{\DUrole{n}{attitude}\DUrole{p}{:}\DUrole{w}{ }\DUrole{n}{ndarray}}\sphinxparamcomma \sphinxparam{\DUrole{n}{angle\_rates}\DUrole{p}{:}\DUrole{w}{ }\DUrole{n}{ndarray}}}{}
\pysigstopsignatures
\sphinxAtStartPar
A record class for trajectory ground truth a given time.
This class contains the ground truth (the true trajectory values) at a
given time. The purpose of this is to supply the “real\sphinxhyphen{}world” data to
sensor models, where simulated noise and errors will then be applied to
produce measurement data.
\index{timestamp (qnav.waypoints.trajectory.GroundTruth attribute)@\spxentry{timestamp}\spxextra{qnav.waypoints.trajectory.GroundTruth attribute}}

\begin{fulllineitems}
\phantomsection\label{\detokenize{modules/waypoints/trajectory:qnav.waypoints.trajectory.GroundTruth.timestamp}}
\pysigstartsignatures
\pysigline{\sphinxbfcode{\sphinxupquote{timestamp}}\sphinxbfcode{\sphinxupquote{\DUrole{p}{:}\DUrole{w}{ }float}}}
\pysigstopsignatures
\sphinxAtStartPar
The time since initialisation (in seconds).

\end{fulllineitems}

\index{position (qnav.waypoints.trajectory.GroundTruth attribute)@\spxentry{position}\spxextra{qnav.waypoints.trajectory.GroundTruth attribute}}

\begin{fulllineitems}
\phantomsection\label{\detokenize{modules/waypoints/trajectory:qnav.waypoints.trajectory.GroundTruth.position}}
\pysigstartsignatures
\pysigline{\sphinxbfcode{\sphinxupquote{position}}\sphinxbfcode{\sphinxupquote{\DUrole{p}{:}\DUrole{w}{ }ndarray}}}
\pysigstopsignatures
\sphinxAtStartPar
The position as latitude\sphinxhyphen{}longitude\sphinxhyphen{}altitude (in deg\sphinxhyphen{}deg\sphinxhyphen{}metres).

\end{fulllineitems}

\index{velocity (qnav.waypoints.trajectory.GroundTruth attribute)@\spxentry{velocity}\spxextra{qnav.waypoints.trajectory.GroundTruth attribute}}

\begin{fulllineitems}
\phantomsection\label{\detokenize{modules/waypoints/trajectory:qnav.waypoints.trajectory.GroundTruth.velocity}}
\pysigstartsignatures
\pysigline{\sphinxbfcode{\sphinxupquote{velocity}}\sphinxbfcode{\sphinxupquote{\DUrole{p}{:}\DUrole{w}{ }ndarray}}}
\pysigstopsignatures
\sphinxAtStartPar
The velocity as along\sphinxhyphen{}across\sphinxhyphen{}down in body axis (in metres/second).

\end{fulllineitems}

\index{acceleration (qnav.waypoints.trajectory.GroundTruth attribute)@\spxentry{acceleration}\spxextra{qnav.waypoints.trajectory.GroundTruth attribute}}

\begin{fulllineitems}
\phantomsection\label{\detokenize{modules/waypoints/trajectory:qnav.waypoints.trajectory.GroundTruth.acceleration}}
\pysigstartsignatures
\pysigline{\sphinxbfcode{\sphinxupquote{acceleration}}\sphinxbfcode{\sphinxupquote{\DUrole{p}{:}\DUrole{w}{ }ndarray}}}
\pysigstopsignatures
\sphinxAtStartPar
The acceleration as along\sphinxhyphen{}across\sphinxhyphen{}down in body axis (in metres/second\textasciicircum{}2).

\end{fulllineitems}

\index{attitude (qnav.waypoints.trajectory.GroundTruth attribute)@\spxentry{attitude}\spxextra{qnav.waypoints.trajectory.GroundTruth attribute}}

\begin{fulllineitems}
\phantomsection\label{\detokenize{modules/waypoints/trajectory:qnav.waypoints.trajectory.GroundTruth.attitude}}
\pysigstartsignatures
\pysigline{\sphinxbfcode{\sphinxupquote{attitude}}\sphinxbfcode{\sphinxupquote{\DUrole{p}{:}\DUrole{w}{ }ndarray}}}
\pysigstopsignatures
\sphinxAtStartPar
The attitude as north\sphinxhyphen{}east\sphinxhyphen{}down in NED axis (in degrees).

\end{fulllineitems}

\index{angle\_rates (qnav.waypoints.trajectory.GroundTruth attribute)@\spxentry{angle\_rates}\spxextra{qnav.waypoints.trajectory.GroundTruth attribute}}

\begin{fulllineitems}
\phantomsection\label{\detokenize{modules/waypoints/trajectory:qnav.waypoints.trajectory.GroundTruth.angle_rates}}
\pysigstartsignatures
\pysigline{\sphinxbfcode{\sphinxupquote{angle\_rates}}\sphinxbfcode{\sphinxupquote{\DUrole{p}{:}\DUrole{w}{ }ndarray}}}
\pysigstopsignatures
\sphinxAtStartPar
The angle rates as north\sphinxhyphen{}east\sphinxhyphen{}down in NED axis (in degrees/second).

\end{fulllineitems}

\index{as\_dict() (qnav.waypoints.trajectory.GroundTruth method)@\spxentry{as\_dict()}\spxextra{qnav.waypoints.trajectory.GroundTruth method}}

\begin{fulllineitems}
\phantomsection\label{\detokenize{modules/waypoints/trajectory:qnav.waypoints.trajectory.GroundTruth.as_dict}}
\pysigstartsignatures
\pysiglinewithargsret{\sphinxbfcode{\sphinxupquote{as\_dict}}}{}{{ $\rightarrow$ dict}}
\pysigstopsignatures
\sphinxAtStartPar
Returns class properties in dictionary form.
\begin{quote}\begin{description}
\sphinxlineitem{Returns}
\sphinxAtStartPar
A dictionary containing ground truth values.

\sphinxlineitem{Return type}
\sphinxAtStartPar
dict

\end{description}\end{quote}

\end{fulllineitems}

\index{as\_numpy() (qnav.waypoints.trajectory.GroundTruth method)@\spxentry{as\_numpy()}\spxextra{qnav.waypoints.trajectory.GroundTruth method}}

\begin{fulllineitems}
\phantomsection\label{\detokenize{modules/waypoints/trajectory:qnav.waypoints.trajectory.GroundTruth.as_numpy}}
\pysigstartsignatures
\pysiglinewithargsret{\sphinxbfcode{\sphinxupquote{as\_numpy}}}{}{{ $\rightarrow$ ndarray}}
\pysigstopsignatures
\sphinxAtStartPar
Returns call properties in 1D numpy array form.
\begin{quote}\begin{description}
\sphinxlineitem{Returns}
\sphinxAtStartPar
A numpy array containing ground truth values.

\sphinxlineitem{Return type}
\sphinxAtStartPar
np.ndarray

\end{description}\end{quote}

\end{fulllineitems}


\end{fulllineitems}

\index{Waypoints (class in qnav.waypoints.trajectory)@\spxentry{Waypoints}\spxextra{class in qnav.waypoints.trajectory}}

\begin{fulllineitems}
\phantomsection\label{\detokenize{modules/waypoints/trajectory:qnav.waypoints.trajectory.Waypoints}}
\pysigstartsignatures
\pysiglinewithargsret{\sphinxbfcode{\sphinxupquote{class\DUrole{w}{ }}}\sphinxcode{\sphinxupquote{qnav.waypoints.trajectory.}}\sphinxbfcode{\sphinxupquote{Waypoints}}}{\sphinxparam{\DUrole{n}{timestamps}\DUrole{p}{:}\DUrole{w}{ }\DUrole{n}{ndarray}}\sphinxparamcomma \sphinxparam{\DUrole{n}{lat\_points}\DUrole{p}{:}\DUrole{w}{ }\DUrole{n}{ndarray}}\sphinxparamcomma \sphinxparam{\DUrole{n}{lon\_points}\DUrole{p}{:}\DUrole{w}{ }\DUrole{n}{ndarray}}\sphinxparamcomma \sphinxparam{\DUrole{n}{alt\_points}\DUrole{p}{:}\DUrole{w}{ }\DUrole{n}{ndarray}}\sphinxparamcomma \sphinxparam{\DUrole{n}{attitudes}\DUrole{p}{:}\DUrole{w}{ }\DUrole{n}{ndarray}\DUrole{w}{ }\DUrole{o}{=}\DUrole{w}{ }\DUrole{default_value}{None}}}{}
\pysigstopsignatures
\sphinxAtStartPar
A container for trajectory waypoint data used to define a course.
This data class will contain the arrival time and positioning records
for core points along a trajectory. These can either represent sparse
base points that a trajectory is to cover, or a series of dense
positions, a fixed frequency, for an already generated trajectory.
Additionally, attitude values can be provided for each position.
\index{as\_array() (qnav.waypoints.trajectory.Waypoints method)@\spxentry{as\_array()}\spxextra{qnav.waypoints.trajectory.Waypoints method}}

\begin{fulllineitems}
\phantomsection\label{\detokenize{modules/waypoints/trajectory:qnav.waypoints.trajectory.Waypoints.as_array}}
\pysigstartsignatures
\pysiglinewithargsret{\sphinxbfcode{\sphinxupquote{as\_array}}}{}{{ $\rightarrow$ ndarray}}
\pysigstopsignatures
\sphinxAtStartPar
Returns required records formatted as column arrays.
Specifically returns an n\sphinxhyphen{}by\sphinxhyphen{}4 array holding the timestamps, latitude,
longitude and altitude records as columns, respectively. Optional
Attitude records are not included in this output.
\begin{quote}\begin{description}
\sphinxlineitem{Returns}
\sphinxAtStartPar
Position data formatted as column array.

\sphinxlineitem{Return type}
\sphinxAtStartPar
np.ndarray

\end{description}\end{quote}

\end{fulllineitems}


\end{fulllineitems}

\index{Trajectory (class in qnav.waypoints.trajectory)@\spxentry{Trajectory}\spxextra{class in qnav.waypoints.trajectory}}

\begin{fulllineitems}
\phantomsection\label{\detokenize{modules/waypoints/trajectory:qnav.waypoints.trajectory.Trajectory}}
\pysigstartsignatures
\pysiglinewithargsret{\sphinxbfcode{\sphinxupquote{class\DUrole{w}{ }}}\sphinxcode{\sphinxupquote{qnav.waypoints.trajectory.}}\sphinxbfcode{\sphinxupquote{Trajectory}}}{\sphinxparam{\DUrole{n}{waypoints}\DUrole{p}{:}\DUrole{w}{ }\DUrole{n}{{\hyperref[\detokenize{modules/waypoints/trajectory:qnav.waypoints.trajectory.Waypoints}]{\sphinxcrossref{Waypoints}}}}}\sphinxparamcomma \sphinxparam{\DUrole{n}{vehicle\_model}\DUrole{p}{:}\DUrole{w}{ }\DUrole{n}{{\hyperref[\detokenize{modules/waypoints/vehicle:qnav.waypoints.vehicle.Vehicle}]{\sphinxcrossref{Vehicle}}}}}\sphinxparamcomma \sphinxparam{\DUrole{n}{gravity\_model}\DUrole{p}{:}\DUrole{w}{ }\DUrole{n}{{\hyperref[\detokenize{modules/gravity/base:qnav.gravity.base.GravityModel}]{\sphinxcrossref{GravityModel}}}}}\sphinxparamcomma \sphinxparam{\DUrole{n}{base\_freq}\DUrole{p}{:}\DUrole{w}{ }\DUrole{n}{float}}\sphinxparamcomma \sphinxparam{\DUrole{n}{inter\_method}\DUrole{p}{:}\DUrole{w}{ }\DUrole{n}{str}\DUrole{w}{ }\DUrole{o}{=}\DUrole{w}{ }\DUrole{default_value}{\textquotesingle{}pchip\textquotesingle{}}}}{}
\pysigstopsignatures
\sphinxAtStartPar
A fully\sphinxhyphen{}generated trajectory with ground truth look\sphinxhyphen{}up.
On initialization, this generates the full trajectory from given base
waypoints, vehicle model and gravity model. The core trajectory is
generated with base points calculated at fixed frequency. Intermediate
values are then interpolated.
\index{get\_record() (qnav.waypoints.trajectory.Trajectory method)@\spxentry{get\_record()}\spxextra{qnav.waypoints.trajectory.Trajectory method}}

\begin{fulllineitems}
\phantomsection\label{\detokenize{modules/waypoints/trajectory:qnav.waypoints.trajectory.Trajectory.get_record}}
\pysigstartsignatures
\pysiglinewithargsret{\sphinxbfcode{\sphinxupquote{get\_record}}}{\sphinxparam{\DUrole{n}{time\_sec}\DUrole{p}{:}\DUrole{w}{ }\DUrole{n}{float}}}{{ $\rightarrow$ {\hyperref[\detokenize{modules/waypoints/trajectory:qnav.waypoints.trajectory.GroundTruth}]{\sphinxcrossref{GroundTruth}}}}}
\pysigstopsignatures
\sphinxAtStartPar
Returns values of the trajectory at given time (in seconds).
Safely interpolates and returns ground truth at requested time.
\begin{quote}\begin{description}
\sphinxlineitem{Parameters}
\sphinxAtStartPar
\sphinxstyleliteralstrong{\sphinxupquote{time\_sec}} (\sphinxstyleliteralemphasis{\sphinxupquote{float}}) \textendash{} A requested time in seconds.

\sphinxlineitem{Returns}
\sphinxAtStartPar
The corresponding ground truth record.

\sphinxlineitem{Return type}
\sphinxAtStartPar
{\hyperref[\detokenize{modules/waypoints/trajectory:qnav.waypoints.trajectory.GroundTruth}]{\sphinxcrossref{GroundTruth}}}

\end{description}\end{quote}

\end{fulllineitems}

\index{get\_records() (qnav.waypoints.trajectory.Trajectory method)@\spxentry{get\_records()}\spxextra{qnav.waypoints.trajectory.Trajectory method}}

\begin{fulllineitems}
\phantomsection\label{\detokenize{modules/waypoints/trajectory:qnav.waypoints.trajectory.Trajectory.get_records}}
\pysigstartsignatures
\pysiglinewithargsret{\sphinxbfcode{\sphinxupquote{get\_records}}}{\sphinxparam{\DUrole{n}{time\_steps}\DUrole{p}{:}\DUrole{w}{ }\DUrole{n}{ndarray}}}{{ $\rightarrow$ dict\DUrole{p}{{[}}str\DUrole{p}{,}\DUrole{w}{ }ndarray\DUrole{p}{{]}}}}
\pysigstopsignatures
\sphinxAtStartPar
Obtains full records of the trajectory for a given times (in seconds).
Essentially performs vectorised interpolation for given time series.
\begin{quote}\begin{description}
\sphinxlineitem{Parameters}
\sphinxAtStartPar
\sphinxstyleliteralstrong{\sphinxupquote{time\_steps}} (\sphinxstyleliteralemphasis{\sphinxupquote{np.ndarray}}) \textendash{} The time steps (in seconds) to obtain records for.

\sphinxlineitem{Returns}
\sphinxAtStartPar
A dictionary holding the records for each time step.

\sphinxlineitem{Return type}
\sphinxAtStartPar
dict{[}str, np.ndarray{]}

\end{description}\end{quote}

\end{fulllineitems}

\index{start\_time (qnav.waypoints.trajectory.Trajectory property)@\spxentry{start\_time}\spxextra{qnav.waypoints.trajectory.Trajectory property}}

\begin{fulllineitems}
\phantomsection\label{\detokenize{modules/waypoints/trajectory:qnav.waypoints.trajectory.Trajectory.start_time}}
\pysigstartsignatures
\pysigline{\sphinxbfcode{\sphinxupquote{property\DUrole{w}{ }}}\sphinxbfcode{\sphinxupquote{start\_time}}\sphinxbfcode{\sphinxupquote{\DUrole{p}{:}\DUrole{w}{ }float}}}
\pysigstopsignatures
\sphinxAtStartPar
Returns the time of the first trajectory record (in seconds).
\begin{quote}\begin{description}
\sphinxlineitem{Returns}
\sphinxAtStartPar
The start time of the trajectory

\sphinxlineitem{Return type}
\sphinxAtStartPar
float

\end{description}\end{quote}

\end{fulllineitems}

\index{end\_time (qnav.waypoints.trajectory.Trajectory property)@\spxentry{end\_time}\spxextra{qnav.waypoints.trajectory.Trajectory property}}

\begin{fulllineitems}
\phantomsection\label{\detokenize{modules/waypoints/trajectory:qnav.waypoints.trajectory.Trajectory.end_time}}
\pysigstartsignatures
\pysigline{\sphinxbfcode{\sphinxupquote{property\DUrole{w}{ }}}\sphinxbfcode{\sphinxupquote{end\_time}}\sphinxbfcode{\sphinxupquote{\DUrole{p}{:}\DUrole{w}{ }float}}}
\pysigstopsignatures
\sphinxAtStartPar
Returns the time of the final trajectory record (in seconds).
\begin{quote}\begin{description}
\sphinxlineitem{Returns}
\sphinxAtStartPar
The end time of the trajectory

\sphinxlineitem{Return type}
\sphinxAtStartPar
float

\end{description}\end{quote}

\end{fulllineitems}

\index{interp\_method (qnav.waypoints.trajectory.Trajectory property)@\spxentry{interp\_method}\spxextra{qnav.waypoints.trajectory.Trajectory property}}

\begin{fulllineitems}
\phantomsection\label{\detokenize{modules/waypoints/trajectory:qnav.waypoints.trajectory.Trajectory.interp_method}}
\pysigstartsignatures
\pysigline{\sphinxbfcode{\sphinxupquote{property\DUrole{w}{ }}}\sphinxbfcode{\sphinxupquote{interp\_method}}\sphinxbfcode{\sphinxupquote{\DUrole{p}{:}\DUrole{w}{ }str}}}
\pysigstopsignatures
\sphinxAtStartPar
The interpolation method used for ground truth generation.
\begin{quote}\begin{description}
\sphinxlineitem{Returns}
\sphinxAtStartPar
Interpolation name.

\sphinxlineitem{Return type}
\sphinxAtStartPar
str

\end{description}\end{quote}

\end{fulllineitems}


\end{fulllineitems}

\index{cumulative\_distance() (in module qnav.waypoints.trajectory)@\spxentry{cumulative\_distance()}\spxextra{in module qnav.waypoints.trajectory}}

\begin{fulllineitems}
\phantomsection\label{\detokenize{modules/waypoints/trajectory:qnav.waypoints.trajectory.cumulative_distance}}
\pysigstartsignatures
\pysiglinewithargsret{\sphinxcode{\sphinxupquote{qnav.waypoints.trajectory.}}\sphinxbfcode{\sphinxupquote{cumulative\_distance}}}{\sphinxparam{\DUrole{n}{lat\_points}\DUrole{p}{:}\DUrole{w}{ }\DUrole{n}{ndarray}}\sphinxparamcomma \sphinxparam{\DUrole{n}{lon\_points}\DUrole{p}{:}\DUrole{w}{ }\DUrole{n}{ndarray}}\sphinxparamcomma \sphinxparam{\DUrole{n}{alt\_points}\DUrole{p}{:}\DUrole{w}{ }\DUrole{n}{ndarray}}}{{ $\rightarrow$ ndarray}}
\pysigstopsignatures
\sphinxAtStartPar
Calculate the cumulative ‘great\sphinxhyphen{}circle’ distance between coordinate points.
For a given series of latitude\sphinxhyphen{}longitude\sphinxhyphen{}altitude points (given in
deg\sphinxhyphen{}deg\sphinxhyphen{}m), this calculates the cumulative distance between them using
the standard WGS\sphinxhyphen{}84 reference ellipsoid.
\begin{quote}\begin{description}
\sphinxlineitem{Parameters}\begin{itemize}
\item {} 
\sphinxAtStartPar
\sphinxstyleliteralstrong{\sphinxupquote{lat\_points}} (\sphinxstyleliteralemphasis{\sphinxupquote{np.ndarray}}\sphinxstyleliteralemphasis{\sphinxupquote{ (}}\sphinxstyleliteralemphasis{\sphinxupquote{n\sphinxhyphen{}elements}}\sphinxstyleliteralemphasis{\sphinxupquote{)}}) \textendash{} A series of latitude positions (in decimal degrees).

\item {} 
\sphinxAtStartPar
\sphinxstyleliteralstrong{\sphinxupquote{lon\_points}} \textendash{} A series of longitude positions (in decimal degrees).

\item {} 
\sphinxAtStartPar
\sphinxstyleliteralstrong{\sphinxupquote{alt\_points}} \textendash{} A series of altitude positions (in decimal degrees).

\end{itemize}

\sphinxlineitem{Returns}
\sphinxAtStartPar
The cumulative distance all points, starting from first.

\sphinxlineitem{Return type}
\sphinxAtStartPar
np.ndarray

\end{description}\end{quote}

\end{fulllineitems}


\sphinxstepscope


\subsubsection{qnav.waypoints.vehicle}
\label{\detokenize{modules/waypoints/vehicle:module-qnav.waypoints.vehicle}}\label{\detokenize{modules/waypoints/vehicle:qnav-waypoints-vehicle}}\label{\detokenize{modules/waypoints/vehicle::doc}}\index{module@\spxentry{module}!qnav.waypoints.vehicle@\spxentry{qnav.waypoints.vehicle}}\index{qnav.waypoints.vehicle@\spxentry{qnav.waypoints.vehicle}!module@\spxentry{module}}

\paragraph{vehicle.py}
\label{\detokenize{modules/waypoints/vehicle:vehicle-py}}\begin{quote}\begin{description}
\sphinxlineitem{summary}
\sphinxAtStartPar
Represents vehicle properties for waypoint generation.
This module contains a template Vehicle tuple listing all the properties
applied during waypoint generation. In addition to this a handful of
integrated vehicle profiles are include for testing. It is recommended
that a custom Vehicle instance is created for finer control.

\sphinxlineitem{authors}
\begin{DUlineblock}{0em}
\item[] Prof. Jason Ralph \sphinxhyphen{} \sphinxhref{mailto:jfralph@liverpool.ac.uk}{jfralph@liverpool.ac.uk}
\item[] Michael Wright \sphinxhyphen{} \sphinxhref{mailto:mjwright@liverpool.ac.uk}{mjwright@liverpool.ac.uk}
\end{DUlineblock}

\sphinxlineitem{version}
\sphinxAtStartPar
1.1.0 \sphinxhyphen{} Added integrated vehicle types for testing.

\end{description}\end{quote}
\index{Vehicle (class in qnav.waypoints.vehicle)@\spxentry{Vehicle}\spxextra{class in qnav.waypoints.vehicle}}

\begin{fulllineitems}
\phantomsection\label{\detokenize{modules/waypoints/vehicle:qnav.waypoints.vehicle.Vehicle}}
\pysigstartsignatures
\pysiglinewithargsret{\sphinxbfcode{\sphinxupquote{class\DUrole{w}{ }}}\sphinxcode{\sphinxupquote{qnav.waypoints.vehicle.}}\sphinxbfcode{\sphinxupquote{Vehicle}}}{\sphinxparam{\DUrole{n}{description}\DUrole{p}{:}\DUrole{w}{ }\DUrole{n}{str}\DUrole{w}{ }\DUrole{o}{=}\DUrole{w}{ }\DUrole{default_value}{\textquotesingle{}Unnamed\textquotesingle{}}}\sphinxparamcomma \sphinxparam{\DUrole{n}{turn\_rate\_max}\DUrole{p}{:}\DUrole{w}{ }\DUrole{n}{float}\DUrole{w}{ }\DUrole{o}{=}\DUrole{w}{ }\DUrole{default_value}{5.0}}\sphinxparamcomma \sphinxparam{\DUrole{n}{acceleration\_max}\DUrole{p}{:}\DUrole{w}{ }\DUrole{n}{float}\DUrole{w}{ }\DUrole{o}{=}\DUrole{w}{ }\DUrole{default_value}{2.0}}\sphinxparamcomma \sphinxparam{\DUrole{n}{deceleration\_max}\DUrole{p}{:}\DUrole{w}{ }\DUrole{n}{float}\DUrole{w}{ }\DUrole{o}{=}\DUrole{w}{ }\DUrole{default_value}{1.0}}\sphinxparamcomma \sphinxparam{\DUrole{n}{time\_delay\_acc}\DUrole{p}{:}\DUrole{w}{ }\DUrole{n}{float}\DUrole{w}{ }\DUrole{o}{=}\DUrole{w}{ }\DUrole{default_value}{1.0}}\sphinxparamcomma \sphinxparam{\DUrole{n}{time\_delay\_turn}\DUrole{p}{:}\DUrole{w}{ }\DUrole{n}{float}\DUrole{w}{ }\DUrole{o}{=}\DUrole{w}{ }\DUrole{default_value}{5.0}}\sphinxparamcomma \sphinxparam{\DUrole{n}{pitch\_constant\_offset}\DUrole{p}{:}\DUrole{w}{ }\DUrole{n}{float}\DUrole{w}{ }\DUrole{o}{=}\DUrole{w}{ }\DUrole{default_value}{0}}\sphinxparamcomma \sphinxparam{\DUrole{n}{pitch\_osc\_period}\DUrole{p}{:}\DUrole{w}{ }\DUrole{n}{float}\DUrole{w}{ }\DUrole{o}{=}\DUrole{w}{ }\DUrole{default_value}{0}}\sphinxparamcomma \sphinxparam{\DUrole{n}{pitch\_osc\_magnitude}\DUrole{p}{:}\DUrole{w}{ }\DUrole{n}{float}\DUrole{w}{ }\DUrole{o}{=}\DUrole{w}{ }\DUrole{default_value}{0}}\sphinxparamcomma \sphinxparam{\DUrole{n}{pitch\_max}\DUrole{p}{:}\DUrole{w}{ }\DUrole{n}{float}\DUrole{w}{ }\DUrole{o}{=}\DUrole{w}{ }\DUrole{default_value}{0}}\sphinxparamcomma \sphinxparam{\DUrole{n}{roll\_constant\_offset}\DUrole{p}{:}\DUrole{w}{ }\DUrole{n}{float}\DUrole{w}{ }\DUrole{o}{=}\DUrole{w}{ }\DUrole{default_value}{0}}\sphinxparamcomma \sphinxparam{\DUrole{n}{roll\_osc\_period}\DUrole{p}{:}\DUrole{w}{ }\DUrole{n}{float}\DUrole{w}{ }\DUrole{o}{=}\DUrole{w}{ }\DUrole{default_value}{0}}\sphinxparamcomma \sphinxparam{\DUrole{n}{roll\_osc\_magnitude}\DUrole{p}{:}\DUrole{w}{ }\DUrole{n}{float}\DUrole{w}{ }\DUrole{o}{=}\DUrole{w}{ }\DUrole{default_value}{0}}\sphinxparamcomma \sphinxparam{\DUrole{n}{roll\_max}\DUrole{p}{:}\DUrole{w}{ }\DUrole{n}{float}\DUrole{w}{ }\DUrole{o}{=}\DUrole{w}{ }\DUrole{default_value}{0}}\sphinxparamcomma \sphinxparam{\DUrole{n}{d\_roll\_d\_turn\_rate}\DUrole{p}{:}\DUrole{w}{ }\DUrole{n}{float}\DUrole{w}{ }\DUrole{o}{=}\DUrole{w}{ }\DUrole{default_value}{0}}\sphinxparamcomma \sphinxparam{\DUrole{n}{vibration\_noise\_acceleration}\DUrole{p}{:}\DUrole{w}{ }\DUrole{n}{float}\DUrole{w}{ }\DUrole{o}{=}\DUrole{w}{ }\DUrole{default_value}{0}}\sphinxparamcomma \sphinxparam{\DUrole{n}{vibration\_noise\_angle\_rates}\DUrole{p}{:}\DUrole{w}{ }\DUrole{n}{float}\DUrole{w}{ }\DUrole{o}{=}\DUrole{w}{ }\DUrole{default_value}{0}}\sphinxparamcomma \sphinxparam{\DUrole{n}{vibration\_damping\_period}\DUrole{p}{:}\DUrole{w}{ }\DUrole{n}{float}\DUrole{w}{ }\DUrole{o}{=}\DUrole{w}{ }\DUrole{default_value}{0}}\sphinxparamcomma \sphinxparam{\DUrole{n}{variation\_per\_second1}\DUrole{p}{:}\DUrole{w}{ }\DUrole{n}{float}\DUrole{w}{ }\DUrole{o}{=}\DUrole{w}{ }\DUrole{default_value}{0.005}}\sphinxparamcomma \sphinxparam{\DUrole{n}{variation\_per\_second2}\DUrole{p}{:}\DUrole{w}{ }\DUrole{n}{float}\DUrole{w}{ }\DUrole{o}{=}\DUrole{w}{ }\DUrole{default_value}{0.01}}\sphinxparamcomma \sphinxparam{\DUrole{n}{rand\_seed}\DUrole{p}{:}\DUrole{w}{ }\DUrole{n}{int}\DUrole{w}{ }\DUrole{o}{=}\DUrole{w}{ }\DUrole{default_value}{0}}}{}
\pysigstopsignatures
\sphinxAtStartPar
A structure used for holding vehicle model information.
This is used to describe limitations, offsets and expected
noise/oscillations. Unlike a normal class, the variables held within this
structure are immutable and cannot be changed. This is done to enforce
restrictions without them being accidentally overwritten later.
\index{description (qnav.waypoints.vehicle.Vehicle attribute)@\spxentry{description}\spxextra{qnav.waypoints.vehicle.Vehicle attribute}}

\begin{fulllineitems}
\phantomsection\label{\detokenize{modules/waypoints/vehicle:qnav.waypoints.vehicle.Vehicle.description}}
\pysigstartsignatures
\pysigline{\sphinxbfcode{\sphinxupquote{description}}\sphinxbfcode{\sphinxupquote{\DUrole{p}{:}\DUrole{w}{ }str}}}
\pysigstopsignatures
\sphinxAtStartPar
Alias for field number 0

\end{fulllineitems}

\index{turn\_rate\_max (qnav.waypoints.vehicle.Vehicle attribute)@\spxentry{turn\_rate\_max}\spxextra{qnav.waypoints.vehicle.Vehicle attribute}}

\begin{fulllineitems}
\phantomsection\label{\detokenize{modules/waypoints/vehicle:qnav.waypoints.vehicle.Vehicle.turn_rate_max}}
\pysigstartsignatures
\pysigline{\sphinxbfcode{\sphinxupquote{turn\_rate\_max}}\sphinxbfcode{\sphinxupquote{\DUrole{p}{:}\DUrole{w}{ }float}}}
\pysigstopsignatures
\sphinxAtStartPar
Alias for field number 1

\end{fulllineitems}

\index{acceleration\_max (qnav.waypoints.vehicle.Vehicle attribute)@\spxentry{acceleration\_max}\spxextra{qnav.waypoints.vehicle.Vehicle attribute}}

\begin{fulllineitems}
\phantomsection\label{\detokenize{modules/waypoints/vehicle:qnav.waypoints.vehicle.Vehicle.acceleration_max}}
\pysigstartsignatures
\pysigline{\sphinxbfcode{\sphinxupquote{acceleration\_max}}\sphinxbfcode{\sphinxupquote{\DUrole{p}{:}\DUrole{w}{ }float}}}
\pysigstopsignatures
\sphinxAtStartPar
Alias for field number 2

\end{fulllineitems}

\index{deceleration\_max (qnav.waypoints.vehicle.Vehicle attribute)@\spxentry{deceleration\_max}\spxextra{qnav.waypoints.vehicle.Vehicle attribute}}

\begin{fulllineitems}
\phantomsection\label{\detokenize{modules/waypoints/vehicle:qnav.waypoints.vehicle.Vehicle.deceleration_max}}
\pysigstartsignatures
\pysigline{\sphinxbfcode{\sphinxupquote{deceleration\_max}}\sphinxbfcode{\sphinxupquote{\DUrole{p}{:}\DUrole{w}{ }float}}}
\pysigstopsignatures
\sphinxAtStartPar
Alias for field number 3

\end{fulllineitems}

\index{time\_delay\_acc (qnav.waypoints.vehicle.Vehicle attribute)@\spxentry{time\_delay\_acc}\spxextra{qnav.waypoints.vehicle.Vehicle attribute}}

\begin{fulllineitems}
\phantomsection\label{\detokenize{modules/waypoints/vehicle:qnav.waypoints.vehicle.Vehicle.time_delay_acc}}
\pysigstartsignatures
\pysigline{\sphinxbfcode{\sphinxupquote{time\_delay\_acc}}\sphinxbfcode{\sphinxupquote{\DUrole{p}{:}\DUrole{w}{ }float}}}
\pysigstopsignatures
\sphinxAtStartPar
Alias for field number 4

\end{fulllineitems}

\index{time\_delay\_turn (qnav.waypoints.vehicle.Vehicle attribute)@\spxentry{time\_delay\_turn}\spxextra{qnav.waypoints.vehicle.Vehicle attribute}}

\begin{fulllineitems}
\phantomsection\label{\detokenize{modules/waypoints/vehicle:qnav.waypoints.vehicle.Vehicle.time_delay_turn}}
\pysigstartsignatures
\pysigline{\sphinxbfcode{\sphinxupquote{time\_delay\_turn}}\sphinxbfcode{\sphinxupquote{\DUrole{p}{:}\DUrole{w}{ }float}}}
\pysigstopsignatures
\sphinxAtStartPar
Alias for field number 5

\end{fulllineitems}

\index{pitch\_constant\_offset (qnav.waypoints.vehicle.Vehicle attribute)@\spxentry{pitch\_constant\_offset}\spxextra{qnav.waypoints.vehicle.Vehicle attribute}}

\begin{fulllineitems}
\phantomsection\label{\detokenize{modules/waypoints/vehicle:qnav.waypoints.vehicle.Vehicle.pitch_constant_offset}}
\pysigstartsignatures
\pysigline{\sphinxbfcode{\sphinxupquote{pitch\_constant\_offset}}\sphinxbfcode{\sphinxupquote{\DUrole{p}{:}\DUrole{w}{ }float}}}
\pysigstopsignatures
\sphinxAtStartPar
Alias for field number 6

\end{fulllineitems}

\index{pitch\_osc\_period (qnav.waypoints.vehicle.Vehicle attribute)@\spxentry{pitch\_osc\_period}\spxextra{qnav.waypoints.vehicle.Vehicle attribute}}

\begin{fulllineitems}
\phantomsection\label{\detokenize{modules/waypoints/vehicle:qnav.waypoints.vehicle.Vehicle.pitch_osc_period}}
\pysigstartsignatures
\pysigline{\sphinxbfcode{\sphinxupquote{pitch\_osc\_period}}\sphinxbfcode{\sphinxupquote{\DUrole{p}{:}\DUrole{w}{ }float}}}
\pysigstopsignatures
\sphinxAtStartPar
Alias for field number 7

\end{fulllineitems}

\index{pitch\_osc\_magnitude (qnav.waypoints.vehicle.Vehicle attribute)@\spxentry{pitch\_osc\_magnitude}\spxextra{qnav.waypoints.vehicle.Vehicle attribute}}

\begin{fulllineitems}
\phantomsection\label{\detokenize{modules/waypoints/vehicle:qnav.waypoints.vehicle.Vehicle.pitch_osc_magnitude}}
\pysigstartsignatures
\pysigline{\sphinxbfcode{\sphinxupquote{pitch\_osc\_magnitude}}\sphinxbfcode{\sphinxupquote{\DUrole{p}{:}\DUrole{w}{ }float}}}
\pysigstopsignatures
\sphinxAtStartPar
Alias for field number 8

\end{fulllineitems}

\index{pitch\_max (qnav.waypoints.vehicle.Vehicle attribute)@\spxentry{pitch\_max}\spxextra{qnav.waypoints.vehicle.Vehicle attribute}}

\begin{fulllineitems}
\phantomsection\label{\detokenize{modules/waypoints/vehicle:qnav.waypoints.vehicle.Vehicle.pitch_max}}
\pysigstartsignatures
\pysigline{\sphinxbfcode{\sphinxupquote{pitch\_max}}\sphinxbfcode{\sphinxupquote{\DUrole{p}{:}\DUrole{w}{ }float}}}
\pysigstopsignatures
\sphinxAtStartPar
Alias for field number 9

\end{fulllineitems}

\index{roll\_constant\_offset (qnav.waypoints.vehicle.Vehicle attribute)@\spxentry{roll\_constant\_offset}\spxextra{qnav.waypoints.vehicle.Vehicle attribute}}

\begin{fulllineitems}
\phantomsection\label{\detokenize{modules/waypoints/vehicle:qnav.waypoints.vehicle.Vehicle.roll_constant_offset}}
\pysigstartsignatures
\pysigline{\sphinxbfcode{\sphinxupquote{roll\_constant\_offset}}\sphinxbfcode{\sphinxupquote{\DUrole{p}{:}\DUrole{w}{ }float}}}
\pysigstopsignatures
\sphinxAtStartPar
Alias for field number 10

\end{fulllineitems}

\index{roll\_osc\_period (qnav.waypoints.vehicle.Vehicle attribute)@\spxentry{roll\_osc\_period}\spxextra{qnav.waypoints.vehicle.Vehicle attribute}}

\begin{fulllineitems}
\phantomsection\label{\detokenize{modules/waypoints/vehicle:qnav.waypoints.vehicle.Vehicle.roll_osc_period}}
\pysigstartsignatures
\pysigline{\sphinxbfcode{\sphinxupquote{roll\_osc\_period}}\sphinxbfcode{\sphinxupquote{\DUrole{p}{:}\DUrole{w}{ }float}}}
\pysigstopsignatures
\sphinxAtStartPar
Alias for field number 11

\end{fulllineitems}

\index{roll\_osc\_magnitude (qnav.waypoints.vehicle.Vehicle attribute)@\spxentry{roll\_osc\_magnitude}\spxextra{qnav.waypoints.vehicle.Vehicle attribute}}

\begin{fulllineitems}
\phantomsection\label{\detokenize{modules/waypoints/vehicle:qnav.waypoints.vehicle.Vehicle.roll_osc_magnitude}}
\pysigstartsignatures
\pysigline{\sphinxbfcode{\sphinxupquote{roll\_osc\_magnitude}}\sphinxbfcode{\sphinxupquote{\DUrole{p}{:}\DUrole{w}{ }float}}}
\pysigstopsignatures
\sphinxAtStartPar
Alias for field number 12

\end{fulllineitems}

\index{roll\_max (qnav.waypoints.vehicle.Vehicle attribute)@\spxentry{roll\_max}\spxextra{qnav.waypoints.vehicle.Vehicle attribute}}

\begin{fulllineitems}
\phantomsection\label{\detokenize{modules/waypoints/vehicle:qnav.waypoints.vehicle.Vehicle.roll_max}}
\pysigstartsignatures
\pysigline{\sphinxbfcode{\sphinxupquote{roll\_max}}\sphinxbfcode{\sphinxupquote{\DUrole{p}{:}\DUrole{w}{ }float}}}
\pysigstopsignatures
\sphinxAtStartPar
Alias for field number 13

\end{fulllineitems}

\index{d\_roll\_d\_turn\_rate (qnav.waypoints.vehicle.Vehicle attribute)@\spxentry{d\_roll\_d\_turn\_rate}\spxextra{qnav.waypoints.vehicle.Vehicle attribute}}

\begin{fulllineitems}
\phantomsection\label{\detokenize{modules/waypoints/vehicle:qnav.waypoints.vehicle.Vehicle.d_roll_d_turn_rate}}
\pysigstartsignatures
\pysigline{\sphinxbfcode{\sphinxupquote{d\_roll\_d\_turn\_rate}}\sphinxbfcode{\sphinxupquote{\DUrole{p}{:}\DUrole{w}{ }float}}}
\pysigstopsignatures
\sphinxAtStartPar
Alias for field number 14

\end{fulllineitems}

\index{vibration\_noise\_acceleration (qnav.waypoints.vehicle.Vehicle attribute)@\spxentry{vibration\_noise\_acceleration}\spxextra{qnav.waypoints.vehicle.Vehicle attribute}}

\begin{fulllineitems}
\phantomsection\label{\detokenize{modules/waypoints/vehicle:qnav.waypoints.vehicle.Vehicle.vibration_noise_acceleration}}
\pysigstartsignatures
\pysigline{\sphinxbfcode{\sphinxupquote{vibration\_noise\_acceleration}}\sphinxbfcode{\sphinxupquote{\DUrole{p}{:}\DUrole{w}{ }float}}}
\pysigstopsignatures
\sphinxAtStartPar
Alias for field number 15

\end{fulllineitems}

\index{vibration\_noise\_angle\_rates (qnav.waypoints.vehicle.Vehicle attribute)@\spxentry{vibration\_noise\_angle\_rates}\spxextra{qnav.waypoints.vehicle.Vehicle attribute}}

\begin{fulllineitems}
\phantomsection\label{\detokenize{modules/waypoints/vehicle:qnav.waypoints.vehicle.Vehicle.vibration_noise_angle_rates}}
\pysigstartsignatures
\pysigline{\sphinxbfcode{\sphinxupquote{vibration\_noise\_angle\_rates}}\sphinxbfcode{\sphinxupquote{\DUrole{p}{:}\DUrole{w}{ }float}}}
\pysigstopsignatures
\sphinxAtStartPar
Alias for field number 16

\end{fulllineitems}

\index{vibration\_damping\_period (qnav.waypoints.vehicle.Vehicle attribute)@\spxentry{vibration\_damping\_period}\spxextra{qnav.waypoints.vehicle.Vehicle attribute}}

\begin{fulllineitems}
\phantomsection\label{\detokenize{modules/waypoints/vehicle:qnav.waypoints.vehicle.Vehicle.vibration_damping_period}}
\pysigstartsignatures
\pysigline{\sphinxbfcode{\sphinxupquote{vibration\_damping\_period}}\sphinxbfcode{\sphinxupquote{\DUrole{p}{:}\DUrole{w}{ }float}}}
\pysigstopsignatures
\sphinxAtStartPar
Alias for field number 17

\end{fulllineitems}

\index{variation\_per\_second1 (qnav.waypoints.vehicle.Vehicle attribute)@\spxentry{variation\_per\_second1}\spxextra{qnav.waypoints.vehicle.Vehicle attribute}}

\begin{fulllineitems}
\phantomsection\label{\detokenize{modules/waypoints/vehicle:qnav.waypoints.vehicle.Vehicle.variation_per_second1}}
\pysigstartsignatures
\pysigline{\sphinxbfcode{\sphinxupquote{variation\_per\_second1}}\sphinxbfcode{\sphinxupquote{\DUrole{p}{:}\DUrole{w}{ }float}}}
\pysigstopsignatures
\sphinxAtStartPar
Alias for field number 18

\end{fulllineitems}

\index{variation\_per\_second2 (qnav.waypoints.vehicle.Vehicle attribute)@\spxentry{variation\_per\_second2}\spxextra{qnav.waypoints.vehicle.Vehicle attribute}}

\begin{fulllineitems}
\phantomsection\label{\detokenize{modules/waypoints/vehicle:qnav.waypoints.vehicle.Vehicle.variation_per_second2}}
\pysigstartsignatures
\pysigline{\sphinxbfcode{\sphinxupquote{variation\_per\_second2}}\sphinxbfcode{\sphinxupquote{\DUrole{p}{:}\DUrole{w}{ }float}}}
\pysigstopsignatures
\sphinxAtStartPar
Alias for field number 19

\end{fulllineitems}

\index{rand\_seed (qnav.waypoints.vehicle.Vehicle attribute)@\spxentry{rand\_seed}\spxextra{qnav.waypoints.vehicle.Vehicle attribute}}

\begin{fulllineitems}
\phantomsection\label{\detokenize{modules/waypoints/vehicle:qnav.waypoints.vehicle.Vehicle.rand_seed}}
\pysigstartsignatures
\pysigline{\sphinxbfcode{\sphinxupquote{rand\_seed}}\sphinxbfcode{\sphinxupquote{\DUrole{p}{:}\DUrole{w}{ }int}}}
\pysigstopsignatures
\sphinxAtStartPar
Alias for field number 20

\end{fulllineitems}


\end{fulllineitems}

\index{get\_default\_profile() (in module qnav.waypoints.vehicle)@\spxentry{get\_default\_profile()}\spxextra{in module qnav.waypoints.vehicle}}

\begin{fulllineitems}
\phantomsection\label{\detokenize{modules/waypoints/vehicle:qnav.waypoints.vehicle.get_default_profile}}
\pysigstartsignatures
\pysiglinewithargsret{\sphinxcode{\sphinxupquote{qnav.waypoints.vehicle.}}\sphinxbfcode{\sphinxupquote{get\_default\_profile}}}{\sphinxparam{\DUrole{n}{rand\_seed}\DUrole{p}{:}\DUrole{w}{ }\DUrole{n}{int}\DUrole{w}{ }\DUrole{o}{=}\DUrole{w}{ }\DUrole{default_value}{0}}}{{ $\rightarrow$ {\hyperref[\detokenize{modules/waypoints/vehicle:qnav.waypoints.vehicle.Vehicle}]{\sphinxcrossref{Vehicle}}}}}
\pysigstopsignatures
\sphinxAtStartPar
Returns a Vehicle instance representing the default profile.
This profile is for a typical aircraft with no noise/vibrations.
\begin{quote}\begin{description}
\sphinxlineitem{Parameters}
\sphinxAtStartPar
\sphinxstyleliteralstrong{\sphinxupquote{rand\_seed}} (\sphinxstyleliteralemphasis{\sphinxupquote{int}}) \textendash{} The seed to use for random number generation.

\sphinxlineitem{Returns}
\sphinxAtStartPar
The default Vehicle profile.

\sphinxlineitem{Return type}
\sphinxAtStartPar
{\hyperref[\detokenize{modules/waypoints/vehicle:qnav.waypoints.vehicle.Vehicle}]{\sphinxcrossref{Vehicle}}}

\end{description}\end{quote}

\end{fulllineitems}

\index{get\_large\_aircraft\_profile() (in module qnav.waypoints.vehicle)@\spxentry{get\_large\_aircraft\_profile()}\spxextra{in module qnav.waypoints.vehicle}}

\begin{fulllineitems}
\phantomsection\label{\detokenize{modules/waypoints/vehicle:qnav.waypoints.vehicle.get_large_aircraft_profile}}
\pysigstartsignatures
\pysiglinewithargsret{\sphinxcode{\sphinxupquote{qnav.waypoints.vehicle.}}\sphinxbfcode{\sphinxupquote{get\_large\_aircraft\_profile}}}{\sphinxparam{\DUrole{n}{rand\_seed}\DUrole{p}{:}\DUrole{w}{ }\DUrole{n}{int}\DUrole{w}{ }\DUrole{o}{=}\DUrole{w}{ }\DUrole{default_value}{0}}}{{ $\rightarrow$ {\hyperref[\detokenize{modules/waypoints/vehicle:qnav.waypoints.vehicle.Vehicle}]{\sphinxcrossref{Vehicle}}}}}
\pysigstopsignatures
\sphinxAtStartPar
Returns a Vehicle instance representing that of a large aircraft.
This profile supplies the parameters for a 737 type of air\sphinxhyphen{}vehicle,
stating slow acceleration and deceleration with wide turning.
\begin{quote}\begin{description}
\sphinxlineitem{Parameters}
\sphinxAtStartPar
\sphinxstyleliteralstrong{\sphinxupquote{rand\_seed}} (\sphinxstyleliteralemphasis{\sphinxupquote{int}}) \textendash{} The seed to use for random number generation.

\sphinxlineitem{Returns}
\sphinxAtStartPar
The Vehicle profile for a large aircraft.

\sphinxlineitem{Return type}
\sphinxAtStartPar
{\hyperref[\detokenize{modules/waypoints/vehicle:qnav.waypoints.vehicle.Vehicle}]{\sphinxcrossref{Vehicle}}}

\end{description}\end{quote}

\end{fulllineitems}

\index{get\_large\_ship\_profile() (in module qnav.waypoints.vehicle)@\spxentry{get\_large\_ship\_profile()}\spxextra{in module qnav.waypoints.vehicle}}

\begin{fulllineitems}
\phantomsection\label{\detokenize{modules/waypoints/vehicle:qnav.waypoints.vehicle.get_large_ship_profile}}
\pysigstartsignatures
\pysiglinewithargsret{\sphinxcode{\sphinxupquote{qnav.waypoints.vehicle.}}\sphinxbfcode{\sphinxupquote{get\_large\_ship\_profile}}}{\sphinxparam{\DUrole{n}{rand\_seed}\DUrole{p}{:}\DUrole{w}{ }\DUrole{n}{int}\DUrole{w}{ }\DUrole{o}{=}\DUrole{w}{ }\DUrole{default_value}{0}}}{{ $\rightarrow$ {\hyperref[\detokenize{modules/waypoints/vehicle:qnav.waypoints.vehicle.Vehicle}]{\sphinxcrossref{Vehicle}}}}}
\pysigstopsignatures
\sphinxAtStartPar
Returns a Vehicle instance representing that of a large aircraft.
This profile supplies the parameters for a large ship, such as a cargo
vessel.
\begin{quote}\begin{description}
\sphinxlineitem{Parameters}
\sphinxAtStartPar
\sphinxstyleliteralstrong{\sphinxupquote{rand\_seed}} (\sphinxstyleliteralemphasis{\sphinxupquote{int}}) \textendash{} The seed to use for random number generation.

\sphinxlineitem{Returns}
\sphinxAtStartPar
The Vehicle profile for a large ship.

\sphinxlineitem{Return type}
\sphinxAtStartPar
{\hyperref[\detokenize{modules/waypoints/vehicle:qnav.waypoints.vehicle.Vehicle}]{\sphinxcrossref{Vehicle}}}

\end{description}\end{quote}

\end{fulllineitems}

\index{get\_large\_van\_profile() (in module qnav.waypoints.vehicle)@\spxentry{get\_large\_van\_profile()}\spxextra{in module qnav.waypoints.vehicle}}

\begin{fulllineitems}
\phantomsection\label{\detokenize{modules/waypoints/vehicle:qnav.waypoints.vehicle.get_large_van_profile}}
\pysigstartsignatures
\pysiglinewithargsret{\sphinxcode{\sphinxupquote{qnav.waypoints.vehicle.}}\sphinxbfcode{\sphinxupquote{get\_large\_van\_profile}}}{\sphinxparam{\DUrole{n}{rand\_seed}\DUrole{p}{:}\DUrole{w}{ }\DUrole{n}{int}\DUrole{w}{ }\DUrole{o}{=}\DUrole{w}{ }\DUrole{default_value}{0}}}{{ $\rightarrow$ {\hyperref[\detokenize{modules/waypoints/vehicle:qnav.waypoints.vehicle.Vehicle}]{\sphinxcrossref{Vehicle}}}}}
\pysigstopsignatures
\sphinxAtStartPar
Returns a Vehicle instance representing that of a large van.
This profile supplies the parameters for a large ground vehicle, such as
a heavy van vessel, slow acceleration/deceleration and little turning
limitations.
\begin{quote}\begin{description}
\sphinxlineitem{Parameters}
\sphinxAtStartPar
\sphinxstyleliteralstrong{\sphinxupquote{rand\_seed}} (\sphinxstyleliteralemphasis{\sphinxupquote{int}}) \textendash{} The seed to use for random number generation.

\sphinxlineitem{Returns}
\sphinxAtStartPar
The Vehicle profile for a large van.

\sphinxlineitem{Return type}
\sphinxAtStartPar
{\hyperref[\detokenize{modules/waypoints/vehicle:qnav.waypoints.vehicle.Vehicle}]{\sphinxcrossref{Vehicle}}}

\end{description}\end{quote}

\end{fulllineitems}

\index{get\_submarine\_profile() (in module qnav.waypoints.vehicle)@\spxentry{get\_submarine\_profile()}\spxextra{in module qnav.waypoints.vehicle}}

\begin{fulllineitems}
\phantomsection\label{\detokenize{modules/waypoints/vehicle:qnav.waypoints.vehicle.get_submarine_profile}}
\pysigstartsignatures
\pysiglinewithargsret{\sphinxcode{\sphinxupquote{qnav.waypoints.vehicle.}}\sphinxbfcode{\sphinxupquote{get\_submarine\_profile}}}{\sphinxparam{\DUrole{n}{rand\_seed}\DUrole{p}{:}\DUrole{w}{ }\DUrole{n}{int}\DUrole{w}{ }\DUrole{o}{=}\DUrole{w}{ }\DUrole{default_value}{0}}}{{ $\rightarrow$ {\hyperref[\detokenize{modules/waypoints/vehicle:qnav.waypoints.vehicle.Vehicle}]{\sphinxcrossref{Vehicle}}}}}
\pysigstopsignatures
\sphinxAtStartPar
Returns a Vehicle instance representing that of a submarine.
This profile supplies the parameters for a submarine vessel.
\begin{quote}\begin{description}
\sphinxlineitem{Parameters}
\sphinxAtStartPar
\sphinxstyleliteralstrong{\sphinxupquote{rand\_seed}} (\sphinxstyleliteralemphasis{\sphinxupquote{int}}) \textendash{} The seed to use for random number generation.

\sphinxlineitem{Returns}
\sphinxAtStartPar
The Vehicle profile for a submarine vessel.

\sphinxlineitem{Return type}
\sphinxAtStartPar
{\hyperref[\detokenize{modules/waypoints/vehicle:qnav.waypoints.vehicle.Vehicle}]{\sphinxcrossref{Vehicle}}}

\end{description}\end{quote}

\end{fulllineitems}

\index{get\_unrestricted\_profile() (in module qnav.waypoints.vehicle)@\spxentry{get\_unrestricted\_profile()}\spxextra{in module qnav.waypoints.vehicle}}

\begin{fulllineitems}
\phantomsection\label{\detokenize{modules/waypoints/vehicle:qnav.waypoints.vehicle.get_unrestricted_profile}}
\pysigstartsignatures
\pysiglinewithargsret{\sphinxcode{\sphinxupquote{qnav.waypoints.vehicle.}}\sphinxbfcode{\sphinxupquote{get\_unrestricted\_profile}}}{}{{ $\rightarrow$ {\hyperref[\detokenize{modules/waypoints/vehicle:qnav.waypoints.vehicle.Vehicle}]{\sphinxcrossref{Vehicle}}}}}
\pysigstopsignatures
\sphinxAtStartPar
Returns a Vehicle instance with no movement restrictions
This profile is solely used for following tightly placed waypoints in
situations where placing vehicle limitations would be impractical. No
vehicle noise or vibrations are contained in this profile.
\begin{quote}\begin{description}
\sphinxlineitem{Returns}
\sphinxAtStartPar
The Vehicle profile with no limitations or noise.

\sphinxlineitem{Return type}
\sphinxAtStartPar
{\hyperref[\detokenize{modules/waypoints/vehicle:qnav.waypoints.vehicle.Vehicle}]{\sphinxcrossref{Vehicle}}}

\end{description}\end{quote}

\end{fulllineitems}


\sphinxstepscope


\subsubsection{qnav.waypoints.tracks}
\label{\detokenize{modules/waypoints/tracks:module-qnav.waypoints.tracks}}\label{\detokenize{modules/waypoints/tracks:qnav-waypoints-tracks}}\label{\detokenize{modules/waypoints/tracks::doc}}\index{module@\spxentry{module}!qnav.waypoints.tracks@\spxentry{qnav.waypoints.tracks}}\index{qnav.waypoints.tracks@\spxentry{qnav.waypoints.tracks}!module@\spxentry{module}}

\paragraph{tracks.py}
\label{\detokenize{modules/waypoints/tracks:tracks-py}}\begin{quote}\begin{description}
\sphinxlineitem{summary}
\sphinxAtStartPar
Provides default waypoint generation for testing tracks.
This module contains various functions for generating different default
waypoint tracks. This is an alternative instead of using user supplied
trajectories. This is namely used for testing, but can also be used for
rapid experimentation when a custom trajectory is unavailable.

\sphinxlineitem{authors}
\begin{DUlineblock}{0em}
\item[] Prof. Jason Ralph \sphinxhyphen{} \sphinxhref{mailto:jfralph@liverpool.ac.uk}{jfralph@liverpool.ac.uk}
\item[] Michael Wright \sphinxhyphen{} \sphinxhref{mailto:mjwright@liverpool.ac.uk}{mjwright@liverpool.ac.uk}
\end{DUlineblock}

\sphinxlineitem{version}
\sphinxAtStartPar
1.0.0 \sphinxhyphen{} First full implementation of module.

\end{description}\end{quote}
\index{generate\_racetrack() (in module qnav.waypoints.tracks)@\spxentry{generate\_racetrack()}\spxextra{in module qnav.waypoints.tracks}}

\begin{fulllineitems}
\phantomsection\label{\detokenize{modules/waypoints/tracks:qnav.waypoints.tracks.generate_racetrack}}
\pysigstartsignatures
\pysiglinewithargsret{\sphinxcode{\sphinxupquote{qnav.waypoints.tracks.}}\sphinxbfcode{\sphinxupquote{generate\_racetrack}}}{\sphinxparam{\DUrole{n}{frequency}\DUrole{p}{:}\DUrole{w}{ }\DUrole{n}{float}\DUrole{w}{ }\DUrole{o}{=}\DUrole{w}{ }\DUrole{default_value}{1000.0}}\sphinxparamcomma \sphinxparam{\DUrole{n}{max\_time}\DUrole{p}{:}\DUrole{w}{ }\DUrole{n}{float}\DUrole{w}{ }\DUrole{o}{=}\DUrole{w}{ }\DUrole{default_value}{500}}\sphinxparamcomma \sphinxparam{\DUrole{n}{lap\_time}\DUrole{p}{:}\DUrole{w}{ }\DUrole{n}{float}\DUrole{w}{ }\DUrole{o}{=}\DUrole{w}{ }\DUrole{default_value}{500}}\sphinxparamcomma \sphinxparam{\DUrole{n}{radius\_x}\DUrole{p}{:}\DUrole{w}{ }\DUrole{n}{float}\DUrole{w}{ }\DUrole{o}{=}\DUrole{w}{ }\DUrole{default_value}{1000.0}}\sphinxparamcomma \sphinxparam{\DUrole{n}{radius\_y}\DUrole{p}{:}\DUrole{w}{ }\DUrole{n}{float}\DUrole{w}{ }\DUrole{o}{=}\DUrole{w}{ }\DUrole{default_value}{1000.0}}\sphinxparamcomma \sphinxparam{\DUrole{n}{radius\_z}\DUrole{p}{:}\DUrole{w}{ }\DUrole{n}{float}\DUrole{w}{ }\DUrole{o}{=}\DUrole{w}{ }\DUrole{default_value}{5.0}}\sphinxparamcomma \sphinxparam{\DUrole{n}{theta\_x}\DUrole{p}{:}\DUrole{w}{ }\DUrole{n}{float}\DUrole{w}{ }\DUrole{o}{=}\DUrole{w}{ }\DUrole{default_value}{0.0}}\sphinxparamcomma \sphinxparam{\DUrole{n}{theta\_y}\DUrole{p}{:}\DUrole{w}{ }\DUrole{n}{float}\DUrole{w}{ }\DUrole{o}{=}\DUrole{w}{ }\DUrole{default_value}{0.0}}\sphinxparamcomma \sphinxparam{\DUrole{n}{theta\_z}\DUrole{p}{:}\DUrole{w}{ }\DUrole{n}{float}\DUrole{w}{ }\DUrole{o}{=}\DUrole{w}{ }\DUrole{default_value}{0.0}}\sphinxparamcomma \sphinxparam{\DUrole{n}{ref\_lat}\DUrole{p}{:}\DUrole{w}{ }\DUrole{n}{float}\DUrole{w}{ }\DUrole{o}{=}\DUrole{w}{ }\DUrole{default_value}{0.0}}\sphinxparamcomma \sphinxparam{\DUrole{n}{ref\_lon}\DUrole{p}{:}\DUrole{w}{ }\DUrole{n}{float}\DUrole{w}{ }\DUrole{o}{=}\DUrole{w}{ }\DUrole{default_value}{0.0}}\sphinxparamcomma \sphinxparam{\DUrole{n}{ref\_alt}\DUrole{p}{:}\DUrole{w}{ }\DUrole{n}{float}\DUrole{w}{ }\DUrole{o}{=}\DUrole{w}{ }\DUrole{default_value}{0.0}}\sphinxparamcomma \sphinxparam{\DUrole{n}{num\_var\_y}\DUrole{p}{:}\DUrole{w}{ }\DUrole{n}{float}\DUrole{w}{ }\DUrole{o}{=}\DUrole{w}{ }\DUrole{default_value}{2}}\sphinxparamcomma \sphinxparam{\DUrole{n}{num\_var\_z}\DUrole{p}{:}\DUrole{w}{ }\DUrole{n}{float}\DUrole{w}{ }\DUrole{o}{=}\DUrole{w}{ }\DUrole{default_value}{3}}}{{ $\rightarrow$ ndarray}}
\pysigstopsignatures
\sphinxAtStartPar
Generates waypoints for a racetrack of requested dimensions.
This function produces waypoints at a requested frequency that represent
a simple looping path. The produced waypoints can be then used with
previously implemented waypoint functions, essentially (efficiently)
pre\sphinxhyphen{}calculating the values at each position along the racetrack.
\begin{quote}\begin{description}
\sphinxlineitem{Parameters}\begin{itemize}
\item {} 
\sphinxAtStartPar
\sphinxstyleliteralstrong{\sphinxupquote{frequency}} (\sphinxstyleliteralemphasis{\sphinxupquote{float}}) \textendash{} The simulation’s update frequency (in Hz). This is used
for calculating the number of intermediate points and total number of
waypoints to return.

\item {} 
\sphinxAtStartPar
\sphinxstyleliteralstrong{\sphinxupquote{max\_time}} (\sphinxstyleliteralemphasis{\sphinxupquote{float}}) \textendash{} The time at which the simulation is to end (in seconds).
No waypoints will be generated for points after this time

\item {} 
\sphinxAtStartPar
\sphinxstyleliteralstrong{\sphinxupquote{lap\_time}} (\sphinxstyleliteralemphasis{\sphinxupquote{float}}) \textendash{} The time period (in seconds) for one circuit of the track.
This is used to calculate how fast the vehicle is moving around the
generated racetrack.

\item {} 
\sphinxAtStartPar
\sphinxstyleliteralstrong{\sphinxupquote{radius\_x}} (\sphinxstyleliteralemphasis{\sphinxupquote{float}}) \textendash{} The size of the racetrack along the x\sphinxhyphen{}direction (in metres).

\item {} 
\sphinxAtStartPar
\sphinxstyleliteralstrong{\sphinxupquote{radius\_y}} (\sphinxstyleliteralemphasis{\sphinxupquote{float}}) \textendash{} The size of the racetrack along the y\sphinxhyphen{}direction (in metres).

\item {} 
\sphinxAtStartPar
\sphinxstyleliteralstrong{\sphinxupquote{radius\_z}} (\sphinxstyleliteralemphasis{\sphinxupquote{float}}) \textendash{} The height variations of the racetrack (in metres).

\item {} 
\sphinxAtStartPar
\sphinxstyleliteralstrong{\sphinxupquote{theta\_x}} (\sphinxstyleliteralemphasis{\sphinxupquote{float}}) \textendash{} The starting X position on the circuit.

\item {} 
\sphinxAtStartPar
\sphinxstyleliteralstrong{\sphinxupquote{theta\_y}} (\sphinxstyleliteralemphasis{\sphinxupquote{float}}) \textendash{} The starting Y position on the circuit.

\item {} 
\sphinxAtStartPar
\sphinxstyleliteralstrong{\sphinxupquote{theta\_z}} (\sphinxstyleliteralemphasis{\sphinxupquote{float}}) \textendash{} The starting Z position on the circuit.

\item {} 
\sphinxAtStartPar
\sphinxstyleliteralstrong{\sphinxupquote{ref\_lat}} (\sphinxstyleliteralemphasis{\sphinxupquote{float}}) \textendash{} The reference latitude of the centre position of the
circuit (in degrees).

\item {} 
\sphinxAtStartPar
\sphinxstyleliteralstrong{\sphinxupquote{ref\_lon}} (\sphinxstyleliteralemphasis{\sphinxupquote{float}}) \textendash{} The reference longitude of the centre position of the
circuit (in degrees).

\item {} 
\sphinxAtStartPar
\sphinxstyleliteralstrong{\sphinxupquote{ref\_alt}} (\sphinxstyleliteralemphasis{\sphinxupquote{float}}) \textendash{} The reference altitude of the centre position of the
circuit (in degrees).

\item {} 
\sphinxAtStartPar
\sphinxstyleliteralstrong{\sphinxupquote{num\_var\_y}} (\sphinxstyleliteralemphasis{\sphinxupquote{float}}) \textendash{} Number of periods of y variations for each circuit in x.

\item {} 
\sphinxAtStartPar
\sphinxstyleliteralstrong{\sphinxupquote{num\_var\_z}} (\sphinxstyleliteralemphasis{\sphinxupquote{float}}) \textendash{} Number of periods of height variations for each circuit in x.

\end{itemize}

\sphinxlineitem{Returns}
\sphinxAtStartPar
The calculated latitude, longitude and altitude for each time
step from 0 to the maximum time requested.

\sphinxlineitem{Return type}
\sphinxAtStartPar
numpy.ndarray (n\sphinxhyphen{}by\sphinxhyphen{}4 elements)

\end{description}\end{quote}

\end{fulllineitems}


\sphinxstepscope


\section{System Requirements}
\label{\detokenize{usage/requirements:system-requirements}}\label{\detokenize{usage/requirements::doc}}
\sphinxAtStartPar
The finalised toolbox requires various hardware and software requirements to be efficiently used.
The minimum and suggested requirements have been acquired based on experimentation and auditing of the toolbox.


\subsection{Hardware Requirements}
\label{\detokenize{usage/requirements:hardware-requirements}}
\sphinxAtStartPar
The hardware requirements are subject to the requested input. The resources needed can be reduced by decreasing the
requested frequency/resolution and increasing the output down\sphinxhyphen{}sampling via the configuration file.


\subsubsection{Minimum Specifications}
\label{\detokenize{usage/requirements:minimum-specifications}}
\sphinxAtStartPar
The minimum requirements to use the toolbox are outlined in the table below. The host system must at a minimum have
these hardware specifications to use the toolbox at a basic (limited) level. Updates and potential additions to the
toolbox may increase these minimum requirements.


\begin{savenotes}\sphinxattablestart
\sphinxthistablewithglobalstyle
\centering
\begin{tabulary}{\linewidth}[t]{TT}
\sphinxtoprule
\sphinxstyletheadfamily 
\sphinxAtStartPar
Hardware
&\sphinxstyletheadfamily 
\sphinxAtStartPar
Minimum Resources Required
\\
\sphinxmidrule
\sphinxtableatstartofbodyhook
\sphinxAtStartPar
Processor (CPU)
&
\sphinxAtStartPar
x86 32\sphinxhyphen{}bit processor 1.00 GHz (or better)
\\
\sphinxhline
\sphinxAtStartPar
Physical Memory (RAM)
&
\sphinxAtStartPar
2.00 GB of memory (or larger)
\\
\sphinxhline
\sphinxAtStartPar
Disk Storage Space (HDD)
&
\sphinxAtStartPar
20.00 GB available (or more)
\\
\sphinxbottomrule
\end{tabulary}
\sphinxtableafterendhook\par
\sphinxattableend\end{savenotes}


\subsubsection{Recommended Specifications}
\label{\detokenize{usage/requirements:recommended-specifications}}
\sphinxAtStartPar
The recommended requirements to fully use the toolbox are outlined in the table below. If the host system has these
hardware specifications or higher it is guaranteed to use the software to its full extent.


\begin{savenotes}\sphinxattablestart
\sphinxthistablewithglobalstyle
\centering
\begin{tabulary}{\linewidth}[t]{TT}
\sphinxtoprule
\sphinxstyletheadfamily 
\sphinxAtStartPar
Hardware
&\sphinxstyletheadfamily 
\sphinxAtStartPar
Recommended Resources Required
\\
\sphinxmidrule
\sphinxtableatstartofbodyhook
\sphinxAtStartPar
Processor (CPU)
&
\sphinxAtStartPar
x86 64\sphinxhyphen{}bit processor 2.00 GHz (or better)
\\
\sphinxhline
\sphinxAtStartPar
Physical Memory (RAM)
&
\sphinxAtStartPar
8.00 GB of memory (or larger)
\\
\sphinxhline
\sphinxAtStartPar
Disk Storage Space (HDD)
&
\sphinxAtStartPar
100.00 GB available (or more)
\\
\sphinxbottomrule
\end{tabulary}
\sphinxtableafterendhook\par
\sphinxattableend\end{savenotes}


\subsubsection{Note About Memory and Storage}
\label{\detokenize{usage/requirements:note-about-memory-and-storage}}
\sphinxAtStartPar
Executing simulations can demand a large amount of RAM. This is especially true if high frequencies and/or a large
amount of waypoint data is used. For example, a single simulation at 1,000 Hz is expected to require around 1 GB of RAM
for each simulated hour due to the amount of data that needs to be produced and recorded. To significantly reduce this
resource requirement various settings can be changed.

\sphinxAtStartPar
Firstly, virtual memory can be used for offloading waypoint and estimation results to disk. This is enabled under
the \sphinxstylestrong{Waypoint} and \sphinxstylestrong{Estimation} sections in the configuration file via the \sphinxcode{\sphinxupquote{useVirtualMemory}} options. In
addition, an \sphinxcode{\sphinxupquote{outputDownSampleRate}} can be specified under the \sphinxstylestrong{Output} section, meaning that only a sub\sphinxhyphen{}sampled
amount of the estimation results will be collected and saved. Lastly, at a worst case scenario the waypoint and
measurement frequencies can be reduced.

\begin{sphinxadmonition}{warning}{Warning:}
\sphinxAtStartPar
The use of virtual memory on 32\sphinxhyphen{}bit systems is limited to at most 2 GB. This means that waypoints and
simulation results cannot use more than 2 GB if the use of virtual memory is enabled on 32\sphinxhyphen{}bit systems. If this
limitation causes significant problems requested updates to the toolbox can be made to overcome it.
\end{sphinxadmonition}

\sphinxAtStartPar
In the event that the output data needs to be saved in multiple formats, but there is insufficient storage space,
a single format (such as simple binary) can used which can later be converted. Using the read and write modules inside
the output sub\sphinxhyphen{}package, the single generated file can be read and its content can then be written into the other
needed formats.

\begin{DUlineblock}{0em}
\item[] 
\end{DUlineblock}


\subsection{Software Requirements}
\label{\detokenize{usage/requirements:software-requirements}}

\subsubsection{Operating System}
\label{\detokenize{usage/requirements:operating-system}}
\sphinxAtStartPar
The toolbox has been developed with platform independence constantly in mind, taking fully advantage of Python’s
flexibility.

\sphinxAtStartPar
The handling of file paths and managing local resources has been done without replaying on any operating system
specific features/formats. Likewise, the implementation of multi\sphinxhyphen{}processing functions has been done in a way that
acknowledges adapts to the different ways operating system handle communication between multiple threads/processes.
As a result, the same project build can be used on any major platform without any alteration to its source code.

\sphinxAtStartPar
The toolbox has been tested and proved to function flawlessly on Windows, Mac and Linux systems. Specifically:
\begin{itemize}
\item {} 
\sphinxAtStartPar
Windows 10 \sphinxhyphen{} Enterprise

\item {} 
\sphinxAtStartPar
Linux \sphinxhyphen{} Ubuntu LTS 22.04

\item {} 
\sphinxAtStartPar
MacOS \sphinxhyphen{} Sonoma 14.6.1

\item {} 
\sphinxAtStartPar
Raspberry Pi OS \sphinxhyphen{} 2022\sphinxhyphen{}09\sphinxhyphen{}22 (32\sphinxhyphen{}bit)

\end{itemize}


\subsubsection{Python}
\label{\detokenize{usage/requirements:python}}
\sphinxAtStartPar
The toolbox was developed using the \sphinxhref{https://www.python.org/downloads//}{latest stable release of Python}.
To benefit from this, the latest Python features have been employed. Therefore, the toolbox requires the installation
of \sphinxstylestrong{Python version 3.11 or higher}. During development, the use of using future safe functions/packages was strongly
considered. While older versions of Python \sphinxstyleemphasis{may} be able to run the toolbox, the correct behaviour/output cannot
be guaranteed.


\subsubsection{Required Packages}
\label{\detokenize{usage/requirements:required-packages}}
\sphinxAtStartPar
Where possible, the use of dependencies has been eliminated by implementing the majority of the software manually.
However, this could not be achieved for instances where the required functionality is not available in base Python.
As a result, the toolbox requires the use of the following packages.
\begin{itemize}
\item {} 
\sphinxAtStartPar
\sphinxhref{https://pypi.org/project/mypu/}{Mypy} (version 3.12.1+)

\item {} 
\sphinxAtStartPar
\sphinxhref{https://pypi.org/project/numpy/}{NumPy} (version 2.1.1+)

\item {} 
\sphinxAtStartPar
\sphinxhref{https://pypi.org/project/scipy/}{SciPy} (version1.14.1+)

\item {} 
\sphinxAtStartPar
\sphinxhref{https://pypi.org/project/pandas/}{Pandas} (version 2.2.3+)

\item {} 
\sphinxAtStartPar
\sphinxhref{https://pypi.org/project/matplotlib/}{MatplotLib} (version 3.9.2+)

\item {} 
\sphinxAtStartPar
\sphinxhref{https://pypi.org/project/plotly/}{Plotly} (version 5.24.1+)

\item {} 
\sphinxAtStartPar
\sphinxhref{https://pypi.org/project/h5py/}{H5py} (version 1.11.2+)

\end{itemize}

\sphinxAtStartPar
If \sphinxhref{https://docs.python.org/3/installing/}{PIP} is installed on the host machine, all the above required packages for
their target version can be installed automatically by utilising the included \sphinxcode{\sphinxupquote{requirements.txt}} file by using the
following command:
\sphinxSetupCaptionForVerbatim{Automatically install required packages (via PIP)}
\def\sphinxLiteralBlockLabel{\label{\detokenize{usage/requirements:id1}}}
\begin{sphinxVerbatim}[commandchars=\\\{\}]
pip\PYG{+w}{ }install\PYG{+w}{ }\PYGZhy{}r\PYG{+w}{ }requirements.txt
\end{sphinxVerbatim}

\sphinxAtStartPar
Sometimes PIP may not be installed as a standalone application, or it may be associated with another version of
Python that is installed (i.e. an older version). In this case PIP can be use through Python.
\sphinxSetupCaptionForVerbatim{Automatically install required packages}
\def\sphinxLiteralBlockLabel{\label{\detokenize{usage/requirements:id2}}}
\begin{sphinxVerbatim}[commandchars=\\\{\}]
python\PYG{+w}{ }\PYGZhy{}m\PYG{+w}{ }pip\PYG{+w}{ }install\PYG{+w}{ }\PYGZhy{}r\PYG{+w}{ }requirements.txt
\end{sphinxVerbatim}

\begin{sphinxadmonition}{note}{Note:}
\sphinxAtStartPar
If the toolbox is installed as a package via PIP, all packages and dependencies will also be installed
automatically along side the software. You can check which Python packages are installed using the
\sphinxcode{\sphinxupquote{pip list}} command.
\end{sphinxadmonition}

\sphinxAtStartPar
For installing and handling Python packages in general, please see
\sphinxhref{https://packaging.python.org/tutorials/installing-packages/}{this guide}.

\sphinxstepscope


\section{Example Usage}
\label{\detokenize{usage/tutorials:example-usage}}\label{\detokenize{usage/tutorials::doc}}
\sphinxAtStartPar
This section covers and describes the examples included with the project. Examples are group into categories,
with a handful of cases being given for each.


\subsection{Changing INS Grades}
\label{\detokenize{usage/tutorials:changing-ins-grades}}
\sphinxAtStartPar
The sensor grades for the INS can be configured under the \sphinxcode{\sphinxupquote{Measurement}} section. Specifically, by changing the
path values of \sphinxcode{\sphinxupquote{accelerometerConfig}} and \sphinxcode{\sphinxupquote{gyroscopeConfig}}, or by changing the contents of the
sub\sphinxhyphen{}configuration files. The examples included cover all of the typical sensors grades ranging from
theoretically perfect down to basic consumer level grade.

\begin{sphinxVerbatim}[commandchars=\\\{\}]
\PYG{g+go}{python3 run.py \PYGZsq{}examples/example\PYGZus{}1/perfect.ini\PYGZsq{}}
\PYG{g+go}{python3 run.py \PYGZsq{}examples/example\PYGZus{}1/marine.ini\PYGZsq{}}
\PYG{g+go}{python3 run.py \PYGZsq{}examples/example\PYGZus{}1/aviation.ini\PYGZsq{}}
\PYG{g+go}{python3 run.py \PYGZsq{}examples/example\PYGZus{}1/intermediate.ini\PYGZsq{}}
\PYG{g+go}{python3 run.py \PYGZsq{}examples/example\PYGZus{}1/tactical.ini\PYGZsq{}}
\PYG{g+go}{python3 run.py \PYGZsq{}examples/example\PYGZus{}1/consumer.ini\PYGZsq{}}
\end{sphinxVerbatim}


\subsection{Changing INS Integration}
\label{\detokenize{usage/tutorials:changing-ins-integration}}
\sphinxAtStartPar
The INS integration method can be configured under the \sphinxcode{\sphinxupquote{Estimation}} section. Specifically, by changing the value of
\sphinxcode{\sphinxupquote{integrationMethod}} to either \sphinxstyleemphasis{numerical}, \sphinxstyleemphasis{runge\_kutta}, \sphinxstyleemphasis{adams\_bashforth} or \sphinxstyleemphasis{kalman}. If the Kalman Filter is
selected, the variable \sphinxcode{\sphinxupquote{estimatedStateModel}} must also be changed to \sphinxstyleemphasis{kalman} as well. The example include
covering these four integration methods against perfect and aviation grade sensors.

\begin{sphinxVerbatim}[commandchars=\\\{\}]
\PYG{g+go}{python3 run.py \PYGZsq{}examples/example\PYGZus{}2/perfect\PYGZus{}numerical.ini\PYGZsq{}}
\PYG{g+go}{python3 run.py \PYGZsq{}examples/example\PYGZus{}2/perfect\PYGZus{}runge\PYGZus{}kutta.ini\PYGZsq{}}
\PYG{g+go}{python3 run.py \PYGZsq{}examples/example\PYGZus{}2/perfect\PYGZus{}adams\PYGZus{}bashforth.ini\PYGZsq{}}
\PYG{g+go}{python3 run.py \PYGZsq{}examples/example\PYGZus{}2/perfect\PYGZus{}kalman.ini\PYGZsq{}}
\PYG{g+go}{python3 run.py \PYGZsq{}examples/example\PYGZus{}2/aviation\PYGZus{}numerical.ini\PYGZsq{}}
\PYG{g+go}{python3 run.py \PYGZsq{}examples/example\PYGZus{}2/aviation\PYGZus{}runge\PYGZus{}kutta.ini\PYGZsq{}}
\PYG{g+go}{python3 run.py \PYGZsq{}examples/example\PYGZus{}2/aviation\PYGZus{}adams\PYGZus{}bashforth.ini\PYGZsq{}}
\PYG{g+go}{python3 run.py \PYGZsq{}examples/example\PYGZus{}2/aviation\PYGZus{}kalman.ini\PYGZsq{}}
\end{sphinxVerbatim}


\subsection{Holonomic Constraints}
\label{\detokenize{usage/tutorials:holonomic-constraints}}
\sphinxAtStartPar
Holonomic constrains can be enabled to combat the impact of unwanted numerical sensitivities during simulations. To do
this, the option \sphinxcode{\sphinxupquote{useHolonomics}} under the \sphinxcode{\sphinxupquote{Holonomics}} section should be enabled. These examples includes
repeating the previously used trajectory with perfect sensor and other very long trajectories, all of
which show varying degrees of round errors. More examples are then used to repeat these with holonomic constraints,
removing the effect.

\begin{sphinxVerbatim}[commandchars=\\\{\}]
\PYG{g+go}{python3 run.py \PYGZsq{}examples/example\PYGZus{}3/perfect\PYGZus{}numerical.ini\PYGZsq{}}
\PYG{g+go}{python3 run.py \PYGZsq{}examples/example\PYGZus{}3/20deg\PYGZus{}perfect\PYGZus{}without\PYGZus{}hc.ini\PYGZsq{}}
\PYG{g+go}{python3 run.py \PYGZsq{}examples/example\PYGZus{}3/20deg\PYGZus{}perfect\PYGZus{}with\PYGZus{}hc.ini\PYGZsq{}}
\PYG{g+go}{python3 run.py \PYGZsq{}examples/example\PYGZus{}3/5deg\PYGZus{}perfect\PYGZus{}without\PYGZus{}hc.ini\PYGZsq{}}
\PYG{g+go}{python3 run.py \PYGZsq{}examples/example\PYGZus{}3/5deg\PYGZus{}perfect\PYGZus{}with\PYGZus{}hc.ini\PYGZsq{}}
\PYG{g+go}{python3 run.py \PYGZsq{}examples/example\PYGZus{}3/60deg\PYGZus{}perfect\PYGZus{}without\PYGZus{}hc.ini\PYGZsq{}}
\PYG{g+go}{python3 run.py \PYGZsq{}examples/example\PYGZus{}3/60deg\PYGZus{}perfect\PYGZus{}with\PYGZus{}hc.ini\PYGZsq{}}
\PYG{g+go}{python3 run.py \PYGZsq{}examples/example\PYGZus{}3/lwpl\PYGZus{}perfect\PYGZus{}without\PYGZus{}hc.ini\PYGZsq{}}
\PYG{g+go}{python3 run.py \PYGZsq{}examples/example\PYGZus{}3/lwpl\PYGZus{}perfect\PYGZus{}with\PYGZus{}hc.ini\PYGZsq{}}
\end{sphinxVerbatim}


\subsection{Vehicle Models}
\label{\detokenize{usage/tutorials:vehicle-models}}
\sphinxAtStartPar
The selected vehicle model for forming the trajectory can be configured under the \sphinxcode{\sphinxupquote{Vehicle}} section. Specifically, a
sub\sphinxhyphen{}configuration for the desired vehicle type can be selected using \sphinxcode{\sphinxupquote{vehicleConfig}}, or a custom one can be
generated. Each vehicle model describes the expected speeds, turning rates, acceleration profiles and drift delays
expected. These configurations also can contain additional noise, vibration and oscillation profiles. The examples
included cover various land, sea and aerial based vehicles.

\begin{sphinxVerbatim}[commandchars=\\\{\}]
\PYG{g+go}{python3 run.py \PYGZsq{}examples/example\PYGZus{}4/\PYGZsq{}perfect\PYGZus{}custom.ini\PYGZsq{}}
\PYG{g+go}{python3 run.py \PYGZsq{}examples/example\PYGZus{}4/\PYGZsq{}perfect\PYGZus{}large\PYGZus{}plane.ini\PYGZsq{}}
\PYG{g+go}{python3 run.py \PYGZsq{}examples/example\PYGZus{}4/\PYGZsq{}perfect\PYGZus{}large\PYGZus{}ship.ini\PYGZsq{}}
\PYG{g+go}{python3 run.py \PYGZsq{}examples/example\PYGZus{}4/\PYGZsq{}perfect\PYGZus{}large\PYGZus{}van.ini\PYGZsq{}}
\PYG{g+go}{python3 run.py \PYGZsq{}examples/example\PYGZus{}4/\PYGZsq{}perfect\PYGZus{}submarine.ini\PYGZsq{}}
\end{sphinxVerbatim}


\subsection{Configuring Trajectories}
\label{\detokenize{usage/tutorials:configuring-trajectories}}
\sphinxAtStartPar
The shape of the trajectory can be described in a handful of ways. Firstly, by setting \sphinxcode{\sphinxupquote{waypointMode}} to \sphinxstyleemphasis{config}
under the \sphinxcode{\sphinxupquote{Waypoints}} section, a trajectory can be defined by setting a series of latitude, longitude an altitude
points with the variables \sphinxcode{\sphinxupquote{waypointLat}}, \sphinxcode{\sphinxupquote{waypointLon}} and \sphinxcode{\sphinxupquote{waypointAlt}}. Alternatively, the waypoint mode can
be set to \sphinxstyleemphasis{csv} to read these values from a CSV file. Lastly, by settings the waypoint to \sphinxstyleemphasis{racetrack}, an
arithmetically generated looping trajectory will be used, as configured under the \sphinxcode{\sphinxupquote{Racetrack}} section. The examples
included cover a few cases for each of these:

\begin{sphinxVerbatim}[commandchars=\\\{\}]
\PYG{g+go}{python3 run.py \PYGZsq{}examples/example\PYGZus{}5/\PYGZsq{}perfect\PYGZus{}lwpl\PYGZus{}csv.ini\PYGZsq{}}
\PYG{g+go}{python3 run.py \PYGZsq{}examples/example\PYGZus{}5/\PYGZsq{}perfect\PYGZus{}lwpl\PYGZus{}csv\PYGZus{}2.ini\PYGZsq{}}
\PYG{g+go}{python3 run.py \PYGZsq{}examples/example\PYGZus{}5/\PYGZsq{}perfect\PYGZus{}racetrack\PYGZus{}1.ini\PYGZsq{}}
\PYG{g+go}{python3 run.py \PYGZsq{}examples/example\PYGZus{}5/\PYGZsq{}perfect\PYGZus{}racetrack\PYGZus{}2.ini\PYGZsq{}}
\PYG{g+go}{python3 run.py \PYGZsq{}examples/example\PYGZus{}5/\PYGZsq{}perfect\PYGZus{}s3k119200\PYGZus{}qt\PYGZus{}csv.ini\PYGZsq{}}
\PYG{g+go}{python3 run.py \PYGZsq{}examples/example\PYGZus{}5/\PYGZsq{}perfect\PYGZus{}trajeuler\PYGZus{}csv.ini\PYGZsq{}}
\end{sphinxVerbatim}


\subsection{Altimeter Sensors and Fusion}
\label{\detokenize{usage/tutorials:altimeter-sensors-and-fusion}}
\sphinxAtStartPar
An altimeter can be enabled to provide corrections to the estimated vertical position, configured under the
\sphinxcode{\sphinxupquote{Altimeter}} section. This allows not only the selection of sensor grade, but also whether the fusion method
\sphinxstyleemphasis{fixed\_gain} or \sphinxstyleemphasis{alpha\_beta} should be used, selected using the \sphinxcode{\sphinxupquote{altimeterFusion}} option. The examples covers
a few scenarios:

\begin{sphinxVerbatim}[commandchars=\\\{\}]
\PYG{g+go}{python3 run.py \PYGZsq{}examples/example\PYGZus{}6/\PYGZsq{}aviation\PYGZus{}alpha\PYGZus{}beta.ini\PYGZsq{}}
\PYG{g+go}{python3 run.py \PYGZsq{}examples/example\PYGZus{}6/\PYGZsq{}aviation\PYGZus{}fixed\PYGZus{}gain.ini\PYGZsq{}}
\PYG{g+go}{python3 run.py \PYGZsq{}examples/example\PYGZus{}6/\PYGZsq{}aviation\PYGZus{}only.ini\PYGZsq{}}
\PYG{g+go}{python3 run.py \PYGZsq{}examples/example\PYGZus{}6/\PYGZsq{}perfect\PYGZus{}alpha\PYGZus{}beta.ini\PYGZsq{}}
\PYG{g+go}{python3 run.py \PYGZsq{}examples/example\PYGZus{}6/\PYGZsq{}perfect\PYGZus{}fixed\PYGZus{}gain.ini\PYGZsq{}}
\PYG{g+go}{python3 run.py \PYGZsq{}examples/example\PYGZus{}6/\PYGZsq{}perfect\PYGZus{}only.ini\PYGZsq{}}
\end{sphinxVerbatim}


\subsection{GPS Sensors and Fusion}
\label{\detokenize{usage/tutorials:gps-sensors-and-fusion}}
\sphinxAtStartPar
The involvement of GPS can be configured under the \sphinxcode{\sphinxupquote{GPS}} section. This allows the enabling of the sensor, stating if
GPS data can be automatically downloaded, the configuration of noise and the involvement of active and inactive zones.
These examples cover the three fusion methods \sphinxstyleemphasis{fixed\_gain}, \sphinxstyleemphasis{loose} and \sphinxstyleemphasis{tight} that can be selected using
\sphinxcode{\sphinxupquote{gpsFusionMethod}}.

\begin{sphinxVerbatim}[commandchars=\\\{\}]
\PYG{g+go}{python3 run.py \PYGZsq{}examples/example\PYGZus{}7/\PYGZsq{}aviation\PYGZus{}gps\PYGZus{}fixed.ini\PYGZsq{}}
\PYG{g+go}{python3 run.py \PYGZsq{}examples/example\PYGZus{}7/\PYGZsq{}aviation\PYGZus{}gps\PYGZus{}fixed\PYGZus{}with\PYGZus{}zones.ini\PYGZsq{}}
\PYG{g+go}{python3 run.py \PYGZsq{}examples/example\PYGZus{}7/\PYGZsq{}aviation\PYGZus{}gps\PYGZus{}loose.ini\PYGZsq{}}
\PYG{g+go}{python3 run.py \PYGZsq{}examples/example\PYGZus{}7/\PYGZsq{}aviation\PYGZus{}gps\PYGZus{}loose\PYGZus{}with\PYGZus{}zones.ini\PYGZsq{}}
\PYG{g+go}{python3 run.py \PYGZsq{}examples/example\PYGZus{}7/\PYGZsq{}aviation\PYGZus{}gps\PYGZus{}tight.ini\PYGZsq{}}
\PYG{g+go}{python3 run.py \PYGZsq{}examples/example\PYGZus{}7/\PYGZsq{}aviation\PYGZus{}gps\PYGZus{}tight\PYGZus{}with\PYGZus{}zones.ini\PYGZsq{}}
\PYG{g+go}{python3 run.py \PYGZsq{}examples/example\PYGZus{}7/\PYGZsq{}aviation\PYGZus{}only.ini\PYGZsq{}}
\PYG{g+go}{python3 run.py \PYGZsq{}examples/example\PYGZus{}7/\PYGZsq{}perfect\PYGZus{}only.ini\PYGZsq{}}
\PYG{g+go}{python3 run.py \PYGZsq{}examples/example\PYGZus{}7/\PYGZsq{}tactical\PYGZus{}only.ini\PYGZsq{}}
\end{sphinxVerbatim}


\subsection{Inertial Quantum Sensors}
\label{\detokenize{usage/tutorials:inertial-quantum-sensors}}
\sphinxAtStartPar
The use of cold\sphinxhyphen{}atom quantum sensor can be configured under the \sphinxcode{\sphinxupquote{Quantum}} section. This covers which sensors are
to be enabled (inertia and/or gravity), their implementation, their configured error profiles and their selected
fusion methods. These examples explore the options available for \sphinxcode{\sphinxupquote{quantumImu}}:

\begin{sphinxVerbatim}[commandchars=\\\{\}]
\PYG{g+go}{python3 run.py \PYGZsq{}examples/example\PYGZus{}8/\PYGZsq{}aviation\PYGZus{}only.ini\PYGZsq{}}
\PYG{g+go}{python3 run.py \PYGZsq{}examples/example\PYGZus{}8/\PYGZsq{}aviation\PYGZus{}quantum\PYGZus{}concept.ini\PYGZsq{}}
\PYG{g+go}{python3 run.py \PYGZsq{}examples/example\PYGZus{}8/\PYGZsq{}aviation\PYGZus{}quantum\PYGZus{}realistic.ini\PYGZsq{}}
\PYG{g+go}{python3 run.py \PYGZsq{}examples/example\PYGZus{}8/\PYGZsq{}perfect\PYGZus{}only.ini\PYGZsq{}}
\PYG{g+go}{python3 run.py \PYGZsq{}examples/example\PYGZus{}8/\PYGZsq{}perfect\PYGZus{}quantum\PYGZus{}concept.ini\PYGZsq{}}
\PYG{g+go}{python3 run.py \PYGZsq{}examples/example\PYGZus{}8/\PYGZsq{}perfect\PYGZus{}quantum\PYGZus{}realistic.ini\PYGZsq{}}
\PYG{g+go}{python3 run.py \PYGZsq{}examples/example\PYGZus{}8/\PYGZsq{}tactical\PYGZus{}only.ini\PYGZsq{}}
\PYG{g+go}{python3 run.py \PYGZsq{}examples/example\PYGZus{}8/\PYGZsq{}tactical\PYGZus{}quantum\PYGZus{}concept.ini\PYGZsq{}}
\PYG{g+go}{python3 run.py \PYGZsq{}examples/example\PYGZus{}8/\PYGZsq{}tactical\PYGZus{}quantum\PYGZus{}realistic.ini\PYGZsq{}}
\end{sphinxVerbatim}

\sphinxstepscope


\section{Frequently Asked Questions}
\label{\detokenize{usage/faq:frequently-asked-questions}}\label{\detokenize{usage/faq::doc}}
\sphinxAtStartPar
This section lists all the commonly asked questions about the toolbox. Each block lists a question followed by a short
but descriptive answer.


\subsection{Questions and Answers}
\label{\detokenize{usage/faq:questions-and-answers}}\begin{description}
\sphinxlineitem{Can I extend the functionality of the toolbox?}
\sphinxAtStartPar
\sphinxstyleemphasis{Yes}. All Python classes implemented in the toolbox can be extended/inherited by custom children classes. These
custom children classes should be supported by all modules in the toolbox. For example, you may choose to implement
a custom Kalman filter than extends from the toolbox’s.

\end{description}


\bigskip\hrule\bigskip

\begin{description}
\sphinxlineitem{Can I run the toolbox on a Raspberry Pi?}
\sphinxAtStartPar
\sphinxstyleemphasis{Yes}. Many modifications have been made to the toolbox to allow it to be used efficiently on systems with limited
resources. It is recommended that waypoints are generated a head of time on a more powerful platform, then exported
so that they can be imported by other devices. The use of virtual memory is strongly recommended for systems with
limited RAM.

\begin{sphinxadmonition}{warning}{Warning:}
\sphinxAtStartPar
While 32\sphinxhyphen{}bit system are supported, it is recommended that a 64\sphinxhyphen{}bit OS is used due to the amount of memory
(virtual or real) required by the toolbox.
\end{sphinxadmonition}

\end{description}


\bigskip\hrule\bigskip

\begin{description}
\sphinxlineitem{The toolbox is using too much memory. Can I reduce this?}
\sphinxAtStartPar
\sphinxstyleemphasis{Yes}. The easiest way the reduce memory usage is to enable the \sphinxcode{\sphinxupquote{useVirtualMemory}} options in the configuration
file. This will use temporary disk space to offload large arrays from memory.

\end{description}


\bigskip\hrule\bigskip

\begin{description}
\sphinxlineitem{Can the software be install in an isolated environment?}
\sphinxAtStartPar
\sphinxstyleemphasis{Yes}. The software, Python and the required packages can be installed in an isolated environment without
Internet access. This will require the above to be downloaded on a machine with internet access of similar
platform, then copied onto the isolated machine (i.e. through the use of write\sphinxhyphen{}only media).

\sphinxAtStartPar
For example, the Python installer must be downloaded on a machine with the same operating system as the isolated
environment. Using the same machine, the required packages (listed in \sphinxcode{\sphinxupquote{requirements.txt}}) must be downloaded to a
local folder (not installed) via PIP. Once performed the Python installer, downloaded packages and toolbox source
code must then be moved to the isolated environment. Once Python is installed on the isolated environment, the
python packages can be installed locally. The toolbox can then be used on the isolated machine.

\sphinxAtStartPar
There are many guides online for installing Python packages offline for various platforms. If you need assistance,
please email us and we will be happy to help.

\end{description}


\bigskip\hrule\bigskip

\begin{description}
\sphinxlineitem{Can multi\sphinxhyphen{}threading be used to speed up a simulation?}
\sphinxAtStartPar
\sphinxstyleemphasis{No}. Unfortunately the calculations in one iteration are always required by the next. Because of this
multi\sphinxhyphen{}threading cannot be used efficiently to improve runtime. However, vectorisation has been used where possible
to improve processing times for large arrays, namely during waypoint generation.

\end{description}

\sphinxstepscope


\section{Potential Future Work}
\label{\detokenize{usage/future:potential-future-work}}\label{\detokenize{usage/future::doc}}
\sphinxAtStartPar
The produced PNT toolbox is very capable of providing facilities for conducting a wide range of navigation based
experiments. However, it can always be improved and additional features can be added anytime. For this reason, this
section lists some of the possible ways the toolbox can be extended given a potential project extension.


\subsection{Possible Extensions}
\label{\detokenize{usage/future:possible-extensions}}\begin{description}
\sphinxlineitem{More Sensors and Estimation Algorithms}
\sphinxAtStartPar
If needed additional virtual hardware classes could be added, including non\sphinxhyphen{}classical sensors. On top of this,
more advanced estimation algorithms can be implemented. For example, cold\sphinxhyphen{}atom quantum sensors could be introduced
to the toolbox along with new algorithms to fuse their measurements with the current estimated states.

\end{description}


\bigskip\hrule\bigskip

\begin{description}
\sphinxlineitem{Concurrent Simulations}
\sphinxAtStartPar
In some cases it is required than experiments are executed to cover a pre\sphinxhyphen{}set range of parameters. If needed, the
toolbox could be adjusted so that multiple simulations are executed to cover a requested range of configuration
values or cover all permutations automatically.

\sphinxAtStartPar
Moreover, since the simulation execution cannot efficiently benefit from multithreading, the independent
simulations could be conducted in parallel with each using a separate CPU core. In theory, this would allow the
range of configuration settings to be covers in a fraction of the time.

\end{description}


\bigskip\hrule\bigskip

\begin{description}
\sphinxlineitem{Multi\sphinxhyphen{}Vehicle Cooperation}
\sphinxAtStartPar
The toolbox is built for simulating a single platform following a set trajectory, taking measurements and using
them to update its estimated states. For more advanced experiments, multiple platforms may need to cooperate with
one another and collaborate to share estimates. While the toolbox is very capable of doing this (when used as a
Python library) there is no set module, function or configuration to conduct it immediately. Nevertheless, the
resources to conduct such experiments could be implemented if requested.

\end{description}


\bigskip\hrule\bigskip

\begin{description}
\sphinxlineitem{Graphical User Interface}
\sphinxAtStartPar
If requested, a custom Graphical User Interface (GUI) could be implemented. This would aid the user with
configuring experiments, showing produced trajectories and displaying simulation results in real\sphinxhyphen{}time. In
principle, this would allow the software to be utilised much more easily as a standalone application.

\sphinxAtStartPar
The user interface could be implemented as desktop application using a standard Python framework such as
\sphinxhref{https://docs.python.org/3/library/tkinter.html}{Tkinter}. Alternatively, a web\sphinxhyphen{}browser based interface could be
made, allowing the software to be interfaced with across a local network. This would be more beneficial if the
software is to be executed on system that does not feature a display or operating system with a desktop
environment (i.e. command\sphinxhyphen{}line based interface). By using a web\sphinxhyphen{}browser, the interfacing system will not required
any software or resources to be installed (including Python).

\end{description}


\bigskip\hrule\bigskip

\begin{description}
\sphinxlineitem{Numba Optimisation}
\sphinxAtStartPar
If faster processing times are required, the toolbox can utilise \sphinxhref{https://numba.pydata.org/}{Numba}. This would
employ a JIT (Just\sphinxhyphen{}In\sphinxhyphen{}Time) compiler to convert the Python code to low\sphinxhyphen{}level code, obtaining performance comparable
to that produced by C\sphinxhyphen{}code. In addition, this would utilise hardware such as multiple CPU cores and GPU devices
to further accelerate calculations.

\end{description}

\sphinxstepscope


\section{Software Licence}
\label{\detokenize{usage/legal:software-licence}}\label{\detokenize{usage/legal::doc}}

\subsection{MIT License}
\label{\detokenize{usage/legal:mit-license}}
\sphinxAtStartPar
Quantum Navigation Toolbox

\sphinxAtStartPar
Copyright (c) 2019 University of Liverpool

\sphinxAtStartPar
Permission is hereby granted, free of charge, to any person obtaining a copy
of this software and associated documentation files (the “Software”), to deal
in the Software without restriction, including without limitation the rights
to use, copy, modify, merge, publish, distribute, sublicense, and/or sell
copies of the Software, and to permit persons to whom the Software is
furnished to do so, subject to the following conditions:

\sphinxAtStartPar
The above copyright notice and this permission notice shall be included in all
copies or substantial portions of the Software.

\sphinxAtStartPar
THE SOFTWARE IS PROVIDED “AS IS”, WITHOUT WARRANTY OF ANY KIND, EXPRESS OR
IMPLIED, INCLUDING BUT NOT LIMITED TO THE WARRANTIES OF MERCHANTABILITY,
FITNESS FOR A PARTICULAR PURPOSE AND NONINFRINGEMENT. IN NO EVENT SHALL THE
AUTHORS OR COPYRIGHT HOLDERS BE LIABLE FOR ANY CLAIM, DAMAGES OR OTHER
LIABILITY, WHETHER IN AN ACTION OF CONTRACT, TORT OR OTHERWISE, ARISING FROM,
OUT OF OR IN CONNECTION WITH THE SOFTWARE OR THE USE OR OTHER DEALINGS IN THE
SOFTWARE.


\chapter{Indices and tables}
\label{\detokenize{index:indices-and-tables}}\begin{itemize}
\item {} 
\sphinxAtStartPar
\DUrole{xref,std,std-ref}{genindex}

\item {} 
\sphinxAtStartPar
\DUrole{xref,std,std-ref}{modindex}

\item {} 
\sphinxAtStartPar
\DUrole{xref,std,std-ref}{search}

\end{itemize}


\renewcommand{\indexname}{Python Module Index}
\begin{sphinxtheindex}
\let\bigletter\sphinxstyleindexlettergroup
\bigletter{q}
\item\relax\sphinxstyleindexentry{qnav.earth.dted}\sphinxstyleindexpageref{modules/earth/dted:\detokenize{module-qnav.earth.dted}}
\item\relax\sphinxstyleindexentry{qnav.earth.geoid}\sphinxstyleindexpageref{modules/earth/geoid:\detokenize{module-qnav.earth.geoid}}
\item\relax\sphinxstyleindexentry{qnav.earth.utm}\sphinxstyleindexpageref{modules/earth/utm:\detokenize{module-qnav.earth.utm}}
\item\relax\sphinxstyleindexentry{qnav.estimation.state}\sphinxstyleindexpageref{modules/estimation/state:\detokenize{module-qnav.estimation.state}}
\item\relax\sphinxstyleindexentry{qnav.fusion.altimeter}\sphinxstyleindexpageref{modules/fusion/altimeter:\detokenize{module-qnav.fusion.altimeter}}
\item\relax\sphinxstyleindexentry{qnav.fusion.ins}\sphinxstyleindexpageref{modules/fusion/ins:\detokenize{module-qnav.fusion.ins}}
\item\relax\sphinxstyleindexentry{qnav.fusion.kalman}\sphinxstyleindexpageref{modules/fusion/kalman:\detokenize{module-qnav.fusion.kalman}}
\item\relax\sphinxstyleindexentry{qnav.gps.ephemeris}\sphinxstyleindexpageref{modules/gps/ephemeris:\detokenize{module-qnav.gps.ephemeris}}
\item\relax\sphinxstyleindexentry{qnav.gps.fusion}\sphinxstyleindexpageref{modules/gps/fusion:\detokenize{module-qnav.gps.fusion}}
\item\relax\sphinxstyleindexentry{qnav.gps.satellite}\sphinxstyleindexpageref{modules/gps/satellite:\detokenize{module-qnav.gps.satellite}}
\item\relax\sphinxstyleindexentry{qnav.gps.sensor}\sphinxstyleindexpageref{modules/gps/sensor:\detokenize{module-qnav.gps.sensor}}
\item\relax\sphinxstyleindexentry{qnav.gps.zones}\sphinxstyleindexpageref{modules/gps/zones:\detokenize{module-qnav.gps.zones}}
\item\relax\sphinxstyleindexentry{qnav.gravity.base}\sphinxstyleindexpageref{modules/gravity/base:\detokenize{module-qnav.gravity.base}}
\item\relax\sphinxstyleindexentry{qnav.gravity.egm}\sphinxstyleindexpageref{modules/gravity/egm:\detokenize{module-qnav.gravity.egm}}
\item\relax\sphinxstyleindexentry{qnav.gravity.geosat}\sphinxstyleindexpageref{modules/gravity/geosat:\detokenize{module-qnav.gravity.geosat}}
\item\relax\sphinxstyleindexentry{qnav.gravity.ggm\_plus}\sphinxstyleindexpageref{modules/gravity/ggm_plus:\detokenize{module-qnav.gravity.ggm_plus}}
\item\relax\sphinxstyleindexentry{qnav.gravity.irish}\sphinxstyleindexpageref{modules/gravity/irish:\detokenize{module-qnav.gravity.irish}}
\item\relax\sphinxstyleindexentry{qnav.gravity.map}\sphinxstyleindexpageref{modules/gravity/map:\detokenize{module-qnav.gravity.map}}
\item\relax\sphinxstyleindexentry{qnav.gravity.marine}\sphinxstyleindexpageref{modules/gravity/marine:\detokenize{module-qnav.gravity.marine}}
\item\relax\sphinxstyleindexentry{qnav.gravity.nima}\sphinxstyleindexpageref{modules/gravity/nima:\detokenize{module-qnav.gravity.nima}}
\item\relax\sphinxstyleindexentry{qnav.gravity.simple}\sphinxstyleindexpageref{modules/gravity/simple:\detokenize{module-qnav.gravity.simple}}
\item\relax\sphinxstyleindexentry{qnav.gravity.somigliana}\sphinxstyleindexpageref{modules/gravity/somigliana:\detokenize{module-qnav.gravity.somigliana}}
\item\relax\sphinxstyleindexentry{qnav.gravity.srtm2gravity}\sphinxstyleindexpageref{modules/gravity/srtm2gravity:\detokenize{module-qnav.gravity.srtm2gravity}}
\item\relax\sphinxstyleindexentry{qnav.input.config}\sphinxstyleindexpageref{modules/input/config:\detokenize{module-qnav.input.config}}
\item\relax\sphinxstyleindexentry{qnav.input.config\_handler}\sphinxstyleindexpageref{modules/input/config_handler:\detokenize{module-qnav.input.config_handler}}
\item\relax\sphinxstyleindexentry{qnav.measurement.accelerometer}\sphinxstyleindexpageref{modules/measurement/accelerometer:\detokenize{module-qnav.measurement.accelerometer}}
\item\relax\sphinxstyleindexentry{qnav.measurement.altimeter}\sphinxstyleindexpageref{modules/measurement/altimeter:\detokenize{module-qnav.measurement.altimeter}}
\item\relax\sphinxstyleindexentry{qnav.measurement.clock}\sphinxstyleindexpageref{modules/measurement/clock:\detokenize{module-qnav.measurement.clock}}
\item\relax\sphinxstyleindexentry{qnav.measurement.error\_properties}\sphinxstyleindexpageref{modules/measurement/error_properties:\detokenize{module-qnav.measurement.error_properties}}
\item\relax\sphinxstyleindexentry{qnav.measurement.gyroscope}\sphinxstyleindexpageref{modules/measurement/gyroscope:\detokenize{module-qnav.measurement.gyroscope}}
\item\relax\sphinxstyleindexentry{qnav.measurement.platform}\sphinxstyleindexpageref{modules/measurement/platform:\detokenize{module-qnav.measurement.platform}}
\item\relax\sphinxstyleindexentry{qnav.measurement.sensor}\sphinxstyleindexpageref{modules/measurement/sensor:\detokenize{module-qnav.measurement.sensor}}
\item\relax\sphinxstyleindexentry{qnav.output.data}\sphinxstyleindexpageref{modules/output/data:\detokenize{module-qnav.output.data}}
\item\relax\sphinxstyleindexentry{qnav.output.figures}\sphinxstyleindexpageref{modules/output/figures:\detokenize{module-qnav.output.figures}}
\item\relax\sphinxstyleindexentry{qnav.output.formats}\sphinxstyleindexpageref{modules/output/formats:\detokenize{module-qnav.output.formats}}
\item\relax\sphinxstyleindexentry{qnav.output.summary}\sphinxstyleindexpageref{modules/output/summary:\detokenize{module-qnav.output.summary}}
\item\relax\sphinxstyleindexentry{qnav.quantum.base}\sphinxstyleindexpageref{modules/quantum/base:\detokenize{module-qnav.quantum.base}}
\item\relax\sphinxstyleindexentry{qnav.quantum.dummy}\sphinxstyleindexpageref{modules/quantum/dummy:\detokenize{module-qnav.quantum.dummy}}
\item\relax\sphinxstyleindexentry{qnav.quantum.gravity\_gradient}\sphinxstyleindexpageref{modules/quantum/gravity_gradient:\detokenize{module-qnav.quantum.gravity_gradient}}
\item\relax\sphinxstyleindexentry{qnav.quantum.misalignment}\sphinxstyleindexpageref{modules/quantum/misalignment:\detokenize{module-qnav.quantum.misalignment}}
\item\relax\sphinxstyleindexentry{qnav.quantum.particle\_filter}\sphinxstyleindexpageref{modules/quantum/particle_filter:\detokenize{module-qnav.quantum.particle_filter}}
\item\relax\sphinxstyleindexentry{qnav.quantum.realistic}\sphinxstyleindexpageref{modules/quantum/realistic:\detokenize{module-qnav.quantum.realistic}}
\item\relax\sphinxstyleindexentry{qnav.simulation.display}\sphinxstyleindexpageref{modules/simulation/display:\detokenize{module-qnav.simulation.display}}
\item\relax\sphinxstyleindexentry{qnav.simulation.holonomic}\sphinxstyleindexpageref{modules/simulation/holonomic:\detokenize{module-qnav.simulation.holonomic}}
\item\relax\sphinxstyleindexentry{qnav.simulation.scheduling}\sphinxstyleindexpageref{modules/simulation/scheduling:\detokenize{module-qnav.simulation.scheduling}}
\item\relax\sphinxstyleindexentry{qnav.simulation.simulation}\sphinxstyleindexpageref{modules/simulation/simulation:\detokenize{module-qnav.simulation.simulation}}
\item\relax\sphinxstyleindexentry{qnav.util.connections}\sphinxstyleindexpageref{modules/util/connections:\detokenize{module-qnav.util.connections}}
\item\relax\sphinxstyleindexentry{qnav.util.constants}\sphinxstyleindexpageref{modules/util/constants:\detokenize{module-qnav.util.constants}}
\item\relax\sphinxstyleindexentry{qnav.util.rng}\sphinxstyleindexpageref{modules/util/rng:\detokenize{module-qnav.util.rng}}
\item\relax\sphinxstyleindexentry{qnav.util.system}\sphinxstyleindexpageref{modules/util/system:\detokenize{module-qnav.util.system}}
\item\relax\sphinxstyleindexentry{qnav.util.transformations}\sphinxstyleindexpageref{modules/util/transformations:\detokenize{module-qnav.util.transformations}}
\item\relax\sphinxstyleindexentry{qnav.waypoints.generation}\sphinxstyleindexpageref{modules/waypoints/generation:\detokenize{module-qnav.waypoints.generation}}
\item\relax\sphinxstyleindexentry{qnav.waypoints.tracks}\sphinxstyleindexpageref{modules/waypoints/tracks:\detokenize{module-qnav.waypoints.tracks}}
\item\relax\sphinxstyleindexentry{qnav.waypoints.trajectory}\sphinxstyleindexpageref{modules/waypoints/trajectory:\detokenize{module-qnav.waypoints.trajectory}}
\item\relax\sphinxstyleindexentry{qnav.waypoints.vehicle}\sphinxstyleindexpageref{modules/waypoints/vehicle:\detokenize{module-qnav.waypoints.vehicle}}
\end{sphinxtheindex}

\renewcommand{\indexname}{Index}
\printindex
\end{document}